% ============================================================
%  tikz-chemlab.tex
%  化学实验器材 TikZ 绘图包 - 主实现 (参数系统 + 高级命令)
%  原作者 : D. Girou / C. Jorsen / M. Luque / H. Voss
%  TikZ 改写 : Aihua Wang(王爱华) (原 PSTricks 包 pst-labo)
% ============================================================
%
%  依赖 :
%    \usepackage{tikz}
%    \usetikzlibrary{arrows.meta,calc,patterns,patterns.meta,shadings}
%    % ============================================================
%  tikz-chemlabObj.tex
%  化学实验器材 TikZ 绘图包 - 形状宏库
%  原作者 : D. Girou / C. Jorsen / M. Luque / H. Voss
%  TikZ 改写 : Aihua Wang(王爱华) (原 PSTricks 包 pst-laboObj)
% ============================================================
%
%  使用说明 :
%    \usepackage{tikz}
%    \usetikzlibrary{arrows.meta,calc,patterns,patterns.meta,shadings}
%    % ============================================================
%  tikz-chemlabObj.tex
%  化学实验器材 TikZ 绘图包 - 形状宏库
%  原作者 : D. Girou / C. Jorsen / M. Luque / H. Voss
%  TikZ 改写 : Aihua Wang(王爱华) (原 PSTricks 包 pst-laboObj)
% ============================================================
%
%  使用说明 :
%    \usepackage{tikz}
%    \usetikzlibrary{arrows.meta,calc,patterns,patterns.meta,shadings}
%    % ============================================================
%  tikz-chemlabObj.tex
%  化学实验器材 TikZ 绘图包 - 形状宏库
%  原作者 : D. Girou / C. Jorsen / M. Luque / H. Voss
%  TikZ 改写 : Aihua Wang(王爱华) (原 PSTricks 包 pst-laboObj)
% ============================================================
%
%  使用说明 :
%    \usepackage{tikz}
%    \usetikzlibrary{arrows.meta,calc,patterns,patterns.meta,shadings}
%    \input{tikz-chemlabObj}
%    \begin{tikzpicture}
%       \BuretteCorps  ... 等等
%    \end{tikzpicture}
%
%  映射规则 :
%    \psline(A)(B)(C)              -> \draw (A) -- (B) -- (C);
%    \psline{->}(A)(B)             -> \draw[->] (A) -- (B);
%    \pspolygon(A)(B)(C)           -> \filldraw (A) -- (B) -- (C) -- cycle;
%    \psframe(x1,y1)(x2,y2)        -> \draw (x1,y1) rectangle (x2,y2);
%    \psframe*[linecolor=c]        -> \fill[c] (..) rectangle (..);
%    \psframe[framearc=r]          -> \draw[rounded corners=r] (...) rectangle (...);
%    \pscircle(c){r}               -> \draw (c) circle (r);
%    \pscircle*[linecolor=c](c){r} -> \fill[c] (c) circle (r);
%    \psellipse(c)(a,b)            -> \draw (c) ellipse (a and b);
%    \psarc(c){r}{a1}{a2}          -> \draw ([shift={(c)}]a1:r) arc[start angle=a1,end angle=a2,radius=r];
%    \psarcn(c){r}{a1}{a2}         -> \draw ([shift={(c)}]a1:r) arc[start angle=a1,end angle=a2,radius=r];   (注意 : PSTricks 中 arcn 是逆时针 / 角度递减)
%    \psbezier(A)(B)(C)(D)         -> \draw (A) .. controls (B) and (C) .. (D);
%    \pscurve(P1)(P2)...(Pn)       -> \draw plot[smooth,tension=0.5] coordinates {(P1) (P2) ... (Pn)};
%    \pscustom{ ... \fill}         -> \filldraw / \path[fill] (整段路径) -- cycle;
%    \rput(x,y){...}               -> \begin{scope}[shift={(x,y)}] ... \end{scope}
%    \rput{angle}(x,y){...}        -> \begin{scope}[shift={(x,y)},rotate=angle] ... \end{scope}
%    \multido{\i=a+b}{n}{...}      -> \foreach \i in {a,a+b,...,a+(n-1)*b} {...}
%    \psclip{<path>} ... \endpsclip-> \begin{scope} \clip <path>; ... \end{scope}
%    fillstyle=solid,fillcolor=c   -> fill=c
%    fillstyle=gradient,gradbegin=..,gradend=.. -> \shade[left color=..,right color=..]
%    fillstyle=hlines              -> pattern=horizontal lines
%    doubleline=true               -> double
%    linestyle=dashed              -> dashed
%    framearc=r                    -> rounded corners=r (近似)
%    linearc=r                     -> rounded corners=r
% ============================================================

\NeedsTeXFormat{LaTeX2e}
\ProvidesFile{tikz-chemlabObj.tex}[2026/08/03 v1.0.0]

\RequirePackage{tikz}
\usetikzlibrary{arrows.meta,calc,patterns,patterns.meta,shadings,decorations.pathmorphing}

% ------------------------------------------------------------
% 颜色定义 (对应原包中使用的颜色)
% ------------------------------------------------------------
\definecolor{LightGray}{gray}{0.80}      % 浅灰
\definecolor{VeryLightGray}{cmyk}{0,0,0,0.1}  % 极浅灰 (label 标签背景)
\definecolor{LightBlue}{rgb}{0.68,0.85,0.90}
\definecolor{LightBlue1}{rgb}{0.60,0.80,1.00}
\definecolor{VeryPaleOrange}{cmyk}{0,0.1,0.3,0}  % 滴定管液体默认色
\definecolor{PaleOrange}{rgb}{1.00,0.78,0.55}
\definecolor{Brown}{rgb}{0.50,0.25,0.00}
\definecolor{woodColor}{rgb}{0.957,0.84,0.69}
\definecolor{darkWoodColor}{rgb}{0.50,0.25,0.00}
\definecolor{flameAndGrill}{rgb}{0.95,0.55,0.10}  % 估算颜色
\definecolor{Beige}{rgb}{0.96,0.96,0.86}  % Champagne 液体默认色

% ------------------------------------------------------------
% 公共缺省值 (供调用方按需修改)
% ------------------------------------------------------------
\tikzset{
  Distillation/.cd,
  DistillateColor/.store in=\DistillationDistillateColor,
  DistillateColor=LightBlue1,
  MixtureStyle/.style={fill=LightBlue},
}
\newcommand{\tkzDistillationDistillateColor}{LightBlue1}
\newcommand{\tkzDistillationMixtureStyle}{LightBlue}
\newcommand{\tkzTestTubeLiquidLevelI}{11}     % liquidLevel1
\newcommand{\tkzTestTubeLiquidLevelII}{12}    % liquidLevel2
\newif\ifTkzTapOpen
\newif\ifTkzTestTubeStopper

% ============================================================
% 1. 滴定管 (Burette)
% ============================================================
\newcommand{\tkzBuretteBody}{%
  \draw (-0.5,11.5) -- (-0.25,11) -- (-0.25,3) -- (0,2.75) -- (0,0.5);
  \draw (0.2,0.7) -- (0.2,2.75) -- (0.5,3) -- (0.5,11) -- (0.75,11.5);
}

\newcommand{\tkzBuretteStopcock}{%
  % pst-labo: \psframe[fillstyle=solid] → 填充默认白色 + 黑色边框
  \filldraw[fill=white] (-.1,1.6) rectangle (.3,2.1);
  \draw (-.1,1.5) rectangle (-.2,2.2);
  \draw (.3,1.5) rectangle (.4,2.2);
  \draw (.4,1.7) -- (.5,1.7) -- (.6,1.4) -- (.8,1.4)
        -- (.8,2.3) -- (.6,2.3) -- (.5,2) -- (.4,2) -- cycle;
  % 左侧旋钮: pst-labo \pspolygon + \psframe with fillstyle=solid,fillcolor=LightBlue
  % 需要 LightBlue 填充 + 黑色外轮廓线
  \begin{scope}[fill=LightBlue,draw=black]
    \filldraw (-.5,1.5) -- (-.75,1.7) -- (-.75,2) -- (-.5,2.2)
          -- (-.3,2.2) -- (-.3,1.5) -- cycle;
    \filldraw (-.5,1.5) rectangle (-.2,2.2);
  \end{scope}
}

\newcommand{\tkzBuretteGraduation}{%
  % 滴定管刻度 (原 pst-labo 无显式线宽, 这里用细刻度线宽)
  \foreach \n in {3.0,3.32,...,10.68} { \draw[line width=\chemlabGraduationLineWidth] (-0.25,\n) -- (0,\n); }
}

% ============================================================
% 2. pH 计 (pH-mètre)
% ============================================================
\newcommand{\tkzTitrationPHMeter}{%
  % boitier
  \filldraw[fill=white,draw=black,rounded corners=0.125cm] (-1,0) rectangle (1,3);
  \fill[LightGray] (-0.75,2) rectangle (0.75,2.75);
  \node[white] at (0,2.375) {PH};
  \draw (0,1.25) circle (0.3);
  \draw[->] (0,1.25) -- (0,1.5);
  \node at (0.65,1.5) {$^\circ$C};
  \draw (-0.5,0.5) circle (0.2);
  \draw (0.5,0.5) circle (0.2);
  % électrode — 完全透明填充 (不遮挡烧杯刻度线)
  % pst-labo: \psframe[framearc=0.5](2.25,2.5)(2.75,4.5) 外框圆角
  %   顶部圆角被下方 \psframe[LightGray](sharp) 覆盖, 底部圆角被 \pscircle 覆盖
  %   实际可见的外框是直角矩形, 此处用 sharp corners 匹配 pst-labo 可见效果
  \draw (2.25,2.5) rectangle (2.75,4.5);
  % 顶部 LightGray 填充 + 描边 (sharp corners)
  \filldraw[fill=LightGray,draw=black] (2.25,4) rectangle (2.75,4.5);
  \draw (2.5,2.25) circle (0.25);
  % fil
  \draw[line width=0.05cm,LightBlue,rounded corners=0.2cm]
        (2.5,4.5) -- (2.5,5) -- (1.25,4.75) -- (0,5) -- (0,3);
}

% ============================================================
% 3. 滴定支架 (Support dosage)
% ============================================================
\newcommand{\tkzTitrationStand}{%
  % pst-labo: \psframe(-1.5,0)(1.5,1.5) + {\psset{doubleline=true} ...}
  % 线宽: pst-labo 在 \tkzTitration 中设 linewidth=0.053 (×0.4cm = 0.0212cm ≈ 0.6pt)
  \chemlabUpdateLineWidths
  % 平台: 无填充, 只有边框
  \draw[line width=\chemlabLineWidth] (-1.5,0) rectangle (1.5,1.5);
  % 红色按钮: \pscircle*[linecolor=red] → 红色填充, 双线边框也是红色 (不可见)
  \fill[red] (0,0.75) circle (0.2);
  % 回字型矩形按钮: 无填充, 双线边框
  % PSTricks doubleline: 指定坐标在内外线之间
  % 实测: 外线向外偏移 0.085 units, 内线向内偏移 0.034 units
  \draw[line width=\chemlabLineWidth] ({-1.25-0.085},{0.3-0.085}) rectangle ({-0.75+0.085},{1.2+0.085});
  \draw[line width=\chemlabLineWidth] ({-1.25+0.034},{0.3+0.034}) rectangle ({-0.75-0.034},{1.2-0.034});
  % 环形旋钮: 无填充, 双线边框
  % 外线向外偏移, 内线向内偏移
  \draw[line width=\chemlabLineWidth] (1,0.75) circle ({0.25+0.085});
  \draw[line width=\chemlabLineWidth] (1,0.75) circle ({0.25-0.034});
}

% ============================================================
% 4. 磁力搅拌子 (Aimant)
% ============================================================
\newcommand{\tkzTitrationStirBar}{%
  % pst-labo: \psset{linewidth=0.053} 在外层 \tkzTitration 中设置, 随 scale 缩放
  \chemlabUpdateLineWidths
  \draw[line width=\chemlabLineWidth,rounded corners=0.075cm,fill=white] (-0.5,0.1) rectangle (0.5,0.4);
  \draw[line width=\chemlabLineWidth,rounded corners=0.025cm,fill=white] (-0.05,0) rectangle (0.05,0.5);
}

% ============================================================
% 5. 移液管 (Pipette)
% ============================================================
\newcommand{\tkzPipetteBody}{%
  \chemlabUpdateLineWidths\begin{scope}[line width=\chemlabLineWidth,rounded corners=0.05cm]
    \draw (-0.05,0) -- (-0.15,0.2) -- (-0.15,6)
          -- (0.15,6) -- (0.15,0.2) -- (0.05,0) -- cycle;
    \fill[white] (0,6) ellipse (0.15 and 0.05);
    \foreach \Level in {0.5,1.0,...,5.0} {%
      \draw[line width=\chemlabGraduationLineWidth] (0,\Level) -- (0.1,\Level);}
  \end{scope}
}

\newcommand{\tkzLargePipetteBody}{%
  \chemlabUpdateLineWidths\begin{scope}[line width=\chemlabLineWidth,rounded corners=0.05cm]
    \draw (0,0) -- (0,1) -- (-0.23,2.5) -- (-0.23,22.5)
          arc[start angle=-80,end angle=80,radius=0.45];
  \end{scope}
}

% ============================================================
% 6. 液滴 (Goutte)
% ============================================================
% 原文件中 \pst@goutte 被重复定义了两次, 此处使用后者 (unit=0.6 的版本)
% pst-labo: \psset{unit=0.6}\psbezier[linewidth=0.017] — linewidth=0.017 是
%   0.017 用户单位 = 0.017*0.6 = 0.0102cm (随 goutte 的 LOCAL unit=0.6 缩放, 
%   不随外层 unit 变化, 因为 \psset{unit=0.6} 覆盖了外层设置)
% TikZ: line width=0.0102cm (ABSOLUTE, 匹配 pst-labo 的固定 0.0102cm)
\newcommand{\tkzDrop}{%
  \begin{scope}[scale=0.6]
    \draw[line width=0.0102cm]
      (0,0) .. controls (0.25,-.4) and (-.25,-.4) .. (0,0);
  \end{scope}
}

% ============================================================
% 7. 量筒 (Eprouvette)
% ============================================================
\newcommand{\tkzGraduatedCylinderBody}{%
  \draw (-0.5,5) -- (-0.5,0) -- (0.5,0) -- (0.5,4.9) -- (0.6,4.93);
}

\newcommand{\tkzGraduatedCylinderRim}{%
  \draw plot[smooth,tension=0.5] coordinates {%
    (0.6,-0.07) (0.498,0.008) (0.483,0.026) (0.433,0.05)
    (0.353,0.0707) (0.25,0.0866) (0.13,0.0966) (0,0.1)
    (-0.13,0.0966) (-0.25,0.0866) (-0.353,0.0707) (-0.433,0.05)
    (-0.483,0.026) (-0.5,0) (-0.483,-0.026) (-0.353,-0.0707)
    (-0.25,-0.0866) (-0.13,-0.096) (0,-0.1) (0.13,-0.096)
    (0.25,-0.0866) (0.353,-0.0707) (0.433,-0.05)
    (0.498,-0.026) (0.6,-0.07)};
}

\newcommand{\tkzGraduatedCylinderBase}{%
  \fill[white]
    ([shift={(-0.95,0.5)}]0:0.45) arc[start angle=0,end angle=-90,radius=0.45]
    -- ([shift={(-0.95,0.025)}]90:0.025) arc[start angle=90,end angle=270,radius=0.025]
    -- (-0.95,0) -- (0.95,0)
    -- ([shift={(0.95,0.025)}]-90:0.025) arc[start angle=-90,end angle=90,radius=0.025]
    -- ([shift={(0.95,0.5)}]-90:0.45) arc[start angle=-90,end angle=-180,radius=0.45]
    -- plot[smooth,tension=0.5] coordinates {(0.5,0.4) (0,0.25) (-0.5,0.4)}
    -- cycle;
  \draw
    ([shift={(-0.95,0.5)}]0:0.45) arc[start angle=0,end angle=-90,radius=0.45]
    -- ([shift={(-0.95,0.025)}]90:0.025) arc[start angle=90,end angle=270,radius=0.025]
    -- (-0.95,0) -- (0.95,0)
    -- ([shift={(0.95,0.025)}]-90:0.025) arc[start angle=-90,end angle=90,radius=0.025]
    -- ([shift={(0.95,0.5)}]-90:0.45) arc[start angle=-90,end angle=-180,radius=0.45]
    -- plot[smooth,tension=0.5] coordinates {(0.5,0.4) (0,0.25) (-0.5,0.4)};
}

\newcommand{\tkzGraduatedCylinderGraduation}{%
  % 量筒刻度 : 主刻度 1.5× 细刻度 (与烧杯一致)
  \pgfmathsetlengthmacro{\chemlab@eprouvette@major}{1.5*\chemlabGraduationLineWidth}%
  \foreach \nyD in {0.5,1.0,...,4.5} { \draw[line width=\chemlab@eprouvette@major] (-0.4,\nyD) -- (-0.1,\nyD); }
  \foreach \nyd in {0.5,0.6,...,4.4} { \draw[line width=\chemlabGraduationLineWidth] (-0.4,\nyd) -- (-0.25,\nyd); }
}

% ============================================================
% 8. 分馏装置 (Distillation fractionnée)
% ============================================================
% \pst@goutte 这里又被定义了一次, 在 TikZ 版中沿用同一个 \tkzDrop

\newcommand{\tkzColumn}{%
  % pst-labo: linearc=0.3 (×unit, 不是绝对值) → TikZ: rounded corners=0.12cm (= 0.3 × 0.4cm)
  % 尖刺 (pointeG/D) 先画并用 clip 裁剪到蒸馏柱内部, 再画轮廓线覆盖切边。
  %   pst-labo 中 linearc=0.3 (R=0.3 user units), 尖刺端点恰在圆弧切点附近, 几乎
  %   不可见。TikZ 的 rounded corners 实际 R 比标称值大 ~2×, 圆弧切点更低,
  %   尖刺端点 (在 y_col=8 和 7.9) 会突出到圆弧外侧形成多余线头。用 clip 将
  %   尖刺裁剪到柱内, 再画轮廓线覆盖切边, 消除突出线头。
  \begin{scope}
    \clip[rounded corners=0.12cm]
      (-.25,10.5) -- (-.25,8.5) -- (-.75,8) -- (-.75,2.5) -- (-.25,2) -- (-.25,0)
      -- (0.25,0) -- (0.25,2) -- (.75,2.5) -- (.75,8) -- (0.25,8.5) -- (0.25,9.36)
      -- (0.25,10.5) -- cycle;
    % 左侧尖刺 (pointeG)
    \foreach \i in {3,4,...,8} {%
      \draw[line width=0.4\pgflinewidth] (-.75,\i) -- (-.25,\i-0.5)
          -- (-.75,\i-0.1);}
    % 右侧尖刺 (pointeD)
    \foreach \i in {3,4,...,8} {%
      \draw[line width=0.4\pgflinewidth] (.75,\i) -- (.25,\i-0.5)
          -- (.75,\i-0.1);}
  \end{scope}
  \draw[rounded corners=0.12cm]
    (-.25,10.5) -- (-.25,8.5) -- (-.75,8) -- (-.75,2.5)
    -- (-.25,2) -- (-.25,0);
  \draw[rounded corners=0.12cm]
    (0.25,0) -- (0.25,2) -- (.75,2.5) -- (.75,8)
    -- (0.25,8.5) -- (0.25,9.36);
  % 上部断口 (line 3a/3b):
  %   pst-labo 源码 line 3a: \psline(0.25,8.5)(0.25,9.30) 用于填充 linearc 圆角
  %   留下的间隙。实测发现 TikZ 的 rounded corners 实际圆弧半径比标称值大
  %   ~2× (在 scale=0.6 下 rounded corners=0.3cm 给出 R≈0.62cm), 圆弧已覆盖到
  %   x=0.25 处 y_col=8.93 附近, 故 line 3a 在任何 scale/line width 下均被圆弧
  %   覆盖, 不可见。为避免在编号图等彩色调试场景中残留多余线段, 此处不再
  %   单独绘制 line 3a。
  \draw (0.25,9.5) -- (0.25,10.5);
  % 填充 / 阴影
  \begin{scope}[pattern={Lines[angle=45,distance={1pt/sqrt(2)}]},pattern color=gray]
    \fill (-.25,10) rectangle (-.05,10.5);
    \fill (.25,10)  rectangle (0.05,10.5);
  \end{scope}
  \foreach \n in {2.5,3.5,...,7.5} {%
    \begin{scope}[shift={(-.25,\n)}]\tkzDrop\end{scope}
    \begin{scope}[shift={(.25,\n)}]\tkzDrop\end{scope}}
  \begin{scope}[shift={(-0.25,-.1)}]\tkzDrop\end{scope}
  \begin{scope}[shift={(0.25,-.1)}]\tkzDrop\end{scope}
}

% 公共别名: pst-labo 中 \tkzColumn 是用户可调用的公共宏 (小写)
% tikz-chemlab 内部用 \tkzColumn (驼峰), 这里注册公共别名保持兼容
\let\tkzColumn\tkzColumn

% réfrigérant (condenseur)
\newcommand{\tkzCondenser}{%
  \draw ([shift={(3,0)}]90:0.5) arc[start angle=90,end angle=270,radius=0.5];
  % pst-labo: \psarc(7,0){0.5}{270}{90} → 逆时针经过 0°/右侧
  \draw ([shift={(7,0)}]270:0.5) arc[start angle=270,end angle=450,radius=0.5];
  \fill[white] (0,-.05) rectangle (9,0.05);
  \begin{scope}[line width=0.4pt]
    \draw (-.075,.05) -- (9,0.05);
    \draw (0.02,-.05) -- (9,-.05);
  \end{scope}
  \draw (3,0.5) -- (3.5,0.5);
  \draw (3,-.5) -- (6.5,-.5);
  \draw (3.6,0.5) -- (7,0.5);
  \begin{scope}[line width=0.4pt]
    \draw (3.5,0.5) -- (3.5,1.5);
    \draw (3.6,0.5) -- (3.6,1.5);
    \draw (6.6,-.5) -- (6.6,-1.5);
    \draw (6.5,-.5) -- (6.5,-1.5);
  \end{scope}
  \draw (6.6,-.5) -- (7,-.5);
  \draw[->] (6.55,-2.5) -- (6.55,-1.5);
  \draw[->] (3.55,1.5)  -- (3.55,2.5);
}

\newcommand{\tkzThermometer}{%
  % pst-labo: \psline[linearc=0.05](-.05,.25)(-.05,2.75)(0.05,2.75)(0.05,0.25)
  % linearc=0.05 × unit, TikZ: rounded corners=0.05 (受 scale 缩放)
  \draw[rounded corners=0.025cm]
    (-.05,.25) -- (-.05,2.75) -- (0.05,2.75) -- (0.05,0.25) -- cycle;
  % pst-labo: \psframe[framearc=0.8,fillstyle=solid,fillcolor=red](-.15,-.25)(0.15,0.25)
  % framearc=0.8: 80% of half-min(0.3,0.5)=0.15 → radius=0.12
  \filldraw[fill=red,draw=black,rounded corners=0.06cm] (-.15,-.25) rectangle (0.15,0.25);
  % pst-labo: \psline[linecolor=red](0,2) — 红色竖线, 使用默认线宽
  % 在 pst-labo 渲染中较粗, 看起来像红色填充
  \draw[red,line width=0.8pt] (0,0) -- (0,2);
}

\newcommand{\tkzHeater}{%
  \draw (1.17,2.5) -- (0,2.5) -- (0,0) -- (5.5,0) -- (5.5,2.5) -- (4.33,2.5);
}

% 内部小量筒轮廓 - 路径版本 (供 clip 使用)
\newcommand{\tkzGraduatedCylinderPath}{%
  (-0.5,6) -- (-0.5,.5)
  -- ([shift={(-1,0.5)}]0:0.5) arc[start angle=0,end angle=-90,radius=0.5]
  -- (-1,0) -- (1,0)
  -- ([shift={(1,0.5)}]270:0.5) arc[start angle=270,end angle=180,radius=0.5]
  -- (0.5,0.5) -- (0.5,6) -- cycle
}

% 内部小量筒 (描边)
% pst-labo: \psset{linewidth=1.5\pslinewidth} — \pslinewidth 默认 0.8pt (ABSOLUTE)
% → 1.5*0.8pt = 1.2pt (绝对值, 不随 unit/scale 缩放)
% 注意: TikZ 的 [line width=N\pgflinewidth] 有循环引用 bug — \pgflinewidth 在
% 乘法前被更新, 导致 line width=3\pgflinewidth 实际给出 3.6pt 而非 1.2pt。
% 必须用显式 pt 值。
\newcommand{\tkzGraduatedCylinderSmall}{%
  \begin{scope}[line width=0.72pt,line cap=round,line join=round]
    \draw (-0.5,6) -- (-0.5,.5)
      -- ([shift={(-1,0.5)}]0:0.5) arc[start angle=0,end angle=-90,radius=0.5]
      -- (-1,0) -- (1,0)
      -- ([shift={(1,0.5)}]270:0.5) arc[start angle=270,end angle=180,radius=0.5]
      -- (0.5,0.5) -- (0.5,6);
  \end{scope}
}

% 带液体的量筒 (内部使用, 重命名以避免与主文件 \pstEproulette 冲突)
\newcommand{\tkzGraduatedCylinderWithLiquid}{%
  \begin{scope}
    \clip \tkzGraduatedCylinderPath;
    % pst-labo: \psframe*[linecolor=\psk@Distillation@DistillateColor]
    \fill[\tkzDistillateColor] (-1.5,-1) rectangle (1.5,1);
  \end{scope}
  \tkzGraduatedCylinderSmall
}

% (注 : 原 PSTricks 中 \pst@roundFlask 已合并到 \tkzRoundFlaskSimple, 此处不再单独定义,
%  避免与 tikz-chemlab.tex 中的用户命令 \tkzRoundFlask 冲突)

% 自定义密集斜线图案 (匹配 pst-labo hlines 的密度)
{\pgfpointorigin}{\pgfpoint{0.012cm}{0.012cm}}{\pgfpoint{0.012cm}{0.012cm}}
{\pgfsetlinewidth{0.0028cm}\pgfpathmoveto{\pgfpoint{0cm}{0cm}}\pgfpathlineto{\pgfpoint{0.012cm}{0.012cm}}\pgfusepath{stroke}}
{\pgfpointorigin}{\pgfpoint{0.012cm}{0.012cm}}{\pgfpoint{0.012cm}{0.012cm}}
{\pgfsetlinewidth{0.004cm}\pgfpathmoveto{\pgfpoint{0cm}{0cm}}\pgfpathlineto{\pgfpoint{0.012cm}{0.012cm}}\pgfusepath{stroke}}
% \pst@DistillationFractionnee 主图
\newcommand{\tkzDistillationFractionated}{%
  % --- roundFlask ---
  \def\flaskLocal{%
    % pst-labo: fillstyle=hlines,hatchwidth=0.007,hatchsep=0.03 → 水平线填充
    \begin{scope}[pattern={Lines[angle=45,distance={1pt/sqrt(2)}]},pattern color=gray]
      \fill (-.5,2.25) rectangle (-.25,3);
      \fill (0.5,2.25)  rectangle (0.25,3);
    \end{scope}
    % pst-labo 线宽: neck=0.053, arcs=0.03 (都是 ×unit 用户单位, 随 scale 缩放)
    % \chemlabUpdateLineWidths 已设置 \chemlabLineWidth = 0.053*scale
    \chemlabUpdateLineWidths
    \draw[line width=\chemlabLineWidth] (0.5,1.5) -- (0.5,3);
    \draw[line width=\chemlabLineWidth] (-.5,1.5) -- (-.5,3);
    % pst-labo: \psarc(0,0){1.581}{108.435}{71.565} 逆时针走长路 (经底部),
    %   左右侧弧可见, 底部被液体填充覆盖, 顶部无可见弧线。此处拆分为
    %   左右侧两段 (底部由下方轮廓弧绘制, 顶部不画)。
    % 右上侧弧: 71° → 0° (连接右颈底部到球体最右点)
    \draw[line width=\chemlabLineWidth] ([shift={(0,0)}]71.565:1.581) arc[start angle=71.565,end angle=0,radius=1.581];
    % 左上侧弧: 180° → 108° (连接球体最左点到左颈底部)
    \draw[line width=\chemlabLineWidth] ([shift={(0,0)}]180:1.581) arc[start angle=180,end angle=108.435,radius=1.581];
    % pst-labo: \psset{style=\psk@Distillation@MixtureStyle}\fill
    % 需要用 \shade (不是 \fill) 来产生渐变效果
    \shade[\tkzMixtureStyle] plot[smooth,tension=0.5] coordinates {%
      (-1.581,0) (-1.25,0.1) (-.75,-.1) (0,.1) (0.75,-.1) (1.25,.1) (1.581,0)}
      arc[start angle=0,end angle=-180,radius=1.581];
    % 底部轮廓弧: pst-labo \psarc[linewidth=0.03] → 0.03*unit, 用 \chemlabLineWidth 0.053
    % (pst-labo 中 0.03 vs 0.053 差异微小, 统一用 0.053*scale 避免额外变量)
    \draw[line width=\chemlabLineWidth] (0,0) ++(180:1.581) arc[start angle=180,end angle=360,radius=1.581];
  }%
  % --- 主图整体平移 ---
  \begin{scope}[shift={(0,-3)}]
    \begin{scope}[shift={(2.75,13.5)}]\tkzThermometer\end{scope}
    \begin{scope}[shift={(2.75,2.5)}]\flaskLocal\end{scope}
    \begin{scope}[shift={(2.75,4)}]\tkzColumn\end{scope}
    \begin{scope}[shift={(3,13.4)},rotate=-45]\tkzCondenser\end{scope}
    \begin{scope}[shift={(0,0)}]\tkzHeater\end{scope}
    \begin{scope}[shift={(9.4,0)}]\tkzGraduatedCylinderWithLiquid
      % 蒸馏装置量筒刻度 : 主刻度 1.5× 细刻度
      \pgfmathsetlengthmacro{\chemlab@distill@major}{1.5*\chemlabGraduationLineWidth}%
      \foreach \Level in {1,2,...,5} {\draw[line width=\chemlab@distill@major] (-.4,\Level) -- (-.1,\Level);}
      \foreach \Level in {0.2,0.4,...,4.8} {\draw[line width=\chemlabGraduationLineWidth] (-.4,\Level) -- (-.2,\Level);}
    \end{scope}
    \begin{scope}[shift={(9.37,7)}]\tkzDrop\end{scope}
    \begin{scope}[shift={(9.37,6)}]\tkzDrop\end{scope}
    \begin{scope}[shift={(9.37,4)}]\tkzDrop\end{scope}
    \begin{scope}[shift={(9.37,1.2)}]\tkzDrop\end{scope}
    % 加热盘 (搅拌器)
    \begin{scope}[shift={(0.75,.75)}]
      \draw (0,0) circle (.5);
      \foreach \i in {0,30,...,180} {\draw (\i:0.55) -- (\i:0.65);}
      \draw[->,line width=0.035cm] (0,0) -- (0,0.45);
    \end{scope}
  \end{scope}
}

% ============================================================
% 9. 本生灯 (Bec Bunsen) - 不带 grille
% ============================================================
\newcommand{\tkzBunsenBurner}{%
  \chemlabUpdateLineWidths\begin{scope}[line width=\chemlabLineWidth,line cap=round,line join=round]
    \draw (-1.25,0) rectangle (1.25,0.25);
    \draw (-.5,1.25) rectangle (0.5,2.25);
    \foreach \n in {-0.3,0.0,0.3} {\draw (\n,1.75) circle (0.1);}
    \draw (-.25,2.25) rectangle (0.25,4.25);
    \draw (0.25,1.25) -- (0.25,0.5) -- (1.25,0.25);
    \draw (-1.25,0.25) -- (-.25,0.5) -- (-0.25,0.75);
    \draw (-2.25,0.75) -- (-.25,0.75);
    \draw (-2.25,1) -- (-.25,1);
    \draw (-.25,0.875) ellipse (0.1 and 0.125);
    \fill[white] (-2.25,0.75) rectangle (-0.25,1);
    \draw (-2.25,0.75) -- (-0.25,0.75);
    \draw (-2.25,1) -- (-0.25,1) -- (-.25,1.25);
    \draw plot[smooth,tension=0.5] coordinates {(-0.25,0.5) (0,0.4) (0.25,0.5)};
    % flamme
    % pst-labo 源码: \rput(0,4.25){\psclip{\psbezier[...gradient...] (12点)}\pspolygon[...gradient...](-0.25,0)(0.25,0)(0,1)\endpsclip}
    % \psbezier 12个点 + BezierNArray padding (0,0) = 13点 = 4段贝塞尔
    % 经像素级验证, 正确路径从 (0,0) 出发, 经4段贝塞尔回到 (0,0), 无闭合直线:
    %   Seg 1: (0,0) → c1=P1=(-0.25,0), c2=P2=(-0.35,0.5) → P3=(-0.4,0.75) [左下→左中]
    %   Seg 2: P3=(-0.4,0.75) → c1=P4=(-0.35,1), c2=P5=(-0.25,1.5) → P6=(0.5,2) [左中→顶]
    %   Seg 3: P6=(0.5,2) → c1=P7=(0.25,1.5), c2=P8=(0.35,1) → P9=(0.4,0.75) [顶→右中]
    %   Seg 4: P9=(0.4,0.75) → c1=P10=(0.35,0.5), c2=P11=(0.25,0) → (0,0) [右中→底]
    % 渐变: gradbegin=PaleOrange (底), gradend=yellow (顶), gradmidpoint=0, gradangle=0 (垂直)
    \begin{scope}[shift={(0,4.25)}]
      % 外焰填充: PaleOrange (底) -> yellow (顶) 线性垂直渐变
      \shade[top color=yellow, bottom color=PaleOrange]
        (0,0) .. controls (-0.25,0) and (-0.35,0.5)
        .. (-0.4,0.75) .. controls (-0.35,1) and (-0.25,1.5)
        .. (0.5,2) .. controls (0.25,1.5) and (0.35,1)
        .. (0.4,0.75) .. controls (0.35,0.5) and (0.25,0)
        .. (0,0);
      % 内焰: 三角形 cyan(底)->white(顶) 渐变, 裁剪到外焰形状内
      \begin{scope}
        \clip
          (0,0) .. controls (-0.25,0) and (-0.35,0.5)
          .. (-0.4,0.75) .. controls (-0.35,1) and (-0.25,1.5)
          .. (0.5,2) .. controls (0.25,1.5) and (0.35,1)
          .. (0.4,0.75) .. controls (0.35,0.5) and (0.25,0)
          .. (0,0);
        \shade[top color=white, bottom color=cyan]
          (-0.25,0) -- (0.25,0) -- (0,1) -- cycle;
      \end{scope}
    \end{scope}
  \end{scope}
}

% 火焰 + 网格 (FlammeGrille 形状) - 路径版本 (供 clip 使用)
\newcommand{\tkzFlameGrillPath}{%
  ([shift={(0,0.75)}]270:0.75) arc[start angle=270,end angle=360,radius=0.75]
  -- ([shift={(1.5,0.75)}]180:0.75) arc[start angle=180,end angle=90,radius=0.75]
  -- (1.5,1.5) -- (-1.5,1.5)
  -- ([shift={(-1.5,0.75)}]90:0.75) arc[start angle=90,end angle=0,radius=0.75]
  -- ([shift={(0,0.75)}]180:0.75) arc[start angle=180,end angle=270,radius=0.75]
  -- cycle
}

\newcommand{\tkzFlameGrill}{%
  \fill[flameAndGrill] \tkzFlameGrillPath;
}

% ============================================================
% 10. 本生灯 - 带 grille
% ============================================================
\newcommand{\tkzBunsenBurnerGrill}{%
  \chemlabUpdateLineWidths\begin{scope}[line width=\chemlabLineWidth,line cap=round,line join=round]
    \draw (-1.25,0) rectangle (1.25,0.25);
    \draw (-.5,1.25) rectangle (0.5,2.25);
    \foreach \n in {-0.3,0.0,0.3} {\draw (\n,1.75) circle (0.1);}
    \draw (-.25,2.25) rectangle (0.25,4.25);
    \draw (0.25,1.25) -- (0.25,0.5) -- (1.25,0.25);
    \draw (-1.25,0.25) -- (-.25,0.5) -- (-0.25,0.75);
    \draw (-2.25,0.75) -- (-.25,0.75);
    \draw (-2.25,1) -- (-.25,1);
    \draw (-.25,0.875) ellipse (0.1 and 0.125);
    \fill[white] (-2.25,0.75) rectangle (-0.25,1);
    \draw (-2.25,0.75) -- (-0.25,0.75);
    \draw (-2.25,1) -- (-0.25,1) -- (-.25,1.25);
    \draw plot[smooth,tension=0.5] coordinates {(-0.25,0.5) (0,0.4) (0.25,0.5)};
    % flamme (clip 到 FlammeGrille)
    % 原 pst-labo 火焰分两层:
    %   外焰: \pst@FlammeGrille 整体填充 flameAndGrill (PaleOrange->yellow 渐变)
    %   内焰: 三角形 cyan->white 渐变 (在内焰区域)
    % 颜色和大小可通过 \chemlabset 调整
    \begin{scope}[shift={(0,4.25)},scale=\chemlabFlameScale]
      \begin{scope}
        \clip \tkzFlameGrillPath;
        % 外焰: PaleOrange (底) -> yellow (顶) 线性垂直渐变
        % 对应 pst-labo: fillstyle=gradient, gradbegin=PaleOrange, gradend=yellow
        % 渐变范围与火焰形状 (y=0..1.5) 一致, 避免底部颜色被稀释
        \shade[top color=\chemlabFlameOuterEndColor, bottom color=\chemlabFlameOuterColor]
          (-1.5,0) rectangle (1.5,1.5);
        % 内焰: cyan (底) -> white (顶) 线性垂直渐变 (三角形)
        \shade[top color=\chemlabFlameInnerEndColor, bottom color=\chemlabFlameInnerColor]
          (-0.25,0) -- (0.25,0) -- (0,1) -- cycle;
      \end{scope}
    \end{scope}
    \draw[dashed,line width=\chemlabEllipseLineWidth] (-2,5.75) -- (2,5.75);
  \end{scope}
}

% ============================================================
% 11. 烧瓶 (Ballon)
% ============================================================
% 路径版本 - 供 clip 使用 (闭合轮廓)
\newcommand{\tkzRoundFlaskSimplePath}{%
  ([shift={(-0.7,3.4)}]90:0.2) arc[start angle=90,end angle=0,radius=0.2]
  -- (-0.5,3.4) -- (-0.5,2.4)
  -- (-0.489,2.401)
  arc[start angle=113,end angle=427,radius=1.25]
  -- (0.5,2.4) -- (0.5,3.4)
  -- ([shift={(0.7,3.4)}]180:0.2) arc[start angle=180,end angle=90,radius=0.2]
  -- (0,3.5) -- cycle
}

\newcommand{\tkzRoundFlaskSimple}{%
  \chemlabUpdateLineWidths\begin{scope}[line width=\chemlabLineWidth,line cap=round,line join=round]
    % 左侧: 瓶口圆弧 + 颈部直线 + 球体大圆弧, 合并为单条路径以平滑过渡
    \draw ([shift={(-0.7,3.4)}]90:0.2) arc[start angle=90,end angle=0,radius=0.2]
          -- (-0.5,3.4) -- (-0.5,2.4)
          -- (-0.489,2.401) arc[start angle=113,end angle=427,radius=1.25];
    % 右侧: 球体终点(0.489,2.401) -> 颈部直线 -> 瓶口圆弧
    \draw (0.489,2.401) -- (0.5,2.4) -- (0.5,3.4)
          -- ([shift={(0.7,3.4)}]180:0.2) arc[start angle=180,end angle=90,radius=0.2];
    % 瓶口椭圆: 白色填充覆盖颈部上端圆头, 避免穿入椭圆
    % pst-labo: ballon椭圆用默认线宽 0.053 (与主轮廓一致, 不是 0.07)
    \filldraw[fill=white,draw=black,line width=\chemlabLineWidth] (0,3.6) ellipse (0.7 and 0.1);
  \end{scope}
}

\newcommand{\tkzRoundFlaskGrill}{%
  \chemlabUpdateLineWidths\begin{scope}[line width=\chemlabLineWidth,line cap=round,line join=round]
    % 左侧: 瓶口圆弧 + 颈部直线 + 球体大圆弧, 合并为单条路径
    \draw ([shift={(-0.7,2.15)}]90:0.2) arc[start angle=90,end angle=0,radius=0.2]
          -- (-0.5,2.15) -- (-0.5,1.4)
          -- (-0.489,1.401) arc[start angle=113,end angle=427,radius=1.25];
    % 右侧: 球体终点 -> 颈部直线 -> 瓶口圆弧
    \draw (0.489,1.401) -- (0.5,1.15) -- (0.5,3.4)
          -- ([shift={(0.7,2.15)}]180:0.2) arc[start angle=180,end angle=90,radius=0.2];
    % 瓶口椭圆: 白色填充覆盖颈部上端圆头, 避免穿入椭圆
    % pst-labo: ballon椭圆用默认线宽 0.053 (与主轮廓一致)
    \filldraw[fill=white,draw=black,line width=\chemlabLineWidth] (0,2.35) ellipse (0.7 and 0.1);
  \end{scope}
}

% ============================================================
% 12. 试管 (TubeEssais)
% ============================================================
\newcommand{\tkzTestTubeStopper}{%
  \chemlabUpdateLineWidths
  \begin{scope}[fill=LightGray,draw=black,line width=\chemlabLineWidth,line cap=round,line join=round]
    \filldraw (0,-0.3) ellipse (0.4 and 0.1);
    \filldraw (-0.4,-0.3) -- (-0.6,0.5) -- (0.6,0.5) -- (0.4,-0.3) -- cycle;
    \fill[draw=none] (0,-0.3) ellipse (0.4 and 0.1);
    \filldraw (0,0.5) ellipse (0.6 and 0.1);
  \end{scope}
  \draw[line width=\chemlabEllipseLineWidth] plot[smooth,tension=0.5]
        coordinates {(-0.5,0) (0,-0.07) (0.5,0)};
}

\newcommand{\tkzTestTubeBody}{%
  \chemlabUpdateLineWidths\begin{scope}[line width=\chemlabLineWidth,line cap=round,line join=round]
    % 合并左侧直线 + 底部圆弧 + 右侧直线为单条路径, 平滑转角
    \draw (-0.5,3.5) -- (-0.5,0.5)
          arc[start angle=180,end angle=360,radius=0.5]
          -- (0.5,3.5);
    % 试管口椭圆 (原 pst-labo 用 0.08cm, 比其他椭圆粗)
    \draw[line width=\chemlabTubeEllipseLineWidth] (0,3.5) ellipse (0.5 and 0.1);
  \end{scope}
}

\newcommand{\tkzTestTubeStraightTube}{%
  \draw[line width=0.5\pgflinewidth] (0.05,-1) -- (0.05,5);
  \draw[line width=0.5\pgflinewidth] (-0.05,-1) -- (-0.05,5);
  \tkzTestTubeStopper
  \draw[fill=white,rounded corners=0.5,line width=0.5\pgflinewidth]
       (-0.05,5) rectangle (0.05,0.5);
  \draw[white,line width=0.14cm] (-0.1,5) -- (0.1,5);
}

\newcommand{\tkzTestTubeElbowTube}{%
  \draw[line width=0.5\pgflinewidth] (0.05,-1) -- (0.05,0);
  \draw[line width=0.5\pgflinewidth] (-0.05,-1) -- (-0.05,0);
  \tkzTestTubeStopper
  % 弯管路径: pst-labo \pspolygon(...)(8,2)(8,1.9)...
  \filldraw[fill=white,rounded corners=0.05cm,line width=0.5\pgflinewidth]
       (-0.05,0.5) -- (-0.05,2) -- (8,2) -- (8,1.9)
       -- (0.05,1.9) -- (0.05,0.5) -- cycle;
  \draw[white,line width=0.14cm] (7.95,1.9) -- (8.05,2);
}

\newcommand{\tkzTestTubeUTube}{%
  \draw[line width=0.5\pgflinewidth] (0.05,-1) -- (0.05,0);
  \draw[line width=0.5\pgflinewidth] (-0.05,-1) -- (-0.05,0);
  \tkzTestTubeStopper
  \filldraw[fill=white,rounded corners=0.05cm,line width=0.5\pgflinewidth]
       (-0.05,0.5) -- (-0.05,2) -- (4.5,2) -- (4.5,-3)
       -- (4.4,-3) -- (4.4,1.9) -- (0.05,1.9) -- (0.05,0.5) -- cycle;
  \draw[white,line width=0.14cm] (4.3,-3) -- (4.6,-3);
}

\newcommand{\tkzTestTubeUTubeLong}{%
  \draw[line width=0.5\pgflinewidth] (0.05,-1) -- (0.05,0);
  \draw[line width=0.5\pgflinewidth] (-0.05,-1) -- (-0.05,0);
  \tkzTestTubeStopper
  \filldraw[fill=white,rounded corners=0.05cm,line width=0.5\pgflinewidth]
       (-0.05,0.5) -- (-0.05,2) -- (4.5,2) -- (4.5,-9)
       -- (4.4,-9) -- (4.4,1.9) -- (0.05,1.9) -- (0.05,0.5) -- cycle;
  \draw[white,line width=0.14cm] (4.3,-9) -- (4.6,-9);
}

\newcommand{\tkzCurvedTube}{%
  \draw[line width=0.5\pgflinewidth] (0.05,-1) -- (0.05,0);
  \draw[line width=0.5\pgflinewidth] (-0.05,-1) -- (-0.05,0);
  \tkzTestTubeStopper
  \begin{scope}[line width=0.4pt]
    \draw (-0.05,0.5) -- (-0.05,1);
    \draw ([shift={(0.5,1)}]180:0.55) arc[start angle=180,end angle=90,radius=0.55];
    \draw (0.5,1.55) -- (5.5,1.55);
    \draw ([shift={(5.5,1)}]90:0.55) arc[start angle=90,end angle=0,radius=0.55];
    \draw (6.05,1) -- (6.05,-8.8);
    \draw ([shift={(6.5,-8.8)}]180:0.45) arc[start angle=180,end angle=360,radius=0.45];
    \draw (6.95,-8.8) -- (6.95,-8.3);
    \draw (7.05,-8.3) -- (7.05,-8.8);
    \draw ([shift={(6.5,-8.8)}]0:0.55) arc[start angle=0,end angle=-180,radius=0.55];
    \draw (5.95,-8.8) -- (5.95,1);
    \draw ([shift={(5.5,1)}]0:0.45) arc[start angle=0,end angle=90,radius=0.45];
    \draw (5.5,1.45) -- (0.5,1.45);
    \draw ([shift={(0.5,1)}]90:0.45) arc[start angle=90,end angle=180,radius=0.45];
    \draw (0.05,1) -- (0.05,0.5);
    \draw (-0.05,0.5) -- (0.05,0.5);
  \end{scope}
  \fill[white] (-0.025,0.52) rectangle (0.020,0.6);
}

\newcommand{\tkzCurvedTubeShort}{%
  \draw[line width=0.5\pgflinewidth] (0.05,-1) -- (0.05,0);
  \draw[line width=0.5\pgflinewidth] (-0.05,-1) -- (-0.05,0);
  \tkzTestTubeStopper
  \draw (-0.05,0.5) -- (-0.05,1);
  \begin{scope}[line width=0.4pt]
    \draw ([shift={(0.5,1)}]180:0.55) arc[start angle=180,end angle=90,radius=0.55];
    \draw (0.5,1.55) -- (3.5,1.55);
    \draw ([shift={(3.5,1)}]90:0.55) arc[start angle=90,end angle=0,radius=0.55];
    \draw (4.05,1) -- (4.05,-2.8);
    \draw ([shift={(4.5,-2.8)}]180:0.45) arc[start angle=180,end angle=360,radius=0.45];
    \draw (4.95,-2.8) -- (4.95,-2.3);
    \draw (5.05,-2.3) -- (5.05,-2.8);
    \draw ([shift={(4.5,-2.8)}]0:0.55) arc[start angle=0,end angle=-180,radius=0.55];
    \draw (3.95,-2.8) -- (3.95,1);
    \draw ([shift={(3.5,1)}]0:0.45) arc[start angle=0,end angle=90,radius=0.45];
    \draw (3.5,1.45) -- (0.5,1.45);
    \draw ([shift={(0.5,1)}]90:0.45) arc[start angle=90,end angle=180,radius=0.45];
    \draw (0.05,1) -- (0.05,0.5);
    \draw (-0.05,0.5) -- (0.05,0.5);
  \end{scope}
  \fill[white] (-0.025,0.52) rectangle (0.020,0.6);
}

\newcommand{\tkzCurvedTubeShortDoubleTube}{%
  \begin{scope}[line width=0.4pt]
    \draw (-0.05,0.5) -- (-0.05,1);
    \draw ([shift={(0.5,1)}]180:0.55) arc[start angle=180,end angle=90,radius=0.55];
    \draw (0.5,1.55) -- (3.5,1.55);
    \draw ([shift={(3.5,1)}]90:0.55) arc[start angle=90,end angle=0,radius=0.55];
    \draw (4.05,1) -- (4.05,-2.8);
    \draw ([shift={(4.5,-2.8)}]180:0.45) arc[start angle=180,end angle=360,radius=0.45];
    \draw (4.95,-2.8) -- (4.95,-2.3);
    \draw (5.05,-2.3) -- (5.05,-2.8);
    \draw ([shift={(4.5,-2.8)}]0:0.55) arc[start angle=0,end angle=-180,radius=0.55];
    \draw (3.95,-2.8) -- (3.95,1);
    \draw ([shift={(3.5,1)}]0:0.45) arc[start angle=0,end angle=90,radius=0.45];
    \draw (3.5,1.45) -- (0.5,1.45);
    \draw ([shift={(0.5,1)}]90:0.45) arc[start angle=90,end angle=180,radius=0.45];
    \draw (0.05,1) -- (0.05,0.5);
    \draw (-0.05,0.5) -- (0.05,0.5);
  \end{scope}
  \fill[white] (-0.025,0.52) rectangle (0.018,0.8);
}

\newcommand{\tkzTestTubeDoubleTube}{%
  \begin{scope}[line width=0.4pt]
    \draw (-0.3,-3) -- (-0.3,0);
    \draw (-.2,-3) -- (-0.2,0);
    \draw (0.25,-1) -- (0.25,0);
    \draw (.15,-1) -- (.15,0);
  \end{scope}
  \tkzTestTubeStopper
  \begin{scope}[shift={(0.2,0)}]\tkzCurvedTubeShortDoubleTube\end{scope}
  \begin{scope}[shift={(0,0.5)}]
    \begin{scope}[line width=0.4pt]
      \filldraw[fill=white,rounded corners=0.05cm]
        (-.3,0) -- (-0.3,2) -- (-0.5,2.2) -- (-0.5,3) -- (0,3)
                -- (0,2.2) -- (-0.2,2) -- (-0.2,0) -- cycle;
      \draw (0.1,1.55) -- (-.5,1.55) -- (-0.5,1.7) -- (-.55,1.7)
            -- (-.55,1.3) -- (-0.5,1.3) -- (-0.5,1.55) -- (-.5,1.45) -- (.1,1.45);
      \filldraw[fill=white,rounded corners=0.2] (-0.4,1.3) rectangle (-0.1,1.7);
    \end{scope}
    \fill[white] (-0.25,3) ellipse (0.3 and 0.1);
  \end{scope}
}

% 试管夹 (Pince) - 几何与原 pst-labo pst-laboObj.tex line 483-508 一致
% 根据原 pst-labo 实际渲染 (5.3x 截图) 校准:
%   - 上边缘凹弧较浅 (原 \pscurve 3点几乎共线, 浅弧)
%   - 顶部小凸起为圆弧形 (\psarcn 圆弧)
%   - 所有转角为圆角 (原 \pscustom 默认 linejoin=0 即 round)
\newcommand{\tkzTestTubeClamp}{%
  % 第1块: 上部夹爪
  % 上边缘(1a, 夹合面): 下凹弧贴合试管顶部弧面
  %   试管半径 0.5, 在 y=3.2 处试管顶部弧面下凹约 0.1, 用贝塞尔曲线模拟
  %   从 (0.5,0) 经 (0,-0.1) 到 (-0.5,0), 控制点 y=-0.13 产生平滑下凹
  \filldraw[fill=PaleOrange,draw=black,line cap=round,line join=round]
    (0.5,0) .. controls (0.25,-0.13) and (-0.25,-0.13) .. (-0.5,0)
    -- (-0.8,0)
    -- ([shift={(-0.5,0)}]180:0.3) arc[start angle=180,end angle=270,radius=0.3]
    -- (-.5,-.3) -- (2.5,-.3) -- (2.4,-.2) -- (1,0) -- (0.5,0) -- cycle;
  % 第2块: 下部夹爪
  % 注: 原 \psarcn(-.5,0){0.3}{-90}{180} 是顺时针从 -90°(=270°) 到 180°, 跨度 90° (四分之一圆弧)
  % TikZ arc 当 end<start 时顺时针, 需用 end=-180 让其走短弧 (避免四分之三圆弧)
  \filldraw[fill=PaleOrange,draw=black,line cap=round,line join=round]
    (-0.8,0) -- (-0.8,-0.3)
    -- ([shift={(-0.5,-0.3)}]180:0.3) arc[start angle=180,end angle=270,radius=0.3]
    -- (-0.5,-0.6) -- (2.5,-.6) -- (2.5,-.3) -- (-.5,-.3)
    -- ([shift={(-.5,0)}]-90:0.3) arc[start angle=-90,end angle=-180,radius=0.3] -- cycle;
  % 第3块: 顶部小凸起 - 原 \psarcn 顺时针 180°->90°, 圆弧形凸起
  % 去掉 (-0.8,0.05)->(-0.75,0.05) 的多余水平线段, 直接从 (-0.5,0.05) 到圆弧起点
  % 使凸起呈纯圆弧形 (与原 pst-labo 渲染一致)
  \filldraw[fill=PaleOrange,draw=black,line cap=round,line join=round]
    (-.5,0.05) -- ([shift={(-0.5,0.05)}]180:0.25)
    arc[start angle=180,end angle=90,radius=0.25]
    -- (-.5,0.3) -- (-0.5,0.05) -- cycle;
  % 第4块: 夹臂上部
  \filldraw[fill=PaleOrange,draw=black,line cap=round,line join=round]
    (0.5,0.3) -- (4.9,0.3) -- (5,0.2) -- (1,0) -- (0.5,0) -- (0.5,0.3) -- cycle;
  % 第5块: 夹臂下部 (Brown 棕色)
  \filldraw[fill=Brown,draw=black,line cap=round,line join=round]
    (0.5,0) -- (5,0.2) -- (5,0) -- (2,-0.2) -- (1,0) -- cycle;
  % 螺丝椭圆
  \draw[line width=1.5\pgflinewidth,line cap=round] (1,0) ellipse (0.3 and 0.1);
  % 螺丝下方的杆
  \draw[line width=1.5\pgflinewidth,line cap=round,line join=round]
        (0.7,0) -- (0.4,-.3) -- (0.4,-.6);
}

% ============================================================
% 13. 容量瓶 (Fiole jaugée)
% ============================================================
\newcommand{\tkzVolumetricFlask}{%
  \chemlabUpdateLineWidths
  \draw[rounded corners=0.1cm,line width=\chemlabLineWidth,line cap=round,line join=round]
       (-0.21,3.5) -- (-0.1,3.5) -- (-0.1,1.5)
       -- (-0.75,0) -- (0.75,0) -- (0.1,1.5) -- (0.1,3.5) -- (0.21,3.5);
  \draw[line width=\chemlabGraduationLineWidth,line cap=round] (-0.1,2.5) -- (0.1,2.5);
}

% ============================================================
% 14. 结晶皿 (Cristallisoir)
% ============================================================
\newcommand{\Cristallisoir}{%
  \chemlabUpdateLineWidths
  \begin{scope}[line width=\chemlabLineWidth,line cap=round,line join=round]
  \draw (-1.9,2) -- (-2,2);
  \draw ([shift={(-2,1.95)}]90:0.05) arc[start angle=90,end angle=270,radius=0.05];
  \draw (-2,1.9) -- (-2,1);
  \draw ([shift={(-1,1)}]180:1) arc[start angle=180,end angle=270,radius=1];
  \draw (-1,0) -- (1,0);
  \draw ([shift={(1,1)}]270:1) arc[start angle=270,end angle=360,radius=1];
  \draw (2,1) -- (2,1.9);
  \draw ([shift={(2,1.95)}]270:0.05) arc[start angle=270,end angle=90,radius=0.05];
  \draw (2,2) -- (1.9,2);
  \draw (1.9,2) -- (1.9,1);
  \draw ([shift={(1,1)}]0:0.9) arc[start angle=0,end angle=-90,radius=0.9];
  \draw (1,0.1) -- (-1,0.1);
  \draw ([shift={(-1,1)}]270:0.9) arc[start angle=270,end angle=180,radius=0.9];
  \draw (-1.9,1) -- (-1.9,2);
  \end{scope}
}

\newcommand{\InterieurCristallisoir}{%
  (1.9,2) -- (1.9,1)
  -- ([shift={(1,1)}]0:0.9) arc[start angle=0,end angle=-90,radius=0.9]
  -- (1,0.1) -- (-1,0.1)
  -- ([shift={(-1,1)}]270:0.9) arc[start angle=270,end angle=180,radius=0.9]
  -- (-1.9,1) -- (-1.9,2) -- cycle
}

\newcommand{\tkzCrystallizingDish}{%
  \begin{scope}
    \clip \InterieurCristallisoir;
    \fill[LightBlue1] (-3,-1) rectangle (3,1.5);
  \end{scope}
  \Cristallisoir
}

% ============================================================
% 15. 倒置试管 (TubeRenverse)
% ============================================================
\newcommand{\tkzInvertedTube}{%
  \chemlabUpdateLineWidths
  \begin{scope}[line width=\chemlabEllipseLineWidth,line cap=round,line join=round]
    \fill[LightBlue1] (-0.5,0) rectangle (0.5,5.5);
    % 合并左侧 + 顶部圆弧 + 右侧为单条路径
    % pst-labo: \psarcn(0,6){0.5}{180}{0} = 从180°顺时针到0°, 经过顶部90°
    % TikZ: arc[start=180, end=0] = 从180°逆时针(减小)到0°, 经过顶部90°
    \draw (-0.5,0) -- (-0.5,6)
          arc[start angle=180,end angle=0,radius=0.5]
          -- (0.5,0);
  \end{scope}
}

% ============================================================
% 16. 试管架 (ReposeTube)
% ============================================================
\newcommand{\tkzTestTubeRack}{%
  \begin{scope}[line width=0.1cm,draw=lightgray]
    \draw ([shift={(0,-0.25)}]100:1) arc[start angle=100,end angle=140,radius=1];
    \draw ([shift={(0,-.25)}]150:1) arc[start angle=150,end angle=160,radius=1];
    \draw ([shift={(0,-0.25)}]20:1) arc[start angle=20,end angle=80,radius=1];
  \end{scope}
}

% ============================================================
% 17. 烧杯边缘 + 刻度 (RebordBecher)
% ============================================================
\newcommand{\tkzBeakerRim}{%
  \chemlabUpdateLineWidths
  \draw[line width=\chemlabLineWidth] (0,2.55) ellipse (1.1 and .1);
  % 主刻度 : 粗细为细刻度线的 1.5 倍 (原 pst-labo: 0.03cm = 1.5 × 0.02cm)
  % 用 \pgfmathsetmacro 预算尺寸, 避免 \dimexpr 在 pgfkeys 中的兼容问题
  \pgfmathsetlengthmacro{\chemlab@majorGrad}{1.5*\chemlabGraduationLineWidth}%
  \foreach \nDiv in {0.5,1.0,...,2.0} {%
    \draw[line width=\chemlab@majorGrad] (-.75,\nDiv) -- (-.4,\nDiv);}
  % 细分刻度
  \foreach \nSubDiv in {0.5,0.6,...,2.0} {%
    \draw[line width=\chemlabGraduationLineWidth] (-0.75,\nSubDiv) -- (-.5,\nSubDiv);}
}

% ============================================================
% 18. 烧杯 (Becher)
% ============================================================
% 注 : 原 PSTricks 使用 \psline[linearc=0.5](-1,2.5)(-1,0)(1,0)(1,2.5)
% TikZ 的 rounded corners 与 PSTricks linearc 语义略有不同,
% 因此改用显式圆弧精确模拟 PSTricks linearc=0.5 的行为 :
% 左下角圆弧 : 圆心 (-0.5,0.5), 半径 0.5, 从 180° 到 270°
% 右下角圆弧 : 圆心 (0.5,0.5),  半径 0.5, 从 270° 到 360°
\newcommand{\tkzBeakerBody}{%
  \chemlabUpdateLineWidths
  \draw[line width=\chemlabLineWidth,line cap=round,line join=round]
       (-1,2.5) -- (-1,0.5)
       -- ([shift={(-0.5,0.5)}]180:0.5) arc[start angle=180,end angle=270,radius=0.5]
       -- (0.5,0)
       -- ([shift={(0.5,0.5)}]270:0.5) arc[start angle=270,end angle=360,radius=0.5]
       -- (1,2.5);
}

% ============================================================
% 19. 锥形瓶 (Erlenmeyer)
% ============================================================
% 注 : 原 PSTricks 使用 \psline[linearc=0.3](-0.5,3.5)(-0.5,2.5)(-1.5,0)(1.5,0)(0.5,2.5)(0.5,3.5)
% TikZ 的 rounded corners 与 PSTricks linearc 语义不同,
% 因此改用显式圆弧精确模拟 linearc=0.3 (由 scripts/erlen_arcs5.py 计算) :
%   肩部 (凹转角, 内侧角 21.8°) : 半径 0.3 太大无法相切 (切点距 1.558 > 颈部线段长 1.0), 保留尖角
%   左下 (凸转角, 内侧角 68.2°) : center=(-1.0569,0.3), r=0.3, 158.20° -> 270° (ccw)
%   右下 (凸转角, 内侧角 68.2°) : center=(1.0569,0.3),  r=0.3, -90° -> 21.80° (ccw)
\newcommand{\tkzErlenmeyerBody}{%
  \chemlabUpdateLineWidths
  \draw[line width=\chemlabLineWidth,line cap=round,line join=round]
       (-0.5,3.5) -- (-0.5,2.5)
       -- (-1.3354,0.4114)
       -- ([shift={(-1.0569,0.3)}]158.20:0.3) arc[start angle=158.20,end angle=270,radius=0.3]
       -- (1.0569,0)
       -- ([shift={(1.0569,0.3)}]-90:0.3) arc[start angle=-90,end angle=21.80,radius=0.3]
       -- (0.5,2.5) -- (0.5,3.5);
}

% ============================================================
% 20. 试剂瓶 (Flacon)
% ============================================================
\newcommand{\tkzReagentBottleBody}{%
  \chemlabUpdateLineWidths\begin{scope}[line width=\chemlabLineWidth,line cap=round,line join=round]
    % 左侧瓶口上角小弧 (psarcn 90->0, 顺时针)
    \draw ([shift={(-0.6,3.4)}]90:0.1) arc[start angle=90,end angle=0,radius=0.1];
    % 左侧颈部直线
    \draw (-0.5,3.4) -- (-0.5,3.1);
    % 左侧瓶口下角凹弧 (psarcn 0->-90, 顺时针; TikZ end<start 自动顺时针)
    \draw ([shift={(-0.6,3.1)}]0:0.1) arc[start angle=0,end angle=-90,radius=0.1];
    % 左侧肩部大凸弧 (psarc 90->180, 逆时针)
    \draw ([shift={(-0.6,2.6)}]90:0.4) arc[start angle=90,end angle=180,radius=0.4];
    % 左侧瓶身直线
    \draw (-1,2.6) -- (-1,0.2);
    % 左侧底部圆角 (psarc 180->270, 逆时针)
    \draw ([shift={(-.8,0.2)}]180:0.2) arc[start angle=180,end angle=270,radius=0.2];
    % 底部
    \draw (-.8,0) -- (.8,0);
    % 右侧底部圆角 (psarc -90->0, 逆时针)
    \draw ([shift={(.8,0.2)}]-90:0.2) arc[start angle=-90,end angle=0,radius=0.2];
    % 右侧瓶身直线
    \draw (1,0.2) -- (1,2.6);
    % 右侧肩部大凸弧 (psarc 0->90, 逆时针)
    \draw ([shift={(0.6,2.6)}]0:0.4) arc[start angle=0,end angle=90,radius=0.4];
    % 右侧瓶口下角凹弧 (psarcn -90->180, 顺时针短弧; TikZ 需 end=-180 避免 270° 长弧)
    \draw ([shift={(0.6,3.1)}]-90:0.1) arc[start angle=-90,end angle=-180,radius=0.1];
    % 右侧颈部直线
    \draw (0.5,3.1) -- (0.5,3.4);
    % 右侧瓶口上角小弧 (psarcn 180->90, 顺时针)
    \draw ([shift={(0.6,3.4)}]180:0.1) arc[start angle=180,end angle=90,radius=0.1];
  \end{scope}
}

% ============================================================
% 21. 玻璃杯 (Verre) - F. Bréal
% ============================================================
\newcommand{\tkzWatchGlassBody}{%
  \chemlabUpdateLineWidths\begin{scope}[line width=\chemlabLineWidth,line cap=round,line join=round]
    % 合并左斜线 + 底部圆弧 + 右斜线为单条路径, 平滑转角
    \draw (-1,3.5) -- (-0.5,1)
          arc[start angle=180,end angle=360,radius=0.5]
          -- (1,3.5);
  \end{scope}
}

\newcommand{\tkzWatchGlassBase}{%
  \chemlabUpdateLineWidths\begin{scope}[line width=\chemlabLineWidth,rounded corners=0.03cm,line cap=round,line join=round]
    \draw (-0.5,1) -- (-0.34,0.2) -- (-1.2,0) -- (1.2,0) -- (0.34,0.2) -- (0.5,1);
  \end{scope}
}

% ============================================================
% 22. 漏斗 (Entonnoir)
% ============================================================
\newcommand{\tkzFunnelOutline}{%
  \begin{scope}[line cap=round,line join=round]
    \draw (-0.07,0) -- (-0.07,1) -- (-1,2.5);
    \draw (1,2.5) -- (0.07,1) -- (0.07,0.1);
  \end{scope}
}

\newcommand{\tkzFunnelBody}{%
  \begin{scope}
    \clip (-0.07,0) -- (-0.07,1) -- (-1,2.5)
          -- (1,2.5) -- (0.07,1) -- (0.07,0.1) -- cycle;
    \fill[white] (-2,-1) rectangle (2,3);
  \end{scope}
  \begin{scope}[rounded corners=0.25cm]
    \draw (-0.07,0) -- (-0.07,1) -- (-1,2.5);
    \draw (1,2.5) -- (0.07,1) -- (0.07,0.1);
  \end{scope}
  \draw (0,2.5) ellipse (1 and 0.1);
}

% ============================================================
% 23. 球形冷凝器 (Réfrigérant à boulles)
% ============================================================
\newcommand{\tkzCondenserStem}{%
  % 整个 TigeBoulles 作为一个连续路径 (匹配 pst-labo \pscustom)
  % pst-labo: \pscustom[linestyle=none]{...\fill} 填充 + 单独描边 (无闭合线)
  % TikZ: 分开 \fill (带 cycle 闭合填充) 和 \draw (无 cycle, 顶部敞口)
  % 顶部是敞口的 — 两端弧之间没有连线 (匹配 pst-labo)
  %
  % 填充: 带闭合线 (fill 需要 cycle)
  \fill[white]
    % 顶部圆弧与颈部 - 左侧
    ({-0.35+0.1*cos(90)},{7.5+0.1*sin(90)}) arc[start angle=90,end angle=0,radius=0.1]
    -- (-.25,7.5) -- (-.25,6.6)
    -- ({-.15+0.1*cos(180)},{6.6+0.1*sin(180)}) arc[start angle=180,end angle=270,radius=0.1]
    -- ({-.15+0.1*cos(90)},{6.4+0.1*sin(90)}) arc[start angle=90,end angle=0,radius=0.1]
    -- (-.05,6.4) -- (-0.05,6)
    % 左侧 7 个球状突起 (中心 x=-0.05, 半径 0.25, arc 90→270 向左凸)
    (-0.05,6.0) arc[start angle=90,end angle=270,radius=0.25] -- (-0.05,5.50) -- (-0.05,5.25)
    arc[start angle=90,end angle=270,radius=0.25] -- (-0.05,4.75) -- (-0.05,4.50)
    arc[start angle=90,end angle=270,radius=0.25] -- (-0.05,4.00) -- (-0.05,3.75)
    arc[start angle=90,end angle=270,radius=0.25] -- (-0.05,3.25) -- (-0.05,3.00)
    arc[start angle=90,end angle=270,radius=0.25] -- (-0.05,2.50) -- (-0.05,2.25)
    arc[start angle=90,end angle=270,radius=0.25] -- (-0.05,1.75) -- (-0.05,1.50)
    arc[start angle=90,end angle=270,radius=0.25] -- (-0.05,1.00) -- (-0.05,0.75)
    % 左侧底部到下管口
    -- (-0.05,-1.5) -- (0.05,-1.5) -- (0.05,1)
    % 右侧 7 个球状突起 (中心 x=0.05, 半径 0.25, arc 270→450 向右凸)
    arc[start angle=270,end angle=450,radius=0.25] -- (0.05,1.50) -- (0.05,1.75)
    arc[start angle=270,end angle=450,radius=0.25] -- (0.05,2.25) -- (0.05,2.50)
    arc[start angle=270,end angle=450,radius=0.25] -- (0.05,3.00) -- (0.05,3.25)
    arc[start angle=270,end angle=450,radius=0.25] -- (0.05,3.75) -- (0.05,4.00)
    arc[start angle=270,end angle=450,radius=0.25] -- (0.05,4.50) -- (0.05,4.75)
    arc[start angle=270,end angle=450,radius=0.25] -- (0.05,5.25) -- (0.05,5.50)
    arc[start angle=270,end angle=450,radius=0.25] -- (0.05,6.00) -- (0.05,6.40)
    % 顶部圆弧与颈部 - 右侧
    -- ({0.15+0.1*cos(180)},{6.4+0.1*sin(180)}) arc[start angle=180,end angle=90,radius=0.1]
    -- ({0.15+0.1*cos(-90)},{6.6+0.1*sin(-90)}) arc[start angle=-90,end angle=0,radius=0.1]
    -- (0.25,6.6) -- (0.25,7.5)
    -- ({0.35+0.1*cos(180)},{7.5+0.1*sin(180)}) arc[start angle=180,end angle=90,radius=0.1]
    -- cycle;
  %
  % 描边: 底部敞口 (匹配 pst-labo: 左右侧分别 \psline, 无底部连线)
  \draw
    % 顶部圆弧与颈部 - 左侧
    ({-0.35+0.1*cos(90)},{7.5+0.1*sin(90)}) arc[start angle=90,end angle=0,radius=0.1]
    -- (-.25,7.5) -- (-.25,6.6)
    -- ({-.15+0.1*cos(180)},{6.6+0.1*sin(180)}) arc[start angle=180,end angle=270,radius=0.1]
    -- ({-.15+0.1*cos(90)},{6.4+0.1*sin(90)}) arc[start angle=90,end angle=0,radius=0.1]
    -- (-.05,6.4) -- (-0.05,6)
    % 左侧 7 个球状突起
    (-0.05,6.0) arc[start angle=90,end angle=270,radius=0.25] -- (-0.05,5.50) -- (-0.05,5.25)
    arc[start angle=90,end angle=270,radius=0.25] -- (-0.05,4.75) -- (-0.05,4.50)
    arc[start angle=90,end angle=270,radius=0.25] -- (-0.05,4.00) -- (-0.05,3.75)
    arc[start angle=90,end angle=270,radius=0.25] -- (-0.05,3.25) -- (-0.05,3.00)
    arc[start angle=90,end angle=270,radius=0.25] -- (-0.05,2.50) -- (-0.05,2.25)
    arc[start angle=90,end angle=270,radius=0.25] -- (-0.05,1.75) -- (-0.05,1.50)
    arc[start angle=90,end angle=270,radius=0.25] -- (-0.05,1.00) -- (-0.05,0.75)
    -- (-0.05,-1.5);
  % 右侧 (单独 \draw, 底部不连接左侧)
  \draw
    (0.05,-1.5) -- (0.05,1)
    % 右侧 7 个球状突起
    arc[start angle=270,end angle=450,radius=0.25] -- (0.05,1.50) -- (0.05,1.75)
    arc[start angle=270,end angle=450,radius=0.25] -- (0.05,2.25) -- (0.05,2.50)
    arc[start angle=270,end angle=450,radius=0.25] -- (0.05,3.00) -- (0.05,3.25)
    arc[start angle=270,end angle=450,radius=0.25] -- (0.05,3.75) -- (0.05,4.00)
    arc[start angle=270,end angle=450,radius=0.25] -- (0.05,4.50) -- (0.05,4.75)
    arc[start angle=270,end angle=450,radius=0.25] -- (0.05,5.25) -- (0.05,5.50)
    arc[start angle=270,end angle=450,radius=0.25] -- (0.05,6.00) -- (0.05,6.40)
    -- ({0.15+0.1*cos(180)},{6.4+0.1*sin(180)}) arc[start angle=180,end angle=90,radius=0.1]
    -- ({0.15+0.1*cos(-90)},{6.6+0.1*sin(-90)}) arc[start angle=-90,end angle=0,radius=0.1]
    -- (0.25,6.6) -- (0.25,7.5)
    -- ({0.35+0.1*cos(180)},{7.5+0.1*sin(180)}) arc[start angle=180,end angle=90,radius=0.1];
  % 底部液滴
  \begin{scope}[shift={(0,-1.6)}]\tkzDrop\end{scope}
}

\newcommand{\tkzCondenserBody}{%
  % 渐变填充: pst-labo gradbegin=cyan (底), gradend=PaleOrange (顶), gradmidpoint=0
  % TikZ: bottom color=cyan (底), top color=PaleOrange (顶)
  \shade[top color=PaleOrange,bottom color=cyan]
    ([shift={(0,5.75)}]0:0.5) arc[start angle=0,end angle=180,radius=0.5]
    -- (-0.5,5.75) -- (-0.5,5.5) -- (-1.5,5.5)
    -- (-1.5,5.3) -- (-.5,5.3) -- (-.5,1.7) -- (-1.5,1.7)
    -- (-1.5,1.5) -- (-.5,1.5) -- (-0.5,1)
    -- ([shift={(0,1)}]180:0.5) arc[start angle=180,end angle=360,radius=0.5]
    -- (0.5,1) -- (0.5,5.75) -- cycle;
  % 描边
  \begin{scope}[line width=0.88pt]
    \draw ([shift={(0,5.75)}]0:0.5) arc[start angle=0,end angle=180,radius=0.5];
    \draw (-0.5,5.75) -- (-0.5,5.5) -- (-1.5,5.5);
    \draw (-0.5,5.75) -- (-0.5,5.5) -- (-1.5,5.5);
    \draw (-1.5,5.3) -- (-.5,5.3) -- (-.5,1.7) -- (-1.5,1.7);
    \draw (-1.5,1.5) -- (-.5,1.5) -- (-0.5,1);
    \draw ([shift={(0,1)}]180:0.5) arc[start angle=180,end angle=360,radius=0.5];
    \draw (0.5,1) -- (0.5,5.75);
  \end{scope}
  \draw[->,cyan] (-2.5,1.6) -- (-1.6,1.6);
  \draw[<-,PaleOrange] (-2.5,5.4) -- (-1.6,5.4);
}

\newcommand{\tkzCondenserBulbs}{%
  \begin{scope}[shift={(0,0.2)}]
    \tkzCondenserBody
    \tkzCondenserStem
  \end{scope}
  \tkzTestTubeStopper
  \fill[white,rounded corners=0.025cm]
       (-0.05,0.6) -- (-0.05,0.5) -- (0.05,0.5) -- (0.05,0.6) -- cycle;
}

% ============================================================
% 24. 加热球 (Chauffe-Ballon)
% ============================================================
\newcommand{\tkzHeaterRoundFlask}{%
  \filldraw[fill=white,draw=black] (-2,-1) rectangle (2,1);
  \draw (-1.25,0) circle (0.4);
  \begin{scope}[shift={(-1.25,0)}]
    \foreach \iA in {0,30,...,180} {\draw (\iA:0.45) -- (\iA:0.6);}
    \draw[->] (0,0) -- (60:0.3);
  \end{scope}
  % 按钮区域: pst-labo \psframe[doubleline=true] 产生双线边框
  % PSTricks doubleline: 外框在指定坐标, 内框向内偏移
  % 实测: 外框到内框偏移约 0.13 units (含线宽+间隔, at scale=0.5)
  % 外框 (在指定坐标处)
  \draw[black,line width=0.8pt] (1,-.5) rectangle (1.5,0.5);
  % 内框 (向内偏移, 填充灰色)
  \filldraw[fill=gray,draw=black,line width=0.8pt] ({1+0.14},{-.5+0.14}) rectangle ({1.5-0.14},{0.5-0.14});
  % 红色指示灯
  \fill[red] (1.25,0.7) circle (0.1);
}

% ============================================================
% 25. 回流烧瓶 (Ballon à reflux)
% ============================================================
\newcommand{\tkzRoundFlaskReflux}{%
  \begin{scope}
    \clip \tkzRoundFlaskSimplePath;   % 使用路径版本形成闭合 clip 区域
    \shade[top color=cyan,bottom color=white,draw=cyan]
      plot[smooth,tension=0.5] coordinates {%
        (-1.5,1.5) (-1.375,1.55) (-1.25,1.5) (-1.125,1.45)
        (-1,1.5) (-.875,1.55) (-.75,1.5) (-.625,1.45)
        (-.5,1.5) (-.375,1.55) (-.25,1.5) (-0.125,1.45)
        (0,1.5) (0.125,1.55) (.25,1.5) (0.375,1.45)
        (0.5,1.5) (0.625,1.55) (0.75,1.5) (0.875,1.45)
        (1,1.5) (1.125,1.55) (1.25,1.5) (1.375,1.45) (1.5,1.5)}
      -- (1.5,1.5) -- (1.5,0) -- (-1.5,0) -- (-1.5,1.5) -- cycle;
  \end{scope}
  \tkzRoundFlaskSimple
}

% ============================================================
% 26. 分液漏斗 (Ampoule à décanter / Separate Funnel)
% ============================================================
% 相 : 水相 / 有机相
\definecolor{AqueousPhaseColor}{cmyk}{0.215686,0.062745,0,0}
\definecolor{OrganicPhaseColor}{cmyk}{0.223529,0.094118,0.501961,0}

% pst-labo: \tkzSeparatoryFunnel 设 linewidth=0.04 (用户单位, 随 unit 缩放)
% TikZ: 用 \chemlabLineWidth (基础 0.053cm × scale) 近似, 匹配其他器皿线宽
\newcommand{\tkzSeparatoryFunnelBody}{%
  \chemlabUpdateLineWidths
  \draw[line width=\chemlabLineWidth,line cap=round,line join=round]
    (-0.65,15) -- (-0.55,13.75)
    .. controls (-1.55,13.54) and ( -1.92,12.7) .. (-1.92,11.65)
    .. controls (-1.92,10.6) and (-0.3,9.55) .. (-0.3,8.5)
    -- (-0.3,2)
    .. controls (-0.3,1.5) and (-0.14,1.5) .. (-0.14,1)
    -- (-0.14,-1)
    -- (0.14,-1)
    -- (0.14,1)
    .. controls (0.14,1.5) and (0.3,1.5) .. (0.3,2)
    -- (0.3,8.5)
    .. controls (0.3,9.55) and (1.92,10.6) .. (1.92,11.65)
    .. controls ( 1.92,12.7) and (1.55,13.54) .. (0.55,13.75)
    -- (0.65,15);
}

% 用于 clip 的轮廓
\newcommand{\tkzSeparatoryFunnelOutlineClip}{%
  (-0.65,15) -- (-0.55,13.75)
  .. controls (-1.55,13.54) and ( -1.92,12.7) .. (-1.92,11.65)
  .. controls (-1.92,10.6) and (-0.3,9.55) .. (-0.3,8.5)
  -- (-0.3,2)
  .. controls (-0.3,1.5) and (-0.14,1.5) .. (-0.14,1)
  -- (-0.14,-1) -- (0.14,-1) -- (0.14,1)
  .. controls (0.14,1.5) and (0.3,1.5) .. (0.3,2)
  -- (0.3,8.5)
  .. controls (0.3,9.55) and (1.92,10.6) .. (1.92,11.65)
  .. controls (1.92,12.7) and (1.55,13.54) .. (0.55,13.75)
  -- (0.65,15) -- cycle
}

% 顶盖 (Bouchon) - 继承 linewidth=0.04
\newcommand{\tkzStopper}{%
  \chemlabUpdateLineWidths
  \filldraw[fill=gray!10,draw=gray,line width=\chemlabLineWidth,line cap=round,line join=round]
    (-0.67,15.25) -- (0.67,15.25) -- (0.57,14) -- (-0.57,14) -- cycle;
  \filldraw[fill=gray!5,draw=gray,line width=\chemlabLineWidth,line cap=round,line join=round]
    ([shift={(0.555,15.365)}]-90:0.1145) arc[start angle=-90,end angle=-180,radius=0.1145]
    -- ([shift={(0,15.92)}]-45:0.67) arc[start angle=-45,end angle=225,radius=0.67]
    -- ([shift={(-0.555,15.365)}]0:0.1145) arc[start angle=0,end angle=-90,radius=0.1145]
    -- cycle;
}

% 开启的旋塞 (TapOpen)
% pst-labo: \psframe[fillstyle=solid] 默认 fillcolor=white (PSTricks 默认)
% TikZ: \fill 默认黑色, 必须显式 fill=white
\newcommand{\tkzTapOpen}{%
  \chemlabUpdateLineWidths
  \filldraw[fill=white,line width=\chemlabLineWidth] (-.2,0.5) rectangle (.2,1);
  \draw[line width=\chemlabLineWidth] (-.2,0.4) rectangle (-.3,1.1);
  \draw[line width=\chemlabLineWidth] (.2,0.4) rectangle (.3,1.1);
  \draw[line width=\chemlabLineWidth] (.3,0.6) -- (.4,0.6) -- (.5,0.3) -- (.9,0.3) -- (.9,1.2) -- (.5,1.2)
        -- (.4,0.9) -- (.3,0.9) -- cycle;
  \fill[cyan!50] (-.6,0.4) -- (-.85,0.6) -- (-.85,0.9) -- (-.6,1.1) -- (-.3,1.1) -- (-.3,0.4) -- cycle;
  \draw[line width=\chemlabLineWidth] (-0.6,0.4) -- (-0.6,1.1);
  \draw[line width=\chemlabLineWidth] (0.5,0.3) -- (0.5,1.2);
}

% 关闭的旋塞 (TapClosed)
% pst-labo: \psframe[fillstyle=solid] 默认 fillcolor=white
\newcommand{\tkzTapClosed}{%
  \chemlabUpdateLineWidths
  \filldraw[fill=white,line width=\chemlabLineWidth] (-.2,0.5) rectangle (.2,1);
  \draw[line width=\chemlabLineWidth] (-.2,0.4) rectangle (-.3,1.1);
  \draw[line width=\chemlabLineWidth] (.2,0.4) rectangle (.3,1.1);
  \draw[line width=\chemlabLineWidth] ([shift={(0.35,0.6)}]180:0.05) arc[start angle=180,end angle=90,radius=0.05]
        -- (0.4,0.65) -- (0.9,0.65) -- (0.9,0.85) -- (0.35,0.85)
        -- ([shift={(0.35,0.9)}]-90:0.05) arc[start angle=-90,end angle=-180,radius=0.05]
        -- cycle;
  \draw[line width=\chemlabLineWidth] (0.5,0.65) -- (0.5,0.85);
  \fill[cyan!50] (-.6,0.4) -- (-.85,0.6) -- (-.85,0.9) -- (-.6,1.1) -- (-.3,1.1) -- (-.3,0.4) -- cycle;
  \draw[line width=\chemlabLineWidth] (-0.6,0.4) -- (-0.6,1.1);
}

% \tkzSeparatoryFunnel  (主调用)
% 可选参数 : stopper=true/false, TapOpen=true/false, liquidLevel1, liquidLevel2
\newcommand{\tkzSeparatoryFunnel}[1][]{%
  \begingroup
  \chemlabUpdateLineWidths
  \pgfkeys{
    /tkzSeparatoryFunnel/.cd,
    stopper/.is if=TkzTestTubeStopper,
    TapOpen/.is if=TkzTapOpen,
    liquidLevel1/.store in=\tkzTestTubeLiquidLevelI,
    liquidLevel2/.store in=\tkzTestTubeLiquidLevelII,
    #1}
  % 默认值
  \ifdefined\tkzTestTubeLiquidLevelI\else\def\tkzTestTubeLiquidLevelI{11}\fi
  \ifdefined\tkzTestTubeLiquidLevelII\else\def\tkzTestTubeLiquidLevelII{12}\fi
  % 水相
  \begin{scope}
    \clip \tkzSeparatoryFunnelOutlineClip;
    \fill[AqueousPhaseColor] (-2,1) rectangle (2,\tkzTestTubeLiquidLevelI);
  \end{scope}
  % 有机相
  \begin{scope}
    \clip \tkzSeparatoryFunnelOutlineClip;
    \fill[OrganicPhaseColor] (-2,\tkzTestTubeLiquidLevelI)
          rectangle (2,\tkzTestTubeLiquidLevelII);
  \end{scope}
  \ifTkzTestTubeStopper\tkzStopper\fi
  \tkzSeparatoryFunnelBody
  \ifTkzTapOpen\tkzTapOpen\else\tkzTapClosed\fi
  % 出口管口 (pst-labo: \psellipse[fillstyle=solid,fillcolor=white] 既有白色填充又有黑色轮廓)
  \filldraw[fill=white,line width=\chemlabLineWidth] (0,-1) ellipse (0.15 and 0.25);
  % 瓶口圆点 (填充, 无需线宽)
  \fill[darkgray] (-0.675,15) circle (0.03);
  \fill[darkgray] (0.675,15) circle (0.03);
  \draw[darkgray,line width=\chemlabLineWidth] (-0.65,15) -- (0.65,15);
  % 颈与瓶身接合
  \draw[gray,line width=\chemlabLineWidth] (-0.55,13.75) -- (0.55,13.75);
  \endgroup
}

% ============================================================
% 27. 支架 (pstSupport)
% ============================================================
% pst-labo: \psset{dimen=m,linewidth=0.04} — linewidth=0.04 是用户单位, 随 unit 缩放
% TikZ: 用 \chemlabLineWidth (基础 0.053cm × scale), 与其他器皿一致
\newcommand{\tkzStand}{%
  \chemlabUpdateLineWidths
  \begin{scope}[line width=\chemlabLineWidth]
    \filldraw[fill=woodColor,draw=black,line join=round]
        (-2,-6) rectangle (5,-5);
    \filldraw[fill=woodColor,draw=black,line join=round]
        (2.75,2) rectangle (3.25,16);
    \filldraw[fill=woodColor,draw=black,line join=round]
        ([shift={(2.75,1.5)}]90:0.5) arc[start angle=90,end angle=180,radius=0.5]
        -- (2.25,1.5) -- (2.25,-5) -- (3.75,-5) -- (3.75,1.5)
        -- ([shift={(3.25,1.5)}]0:0.5) arc[start angle=0,end angle=90,radius=0.5]
        -- cycle;
    \draw (2,11.5) -- (1.55,10.5);
    \filldraw[fill=woodColor,draw=black,line join=round]
        (-2.5,11.5) -- (-2,11.5) -- (-1.55,10.5) -- (-2.5,10.5) -- cycle;
    \filldraw[fill=woodColor,draw=black,line join=round]
        (4,11.5) -- (2,11.5) -- (1.55,10.5) -- (4,10.5) -- cycle;
    \draw (-2.5,10.5) rectangle (4,11.5);
    % 旋钮: pst-labo \pscustom[style=woodColor]{...} 填充会自动闭合, 但描边只画实际路径
    % (不画对角闭合边). TikZ: \fill 自动闭合填充, \draw 画开放路径 (无 -- cycle).
    \fill[woodColor]
        (4,11.2) -- (4.5,11.2)
        -- ([shift={(5,11)}]158.2:0.5385) arc[start angle=158.2,end angle=-158.2,radius=0.5385]
        -- (4.5,10.8) -- (4,10.8) -- cycle;
    \draw
        (4,11.2) -- (4.5,11.2)
        -- ([shift={(5,11)}]158.2:0.5385) arc[start angle=158.2,end angle=-158.2,radius=0.5385]
        -- (4.5,10.8) -- (4,10.8);
    \draw (4.5,10.8) -- (4.5,11.2);
  \end{scope}
}

\endinput

%    \begin{tikzpicture}
%       \BuretteCorps  ... 等等
%    \end{tikzpicture}
%
%  映射规则 :
%    \psline(A)(B)(C)              -> \draw (A) -- (B) -- (C);
%    \psline{->}(A)(B)             -> \draw[->] (A) -- (B);
%    \pspolygon(A)(B)(C)           -> \filldraw (A) -- (B) -- (C) -- cycle;
%    \psframe(x1,y1)(x2,y2)        -> \draw (x1,y1) rectangle (x2,y2);
%    \psframe*[linecolor=c]        -> \fill[c] (..) rectangle (..);
%    \psframe[framearc=r]          -> \draw[rounded corners=r] (...) rectangle (...);
%    \pscircle(c){r}               -> \draw (c) circle (r);
%    \pscircle*[linecolor=c](c){r} -> \fill[c] (c) circle (r);
%    \psellipse(c)(a,b)            -> \draw (c) ellipse (a and b);
%    \psarc(c){r}{a1}{a2}          -> \draw ([shift={(c)}]a1:r) arc[start angle=a1,end angle=a2,radius=r];
%    \psarcn(c){r}{a1}{a2}         -> \draw ([shift={(c)}]a1:r) arc[start angle=a1,end angle=a2,radius=r];   (注意 : PSTricks 中 arcn 是逆时针 / 角度递减)
%    \psbezier(A)(B)(C)(D)         -> \draw (A) .. controls (B) and (C) .. (D);
%    \pscurve(P1)(P2)...(Pn)       -> \draw plot[smooth,tension=0.5] coordinates {(P1) (P2) ... (Pn)};
%    \pscustom{ ... \fill}         -> \filldraw / \path[fill] (整段路径) -- cycle;
%    \rput(x,y){...}               -> \begin{scope}[shift={(x,y)}] ... \end{scope}
%    \rput{angle}(x,y){...}        -> \begin{scope}[shift={(x,y)},rotate=angle] ... \end{scope}
%    \multido{\i=a+b}{n}{...}      -> \foreach \i in {a,a+b,...,a+(n-1)*b} {...}
%    \psclip{<path>} ... \endpsclip-> \begin{scope} \clip <path>; ... \end{scope}
%    fillstyle=solid,fillcolor=c   -> fill=c
%    fillstyle=gradient,gradbegin=..,gradend=.. -> \shade[left color=..,right color=..]
%    fillstyle=hlines              -> pattern=horizontal lines
%    doubleline=true               -> double
%    linestyle=dashed              -> dashed
%    framearc=r                    -> rounded corners=r (近似)
%    linearc=r                     -> rounded corners=r
% ============================================================

\NeedsTeXFormat{LaTeX2e}
\ProvidesFile{tikz-chemlabObj.tex}[2026/08/03 v1.0.0]

\RequirePackage{tikz}
\usetikzlibrary{arrows.meta,calc,patterns,patterns.meta,shadings,decorations.pathmorphing}

% ------------------------------------------------------------
% 颜色定义 (对应原包中使用的颜色)
% ------------------------------------------------------------
\definecolor{LightGray}{gray}{0.80}      % 浅灰
\definecolor{VeryLightGray}{cmyk}{0,0,0,0.1}  % 极浅灰 (label 标签背景)
\definecolor{LightBlue}{rgb}{0.68,0.85,0.90}
\definecolor{LightBlue1}{rgb}{0.60,0.80,1.00}
\definecolor{VeryPaleOrange}{cmyk}{0,0.1,0.3,0}  % 滴定管液体默认色
\definecolor{PaleOrange}{rgb}{1.00,0.78,0.55}
\definecolor{Brown}{rgb}{0.50,0.25,0.00}
\definecolor{woodColor}{rgb}{0.957,0.84,0.69}
\definecolor{darkWoodColor}{rgb}{0.50,0.25,0.00}
\definecolor{flameAndGrill}{rgb}{0.95,0.55,0.10}  % 估算颜色
\definecolor{Beige}{rgb}{0.96,0.96,0.86}  % Champagne 液体默认色

% ------------------------------------------------------------
% 公共缺省值 (供调用方按需修改)
% ------------------------------------------------------------
\tikzset{
  Distillation/.cd,
  DistillateColor/.store in=\DistillationDistillateColor,
  DistillateColor=LightBlue1,
  MixtureStyle/.style={fill=LightBlue},
}
\newcommand{\tkzDistillationDistillateColor}{LightBlue1}
\newcommand{\tkzDistillationMixtureStyle}{LightBlue}
\newcommand{\tkzTestTubeLiquidLevelI}{11}     % liquidLevel1
\newcommand{\tkzTestTubeLiquidLevelII}{12}    % liquidLevel2
\newif\ifTkzTapOpen
\newif\ifTkzTestTubeStopper

% ============================================================
% 1. 滴定管 (Burette)
% ============================================================
\newcommand{\tkzBuretteBody}{%
  \draw (-0.5,11.5) -- (-0.25,11) -- (-0.25,3) -- (0,2.75) -- (0,0.5);
  \draw (0.2,0.7) -- (0.2,2.75) -- (0.5,3) -- (0.5,11) -- (0.75,11.5);
}

\newcommand{\tkzBuretteStopcock}{%
  % pst-labo: \psframe[fillstyle=solid] → 填充默认白色 + 黑色边框
  \filldraw[fill=white] (-.1,1.6) rectangle (.3,2.1);
  \draw (-.1,1.5) rectangle (-.2,2.2);
  \draw (.3,1.5) rectangle (.4,2.2);
  \draw (.4,1.7) -- (.5,1.7) -- (.6,1.4) -- (.8,1.4)
        -- (.8,2.3) -- (.6,2.3) -- (.5,2) -- (.4,2) -- cycle;
  % 左侧旋钮: pst-labo \pspolygon + \psframe with fillstyle=solid,fillcolor=LightBlue
  % 需要 LightBlue 填充 + 黑色外轮廓线
  \begin{scope}[fill=LightBlue,draw=black]
    \filldraw (-.5,1.5) -- (-.75,1.7) -- (-.75,2) -- (-.5,2.2)
          -- (-.3,2.2) -- (-.3,1.5) -- cycle;
    \filldraw (-.5,1.5) rectangle (-.2,2.2);
  \end{scope}
}

\newcommand{\tkzBuretteGraduation}{%
  % 滴定管刻度 (原 pst-labo 无显式线宽, 这里用细刻度线宽)
  \foreach \n in {3.0,3.32,...,10.68} { \draw[line width=\chemlabGraduationLineWidth] (-0.25,\n) -- (0,\n); }
}

% ============================================================
% 2. pH 计 (pH-mètre)
% ============================================================
\newcommand{\tkzTitrationPHMeter}{%
  % boitier
  \filldraw[fill=white,draw=black,rounded corners=0.125cm] (-1,0) rectangle (1,3);
  \fill[LightGray] (-0.75,2) rectangle (0.75,2.75);
  \node[white] at (0,2.375) {PH};
  \draw (0,1.25) circle (0.3);
  \draw[->] (0,1.25) -- (0,1.5);
  \node at (0.65,1.5) {$^\circ$C};
  \draw (-0.5,0.5) circle (0.2);
  \draw (0.5,0.5) circle (0.2);
  % électrode — 完全透明填充 (不遮挡烧杯刻度线)
  % pst-labo: \psframe[framearc=0.5](2.25,2.5)(2.75,4.5) 外框圆角
  %   顶部圆角被下方 \psframe[LightGray](sharp) 覆盖, 底部圆角被 \pscircle 覆盖
  %   实际可见的外框是直角矩形, 此处用 sharp corners 匹配 pst-labo 可见效果
  \draw (2.25,2.5) rectangle (2.75,4.5);
  % 顶部 LightGray 填充 + 描边 (sharp corners)
  \filldraw[fill=LightGray,draw=black] (2.25,4) rectangle (2.75,4.5);
  \draw (2.5,2.25) circle (0.25);
  % fil
  \draw[line width=0.05cm,LightBlue,rounded corners=0.2cm]
        (2.5,4.5) -- (2.5,5) -- (1.25,4.75) -- (0,5) -- (0,3);
}

% ============================================================
% 3. 滴定支架 (Support dosage)
% ============================================================
\newcommand{\tkzTitrationStand}{%
  % pst-labo: \psframe(-1.5,0)(1.5,1.5) + {\psset{doubleline=true} ...}
  % 线宽: pst-labo 在 \tkzTitration 中设 linewidth=0.053 (×0.4cm = 0.0212cm ≈ 0.6pt)
  \chemlabUpdateLineWidths
  % 平台: 无填充, 只有边框
  \draw[line width=\chemlabLineWidth] (-1.5,0) rectangle (1.5,1.5);
  % 红色按钮: \pscircle*[linecolor=red] → 红色填充, 双线边框也是红色 (不可见)
  \fill[red] (0,0.75) circle (0.2);
  % 回字型矩形按钮: 无填充, 双线边框
  % PSTricks doubleline: 指定坐标在内外线之间
  % 实测: 外线向外偏移 0.085 units, 内线向内偏移 0.034 units
  \draw[line width=\chemlabLineWidth] ({-1.25-0.085},{0.3-0.085}) rectangle ({-0.75+0.085},{1.2+0.085});
  \draw[line width=\chemlabLineWidth] ({-1.25+0.034},{0.3+0.034}) rectangle ({-0.75-0.034},{1.2-0.034});
  % 环形旋钮: 无填充, 双线边框
  % 外线向外偏移, 内线向内偏移
  \draw[line width=\chemlabLineWidth] (1,0.75) circle ({0.25+0.085});
  \draw[line width=\chemlabLineWidth] (1,0.75) circle ({0.25-0.034});
}

% ============================================================
% 4. 磁力搅拌子 (Aimant)
% ============================================================
\newcommand{\tkzTitrationStirBar}{%
  % pst-labo: \psset{linewidth=0.053} 在外层 \tkzTitration 中设置, 随 scale 缩放
  \chemlabUpdateLineWidths
  \draw[line width=\chemlabLineWidth,rounded corners=0.075cm,fill=white] (-0.5,0.1) rectangle (0.5,0.4);
  \draw[line width=\chemlabLineWidth,rounded corners=0.025cm,fill=white] (-0.05,0) rectangle (0.05,0.5);
}

% ============================================================
% 5. 移液管 (Pipette)
% ============================================================
\newcommand{\tkzPipetteBody}{%
  \chemlabUpdateLineWidths\begin{scope}[line width=\chemlabLineWidth,rounded corners=0.05cm]
    \draw (-0.05,0) -- (-0.15,0.2) -- (-0.15,6)
          -- (0.15,6) -- (0.15,0.2) -- (0.05,0) -- cycle;
    \fill[white] (0,6) ellipse (0.15 and 0.05);
    \foreach \Level in {0.5,1.0,...,5.0} {%
      \draw[line width=\chemlabGraduationLineWidth] (0,\Level) -- (0.1,\Level);}
  \end{scope}
}

\newcommand{\tkzLargePipetteBody}{%
  \chemlabUpdateLineWidths\begin{scope}[line width=\chemlabLineWidth,rounded corners=0.05cm]
    \draw (0,0) -- (0,1) -- (-0.23,2.5) -- (-0.23,22.5)
          arc[start angle=-80,end angle=80,radius=0.45];
  \end{scope}
}

% ============================================================
% 6. 液滴 (Goutte)
% ============================================================
% 原文件中 \pst@goutte 被重复定义了两次, 此处使用后者 (unit=0.6 的版本)
% pst-labo: \psset{unit=0.6}\psbezier[linewidth=0.017] — linewidth=0.017 是
%   0.017 用户单位 = 0.017*0.6 = 0.0102cm (随 goutte 的 LOCAL unit=0.6 缩放, 
%   不随外层 unit 变化, 因为 \psset{unit=0.6} 覆盖了外层设置)
% TikZ: line width=0.0102cm (ABSOLUTE, 匹配 pst-labo 的固定 0.0102cm)
\newcommand{\tkzDrop}{%
  \begin{scope}[scale=0.6]
    \draw[line width=0.0102cm]
      (0,0) .. controls (0.25,-.4) and (-.25,-.4) .. (0,0);
  \end{scope}
}

% ============================================================
% 7. 量筒 (Eprouvette)
% ============================================================
\newcommand{\tkzGraduatedCylinderBody}{%
  \draw (-0.5,5) -- (-0.5,0) -- (0.5,0) -- (0.5,4.9) -- (0.6,4.93);
}

\newcommand{\tkzGraduatedCylinderRim}{%
  \draw plot[smooth,tension=0.5] coordinates {%
    (0.6,-0.07) (0.498,0.008) (0.483,0.026) (0.433,0.05)
    (0.353,0.0707) (0.25,0.0866) (0.13,0.0966) (0,0.1)
    (-0.13,0.0966) (-0.25,0.0866) (-0.353,0.0707) (-0.433,0.05)
    (-0.483,0.026) (-0.5,0) (-0.483,-0.026) (-0.353,-0.0707)
    (-0.25,-0.0866) (-0.13,-0.096) (0,-0.1) (0.13,-0.096)
    (0.25,-0.0866) (0.353,-0.0707) (0.433,-0.05)
    (0.498,-0.026) (0.6,-0.07)};
}

\newcommand{\tkzGraduatedCylinderBase}{%
  \fill[white]
    ([shift={(-0.95,0.5)}]0:0.45) arc[start angle=0,end angle=-90,radius=0.45]
    -- ([shift={(-0.95,0.025)}]90:0.025) arc[start angle=90,end angle=270,radius=0.025]
    -- (-0.95,0) -- (0.95,0)
    -- ([shift={(0.95,0.025)}]-90:0.025) arc[start angle=-90,end angle=90,radius=0.025]
    -- ([shift={(0.95,0.5)}]-90:0.45) arc[start angle=-90,end angle=-180,radius=0.45]
    -- plot[smooth,tension=0.5] coordinates {(0.5,0.4) (0,0.25) (-0.5,0.4)}
    -- cycle;
  \draw
    ([shift={(-0.95,0.5)}]0:0.45) arc[start angle=0,end angle=-90,radius=0.45]
    -- ([shift={(-0.95,0.025)}]90:0.025) arc[start angle=90,end angle=270,radius=0.025]
    -- (-0.95,0) -- (0.95,0)
    -- ([shift={(0.95,0.025)}]-90:0.025) arc[start angle=-90,end angle=90,radius=0.025]
    -- ([shift={(0.95,0.5)}]-90:0.45) arc[start angle=-90,end angle=-180,radius=0.45]
    -- plot[smooth,tension=0.5] coordinates {(0.5,0.4) (0,0.25) (-0.5,0.4)};
}

\newcommand{\tkzGraduatedCylinderGraduation}{%
  % 量筒刻度 : 主刻度 1.5× 细刻度 (与烧杯一致)
  \pgfmathsetlengthmacro{\chemlab@eprouvette@major}{1.5*\chemlabGraduationLineWidth}%
  \foreach \nyD in {0.5,1.0,...,4.5} { \draw[line width=\chemlab@eprouvette@major] (-0.4,\nyD) -- (-0.1,\nyD); }
  \foreach \nyd in {0.5,0.6,...,4.4} { \draw[line width=\chemlabGraduationLineWidth] (-0.4,\nyd) -- (-0.25,\nyd); }
}

% ============================================================
% 8. 分馏装置 (Distillation fractionnée)
% ============================================================
% \pst@goutte 这里又被定义了一次, 在 TikZ 版中沿用同一个 \tkzDrop

\newcommand{\tkzColumn}{%
  % pst-labo: linearc=0.3 (×unit, 不是绝对值) → TikZ: rounded corners=0.12cm (= 0.3 × 0.4cm)
  % 尖刺 (pointeG/D) 先画并用 clip 裁剪到蒸馏柱内部, 再画轮廓线覆盖切边。
  %   pst-labo 中 linearc=0.3 (R=0.3 user units), 尖刺端点恰在圆弧切点附近, 几乎
  %   不可见。TikZ 的 rounded corners 实际 R 比标称值大 ~2×, 圆弧切点更低,
  %   尖刺端点 (在 y_col=8 和 7.9) 会突出到圆弧外侧形成多余线头。用 clip 将
  %   尖刺裁剪到柱内, 再画轮廓线覆盖切边, 消除突出线头。
  \begin{scope}
    \clip[rounded corners=0.12cm]
      (-.25,10.5) -- (-.25,8.5) -- (-.75,8) -- (-.75,2.5) -- (-.25,2) -- (-.25,0)
      -- (0.25,0) -- (0.25,2) -- (.75,2.5) -- (.75,8) -- (0.25,8.5) -- (0.25,9.36)
      -- (0.25,10.5) -- cycle;
    % 左侧尖刺 (pointeG)
    \foreach \i in {3,4,...,8} {%
      \draw[line width=0.4\pgflinewidth] (-.75,\i) -- (-.25,\i-0.5)
          -- (-.75,\i-0.1);}
    % 右侧尖刺 (pointeD)
    \foreach \i in {3,4,...,8} {%
      \draw[line width=0.4\pgflinewidth] (.75,\i) -- (.25,\i-0.5)
          -- (.75,\i-0.1);}
  \end{scope}
  \draw[rounded corners=0.12cm]
    (-.25,10.5) -- (-.25,8.5) -- (-.75,8) -- (-.75,2.5)
    -- (-.25,2) -- (-.25,0);
  \draw[rounded corners=0.12cm]
    (0.25,0) -- (0.25,2) -- (.75,2.5) -- (.75,8)
    -- (0.25,8.5) -- (0.25,9.36);
  % 上部断口 (line 3a/3b):
  %   pst-labo 源码 line 3a: \psline(0.25,8.5)(0.25,9.30) 用于填充 linearc 圆角
  %   留下的间隙。实测发现 TikZ 的 rounded corners 实际圆弧半径比标称值大
  %   ~2× (在 scale=0.6 下 rounded corners=0.3cm 给出 R≈0.62cm), 圆弧已覆盖到
  %   x=0.25 处 y_col=8.93 附近, 故 line 3a 在任何 scale/line width 下均被圆弧
  %   覆盖, 不可见。为避免在编号图等彩色调试场景中残留多余线段, 此处不再
  %   单独绘制 line 3a。
  \draw (0.25,9.5) -- (0.25,10.5);
  % 填充 / 阴影
  \begin{scope}[pattern={Lines[angle=45,distance={1pt/sqrt(2)}]},pattern color=gray]
    \fill (-.25,10) rectangle (-.05,10.5);
    \fill (.25,10)  rectangle (0.05,10.5);
  \end{scope}
  \foreach \n in {2.5,3.5,...,7.5} {%
    \begin{scope}[shift={(-.25,\n)}]\tkzDrop\end{scope}
    \begin{scope}[shift={(.25,\n)}]\tkzDrop\end{scope}}
  \begin{scope}[shift={(-0.25,-.1)}]\tkzDrop\end{scope}
  \begin{scope}[shift={(0.25,-.1)}]\tkzDrop\end{scope}
}

% 公共别名: pst-labo 中 \tkzColumn 是用户可调用的公共宏 (小写)
% tikz-chemlab 内部用 \tkzColumn (驼峰), 这里注册公共别名保持兼容
\let\tkzColumn\tkzColumn

% réfrigérant (condenseur)
\newcommand{\tkzCondenser}{%
  \draw ([shift={(3,0)}]90:0.5) arc[start angle=90,end angle=270,radius=0.5];
  % pst-labo: \psarc(7,0){0.5}{270}{90} → 逆时针经过 0°/右侧
  \draw ([shift={(7,0)}]270:0.5) arc[start angle=270,end angle=450,radius=0.5];
  \fill[white] (0,-.05) rectangle (9,0.05);
  \begin{scope}[line width=0.4pt]
    \draw (-.075,.05) -- (9,0.05);
    \draw (0.02,-.05) -- (9,-.05);
  \end{scope}
  \draw (3,0.5) -- (3.5,0.5);
  \draw (3,-.5) -- (6.5,-.5);
  \draw (3.6,0.5) -- (7,0.5);
  \begin{scope}[line width=0.4pt]
    \draw (3.5,0.5) -- (3.5,1.5);
    \draw (3.6,0.5) -- (3.6,1.5);
    \draw (6.6,-.5) -- (6.6,-1.5);
    \draw (6.5,-.5) -- (6.5,-1.5);
  \end{scope}
  \draw (6.6,-.5) -- (7,-.5);
  \draw[->] (6.55,-2.5) -- (6.55,-1.5);
  \draw[->] (3.55,1.5)  -- (3.55,2.5);
}

\newcommand{\tkzThermometer}{%
  % pst-labo: \psline[linearc=0.05](-.05,.25)(-.05,2.75)(0.05,2.75)(0.05,0.25)
  % linearc=0.05 × unit, TikZ: rounded corners=0.05 (受 scale 缩放)
  \draw[rounded corners=0.025cm]
    (-.05,.25) -- (-.05,2.75) -- (0.05,2.75) -- (0.05,0.25) -- cycle;
  % pst-labo: \psframe[framearc=0.8,fillstyle=solid,fillcolor=red](-.15,-.25)(0.15,0.25)
  % framearc=0.8: 80% of half-min(0.3,0.5)=0.15 → radius=0.12
  \filldraw[fill=red,draw=black,rounded corners=0.06cm] (-.15,-.25) rectangle (0.15,0.25);
  % pst-labo: \psline[linecolor=red](0,2) — 红色竖线, 使用默认线宽
  % 在 pst-labo 渲染中较粗, 看起来像红色填充
  \draw[red,line width=0.8pt] (0,0) -- (0,2);
}

\newcommand{\tkzHeater}{%
  \draw (1.17,2.5) -- (0,2.5) -- (0,0) -- (5.5,0) -- (5.5,2.5) -- (4.33,2.5);
}

% 内部小量筒轮廓 - 路径版本 (供 clip 使用)
\newcommand{\tkzGraduatedCylinderPath}{%
  (-0.5,6) -- (-0.5,.5)
  -- ([shift={(-1,0.5)}]0:0.5) arc[start angle=0,end angle=-90,radius=0.5]
  -- (-1,0) -- (1,0)
  -- ([shift={(1,0.5)}]270:0.5) arc[start angle=270,end angle=180,radius=0.5]
  -- (0.5,0.5) -- (0.5,6) -- cycle
}

% 内部小量筒 (描边)
% pst-labo: \psset{linewidth=1.5\pslinewidth} — \pslinewidth 默认 0.8pt (ABSOLUTE)
% → 1.5*0.8pt = 1.2pt (绝对值, 不随 unit/scale 缩放)
% 注意: TikZ 的 [line width=N\pgflinewidth] 有循环引用 bug — \pgflinewidth 在
% 乘法前被更新, 导致 line width=3\pgflinewidth 实际给出 3.6pt 而非 1.2pt。
% 必须用显式 pt 值。
\newcommand{\tkzGraduatedCylinderSmall}{%
  \begin{scope}[line width=0.72pt,line cap=round,line join=round]
    \draw (-0.5,6) -- (-0.5,.5)
      -- ([shift={(-1,0.5)}]0:0.5) arc[start angle=0,end angle=-90,radius=0.5]
      -- (-1,0) -- (1,0)
      -- ([shift={(1,0.5)}]270:0.5) arc[start angle=270,end angle=180,radius=0.5]
      -- (0.5,0.5) -- (0.5,6);
  \end{scope}
}

% 带液体的量筒 (内部使用, 重命名以避免与主文件 \pstEproulette 冲突)
\newcommand{\tkzGraduatedCylinderWithLiquid}{%
  \begin{scope}
    \clip \tkzGraduatedCylinderPath;
    % pst-labo: \psframe*[linecolor=\psk@Distillation@DistillateColor]
    \fill[\tkzDistillateColor] (-1.5,-1) rectangle (1.5,1);
  \end{scope}
  \tkzGraduatedCylinderSmall
}

% (注 : 原 PSTricks 中 \pst@roundFlask 已合并到 \tkzRoundFlaskSimple, 此处不再单独定义,
%  避免与 tikz-chemlab.tex 中的用户命令 \tkzRoundFlask 冲突)

% 自定义密集斜线图案 (匹配 pst-labo hlines 的密度)
{\pgfpointorigin}{\pgfpoint{0.012cm}{0.012cm}}{\pgfpoint{0.012cm}{0.012cm}}
{\pgfsetlinewidth{0.0028cm}\pgfpathmoveto{\pgfpoint{0cm}{0cm}}\pgfpathlineto{\pgfpoint{0.012cm}{0.012cm}}\pgfusepath{stroke}}
{\pgfpointorigin}{\pgfpoint{0.012cm}{0.012cm}}{\pgfpoint{0.012cm}{0.012cm}}
{\pgfsetlinewidth{0.004cm}\pgfpathmoveto{\pgfpoint{0cm}{0cm}}\pgfpathlineto{\pgfpoint{0.012cm}{0.012cm}}\pgfusepath{stroke}}
% \pst@DistillationFractionnee 主图
\newcommand{\tkzDistillationFractionated}{%
  % --- roundFlask ---
  \def\flaskLocal{%
    % pst-labo: fillstyle=hlines,hatchwidth=0.007,hatchsep=0.03 → 水平线填充
    \begin{scope}[pattern={Lines[angle=45,distance={1pt/sqrt(2)}]},pattern color=gray]
      \fill (-.5,2.25) rectangle (-.25,3);
      \fill (0.5,2.25)  rectangle (0.25,3);
    \end{scope}
    % pst-labo 线宽: neck=0.053, arcs=0.03 (都是 ×unit 用户单位, 随 scale 缩放)
    % \chemlabUpdateLineWidths 已设置 \chemlabLineWidth = 0.053*scale
    \chemlabUpdateLineWidths
    \draw[line width=\chemlabLineWidth] (0.5,1.5) -- (0.5,3);
    \draw[line width=\chemlabLineWidth] (-.5,1.5) -- (-.5,3);
    % pst-labo: \psarc(0,0){1.581}{108.435}{71.565} 逆时针走长路 (经底部),
    %   左右侧弧可见, 底部被液体填充覆盖, 顶部无可见弧线。此处拆分为
    %   左右侧两段 (底部由下方轮廓弧绘制, 顶部不画)。
    % 右上侧弧: 71° → 0° (连接右颈底部到球体最右点)
    \draw[line width=\chemlabLineWidth] ([shift={(0,0)}]71.565:1.581) arc[start angle=71.565,end angle=0,radius=1.581];
    % 左上侧弧: 180° → 108° (连接球体最左点到左颈底部)
    \draw[line width=\chemlabLineWidth] ([shift={(0,0)}]180:1.581) arc[start angle=180,end angle=108.435,radius=1.581];
    % pst-labo: \psset{style=\psk@Distillation@MixtureStyle}\fill
    % 需要用 \shade (不是 \fill) 来产生渐变效果
    \shade[\tkzMixtureStyle] plot[smooth,tension=0.5] coordinates {%
      (-1.581,0) (-1.25,0.1) (-.75,-.1) (0,.1) (0.75,-.1) (1.25,.1) (1.581,0)}
      arc[start angle=0,end angle=-180,radius=1.581];
    % 底部轮廓弧: pst-labo \psarc[linewidth=0.03] → 0.03*unit, 用 \chemlabLineWidth 0.053
    % (pst-labo 中 0.03 vs 0.053 差异微小, 统一用 0.053*scale 避免额外变量)
    \draw[line width=\chemlabLineWidth] (0,0) ++(180:1.581) arc[start angle=180,end angle=360,radius=1.581];
  }%
  % --- 主图整体平移 ---
  \begin{scope}[shift={(0,-3)}]
    \begin{scope}[shift={(2.75,13.5)}]\tkzThermometer\end{scope}
    \begin{scope}[shift={(2.75,2.5)}]\flaskLocal\end{scope}
    \begin{scope}[shift={(2.75,4)}]\tkzColumn\end{scope}
    \begin{scope}[shift={(3,13.4)},rotate=-45]\tkzCondenser\end{scope}
    \begin{scope}[shift={(0,0)}]\tkzHeater\end{scope}
    \begin{scope}[shift={(9.4,0)}]\tkzGraduatedCylinderWithLiquid
      % 蒸馏装置量筒刻度 : 主刻度 1.5× 细刻度
      \pgfmathsetlengthmacro{\chemlab@distill@major}{1.5*\chemlabGraduationLineWidth}%
      \foreach \Level in {1,2,...,5} {\draw[line width=\chemlab@distill@major] (-.4,\Level) -- (-.1,\Level);}
      \foreach \Level in {0.2,0.4,...,4.8} {\draw[line width=\chemlabGraduationLineWidth] (-.4,\Level) -- (-.2,\Level);}
    \end{scope}
    \begin{scope}[shift={(9.37,7)}]\tkzDrop\end{scope}
    \begin{scope}[shift={(9.37,6)}]\tkzDrop\end{scope}
    \begin{scope}[shift={(9.37,4)}]\tkzDrop\end{scope}
    \begin{scope}[shift={(9.37,1.2)}]\tkzDrop\end{scope}
    % 加热盘 (搅拌器)
    \begin{scope}[shift={(0.75,.75)}]
      \draw (0,0) circle (.5);
      \foreach \i in {0,30,...,180} {\draw (\i:0.55) -- (\i:0.65);}
      \draw[->,line width=0.035cm] (0,0) -- (0,0.45);
    \end{scope}
  \end{scope}
}

% ============================================================
% 9. 本生灯 (Bec Bunsen) - 不带 grille
% ============================================================
\newcommand{\tkzBunsenBurner}{%
  \chemlabUpdateLineWidths\begin{scope}[line width=\chemlabLineWidth,line cap=round,line join=round]
    \draw (-1.25,0) rectangle (1.25,0.25);
    \draw (-.5,1.25) rectangle (0.5,2.25);
    \foreach \n in {-0.3,0.0,0.3} {\draw (\n,1.75) circle (0.1);}
    \draw (-.25,2.25) rectangle (0.25,4.25);
    \draw (0.25,1.25) -- (0.25,0.5) -- (1.25,0.25);
    \draw (-1.25,0.25) -- (-.25,0.5) -- (-0.25,0.75);
    \draw (-2.25,0.75) -- (-.25,0.75);
    \draw (-2.25,1) -- (-.25,1);
    \draw (-.25,0.875) ellipse (0.1 and 0.125);
    \fill[white] (-2.25,0.75) rectangle (-0.25,1);
    \draw (-2.25,0.75) -- (-0.25,0.75);
    \draw (-2.25,1) -- (-0.25,1) -- (-.25,1.25);
    \draw plot[smooth,tension=0.5] coordinates {(-0.25,0.5) (0,0.4) (0.25,0.5)};
    % flamme
    % pst-labo 源码: \rput(0,4.25){\psclip{\psbezier[...gradient...] (12点)}\pspolygon[...gradient...](-0.25,0)(0.25,0)(0,1)\endpsclip}
    % \psbezier 12个点 + BezierNArray padding (0,0) = 13点 = 4段贝塞尔
    % 经像素级验证, 正确路径从 (0,0) 出发, 经4段贝塞尔回到 (0,0), 无闭合直线:
    %   Seg 1: (0,0) → c1=P1=(-0.25,0), c2=P2=(-0.35,0.5) → P3=(-0.4,0.75) [左下→左中]
    %   Seg 2: P3=(-0.4,0.75) → c1=P4=(-0.35,1), c2=P5=(-0.25,1.5) → P6=(0.5,2) [左中→顶]
    %   Seg 3: P6=(0.5,2) → c1=P7=(0.25,1.5), c2=P8=(0.35,1) → P9=(0.4,0.75) [顶→右中]
    %   Seg 4: P9=(0.4,0.75) → c1=P10=(0.35,0.5), c2=P11=(0.25,0) → (0,0) [右中→底]
    % 渐变: gradbegin=PaleOrange (底), gradend=yellow (顶), gradmidpoint=0, gradangle=0 (垂直)
    \begin{scope}[shift={(0,4.25)}]
      % 外焰填充: PaleOrange (底) -> yellow (顶) 线性垂直渐变
      \shade[top color=yellow, bottom color=PaleOrange]
        (0,0) .. controls (-0.25,0) and (-0.35,0.5)
        .. (-0.4,0.75) .. controls (-0.35,1) and (-0.25,1.5)
        .. (0.5,2) .. controls (0.25,1.5) and (0.35,1)
        .. (0.4,0.75) .. controls (0.35,0.5) and (0.25,0)
        .. (0,0);
      % 内焰: 三角形 cyan(底)->white(顶) 渐变, 裁剪到外焰形状内
      \begin{scope}
        \clip
          (0,0) .. controls (-0.25,0) and (-0.35,0.5)
          .. (-0.4,0.75) .. controls (-0.35,1) and (-0.25,1.5)
          .. (0.5,2) .. controls (0.25,1.5) and (0.35,1)
          .. (0.4,0.75) .. controls (0.35,0.5) and (0.25,0)
          .. (0,0);
        \shade[top color=white, bottom color=cyan]
          (-0.25,0) -- (0.25,0) -- (0,1) -- cycle;
      \end{scope}
    \end{scope}
  \end{scope}
}

% 火焰 + 网格 (FlammeGrille 形状) - 路径版本 (供 clip 使用)
\newcommand{\tkzFlameGrillPath}{%
  ([shift={(0,0.75)}]270:0.75) arc[start angle=270,end angle=360,radius=0.75]
  -- ([shift={(1.5,0.75)}]180:0.75) arc[start angle=180,end angle=90,radius=0.75]
  -- (1.5,1.5) -- (-1.5,1.5)
  -- ([shift={(-1.5,0.75)}]90:0.75) arc[start angle=90,end angle=0,radius=0.75]
  -- ([shift={(0,0.75)}]180:0.75) arc[start angle=180,end angle=270,radius=0.75]
  -- cycle
}

\newcommand{\tkzFlameGrill}{%
  \fill[flameAndGrill] \tkzFlameGrillPath;
}

% ============================================================
% 10. 本生灯 - 带 grille
% ============================================================
\newcommand{\tkzBunsenBurnerGrill}{%
  \chemlabUpdateLineWidths\begin{scope}[line width=\chemlabLineWidth,line cap=round,line join=round]
    \draw (-1.25,0) rectangle (1.25,0.25);
    \draw (-.5,1.25) rectangle (0.5,2.25);
    \foreach \n in {-0.3,0.0,0.3} {\draw (\n,1.75) circle (0.1);}
    \draw (-.25,2.25) rectangle (0.25,4.25);
    \draw (0.25,1.25) -- (0.25,0.5) -- (1.25,0.25);
    \draw (-1.25,0.25) -- (-.25,0.5) -- (-0.25,0.75);
    \draw (-2.25,0.75) -- (-.25,0.75);
    \draw (-2.25,1) -- (-.25,1);
    \draw (-.25,0.875) ellipse (0.1 and 0.125);
    \fill[white] (-2.25,0.75) rectangle (-0.25,1);
    \draw (-2.25,0.75) -- (-0.25,0.75);
    \draw (-2.25,1) -- (-0.25,1) -- (-.25,1.25);
    \draw plot[smooth,tension=0.5] coordinates {(-0.25,0.5) (0,0.4) (0.25,0.5)};
    % flamme (clip 到 FlammeGrille)
    % 原 pst-labo 火焰分两层:
    %   外焰: \pst@FlammeGrille 整体填充 flameAndGrill (PaleOrange->yellow 渐变)
    %   内焰: 三角形 cyan->white 渐变 (在内焰区域)
    % 颜色和大小可通过 \chemlabset 调整
    \begin{scope}[shift={(0,4.25)},scale=\chemlabFlameScale]
      \begin{scope}
        \clip \tkzFlameGrillPath;
        % 外焰: PaleOrange (底) -> yellow (顶) 线性垂直渐变
        % 对应 pst-labo: fillstyle=gradient, gradbegin=PaleOrange, gradend=yellow
        % 渐变范围与火焰形状 (y=0..1.5) 一致, 避免底部颜色被稀释
        \shade[top color=\chemlabFlameOuterEndColor, bottom color=\chemlabFlameOuterColor]
          (-1.5,0) rectangle (1.5,1.5);
        % 内焰: cyan (底) -> white (顶) 线性垂直渐变 (三角形)
        \shade[top color=\chemlabFlameInnerEndColor, bottom color=\chemlabFlameInnerColor]
          (-0.25,0) -- (0.25,0) -- (0,1) -- cycle;
      \end{scope}
    \end{scope}
    \draw[dashed,line width=\chemlabEllipseLineWidth] (-2,5.75) -- (2,5.75);
  \end{scope}
}

% ============================================================
% 11. 烧瓶 (Ballon)
% ============================================================
% 路径版本 - 供 clip 使用 (闭合轮廓)
\newcommand{\tkzRoundFlaskSimplePath}{%
  ([shift={(-0.7,3.4)}]90:0.2) arc[start angle=90,end angle=0,radius=0.2]
  -- (-0.5,3.4) -- (-0.5,2.4)
  -- (-0.489,2.401)
  arc[start angle=113,end angle=427,radius=1.25]
  -- (0.5,2.4) -- (0.5,3.4)
  -- ([shift={(0.7,3.4)}]180:0.2) arc[start angle=180,end angle=90,radius=0.2]
  -- (0,3.5) -- cycle
}

\newcommand{\tkzRoundFlaskSimple}{%
  \chemlabUpdateLineWidths\begin{scope}[line width=\chemlabLineWidth,line cap=round,line join=round]
    % 左侧: 瓶口圆弧 + 颈部直线 + 球体大圆弧, 合并为单条路径以平滑过渡
    \draw ([shift={(-0.7,3.4)}]90:0.2) arc[start angle=90,end angle=0,radius=0.2]
          -- (-0.5,3.4) -- (-0.5,2.4)
          -- (-0.489,2.401) arc[start angle=113,end angle=427,radius=1.25];
    % 右侧: 球体终点(0.489,2.401) -> 颈部直线 -> 瓶口圆弧
    \draw (0.489,2.401) -- (0.5,2.4) -- (0.5,3.4)
          -- ([shift={(0.7,3.4)}]180:0.2) arc[start angle=180,end angle=90,radius=0.2];
    % 瓶口椭圆: 白色填充覆盖颈部上端圆头, 避免穿入椭圆
    % pst-labo: ballon椭圆用默认线宽 0.053 (与主轮廓一致, 不是 0.07)
    \filldraw[fill=white,draw=black,line width=\chemlabLineWidth] (0,3.6) ellipse (0.7 and 0.1);
  \end{scope}
}

\newcommand{\tkzRoundFlaskGrill}{%
  \chemlabUpdateLineWidths\begin{scope}[line width=\chemlabLineWidth,line cap=round,line join=round]
    % 左侧: 瓶口圆弧 + 颈部直线 + 球体大圆弧, 合并为单条路径
    \draw ([shift={(-0.7,2.15)}]90:0.2) arc[start angle=90,end angle=0,radius=0.2]
          -- (-0.5,2.15) -- (-0.5,1.4)
          -- (-0.489,1.401) arc[start angle=113,end angle=427,radius=1.25];
    % 右侧: 球体终点 -> 颈部直线 -> 瓶口圆弧
    \draw (0.489,1.401) -- (0.5,1.15) -- (0.5,3.4)
          -- ([shift={(0.7,2.15)}]180:0.2) arc[start angle=180,end angle=90,radius=0.2];
    % 瓶口椭圆: 白色填充覆盖颈部上端圆头, 避免穿入椭圆
    % pst-labo: ballon椭圆用默认线宽 0.053 (与主轮廓一致)
    \filldraw[fill=white,draw=black,line width=\chemlabLineWidth] (0,2.35) ellipse (0.7 and 0.1);
  \end{scope}
}

% ============================================================
% 12. 试管 (TubeEssais)
% ============================================================
\newcommand{\tkzTestTubeStopper}{%
  \chemlabUpdateLineWidths
  \begin{scope}[fill=LightGray,draw=black,line width=\chemlabLineWidth,line cap=round,line join=round]
    \filldraw (0,-0.3) ellipse (0.4 and 0.1);
    \filldraw (-0.4,-0.3) -- (-0.6,0.5) -- (0.6,0.5) -- (0.4,-0.3) -- cycle;
    \fill[draw=none] (0,-0.3) ellipse (0.4 and 0.1);
    \filldraw (0,0.5) ellipse (0.6 and 0.1);
  \end{scope}
  \draw[line width=\chemlabEllipseLineWidth] plot[smooth,tension=0.5]
        coordinates {(-0.5,0) (0,-0.07) (0.5,0)};
}

\newcommand{\tkzTestTubeBody}{%
  \chemlabUpdateLineWidths\begin{scope}[line width=\chemlabLineWidth,line cap=round,line join=round]
    % 合并左侧直线 + 底部圆弧 + 右侧直线为单条路径, 平滑转角
    \draw (-0.5,3.5) -- (-0.5,0.5)
          arc[start angle=180,end angle=360,radius=0.5]
          -- (0.5,3.5);
    % 试管口椭圆 (原 pst-labo 用 0.08cm, 比其他椭圆粗)
    \draw[line width=\chemlabTubeEllipseLineWidth] (0,3.5) ellipse (0.5 and 0.1);
  \end{scope}
}

\newcommand{\tkzTestTubeStraightTube}{%
  \draw[line width=0.5\pgflinewidth] (0.05,-1) -- (0.05,5);
  \draw[line width=0.5\pgflinewidth] (-0.05,-1) -- (-0.05,5);
  \tkzTestTubeStopper
  \draw[fill=white,rounded corners=0.5,line width=0.5\pgflinewidth]
       (-0.05,5) rectangle (0.05,0.5);
  \draw[white,line width=0.14cm] (-0.1,5) -- (0.1,5);
}

\newcommand{\tkzTestTubeElbowTube}{%
  \draw[line width=0.5\pgflinewidth] (0.05,-1) -- (0.05,0);
  \draw[line width=0.5\pgflinewidth] (-0.05,-1) -- (-0.05,0);
  \tkzTestTubeStopper
  % 弯管路径: pst-labo \pspolygon(...)(8,2)(8,1.9)...
  \filldraw[fill=white,rounded corners=0.05cm,line width=0.5\pgflinewidth]
       (-0.05,0.5) -- (-0.05,2) -- (8,2) -- (8,1.9)
       -- (0.05,1.9) -- (0.05,0.5) -- cycle;
  \draw[white,line width=0.14cm] (7.95,1.9) -- (8.05,2);
}

\newcommand{\tkzTestTubeUTube}{%
  \draw[line width=0.5\pgflinewidth] (0.05,-1) -- (0.05,0);
  \draw[line width=0.5\pgflinewidth] (-0.05,-1) -- (-0.05,0);
  \tkzTestTubeStopper
  \filldraw[fill=white,rounded corners=0.05cm,line width=0.5\pgflinewidth]
       (-0.05,0.5) -- (-0.05,2) -- (4.5,2) -- (4.5,-3)
       -- (4.4,-3) -- (4.4,1.9) -- (0.05,1.9) -- (0.05,0.5) -- cycle;
  \draw[white,line width=0.14cm] (4.3,-3) -- (4.6,-3);
}

\newcommand{\tkzTestTubeUTubeLong}{%
  \draw[line width=0.5\pgflinewidth] (0.05,-1) -- (0.05,0);
  \draw[line width=0.5\pgflinewidth] (-0.05,-1) -- (-0.05,0);
  \tkzTestTubeStopper
  \filldraw[fill=white,rounded corners=0.05cm,line width=0.5\pgflinewidth]
       (-0.05,0.5) -- (-0.05,2) -- (4.5,2) -- (4.5,-9)
       -- (4.4,-9) -- (4.4,1.9) -- (0.05,1.9) -- (0.05,0.5) -- cycle;
  \draw[white,line width=0.14cm] (4.3,-9) -- (4.6,-9);
}

\newcommand{\tkzCurvedTube}{%
  \draw[line width=0.5\pgflinewidth] (0.05,-1) -- (0.05,0);
  \draw[line width=0.5\pgflinewidth] (-0.05,-1) -- (-0.05,0);
  \tkzTestTubeStopper
  \begin{scope}[line width=0.4pt]
    \draw (-0.05,0.5) -- (-0.05,1);
    \draw ([shift={(0.5,1)}]180:0.55) arc[start angle=180,end angle=90,radius=0.55];
    \draw (0.5,1.55) -- (5.5,1.55);
    \draw ([shift={(5.5,1)}]90:0.55) arc[start angle=90,end angle=0,radius=0.55];
    \draw (6.05,1) -- (6.05,-8.8);
    \draw ([shift={(6.5,-8.8)}]180:0.45) arc[start angle=180,end angle=360,radius=0.45];
    \draw (6.95,-8.8) -- (6.95,-8.3);
    \draw (7.05,-8.3) -- (7.05,-8.8);
    \draw ([shift={(6.5,-8.8)}]0:0.55) arc[start angle=0,end angle=-180,radius=0.55];
    \draw (5.95,-8.8) -- (5.95,1);
    \draw ([shift={(5.5,1)}]0:0.45) arc[start angle=0,end angle=90,radius=0.45];
    \draw (5.5,1.45) -- (0.5,1.45);
    \draw ([shift={(0.5,1)}]90:0.45) arc[start angle=90,end angle=180,radius=0.45];
    \draw (0.05,1) -- (0.05,0.5);
    \draw (-0.05,0.5) -- (0.05,0.5);
  \end{scope}
  \fill[white] (-0.025,0.52) rectangle (0.020,0.6);
}

\newcommand{\tkzCurvedTubeShort}{%
  \draw[line width=0.5\pgflinewidth] (0.05,-1) -- (0.05,0);
  \draw[line width=0.5\pgflinewidth] (-0.05,-1) -- (-0.05,0);
  \tkzTestTubeStopper
  \draw (-0.05,0.5) -- (-0.05,1);
  \begin{scope}[line width=0.4pt]
    \draw ([shift={(0.5,1)}]180:0.55) arc[start angle=180,end angle=90,radius=0.55];
    \draw (0.5,1.55) -- (3.5,1.55);
    \draw ([shift={(3.5,1)}]90:0.55) arc[start angle=90,end angle=0,radius=0.55];
    \draw (4.05,1) -- (4.05,-2.8);
    \draw ([shift={(4.5,-2.8)}]180:0.45) arc[start angle=180,end angle=360,radius=0.45];
    \draw (4.95,-2.8) -- (4.95,-2.3);
    \draw (5.05,-2.3) -- (5.05,-2.8);
    \draw ([shift={(4.5,-2.8)}]0:0.55) arc[start angle=0,end angle=-180,radius=0.55];
    \draw (3.95,-2.8) -- (3.95,1);
    \draw ([shift={(3.5,1)}]0:0.45) arc[start angle=0,end angle=90,radius=0.45];
    \draw (3.5,1.45) -- (0.5,1.45);
    \draw ([shift={(0.5,1)}]90:0.45) arc[start angle=90,end angle=180,radius=0.45];
    \draw (0.05,1) -- (0.05,0.5);
    \draw (-0.05,0.5) -- (0.05,0.5);
  \end{scope}
  \fill[white] (-0.025,0.52) rectangle (0.020,0.6);
}

\newcommand{\tkzCurvedTubeShortDoubleTube}{%
  \begin{scope}[line width=0.4pt]
    \draw (-0.05,0.5) -- (-0.05,1);
    \draw ([shift={(0.5,1)}]180:0.55) arc[start angle=180,end angle=90,radius=0.55];
    \draw (0.5,1.55) -- (3.5,1.55);
    \draw ([shift={(3.5,1)}]90:0.55) arc[start angle=90,end angle=0,radius=0.55];
    \draw (4.05,1) -- (4.05,-2.8);
    \draw ([shift={(4.5,-2.8)}]180:0.45) arc[start angle=180,end angle=360,radius=0.45];
    \draw (4.95,-2.8) -- (4.95,-2.3);
    \draw (5.05,-2.3) -- (5.05,-2.8);
    \draw ([shift={(4.5,-2.8)}]0:0.55) arc[start angle=0,end angle=-180,radius=0.55];
    \draw (3.95,-2.8) -- (3.95,1);
    \draw ([shift={(3.5,1)}]0:0.45) arc[start angle=0,end angle=90,radius=0.45];
    \draw (3.5,1.45) -- (0.5,1.45);
    \draw ([shift={(0.5,1)}]90:0.45) arc[start angle=90,end angle=180,radius=0.45];
    \draw (0.05,1) -- (0.05,0.5);
    \draw (-0.05,0.5) -- (0.05,0.5);
  \end{scope}
  \fill[white] (-0.025,0.52) rectangle (0.018,0.8);
}

\newcommand{\tkzTestTubeDoubleTube}{%
  \begin{scope}[line width=0.4pt]
    \draw (-0.3,-3) -- (-0.3,0);
    \draw (-.2,-3) -- (-0.2,0);
    \draw (0.25,-1) -- (0.25,0);
    \draw (.15,-1) -- (.15,0);
  \end{scope}
  \tkzTestTubeStopper
  \begin{scope}[shift={(0.2,0)}]\tkzCurvedTubeShortDoubleTube\end{scope}
  \begin{scope}[shift={(0,0.5)}]
    \begin{scope}[line width=0.4pt]
      \filldraw[fill=white,rounded corners=0.05cm]
        (-.3,0) -- (-0.3,2) -- (-0.5,2.2) -- (-0.5,3) -- (0,3)
                -- (0,2.2) -- (-0.2,2) -- (-0.2,0) -- cycle;
      \draw (0.1,1.55) -- (-.5,1.55) -- (-0.5,1.7) -- (-.55,1.7)
            -- (-.55,1.3) -- (-0.5,1.3) -- (-0.5,1.55) -- (-.5,1.45) -- (.1,1.45);
      \filldraw[fill=white,rounded corners=0.2] (-0.4,1.3) rectangle (-0.1,1.7);
    \end{scope}
    \fill[white] (-0.25,3) ellipse (0.3 and 0.1);
  \end{scope}
}

% 试管夹 (Pince) - 几何与原 pst-labo pst-laboObj.tex line 483-508 一致
% 根据原 pst-labo 实际渲染 (5.3x 截图) 校准:
%   - 上边缘凹弧较浅 (原 \pscurve 3点几乎共线, 浅弧)
%   - 顶部小凸起为圆弧形 (\psarcn 圆弧)
%   - 所有转角为圆角 (原 \pscustom 默认 linejoin=0 即 round)
\newcommand{\tkzTestTubeClamp}{%
  % 第1块: 上部夹爪
  % 上边缘(1a, 夹合面): 下凹弧贴合试管顶部弧面
  %   试管半径 0.5, 在 y=3.2 处试管顶部弧面下凹约 0.1, 用贝塞尔曲线模拟
  %   从 (0.5,0) 经 (0,-0.1) 到 (-0.5,0), 控制点 y=-0.13 产生平滑下凹
  \filldraw[fill=PaleOrange,draw=black,line cap=round,line join=round]
    (0.5,0) .. controls (0.25,-0.13) and (-0.25,-0.13) .. (-0.5,0)
    -- (-0.8,0)
    -- ([shift={(-0.5,0)}]180:0.3) arc[start angle=180,end angle=270,radius=0.3]
    -- (-.5,-.3) -- (2.5,-.3) -- (2.4,-.2) -- (1,0) -- (0.5,0) -- cycle;
  % 第2块: 下部夹爪
  % 注: 原 \psarcn(-.5,0){0.3}{-90}{180} 是顺时针从 -90°(=270°) 到 180°, 跨度 90° (四分之一圆弧)
  % TikZ arc 当 end<start 时顺时针, 需用 end=-180 让其走短弧 (避免四分之三圆弧)
  \filldraw[fill=PaleOrange,draw=black,line cap=round,line join=round]
    (-0.8,0) -- (-0.8,-0.3)
    -- ([shift={(-0.5,-0.3)}]180:0.3) arc[start angle=180,end angle=270,radius=0.3]
    -- (-0.5,-0.6) -- (2.5,-.6) -- (2.5,-.3) -- (-.5,-.3)
    -- ([shift={(-.5,0)}]-90:0.3) arc[start angle=-90,end angle=-180,radius=0.3] -- cycle;
  % 第3块: 顶部小凸起 - 原 \psarcn 顺时针 180°->90°, 圆弧形凸起
  % 去掉 (-0.8,0.05)->(-0.75,0.05) 的多余水平线段, 直接从 (-0.5,0.05) 到圆弧起点
  % 使凸起呈纯圆弧形 (与原 pst-labo 渲染一致)
  \filldraw[fill=PaleOrange,draw=black,line cap=round,line join=round]
    (-.5,0.05) -- ([shift={(-0.5,0.05)}]180:0.25)
    arc[start angle=180,end angle=90,radius=0.25]
    -- (-.5,0.3) -- (-0.5,0.05) -- cycle;
  % 第4块: 夹臂上部
  \filldraw[fill=PaleOrange,draw=black,line cap=round,line join=round]
    (0.5,0.3) -- (4.9,0.3) -- (5,0.2) -- (1,0) -- (0.5,0) -- (0.5,0.3) -- cycle;
  % 第5块: 夹臂下部 (Brown 棕色)
  \filldraw[fill=Brown,draw=black,line cap=round,line join=round]
    (0.5,0) -- (5,0.2) -- (5,0) -- (2,-0.2) -- (1,0) -- cycle;
  % 螺丝椭圆
  \draw[line width=1.5\pgflinewidth,line cap=round] (1,0) ellipse (0.3 and 0.1);
  % 螺丝下方的杆
  \draw[line width=1.5\pgflinewidth,line cap=round,line join=round]
        (0.7,0) -- (0.4,-.3) -- (0.4,-.6);
}

% ============================================================
% 13. 容量瓶 (Fiole jaugée)
% ============================================================
\newcommand{\tkzVolumetricFlask}{%
  \chemlabUpdateLineWidths
  \draw[rounded corners=0.1cm,line width=\chemlabLineWidth,line cap=round,line join=round]
       (-0.21,3.5) -- (-0.1,3.5) -- (-0.1,1.5)
       -- (-0.75,0) -- (0.75,0) -- (0.1,1.5) -- (0.1,3.5) -- (0.21,3.5);
  \draw[line width=\chemlabGraduationLineWidth,line cap=round] (-0.1,2.5) -- (0.1,2.5);
}

% ============================================================
% 14. 结晶皿 (Cristallisoir)
% ============================================================
\newcommand{\Cristallisoir}{%
  \chemlabUpdateLineWidths
  \begin{scope}[line width=\chemlabLineWidth,line cap=round,line join=round]
  \draw (-1.9,2) -- (-2,2);
  \draw ([shift={(-2,1.95)}]90:0.05) arc[start angle=90,end angle=270,radius=0.05];
  \draw (-2,1.9) -- (-2,1);
  \draw ([shift={(-1,1)}]180:1) arc[start angle=180,end angle=270,radius=1];
  \draw (-1,0) -- (1,0);
  \draw ([shift={(1,1)}]270:1) arc[start angle=270,end angle=360,radius=1];
  \draw (2,1) -- (2,1.9);
  \draw ([shift={(2,1.95)}]270:0.05) arc[start angle=270,end angle=90,radius=0.05];
  \draw (2,2) -- (1.9,2);
  \draw (1.9,2) -- (1.9,1);
  \draw ([shift={(1,1)}]0:0.9) arc[start angle=0,end angle=-90,radius=0.9];
  \draw (1,0.1) -- (-1,0.1);
  \draw ([shift={(-1,1)}]270:0.9) arc[start angle=270,end angle=180,radius=0.9];
  \draw (-1.9,1) -- (-1.9,2);
  \end{scope}
}

\newcommand{\InterieurCristallisoir}{%
  (1.9,2) -- (1.9,1)
  -- ([shift={(1,1)}]0:0.9) arc[start angle=0,end angle=-90,radius=0.9]
  -- (1,0.1) -- (-1,0.1)
  -- ([shift={(-1,1)}]270:0.9) arc[start angle=270,end angle=180,radius=0.9]
  -- (-1.9,1) -- (-1.9,2) -- cycle
}

\newcommand{\tkzCrystallizingDish}{%
  \begin{scope}
    \clip \InterieurCristallisoir;
    \fill[LightBlue1] (-3,-1) rectangle (3,1.5);
  \end{scope}
  \Cristallisoir
}

% ============================================================
% 15. 倒置试管 (TubeRenverse)
% ============================================================
\newcommand{\tkzInvertedTube}{%
  \chemlabUpdateLineWidths
  \begin{scope}[line width=\chemlabEllipseLineWidth,line cap=round,line join=round]
    \fill[LightBlue1] (-0.5,0) rectangle (0.5,5.5);
    % 合并左侧 + 顶部圆弧 + 右侧为单条路径
    % pst-labo: \psarcn(0,6){0.5}{180}{0} = 从180°顺时针到0°, 经过顶部90°
    % TikZ: arc[start=180, end=0] = 从180°逆时针(减小)到0°, 经过顶部90°
    \draw (-0.5,0) -- (-0.5,6)
          arc[start angle=180,end angle=0,radius=0.5]
          -- (0.5,0);
  \end{scope}
}

% ============================================================
% 16. 试管架 (ReposeTube)
% ============================================================
\newcommand{\tkzTestTubeRack}{%
  \begin{scope}[line width=0.1cm,draw=lightgray]
    \draw ([shift={(0,-0.25)}]100:1) arc[start angle=100,end angle=140,radius=1];
    \draw ([shift={(0,-.25)}]150:1) arc[start angle=150,end angle=160,radius=1];
    \draw ([shift={(0,-0.25)}]20:1) arc[start angle=20,end angle=80,radius=1];
  \end{scope}
}

% ============================================================
% 17. 烧杯边缘 + 刻度 (RebordBecher)
% ============================================================
\newcommand{\tkzBeakerRim}{%
  \chemlabUpdateLineWidths
  \draw[line width=\chemlabLineWidth] (0,2.55) ellipse (1.1 and .1);
  % 主刻度 : 粗细为细刻度线的 1.5 倍 (原 pst-labo: 0.03cm = 1.5 × 0.02cm)
  % 用 \pgfmathsetmacro 预算尺寸, 避免 \dimexpr 在 pgfkeys 中的兼容问题
  \pgfmathsetlengthmacro{\chemlab@majorGrad}{1.5*\chemlabGraduationLineWidth}%
  \foreach \nDiv in {0.5,1.0,...,2.0} {%
    \draw[line width=\chemlab@majorGrad] (-.75,\nDiv) -- (-.4,\nDiv);}
  % 细分刻度
  \foreach \nSubDiv in {0.5,0.6,...,2.0} {%
    \draw[line width=\chemlabGraduationLineWidth] (-0.75,\nSubDiv) -- (-.5,\nSubDiv);}
}

% ============================================================
% 18. 烧杯 (Becher)
% ============================================================
% 注 : 原 PSTricks 使用 \psline[linearc=0.5](-1,2.5)(-1,0)(1,0)(1,2.5)
% TikZ 的 rounded corners 与 PSTricks linearc 语义略有不同,
% 因此改用显式圆弧精确模拟 PSTricks linearc=0.5 的行为 :
% 左下角圆弧 : 圆心 (-0.5,0.5), 半径 0.5, 从 180° 到 270°
% 右下角圆弧 : 圆心 (0.5,0.5),  半径 0.5, 从 270° 到 360°
\newcommand{\tkzBeakerBody}{%
  \chemlabUpdateLineWidths
  \draw[line width=\chemlabLineWidth,line cap=round,line join=round]
       (-1,2.5) -- (-1,0.5)
       -- ([shift={(-0.5,0.5)}]180:0.5) arc[start angle=180,end angle=270,radius=0.5]
       -- (0.5,0)
       -- ([shift={(0.5,0.5)}]270:0.5) arc[start angle=270,end angle=360,radius=0.5]
       -- (1,2.5);
}

% ============================================================
% 19. 锥形瓶 (Erlenmeyer)
% ============================================================
% 注 : 原 PSTricks 使用 \psline[linearc=0.3](-0.5,3.5)(-0.5,2.5)(-1.5,0)(1.5,0)(0.5,2.5)(0.5,3.5)
% TikZ 的 rounded corners 与 PSTricks linearc 语义不同,
% 因此改用显式圆弧精确模拟 linearc=0.3 (由 scripts/erlen_arcs5.py 计算) :
%   肩部 (凹转角, 内侧角 21.8°) : 半径 0.3 太大无法相切 (切点距 1.558 > 颈部线段长 1.0), 保留尖角
%   左下 (凸转角, 内侧角 68.2°) : center=(-1.0569,0.3), r=0.3, 158.20° -> 270° (ccw)
%   右下 (凸转角, 内侧角 68.2°) : center=(1.0569,0.3),  r=0.3, -90° -> 21.80° (ccw)
\newcommand{\tkzErlenmeyerBody}{%
  \chemlabUpdateLineWidths
  \draw[line width=\chemlabLineWidth,line cap=round,line join=round]
       (-0.5,3.5) -- (-0.5,2.5)
       -- (-1.3354,0.4114)
       -- ([shift={(-1.0569,0.3)}]158.20:0.3) arc[start angle=158.20,end angle=270,radius=0.3]
       -- (1.0569,0)
       -- ([shift={(1.0569,0.3)}]-90:0.3) arc[start angle=-90,end angle=21.80,radius=0.3]
       -- (0.5,2.5) -- (0.5,3.5);
}

% ============================================================
% 20. 试剂瓶 (Flacon)
% ============================================================
\newcommand{\tkzReagentBottleBody}{%
  \chemlabUpdateLineWidths\begin{scope}[line width=\chemlabLineWidth,line cap=round,line join=round]
    % 左侧瓶口上角小弧 (psarcn 90->0, 顺时针)
    \draw ([shift={(-0.6,3.4)}]90:0.1) arc[start angle=90,end angle=0,radius=0.1];
    % 左侧颈部直线
    \draw (-0.5,3.4) -- (-0.5,3.1);
    % 左侧瓶口下角凹弧 (psarcn 0->-90, 顺时针; TikZ end<start 自动顺时针)
    \draw ([shift={(-0.6,3.1)}]0:0.1) arc[start angle=0,end angle=-90,radius=0.1];
    % 左侧肩部大凸弧 (psarc 90->180, 逆时针)
    \draw ([shift={(-0.6,2.6)}]90:0.4) arc[start angle=90,end angle=180,radius=0.4];
    % 左侧瓶身直线
    \draw (-1,2.6) -- (-1,0.2);
    % 左侧底部圆角 (psarc 180->270, 逆时针)
    \draw ([shift={(-.8,0.2)}]180:0.2) arc[start angle=180,end angle=270,radius=0.2];
    % 底部
    \draw (-.8,0) -- (.8,0);
    % 右侧底部圆角 (psarc -90->0, 逆时针)
    \draw ([shift={(.8,0.2)}]-90:0.2) arc[start angle=-90,end angle=0,radius=0.2];
    % 右侧瓶身直线
    \draw (1,0.2) -- (1,2.6);
    % 右侧肩部大凸弧 (psarc 0->90, 逆时针)
    \draw ([shift={(0.6,2.6)}]0:0.4) arc[start angle=0,end angle=90,radius=0.4];
    % 右侧瓶口下角凹弧 (psarcn -90->180, 顺时针短弧; TikZ 需 end=-180 避免 270° 长弧)
    \draw ([shift={(0.6,3.1)}]-90:0.1) arc[start angle=-90,end angle=-180,radius=0.1];
    % 右侧颈部直线
    \draw (0.5,3.1) -- (0.5,3.4);
    % 右侧瓶口上角小弧 (psarcn 180->90, 顺时针)
    \draw ([shift={(0.6,3.4)}]180:0.1) arc[start angle=180,end angle=90,radius=0.1];
  \end{scope}
}

% ============================================================
% 21. 玻璃杯 (Verre) - F. Bréal
% ============================================================
\newcommand{\tkzWatchGlassBody}{%
  \chemlabUpdateLineWidths\begin{scope}[line width=\chemlabLineWidth,line cap=round,line join=round]
    % 合并左斜线 + 底部圆弧 + 右斜线为单条路径, 平滑转角
    \draw (-1,3.5) -- (-0.5,1)
          arc[start angle=180,end angle=360,radius=0.5]
          -- (1,3.5);
  \end{scope}
}

\newcommand{\tkzWatchGlassBase}{%
  \chemlabUpdateLineWidths\begin{scope}[line width=\chemlabLineWidth,rounded corners=0.03cm,line cap=round,line join=round]
    \draw (-0.5,1) -- (-0.34,0.2) -- (-1.2,0) -- (1.2,0) -- (0.34,0.2) -- (0.5,1);
  \end{scope}
}

% ============================================================
% 22. 漏斗 (Entonnoir)
% ============================================================
\newcommand{\tkzFunnelOutline}{%
  \begin{scope}[line cap=round,line join=round]
    \draw (-0.07,0) -- (-0.07,1) -- (-1,2.5);
    \draw (1,2.5) -- (0.07,1) -- (0.07,0.1);
  \end{scope}
}

\newcommand{\tkzFunnelBody}{%
  \begin{scope}
    \clip (-0.07,0) -- (-0.07,1) -- (-1,2.5)
          -- (1,2.5) -- (0.07,1) -- (0.07,0.1) -- cycle;
    \fill[white] (-2,-1) rectangle (2,3);
  \end{scope}
  \begin{scope}[rounded corners=0.25cm]
    \draw (-0.07,0) -- (-0.07,1) -- (-1,2.5);
    \draw (1,2.5) -- (0.07,1) -- (0.07,0.1);
  \end{scope}
  \draw (0,2.5) ellipse (1 and 0.1);
}

% ============================================================
% 23. 球形冷凝器 (Réfrigérant à boulles)
% ============================================================
\newcommand{\tkzCondenserStem}{%
  % 整个 TigeBoulles 作为一个连续路径 (匹配 pst-labo \pscustom)
  % pst-labo: \pscustom[linestyle=none]{...\fill} 填充 + 单独描边 (无闭合线)
  % TikZ: 分开 \fill (带 cycle 闭合填充) 和 \draw (无 cycle, 顶部敞口)
  % 顶部是敞口的 — 两端弧之间没有连线 (匹配 pst-labo)
  %
  % 填充: 带闭合线 (fill 需要 cycle)
  \fill[white]
    % 顶部圆弧与颈部 - 左侧
    ({-0.35+0.1*cos(90)},{7.5+0.1*sin(90)}) arc[start angle=90,end angle=0,radius=0.1]
    -- (-.25,7.5) -- (-.25,6.6)
    -- ({-.15+0.1*cos(180)},{6.6+0.1*sin(180)}) arc[start angle=180,end angle=270,radius=0.1]
    -- ({-.15+0.1*cos(90)},{6.4+0.1*sin(90)}) arc[start angle=90,end angle=0,radius=0.1]
    -- (-.05,6.4) -- (-0.05,6)
    % 左侧 7 个球状突起 (中心 x=-0.05, 半径 0.25, arc 90→270 向左凸)
    (-0.05,6.0) arc[start angle=90,end angle=270,radius=0.25] -- (-0.05,5.50) -- (-0.05,5.25)
    arc[start angle=90,end angle=270,radius=0.25] -- (-0.05,4.75) -- (-0.05,4.50)
    arc[start angle=90,end angle=270,radius=0.25] -- (-0.05,4.00) -- (-0.05,3.75)
    arc[start angle=90,end angle=270,radius=0.25] -- (-0.05,3.25) -- (-0.05,3.00)
    arc[start angle=90,end angle=270,radius=0.25] -- (-0.05,2.50) -- (-0.05,2.25)
    arc[start angle=90,end angle=270,radius=0.25] -- (-0.05,1.75) -- (-0.05,1.50)
    arc[start angle=90,end angle=270,radius=0.25] -- (-0.05,1.00) -- (-0.05,0.75)
    % 左侧底部到下管口
    -- (-0.05,-1.5) -- (0.05,-1.5) -- (0.05,1)
    % 右侧 7 个球状突起 (中心 x=0.05, 半径 0.25, arc 270→450 向右凸)
    arc[start angle=270,end angle=450,radius=0.25] -- (0.05,1.50) -- (0.05,1.75)
    arc[start angle=270,end angle=450,radius=0.25] -- (0.05,2.25) -- (0.05,2.50)
    arc[start angle=270,end angle=450,radius=0.25] -- (0.05,3.00) -- (0.05,3.25)
    arc[start angle=270,end angle=450,radius=0.25] -- (0.05,3.75) -- (0.05,4.00)
    arc[start angle=270,end angle=450,radius=0.25] -- (0.05,4.50) -- (0.05,4.75)
    arc[start angle=270,end angle=450,radius=0.25] -- (0.05,5.25) -- (0.05,5.50)
    arc[start angle=270,end angle=450,radius=0.25] -- (0.05,6.00) -- (0.05,6.40)
    % 顶部圆弧与颈部 - 右侧
    -- ({0.15+0.1*cos(180)},{6.4+0.1*sin(180)}) arc[start angle=180,end angle=90,radius=0.1]
    -- ({0.15+0.1*cos(-90)},{6.6+0.1*sin(-90)}) arc[start angle=-90,end angle=0,radius=0.1]
    -- (0.25,6.6) -- (0.25,7.5)
    -- ({0.35+0.1*cos(180)},{7.5+0.1*sin(180)}) arc[start angle=180,end angle=90,radius=0.1]
    -- cycle;
  %
  % 描边: 底部敞口 (匹配 pst-labo: 左右侧分别 \psline, 无底部连线)
  \draw
    % 顶部圆弧与颈部 - 左侧
    ({-0.35+0.1*cos(90)},{7.5+0.1*sin(90)}) arc[start angle=90,end angle=0,radius=0.1]
    -- (-.25,7.5) -- (-.25,6.6)
    -- ({-.15+0.1*cos(180)},{6.6+0.1*sin(180)}) arc[start angle=180,end angle=270,radius=0.1]
    -- ({-.15+0.1*cos(90)},{6.4+0.1*sin(90)}) arc[start angle=90,end angle=0,radius=0.1]
    -- (-.05,6.4) -- (-0.05,6)
    % 左侧 7 个球状突起
    (-0.05,6.0) arc[start angle=90,end angle=270,radius=0.25] -- (-0.05,5.50) -- (-0.05,5.25)
    arc[start angle=90,end angle=270,radius=0.25] -- (-0.05,4.75) -- (-0.05,4.50)
    arc[start angle=90,end angle=270,radius=0.25] -- (-0.05,4.00) -- (-0.05,3.75)
    arc[start angle=90,end angle=270,radius=0.25] -- (-0.05,3.25) -- (-0.05,3.00)
    arc[start angle=90,end angle=270,radius=0.25] -- (-0.05,2.50) -- (-0.05,2.25)
    arc[start angle=90,end angle=270,radius=0.25] -- (-0.05,1.75) -- (-0.05,1.50)
    arc[start angle=90,end angle=270,radius=0.25] -- (-0.05,1.00) -- (-0.05,0.75)
    -- (-0.05,-1.5);
  % 右侧 (单独 \draw, 底部不连接左侧)
  \draw
    (0.05,-1.5) -- (0.05,1)
    % 右侧 7 个球状突起
    arc[start angle=270,end angle=450,radius=0.25] -- (0.05,1.50) -- (0.05,1.75)
    arc[start angle=270,end angle=450,radius=0.25] -- (0.05,2.25) -- (0.05,2.50)
    arc[start angle=270,end angle=450,radius=0.25] -- (0.05,3.00) -- (0.05,3.25)
    arc[start angle=270,end angle=450,radius=0.25] -- (0.05,3.75) -- (0.05,4.00)
    arc[start angle=270,end angle=450,radius=0.25] -- (0.05,4.50) -- (0.05,4.75)
    arc[start angle=270,end angle=450,radius=0.25] -- (0.05,5.25) -- (0.05,5.50)
    arc[start angle=270,end angle=450,radius=0.25] -- (0.05,6.00) -- (0.05,6.40)
    -- ({0.15+0.1*cos(180)},{6.4+0.1*sin(180)}) arc[start angle=180,end angle=90,radius=0.1]
    -- ({0.15+0.1*cos(-90)},{6.6+0.1*sin(-90)}) arc[start angle=-90,end angle=0,radius=0.1]
    -- (0.25,6.6) -- (0.25,7.5)
    -- ({0.35+0.1*cos(180)},{7.5+0.1*sin(180)}) arc[start angle=180,end angle=90,radius=0.1];
  % 底部液滴
  \begin{scope}[shift={(0,-1.6)}]\tkzDrop\end{scope}
}

\newcommand{\tkzCondenserBody}{%
  % 渐变填充: pst-labo gradbegin=cyan (底), gradend=PaleOrange (顶), gradmidpoint=0
  % TikZ: bottom color=cyan (底), top color=PaleOrange (顶)
  \shade[top color=PaleOrange,bottom color=cyan]
    ([shift={(0,5.75)}]0:0.5) arc[start angle=0,end angle=180,radius=0.5]
    -- (-0.5,5.75) -- (-0.5,5.5) -- (-1.5,5.5)
    -- (-1.5,5.3) -- (-.5,5.3) -- (-.5,1.7) -- (-1.5,1.7)
    -- (-1.5,1.5) -- (-.5,1.5) -- (-0.5,1)
    -- ([shift={(0,1)}]180:0.5) arc[start angle=180,end angle=360,radius=0.5]
    -- (0.5,1) -- (0.5,5.75) -- cycle;
  % 描边
  \begin{scope}[line width=0.88pt]
    \draw ([shift={(0,5.75)}]0:0.5) arc[start angle=0,end angle=180,radius=0.5];
    \draw (-0.5,5.75) -- (-0.5,5.5) -- (-1.5,5.5);
    \draw (-0.5,5.75) -- (-0.5,5.5) -- (-1.5,5.5);
    \draw (-1.5,5.3) -- (-.5,5.3) -- (-.5,1.7) -- (-1.5,1.7);
    \draw (-1.5,1.5) -- (-.5,1.5) -- (-0.5,1);
    \draw ([shift={(0,1)}]180:0.5) arc[start angle=180,end angle=360,radius=0.5];
    \draw (0.5,1) -- (0.5,5.75);
  \end{scope}
  \draw[->,cyan] (-2.5,1.6) -- (-1.6,1.6);
  \draw[<-,PaleOrange] (-2.5,5.4) -- (-1.6,5.4);
}

\newcommand{\tkzCondenserBulbs}{%
  \begin{scope}[shift={(0,0.2)}]
    \tkzCondenserBody
    \tkzCondenserStem
  \end{scope}
  \tkzTestTubeStopper
  \fill[white,rounded corners=0.025cm]
       (-0.05,0.6) -- (-0.05,0.5) -- (0.05,0.5) -- (0.05,0.6) -- cycle;
}

% ============================================================
% 24. 加热球 (Chauffe-Ballon)
% ============================================================
\newcommand{\tkzHeaterRoundFlask}{%
  \filldraw[fill=white,draw=black] (-2,-1) rectangle (2,1);
  \draw (-1.25,0) circle (0.4);
  \begin{scope}[shift={(-1.25,0)}]
    \foreach \iA in {0,30,...,180} {\draw (\iA:0.45) -- (\iA:0.6);}
    \draw[->] (0,0) -- (60:0.3);
  \end{scope}
  % 按钮区域: pst-labo \psframe[doubleline=true] 产生双线边框
  % PSTricks doubleline: 外框在指定坐标, 内框向内偏移
  % 实测: 外框到内框偏移约 0.13 units (含线宽+间隔, at scale=0.5)
  % 外框 (在指定坐标处)
  \draw[black,line width=0.8pt] (1,-.5) rectangle (1.5,0.5);
  % 内框 (向内偏移, 填充灰色)
  \filldraw[fill=gray,draw=black,line width=0.8pt] ({1+0.14},{-.5+0.14}) rectangle ({1.5-0.14},{0.5-0.14});
  % 红色指示灯
  \fill[red] (1.25,0.7) circle (0.1);
}

% ============================================================
% 25. 回流烧瓶 (Ballon à reflux)
% ============================================================
\newcommand{\tkzRoundFlaskReflux}{%
  \begin{scope}
    \clip \tkzRoundFlaskSimplePath;   % 使用路径版本形成闭合 clip 区域
    \shade[top color=cyan,bottom color=white,draw=cyan]
      plot[smooth,tension=0.5] coordinates {%
        (-1.5,1.5) (-1.375,1.55) (-1.25,1.5) (-1.125,1.45)
        (-1,1.5) (-.875,1.55) (-.75,1.5) (-.625,1.45)
        (-.5,1.5) (-.375,1.55) (-.25,1.5) (-0.125,1.45)
        (0,1.5) (0.125,1.55) (.25,1.5) (0.375,1.45)
        (0.5,1.5) (0.625,1.55) (0.75,1.5) (0.875,1.45)
        (1,1.5) (1.125,1.55) (1.25,1.5) (1.375,1.45) (1.5,1.5)}
      -- (1.5,1.5) -- (1.5,0) -- (-1.5,0) -- (-1.5,1.5) -- cycle;
  \end{scope}
  \tkzRoundFlaskSimple
}

% ============================================================
% 26. 分液漏斗 (Ampoule à décanter / Separate Funnel)
% ============================================================
% 相 : 水相 / 有机相
\definecolor{AqueousPhaseColor}{cmyk}{0.215686,0.062745,0,0}
\definecolor{OrganicPhaseColor}{cmyk}{0.223529,0.094118,0.501961,0}

% pst-labo: \tkzSeparatoryFunnel 设 linewidth=0.04 (用户单位, 随 unit 缩放)
% TikZ: 用 \chemlabLineWidth (基础 0.053cm × scale) 近似, 匹配其他器皿线宽
\newcommand{\tkzSeparatoryFunnelBody}{%
  \chemlabUpdateLineWidths
  \draw[line width=\chemlabLineWidth,line cap=round,line join=round]
    (-0.65,15) -- (-0.55,13.75)
    .. controls (-1.55,13.54) and ( -1.92,12.7) .. (-1.92,11.65)
    .. controls (-1.92,10.6) and (-0.3,9.55) .. (-0.3,8.5)
    -- (-0.3,2)
    .. controls (-0.3,1.5) and (-0.14,1.5) .. (-0.14,1)
    -- (-0.14,-1)
    -- (0.14,-1)
    -- (0.14,1)
    .. controls (0.14,1.5) and (0.3,1.5) .. (0.3,2)
    -- (0.3,8.5)
    .. controls (0.3,9.55) and (1.92,10.6) .. (1.92,11.65)
    .. controls ( 1.92,12.7) and (1.55,13.54) .. (0.55,13.75)
    -- (0.65,15);
}

% 用于 clip 的轮廓
\newcommand{\tkzSeparatoryFunnelOutlineClip}{%
  (-0.65,15) -- (-0.55,13.75)
  .. controls (-1.55,13.54) and ( -1.92,12.7) .. (-1.92,11.65)
  .. controls (-1.92,10.6) and (-0.3,9.55) .. (-0.3,8.5)
  -- (-0.3,2)
  .. controls (-0.3,1.5) and (-0.14,1.5) .. (-0.14,1)
  -- (-0.14,-1) -- (0.14,-1) -- (0.14,1)
  .. controls (0.14,1.5) and (0.3,1.5) .. (0.3,2)
  -- (0.3,8.5)
  .. controls (0.3,9.55) and (1.92,10.6) .. (1.92,11.65)
  .. controls (1.92,12.7) and (1.55,13.54) .. (0.55,13.75)
  -- (0.65,15) -- cycle
}

% 顶盖 (Bouchon) - 继承 linewidth=0.04
\newcommand{\tkzStopper}{%
  \chemlabUpdateLineWidths
  \filldraw[fill=gray!10,draw=gray,line width=\chemlabLineWidth,line cap=round,line join=round]
    (-0.67,15.25) -- (0.67,15.25) -- (0.57,14) -- (-0.57,14) -- cycle;
  \filldraw[fill=gray!5,draw=gray,line width=\chemlabLineWidth,line cap=round,line join=round]
    ([shift={(0.555,15.365)}]-90:0.1145) arc[start angle=-90,end angle=-180,radius=0.1145]
    -- ([shift={(0,15.92)}]-45:0.67) arc[start angle=-45,end angle=225,radius=0.67]
    -- ([shift={(-0.555,15.365)}]0:0.1145) arc[start angle=0,end angle=-90,radius=0.1145]
    -- cycle;
}

% 开启的旋塞 (TapOpen)
% pst-labo: \psframe[fillstyle=solid] 默认 fillcolor=white (PSTricks 默认)
% TikZ: \fill 默认黑色, 必须显式 fill=white
\newcommand{\tkzTapOpen}{%
  \chemlabUpdateLineWidths
  \filldraw[fill=white,line width=\chemlabLineWidth] (-.2,0.5) rectangle (.2,1);
  \draw[line width=\chemlabLineWidth] (-.2,0.4) rectangle (-.3,1.1);
  \draw[line width=\chemlabLineWidth] (.2,0.4) rectangle (.3,1.1);
  \draw[line width=\chemlabLineWidth] (.3,0.6) -- (.4,0.6) -- (.5,0.3) -- (.9,0.3) -- (.9,1.2) -- (.5,1.2)
        -- (.4,0.9) -- (.3,0.9) -- cycle;
  \fill[cyan!50] (-.6,0.4) -- (-.85,0.6) -- (-.85,0.9) -- (-.6,1.1) -- (-.3,1.1) -- (-.3,0.4) -- cycle;
  \draw[line width=\chemlabLineWidth] (-0.6,0.4) -- (-0.6,1.1);
  \draw[line width=\chemlabLineWidth] (0.5,0.3) -- (0.5,1.2);
}

% 关闭的旋塞 (TapClosed)
% pst-labo: \psframe[fillstyle=solid] 默认 fillcolor=white
\newcommand{\tkzTapClosed}{%
  \chemlabUpdateLineWidths
  \filldraw[fill=white,line width=\chemlabLineWidth] (-.2,0.5) rectangle (.2,1);
  \draw[line width=\chemlabLineWidth] (-.2,0.4) rectangle (-.3,1.1);
  \draw[line width=\chemlabLineWidth] (.2,0.4) rectangle (.3,1.1);
  \draw[line width=\chemlabLineWidth] ([shift={(0.35,0.6)}]180:0.05) arc[start angle=180,end angle=90,radius=0.05]
        -- (0.4,0.65) -- (0.9,0.65) -- (0.9,0.85) -- (0.35,0.85)
        -- ([shift={(0.35,0.9)}]-90:0.05) arc[start angle=-90,end angle=-180,radius=0.05]
        -- cycle;
  \draw[line width=\chemlabLineWidth] (0.5,0.65) -- (0.5,0.85);
  \fill[cyan!50] (-.6,0.4) -- (-.85,0.6) -- (-.85,0.9) -- (-.6,1.1) -- (-.3,1.1) -- (-.3,0.4) -- cycle;
  \draw[line width=\chemlabLineWidth] (-0.6,0.4) -- (-0.6,1.1);
}

% \tkzSeparatoryFunnel  (主调用)
% 可选参数 : stopper=true/false, TapOpen=true/false, liquidLevel1, liquidLevel2
\newcommand{\tkzSeparatoryFunnel}[1][]{%
  \begingroup
  \chemlabUpdateLineWidths
  \pgfkeys{
    /tkzSeparatoryFunnel/.cd,
    stopper/.is if=TkzTestTubeStopper,
    TapOpen/.is if=TkzTapOpen,
    liquidLevel1/.store in=\tkzTestTubeLiquidLevelI,
    liquidLevel2/.store in=\tkzTestTubeLiquidLevelII,
    #1}
  % 默认值
  \ifdefined\tkzTestTubeLiquidLevelI\else\def\tkzTestTubeLiquidLevelI{11}\fi
  \ifdefined\tkzTestTubeLiquidLevelII\else\def\tkzTestTubeLiquidLevelII{12}\fi
  % 水相
  \begin{scope}
    \clip \tkzSeparatoryFunnelOutlineClip;
    \fill[AqueousPhaseColor] (-2,1) rectangle (2,\tkzTestTubeLiquidLevelI);
  \end{scope}
  % 有机相
  \begin{scope}
    \clip \tkzSeparatoryFunnelOutlineClip;
    \fill[OrganicPhaseColor] (-2,\tkzTestTubeLiquidLevelI)
          rectangle (2,\tkzTestTubeLiquidLevelII);
  \end{scope}
  \ifTkzTestTubeStopper\tkzStopper\fi
  \tkzSeparatoryFunnelBody
  \ifTkzTapOpen\tkzTapOpen\else\tkzTapClosed\fi
  % 出口管口 (pst-labo: \psellipse[fillstyle=solid,fillcolor=white] 既有白色填充又有黑色轮廓)
  \filldraw[fill=white,line width=\chemlabLineWidth] (0,-1) ellipse (0.15 and 0.25);
  % 瓶口圆点 (填充, 无需线宽)
  \fill[darkgray] (-0.675,15) circle (0.03);
  \fill[darkgray] (0.675,15) circle (0.03);
  \draw[darkgray,line width=\chemlabLineWidth] (-0.65,15) -- (0.65,15);
  % 颈与瓶身接合
  \draw[gray,line width=\chemlabLineWidth] (-0.55,13.75) -- (0.55,13.75);
  \endgroup
}

% ============================================================
% 27. 支架 (pstSupport)
% ============================================================
% pst-labo: \psset{dimen=m,linewidth=0.04} — linewidth=0.04 是用户单位, 随 unit 缩放
% TikZ: 用 \chemlabLineWidth (基础 0.053cm × scale), 与其他器皿一致
\newcommand{\tkzStand}{%
  \chemlabUpdateLineWidths
  \begin{scope}[line width=\chemlabLineWidth]
    \filldraw[fill=woodColor,draw=black,line join=round]
        (-2,-6) rectangle (5,-5);
    \filldraw[fill=woodColor,draw=black,line join=round]
        (2.75,2) rectangle (3.25,16);
    \filldraw[fill=woodColor,draw=black,line join=round]
        ([shift={(2.75,1.5)}]90:0.5) arc[start angle=90,end angle=180,radius=0.5]
        -- (2.25,1.5) -- (2.25,-5) -- (3.75,-5) -- (3.75,1.5)
        -- ([shift={(3.25,1.5)}]0:0.5) arc[start angle=0,end angle=90,radius=0.5]
        -- cycle;
    \draw (2,11.5) -- (1.55,10.5);
    \filldraw[fill=woodColor,draw=black,line join=round]
        (-2.5,11.5) -- (-2,11.5) -- (-1.55,10.5) -- (-2.5,10.5) -- cycle;
    \filldraw[fill=woodColor,draw=black,line join=round]
        (4,11.5) -- (2,11.5) -- (1.55,10.5) -- (4,10.5) -- cycle;
    \draw (-2.5,10.5) rectangle (4,11.5);
    % 旋钮: pst-labo \pscustom[style=woodColor]{...} 填充会自动闭合, 但描边只画实际路径
    % (不画对角闭合边). TikZ: \fill 自动闭合填充, \draw 画开放路径 (无 -- cycle).
    \fill[woodColor]
        (4,11.2) -- (4.5,11.2)
        -- ([shift={(5,11)}]158.2:0.5385) arc[start angle=158.2,end angle=-158.2,radius=0.5385]
        -- (4.5,10.8) -- (4,10.8) -- cycle;
    \draw
        (4,11.2) -- (4.5,11.2)
        -- ([shift={(5,11)}]158.2:0.5385) arc[start angle=158.2,end angle=-158.2,radius=0.5385]
        -- (4.5,10.8) -- (4,10.8);
    \draw (4.5,10.8) -- (4.5,11.2);
  \end{scope}
}

\endinput

%    \begin{tikzpicture}
%       \BuretteCorps  ... 等等
%    \end{tikzpicture}
%
%  映射规则 :
%    \psline(A)(B)(C)              -> \draw (A) -- (B) -- (C);
%    \psline{->}(A)(B)             -> \draw[->] (A) -- (B);
%    \pspolygon(A)(B)(C)           -> \filldraw (A) -- (B) -- (C) -- cycle;
%    \psframe(x1,y1)(x2,y2)        -> \draw (x1,y1) rectangle (x2,y2);
%    \psframe*[linecolor=c]        -> \fill[c] (..) rectangle (..);
%    \psframe[framearc=r]          -> \draw[rounded corners=r] (...) rectangle (...);
%    \pscircle(c){r}               -> \draw (c) circle (r);
%    \pscircle*[linecolor=c](c){r} -> \fill[c] (c) circle (r);
%    \psellipse(c)(a,b)            -> \draw (c) ellipse (a and b);
%    \psarc(c){r}{a1}{a2}          -> \draw ([shift={(c)}]a1:r) arc[start angle=a1,end angle=a2,radius=r];
%    \psarcn(c){r}{a1}{a2}         -> \draw ([shift={(c)}]a1:r) arc[start angle=a1,end angle=a2,radius=r];   (注意 : PSTricks 中 arcn 是逆时针 / 角度递减)
%    \psbezier(A)(B)(C)(D)         -> \draw (A) .. controls (B) and (C) .. (D);
%    \pscurve(P1)(P2)...(Pn)       -> \draw plot[smooth,tension=0.5] coordinates {(P1) (P2) ... (Pn)};
%    \pscustom{ ... \fill}         -> \filldraw / \path[fill] (整段路径) -- cycle;
%    \rput(x,y){...}               -> \begin{scope}[shift={(x,y)}] ... \end{scope}
%    \rput{angle}(x,y){...}        -> \begin{scope}[shift={(x,y)},rotate=angle] ... \end{scope}
%    \multido{\i=a+b}{n}{...}      -> \foreach \i in {a,a+b,...,a+(n-1)*b} {...}
%    \psclip{<path>} ... \endpsclip-> \begin{scope} \clip <path>; ... \end{scope}
%    fillstyle=solid,fillcolor=c   -> fill=c
%    fillstyle=gradient,gradbegin=..,gradend=.. -> \shade[left color=..,right color=..]
%    fillstyle=hlines              -> pattern=horizontal lines
%    doubleline=true               -> double
%    linestyle=dashed              -> dashed
%    framearc=r                    -> rounded corners=r (近似)
%    linearc=r                     -> rounded corners=r
% ============================================================

\NeedsTeXFormat{LaTeX2e}
\ProvidesFile{tikz-chemlabObj.tex}[2026/08/03 v1.0.0]

\RequirePackage{tikz}
\usetikzlibrary{arrows.meta,calc,patterns,patterns.meta,shadings,decorations.pathmorphing}

% ------------------------------------------------------------
% 颜色定义 (对应原包中使用的颜色)
% ------------------------------------------------------------
\definecolor{LightGray}{gray}{0.80}      % 浅灰
\definecolor{VeryLightGray}{cmyk}{0,0,0,0.1}  % 极浅灰 (label 标签背景)
\definecolor{LightBlue}{rgb}{0.68,0.85,0.90}
\definecolor{LightBlue1}{rgb}{0.60,0.80,1.00}
\definecolor{VeryPaleOrange}{cmyk}{0,0.1,0.3,0}  % 滴定管液体默认色
\definecolor{PaleOrange}{rgb}{1.00,0.78,0.55}
\definecolor{Brown}{rgb}{0.50,0.25,0.00}
\definecolor{woodColor}{rgb}{0.957,0.84,0.69}
\definecolor{darkWoodColor}{rgb}{0.50,0.25,0.00}
\definecolor{flameAndGrill}{rgb}{0.95,0.55,0.10}  % 估算颜色
\definecolor{Beige}{rgb}{0.96,0.96,0.86}  % Champagne 液体默认色

% ------------------------------------------------------------
% 公共缺省值 (供调用方按需修改)
% ------------------------------------------------------------
\tikzset{
  Distillation/.cd,
  DistillateColor/.store in=\DistillationDistillateColor,
  DistillateColor=LightBlue1,
  MixtureStyle/.style={fill=LightBlue},
}
\newcommand{\tkzDistillationDistillateColor}{LightBlue1}
\newcommand{\tkzDistillationMixtureStyle}{LightBlue}
\newcommand{\tkzTestTubeLiquidLevelI}{11}     % liquidLevel1
\newcommand{\tkzTestTubeLiquidLevelII}{12}    % liquidLevel2
\newif\ifTkzTapOpen
\newif\ifTkzTestTubeStopper

% ============================================================
% 1. 滴定管 (Burette)
% ============================================================
\newcommand{\tkzBuretteBody}{%
  \draw (-0.5,11.5) -- (-0.25,11) -- (-0.25,3) -- (0,2.75) -- (0,0.5);
  \draw (0.2,0.7) -- (0.2,2.75) -- (0.5,3) -- (0.5,11) -- (0.75,11.5);
}

\newcommand{\tkzBuretteStopcock}{%
  % pst-labo: \psframe[fillstyle=solid] → 填充默认白色 + 黑色边框
  \filldraw[fill=white] (-.1,1.6) rectangle (.3,2.1);
  \draw (-.1,1.5) rectangle (-.2,2.2);
  \draw (.3,1.5) rectangle (.4,2.2);
  \draw (.4,1.7) -- (.5,1.7) -- (.6,1.4) -- (.8,1.4)
        -- (.8,2.3) -- (.6,2.3) -- (.5,2) -- (.4,2) -- cycle;
  % 左侧旋钮: pst-labo \pspolygon + \psframe with fillstyle=solid,fillcolor=LightBlue
  % 需要 LightBlue 填充 + 黑色外轮廓线
  \begin{scope}[fill=LightBlue,draw=black]
    \filldraw (-.5,1.5) -- (-.75,1.7) -- (-.75,2) -- (-.5,2.2)
          -- (-.3,2.2) -- (-.3,1.5) -- cycle;
    \filldraw (-.5,1.5) rectangle (-.2,2.2);
  \end{scope}
}

\newcommand{\tkzBuretteGraduation}{%
  % 滴定管刻度 (原 pst-labo 无显式线宽, 这里用细刻度线宽)
  \foreach \n in {3.0,3.32,...,10.68} { \draw[line width=\chemlabGraduationLineWidth] (-0.25,\n) -- (0,\n); }
}

% ============================================================
% 2. pH 计 (pH-mètre)
% ============================================================
\newcommand{\tkzTitrationPHMeter}{%
  % boitier
  \filldraw[fill=white,draw=black,rounded corners=0.125cm] (-1,0) rectangle (1,3);
  \fill[LightGray] (-0.75,2) rectangle (0.75,2.75);
  \node[white] at (0,2.375) {PH};
  \draw (0,1.25) circle (0.3);
  \draw[->] (0,1.25) -- (0,1.5);
  \node at (0.65,1.5) {$^\circ$C};
  \draw (-0.5,0.5) circle (0.2);
  \draw (0.5,0.5) circle (0.2);
  % électrode — 完全透明填充 (不遮挡烧杯刻度线)
  % pst-labo: \psframe[framearc=0.5](2.25,2.5)(2.75,4.5) 外框圆角
  %   顶部圆角被下方 \psframe[LightGray](sharp) 覆盖, 底部圆角被 \pscircle 覆盖
  %   实际可见的外框是直角矩形, 此处用 sharp corners 匹配 pst-labo 可见效果
  \draw (2.25,2.5) rectangle (2.75,4.5);
  % 顶部 LightGray 填充 + 描边 (sharp corners)
  \filldraw[fill=LightGray,draw=black] (2.25,4) rectangle (2.75,4.5);
  \draw (2.5,2.25) circle (0.25);
  % fil
  \draw[line width=0.05cm,LightBlue,rounded corners=0.2cm]
        (2.5,4.5) -- (2.5,5) -- (1.25,4.75) -- (0,5) -- (0,3);
}

% ============================================================
% 3. 滴定支架 (Support dosage)
% ============================================================
\newcommand{\tkzTitrationStand}{%
  % pst-labo: \psframe(-1.5,0)(1.5,1.5) + {\psset{doubleline=true} ...}
  % 线宽: pst-labo 在 \tkzTitration 中设 linewidth=0.053 (×0.4cm = 0.0212cm ≈ 0.6pt)
  \chemlabUpdateLineWidths
  % 平台: 无填充, 只有边框
  \draw[line width=\chemlabLineWidth] (-1.5,0) rectangle (1.5,1.5);
  % 红色按钮: \pscircle*[linecolor=red] → 红色填充, 双线边框也是红色 (不可见)
  \fill[red] (0,0.75) circle (0.2);
  % 回字型矩形按钮: 无填充, 双线边框
  % PSTricks doubleline: 指定坐标在内外线之间
  % 实测: 外线向外偏移 0.085 units, 内线向内偏移 0.034 units
  \draw[line width=\chemlabLineWidth] ({-1.25-0.085},{0.3-0.085}) rectangle ({-0.75+0.085},{1.2+0.085});
  \draw[line width=\chemlabLineWidth] ({-1.25+0.034},{0.3+0.034}) rectangle ({-0.75-0.034},{1.2-0.034});
  % 环形旋钮: 无填充, 双线边框
  % 外线向外偏移, 内线向内偏移
  \draw[line width=\chemlabLineWidth] (1,0.75) circle ({0.25+0.085});
  \draw[line width=\chemlabLineWidth] (1,0.75) circle ({0.25-0.034});
}

% ============================================================
% 4. 磁力搅拌子 (Aimant)
% ============================================================
\newcommand{\tkzTitrationStirBar}{%
  % pst-labo: \psset{linewidth=0.053} 在外层 \tkzTitration 中设置, 随 scale 缩放
  \chemlabUpdateLineWidths
  \draw[line width=\chemlabLineWidth,rounded corners=0.075cm,fill=white] (-0.5,0.1) rectangle (0.5,0.4);
  \draw[line width=\chemlabLineWidth,rounded corners=0.025cm,fill=white] (-0.05,0) rectangle (0.05,0.5);
}

% ============================================================
% 5. 移液管 (Pipette)
% ============================================================
\newcommand{\tkzPipetteBody}{%
  \chemlabUpdateLineWidths\begin{scope}[line width=\chemlabLineWidth,rounded corners=0.05cm]
    \draw (-0.05,0) -- (-0.15,0.2) -- (-0.15,6)
          -- (0.15,6) -- (0.15,0.2) -- (0.05,0) -- cycle;
    \fill[white] (0,6) ellipse (0.15 and 0.05);
    \foreach \Level in {0.5,1.0,...,5.0} {%
      \draw[line width=\chemlabGraduationLineWidth] (0,\Level) -- (0.1,\Level);}
  \end{scope}
}

\newcommand{\tkzLargePipetteBody}{%
  \chemlabUpdateLineWidths\begin{scope}[line width=\chemlabLineWidth,rounded corners=0.05cm]
    \draw (0,0) -- (0,1) -- (-0.23,2.5) -- (-0.23,22.5)
          arc[start angle=-80,end angle=80,radius=0.45];
  \end{scope}
}

% ============================================================
% 6. 液滴 (Goutte)
% ============================================================
% 原文件中 \pst@goutte 被重复定义了两次, 此处使用后者 (unit=0.6 的版本)
% pst-labo: \psset{unit=0.6}\psbezier[linewidth=0.017] — linewidth=0.017 是
%   0.017 用户单位 = 0.017*0.6 = 0.0102cm (随 goutte 的 LOCAL unit=0.6 缩放, 
%   不随外层 unit 变化, 因为 \psset{unit=0.6} 覆盖了外层设置)
% TikZ: line width=0.0102cm (ABSOLUTE, 匹配 pst-labo 的固定 0.0102cm)
\newcommand{\tkzDrop}{%
  \begin{scope}[scale=0.6]
    \draw[line width=0.0102cm]
      (0,0) .. controls (0.25,-.4) and (-.25,-.4) .. (0,0);
  \end{scope}
}

% ============================================================
% 7. 量筒 (Eprouvette)
% ============================================================
\newcommand{\tkzGraduatedCylinderBody}{%
  \draw (-0.5,5) -- (-0.5,0) -- (0.5,0) -- (0.5,4.9) -- (0.6,4.93);
}

\newcommand{\tkzGraduatedCylinderRim}{%
  \draw plot[smooth,tension=0.5] coordinates {%
    (0.6,-0.07) (0.498,0.008) (0.483,0.026) (0.433,0.05)
    (0.353,0.0707) (0.25,0.0866) (0.13,0.0966) (0,0.1)
    (-0.13,0.0966) (-0.25,0.0866) (-0.353,0.0707) (-0.433,0.05)
    (-0.483,0.026) (-0.5,0) (-0.483,-0.026) (-0.353,-0.0707)
    (-0.25,-0.0866) (-0.13,-0.096) (0,-0.1) (0.13,-0.096)
    (0.25,-0.0866) (0.353,-0.0707) (0.433,-0.05)
    (0.498,-0.026) (0.6,-0.07)};
}

\newcommand{\tkzGraduatedCylinderBase}{%
  \fill[white]
    ([shift={(-0.95,0.5)}]0:0.45) arc[start angle=0,end angle=-90,radius=0.45]
    -- ([shift={(-0.95,0.025)}]90:0.025) arc[start angle=90,end angle=270,radius=0.025]
    -- (-0.95,0) -- (0.95,0)
    -- ([shift={(0.95,0.025)}]-90:0.025) arc[start angle=-90,end angle=90,radius=0.025]
    -- ([shift={(0.95,0.5)}]-90:0.45) arc[start angle=-90,end angle=-180,radius=0.45]
    -- plot[smooth,tension=0.5] coordinates {(0.5,0.4) (0,0.25) (-0.5,0.4)}
    -- cycle;
  \draw
    ([shift={(-0.95,0.5)}]0:0.45) arc[start angle=0,end angle=-90,radius=0.45]
    -- ([shift={(-0.95,0.025)}]90:0.025) arc[start angle=90,end angle=270,radius=0.025]
    -- (-0.95,0) -- (0.95,0)
    -- ([shift={(0.95,0.025)}]-90:0.025) arc[start angle=-90,end angle=90,radius=0.025]
    -- ([shift={(0.95,0.5)}]-90:0.45) arc[start angle=-90,end angle=-180,radius=0.45]
    -- plot[smooth,tension=0.5] coordinates {(0.5,0.4) (0,0.25) (-0.5,0.4)};
}

\newcommand{\tkzGraduatedCylinderGraduation}{%
  % 量筒刻度 : 主刻度 1.5× 细刻度 (与烧杯一致)
  \pgfmathsetlengthmacro{\chemlab@eprouvette@major}{1.5*\chemlabGraduationLineWidth}%
  \foreach \nyD in {0.5,1.0,...,4.5} { \draw[line width=\chemlab@eprouvette@major] (-0.4,\nyD) -- (-0.1,\nyD); }
  \foreach \nyd in {0.5,0.6,...,4.4} { \draw[line width=\chemlabGraduationLineWidth] (-0.4,\nyd) -- (-0.25,\nyd); }
}

% ============================================================
% 8. 分馏装置 (Distillation fractionnée)
% ============================================================
% \pst@goutte 这里又被定义了一次, 在 TikZ 版中沿用同一个 \tkzDrop

\newcommand{\tkzColumn}{%
  % pst-labo: linearc=0.3 (×unit, 不是绝对值) → TikZ: rounded corners=0.12cm (= 0.3 × 0.4cm)
  % 尖刺 (pointeG/D) 先画并用 clip 裁剪到蒸馏柱内部, 再画轮廓线覆盖切边。
  %   pst-labo 中 linearc=0.3 (R=0.3 user units), 尖刺端点恰在圆弧切点附近, 几乎
  %   不可见。TikZ 的 rounded corners 实际 R 比标称值大 ~2×, 圆弧切点更低,
  %   尖刺端点 (在 y_col=8 和 7.9) 会突出到圆弧外侧形成多余线头。用 clip 将
  %   尖刺裁剪到柱内, 再画轮廓线覆盖切边, 消除突出线头。
  \begin{scope}
    \clip[rounded corners=0.12cm]
      (-.25,10.5) -- (-.25,8.5) -- (-.75,8) -- (-.75,2.5) -- (-.25,2) -- (-.25,0)
      -- (0.25,0) -- (0.25,2) -- (.75,2.5) -- (.75,8) -- (0.25,8.5) -- (0.25,9.36)
      -- (0.25,10.5) -- cycle;
    % 左侧尖刺 (pointeG)
    \foreach \i in {3,4,...,8} {%
      \draw[line width=0.4\pgflinewidth] (-.75,\i) -- (-.25,\i-0.5)
          -- (-.75,\i-0.1);}
    % 右侧尖刺 (pointeD)
    \foreach \i in {3,4,...,8} {%
      \draw[line width=0.4\pgflinewidth] (.75,\i) -- (.25,\i-0.5)
          -- (.75,\i-0.1);}
  \end{scope}
  \draw[rounded corners=0.12cm]
    (-.25,10.5) -- (-.25,8.5) -- (-.75,8) -- (-.75,2.5)
    -- (-.25,2) -- (-.25,0);
  \draw[rounded corners=0.12cm]
    (0.25,0) -- (0.25,2) -- (.75,2.5) -- (.75,8)
    -- (0.25,8.5) -- (0.25,9.36);
  % 上部断口 (line 3a/3b):
  %   pst-labo 源码 line 3a: \psline(0.25,8.5)(0.25,9.30) 用于填充 linearc 圆角
  %   留下的间隙。实测发现 TikZ 的 rounded corners 实际圆弧半径比标称值大
  %   ~2× (在 scale=0.6 下 rounded corners=0.3cm 给出 R≈0.62cm), 圆弧已覆盖到
  %   x=0.25 处 y_col=8.93 附近, 故 line 3a 在任何 scale/line width 下均被圆弧
  %   覆盖, 不可见。为避免在编号图等彩色调试场景中残留多余线段, 此处不再
  %   单独绘制 line 3a。
  \draw (0.25,9.5) -- (0.25,10.5);
  % 填充 / 阴影
  \begin{scope}[pattern={Lines[angle=45,distance={1pt/sqrt(2)}]},pattern color=gray]
    \fill (-.25,10) rectangle (-.05,10.5);
    \fill (.25,10)  rectangle (0.05,10.5);
  \end{scope}
  \foreach \n in {2.5,3.5,...,7.5} {%
    \begin{scope}[shift={(-.25,\n)}]\tkzDrop\end{scope}
    \begin{scope}[shift={(.25,\n)}]\tkzDrop\end{scope}}
  \begin{scope}[shift={(-0.25,-.1)}]\tkzDrop\end{scope}
  \begin{scope}[shift={(0.25,-.1)}]\tkzDrop\end{scope}
}

% 公共别名: pst-labo 中 \tkzColumn 是用户可调用的公共宏 (小写)
% tikz-chemlab 内部用 \tkzColumn (驼峰), 这里注册公共别名保持兼容
\let\tkzColumn\tkzColumn

% réfrigérant (condenseur)
\newcommand{\tkzCondenser}{%
  \draw ([shift={(3,0)}]90:0.5) arc[start angle=90,end angle=270,radius=0.5];
  % pst-labo: \psarc(7,0){0.5}{270}{90} → 逆时针经过 0°/右侧
  \draw ([shift={(7,0)}]270:0.5) arc[start angle=270,end angle=450,radius=0.5];
  \fill[white] (0,-.05) rectangle (9,0.05);
  \begin{scope}[line width=0.4pt]
    \draw (-.075,.05) -- (9,0.05);
    \draw (0.02,-.05) -- (9,-.05);
  \end{scope}
  \draw (3,0.5) -- (3.5,0.5);
  \draw (3,-.5) -- (6.5,-.5);
  \draw (3.6,0.5) -- (7,0.5);
  \begin{scope}[line width=0.4pt]
    \draw (3.5,0.5) -- (3.5,1.5);
    \draw (3.6,0.5) -- (3.6,1.5);
    \draw (6.6,-.5) -- (6.6,-1.5);
    \draw (6.5,-.5) -- (6.5,-1.5);
  \end{scope}
  \draw (6.6,-.5) -- (7,-.5);
  \draw[->] (6.55,-2.5) -- (6.55,-1.5);
  \draw[->] (3.55,1.5)  -- (3.55,2.5);
}

\newcommand{\tkzThermometer}{%
  % pst-labo: \psline[linearc=0.05](-.05,.25)(-.05,2.75)(0.05,2.75)(0.05,0.25)
  % linearc=0.05 × unit, TikZ: rounded corners=0.05 (受 scale 缩放)
  \draw[rounded corners=0.025cm]
    (-.05,.25) -- (-.05,2.75) -- (0.05,2.75) -- (0.05,0.25) -- cycle;
  % pst-labo: \psframe[framearc=0.8,fillstyle=solid,fillcolor=red](-.15,-.25)(0.15,0.25)
  % framearc=0.8: 80% of half-min(0.3,0.5)=0.15 → radius=0.12
  \filldraw[fill=red,draw=black,rounded corners=0.06cm] (-.15,-.25) rectangle (0.15,0.25);
  % pst-labo: \psline[linecolor=red](0,2) — 红色竖线, 使用默认线宽
  % 在 pst-labo 渲染中较粗, 看起来像红色填充
  \draw[red,line width=0.8pt] (0,0) -- (0,2);
}

\newcommand{\tkzHeater}{%
  \draw (1.17,2.5) -- (0,2.5) -- (0,0) -- (5.5,0) -- (5.5,2.5) -- (4.33,2.5);
}

% 内部小量筒轮廓 - 路径版本 (供 clip 使用)
\newcommand{\tkzGraduatedCylinderPath}{%
  (-0.5,6) -- (-0.5,.5)
  -- ([shift={(-1,0.5)}]0:0.5) arc[start angle=0,end angle=-90,radius=0.5]
  -- (-1,0) -- (1,0)
  -- ([shift={(1,0.5)}]270:0.5) arc[start angle=270,end angle=180,radius=0.5]
  -- (0.5,0.5) -- (0.5,6) -- cycle
}

% 内部小量筒 (描边)
% pst-labo: \psset{linewidth=1.5\pslinewidth} — \pslinewidth 默认 0.8pt (ABSOLUTE)
% → 1.5*0.8pt = 1.2pt (绝对值, 不随 unit/scale 缩放)
% 注意: TikZ 的 [line width=N\pgflinewidth] 有循环引用 bug — \pgflinewidth 在
% 乘法前被更新, 导致 line width=3\pgflinewidth 实际给出 3.6pt 而非 1.2pt。
% 必须用显式 pt 值。
\newcommand{\tkzGraduatedCylinderSmall}{%
  \begin{scope}[line width=0.72pt,line cap=round,line join=round]
    \draw (-0.5,6) -- (-0.5,.5)
      -- ([shift={(-1,0.5)}]0:0.5) arc[start angle=0,end angle=-90,radius=0.5]
      -- (-1,0) -- (1,0)
      -- ([shift={(1,0.5)}]270:0.5) arc[start angle=270,end angle=180,radius=0.5]
      -- (0.5,0.5) -- (0.5,6);
  \end{scope}
}

% 带液体的量筒 (内部使用, 重命名以避免与主文件 \pstEproulette 冲突)
\newcommand{\tkzGraduatedCylinderWithLiquid}{%
  \begin{scope}
    \clip \tkzGraduatedCylinderPath;
    % pst-labo: \psframe*[linecolor=\psk@Distillation@DistillateColor]
    \fill[\tkzDistillateColor] (-1.5,-1) rectangle (1.5,1);
  \end{scope}
  \tkzGraduatedCylinderSmall
}

% (注 : 原 PSTricks 中 \pst@roundFlask 已合并到 \tkzRoundFlaskSimple, 此处不再单独定义,
%  避免与 tikz-chemlab.tex 中的用户命令 \tkzRoundFlask 冲突)

% 自定义密集斜线图案 (匹配 pst-labo hlines 的密度)
{\pgfpointorigin}{\pgfpoint{0.012cm}{0.012cm}}{\pgfpoint{0.012cm}{0.012cm}}
{\pgfsetlinewidth{0.0028cm}\pgfpathmoveto{\pgfpoint{0cm}{0cm}}\pgfpathlineto{\pgfpoint{0.012cm}{0.012cm}}\pgfusepath{stroke}}
{\pgfpointorigin}{\pgfpoint{0.012cm}{0.012cm}}{\pgfpoint{0.012cm}{0.012cm}}
{\pgfsetlinewidth{0.004cm}\pgfpathmoveto{\pgfpoint{0cm}{0cm}}\pgfpathlineto{\pgfpoint{0.012cm}{0.012cm}}\pgfusepath{stroke}}
% \pst@DistillationFractionnee 主图
\newcommand{\tkzDistillationFractionated}{%
  % --- roundFlask ---
  \def\flaskLocal{%
    % pst-labo: fillstyle=hlines,hatchwidth=0.007,hatchsep=0.03 → 水平线填充
    \begin{scope}[pattern={Lines[angle=45,distance={1pt/sqrt(2)}]},pattern color=gray]
      \fill (-.5,2.25) rectangle (-.25,3);
      \fill (0.5,2.25)  rectangle (0.25,3);
    \end{scope}
    % pst-labo 线宽: neck=0.053, arcs=0.03 (都是 ×unit 用户单位, 随 scale 缩放)
    % \chemlabUpdateLineWidths 已设置 \chemlabLineWidth = 0.053*scale
    \chemlabUpdateLineWidths
    \draw[line width=\chemlabLineWidth] (0.5,1.5) -- (0.5,3);
    \draw[line width=\chemlabLineWidth] (-.5,1.5) -- (-.5,3);
    % pst-labo: \psarc(0,0){1.581}{108.435}{71.565} 逆时针走长路 (经底部),
    %   左右侧弧可见, 底部被液体填充覆盖, 顶部无可见弧线。此处拆分为
    %   左右侧两段 (底部由下方轮廓弧绘制, 顶部不画)。
    % 右上侧弧: 71° → 0° (连接右颈底部到球体最右点)
    \draw[line width=\chemlabLineWidth] ([shift={(0,0)}]71.565:1.581) arc[start angle=71.565,end angle=0,radius=1.581];
    % 左上侧弧: 180° → 108° (连接球体最左点到左颈底部)
    \draw[line width=\chemlabLineWidth] ([shift={(0,0)}]180:1.581) arc[start angle=180,end angle=108.435,radius=1.581];
    % pst-labo: \psset{style=\psk@Distillation@MixtureStyle}\fill
    % 需要用 \shade (不是 \fill) 来产生渐变效果
    \shade[\tkzMixtureStyle] plot[smooth,tension=0.5] coordinates {%
      (-1.581,0) (-1.25,0.1) (-.75,-.1) (0,.1) (0.75,-.1) (1.25,.1) (1.581,0)}
      arc[start angle=0,end angle=-180,radius=1.581];
    % 底部轮廓弧: pst-labo \psarc[linewidth=0.03] → 0.03*unit, 用 \chemlabLineWidth 0.053
    % (pst-labo 中 0.03 vs 0.053 差异微小, 统一用 0.053*scale 避免额外变量)
    \draw[line width=\chemlabLineWidth] (0,0) ++(180:1.581) arc[start angle=180,end angle=360,radius=1.581];
  }%
  % --- 主图整体平移 ---
  \begin{scope}[shift={(0,-3)}]
    \begin{scope}[shift={(2.75,13.5)}]\tkzThermometer\end{scope}
    \begin{scope}[shift={(2.75,2.5)}]\flaskLocal\end{scope}
    \begin{scope}[shift={(2.75,4)}]\tkzColumn\end{scope}
    \begin{scope}[shift={(3,13.4)},rotate=-45]\tkzCondenser\end{scope}
    \begin{scope}[shift={(0,0)}]\tkzHeater\end{scope}
    \begin{scope}[shift={(9.4,0)}]\tkzGraduatedCylinderWithLiquid
      % 蒸馏装置量筒刻度 : 主刻度 1.5× 细刻度
      \pgfmathsetlengthmacro{\chemlab@distill@major}{1.5*\chemlabGraduationLineWidth}%
      \foreach \Level in {1,2,...,5} {\draw[line width=\chemlab@distill@major] (-.4,\Level) -- (-.1,\Level);}
      \foreach \Level in {0.2,0.4,...,4.8} {\draw[line width=\chemlabGraduationLineWidth] (-.4,\Level) -- (-.2,\Level);}
    \end{scope}
    \begin{scope}[shift={(9.37,7)}]\tkzDrop\end{scope}
    \begin{scope}[shift={(9.37,6)}]\tkzDrop\end{scope}
    \begin{scope}[shift={(9.37,4)}]\tkzDrop\end{scope}
    \begin{scope}[shift={(9.37,1.2)}]\tkzDrop\end{scope}
    % 加热盘 (搅拌器)
    \begin{scope}[shift={(0.75,.75)}]
      \draw (0,0) circle (.5);
      \foreach \i in {0,30,...,180} {\draw (\i:0.55) -- (\i:0.65);}
      \draw[->,line width=0.035cm] (0,0) -- (0,0.45);
    \end{scope}
  \end{scope}
}

% ============================================================
% 9. 本生灯 (Bec Bunsen) - 不带 grille
% ============================================================
\newcommand{\tkzBunsenBurner}{%
  \chemlabUpdateLineWidths\begin{scope}[line width=\chemlabLineWidth,line cap=round,line join=round]
    \draw (-1.25,0) rectangle (1.25,0.25);
    \draw (-.5,1.25) rectangle (0.5,2.25);
    \foreach \n in {-0.3,0.0,0.3} {\draw (\n,1.75) circle (0.1);}
    \draw (-.25,2.25) rectangle (0.25,4.25);
    \draw (0.25,1.25) -- (0.25,0.5) -- (1.25,0.25);
    \draw (-1.25,0.25) -- (-.25,0.5) -- (-0.25,0.75);
    \draw (-2.25,0.75) -- (-.25,0.75);
    \draw (-2.25,1) -- (-.25,1);
    \draw (-.25,0.875) ellipse (0.1 and 0.125);
    \fill[white] (-2.25,0.75) rectangle (-0.25,1);
    \draw (-2.25,0.75) -- (-0.25,0.75);
    \draw (-2.25,1) -- (-0.25,1) -- (-.25,1.25);
    \draw plot[smooth,tension=0.5] coordinates {(-0.25,0.5) (0,0.4) (0.25,0.5)};
    % flamme
    % pst-labo 源码: \rput(0,4.25){\psclip{\psbezier[...gradient...] (12点)}\pspolygon[...gradient...](-0.25,0)(0.25,0)(0,1)\endpsclip}
    % \psbezier 12个点 + BezierNArray padding (0,0) = 13点 = 4段贝塞尔
    % 经像素级验证, 正确路径从 (0,0) 出发, 经4段贝塞尔回到 (0,0), 无闭合直线:
    %   Seg 1: (0,0) → c1=P1=(-0.25,0), c2=P2=(-0.35,0.5) → P3=(-0.4,0.75) [左下→左中]
    %   Seg 2: P3=(-0.4,0.75) → c1=P4=(-0.35,1), c2=P5=(-0.25,1.5) → P6=(0.5,2) [左中→顶]
    %   Seg 3: P6=(0.5,2) → c1=P7=(0.25,1.5), c2=P8=(0.35,1) → P9=(0.4,0.75) [顶→右中]
    %   Seg 4: P9=(0.4,0.75) → c1=P10=(0.35,0.5), c2=P11=(0.25,0) → (0,0) [右中→底]
    % 渐变: gradbegin=PaleOrange (底), gradend=yellow (顶), gradmidpoint=0, gradangle=0 (垂直)
    \begin{scope}[shift={(0,4.25)}]
      % 外焰填充: PaleOrange (底) -> yellow (顶) 线性垂直渐变
      \shade[top color=yellow, bottom color=PaleOrange]
        (0,0) .. controls (-0.25,0) and (-0.35,0.5)
        .. (-0.4,0.75) .. controls (-0.35,1) and (-0.25,1.5)
        .. (0.5,2) .. controls (0.25,1.5) and (0.35,1)
        .. (0.4,0.75) .. controls (0.35,0.5) and (0.25,0)
        .. (0,0);
      % 内焰: 三角形 cyan(底)->white(顶) 渐变, 裁剪到外焰形状内
      \begin{scope}
        \clip
          (0,0) .. controls (-0.25,0) and (-0.35,0.5)
          .. (-0.4,0.75) .. controls (-0.35,1) and (-0.25,1.5)
          .. (0.5,2) .. controls (0.25,1.5) and (0.35,1)
          .. (0.4,0.75) .. controls (0.35,0.5) and (0.25,0)
          .. (0,0);
        \shade[top color=white, bottom color=cyan]
          (-0.25,0) -- (0.25,0) -- (0,1) -- cycle;
      \end{scope}
    \end{scope}
  \end{scope}
}

% 火焰 + 网格 (FlammeGrille 形状) - 路径版本 (供 clip 使用)
\newcommand{\tkzFlameGrillPath}{%
  ([shift={(0,0.75)}]270:0.75) arc[start angle=270,end angle=360,radius=0.75]
  -- ([shift={(1.5,0.75)}]180:0.75) arc[start angle=180,end angle=90,radius=0.75]
  -- (1.5,1.5) -- (-1.5,1.5)
  -- ([shift={(-1.5,0.75)}]90:0.75) arc[start angle=90,end angle=0,radius=0.75]
  -- ([shift={(0,0.75)}]180:0.75) arc[start angle=180,end angle=270,radius=0.75]
  -- cycle
}

\newcommand{\tkzFlameGrill}{%
  \fill[flameAndGrill] \tkzFlameGrillPath;
}

% ============================================================
% 10. 本生灯 - 带 grille
% ============================================================
\newcommand{\tkzBunsenBurnerGrill}{%
  \chemlabUpdateLineWidths\begin{scope}[line width=\chemlabLineWidth,line cap=round,line join=round]
    \draw (-1.25,0) rectangle (1.25,0.25);
    \draw (-.5,1.25) rectangle (0.5,2.25);
    \foreach \n in {-0.3,0.0,0.3} {\draw (\n,1.75) circle (0.1);}
    \draw (-.25,2.25) rectangle (0.25,4.25);
    \draw (0.25,1.25) -- (0.25,0.5) -- (1.25,0.25);
    \draw (-1.25,0.25) -- (-.25,0.5) -- (-0.25,0.75);
    \draw (-2.25,0.75) -- (-.25,0.75);
    \draw (-2.25,1) -- (-.25,1);
    \draw (-.25,0.875) ellipse (0.1 and 0.125);
    \fill[white] (-2.25,0.75) rectangle (-0.25,1);
    \draw (-2.25,0.75) -- (-0.25,0.75);
    \draw (-2.25,1) -- (-0.25,1) -- (-.25,1.25);
    \draw plot[smooth,tension=0.5] coordinates {(-0.25,0.5) (0,0.4) (0.25,0.5)};
    % flamme (clip 到 FlammeGrille)
    % 原 pst-labo 火焰分两层:
    %   外焰: \pst@FlammeGrille 整体填充 flameAndGrill (PaleOrange->yellow 渐变)
    %   内焰: 三角形 cyan->white 渐变 (在内焰区域)
    % 颜色和大小可通过 \chemlabset 调整
    \begin{scope}[shift={(0,4.25)},scale=\chemlabFlameScale]
      \begin{scope}
        \clip \tkzFlameGrillPath;
        % 外焰: PaleOrange (底) -> yellow (顶) 线性垂直渐变
        % 对应 pst-labo: fillstyle=gradient, gradbegin=PaleOrange, gradend=yellow
        % 渐变范围与火焰形状 (y=0..1.5) 一致, 避免底部颜色被稀释
        \shade[top color=\chemlabFlameOuterEndColor, bottom color=\chemlabFlameOuterColor]
          (-1.5,0) rectangle (1.5,1.5);
        % 内焰: cyan (底) -> white (顶) 线性垂直渐变 (三角形)
        \shade[top color=\chemlabFlameInnerEndColor, bottom color=\chemlabFlameInnerColor]
          (-0.25,0) -- (0.25,0) -- (0,1) -- cycle;
      \end{scope}
    \end{scope}
    \draw[dashed,line width=\chemlabEllipseLineWidth] (-2,5.75) -- (2,5.75);
  \end{scope}
}

% ============================================================
% 11. 烧瓶 (Ballon)
% ============================================================
% 路径版本 - 供 clip 使用 (闭合轮廓)
\newcommand{\tkzRoundFlaskSimplePath}{%
  ([shift={(-0.7,3.4)}]90:0.2) arc[start angle=90,end angle=0,radius=0.2]
  -- (-0.5,3.4) -- (-0.5,2.4)
  -- (-0.489,2.401)
  arc[start angle=113,end angle=427,radius=1.25]
  -- (0.5,2.4) -- (0.5,3.4)
  -- ([shift={(0.7,3.4)}]180:0.2) arc[start angle=180,end angle=90,radius=0.2]
  -- (0,3.5) -- cycle
}

\newcommand{\tkzRoundFlaskSimple}{%
  \chemlabUpdateLineWidths\begin{scope}[line width=\chemlabLineWidth,line cap=round,line join=round]
    % 左侧: 瓶口圆弧 + 颈部直线 + 球体大圆弧, 合并为单条路径以平滑过渡
    \draw ([shift={(-0.7,3.4)}]90:0.2) arc[start angle=90,end angle=0,radius=0.2]
          -- (-0.5,3.4) -- (-0.5,2.4)
          -- (-0.489,2.401) arc[start angle=113,end angle=427,radius=1.25];
    % 右侧: 球体终点(0.489,2.401) -> 颈部直线 -> 瓶口圆弧
    \draw (0.489,2.401) -- (0.5,2.4) -- (0.5,3.4)
          -- ([shift={(0.7,3.4)}]180:0.2) arc[start angle=180,end angle=90,radius=0.2];
    % 瓶口椭圆: 白色填充覆盖颈部上端圆头, 避免穿入椭圆
    % pst-labo: ballon椭圆用默认线宽 0.053 (与主轮廓一致, 不是 0.07)
    \filldraw[fill=white,draw=black,line width=\chemlabLineWidth] (0,3.6) ellipse (0.7 and 0.1);
  \end{scope}
}

\newcommand{\tkzRoundFlaskGrill}{%
  \chemlabUpdateLineWidths\begin{scope}[line width=\chemlabLineWidth,line cap=round,line join=round]
    % 左侧: 瓶口圆弧 + 颈部直线 + 球体大圆弧, 合并为单条路径
    \draw ([shift={(-0.7,2.15)}]90:0.2) arc[start angle=90,end angle=0,radius=0.2]
          -- (-0.5,2.15) -- (-0.5,1.4)
          -- (-0.489,1.401) arc[start angle=113,end angle=427,radius=1.25];
    % 右侧: 球体终点 -> 颈部直线 -> 瓶口圆弧
    \draw (0.489,1.401) -- (0.5,1.15) -- (0.5,3.4)
          -- ([shift={(0.7,2.15)}]180:0.2) arc[start angle=180,end angle=90,radius=0.2];
    % 瓶口椭圆: 白色填充覆盖颈部上端圆头, 避免穿入椭圆
    % pst-labo: ballon椭圆用默认线宽 0.053 (与主轮廓一致)
    \filldraw[fill=white,draw=black,line width=\chemlabLineWidth] (0,2.35) ellipse (0.7 and 0.1);
  \end{scope}
}

% ============================================================
% 12. 试管 (TubeEssais)
% ============================================================
\newcommand{\tkzTestTubeStopper}{%
  \chemlabUpdateLineWidths
  \begin{scope}[fill=LightGray,draw=black,line width=\chemlabLineWidth,line cap=round,line join=round]
    \filldraw (0,-0.3) ellipse (0.4 and 0.1);
    \filldraw (-0.4,-0.3) -- (-0.6,0.5) -- (0.6,0.5) -- (0.4,-0.3) -- cycle;
    \fill[draw=none] (0,-0.3) ellipse (0.4 and 0.1);
    \filldraw (0,0.5) ellipse (0.6 and 0.1);
  \end{scope}
  \draw[line width=\chemlabEllipseLineWidth] plot[smooth,tension=0.5]
        coordinates {(-0.5,0) (0,-0.07) (0.5,0)};
}

\newcommand{\tkzTestTubeBody}{%
  \chemlabUpdateLineWidths\begin{scope}[line width=\chemlabLineWidth,line cap=round,line join=round]
    % 合并左侧直线 + 底部圆弧 + 右侧直线为单条路径, 平滑转角
    \draw (-0.5,3.5) -- (-0.5,0.5)
          arc[start angle=180,end angle=360,radius=0.5]
          -- (0.5,3.5);
    % 试管口椭圆 (原 pst-labo 用 0.08cm, 比其他椭圆粗)
    \draw[line width=\chemlabTubeEllipseLineWidth] (0,3.5) ellipse (0.5 and 0.1);
  \end{scope}
}

\newcommand{\tkzTestTubeStraightTube}{%
  \draw[line width=0.5\pgflinewidth] (0.05,-1) -- (0.05,5);
  \draw[line width=0.5\pgflinewidth] (-0.05,-1) -- (-0.05,5);
  \tkzTestTubeStopper
  \draw[fill=white,rounded corners=0.5,line width=0.5\pgflinewidth]
       (-0.05,5) rectangle (0.05,0.5);
  \draw[white,line width=0.14cm] (-0.1,5) -- (0.1,5);
}

\newcommand{\tkzTestTubeElbowTube}{%
  \draw[line width=0.5\pgflinewidth] (0.05,-1) -- (0.05,0);
  \draw[line width=0.5\pgflinewidth] (-0.05,-1) -- (-0.05,0);
  \tkzTestTubeStopper
  % 弯管路径: pst-labo \pspolygon(...)(8,2)(8,1.9)...
  \filldraw[fill=white,rounded corners=0.05cm,line width=0.5\pgflinewidth]
       (-0.05,0.5) -- (-0.05,2) -- (8,2) -- (8,1.9)
       -- (0.05,1.9) -- (0.05,0.5) -- cycle;
  \draw[white,line width=0.14cm] (7.95,1.9) -- (8.05,2);
}

\newcommand{\tkzTestTubeUTube}{%
  \draw[line width=0.5\pgflinewidth] (0.05,-1) -- (0.05,0);
  \draw[line width=0.5\pgflinewidth] (-0.05,-1) -- (-0.05,0);
  \tkzTestTubeStopper
  \filldraw[fill=white,rounded corners=0.05cm,line width=0.5\pgflinewidth]
       (-0.05,0.5) -- (-0.05,2) -- (4.5,2) -- (4.5,-3)
       -- (4.4,-3) -- (4.4,1.9) -- (0.05,1.9) -- (0.05,0.5) -- cycle;
  \draw[white,line width=0.14cm] (4.3,-3) -- (4.6,-3);
}

\newcommand{\tkzTestTubeUTubeLong}{%
  \draw[line width=0.5\pgflinewidth] (0.05,-1) -- (0.05,0);
  \draw[line width=0.5\pgflinewidth] (-0.05,-1) -- (-0.05,0);
  \tkzTestTubeStopper
  \filldraw[fill=white,rounded corners=0.05cm,line width=0.5\pgflinewidth]
       (-0.05,0.5) -- (-0.05,2) -- (4.5,2) -- (4.5,-9)
       -- (4.4,-9) -- (4.4,1.9) -- (0.05,1.9) -- (0.05,0.5) -- cycle;
  \draw[white,line width=0.14cm] (4.3,-9) -- (4.6,-9);
}

\newcommand{\tkzCurvedTube}{%
  \draw[line width=0.5\pgflinewidth] (0.05,-1) -- (0.05,0);
  \draw[line width=0.5\pgflinewidth] (-0.05,-1) -- (-0.05,0);
  \tkzTestTubeStopper
  \begin{scope}[line width=0.4pt]
    \draw (-0.05,0.5) -- (-0.05,1);
    \draw ([shift={(0.5,1)}]180:0.55) arc[start angle=180,end angle=90,radius=0.55];
    \draw (0.5,1.55) -- (5.5,1.55);
    \draw ([shift={(5.5,1)}]90:0.55) arc[start angle=90,end angle=0,radius=0.55];
    \draw (6.05,1) -- (6.05,-8.8);
    \draw ([shift={(6.5,-8.8)}]180:0.45) arc[start angle=180,end angle=360,radius=0.45];
    \draw (6.95,-8.8) -- (6.95,-8.3);
    \draw (7.05,-8.3) -- (7.05,-8.8);
    \draw ([shift={(6.5,-8.8)}]0:0.55) arc[start angle=0,end angle=-180,radius=0.55];
    \draw (5.95,-8.8) -- (5.95,1);
    \draw ([shift={(5.5,1)}]0:0.45) arc[start angle=0,end angle=90,radius=0.45];
    \draw (5.5,1.45) -- (0.5,1.45);
    \draw ([shift={(0.5,1)}]90:0.45) arc[start angle=90,end angle=180,radius=0.45];
    \draw (0.05,1) -- (0.05,0.5);
    \draw (-0.05,0.5) -- (0.05,0.5);
  \end{scope}
  \fill[white] (-0.025,0.52) rectangle (0.020,0.6);
}

\newcommand{\tkzCurvedTubeShort}{%
  \draw[line width=0.5\pgflinewidth] (0.05,-1) -- (0.05,0);
  \draw[line width=0.5\pgflinewidth] (-0.05,-1) -- (-0.05,0);
  \tkzTestTubeStopper
  \draw (-0.05,0.5) -- (-0.05,1);
  \begin{scope}[line width=0.4pt]
    \draw ([shift={(0.5,1)}]180:0.55) arc[start angle=180,end angle=90,radius=0.55];
    \draw (0.5,1.55) -- (3.5,1.55);
    \draw ([shift={(3.5,1)}]90:0.55) arc[start angle=90,end angle=0,radius=0.55];
    \draw (4.05,1) -- (4.05,-2.8);
    \draw ([shift={(4.5,-2.8)}]180:0.45) arc[start angle=180,end angle=360,radius=0.45];
    \draw (4.95,-2.8) -- (4.95,-2.3);
    \draw (5.05,-2.3) -- (5.05,-2.8);
    \draw ([shift={(4.5,-2.8)}]0:0.55) arc[start angle=0,end angle=-180,radius=0.55];
    \draw (3.95,-2.8) -- (3.95,1);
    \draw ([shift={(3.5,1)}]0:0.45) arc[start angle=0,end angle=90,radius=0.45];
    \draw (3.5,1.45) -- (0.5,1.45);
    \draw ([shift={(0.5,1)}]90:0.45) arc[start angle=90,end angle=180,radius=0.45];
    \draw (0.05,1) -- (0.05,0.5);
    \draw (-0.05,0.5) -- (0.05,0.5);
  \end{scope}
  \fill[white] (-0.025,0.52) rectangle (0.020,0.6);
}

\newcommand{\tkzCurvedTubeShortDoubleTube}{%
  \begin{scope}[line width=0.4pt]
    \draw (-0.05,0.5) -- (-0.05,1);
    \draw ([shift={(0.5,1)}]180:0.55) arc[start angle=180,end angle=90,radius=0.55];
    \draw (0.5,1.55) -- (3.5,1.55);
    \draw ([shift={(3.5,1)}]90:0.55) arc[start angle=90,end angle=0,radius=0.55];
    \draw (4.05,1) -- (4.05,-2.8);
    \draw ([shift={(4.5,-2.8)}]180:0.45) arc[start angle=180,end angle=360,radius=0.45];
    \draw (4.95,-2.8) -- (4.95,-2.3);
    \draw (5.05,-2.3) -- (5.05,-2.8);
    \draw ([shift={(4.5,-2.8)}]0:0.55) arc[start angle=0,end angle=-180,radius=0.55];
    \draw (3.95,-2.8) -- (3.95,1);
    \draw ([shift={(3.5,1)}]0:0.45) arc[start angle=0,end angle=90,radius=0.45];
    \draw (3.5,1.45) -- (0.5,1.45);
    \draw ([shift={(0.5,1)}]90:0.45) arc[start angle=90,end angle=180,radius=0.45];
    \draw (0.05,1) -- (0.05,0.5);
    \draw (-0.05,0.5) -- (0.05,0.5);
  \end{scope}
  \fill[white] (-0.025,0.52) rectangle (0.018,0.8);
}

\newcommand{\tkzTestTubeDoubleTube}{%
  \begin{scope}[line width=0.4pt]
    \draw (-0.3,-3) -- (-0.3,0);
    \draw (-.2,-3) -- (-0.2,0);
    \draw (0.25,-1) -- (0.25,0);
    \draw (.15,-1) -- (.15,0);
  \end{scope}
  \tkzTestTubeStopper
  \begin{scope}[shift={(0.2,0)}]\tkzCurvedTubeShortDoubleTube\end{scope}
  \begin{scope}[shift={(0,0.5)}]
    \begin{scope}[line width=0.4pt]
      \filldraw[fill=white,rounded corners=0.05cm]
        (-.3,0) -- (-0.3,2) -- (-0.5,2.2) -- (-0.5,3) -- (0,3)
                -- (0,2.2) -- (-0.2,2) -- (-0.2,0) -- cycle;
      \draw (0.1,1.55) -- (-.5,1.55) -- (-0.5,1.7) -- (-.55,1.7)
            -- (-.55,1.3) -- (-0.5,1.3) -- (-0.5,1.55) -- (-.5,1.45) -- (.1,1.45);
      \filldraw[fill=white,rounded corners=0.2] (-0.4,1.3) rectangle (-0.1,1.7);
    \end{scope}
    \fill[white] (-0.25,3) ellipse (0.3 and 0.1);
  \end{scope}
}

% 试管夹 (Pince) - 几何与原 pst-labo pst-laboObj.tex line 483-508 一致
% 根据原 pst-labo 实际渲染 (5.3x 截图) 校准:
%   - 上边缘凹弧较浅 (原 \pscurve 3点几乎共线, 浅弧)
%   - 顶部小凸起为圆弧形 (\psarcn 圆弧)
%   - 所有转角为圆角 (原 \pscustom 默认 linejoin=0 即 round)
\newcommand{\tkzTestTubeClamp}{%
  % 第1块: 上部夹爪
  % 上边缘(1a, 夹合面): 下凹弧贴合试管顶部弧面
  %   试管半径 0.5, 在 y=3.2 处试管顶部弧面下凹约 0.1, 用贝塞尔曲线模拟
  %   从 (0.5,0) 经 (0,-0.1) 到 (-0.5,0), 控制点 y=-0.13 产生平滑下凹
  \filldraw[fill=PaleOrange,draw=black,line cap=round,line join=round]
    (0.5,0) .. controls (0.25,-0.13) and (-0.25,-0.13) .. (-0.5,0)
    -- (-0.8,0)
    -- ([shift={(-0.5,0)}]180:0.3) arc[start angle=180,end angle=270,radius=0.3]
    -- (-.5,-.3) -- (2.5,-.3) -- (2.4,-.2) -- (1,0) -- (0.5,0) -- cycle;
  % 第2块: 下部夹爪
  % 注: 原 \psarcn(-.5,0){0.3}{-90}{180} 是顺时针从 -90°(=270°) 到 180°, 跨度 90° (四分之一圆弧)
  % TikZ arc 当 end<start 时顺时针, 需用 end=-180 让其走短弧 (避免四分之三圆弧)
  \filldraw[fill=PaleOrange,draw=black,line cap=round,line join=round]
    (-0.8,0) -- (-0.8,-0.3)
    -- ([shift={(-0.5,-0.3)}]180:0.3) arc[start angle=180,end angle=270,radius=0.3]
    -- (-0.5,-0.6) -- (2.5,-.6) -- (2.5,-.3) -- (-.5,-.3)
    -- ([shift={(-.5,0)}]-90:0.3) arc[start angle=-90,end angle=-180,radius=0.3] -- cycle;
  % 第3块: 顶部小凸起 - 原 \psarcn 顺时针 180°->90°, 圆弧形凸起
  % 去掉 (-0.8,0.05)->(-0.75,0.05) 的多余水平线段, 直接从 (-0.5,0.05) 到圆弧起点
  % 使凸起呈纯圆弧形 (与原 pst-labo 渲染一致)
  \filldraw[fill=PaleOrange,draw=black,line cap=round,line join=round]
    (-.5,0.05) -- ([shift={(-0.5,0.05)}]180:0.25)
    arc[start angle=180,end angle=90,radius=0.25]
    -- (-.5,0.3) -- (-0.5,0.05) -- cycle;
  % 第4块: 夹臂上部
  \filldraw[fill=PaleOrange,draw=black,line cap=round,line join=round]
    (0.5,0.3) -- (4.9,0.3) -- (5,0.2) -- (1,0) -- (0.5,0) -- (0.5,0.3) -- cycle;
  % 第5块: 夹臂下部 (Brown 棕色)
  \filldraw[fill=Brown,draw=black,line cap=round,line join=round]
    (0.5,0) -- (5,0.2) -- (5,0) -- (2,-0.2) -- (1,0) -- cycle;
  % 螺丝椭圆
  \draw[line width=1.5\pgflinewidth,line cap=round] (1,0) ellipse (0.3 and 0.1);
  % 螺丝下方的杆
  \draw[line width=1.5\pgflinewidth,line cap=round,line join=round]
        (0.7,0) -- (0.4,-.3) -- (0.4,-.6);
}

% ============================================================
% 13. 容量瓶 (Fiole jaugée)
% ============================================================
\newcommand{\tkzVolumetricFlask}{%
  \chemlabUpdateLineWidths
  \draw[rounded corners=0.1cm,line width=\chemlabLineWidth,line cap=round,line join=round]
       (-0.21,3.5) -- (-0.1,3.5) -- (-0.1,1.5)
       -- (-0.75,0) -- (0.75,0) -- (0.1,1.5) -- (0.1,3.5) -- (0.21,3.5);
  \draw[line width=\chemlabGraduationLineWidth,line cap=round] (-0.1,2.5) -- (0.1,2.5);
}

% ============================================================
% 14. 结晶皿 (Cristallisoir)
% ============================================================
\newcommand{\Cristallisoir}{%
  \chemlabUpdateLineWidths
  \begin{scope}[line width=\chemlabLineWidth,line cap=round,line join=round]
  \draw (-1.9,2) -- (-2,2);
  \draw ([shift={(-2,1.95)}]90:0.05) arc[start angle=90,end angle=270,radius=0.05];
  \draw (-2,1.9) -- (-2,1);
  \draw ([shift={(-1,1)}]180:1) arc[start angle=180,end angle=270,radius=1];
  \draw (-1,0) -- (1,0);
  \draw ([shift={(1,1)}]270:1) arc[start angle=270,end angle=360,radius=1];
  \draw (2,1) -- (2,1.9);
  \draw ([shift={(2,1.95)}]270:0.05) arc[start angle=270,end angle=90,radius=0.05];
  \draw (2,2) -- (1.9,2);
  \draw (1.9,2) -- (1.9,1);
  \draw ([shift={(1,1)}]0:0.9) arc[start angle=0,end angle=-90,radius=0.9];
  \draw (1,0.1) -- (-1,0.1);
  \draw ([shift={(-1,1)}]270:0.9) arc[start angle=270,end angle=180,radius=0.9];
  \draw (-1.9,1) -- (-1.9,2);
  \end{scope}
}

\newcommand{\InterieurCristallisoir}{%
  (1.9,2) -- (1.9,1)
  -- ([shift={(1,1)}]0:0.9) arc[start angle=0,end angle=-90,radius=0.9]
  -- (1,0.1) -- (-1,0.1)
  -- ([shift={(-1,1)}]270:0.9) arc[start angle=270,end angle=180,radius=0.9]
  -- (-1.9,1) -- (-1.9,2) -- cycle
}

\newcommand{\tkzCrystallizingDish}{%
  \begin{scope}
    \clip \InterieurCristallisoir;
    \fill[LightBlue1] (-3,-1) rectangle (3,1.5);
  \end{scope}
  \Cristallisoir
}

% ============================================================
% 15. 倒置试管 (TubeRenverse)
% ============================================================
\newcommand{\tkzInvertedTube}{%
  \chemlabUpdateLineWidths
  \begin{scope}[line width=\chemlabEllipseLineWidth,line cap=round,line join=round]
    \fill[LightBlue1] (-0.5,0) rectangle (0.5,5.5);
    % 合并左侧 + 顶部圆弧 + 右侧为单条路径
    % pst-labo: \psarcn(0,6){0.5}{180}{0} = 从180°顺时针到0°, 经过顶部90°
    % TikZ: arc[start=180, end=0] = 从180°逆时针(减小)到0°, 经过顶部90°
    \draw (-0.5,0) -- (-0.5,6)
          arc[start angle=180,end angle=0,radius=0.5]
          -- (0.5,0);
  \end{scope}
}

% ============================================================
% 16. 试管架 (ReposeTube)
% ============================================================
\newcommand{\tkzTestTubeRack}{%
  \begin{scope}[line width=0.1cm,draw=lightgray]
    \draw ([shift={(0,-0.25)}]100:1) arc[start angle=100,end angle=140,radius=1];
    \draw ([shift={(0,-.25)}]150:1) arc[start angle=150,end angle=160,radius=1];
    \draw ([shift={(0,-0.25)}]20:1) arc[start angle=20,end angle=80,radius=1];
  \end{scope}
}

% ============================================================
% 17. 烧杯边缘 + 刻度 (RebordBecher)
% ============================================================
\newcommand{\tkzBeakerRim}{%
  \chemlabUpdateLineWidths
  \draw[line width=\chemlabLineWidth] (0,2.55) ellipse (1.1 and .1);
  % 主刻度 : 粗细为细刻度线的 1.5 倍 (原 pst-labo: 0.03cm = 1.5 × 0.02cm)
  % 用 \pgfmathsetmacro 预算尺寸, 避免 \dimexpr 在 pgfkeys 中的兼容问题
  \pgfmathsetlengthmacro{\chemlab@majorGrad}{1.5*\chemlabGraduationLineWidth}%
  \foreach \nDiv in {0.5,1.0,...,2.0} {%
    \draw[line width=\chemlab@majorGrad] (-.75,\nDiv) -- (-.4,\nDiv);}
  % 细分刻度
  \foreach \nSubDiv in {0.5,0.6,...,2.0} {%
    \draw[line width=\chemlabGraduationLineWidth] (-0.75,\nSubDiv) -- (-.5,\nSubDiv);}
}

% ============================================================
% 18. 烧杯 (Becher)
% ============================================================
% 注 : 原 PSTricks 使用 \psline[linearc=0.5](-1,2.5)(-1,0)(1,0)(1,2.5)
% TikZ 的 rounded corners 与 PSTricks linearc 语义略有不同,
% 因此改用显式圆弧精确模拟 PSTricks linearc=0.5 的行为 :
% 左下角圆弧 : 圆心 (-0.5,0.5), 半径 0.5, 从 180° 到 270°
% 右下角圆弧 : 圆心 (0.5,0.5),  半径 0.5, 从 270° 到 360°
\newcommand{\tkzBeakerBody}{%
  \chemlabUpdateLineWidths
  \draw[line width=\chemlabLineWidth,line cap=round,line join=round]
       (-1,2.5) -- (-1,0.5)
       -- ([shift={(-0.5,0.5)}]180:0.5) arc[start angle=180,end angle=270,radius=0.5]
       -- (0.5,0)
       -- ([shift={(0.5,0.5)}]270:0.5) arc[start angle=270,end angle=360,radius=0.5]
       -- (1,2.5);
}

% ============================================================
% 19. 锥形瓶 (Erlenmeyer)
% ============================================================
% 注 : 原 PSTricks 使用 \psline[linearc=0.3](-0.5,3.5)(-0.5,2.5)(-1.5,0)(1.5,0)(0.5,2.5)(0.5,3.5)
% TikZ 的 rounded corners 与 PSTricks linearc 语义不同,
% 因此改用显式圆弧精确模拟 linearc=0.3 (由 scripts/erlen_arcs5.py 计算) :
%   肩部 (凹转角, 内侧角 21.8°) : 半径 0.3 太大无法相切 (切点距 1.558 > 颈部线段长 1.0), 保留尖角
%   左下 (凸转角, 内侧角 68.2°) : center=(-1.0569,0.3), r=0.3, 158.20° -> 270° (ccw)
%   右下 (凸转角, 内侧角 68.2°) : center=(1.0569,0.3),  r=0.3, -90° -> 21.80° (ccw)
\newcommand{\tkzErlenmeyerBody}{%
  \chemlabUpdateLineWidths
  \draw[line width=\chemlabLineWidth,line cap=round,line join=round]
       (-0.5,3.5) -- (-0.5,2.5)
       -- (-1.3354,0.4114)
       -- ([shift={(-1.0569,0.3)}]158.20:0.3) arc[start angle=158.20,end angle=270,radius=0.3]
       -- (1.0569,0)
       -- ([shift={(1.0569,0.3)}]-90:0.3) arc[start angle=-90,end angle=21.80,radius=0.3]
       -- (0.5,2.5) -- (0.5,3.5);
}

% ============================================================
% 20. 试剂瓶 (Flacon)
% ============================================================
\newcommand{\tkzReagentBottleBody}{%
  \chemlabUpdateLineWidths\begin{scope}[line width=\chemlabLineWidth,line cap=round,line join=round]
    % 左侧瓶口上角小弧 (psarcn 90->0, 顺时针)
    \draw ([shift={(-0.6,3.4)}]90:0.1) arc[start angle=90,end angle=0,radius=0.1];
    % 左侧颈部直线
    \draw (-0.5,3.4) -- (-0.5,3.1);
    % 左侧瓶口下角凹弧 (psarcn 0->-90, 顺时针; TikZ end<start 自动顺时针)
    \draw ([shift={(-0.6,3.1)}]0:0.1) arc[start angle=0,end angle=-90,radius=0.1];
    % 左侧肩部大凸弧 (psarc 90->180, 逆时针)
    \draw ([shift={(-0.6,2.6)}]90:0.4) arc[start angle=90,end angle=180,radius=0.4];
    % 左侧瓶身直线
    \draw (-1,2.6) -- (-1,0.2);
    % 左侧底部圆角 (psarc 180->270, 逆时针)
    \draw ([shift={(-.8,0.2)}]180:0.2) arc[start angle=180,end angle=270,radius=0.2];
    % 底部
    \draw (-.8,0) -- (.8,0);
    % 右侧底部圆角 (psarc -90->0, 逆时针)
    \draw ([shift={(.8,0.2)}]-90:0.2) arc[start angle=-90,end angle=0,radius=0.2];
    % 右侧瓶身直线
    \draw (1,0.2) -- (1,2.6);
    % 右侧肩部大凸弧 (psarc 0->90, 逆时针)
    \draw ([shift={(0.6,2.6)}]0:0.4) arc[start angle=0,end angle=90,radius=0.4];
    % 右侧瓶口下角凹弧 (psarcn -90->180, 顺时针短弧; TikZ 需 end=-180 避免 270° 长弧)
    \draw ([shift={(0.6,3.1)}]-90:0.1) arc[start angle=-90,end angle=-180,radius=0.1];
    % 右侧颈部直线
    \draw (0.5,3.1) -- (0.5,3.4);
    % 右侧瓶口上角小弧 (psarcn 180->90, 顺时针)
    \draw ([shift={(0.6,3.4)}]180:0.1) arc[start angle=180,end angle=90,radius=0.1];
  \end{scope}
}

% ============================================================
% 21. 玻璃杯 (Verre) - F. Bréal
% ============================================================
\newcommand{\tkzWatchGlassBody}{%
  \chemlabUpdateLineWidths\begin{scope}[line width=\chemlabLineWidth,line cap=round,line join=round]
    % 合并左斜线 + 底部圆弧 + 右斜线为单条路径, 平滑转角
    \draw (-1,3.5) -- (-0.5,1)
          arc[start angle=180,end angle=360,radius=0.5]
          -- (1,3.5);
  \end{scope}
}

\newcommand{\tkzWatchGlassBase}{%
  \chemlabUpdateLineWidths\begin{scope}[line width=\chemlabLineWidth,rounded corners=0.03cm,line cap=round,line join=round]
    \draw (-0.5,1) -- (-0.34,0.2) -- (-1.2,0) -- (1.2,0) -- (0.34,0.2) -- (0.5,1);
  \end{scope}
}

% ============================================================
% 22. 漏斗 (Entonnoir)
% ============================================================
\newcommand{\tkzFunnelOutline}{%
  \begin{scope}[line cap=round,line join=round]
    \draw (-0.07,0) -- (-0.07,1) -- (-1,2.5);
    \draw (1,2.5) -- (0.07,1) -- (0.07,0.1);
  \end{scope}
}

\newcommand{\tkzFunnelBody}{%
  \begin{scope}
    \clip (-0.07,0) -- (-0.07,1) -- (-1,2.5)
          -- (1,2.5) -- (0.07,1) -- (0.07,0.1) -- cycle;
    \fill[white] (-2,-1) rectangle (2,3);
  \end{scope}
  \begin{scope}[rounded corners=0.25cm]
    \draw (-0.07,0) -- (-0.07,1) -- (-1,2.5);
    \draw (1,2.5) -- (0.07,1) -- (0.07,0.1);
  \end{scope}
  \draw (0,2.5) ellipse (1 and 0.1);
}

% ============================================================
% 23. 球形冷凝器 (Réfrigérant à boulles)
% ============================================================
\newcommand{\tkzCondenserStem}{%
  % 整个 TigeBoulles 作为一个连续路径 (匹配 pst-labo \pscustom)
  % pst-labo: \pscustom[linestyle=none]{...\fill} 填充 + 单独描边 (无闭合线)
  % TikZ: 分开 \fill (带 cycle 闭合填充) 和 \draw (无 cycle, 顶部敞口)
  % 顶部是敞口的 — 两端弧之间没有连线 (匹配 pst-labo)
  %
  % 填充: 带闭合线 (fill 需要 cycle)
  \fill[white]
    % 顶部圆弧与颈部 - 左侧
    ({-0.35+0.1*cos(90)},{7.5+0.1*sin(90)}) arc[start angle=90,end angle=0,radius=0.1]
    -- (-.25,7.5) -- (-.25,6.6)
    -- ({-.15+0.1*cos(180)},{6.6+0.1*sin(180)}) arc[start angle=180,end angle=270,radius=0.1]
    -- ({-.15+0.1*cos(90)},{6.4+0.1*sin(90)}) arc[start angle=90,end angle=0,radius=0.1]
    -- (-.05,6.4) -- (-0.05,6)
    % 左侧 7 个球状突起 (中心 x=-0.05, 半径 0.25, arc 90→270 向左凸)
    (-0.05,6.0) arc[start angle=90,end angle=270,radius=0.25] -- (-0.05,5.50) -- (-0.05,5.25)
    arc[start angle=90,end angle=270,radius=0.25] -- (-0.05,4.75) -- (-0.05,4.50)
    arc[start angle=90,end angle=270,radius=0.25] -- (-0.05,4.00) -- (-0.05,3.75)
    arc[start angle=90,end angle=270,radius=0.25] -- (-0.05,3.25) -- (-0.05,3.00)
    arc[start angle=90,end angle=270,radius=0.25] -- (-0.05,2.50) -- (-0.05,2.25)
    arc[start angle=90,end angle=270,radius=0.25] -- (-0.05,1.75) -- (-0.05,1.50)
    arc[start angle=90,end angle=270,radius=0.25] -- (-0.05,1.00) -- (-0.05,0.75)
    % 左侧底部到下管口
    -- (-0.05,-1.5) -- (0.05,-1.5) -- (0.05,1)
    % 右侧 7 个球状突起 (中心 x=0.05, 半径 0.25, arc 270→450 向右凸)
    arc[start angle=270,end angle=450,radius=0.25] -- (0.05,1.50) -- (0.05,1.75)
    arc[start angle=270,end angle=450,radius=0.25] -- (0.05,2.25) -- (0.05,2.50)
    arc[start angle=270,end angle=450,radius=0.25] -- (0.05,3.00) -- (0.05,3.25)
    arc[start angle=270,end angle=450,radius=0.25] -- (0.05,3.75) -- (0.05,4.00)
    arc[start angle=270,end angle=450,radius=0.25] -- (0.05,4.50) -- (0.05,4.75)
    arc[start angle=270,end angle=450,radius=0.25] -- (0.05,5.25) -- (0.05,5.50)
    arc[start angle=270,end angle=450,radius=0.25] -- (0.05,6.00) -- (0.05,6.40)
    % 顶部圆弧与颈部 - 右侧
    -- ({0.15+0.1*cos(180)},{6.4+0.1*sin(180)}) arc[start angle=180,end angle=90,radius=0.1]
    -- ({0.15+0.1*cos(-90)},{6.6+0.1*sin(-90)}) arc[start angle=-90,end angle=0,radius=0.1]
    -- (0.25,6.6) -- (0.25,7.5)
    -- ({0.35+0.1*cos(180)},{7.5+0.1*sin(180)}) arc[start angle=180,end angle=90,radius=0.1]
    -- cycle;
  %
  % 描边: 底部敞口 (匹配 pst-labo: 左右侧分别 \psline, 无底部连线)
  \draw
    % 顶部圆弧与颈部 - 左侧
    ({-0.35+0.1*cos(90)},{7.5+0.1*sin(90)}) arc[start angle=90,end angle=0,radius=0.1]
    -- (-.25,7.5) -- (-.25,6.6)
    -- ({-.15+0.1*cos(180)},{6.6+0.1*sin(180)}) arc[start angle=180,end angle=270,radius=0.1]
    -- ({-.15+0.1*cos(90)},{6.4+0.1*sin(90)}) arc[start angle=90,end angle=0,radius=0.1]
    -- (-.05,6.4) -- (-0.05,6)
    % 左侧 7 个球状突起
    (-0.05,6.0) arc[start angle=90,end angle=270,radius=0.25] -- (-0.05,5.50) -- (-0.05,5.25)
    arc[start angle=90,end angle=270,radius=0.25] -- (-0.05,4.75) -- (-0.05,4.50)
    arc[start angle=90,end angle=270,radius=0.25] -- (-0.05,4.00) -- (-0.05,3.75)
    arc[start angle=90,end angle=270,radius=0.25] -- (-0.05,3.25) -- (-0.05,3.00)
    arc[start angle=90,end angle=270,radius=0.25] -- (-0.05,2.50) -- (-0.05,2.25)
    arc[start angle=90,end angle=270,radius=0.25] -- (-0.05,1.75) -- (-0.05,1.50)
    arc[start angle=90,end angle=270,radius=0.25] -- (-0.05,1.00) -- (-0.05,0.75)
    -- (-0.05,-1.5);
  % 右侧 (单独 \draw, 底部不连接左侧)
  \draw
    (0.05,-1.5) -- (0.05,1)
    % 右侧 7 个球状突起
    arc[start angle=270,end angle=450,radius=0.25] -- (0.05,1.50) -- (0.05,1.75)
    arc[start angle=270,end angle=450,radius=0.25] -- (0.05,2.25) -- (0.05,2.50)
    arc[start angle=270,end angle=450,radius=0.25] -- (0.05,3.00) -- (0.05,3.25)
    arc[start angle=270,end angle=450,radius=0.25] -- (0.05,3.75) -- (0.05,4.00)
    arc[start angle=270,end angle=450,radius=0.25] -- (0.05,4.50) -- (0.05,4.75)
    arc[start angle=270,end angle=450,radius=0.25] -- (0.05,5.25) -- (0.05,5.50)
    arc[start angle=270,end angle=450,radius=0.25] -- (0.05,6.00) -- (0.05,6.40)
    -- ({0.15+0.1*cos(180)},{6.4+0.1*sin(180)}) arc[start angle=180,end angle=90,radius=0.1]
    -- ({0.15+0.1*cos(-90)},{6.6+0.1*sin(-90)}) arc[start angle=-90,end angle=0,radius=0.1]
    -- (0.25,6.6) -- (0.25,7.5)
    -- ({0.35+0.1*cos(180)},{7.5+0.1*sin(180)}) arc[start angle=180,end angle=90,radius=0.1];
  % 底部液滴
  \begin{scope}[shift={(0,-1.6)}]\tkzDrop\end{scope}
}

\newcommand{\tkzCondenserBody}{%
  % 渐变填充: pst-labo gradbegin=cyan (底), gradend=PaleOrange (顶), gradmidpoint=0
  % TikZ: bottom color=cyan (底), top color=PaleOrange (顶)
  \shade[top color=PaleOrange,bottom color=cyan]
    ([shift={(0,5.75)}]0:0.5) arc[start angle=0,end angle=180,radius=0.5]
    -- (-0.5,5.75) -- (-0.5,5.5) -- (-1.5,5.5)
    -- (-1.5,5.3) -- (-.5,5.3) -- (-.5,1.7) -- (-1.5,1.7)
    -- (-1.5,1.5) -- (-.5,1.5) -- (-0.5,1)
    -- ([shift={(0,1)}]180:0.5) arc[start angle=180,end angle=360,radius=0.5]
    -- (0.5,1) -- (0.5,5.75) -- cycle;
  % 描边
  \begin{scope}[line width=0.88pt]
    \draw ([shift={(0,5.75)}]0:0.5) arc[start angle=0,end angle=180,radius=0.5];
    \draw (-0.5,5.75) -- (-0.5,5.5) -- (-1.5,5.5);
    \draw (-0.5,5.75) -- (-0.5,5.5) -- (-1.5,5.5);
    \draw (-1.5,5.3) -- (-.5,5.3) -- (-.5,1.7) -- (-1.5,1.7);
    \draw (-1.5,1.5) -- (-.5,1.5) -- (-0.5,1);
    \draw ([shift={(0,1)}]180:0.5) arc[start angle=180,end angle=360,radius=0.5];
    \draw (0.5,1) -- (0.5,5.75);
  \end{scope}
  \draw[->,cyan] (-2.5,1.6) -- (-1.6,1.6);
  \draw[<-,PaleOrange] (-2.5,5.4) -- (-1.6,5.4);
}

\newcommand{\tkzCondenserBulbs}{%
  \begin{scope}[shift={(0,0.2)}]
    \tkzCondenserBody
    \tkzCondenserStem
  \end{scope}
  \tkzTestTubeStopper
  \fill[white,rounded corners=0.025cm]
       (-0.05,0.6) -- (-0.05,0.5) -- (0.05,0.5) -- (0.05,0.6) -- cycle;
}

% ============================================================
% 24. 加热球 (Chauffe-Ballon)
% ============================================================
\newcommand{\tkzHeaterRoundFlask}{%
  \filldraw[fill=white,draw=black] (-2,-1) rectangle (2,1);
  \draw (-1.25,0) circle (0.4);
  \begin{scope}[shift={(-1.25,0)}]
    \foreach \iA in {0,30,...,180} {\draw (\iA:0.45) -- (\iA:0.6);}
    \draw[->] (0,0) -- (60:0.3);
  \end{scope}
  % 按钮区域: pst-labo \psframe[doubleline=true] 产生双线边框
  % PSTricks doubleline: 外框在指定坐标, 内框向内偏移
  % 实测: 外框到内框偏移约 0.13 units (含线宽+间隔, at scale=0.5)
  % 外框 (在指定坐标处)
  \draw[black,line width=0.8pt] (1,-.5) rectangle (1.5,0.5);
  % 内框 (向内偏移, 填充灰色)
  \filldraw[fill=gray,draw=black,line width=0.8pt] ({1+0.14},{-.5+0.14}) rectangle ({1.5-0.14},{0.5-0.14});
  % 红色指示灯
  \fill[red] (1.25,0.7) circle (0.1);
}

% ============================================================
% 25. 回流烧瓶 (Ballon à reflux)
% ============================================================
\newcommand{\tkzRoundFlaskReflux}{%
  \begin{scope}
    \clip \tkzRoundFlaskSimplePath;   % 使用路径版本形成闭合 clip 区域
    \shade[top color=cyan,bottom color=white,draw=cyan]
      plot[smooth,tension=0.5] coordinates {%
        (-1.5,1.5) (-1.375,1.55) (-1.25,1.5) (-1.125,1.45)
        (-1,1.5) (-.875,1.55) (-.75,1.5) (-.625,1.45)
        (-.5,1.5) (-.375,1.55) (-.25,1.5) (-0.125,1.45)
        (0,1.5) (0.125,1.55) (.25,1.5) (0.375,1.45)
        (0.5,1.5) (0.625,1.55) (0.75,1.5) (0.875,1.45)
        (1,1.5) (1.125,1.55) (1.25,1.5) (1.375,1.45) (1.5,1.5)}
      -- (1.5,1.5) -- (1.5,0) -- (-1.5,0) -- (-1.5,1.5) -- cycle;
  \end{scope}
  \tkzRoundFlaskSimple
}

% ============================================================
% 26. 分液漏斗 (Ampoule à décanter / Separate Funnel)
% ============================================================
% 相 : 水相 / 有机相
\definecolor{AqueousPhaseColor}{cmyk}{0.215686,0.062745,0,0}
\definecolor{OrganicPhaseColor}{cmyk}{0.223529,0.094118,0.501961,0}

% pst-labo: \tkzSeparatoryFunnel 设 linewidth=0.04 (用户单位, 随 unit 缩放)
% TikZ: 用 \chemlabLineWidth (基础 0.053cm × scale) 近似, 匹配其他器皿线宽
\newcommand{\tkzSeparatoryFunnelBody}{%
  \chemlabUpdateLineWidths
  \draw[line width=\chemlabLineWidth,line cap=round,line join=round]
    (-0.65,15) -- (-0.55,13.75)
    .. controls (-1.55,13.54) and ( -1.92,12.7) .. (-1.92,11.65)
    .. controls (-1.92,10.6) and (-0.3,9.55) .. (-0.3,8.5)
    -- (-0.3,2)
    .. controls (-0.3,1.5) and (-0.14,1.5) .. (-0.14,1)
    -- (-0.14,-1)
    -- (0.14,-1)
    -- (0.14,1)
    .. controls (0.14,1.5) and (0.3,1.5) .. (0.3,2)
    -- (0.3,8.5)
    .. controls (0.3,9.55) and (1.92,10.6) .. (1.92,11.65)
    .. controls ( 1.92,12.7) and (1.55,13.54) .. (0.55,13.75)
    -- (0.65,15);
}

% 用于 clip 的轮廓
\newcommand{\tkzSeparatoryFunnelOutlineClip}{%
  (-0.65,15) -- (-0.55,13.75)
  .. controls (-1.55,13.54) and ( -1.92,12.7) .. (-1.92,11.65)
  .. controls (-1.92,10.6) and (-0.3,9.55) .. (-0.3,8.5)
  -- (-0.3,2)
  .. controls (-0.3,1.5) and (-0.14,1.5) .. (-0.14,1)
  -- (-0.14,-1) -- (0.14,-1) -- (0.14,1)
  .. controls (0.14,1.5) and (0.3,1.5) .. (0.3,2)
  -- (0.3,8.5)
  .. controls (0.3,9.55) and (1.92,10.6) .. (1.92,11.65)
  .. controls (1.92,12.7) and (1.55,13.54) .. (0.55,13.75)
  -- (0.65,15) -- cycle
}

% 顶盖 (Bouchon) - 继承 linewidth=0.04
\newcommand{\tkzStopper}{%
  \chemlabUpdateLineWidths
  \filldraw[fill=gray!10,draw=gray,line width=\chemlabLineWidth,line cap=round,line join=round]
    (-0.67,15.25) -- (0.67,15.25) -- (0.57,14) -- (-0.57,14) -- cycle;
  \filldraw[fill=gray!5,draw=gray,line width=\chemlabLineWidth,line cap=round,line join=round]
    ([shift={(0.555,15.365)}]-90:0.1145) arc[start angle=-90,end angle=-180,radius=0.1145]
    -- ([shift={(0,15.92)}]-45:0.67) arc[start angle=-45,end angle=225,radius=0.67]
    -- ([shift={(-0.555,15.365)}]0:0.1145) arc[start angle=0,end angle=-90,radius=0.1145]
    -- cycle;
}

% 开启的旋塞 (TapOpen)
% pst-labo: \psframe[fillstyle=solid] 默认 fillcolor=white (PSTricks 默认)
% TikZ: \fill 默认黑色, 必须显式 fill=white
\newcommand{\tkzTapOpen}{%
  \chemlabUpdateLineWidths
  \filldraw[fill=white,line width=\chemlabLineWidth] (-.2,0.5) rectangle (.2,1);
  \draw[line width=\chemlabLineWidth] (-.2,0.4) rectangle (-.3,1.1);
  \draw[line width=\chemlabLineWidth] (.2,0.4) rectangle (.3,1.1);
  \draw[line width=\chemlabLineWidth] (.3,0.6) -- (.4,0.6) -- (.5,0.3) -- (.9,0.3) -- (.9,1.2) -- (.5,1.2)
        -- (.4,0.9) -- (.3,0.9) -- cycle;
  \fill[cyan!50] (-.6,0.4) -- (-.85,0.6) -- (-.85,0.9) -- (-.6,1.1) -- (-.3,1.1) -- (-.3,0.4) -- cycle;
  \draw[line width=\chemlabLineWidth] (-0.6,0.4) -- (-0.6,1.1);
  \draw[line width=\chemlabLineWidth] (0.5,0.3) -- (0.5,1.2);
}

% 关闭的旋塞 (TapClosed)
% pst-labo: \psframe[fillstyle=solid] 默认 fillcolor=white
\newcommand{\tkzTapClosed}{%
  \chemlabUpdateLineWidths
  \filldraw[fill=white,line width=\chemlabLineWidth] (-.2,0.5) rectangle (.2,1);
  \draw[line width=\chemlabLineWidth] (-.2,0.4) rectangle (-.3,1.1);
  \draw[line width=\chemlabLineWidth] (.2,0.4) rectangle (.3,1.1);
  \draw[line width=\chemlabLineWidth] ([shift={(0.35,0.6)}]180:0.05) arc[start angle=180,end angle=90,radius=0.05]
        -- (0.4,0.65) -- (0.9,0.65) -- (0.9,0.85) -- (0.35,0.85)
        -- ([shift={(0.35,0.9)}]-90:0.05) arc[start angle=-90,end angle=-180,radius=0.05]
        -- cycle;
  \draw[line width=\chemlabLineWidth] (0.5,0.65) -- (0.5,0.85);
  \fill[cyan!50] (-.6,0.4) -- (-.85,0.6) -- (-.85,0.9) -- (-.6,1.1) -- (-.3,1.1) -- (-.3,0.4) -- cycle;
  \draw[line width=\chemlabLineWidth] (-0.6,0.4) -- (-0.6,1.1);
}

% \tkzSeparatoryFunnel  (主调用)
% 可选参数 : stopper=true/false, TapOpen=true/false, liquidLevel1, liquidLevel2
\newcommand{\tkzSeparatoryFunnel}[1][]{%
  \begingroup
  \chemlabUpdateLineWidths
  \pgfkeys{
    /tkzSeparatoryFunnel/.cd,
    stopper/.is if=TkzTestTubeStopper,
    TapOpen/.is if=TkzTapOpen,
    liquidLevel1/.store in=\tkzTestTubeLiquidLevelI,
    liquidLevel2/.store in=\tkzTestTubeLiquidLevelII,
    #1}
  % 默认值
  \ifdefined\tkzTestTubeLiquidLevelI\else\def\tkzTestTubeLiquidLevelI{11}\fi
  \ifdefined\tkzTestTubeLiquidLevelII\else\def\tkzTestTubeLiquidLevelII{12}\fi
  % 水相
  \begin{scope}
    \clip \tkzSeparatoryFunnelOutlineClip;
    \fill[AqueousPhaseColor] (-2,1) rectangle (2,\tkzTestTubeLiquidLevelI);
  \end{scope}
  % 有机相
  \begin{scope}
    \clip \tkzSeparatoryFunnelOutlineClip;
    \fill[OrganicPhaseColor] (-2,\tkzTestTubeLiquidLevelI)
          rectangle (2,\tkzTestTubeLiquidLevelII);
  \end{scope}
  \ifTkzTestTubeStopper\tkzStopper\fi
  \tkzSeparatoryFunnelBody
  \ifTkzTapOpen\tkzTapOpen\else\tkzTapClosed\fi
  % 出口管口 (pst-labo: \psellipse[fillstyle=solid,fillcolor=white] 既有白色填充又有黑色轮廓)
  \filldraw[fill=white,line width=\chemlabLineWidth] (0,-1) ellipse (0.15 and 0.25);
  % 瓶口圆点 (填充, 无需线宽)
  \fill[darkgray] (-0.675,15) circle (0.03);
  \fill[darkgray] (0.675,15) circle (0.03);
  \draw[darkgray,line width=\chemlabLineWidth] (-0.65,15) -- (0.65,15);
  % 颈与瓶身接合
  \draw[gray,line width=\chemlabLineWidth] (-0.55,13.75) -- (0.55,13.75);
  \endgroup
}

% ============================================================
% 27. 支架 (pstSupport)
% ============================================================
% pst-labo: \psset{dimen=m,linewidth=0.04} — linewidth=0.04 是用户单位, 随 unit 缩放
% TikZ: 用 \chemlabLineWidth (基础 0.053cm × scale), 与其他器皿一致
\newcommand{\tkzStand}{%
  \chemlabUpdateLineWidths
  \begin{scope}[line width=\chemlabLineWidth]
    \filldraw[fill=woodColor,draw=black,line join=round]
        (-2,-6) rectangle (5,-5);
    \filldraw[fill=woodColor,draw=black,line join=round]
        (2.75,2) rectangle (3.25,16);
    \filldraw[fill=woodColor,draw=black,line join=round]
        ([shift={(2.75,1.5)}]90:0.5) arc[start angle=90,end angle=180,radius=0.5]
        -- (2.25,1.5) -- (2.25,-5) -- (3.75,-5) -- (3.75,1.5)
        -- ([shift={(3.25,1.5)}]0:0.5) arc[start angle=0,end angle=90,radius=0.5]
        -- cycle;
    \draw (2,11.5) -- (1.55,10.5);
    \filldraw[fill=woodColor,draw=black,line join=round]
        (-2.5,11.5) -- (-2,11.5) -- (-1.55,10.5) -- (-2.5,10.5) -- cycle;
    \filldraw[fill=woodColor,draw=black,line join=round]
        (4,11.5) -- (2,11.5) -- (1.55,10.5) -- (4,10.5) -- cycle;
    \draw (-2.5,10.5) rectangle (4,11.5);
    % 旋钮: pst-labo \pscustom[style=woodColor]{...} 填充会自动闭合, 但描边只画实际路径
    % (不画对角闭合边). TikZ: \fill 自动闭合填充, \draw 画开放路径 (无 -- cycle).
    \fill[woodColor]
        (4,11.2) -- (4.5,11.2)
        -- ([shift={(5,11)}]158.2:0.5385) arc[start angle=158.2,end angle=-158.2,radius=0.5385]
        -- (4.5,10.8) -- (4,10.8) -- cycle;
    \draw
        (4,11.2) -- (4.5,11.2)
        -- ([shift={(5,11)}]158.2:0.5385) arc[start angle=158.2,end angle=-158.2,radius=0.5385]
        -- (4.5,10.8) -- (4,10.8);
    \draw (4.5,10.8) -- (4.5,11.2);
  \end{scope}
}

\endinput
    % 形状宏库 (必须先加载)
%    % ============================================================
%  tikz-chemlab.tex
%  化学实验器材 TikZ 绘图包 - 主实现 (参数系统 + 高级命令)
%  原作者 : D. Girou / C. Jorsen / M. Luque / H. Voss
%  TikZ 改写 : Aihua Wang(王爱华) (原 PSTricks 包 pst-labo)
% ============================================================
%
%  依赖 :
%    \usepackage{tikz}
%    \usetikzlibrary{arrows.meta,calc,patterns,patterns.meta,shadings}
%    % ============================================================
%  tikz-chemlabObj.tex
%  化学实验器材 TikZ 绘图包 - 形状宏库
%  原作者 : D. Girou / C. Jorsen / M. Luque / H. Voss
%  TikZ 改写 : Aihua Wang(王爱华) (原 PSTricks 包 pst-laboObj)
% ============================================================
%
%  使用说明 :
%    \usepackage{tikz}
%    \usetikzlibrary{arrows.meta,calc,patterns,patterns.meta,shadings}
%    % ============================================================
%  tikz-chemlabObj.tex
%  化学实验器材 TikZ 绘图包 - 形状宏库
%  原作者 : D. Girou / C. Jorsen / M. Luque / H. Voss
%  TikZ 改写 : Aihua Wang(王爱华) (原 PSTricks 包 pst-laboObj)
% ============================================================
%
%  使用说明 :
%    \usepackage{tikz}
%    \usetikzlibrary{arrows.meta,calc,patterns,patterns.meta,shadings}
%    \input{tikz-chemlabObj}
%    \begin{tikzpicture}
%       \BuretteCorps  ... 等等
%    \end{tikzpicture}
%
%  映射规则 :
%    \psline(A)(B)(C)              -> \draw (A) -- (B) -- (C);
%    \psline{->}(A)(B)             -> \draw[->] (A) -- (B);
%    \pspolygon(A)(B)(C)           -> \filldraw (A) -- (B) -- (C) -- cycle;
%    \psframe(x1,y1)(x2,y2)        -> \draw (x1,y1) rectangle (x2,y2);
%    \psframe*[linecolor=c]        -> \fill[c] (..) rectangle (..);
%    \psframe[framearc=r]          -> \draw[rounded corners=r] (...) rectangle (...);
%    \pscircle(c){r}               -> \draw (c) circle (r);
%    \pscircle*[linecolor=c](c){r} -> \fill[c] (c) circle (r);
%    \psellipse(c)(a,b)            -> \draw (c) ellipse (a and b);
%    \psarc(c){r}{a1}{a2}          -> \draw ([shift={(c)}]a1:r) arc[start angle=a1,end angle=a2,radius=r];
%    \psarcn(c){r}{a1}{a2}         -> \draw ([shift={(c)}]a1:r) arc[start angle=a1,end angle=a2,radius=r];   (注意 : PSTricks 中 arcn 是逆时针 / 角度递减)
%    \psbezier(A)(B)(C)(D)         -> \draw (A) .. controls (B) and (C) .. (D);
%    \pscurve(P1)(P2)...(Pn)       -> \draw plot[smooth,tension=0.5] coordinates {(P1) (P2) ... (Pn)};
%    \pscustom{ ... \fill}         -> \filldraw / \path[fill] (整段路径) -- cycle;
%    \rput(x,y){...}               -> \begin{scope}[shift={(x,y)}] ... \end{scope}
%    \rput{angle}(x,y){...}        -> \begin{scope}[shift={(x,y)},rotate=angle] ... \end{scope}
%    \multido{\i=a+b}{n}{...}      -> \foreach \i in {a,a+b,...,a+(n-1)*b} {...}
%    \psclip{<path>} ... \endpsclip-> \begin{scope} \clip <path>; ... \end{scope}
%    fillstyle=solid,fillcolor=c   -> fill=c
%    fillstyle=gradient,gradbegin=..,gradend=.. -> \shade[left color=..,right color=..]
%    fillstyle=hlines              -> pattern=horizontal lines
%    doubleline=true               -> double
%    linestyle=dashed              -> dashed
%    framearc=r                    -> rounded corners=r (近似)
%    linearc=r                     -> rounded corners=r
% ============================================================

\NeedsTeXFormat{LaTeX2e}
\ProvidesFile{tikz-chemlabObj.tex}[2026/08/03 v1.0.0]

\RequirePackage{tikz}
\usetikzlibrary{arrows.meta,calc,patterns,patterns.meta,shadings,decorations.pathmorphing}

% ------------------------------------------------------------
% 颜色定义 (对应原包中使用的颜色)
% ------------------------------------------------------------
\definecolor{LightGray}{gray}{0.80}      % 浅灰
\definecolor{VeryLightGray}{cmyk}{0,0,0,0.1}  % 极浅灰 (label 标签背景)
\definecolor{LightBlue}{rgb}{0.68,0.85,0.90}
\definecolor{LightBlue1}{rgb}{0.60,0.80,1.00}
\definecolor{VeryPaleOrange}{cmyk}{0,0.1,0.3,0}  % 滴定管液体默认色
\definecolor{PaleOrange}{rgb}{1.00,0.78,0.55}
\definecolor{Brown}{rgb}{0.50,0.25,0.00}
\definecolor{woodColor}{rgb}{0.957,0.84,0.69}
\definecolor{darkWoodColor}{rgb}{0.50,0.25,0.00}
\definecolor{flameAndGrill}{rgb}{0.95,0.55,0.10}  % 估算颜色
\definecolor{Beige}{rgb}{0.96,0.96,0.86}  % Champagne 液体默认色

% ------------------------------------------------------------
% 公共缺省值 (供调用方按需修改)
% ------------------------------------------------------------
\tikzset{
  Distillation/.cd,
  DistillateColor/.store in=\DistillationDistillateColor,
  DistillateColor=LightBlue1,
  MixtureStyle/.style={fill=LightBlue},
}
\newcommand{\tkzDistillationDistillateColor}{LightBlue1}
\newcommand{\tkzDistillationMixtureStyle}{LightBlue}
\newcommand{\tkzTestTubeLiquidLevelI}{11}     % liquidLevel1
\newcommand{\tkzTestTubeLiquidLevelII}{12}    % liquidLevel2
\newif\ifTkzTapOpen
\newif\ifTkzTestTubeStopper

% ============================================================
% 1. 滴定管 (Burette)
% ============================================================
\newcommand{\tkzBuretteBody}{%
  \draw (-0.5,11.5) -- (-0.25,11) -- (-0.25,3) -- (0,2.75) -- (0,0.5);
  \draw (0.2,0.7) -- (0.2,2.75) -- (0.5,3) -- (0.5,11) -- (0.75,11.5);
}

\newcommand{\tkzBuretteStopcock}{%
  % pst-labo: \psframe[fillstyle=solid] → 填充默认白色 + 黑色边框
  \filldraw[fill=white] (-.1,1.6) rectangle (.3,2.1);
  \draw (-.1,1.5) rectangle (-.2,2.2);
  \draw (.3,1.5) rectangle (.4,2.2);
  \draw (.4,1.7) -- (.5,1.7) -- (.6,1.4) -- (.8,1.4)
        -- (.8,2.3) -- (.6,2.3) -- (.5,2) -- (.4,2) -- cycle;
  % 左侧旋钮: pst-labo \pspolygon + \psframe with fillstyle=solid,fillcolor=LightBlue
  % 需要 LightBlue 填充 + 黑色外轮廓线
  \begin{scope}[fill=LightBlue,draw=black]
    \filldraw (-.5,1.5) -- (-.75,1.7) -- (-.75,2) -- (-.5,2.2)
          -- (-.3,2.2) -- (-.3,1.5) -- cycle;
    \filldraw (-.5,1.5) rectangle (-.2,2.2);
  \end{scope}
}

\newcommand{\tkzBuretteGraduation}{%
  % 滴定管刻度 (原 pst-labo 无显式线宽, 这里用细刻度线宽)
  \foreach \n in {3.0,3.32,...,10.68} { \draw[line width=\chemlabGraduationLineWidth] (-0.25,\n) -- (0,\n); }
}

% ============================================================
% 2. pH 计 (pH-mètre)
% ============================================================
\newcommand{\tkzTitrationPHMeter}{%
  % boitier
  \filldraw[fill=white,draw=black,rounded corners=0.125cm] (-1,0) rectangle (1,3);
  \fill[LightGray] (-0.75,2) rectangle (0.75,2.75);
  \node[white] at (0,2.375) {PH};
  \draw (0,1.25) circle (0.3);
  \draw[->] (0,1.25) -- (0,1.5);
  \node at (0.65,1.5) {$^\circ$C};
  \draw (-0.5,0.5) circle (0.2);
  \draw (0.5,0.5) circle (0.2);
  % électrode — 完全透明填充 (不遮挡烧杯刻度线)
  % pst-labo: \psframe[framearc=0.5](2.25,2.5)(2.75,4.5) 外框圆角
  %   顶部圆角被下方 \psframe[LightGray](sharp) 覆盖, 底部圆角被 \pscircle 覆盖
  %   实际可见的外框是直角矩形, 此处用 sharp corners 匹配 pst-labo 可见效果
  \draw (2.25,2.5) rectangle (2.75,4.5);
  % 顶部 LightGray 填充 + 描边 (sharp corners)
  \filldraw[fill=LightGray,draw=black] (2.25,4) rectangle (2.75,4.5);
  \draw (2.5,2.25) circle (0.25);
  % fil
  \draw[line width=0.05cm,LightBlue,rounded corners=0.2cm]
        (2.5,4.5) -- (2.5,5) -- (1.25,4.75) -- (0,5) -- (0,3);
}

% ============================================================
% 3. 滴定支架 (Support dosage)
% ============================================================
\newcommand{\tkzTitrationStand}{%
  % pst-labo: \psframe(-1.5,0)(1.5,1.5) + {\psset{doubleline=true} ...}
  % 线宽: pst-labo 在 \tkzTitration 中设 linewidth=0.053 (×0.4cm = 0.0212cm ≈ 0.6pt)
  \chemlabUpdateLineWidths
  % 平台: 无填充, 只有边框
  \draw[line width=\chemlabLineWidth] (-1.5,0) rectangle (1.5,1.5);
  % 红色按钮: \pscircle*[linecolor=red] → 红色填充, 双线边框也是红色 (不可见)
  \fill[red] (0,0.75) circle (0.2);
  % 回字型矩形按钮: 无填充, 双线边框
  % PSTricks doubleline: 指定坐标在内外线之间
  % 实测: 外线向外偏移 0.085 units, 内线向内偏移 0.034 units
  \draw[line width=\chemlabLineWidth] ({-1.25-0.085},{0.3-0.085}) rectangle ({-0.75+0.085},{1.2+0.085});
  \draw[line width=\chemlabLineWidth] ({-1.25+0.034},{0.3+0.034}) rectangle ({-0.75-0.034},{1.2-0.034});
  % 环形旋钮: 无填充, 双线边框
  % 外线向外偏移, 内线向内偏移
  \draw[line width=\chemlabLineWidth] (1,0.75) circle ({0.25+0.085});
  \draw[line width=\chemlabLineWidth] (1,0.75) circle ({0.25-0.034});
}

% ============================================================
% 4. 磁力搅拌子 (Aimant)
% ============================================================
\newcommand{\tkzTitrationStirBar}{%
  % pst-labo: \psset{linewidth=0.053} 在外层 \tkzTitration 中设置, 随 scale 缩放
  \chemlabUpdateLineWidths
  \draw[line width=\chemlabLineWidth,rounded corners=0.075cm,fill=white] (-0.5,0.1) rectangle (0.5,0.4);
  \draw[line width=\chemlabLineWidth,rounded corners=0.025cm,fill=white] (-0.05,0) rectangle (0.05,0.5);
}

% ============================================================
% 5. 移液管 (Pipette)
% ============================================================
\newcommand{\tkzPipetteBody}{%
  \chemlabUpdateLineWidths\begin{scope}[line width=\chemlabLineWidth,rounded corners=0.05cm]
    \draw (-0.05,0) -- (-0.15,0.2) -- (-0.15,6)
          -- (0.15,6) -- (0.15,0.2) -- (0.05,0) -- cycle;
    \fill[white] (0,6) ellipse (0.15 and 0.05);
    \foreach \Level in {0.5,1.0,...,5.0} {%
      \draw[line width=\chemlabGraduationLineWidth] (0,\Level) -- (0.1,\Level);}
  \end{scope}
}

\newcommand{\tkzLargePipetteBody}{%
  \chemlabUpdateLineWidths\begin{scope}[line width=\chemlabLineWidth,rounded corners=0.05cm]
    \draw (0,0) -- (0,1) -- (-0.23,2.5) -- (-0.23,22.5)
          arc[start angle=-80,end angle=80,radius=0.45];
  \end{scope}
}

% ============================================================
% 6. 液滴 (Goutte)
% ============================================================
% 原文件中 \pst@goutte 被重复定义了两次, 此处使用后者 (unit=0.6 的版本)
% pst-labo: \psset{unit=0.6}\psbezier[linewidth=0.017] — linewidth=0.017 是
%   0.017 用户单位 = 0.017*0.6 = 0.0102cm (随 goutte 的 LOCAL unit=0.6 缩放, 
%   不随外层 unit 变化, 因为 \psset{unit=0.6} 覆盖了外层设置)
% TikZ: line width=0.0102cm (ABSOLUTE, 匹配 pst-labo 的固定 0.0102cm)
\newcommand{\tkzDrop}{%
  \begin{scope}[scale=0.6]
    \draw[line width=0.0102cm]
      (0,0) .. controls (0.25,-.4) and (-.25,-.4) .. (0,0);
  \end{scope}
}

% ============================================================
% 7. 量筒 (Eprouvette)
% ============================================================
\newcommand{\tkzGraduatedCylinderBody}{%
  \draw (-0.5,5) -- (-0.5,0) -- (0.5,0) -- (0.5,4.9) -- (0.6,4.93);
}

\newcommand{\tkzGraduatedCylinderRim}{%
  \draw plot[smooth,tension=0.5] coordinates {%
    (0.6,-0.07) (0.498,0.008) (0.483,0.026) (0.433,0.05)
    (0.353,0.0707) (0.25,0.0866) (0.13,0.0966) (0,0.1)
    (-0.13,0.0966) (-0.25,0.0866) (-0.353,0.0707) (-0.433,0.05)
    (-0.483,0.026) (-0.5,0) (-0.483,-0.026) (-0.353,-0.0707)
    (-0.25,-0.0866) (-0.13,-0.096) (0,-0.1) (0.13,-0.096)
    (0.25,-0.0866) (0.353,-0.0707) (0.433,-0.05)
    (0.498,-0.026) (0.6,-0.07)};
}

\newcommand{\tkzGraduatedCylinderBase}{%
  \fill[white]
    ([shift={(-0.95,0.5)}]0:0.45) arc[start angle=0,end angle=-90,radius=0.45]
    -- ([shift={(-0.95,0.025)}]90:0.025) arc[start angle=90,end angle=270,radius=0.025]
    -- (-0.95,0) -- (0.95,0)
    -- ([shift={(0.95,0.025)}]-90:0.025) arc[start angle=-90,end angle=90,radius=0.025]
    -- ([shift={(0.95,0.5)}]-90:0.45) arc[start angle=-90,end angle=-180,radius=0.45]
    -- plot[smooth,tension=0.5] coordinates {(0.5,0.4) (0,0.25) (-0.5,0.4)}
    -- cycle;
  \draw
    ([shift={(-0.95,0.5)}]0:0.45) arc[start angle=0,end angle=-90,radius=0.45]
    -- ([shift={(-0.95,0.025)}]90:0.025) arc[start angle=90,end angle=270,radius=0.025]
    -- (-0.95,0) -- (0.95,0)
    -- ([shift={(0.95,0.025)}]-90:0.025) arc[start angle=-90,end angle=90,radius=0.025]
    -- ([shift={(0.95,0.5)}]-90:0.45) arc[start angle=-90,end angle=-180,radius=0.45]
    -- plot[smooth,tension=0.5] coordinates {(0.5,0.4) (0,0.25) (-0.5,0.4)};
}

\newcommand{\tkzGraduatedCylinderGraduation}{%
  % 量筒刻度 : 主刻度 1.5× 细刻度 (与烧杯一致)
  \pgfmathsetlengthmacro{\chemlab@eprouvette@major}{1.5*\chemlabGraduationLineWidth}%
  \foreach \nyD in {0.5,1.0,...,4.5} { \draw[line width=\chemlab@eprouvette@major] (-0.4,\nyD) -- (-0.1,\nyD); }
  \foreach \nyd in {0.5,0.6,...,4.4} { \draw[line width=\chemlabGraduationLineWidth] (-0.4,\nyd) -- (-0.25,\nyd); }
}

% ============================================================
% 8. 分馏装置 (Distillation fractionnée)
% ============================================================
% \pst@goutte 这里又被定义了一次, 在 TikZ 版中沿用同一个 \tkzDrop

\newcommand{\tkzColumn}{%
  % pst-labo: linearc=0.3 (×unit, 不是绝对值) → TikZ: rounded corners=0.12cm (= 0.3 × 0.4cm)
  % 尖刺 (pointeG/D) 先画并用 clip 裁剪到蒸馏柱内部, 再画轮廓线覆盖切边。
  %   pst-labo 中 linearc=0.3 (R=0.3 user units), 尖刺端点恰在圆弧切点附近, 几乎
  %   不可见。TikZ 的 rounded corners 实际 R 比标称值大 ~2×, 圆弧切点更低,
  %   尖刺端点 (在 y_col=8 和 7.9) 会突出到圆弧外侧形成多余线头。用 clip 将
  %   尖刺裁剪到柱内, 再画轮廓线覆盖切边, 消除突出线头。
  \begin{scope}
    \clip[rounded corners=0.12cm]
      (-.25,10.5) -- (-.25,8.5) -- (-.75,8) -- (-.75,2.5) -- (-.25,2) -- (-.25,0)
      -- (0.25,0) -- (0.25,2) -- (.75,2.5) -- (.75,8) -- (0.25,8.5) -- (0.25,9.36)
      -- (0.25,10.5) -- cycle;
    % 左侧尖刺 (pointeG)
    \foreach \i in {3,4,...,8} {%
      \draw[line width=0.4\pgflinewidth] (-.75,\i) -- (-.25,\i-0.5)
          -- (-.75,\i-0.1);}
    % 右侧尖刺 (pointeD)
    \foreach \i in {3,4,...,8} {%
      \draw[line width=0.4\pgflinewidth] (.75,\i) -- (.25,\i-0.5)
          -- (.75,\i-0.1);}
  \end{scope}
  \draw[rounded corners=0.12cm]
    (-.25,10.5) -- (-.25,8.5) -- (-.75,8) -- (-.75,2.5)
    -- (-.25,2) -- (-.25,0);
  \draw[rounded corners=0.12cm]
    (0.25,0) -- (0.25,2) -- (.75,2.5) -- (.75,8)
    -- (0.25,8.5) -- (0.25,9.36);
  % 上部断口 (line 3a/3b):
  %   pst-labo 源码 line 3a: \psline(0.25,8.5)(0.25,9.30) 用于填充 linearc 圆角
  %   留下的间隙。实测发现 TikZ 的 rounded corners 实际圆弧半径比标称值大
  %   ~2× (在 scale=0.6 下 rounded corners=0.3cm 给出 R≈0.62cm), 圆弧已覆盖到
  %   x=0.25 处 y_col=8.93 附近, 故 line 3a 在任何 scale/line width 下均被圆弧
  %   覆盖, 不可见。为避免在编号图等彩色调试场景中残留多余线段, 此处不再
  %   单独绘制 line 3a。
  \draw (0.25,9.5) -- (0.25,10.5);
  % 填充 / 阴影
  \begin{scope}[pattern={Lines[angle=45,distance={1pt/sqrt(2)}]},pattern color=gray]
    \fill (-.25,10) rectangle (-.05,10.5);
    \fill (.25,10)  rectangle (0.05,10.5);
  \end{scope}
  \foreach \n in {2.5,3.5,...,7.5} {%
    \begin{scope}[shift={(-.25,\n)}]\tkzDrop\end{scope}
    \begin{scope}[shift={(.25,\n)}]\tkzDrop\end{scope}}
  \begin{scope}[shift={(-0.25,-.1)}]\tkzDrop\end{scope}
  \begin{scope}[shift={(0.25,-.1)}]\tkzDrop\end{scope}
}

% 公共别名: pst-labo 中 \tkzColumn 是用户可调用的公共宏 (小写)
% tikz-chemlab 内部用 \tkzColumn (驼峰), 这里注册公共别名保持兼容
\let\tkzColumn\tkzColumn

% réfrigérant (condenseur)
\newcommand{\tkzCondenser}{%
  \draw ([shift={(3,0)}]90:0.5) arc[start angle=90,end angle=270,radius=0.5];
  % pst-labo: \psarc(7,0){0.5}{270}{90} → 逆时针经过 0°/右侧
  \draw ([shift={(7,0)}]270:0.5) arc[start angle=270,end angle=450,radius=0.5];
  \fill[white] (0,-.05) rectangle (9,0.05);
  \begin{scope}[line width=0.4pt]
    \draw (-.075,.05) -- (9,0.05);
    \draw (0.02,-.05) -- (9,-.05);
  \end{scope}
  \draw (3,0.5) -- (3.5,0.5);
  \draw (3,-.5) -- (6.5,-.5);
  \draw (3.6,0.5) -- (7,0.5);
  \begin{scope}[line width=0.4pt]
    \draw (3.5,0.5) -- (3.5,1.5);
    \draw (3.6,0.5) -- (3.6,1.5);
    \draw (6.6,-.5) -- (6.6,-1.5);
    \draw (6.5,-.5) -- (6.5,-1.5);
  \end{scope}
  \draw (6.6,-.5) -- (7,-.5);
  \draw[->] (6.55,-2.5) -- (6.55,-1.5);
  \draw[->] (3.55,1.5)  -- (3.55,2.5);
}

\newcommand{\tkzThermometer}{%
  % pst-labo: \psline[linearc=0.05](-.05,.25)(-.05,2.75)(0.05,2.75)(0.05,0.25)
  % linearc=0.05 × unit, TikZ: rounded corners=0.05 (受 scale 缩放)
  \draw[rounded corners=0.025cm]
    (-.05,.25) -- (-.05,2.75) -- (0.05,2.75) -- (0.05,0.25) -- cycle;
  % pst-labo: \psframe[framearc=0.8,fillstyle=solid,fillcolor=red](-.15,-.25)(0.15,0.25)
  % framearc=0.8: 80% of half-min(0.3,0.5)=0.15 → radius=0.12
  \filldraw[fill=red,draw=black,rounded corners=0.06cm] (-.15,-.25) rectangle (0.15,0.25);
  % pst-labo: \psline[linecolor=red](0,2) — 红色竖线, 使用默认线宽
  % 在 pst-labo 渲染中较粗, 看起来像红色填充
  \draw[red,line width=0.8pt] (0,0) -- (0,2);
}

\newcommand{\tkzHeater}{%
  \draw (1.17,2.5) -- (0,2.5) -- (0,0) -- (5.5,0) -- (5.5,2.5) -- (4.33,2.5);
}

% 内部小量筒轮廓 - 路径版本 (供 clip 使用)
\newcommand{\tkzGraduatedCylinderPath}{%
  (-0.5,6) -- (-0.5,.5)
  -- ([shift={(-1,0.5)}]0:0.5) arc[start angle=0,end angle=-90,radius=0.5]
  -- (-1,0) -- (1,0)
  -- ([shift={(1,0.5)}]270:0.5) arc[start angle=270,end angle=180,radius=0.5]
  -- (0.5,0.5) -- (0.5,6) -- cycle
}

% 内部小量筒 (描边)
% pst-labo: \psset{linewidth=1.5\pslinewidth} — \pslinewidth 默认 0.8pt (ABSOLUTE)
% → 1.5*0.8pt = 1.2pt (绝对值, 不随 unit/scale 缩放)
% 注意: TikZ 的 [line width=N\pgflinewidth] 有循环引用 bug — \pgflinewidth 在
% 乘法前被更新, 导致 line width=3\pgflinewidth 实际给出 3.6pt 而非 1.2pt。
% 必须用显式 pt 值。
\newcommand{\tkzGraduatedCylinderSmall}{%
  \begin{scope}[line width=0.72pt,line cap=round,line join=round]
    \draw (-0.5,6) -- (-0.5,.5)
      -- ([shift={(-1,0.5)}]0:0.5) arc[start angle=0,end angle=-90,radius=0.5]
      -- (-1,0) -- (1,0)
      -- ([shift={(1,0.5)}]270:0.5) arc[start angle=270,end angle=180,radius=0.5]
      -- (0.5,0.5) -- (0.5,6);
  \end{scope}
}

% 带液体的量筒 (内部使用, 重命名以避免与主文件 \pstEproulette 冲突)
\newcommand{\tkzGraduatedCylinderWithLiquid}{%
  \begin{scope}
    \clip \tkzGraduatedCylinderPath;
    % pst-labo: \psframe*[linecolor=\psk@Distillation@DistillateColor]
    \fill[\tkzDistillateColor] (-1.5,-1) rectangle (1.5,1);
  \end{scope}
  \tkzGraduatedCylinderSmall
}

% (注 : 原 PSTricks 中 \pst@roundFlask 已合并到 \tkzRoundFlaskSimple, 此处不再单独定义,
%  避免与 tikz-chemlab.tex 中的用户命令 \tkzRoundFlask 冲突)

% 自定义密集斜线图案 (匹配 pst-labo hlines 的密度)
{\pgfpointorigin}{\pgfpoint{0.012cm}{0.012cm}}{\pgfpoint{0.012cm}{0.012cm}}
{\pgfsetlinewidth{0.0028cm}\pgfpathmoveto{\pgfpoint{0cm}{0cm}}\pgfpathlineto{\pgfpoint{0.012cm}{0.012cm}}\pgfusepath{stroke}}
{\pgfpointorigin}{\pgfpoint{0.012cm}{0.012cm}}{\pgfpoint{0.012cm}{0.012cm}}
{\pgfsetlinewidth{0.004cm}\pgfpathmoveto{\pgfpoint{0cm}{0cm}}\pgfpathlineto{\pgfpoint{0.012cm}{0.012cm}}\pgfusepath{stroke}}
% \pst@DistillationFractionnee 主图
\newcommand{\tkzDistillationFractionated}{%
  % --- roundFlask ---
  \def\flaskLocal{%
    % pst-labo: fillstyle=hlines,hatchwidth=0.007,hatchsep=0.03 → 水平线填充
    \begin{scope}[pattern={Lines[angle=45,distance={1pt/sqrt(2)}]},pattern color=gray]
      \fill (-.5,2.25) rectangle (-.25,3);
      \fill (0.5,2.25)  rectangle (0.25,3);
    \end{scope}
    % pst-labo 线宽: neck=0.053, arcs=0.03 (都是 ×unit 用户单位, 随 scale 缩放)
    % \chemlabUpdateLineWidths 已设置 \chemlabLineWidth = 0.053*scale
    \chemlabUpdateLineWidths
    \draw[line width=\chemlabLineWidth] (0.5,1.5) -- (0.5,3);
    \draw[line width=\chemlabLineWidth] (-.5,1.5) -- (-.5,3);
    % pst-labo: \psarc(0,0){1.581}{108.435}{71.565} 逆时针走长路 (经底部),
    %   左右侧弧可见, 底部被液体填充覆盖, 顶部无可见弧线。此处拆分为
    %   左右侧两段 (底部由下方轮廓弧绘制, 顶部不画)。
    % 右上侧弧: 71° → 0° (连接右颈底部到球体最右点)
    \draw[line width=\chemlabLineWidth] ([shift={(0,0)}]71.565:1.581) arc[start angle=71.565,end angle=0,radius=1.581];
    % 左上侧弧: 180° → 108° (连接球体最左点到左颈底部)
    \draw[line width=\chemlabLineWidth] ([shift={(0,0)}]180:1.581) arc[start angle=180,end angle=108.435,radius=1.581];
    % pst-labo: \psset{style=\psk@Distillation@MixtureStyle}\fill
    % 需要用 \shade (不是 \fill) 来产生渐变效果
    \shade[\tkzMixtureStyle] plot[smooth,tension=0.5] coordinates {%
      (-1.581,0) (-1.25,0.1) (-.75,-.1) (0,.1) (0.75,-.1) (1.25,.1) (1.581,0)}
      arc[start angle=0,end angle=-180,radius=1.581];
    % 底部轮廓弧: pst-labo \psarc[linewidth=0.03] → 0.03*unit, 用 \chemlabLineWidth 0.053
    % (pst-labo 中 0.03 vs 0.053 差异微小, 统一用 0.053*scale 避免额外变量)
    \draw[line width=\chemlabLineWidth] (0,0) ++(180:1.581) arc[start angle=180,end angle=360,radius=1.581];
  }%
  % --- 主图整体平移 ---
  \begin{scope}[shift={(0,-3)}]
    \begin{scope}[shift={(2.75,13.5)}]\tkzThermometer\end{scope}
    \begin{scope}[shift={(2.75,2.5)}]\flaskLocal\end{scope}
    \begin{scope}[shift={(2.75,4)}]\tkzColumn\end{scope}
    \begin{scope}[shift={(3,13.4)},rotate=-45]\tkzCondenser\end{scope}
    \begin{scope}[shift={(0,0)}]\tkzHeater\end{scope}
    \begin{scope}[shift={(9.4,0)}]\tkzGraduatedCylinderWithLiquid
      % 蒸馏装置量筒刻度 : 主刻度 1.5× 细刻度
      \pgfmathsetlengthmacro{\chemlab@distill@major}{1.5*\chemlabGraduationLineWidth}%
      \foreach \Level in {1,2,...,5} {\draw[line width=\chemlab@distill@major] (-.4,\Level) -- (-.1,\Level);}
      \foreach \Level in {0.2,0.4,...,4.8} {\draw[line width=\chemlabGraduationLineWidth] (-.4,\Level) -- (-.2,\Level);}
    \end{scope}
    \begin{scope}[shift={(9.37,7)}]\tkzDrop\end{scope}
    \begin{scope}[shift={(9.37,6)}]\tkzDrop\end{scope}
    \begin{scope}[shift={(9.37,4)}]\tkzDrop\end{scope}
    \begin{scope}[shift={(9.37,1.2)}]\tkzDrop\end{scope}
    % 加热盘 (搅拌器)
    \begin{scope}[shift={(0.75,.75)}]
      \draw (0,0) circle (.5);
      \foreach \i in {0,30,...,180} {\draw (\i:0.55) -- (\i:0.65);}
      \draw[->,line width=0.035cm] (0,0) -- (0,0.45);
    \end{scope}
  \end{scope}
}

% ============================================================
% 9. 本生灯 (Bec Bunsen) - 不带 grille
% ============================================================
\newcommand{\tkzBunsenBurner}{%
  \chemlabUpdateLineWidths\begin{scope}[line width=\chemlabLineWidth,line cap=round,line join=round]
    \draw (-1.25,0) rectangle (1.25,0.25);
    \draw (-.5,1.25) rectangle (0.5,2.25);
    \foreach \n in {-0.3,0.0,0.3} {\draw (\n,1.75) circle (0.1);}
    \draw (-.25,2.25) rectangle (0.25,4.25);
    \draw (0.25,1.25) -- (0.25,0.5) -- (1.25,0.25);
    \draw (-1.25,0.25) -- (-.25,0.5) -- (-0.25,0.75);
    \draw (-2.25,0.75) -- (-.25,0.75);
    \draw (-2.25,1) -- (-.25,1);
    \draw (-.25,0.875) ellipse (0.1 and 0.125);
    \fill[white] (-2.25,0.75) rectangle (-0.25,1);
    \draw (-2.25,0.75) -- (-0.25,0.75);
    \draw (-2.25,1) -- (-0.25,1) -- (-.25,1.25);
    \draw plot[smooth,tension=0.5] coordinates {(-0.25,0.5) (0,0.4) (0.25,0.5)};
    % flamme
    % pst-labo 源码: \rput(0,4.25){\psclip{\psbezier[...gradient...] (12点)}\pspolygon[...gradient...](-0.25,0)(0.25,0)(0,1)\endpsclip}
    % \psbezier 12个点 + BezierNArray padding (0,0) = 13点 = 4段贝塞尔
    % 经像素级验证, 正确路径从 (0,0) 出发, 经4段贝塞尔回到 (0,0), 无闭合直线:
    %   Seg 1: (0,0) → c1=P1=(-0.25,0), c2=P2=(-0.35,0.5) → P3=(-0.4,0.75) [左下→左中]
    %   Seg 2: P3=(-0.4,0.75) → c1=P4=(-0.35,1), c2=P5=(-0.25,1.5) → P6=(0.5,2) [左中→顶]
    %   Seg 3: P6=(0.5,2) → c1=P7=(0.25,1.5), c2=P8=(0.35,1) → P9=(0.4,0.75) [顶→右中]
    %   Seg 4: P9=(0.4,0.75) → c1=P10=(0.35,0.5), c2=P11=(0.25,0) → (0,0) [右中→底]
    % 渐变: gradbegin=PaleOrange (底), gradend=yellow (顶), gradmidpoint=0, gradangle=0 (垂直)
    \begin{scope}[shift={(0,4.25)}]
      % 外焰填充: PaleOrange (底) -> yellow (顶) 线性垂直渐变
      \shade[top color=yellow, bottom color=PaleOrange]
        (0,0) .. controls (-0.25,0) and (-0.35,0.5)
        .. (-0.4,0.75) .. controls (-0.35,1) and (-0.25,1.5)
        .. (0.5,2) .. controls (0.25,1.5) and (0.35,1)
        .. (0.4,0.75) .. controls (0.35,0.5) and (0.25,0)
        .. (0,0);
      % 内焰: 三角形 cyan(底)->white(顶) 渐变, 裁剪到外焰形状内
      \begin{scope}
        \clip
          (0,0) .. controls (-0.25,0) and (-0.35,0.5)
          .. (-0.4,0.75) .. controls (-0.35,1) and (-0.25,1.5)
          .. (0.5,2) .. controls (0.25,1.5) and (0.35,1)
          .. (0.4,0.75) .. controls (0.35,0.5) and (0.25,0)
          .. (0,0);
        \shade[top color=white, bottom color=cyan]
          (-0.25,0) -- (0.25,0) -- (0,1) -- cycle;
      \end{scope}
    \end{scope}
  \end{scope}
}

% 火焰 + 网格 (FlammeGrille 形状) - 路径版本 (供 clip 使用)
\newcommand{\tkzFlameGrillPath}{%
  ([shift={(0,0.75)}]270:0.75) arc[start angle=270,end angle=360,radius=0.75]
  -- ([shift={(1.5,0.75)}]180:0.75) arc[start angle=180,end angle=90,radius=0.75]
  -- (1.5,1.5) -- (-1.5,1.5)
  -- ([shift={(-1.5,0.75)}]90:0.75) arc[start angle=90,end angle=0,radius=0.75]
  -- ([shift={(0,0.75)}]180:0.75) arc[start angle=180,end angle=270,radius=0.75]
  -- cycle
}

\newcommand{\tkzFlameGrill}{%
  \fill[flameAndGrill] \tkzFlameGrillPath;
}

% ============================================================
% 10. 本生灯 - 带 grille
% ============================================================
\newcommand{\tkzBunsenBurnerGrill}{%
  \chemlabUpdateLineWidths\begin{scope}[line width=\chemlabLineWidth,line cap=round,line join=round]
    \draw (-1.25,0) rectangle (1.25,0.25);
    \draw (-.5,1.25) rectangle (0.5,2.25);
    \foreach \n in {-0.3,0.0,0.3} {\draw (\n,1.75) circle (0.1);}
    \draw (-.25,2.25) rectangle (0.25,4.25);
    \draw (0.25,1.25) -- (0.25,0.5) -- (1.25,0.25);
    \draw (-1.25,0.25) -- (-.25,0.5) -- (-0.25,0.75);
    \draw (-2.25,0.75) -- (-.25,0.75);
    \draw (-2.25,1) -- (-.25,1);
    \draw (-.25,0.875) ellipse (0.1 and 0.125);
    \fill[white] (-2.25,0.75) rectangle (-0.25,1);
    \draw (-2.25,0.75) -- (-0.25,0.75);
    \draw (-2.25,1) -- (-0.25,1) -- (-.25,1.25);
    \draw plot[smooth,tension=0.5] coordinates {(-0.25,0.5) (0,0.4) (0.25,0.5)};
    % flamme (clip 到 FlammeGrille)
    % 原 pst-labo 火焰分两层:
    %   外焰: \pst@FlammeGrille 整体填充 flameAndGrill (PaleOrange->yellow 渐变)
    %   内焰: 三角形 cyan->white 渐变 (在内焰区域)
    % 颜色和大小可通过 \chemlabset 调整
    \begin{scope}[shift={(0,4.25)},scale=\chemlabFlameScale]
      \begin{scope}
        \clip \tkzFlameGrillPath;
        % 外焰: PaleOrange (底) -> yellow (顶) 线性垂直渐变
        % 对应 pst-labo: fillstyle=gradient, gradbegin=PaleOrange, gradend=yellow
        % 渐变范围与火焰形状 (y=0..1.5) 一致, 避免底部颜色被稀释
        \shade[top color=\chemlabFlameOuterEndColor, bottom color=\chemlabFlameOuterColor]
          (-1.5,0) rectangle (1.5,1.5);
        % 内焰: cyan (底) -> white (顶) 线性垂直渐变 (三角形)
        \shade[top color=\chemlabFlameInnerEndColor, bottom color=\chemlabFlameInnerColor]
          (-0.25,0) -- (0.25,0) -- (0,1) -- cycle;
      \end{scope}
    \end{scope}
    \draw[dashed,line width=\chemlabEllipseLineWidth] (-2,5.75) -- (2,5.75);
  \end{scope}
}

% ============================================================
% 11. 烧瓶 (Ballon)
% ============================================================
% 路径版本 - 供 clip 使用 (闭合轮廓)
\newcommand{\tkzRoundFlaskSimplePath}{%
  ([shift={(-0.7,3.4)}]90:0.2) arc[start angle=90,end angle=0,radius=0.2]
  -- (-0.5,3.4) -- (-0.5,2.4)
  -- (-0.489,2.401)
  arc[start angle=113,end angle=427,radius=1.25]
  -- (0.5,2.4) -- (0.5,3.4)
  -- ([shift={(0.7,3.4)}]180:0.2) arc[start angle=180,end angle=90,radius=0.2]
  -- (0,3.5) -- cycle
}

\newcommand{\tkzRoundFlaskSimple}{%
  \chemlabUpdateLineWidths\begin{scope}[line width=\chemlabLineWidth,line cap=round,line join=round]
    % 左侧: 瓶口圆弧 + 颈部直线 + 球体大圆弧, 合并为单条路径以平滑过渡
    \draw ([shift={(-0.7,3.4)}]90:0.2) arc[start angle=90,end angle=0,radius=0.2]
          -- (-0.5,3.4) -- (-0.5,2.4)
          -- (-0.489,2.401) arc[start angle=113,end angle=427,radius=1.25];
    % 右侧: 球体终点(0.489,2.401) -> 颈部直线 -> 瓶口圆弧
    \draw (0.489,2.401) -- (0.5,2.4) -- (0.5,3.4)
          -- ([shift={(0.7,3.4)}]180:0.2) arc[start angle=180,end angle=90,radius=0.2];
    % 瓶口椭圆: 白色填充覆盖颈部上端圆头, 避免穿入椭圆
    % pst-labo: ballon椭圆用默认线宽 0.053 (与主轮廓一致, 不是 0.07)
    \filldraw[fill=white,draw=black,line width=\chemlabLineWidth] (0,3.6) ellipse (0.7 and 0.1);
  \end{scope}
}

\newcommand{\tkzRoundFlaskGrill}{%
  \chemlabUpdateLineWidths\begin{scope}[line width=\chemlabLineWidth,line cap=round,line join=round]
    % 左侧: 瓶口圆弧 + 颈部直线 + 球体大圆弧, 合并为单条路径
    \draw ([shift={(-0.7,2.15)}]90:0.2) arc[start angle=90,end angle=0,radius=0.2]
          -- (-0.5,2.15) -- (-0.5,1.4)
          -- (-0.489,1.401) arc[start angle=113,end angle=427,radius=1.25];
    % 右侧: 球体终点 -> 颈部直线 -> 瓶口圆弧
    \draw (0.489,1.401) -- (0.5,1.15) -- (0.5,3.4)
          -- ([shift={(0.7,2.15)}]180:0.2) arc[start angle=180,end angle=90,radius=0.2];
    % 瓶口椭圆: 白色填充覆盖颈部上端圆头, 避免穿入椭圆
    % pst-labo: ballon椭圆用默认线宽 0.053 (与主轮廓一致)
    \filldraw[fill=white,draw=black,line width=\chemlabLineWidth] (0,2.35) ellipse (0.7 and 0.1);
  \end{scope}
}

% ============================================================
% 12. 试管 (TubeEssais)
% ============================================================
\newcommand{\tkzTestTubeStopper}{%
  \chemlabUpdateLineWidths
  \begin{scope}[fill=LightGray,draw=black,line width=\chemlabLineWidth,line cap=round,line join=round]
    \filldraw (0,-0.3) ellipse (0.4 and 0.1);
    \filldraw (-0.4,-0.3) -- (-0.6,0.5) -- (0.6,0.5) -- (0.4,-0.3) -- cycle;
    \fill[draw=none] (0,-0.3) ellipse (0.4 and 0.1);
    \filldraw (0,0.5) ellipse (0.6 and 0.1);
  \end{scope}
  \draw[line width=\chemlabEllipseLineWidth] plot[smooth,tension=0.5]
        coordinates {(-0.5,0) (0,-0.07) (0.5,0)};
}

\newcommand{\tkzTestTubeBody}{%
  \chemlabUpdateLineWidths\begin{scope}[line width=\chemlabLineWidth,line cap=round,line join=round]
    % 合并左侧直线 + 底部圆弧 + 右侧直线为单条路径, 平滑转角
    \draw (-0.5,3.5) -- (-0.5,0.5)
          arc[start angle=180,end angle=360,radius=0.5]
          -- (0.5,3.5);
    % 试管口椭圆 (原 pst-labo 用 0.08cm, 比其他椭圆粗)
    \draw[line width=\chemlabTubeEllipseLineWidth] (0,3.5) ellipse (0.5 and 0.1);
  \end{scope}
}

\newcommand{\tkzTestTubeStraightTube}{%
  \draw[line width=0.5\pgflinewidth] (0.05,-1) -- (0.05,5);
  \draw[line width=0.5\pgflinewidth] (-0.05,-1) -- (-0.05,5);
  \tkzTestTubeStopper
  \draw[fill=white,rounded corners=0.5,line width=0.5\pgflinewidth]
       (-0.05,5) rectangle (0.05,0.5);
  \draw[white,line width=0.14cm] (-0.1,5) -- (0.1,5);
}

\newcommand{\tkzTestTubeElbowTube}{%
  \draw[line width=0.5\pgflinewidth] (0.05,-1) -- (0.05,0);
  \draw[line width=0.5\pgflinewidth] (-0.05,-1) -- (-0.05,0);
  \tkzTestTubeStopper
  % 弯管路径: pst-labo \pspolygon(...)(8,2)(8,1.9)...
  \filldraw[fill=white,rounded corners=0.05cm,line width=0.5\pgflinewidth]
       (-0.05,0.5) -- (-0.05,2) -- (8,2) -- (8,1.9)
       -- (0.05,1.9) -- (0.05,0.5) -- cycle;
  \draw[white,line width=0.14cm] (7.95,1.9) -- (8.05,2);
}

\newcommand{\tkzTestTubeUTube}{%
  \draw[line width=0.5\pgflinewidth] (0.05,-1) -- (0.05,0);
  \draw[line width=0.5\pgflinewidth] (-0.05,-1) -- (-0.05,0);
  \tkzTestTubeStopper
  \filldraw[fill=white,rounded corners=0.05cm,line width=0.5\pgflinewidth]
       (-0.05,0.5) -- (-0.05,2) -- (4.5,2) -- (4.5,-3)
       -- (4.4,-3) -- (4.4,1.9) -- (0.05,1.9) -- (0.05,0.5) -- cycle;
  \draw[white,line width=0.14cm] (4.3,-3) -- (4.6,-3);
}

\newcommand{\tkzTestTubeUTubeLong}{%
  \draw[line width=0.5\pgflinewidth] (0.05,-1) -- (0.05,0);
  \draw[line width=0.5\pgflinewidth] (-0.05,-1) -- (-0.05,0);
  \tkzTestTubeStopper
  \filldraw[fill=white,rounded corners=0.05cm,line width=0.5\pgflinewidth]
       (-0.05,0.5) -- (-0.05,2) -- (4.5,2) -- (4.5,-9)
       -- (4.4,-9) -- (4.4,1.9) -- (0.05,1.9) -- (0.05,0.5) -- cycle;
  \draw[white,line width=0.14cm] (4.3,-9) -- (4.6,-9);
}

\newcommand{\tkzCurvedTube}{%
  \draw[line width=0.5\pgflinewidth] (0.05,-1) -- (0.05,0);
  \draw[line width=0.5\pgflinewidth] (-0.05,-1) -- (-0.05,0);
  \tkzTestTubeStopper
  \begin{scope}[line width=0.4pt]
    \draw (-0.05,0.5) -- (-0.05,1);
    \draw ([shift={(0.5,1)}]180:0.55) arc[start angle=180,end angle=90,radius=0.55];
    \draw (0.5,1.55) -- (5.5,1.55);
    \draw ([shift={(5.5,1)}]90:0.55) arc[start angle=90,end angle=0,radius=0.55];
    \draw (6.05,1) -- (6.05,-8.8);
    \draw ([shift={(6.5,-8.8)}]180:0.45) arc[start angle=180,end angle=360,radius=0.45];
    \draw (6.95,-8.8) -- (6.95,-8.3);
    \draw (7.05,-8.3) -- (7.05,-8.8);
    \draw ([shift={(6.5,-8.8)}]0:0.55) arc[start angle=0,end angle=-180,radius=0.55];
    \draw (5.95,-8.8) -- (5.95,1);
    \draw ([shift={(5.5,1)}]0:0.45) arc[start angle=0,end angle=90,radius=0.45];
    \draw (5.5,1.45) -- (0.5,1.45);
    \draw ([shift={(0.5,1)}]90:0.45) arc[start angle=90,end angle=180,radius=0.45];
    \draw (0.05,1) -- (0.05,0.5);
    \draw (-0.05,0.5) -- (0.05,0.5);
  \end{scope}
  \fill[white] (-0.025,0.52) rectangle (0.020,0.6);
}

\newcommand{\tkzCurvedTubeShort}{%
  \draw[line width=0.5\pgflinewidth] (0.05,-1) -- (0.05,0);
  \draw[line width=0.5\pgflinewidth] (-0.05,-1) -- (-0.05,0);
  \tkzTestTubeStopper
  \draw (-0.05,0.5) -- (-0.05,1);
  \begin{scope}[line width=0.4pt]
    \draw ([shift={(0.5,1)}]180:0.55) arc[start angle=180,end angle=90,radius=0.55];
    \draw (0.5,1.55) -- (3.5,1.55);
    \draw ([shift={(3.5,1)}]90:0.55) arc[start angle=90,end angle=0,radius=0.55];
    \draw (4.05,1) -- (4.05,-2.8);
    \draw ([shift={(4.5,-2.8)}]180:0.45) arc[start angle=180,end angle=360,radius=0.45];
    \draw (4.95,-2.8) -- (4.95,-2.3);
    \draw (5.05,-2.3) -- (5.05,-2.8);
    \draw ([shift={(4.5,-2.8)}]0:0.55) arc[start angle=0,end angle=-180,radius=0.55];
    \draw (3.95,-2.8) -- (3.95,1);
    \draw ([shift={(3.5,1)}]0:0.45) arc[start angle=0,end angle=90,radius=0.45];
    \draw (3.5,1.45) -- (0.5,1.45);
    \draw ([shift={(0.5,1)}]90:0.45) arc[start angle=90,end angle=180,radius=0.45];
    \draw (0.05,1) -- (0.05,0.5);
    \draw (-0.05,0.5) -- (0.05,0.5);
  \end{scope}
  \fill[white] (-0.025,0.52) rectangle (0.020,0.6);
}

\newcommand{\tkzCurvedTubeShortDoubleTube}{%
  \begin{scope}[line width=0.4pt]
    \draw (-0.05,0.5) -- (-0.05,1);
    \draw ([shift={(0.5,1)}]180:0.55) arc[start angle=180,end angle=90,radius=0.55];
    \draw (0.5,1.55) -- (3.5,1.55);
    \draw ([shift={(3.5,1)}]90:0.55) arc[start angle=90,end angle=0,radius=0.55];
    \draw (4.05,1) -- (4.05,-2.8);
    \draw ([shift={(4.5,-2.8)}]180:0.45) arc[start angle=180,end angle=360,radius=0.45];
    \draw (4.95,-2.8) -- (4.95,-2.3);
    \draw (5.05,-2.3) -- (5.05,-2.8);
    \draw ([shift={(4.5,-2.8)}]0:0.55) arc[start angle=0,end angle=-180,radius=0.55];
    \draw (3.95,-2.8) -- (3.95,1);
    \draw ([shift={(3.5,1)}]0:0.45) arc[start angle=0,end angle=90,radius=0.45];
    \draw (3.5,1.45) -- (0.5,1.45);
    \draw ([shift={(0.5,1)}]90:0.45) arc[start angle=90,end angle=180,radius=0.45];
    \draw (0.05,1) -- (0.05,0.5);
    \draw (-0.05,0.5) -- (0.05,0.5);
  \end{scope}
  \fill[white] (-0.025,0.52) rectangle (0.018,0.8);
}

\newcommand{\tkzTestTubeDoubleTube}{%
  \begin{scope}[line width=0.4pt]
    \draw (-0.3,-3) -- (-0.3,0);
    \draw (-.2,-3) -- (-0.2,0);
    \draw (0.25,-1) -- (0.25,0);
    \draw (.15,-1) -- (.15,0);
  \end{scope}
  \tkzTestTubeStopper
  \begin{scope}[shift={(0.2,0)}]\tkzCurvedTubeShortDoubleTube\end{scope}
  \begin{scope}[shift={(0,0.5)}]
    \begin{scope}[line width=0.4pt]
      \filldraw[fill=white,rounded corners=0.05cm]
        (-.3,0) -- (-0.3,2) -- (-0.5,2.2) -- (-0.5,3) -- (0,3)
                -- (0,2.2) -- (-0.2,2) -- (-0.2,0) -- cycle;
      \draw (0.1,1.55) -- (-.5,1.55) -- (-0.5,1.7) -- (-.55,1.7)
            -- (-.55,1.3) -- (-0.5,1.3) -- (-0.5,1.55) -- (-.5,1.45) -- (.1,1.45);
      \filldraw[fill=white,rounded corners=0.2] (-0.4,1.3) rectangle (-0.1,1.7);
    \end{scope}
    \fill[white] (-0.25,3) ellipse (0.3 and 0.1);
  \end{scope}
}

% 试管夹 (Pince) - 几何与原 pst-labo pst-laboObj.tex line 483-508 一致
% 根据原 pst-labo 实际渲染 (5.3x 截图) 校准:
%   - 上边缘凹弧较浅 (原 \pscurve 3点几乎共线, 浅弧)
%   - 顶部小凸起为圆弧形 (\psarcn 圆弧)
%   - 所有转角为圆角 (原 \pscustom 默认 linejoin=0 即 round)
\newcommand{\tkzTestTubeClamp}{%
  % 第1块: 上部夹爪
  % 上边缘(1a, 夹合面): 下凹弧贴合试管顶部弧面
  %   试管半径 0.5, 在 y=3.2 处试管顶部弧面下凹约 0.1, 用贝塞尔曲线模拟
  %   从 (0.5,0) 经 (0,-0.1) 到 (-0.5,0), 控制点 y=-0.13 产生平滑下凹
  \filldraw[fill=PaleOrange,draw=black,line cap=round,line join=round]
    (0.5,0) .. controls (0.25,-0.13) and (-0.25,-0.13) .. (-0.5,0)
    -- (-0.8,0)
    -- ([shift={(-0.5,0)}]180:0.3) arc[start angle=180,end angle=270,radius=0.3]
    -- (-.5,-.3) -- (2.5,-.3) -- (2.4,-.2) -- (1,0) -- (0.5,0) -- cycle;
  % 第2块: 下部夹爪
  % 注: 原 \psarcn(-.5,0){0.3}{-90}{180} 是顺时针从 -90°(=270°) 到 180°, 跨度 90° (四分之一圆弧)
  % TikZ arc 当 end<start 时顺时针, 需用 end=-180 让其走短弧 (避免四分之三圆弧)
  \filldraw[fill=PaleOrange,draw=black,line cap=round,line join=round]
    (-0.8,0) -- (-0.8,-0.3)
    -- ([shift={(-0.5,-0.3)}]180:0.3) arc[start angle=180,end angle=270,radius=0.3]
    -- (-0.5,-0.6) -- (2.5,-.6) -- (2.5,-.3) -- (-.5,-.3)
    -- ([shift={(-.5,0)}]-90:0.3) arc[start angle=-90,end angle=-180,radius=0.3] -- cycle;
  % 第3块: 顶部小凸起 - 原 \psarcn 顺时针 180°->90°, 圆弧形凸起
  % 去掉 (-0.8,0.05)->(-0.75,0.05) 的多余水平线段, 直接从 (-0.5,0.05) 到圆弧起点
  % 使凸起呈纯圆弧形 (与原 pst-labo 渲染一致)
  \filldraw[fill=PaleOrange,draw=black,line cap=round,line join=round]
    (-.5,0.05) -- ([shift={(-0.5,0.05)}]180:0.25)
    arc[start angle=180,end angle=90,radius=0.25]
    -- (-.5,0.3) -- (-0.5,0.05) -- cycle;
  % 第4块: 夹臂上部
  \filldraw[fill=PaleOrange,draw=black,line cap=round,line join=round]
    (0.5,0.3) -- (4.9,0.3) -- (5,0.2) -- (1,0) -- (0.5,0) -- (0.5,0.3) -- cycle;
  % 第5块: 夹臂下部 (Brown 棕色)
  \filldraw[fill=Brown,draw=black,line cap=round,line join=round]
    (0.5,0) -- (5,0.2) -- (5,0) -- (2,-0.2) -- (1,0) -- cycle;
  % 螺丝椭圆
  \draw[line width=1.5\pgflinewidth,line cap=round] (1,0) ellipse (0.3 and 0.1);
  % 螺丝下方的杆
  \draw[line width=1.5\pgflinewidth,line cap=round,line join=round]
        (0.7,0) -- (0.4,-.3) -- (0.4,-.6);
}

% ============================================================
% 13. 容量瓶 (Fiole jaugée)
% ============================================================
\newcommand{\tkzVolumetricFlask}{%
  \chemlabUpdateLineWidths
  \draw[rounded corners=0.1cm,line width=\chemlabLineWidth,line cap=round,line join=round]
       (-0.21,3.5) -- (-0.1,3.5) -- (-0.1,1.5)
       -- (-0.75,0) -- (0.75,0) -- (0.1,1.5) -- (0.1,3.5) -- (0.21,3.5);
  \draw[line width=\chemlabGraduationLineWidth,line cap=round] (-0.1,2.5) -- (0.1,2.5);
}

% ============================================================
% 14. 结晶皿 (Cristallisoir)
% ============================================================
\newcommand{\Cristallisoir}{%
  \chemlabUpdateLineWidths
  \begin{scope}[line width=\chemlabLineWidth,line cap=round,line join=round]
  \draw (-1.9,2) -- (-2,2);
  \draw ([shift={(-2,1.95)}]90:0.05) arc[start angle=90,end angle=270,radius=0.05];
  \draw (-2,1.9) -- (-2,1);
  \draw ([shift={(-1,1)}]180:1) arc[start angle=180,end angle=270,radius=1];
  \draw (-1,0) -- (1,0);
  \draw ([shift={(1,1)}]270:1) arc[start angle=270,end angle=360,radius=1];
  \draw (2,1) -- (2,1.9);
  \draw ([shift={(2,1.95)}]270:0.05) arc[start angle=270,end angle=90,radius=0.05];
  \draw (2,2) -- (1.9,2);
  \draw (1.9,2) -- (1.9,1);
  \draw ([shift={(1,1)}]0:0.9) arc[start angle=0,end angle=-90,radius=0.9];
  \draw (1,0.1) -- (-1,0.1);
  \draw ([shift={(-1,1)}]270:0.9) arc[start angle=270,end angle=180,radius=0.9];
  \draw (-1.9,1) -- (-1.9,2);
  \end{scope}
}

\newcommand{\InterieurCristallisoir}{%
  (1.9,2) -- (1.9,1)
  -- ([shift={(1,1)}]0:0.9) arc[start angle=0,end angle=-90,radius=0.9]
  -- (1,0.1) -- (-1,0.1)
  -- ([shift={(-1,1)}]270:0.9) arc[start angle=270,end angle=180,radius=0.9]
  -- (-1.9,1) -- (-1.9,2) -- cycle
}

\newcommand{\tkzCrystallizingDish}{%
  \begin{scope}
    \clip \InterieurCristallisoir;
    \fill[LightBlue1] (-3,-1) rectangle (3,1.5);
  \end{scope}
  \Cristallisoir
}

% ============================================================
% 15. 倒置试管 (TubeRenverse)
% ============================================================
\newcommand{\tkzInvertedTube}{%
  \chemlabUpdateLineWidths
  \begin{scope}[line width=\chemlabEllipseLineWidth,line cap=round,line join=round]
    \fill[LightBlue1] (-0.5,0) rectangle (0.5,5.5);
    % 合并左侧 + 顶部圆弧 + 右侧为单条路径
    % pst-labo: \psarcn(0,6){0.5}{180}{0} = 从180°顺时针到0°, 经过顶部90°
    % TikZ: arc[start=180, end=0] = 从180°逆时针(减小)到0°, 经过顶部90°
    \draw (-0.5,0) -- (-0.5,6)
          arc[start angle=180,end angle=0,radius=0.5]
          -- (0.5,0);
  \end{scope}
}

% ============================================================
% 16. 试管架 (ReposeTube)
% ============================================================
\newcommand{\tkzTestTubeRack}{%
  \begin{scope}[line width=0.1cm,draw=lightgray]
    \draw ([shift={(0,-0.25)}]100:1) arc[start angle=100,end angle=140,radius=1];
    \draw ([shift={(0,-.25)}]150:1) arc[start angle=150,end angle=160,radius=1];
    \draw ([shift={(0,-0.25)}]20:1) arc[start angle=20,end angle=80,radius=1];
  \end{scope}
}

% ============================================================
% 17. 烧杯边缘 + 刻度 (RebordBecher)
% ============================================================
\newcommand{\tkzBeakerRim}{%
  \chemlabUpdateLineWidths
  \draw[line width=\chemlabLineWidth] (0,2.55) ellipse (1.1 and .1);
  % 主刻度 : 粗细为细刻度线的 1.5 倍 (原 pst-labo: 0.03cm = 1.5 × 0.02cm)
  % 用 \pgfmathsetmacro 预算尺寸, 避免 \dimexpr 在 pgfkeys 中的兼容问题
  \pgfmathsetlengthmacro{\chemlab@majorGrad}{1.5*\chemlabGraduationLineWidth}%
  \foreach \nDiv in {0.5,1.0,...,2.0} {%
    \draw[line width=\chemlab@majorGrad] (-.75,\nDiv) -- (-.4,\nDiv);}
  % 细分刻度
  \foreach \nSubDiv in {0.5,0.6,...,2.0} {%
    \draw[line width=\chemlabGraduationLineWidth] (-0.75,\nSubDiv) -- (-.5,\nSubDiv);}
}

% ============================================================
% 18. 烧杯 (Becher)
% ============================================================
% 注 : 原 PSTricks 使用 \psline[linearc=0.5](-1,2.5)(-1,0)(1,0)(1,2.5)
% TikZ 的 rounded corners 与 PSTricks linearc 语义略有不同,
% 因此改用显式圆弧精确模拟 PSTricks linearc=0.5 的行为 :
% 左下角圆弧 : 圆心 (-0.5,0.5), 半径 0.5, 从 180° 到 270°
% 右下角圆弧 : 圆心 (0.5,0.5),  半径 0.5, 从 270° 到 360°
\newcommand{\tkzBeakerBody}{%
  \chemlabUpdateLineWidths
  \draw[line width=\chemlabLineWidth,line cap=round,line join=round]
       (-1,2.5) -- (-1,0.5)
       -- ([shift={(-0.5,0.5)}]180:0.5) arc[start angle=180,end angle=270,radius=0.5]
       -- (0.5,0)
       -- ([shift={(0.5,0.5)}]270:0.5) arc[start angle=270,end angle=360,radius=0.5]
       -- (1,2.5);
}

% ============================================================
% 19. 锥形瓶 (Erlenmeyer)
% ============================================================
% 注 : 原 PSTricks 使用 \psline[linearc=0.3](-0.5,3.5)(-0.5,2.5)(-1.5,0)(1.5,0)(0.5,2.5)(0.5,3.5)
% TikZ 的 rounded corners 与 PSTricks linearc 语义不同,
% 因此改用显式圆弧精确模拟 linearc=0.3 (由 scripts/erlen_arcs5.py 计算) :
%   肩部 (凹转角, 内侧角 21.8°) : 半径 0.3 太大无法相切 (切点距 1.558 > 颈部线段长 1.0), 保留尖角
%   左下 (凸转角, 内侧角 68.2°) : center=(-1.0569,0.3), r=0.3, 158.20° -> 270° (ccw)
%   右下 (凸转角, 内侧角 68.2°) : center=(1.0569,0.3),  r=0.3, -90° -> 21.80° (ccw)
\newcommand{\tkzErlenmeyerBody}{%
  \chemlabUpdateLineWidths
  \draw[line width=\chemlabLineWidth,line cap=round,line join=round]
       (-0.5,3.5) -- (-0.5,2.5)
       -- (-1.3354,0.4114)
       -- ([shift={(-1.0569,0.3)}]158.20:0.3) arc[start angle=158.20,end angle=270,radius=0.3]
       -- (1.0569,0)
       -- ([shift={(1.0569,0.3)}]-90:0.3) arc[start angle=-90,end angle=21.80,radius=0.3]
       -- (0.5,2.5) -- (0.5,3.5);
}

% ============================================================
% 20. 试剂瓶 (Flacon)
% ============================================================
\newcommand{\tkzReagentBottleBody}{%
  \chemlabUpdateLineWidths\begin{scope}[line width=\chemlabLineWidth,line cap=round,line join=round]
    % 左侧瓶口上角小弧 (psarcn 90->0, 顺时针)
    \draw ([shift={(-0.6,3.4)}]90:0.1) arc[start angle=90,end angle=0,radius=0.1];
    % 左侧颈部直线
    \draw (-0.5,3.4) -- (-0.5,3.1);
    % 左侧瓶口下角凹弧 (psarcn 0->-90, 顺时针; TikZ end<start 自动顺时针)
    \draw ([shift={(-0.6,3.1)}]0:0.1) arc[start angle=0,end angle=-90,radius=0.1];
    % 左侧肩部大凸弧 (psarc 90->180, 逆时针)
    \draw ([shift={(-0.6,2.6)}]90:0.4) arc[start angle=90,end angle=180,radius=0.4];
    % 左侧瓶身直线
    \draw (-1,2.6) -- (-1,0.2);
    % 左侧底部圆角 (psarc 180->270, 逆时针)
    \draw ([shift={(-.8,0.2)}]180:0.2) arc[start angle=180,end angle=270,radius=0.2];
    % 底部
    \draw (-.8,0) -- (.8,0);
    % 右侧底部圆角 (psarc -90->0, 逆时针)
    \draw ([shift={(.8,0.2)}]-90:0.2) arc[start angle=-90,end angle=0,radius=0.2];
    % 右侧瓶身直线
    \draw (1,0.2) -- (1,2.6);
    % 右侧肩部大凸弧 (psarc 0->90, 逆时针)
    \draw ([shift={(0.6,2.6)}]0:0.4) arc[start angle=0,end angle=90,radius=0.4];
    % 右侧瓶口下角凹弧 (psarcn -90->180, 顺时针短弧; TikZ 需 end=-180 避免 270° 长弧)
    \draw ([shift={(0.6,3.1)}]-90:0.1) arc[start angle=-90,end angle=-180,radius=0.1];
    % 右侧颈部直线
    \draw (0.5,3.1) -- (0.5,3.4);
    % 右侧瓶口上角小弧 (psarcn 180->90, 顺时针)
    \draw ([shift={(0.6,3.4)}]180:0.1) arc[start angle=180,end angle=90,radius=0.1];
  \end{scope}
}

% ============================================================
% 21. 玻璃杯 (Verre) - F. Bréal
% ============================================================
\newcommand{\tkzWatchGlassBody}{%
  \chemlabUpdateLineWidths\begin{scope}[line width=\chemlabLineWidth,line cap=round,line join=round]
    % 合并左斜线 + 底部圆弧 + 右斜线为单条路径, 平滑转角
    \draw (-1,3.5) -- (-0.5,1)
          arc[start angle=180,end angle=360,radius=0.5]
          -- (1,3.5);
  \end{scope}
}

\newcommand{\tkzWatchGlassBase}{%
  \chemlabUpdateLineWidths\begin{scope}[line width=\chemlabLineWidth,rounded corners=0.03cm,line cap=round,line join=round]
    \draw (-0.5,1) -- (-0.34,0.2) -- (-1.2,0) -- (1.2,0) -- (0.34,0.2) -- (0.5,1);
  \end{scope}
}

% ============================================================
% 22. 漏斗 (Entonnoir)
% ============================================================
\newcommand{\tkzFunnelOutline}{%
  \begin{scope}[line cap=round,line join=round]
    \draw (-0.07,0) -- (-0.07,1) -- (-1,2.5);
    \draw (1,2.5) -- (0.07,1) -- (0.07,0.1);
  \end{scope}
}

\newcommand{\tkzFunnelBody}{%
  \begin{scope}
    \clip (-0.07,0) -- (-0.07,1) -- (-1,2.5)
          -- (1,2.5) -- (0.07,1) -- (0.07,0.1) -- cycle;
    \fill[white] (-2,-1) rectangle (2,3);
  \end{scope}
  \begin{scope}[rounded corners=0.25cm]
    \draw (-0.07,0) -- (-0.07,1) -- (-1,2.5);
    \draw (1,2.5) -- (0.07,1) -- (0.07,0.1);
  \end{scope}
  \draw (0,2.5) ellipse (1 and 0.1);
}

% ============================================================
% 23. 球形冷凝器 (Réfrigérant à boulles)
% ============================================================
\newcommand{\tkzCondenserStem}{%
  % 整个 TigeBoulles 作为一个连续路径 (匹配 pst-labo \pscustom)
  % pst-labo: \pscustom[linestyle=none]{...\fill} 填充 + 单独描边 (无闭合线)
  % TikZ: 分开 \fill (带 cycle 闭合填充) 和 \draw (无 cycle, 顶部敞口)
  % 顶部是敞口的 — 两端弧之间没有连线 (匹配 pst-labo)
  %
  % 填充: 带闭合线 (fill 需要 cycle)
  \fill[white]
    % 顶部圆弧与颈部 - 左侧
    ({-0.35+0.1*cos(90)},{7.5+0.1*sin(90)}) arc[start angle=90,end angle=0,radius=0.1]
    -- (-.25,7.5) -- (-.25,6.6)
    -- ({-.15+0.1*cos(180)},{6.6+0.1*sin(180)}) arc[start angle=180,end angle=270,radius=0.1]
    -- ({-.15+0.1*cos(90)},{6.4+0.1*sin(90)}) arc[start angle=90,end angle=0,radius=0.1]
    -- (-.05,6.4) -- (-0.05,6)
    % 左侧 7 个球状突起 (中心 x=-0.05, 半径 0.25, arc 90→270 向左凸)
    (-0.05,6.0) arc[start angle=90,end angle=270,radius=0.25] -- (-0.05,5.50) -- (-0.05,5.25)
    arc[start angle=90,end angle=270,radius=0.25] -- (-0.05,4.75) -- (-0.05,4.50)
    arc[start angle=90,end angle=270,radius=0.25] -- (-0.05,4.00) -- (-0.05,3.75)
    arc[start angle=90,end angle=270,radius=0.25] -- (-0.05,3.25) -- (-0.05,3.00)
    arc[start angle=90,end angle=270,radius=0.25] -- (-0.05,2.50) -- (-0.05,2.25)
    arc[start angle=90,end angle=270,radius=0.25] -- (-0.05,1.75) -- (-0.05,1.50)
    arc[start angle=90,end angle=270,radius=0.25] -- (-0.05,1.00) -- (-0.05,0.75)
    % 左侧底部到下管口
    -- (-0.05,-1.5) -- (0.05,-1.5) -- (0.05,1)
    % 右侧 7 个球状突起 (中心 x=0.05, 半径 0.25, arc 270→450 向右凸)
    arc[start angle=270,end angle=450,radius=0.25] -- (0.05,1.50) -- (0.05,1.75)
    arc[start angle=270,end angle=450,radius=0.25] -- (0.05,2.25) -- (0.05,2.50)
    arc[start angle=270,end angle=450,radius=0.25] -- (0.05,3.00) -- (0.05,3.25)
    arc[start angle=270,end angle=450,radius=0.25] -- (0.05,3.75) -- (0.05,4.00)
    arc[start angle=270,end angle=450,radius=0.25] -- (0.05,4.50) -- (0.05,4.75)
    arc[start angle=270,end angle=450,radius=0.25] -- (0.05,5.25) -- (0.05,5.50)
    arc[start angle=270,end angle=450,radius=0.25] -- (0.05,6.00) -- (0.05,6.40)
    % 顶部圆弧与颈部 - 右侧
    -- ({0.15+0.1*cos(180)},{6.4+0.1*sin(180)}) arc[start angle=180,end angle=90,radius=0.1]
    -- ({0.15+0.1*cos(-90)},{6.6+0.1*sin(-90)}) arc[start angle=-90,end angle=0,radius=0.1]
    -- (0.25,6.6) -- (0.25,7.5)
    -- ({0.35+0.1*cos(180)},{7.5+0.1*sin(180)}) arc[start angle=180,end angle=90,radius=0.1]
    -- cycle;
  %
  % 描边: 底部敞口 (匹配 pst-labo: 左右侧分别 \psline, 无底部连线)
  \draw
    % 顶部圆弧与颈部 - 左侧
    ({-0.35+0.1*cos(90)},{7.5+0.1*sin(90)}) arc[start angle=90,end angle=0,radius=0.1]
    -- (-.25,7.5) -- (-.25,6.6)
    -- ({-.15+0.1*cos(180)},{6.6+0.1*sin(180)}) arc[start angle=180,end angle=270,radius=0.1]
    -- ({-.15+0.1*cos(90)},{6.4+0.1*sin(90)}) arc[start angle=90,end angle=0,radius=0.1]
    -- (-.05,6.4) -- (-0.05,6)
    % 左侧 7 个球状突起
    (-0.05,6.0) arc[start angle=90,end angle=270,radius=0.25] -- (-0.05,5.50) -- (-0.05,5.25)
    arc[start angle=90,end angle=270,radius=0.25] -- (-0.05,4.75) -- (-0.05,4.50)
    arc[start angle=90,end angle=270,radius=0.25] -- (-0.05,4.00) -- (-0.05,3.75)
    arc[start angle=90,end angle=270,radius=0.25] -- (-0.05,3.25) -- (-0.05,3.00)
    arc[start angle=90,end angle=270,radius=0.25] -- (-0.05,2.50) -- (-0.05,2.25)
    arc[start angle=90,end angle=270,radius=0.25] -- (-0.05,1.75) -- (-0.05,1.50)
    arc[start angle=90,end angle=270,radius=0.25] -- (-0.05,1.00) -- (-0.05,0.75)
    -- (-0.05,-1.5);
  % 右侧 (单独 \draw, 底部不连接左侧)
  \draw
    (0.05,-1.5) -- (0.05,1)
    % 右侧 7 个球状突起
    arc[start angle=270,end angle=450,radius=0.25] -- (0.05,1.50) -- (0.05,1.75)
    arc[start angle=270,end angle=450,radius=0.25] -- (0.05,2.25) -- (0.05,2.50)
    arc[start angle=270,end angle=450,radius=0.25] -- (0.05,3.00) -- (0.05,3.25)
    arc[start angle=270,end angle=450,radius=0.25] -- (0.05,3.75) -- (0.05,4.00)
    arc[start angle=270,end angle=450,radius=0.25] -- (0.05,4.50) -- (0.05,4.75)
    arc[start angle=270,end angle=450,radius=0.25] -- (0.05,5.25) -- (0.05,5.50)
    arc[start angle=270,end angle=450,radius=0.25] -- (0.05,6.00) -- (0.05,6.40)
    -- ({0.15+0.1*cos(180)},{6.4+0.1*sin(180)}) arc[start angle=180,end angle=90,radius=0.1]
    -- ({0.15+0.1*cos(-90)},{6.6+0.1*sin(-90)}) arc[start angle=-90,end angle=0,radius=0.1]
    -- (0.25,6.6) -- (0.25,7.5)
    -- ({0.35+0.1*cos(180)},{7.5+0.1*sin(180)}) arc[start angle=180,end angle=90,radius=0.1];
  % 底部液滴
  \begin{scope}[shift={(0,-1.6)}]\tkzDrop\end{scope}
}

\newcommand{\tkzCondenserBody}{%
  % 渐变填充: pst-labo gradbegin=cyan (底), gradend=PaleOrange (顶), gradmidpoint=0
  % TikZ: bottom color=cyan (底), top color=PaleOrange (顶)
  \shade[top color=PaleOrange,bottom color=cyan]
    ([shift={(0,5.75)}]0:0.5) arc[start angle=0,end angle=180,radius=0.5]
    -- (-0.5,5.75) -- (-0.5,5.5) -- (-1.5,5.5)
    -- (-1.5,5.3) -- (-.5,5.3) -- (-.5,1.7) -- (-1.5,1.7)
    -- (-1.5,1.5) -- (-.5,1.5) -- (-0.5,1)
    -- ([shift={(0,1)}]180:0.5) arc[start angle=180,end angle=360,radius=0.5]
    -- (0.5,1) -- (0.5,5.75) -- cycle;
  % 描边
  \begin{scope}[line width=0.88pt]
    \draw ([shift={(0,5.75)}]0:0.5) arc[start angle=0,end angle=180,radius=0.5];
    \draw (-0.5,5.75) -- (-0.5,5.5) -- (-1.5,5.5);
    \draw (-0.5,5.75) -- (-0.5,5.5) -- (-1.5,5.5);
    \draw (-1.5,5.3) -- (-.5,5.3) -- (-.5,1.7) -- (-1.5,1.7);
    \draw (-1.5,1.5) -- (-.5,1.5) -- (-0.5,1);
    \draw ([shift={(0,1)}]180:0.5) arc[start angle=180,end angle=360,radius=0.5];
    \draw (0.5,1) -- (0.5,5.75);
  \end{scope}
  \draw[->,cyan] (-2.5,1.6) -- (-1.6,1.6);
  \draw[<-,PaleOrange] (-2.5,5.4) -- (-1.6,5.4);
}

\newcommand{\tkzCondenserBulbs}{%
  \begin{scope}[shift={(0,0.2)}]
    \tkzCondenserBody
    \tkzCondenserStem
  \end{scope}
  \tkzTestTubeStopper
  \fill[white,rounded corners=0.025cm]
       (-0.05,0.6) -- (-0.05,0.5) -- (0.05,0.5) -- (0.05,0.6) -- cycle;
}

% ============================================================
% 24. 加热球 (Chauffe-Ballon)
% ============================================================
\newcommand{\tkzHeaterRoundFlask}{%
  \filldraw[fill=white,draw=black] (-2,-1) rectangle (2,1);
  \draw (-1.25,0) circle (0.4);
  \begin{scope}[shift={(-1.25,0)}]
    \foreach \iA in {0,30,...,180} {\draw (\iA:0.45) -- (\iA:0.6);}
    \draw[->] (0,0) -- (60:0.3);
  \end{scope}
  % 按钮区域: pst-labo \psframe[doubleline=true] 产生双线边框
  % PSTricks doubleline: 外框在指定坐标, 内框向内偏移
  % 实测: 外框到内框偏移约 0.13 units (含线宽+间隔, at scale=0.5)
  % 外框 (在指定坐标处)
  \draw[black,line width=0.8pt] (1,-.5) rectangle (1.5,0.5);
  % 内框 (向内偏移, 填充灰色)
  \filldraw[fill=gray,draw=black,line width=0.8pt] ({1+0.14},{-.5+0.14}) rectangle ({1.5-0.14},{0.5-0.14});
  % 红色指示灯
  \fill[red] (1.25,0.7) circle (0.1);
}

% ============================================================
% 25. 回流烧瓶 (Ballon à reflux)
% ============================================================
\newcommand{\tkzRoundFlaskReflux}{%
  \begin{scope}
    \clip \tkzRoundFlaskSimplePath;   % 使用路径版本形成闭合 clip 区域
    \shade[top color=cyan,bottom color=white,draw=cyan]
      plot[smooth,tension=0.5] coordinates {%
        (-1.5,1.5) (-1.375,1.55) (-1.25,1.5) (-1.125,1.45)
        (-1,1.5) (-.875,1.55) (-.75,1.5) (-.625,1.45)
        (-.5,1.5) (-.375,1.55) (-.25,1.5) (-0.125,1.45)
        (0,1.5) (0.125,1.55) (.25,1.5) (0.375,1.45)
        (0.5,1.5) (0.625,1.55) (0.75,1.5) (0.875,1.45)
        (1,1.5) (1.125,1.55) (1.25,1.5) (1.375,1.45) (1.5,1.5)}
      -- (1.5,1.5) -- (1.5,0) -- (-1.5,0) -- (-1.5,1.5) -- cycle;
  \end{scope}
  \tkzRoundFlaskSimple
}

% ============================================================
% 26. 分液漏斗 (Ampoule à décanter / Separate Funnel)
% ============================================================
% 相 : 水相 / 有机相
\definecolor{AqueousPhaseColor}{cmyk}{0.215686,0.062745,0,0}
\definecolor{OrganicPhaseColor}{cmyk}{0.223529,0.094118,0.501961,0}

% pst-labo: \tkzSeparatoryFunnel 设 linewidth=0.04 (用户单位, 随 unit 缩放)
% TikZ: 用 \chemlabLineWidth (基础 0.053cm × scale) 近似, 匹配其他器皿线宽
\newcommand{\tkzSeparatoryFunnelBody}{%
  \chemlabUpdateLineWidths
  \draw[line width=\chemlabLineWidth,line cap=round,line join=round]
    (-0.65,15) -- (-0.55,13.75)
    .. controls (-1.55,13.54) and ( -1.92,12.7) .. (-1.92,11.65)
    .. controls (-1.92,10.6) and (-0.3,9.55) .. (-0.3,8.5)
    -- (-0.3,2)
    .. controls (-0.3,1.5) and (-0.14,1.5) .. (-0.14,1)
    -- (-0.14,-1)
    -- (0.14,-1)
    -- (0.14,1)
    .. controls (0.14,1.5) and (0.3,1.5) .. (0.3,2)
    -- (0.3,8.5)
    .. controls (0.3,9.55) and (1.92,10.6) .. (1.92,11.65)
    .. controls ( 1.92,12.7) and (1.55,13.54) .. (0.55,13.75)
    -- (0.65,15);
}

% 用于 clip 的轮廓
\newcommand{\tkzSeparatoryFunnelOutlineClip}{%
  (-0.65,15) -- (-0.55,13.75)
  .. controls (-1.55,13.54) and ( -1.92,12.7) .. (-1.92,11.65)
  .. controls (-1.92,10.6) and (-0.3,9.55) .. (-0.3,8.5)
  -- (-0.3,2)
  .. controls (-0.3,1.5) and (-0.14,1.5) .. (-0.14,1)
  -- (-0.14,-1) -- (0.14,-1) -- (0.14,1)
  .. controls (0.14,1.5) and (0.3,1.5) .. (0.3,2)
  -- (0.3,8.5)
  .. controls (0.3,9.55) and (1.92,10.6) .. (1.92,11.65)
  .. controls (1.92,12.7) and (1.55,13.54) .. (0.55,13.75)
  -- (0.65,15) -- cycle
}

% 顶盖 (Bouchon) - 继承 linewidth=0.04
\newcommand{\tkzStopper}{%
  \chemlabUpdateLineWidths
  \filldraw[fill=gray!10,draw=gray,line width=\chemlabLineWidth,line cap=round,line join=round]
    (-0.67,15.25) -- (0.67,15.25) -- (0.57,14) -- (-0.57,14) -- cycle;
  \filldraw[fill=gray!5,draw=gray,line width=\chemlabLineWidth,line cap=round,line join=round]
    ([shift={(0.555,15.365)}]-90:0.1145) arc[start angle=-90,end angle=-180,radius=0.1145]
    -- ([shift={(0,15.92)}]-45:0.67) arc[start angle=-45,end angle=225,radius=0.67]
    -- ([shift={(-0.555,15.365)}]0:0.1145) arc[start angle=0,end angle=-90,radius=0.1145]
    -- cycle;
}

% 开启的旋塞 (TapOpen)
% pst-labo: \psframe[fillstyle=solid] 默认 fillcolor=white (PSTricks 默认)
% TikZ: \fill 默认黑色, 必须显式 fill=white
\newcommand{\tkzTapOpen}{%
  \chemlabUpdateLineWidths
  \filldraw[fill=white,line width=\chemlabLineWidth] (-.2,0.5) rectangle (.2,1);
  \draw[line width=\chemlabLineWidth] (-.2,0.4) rectangle (-.3,1.1);
  \draw[line width=\chemlabLineWidth] (.2,0.4) rectangle (.3,1.1);
  \draw[line width=\chemlabLineWidth] (.3,0.6) -- (.4,0.6) -- (.5,0.3) -- (.9,0.3) -- (.9,1.2) -- (.5,1.2)
        -- (.4,0.9) -- (.3,0.9) -- cycle;
  \fill[cyan!50] (-.6,0.4) -- (-.85,0.6) -- (-.85,0.9) -- (-.6,1.1) -- (-.3,1.1) -- (-.3,0.4) -- cycle;
  \draw[line width=\chemlabLineWidth] (-0.6,0.4) -- (-0.6,1.1);
  \draw[line width=\chemlabLineWidth] (0.5,0.3) -- (0.5,1.2);
}

% 关闭的旋塞 (TapClosed)
% pst-labo: \psframe[fillstyle=solid] 默认 fillcolor=white
\newcommand{\tkzTapClosed}{%
  \chemlabUpdateLineWidths
  \filldraw[fill=white,line width=\chemlabLineWidth] (-.2,0.5) rectangle (.2,1);
  \draw[line width=\chemlabLineWidth] (-.2,0.4) rectangle (-.3,1.1);
  \draw[line width=\chemlabLineWidth] (.2,0.4) rectangle (.3,1.1);
  \draw[line width=\chemlabLineWidth] ([shift={(0.35,0.6)}]180:0.05) arc[start angle=180,end angle=90,radius=0.05]
        -- (0.4,0.65) -- (0.9,0.65) -- (0.9,0.85) -- (0.35,0.85)
        -- ([shift={(0.35,0.9)}]-90:0.05) arc[start angle=-90,end angle=-180,radius=0.05]
        -- cycle;
  \draw[line width=\chemlabLineWidth] (0.5,0.65) -- (0.5,0.85);
  \fill[cyan!50] (-.6,0.4) -- (-.85,0.6) -- (-.85,0.9) -- (-.6,1.1) -- (-.3,1.1) -- (-.3,0.4) -- cycle;
  \draw[line width=\chemlabLineWidth] (-0.6,0.4) -- (-0.6,1.1);
}

% \tkzSeparatoryFunnel  (主调用)
% 可选参数 : stopper=true/false, TapOpen=true/false, liquidLevel1, liquidLevel2
\newcommand{\tkzSeparatoryFunnel}[1][]{%
  \begingroup
  \chemlabUpdateLineWidths
  \pgfkeys{
    /tkzSeparatoryFunnel/.cd,
    stopper/.is if=TkzTestTubeStopper,
    TapOpen/.is if=TkzTapOpen,
    liquidLevel1/.store in=\tkzTestTubeLiquidLevelI,
    liquidLevel2/.store in=\tkzTestTubeLiquidLevelII,
    #1}
  % 默认值
  \ifdefined\tkzTestTubeLiquidLevelI\else\def\tkzTestTubeLiquidLevelI{11}\fi
  \ifdefined\tkzTestTubeLiquidLevelII\else\def\tkzTestTubeLiquidLevelII{12}\fi
  % 水相
  \begin{scope}
    \clip \tkzSeparatoryFunnelOutlineClip;
    \fill[AqueousPhaseColor] (-2,1) rectangle (2,\tkzTestTubeLiquidLevelI);
  \end{scope}
  % 有机相
  \begin{scope}
    \clip \tkzSeparatoryFunnelOutlineClip;
    \fill[OrganicPhaseColor] (-2,\tkzTestTubeLiquidLevelI)
          rectangle (2,\tkzTestTubeLiquidLevelII);
  \end{scope}
  \ifTkzTestTubeStopper\tkzStopper\fi
  \tkzSeparatoryFunnelBody
  \ifTkzTapOpen\tkzTapOpen\else\tkzTapClosed\fi
  % 出口管口 (pst-labo: \psellipse[fillstyle=solid,fillcolor=white] 既有白色填充又有黑色轮廓)
  \filldraw[fill=white,line width=\chemlabLineWidth] (0,-1) ellipse (0.15 and 0.25);
  % 瓶口圆点 (填充, 无需线宽)
  \fill[darkgray] (-0.675,15) circle (0.03);
  \fill[darkgray] (0.675,15) circle (0.03);
  \draw[darkgray,line width=\chemlabLineWidth] (-0.65,15) -- (0.65,15);
  % 颈与瓶身接合
  \draw[gray,line width=\chemlabLineWidth] (-0.55,13.75) -- (0.55,13.75);
  \endgroup
}

% ============================================================
% 27. 支架 (pstSupport)
% ============================================================
% pst-labo: \psset{dimen=m,linewidth=0.04} — linewidth=0.04 是用户单位, 随 unit 缩放
% TikZ: 用 \chemlabLineWidth (基础 0.053cm × scale), 与其他器皿一致
\newcommand{\tkzStand}{%
  \chemlabUpdateLineWidths
  \begin{scope}[line width=\chemlabLineWidth]
    \filldraw[fill=woodColor,draw=black,line join=round]
        (-2,-6) rectangle (5,-5);
    \filldraw[fill=woodColor,draw=black,line join=round]
        (2.75,2) rectangle (3.25,16);
    \filldraw[fill=woodColor,draw=black,line join=round]
        ([shift={(2.75,1.5)}]90:0.5) arc[start angle=90,end angle=180,radius=0.5]
        -- (2.25,1.5) -- (2.25,-5) -- (3.75,-5) -- (3.75,1.5)
        -- ([shift={(3.25,1.5)}]0:0.5) arc[start angle=0,end angle=90,radius=0.5]
        -- cycle;
    \draw (2,11.5) -- (1.55,10.5);
    \filldraw[fill=woodColor,draw=black,line join=round]
        (-2.5,11.5) -- (-2,11.5) -- (-1.55,10.5) -- (-2.5,10.5) -- cycle;
    \filldraw[fill=woodColor,draw=black,line join=round]
        (4,11.5) -- (2,11.5) -- (1.55,10.5) -- (4,10.5) -- cycle;
    \draw (-2.5,10.5) rectangle (4,11.5);
    % 旋钮: pst-labo \pscustom[style=woodColor]{...} 填充会自动闭合, 但描边只画实际路径
    % (不画对角闭合边). TikZ: \fill 自动闭合填充, \draw 画开放路径 (无 -- cycle).
    \fill[woodColor]
        (4,11.2) -- (4.5,11.2)
        -- ([shift={(5,11)}]158.2:0.5385) arc[start angle=158.2,end angle=-158.2,radius=0.5385]
        -- (4.5,10.8) -- (4,10.8) -- cycle;
    \draw
        (4,11.2) -- (4.5,11.2)
        -- ([shift={(5,11)}]158.2:0.5385) arc[start angle=158.2,end angle=-158.2,radius=0.5385]
        -- (4.5,10.8) -- (4,10.8);
    \draw (4.5,10.8) -- (4.5,11.2);
  \end{scope}
}

\endinput

%    \begin{tikzpicture}
%       \BuretteCorps  ... 等等
%    \end{tikzpicture}
%
%  映射规则 :
%    \psline(A)(B)(C)              -> \draw (A) -- (B) -- (C);
%    \psline{->}(A)(B)             -> \draw[->] (A) -- (B);
%    \pspolygon(A)(B)(C)           -> \filldraw (A) -- (B) -- (C) -- cycle;
%    \psframe(x1,y1)(x2,y2)        -> \draw (x1,y1) rectangle (x2,y2);
%    \psframe*[linecolor=c]        -> \fill[c] (..) rectangle (..);
%    \psframe[framearc=r]          -> \draw[rounded corners=r] (...) rectangle (...);
%    \pscircle(c){r}               -> \draw (c) circle (r);
%    \pscircle*[linecolor=c](c){r} -> \fill[c] (c) circle (r);
%    \psellipse(c)(a,b)            -> \draw (c) ellipse (a and b);
%    \psarc(c){r}{a1}{a2}          -> \draw ([shift={(c)}]a1:r) arc[start angle=a1,end angle=a2,radius=r];
%    \psarcn(c){r}{a1}{a2}         -> \draw ([shift={(c)}]a1:r) arc[start angle=a1,end angle=a2,radius=r];   (注意 : PSTricks 中 arcn 是逆时针 / 角度递减)
%    \psbezier(A)(B)(C)(D)         -> \draw (A) .. controls (B) and (C) .. (D);
%    \pscurve(P1)(P2)...(Pn)       -> \draw plot[smooth,tension=0.5] coordinates {(P1) (P2) ... (Pn)};
%    \pscustom{ ... \fill}         -> \filldraw / \path[fill] (整段路径) -- cycle;
%    \rput(x,y){...}               -> \begin{scope}[shift={(x,y)}] ... \end{scope}
%    \rput{angle}(x,y){...}        -> \begin{scope}[shift={(x,y)},rotate=angle] ... \end{scope}
%    \multido{\i=a+b}{n}{...}      -> \foreach \i in {a,a+b,...,a+(n-1)*b} {...}
%    \psclip{<path>} ... \endpsclip-> \begin{scope} \clip <path>; ... \end{scope}
%    fillstyle=solid,fillcolor=c   -> fill=c
%    fillstyle=gradient,gradbegin=..,gradend=.. -> \shade[left color=..,right color=..]
%    fillstyle=hlines              -> pattern=horizontal lines
%    doubleline=true               -> double
%    linestyle=dashed              -> dashed
%    framearc=r                    -> rounded corners=r (近似)
%    linearc=r                     -> rounded corners=r
% ============================================================

\NeedsTeXFormat{LaTeX2e}
\ProvidesFile{tikz-chemlabObj.tex}[2026/08/03 v1.0.0]

\RequirePackage{tikz}
\usetikzlibrary{arrows.meta,calc,patterns,patterns.meta,shadings,decorations.pathmorphing}

% ------------------------------------------------------------
% 颜色定义 (对应原包中使用的颜色)
% ------------------------------------------------------------
\definecolor{LightGray}{gray}{0.80}      % 浅灰
\definecolor{VeryLightGray}{cmyk}{0,0,0,0.1}  % 极浅灰 (label 标签背景)
\definecolor{LightBlue}{rgb}{0.68,0.85,0.90}
\definecolor{LightBlue1}{rgb}{0.60,0.80,1.00}
\definecolor{VeryPaleOrange}{cmyk}{0,0.1,0.3,0}  % 滴定管液体默认色
\definecolor{PaleOrange}{rgb}{1.00,0.78,0.55}
\definecolor{Brown}{rgb}{0.50,0.25,0.00}
\definecolor{woodColor}{rgb}{0.957,0.84,0.69}
\definecolor{darkWoodColor}{rgb}{0.50,0.25,0.00}
\definecolor{flameAndGrill}{rgb}{0.95,0.55,0.10}  % 估算颜色
\definecolor{Beige}{rgb}{0.96,0.96,0.86}  % Champagne 液体默认色

% ------------------------------------------------------------
% 公共缺省值 (供调用方按需修改)
% ------------------------------------------------------------
\tikzset{
  Distillation/.cd,
  DistillateColor/.store in=\DistillationDistillateColor,
  DistillateColor=LightBlue1,
  MixtureStyle/.style={fill=LightBlue},
}
\newcommand{\tkzDistillationDistillateColor}{LightBlue1}
\newcommand{\tkzDistillationMixtureStyle}{LightBlue}
\newcommand{\tkzTestTubeLiquidLevelI}{11}     % liquidLevel1
\newcommand{\tkzTestTubeLiquidLevelII}{12}    % liquidLevel2
\newif\ifTkzTapOpen
\newif\ifTkzTestTubeStopper

% ============================================================
% 1. 滴定管 (Burette)
% ============================================================
\newcommand{\tkzBuretteBody}{%
  \draw (-0.5,11.5) -- (-0.25,11) -- (-0.25,3) -- (0,2.75) -- (0,0.5);
  \draw (0.2,0.7) -- (0.2,2.75) -- (0.5,3) -- (0.5,11) -- (0.75,11.5);
}

\newcommand{\tkzBuretteStopcock}{%
  % pst-labo: \psframe[fillstyle=solid] → 填充默认白色 + 黑色边框
  \filldraw[fill=white] (-.1,1.6) rectangle (.3,2.1);
  \draw (-.1,1.5) rectangle (-.2,2.2);
  \draw (.3,1.5) rectangle (.4,2.2);
  \draw (.4,1.7) -- (.5,1.7) -- (.6,1.4) -- (.8,1.4)
        -- (.8,2.3) -- (.6,2.3) -- (.5,2) -- (.4,2) -- cycle;
  % 左侧旋钮: pst-labo \pspolygon + \psframe with fillstyle=solid,fillcolor=LightBlue
  % 需要 LightBlue 填充 + 黑色外轮廓线
  \begin{scope}[fill=LightBlue,draw=black]
    \filldraw (-.5,1.5) -- (-.75,1.7) -- (-.75,2) -- (-.5,2.2)
          -- (-.3,2.2) -- (-.3,1.5) -- cycle;
    \filldraw (-.5,1.5) rectangle (-.2,2.2);
  \end{scope}
}

\newcommand{\tkzBuretteGraduation}{%
  % 滴定管刻度 (原 pst-labo 无显式线宽, 这里用细刻度线宽)
  \foreach \n in {3.0,3.32,...,10.68} { \draw[line width=\chemlabGraduationLineWidth] (-0.25,\n) -- (0,\n); }
}

% ============================================================
% 2. pH 计 (pH-mètre)
% ============================================================
\newcommand{\tkzTitrationPHMeter}{%
  % boitier
  \filldraw[fill=white,draw=black,rounded corners=0.125cm] (-1,0) rectangle (1,3);
  \fill[LightGray] (-0.75,2) rectangle (0.75,2.75);
  \node[white] at (0,2.375) {PH};
  \draw (0,1.25) circle (0.3);
  \draw[->] (0,1.25) -- (0,1.5);
  \node at (0.65,1.5) {$^\circ$C};
  \draw (-0.5,0.5) circle (0.2);
  \draw (0.5,0.5) circle (0.2);
  % électrode — 完全透明填充 (不遮挡烧杯刻度线)
  % pst-labo: \psframe[framearc=0.5](2.25,2.5)(2.75,4.5) 外框圆角
  %   顶部圆角被下方 \psframe[LightGray](sharp) 覆盖, 底部圆角被 \pscircle 覆盖
  %   实际可见的外框是直角矩形, 此处用 sharp corners 匹配 pst-labo 可见效果
  \draw (2.25,2.5) rectangle (2.75,4.5);
  % 顶部 LightGray 填充 + 描边 (sharp corners)
  \filldraw[fill=LightGray,draw=black] (2.25,4) rectangle (2.75,4.5);
  \draw (2.5,2.25) circle (0.25);
  % fil
  \draw[line width=0.05cm,LightBlue,rounded corners=0.2cm]
        (2.5,4.5) -- (2.5,5) -- (1.25,4.75) -- (0,5) -- (0,3);
}

% ============================================================
% 3. 滴定支架 (Support dosage)
% ============================================================
\newcommand{\tkzTitrationStand}{%
  % pst-labo: \psframe(-1.5,0)(1.5,1.5) + {\psset{doubleline=true} ...}
  % 线宽: pst-labo 在 \tkzTitration 中设 linewidth=0.053 (×0.4cm = 0.0212cm ≈ 0.6pt)
  \chemlabUpdateLineWidths
  % 平台: 无填充, 只有边框
  \draw[line width=\chemlabLineWidth] (-1.5,0) rectangle (1.5,1.5);
  % 红色按钮: \pscircle*[linecolor=red] → 红色填充, 双线边框也是红色 (不可见)
  \fill[red] (0,0.75) circle (0.2);
  % 回字型矩形按钮: 无填充, 双线边框
  % PSTricks doubleline: 指定坐标在内外线之间
  % 实测: 外线向外偏移 0.085 units, 内线向内偏移 0.034 units
  \draw[line width=\chemlabLineWidth] ({-1.25-0.085},{0.3-0.085}) rectangle ({-0.75+0.085},{1.2+0.085});
  \draw[line width=\chemlabLineWidth] ({-1.25+0.034},{0.3+0.034}) rectangle ({-0.75-0.034},{1.2-0.034});
  % 环形旋钮: 无填充, 双线边框
  % 外线向外偏移, 内线向内偏移
  \draw[line width=\chemlabLineWidth] (1,0.75) circle ({0.25+0.085});
  \draw[line width=\chemlabLineWidth] (1,0.75) circle ({0.25-0.034});
}

% ============================================================
% 4. 磁力搅拌子 (Aimant)
% ============================================================
\newcommand{\tkzTitrationStirBar}{%
  % pst-labo: \psset{linewidth=0.053} 在外层 \tkzTitration 中设置, 随 scale 缩放
  \chemlabUpdateLineWidths
  \draw[line width=\chemlabLineWidth,rounded corners=0.075cm,fill=white] (-0.5,0.1) rectangle (0.5,0.4);
  \draw[line width=\chemlabLineWidth,rounded corners=0.025cm,fill=white] (-0.05,0) rectangle (0.05,0.5);
}

% ============================================================
% 5. 移液管 (Pipette)
% ============================================================
\newcommand{\tkzPipetteBody}{%
  \chemlabUpdateLineWidths\begin{scope}[line width=\chemlabLineWidth,rounded corners=0.05cm]
    \draw (-0.05,0) -- (-0.15,0.2) -- (-0.15,6)
          -- (0.15,6) -- (0.15,0.2) -- (0.05,0) -- cycle;
    \fill[white] (0,6) ellipse (0.15 and 0.05);
    \foreach \Level in {0.5,1.0,...,5.0} {%
      \draw[line width=\chemlabGraduationLineWidth] (0,\Level) -- (0.1,\Level);}
  \end{scope}
}

\newcommand{\tkzLargePipetteBody}{%
  \chemlabUpdateLineWidths\begin{scope}[line width=\chemlabLineWidth,rounded corners=0.05cm]
    \draw (0,0) -- (0,1) -- (-0.23,2.5) -- (-0.23,22.5)
          arc[start angle=-80,end angle=80,radius=0.45];
  \end{scope}
}

% ============================================================
% 6. 液滴 (Goutte)
% ============================================================
% 原文件中 \pst@goutte 被重复定义了两次, 此处使用后者 (unit=0.6 的版本)
% pst-labo: \psset{unit=0.6}\psbezier[linewidth=0.017] — linewidth=0.017 是
%   0.017 用户单位 = 0.017*0.6 = 0.0102cm (随 goutte 的 LOCAL unit=0.6 缩放, 
%   不随外层 unit 变化, 因为 \psset{unit=0.6} 覆盖了外层设置)
% TikZ: line width=0.0102cm (ABSOLUTE, 匹配 pst-labo 的固定 0.0102cm)
\newcommand{\tkzDrop}{%
  \begin{scope}[scale=0.6]
    \draw[line width=0.0102cm]
      (0,0) .. controls (0.25,-.4) and (-.25,-.4) .. (0,0);
  \end{scope}
}

% ============================================================
% 7. 量筒 (Eprouvette)
% ============================================================
\newcommand{\tkzGraduatedCylinderBody}{%
  \draw (-0.5,5) -- (-0.5,0) -- (0.5,0) -- (0.5,4.9) -- (0.6,4.93);
}

\newcommand{\tkzGraduatedCylinderRim}{%
  \draw plot[smooth,tension=0.5] coordinates {%
    (0.6,-0.07) (0.498,0.008) (0.483,0.026) (0.433,0.05)
    (0.353,0.0707) (0.25,0.0866) (0.13,0.0966) (0,0.1)
    (-0.13,0.0966) (-0.25,0.0866) (-0.353,0.0707) (-0.433,0.05)
    (-0.483,0.026) (-0.5,0) (-0.483,-0.026) (-0.353,-0.0707)
    (-0.25,-0.0866) (-0.13,-0.096) (0,-0.1) (0.13,-0.096)
    (0.25,-0.0866) (0.353,-0.0707) (0.433,-0.05)
    (0.498,-0.026) (0.6,-0.07)};
}

\newcommand{\tkzGraduatedCylinderBase}{%
  \fill[white]
    ([shift={(-0.95,0.5)}]0:0.45) arc[start angle=0,end angle=-90,radius=0.45]
    -- ([shift={(-0.95,0.025)}]90:0.025) arc[start angle=90,end angle=270,radius=0.025]
    -- (-0.95,0) -- (0.95,0)
    -- ([shift={(0.95,0.025)}]-90:0.025) arc[start angle=-90,end angle=90,radius=0.025]
    -- ([shift={(0.95,0.5)}]-90:0.45) arc[start angle=-90,end angle=-180,radius=0.45]
    -- plot[smooth,tension=0.5] coordinates {(0.5,0.4) (0,0.25) (-0.5,0.4)}
    -- cycle;
  \draw
    ([shift={(-0.95,0.5)}]0:0.45) arc[start angle=0,end angle=-90,radius=0.45]
    -- ([shift={(-0.95,0.025)}]90:0.025) arc[start angle=90,end angle=270,radius=0.025]
    -- (-0.95,0) -- (0.95,0)
    -- ([shift={(0.95,0.025)}]-90:0.025) arc[start angle=-90,end angle=90,radius=0.025]
    -- ([shift={(0.95,0.5)}]-90:0.45) arc[start angle=-90,end angle=-180,radius=0.45]
    -- plot[smooth,tension=0.5] coordinates {(0.5,0.4) (0,0.25) (-0.5,0.4)};
}

\newcommand{\tkzGraduatedCylinderGraduation}{%
  % 量筒刻度 : 主刻度 1.5× 细刻度 (与烧杯一致)
  \pgfmathsetlengthmacro{\chemlab@eprouvette@major}{1.5*\chemlabGraduationLineWidth}%
  \foreach \nyD in {0.5,1.0,...,4.5} { \draw[line width=\chemlab@eprouvette@major] (-0.4,\nyD) -- (-0.1,\nyD); }
  \foreach \nyd in {0.5,0.6,...,4.4} { \draw[line width=\chemlabGraduationLineWidth] (-0.4,\nyd) -- (-0.25,\nyd); }
}

% ============================================================
% 8. 分馏装置 (Distillation fractionnée)
% ============================================================
% \pst@goutte 这里又被定义了一次, 在 TikZ 版中沿用同一个 \tkzDrop

\newcommand{\tkzColumn}{%
  % pst-labo: linearc=0.3 (×unit, 不是绝对值) → TikZ: rounded corners=0.12cm (= 0.3 × 0.4cm)
  % 尖刺 (pointeG/D) 先画并用 clip 裁剪到蒸馏柱内部, 再画轮廓线覆盖切边。
  %   pst-labo 中 linearc=0.3 (R=0.3 user units), 尖刺端点恰在圆弧切点附近, 几乎
  %   不可见。TikZ 的 rounded corners 实际 R 比标称值大 ~2×, 圆弧切点更低,
  %   尖刺端点 (在 y_col=8 和 7.9) 会突出到圆弧外侧形成多余线头。用 clip 将
  %   尖刺裁剪到柱内, 再画轮廓线覆盖切边, 消除突出线头。
  \begin{scope}
    \clip[rounded corners=0.12cm]
      (-.25,10.5) -- (-.25,8.5) -- (-.75,8) -- (-.75,2.5) -- (-.25,2) -- (-.25,0)
      -- (0.25,0) -- (0.25,2) -- (.75,2.5) -- (.75,8) -- (0.25,8.5) -- (0.25,9.36)
      -- (0.25,10.5) -- cycle;
    % 左侧尖刺 (pointeG)
    \foreach \i in {3,4,...,8} {%
      \draw[line width=0.4\pgflinewidth] (-.75,\i) -- (-.25,\i-0.5)
          -- (-.75,\i-0.1);}
    % 右侧尖刺 (pointeD)
    \foreach \i in {3,4,...,8} {%
      \draw[line width=0.4\pgflinewidth] (.75,\i) -- (.25,\i-0.5)
          -- (.75,\i-0.1);}
  \end{scope}
  \draw[rounded corners=0.12cm]
    (-.25,10.5) -- (-.25,8.5) -- (-.75,8) -- (-.75,2.5)
    -- (-.25,2) -- (-.25,0);
  \draw[rounded corners=0.12cm]
    (0.25,0) -- (0.25,2) -- (.75,2.5) -- (.75,8)
    -- (0.25,8.5) -- (0.25,9.36);
  % 上部断口 (line 3a/3b):
  %   pst-labo 源码 line 3a: \psline(0.25,8.5)(0.25,9.30) 用于填充 linearc 圆角
  %   留下的间隙。实测发现 TikZ 的 rounded corners 实际圆弧半径比标称值大
  %   ~2× (在 scale=0.6 下 rounded corners=0.3cm 给出 R≈0.62cm), 圆弧已覆盖到
  %   x=0.25 处 y_col=8.93 附近, 故 line 3a 在任何 scale/line width 下均被圆弧
  %   覆盖, 不可见。为避免在编号图等彩色调试场景中残留多余线段, 此处不再
  %   单独绘制 line 3a。
  \draw (0.25,9.5) -- (0.25,10.5);
  % 填充 / 阴影
  \begin{scope}[pattern={Lines[angle=45,distance={1pt/sqrt(2)}]},pattern color=gray]
    \fill (-.25,10) rectangle (-.05,10.5);
    \fill (.25,10)  rectangle (0.05,10.5);
  \end{scope}
  \foreach \n in {2.5,3.5,...,7.5} {%
    \begin{scope}[shift={(-.25,\n)}]\tkzDrop\end{scope}
    \begin{scope}[shift={(.25,\n)}]\tkzDrop\end{scope}}
  \begin{scope}[shift={(-0.25,-.1)}]\tkzDrop\end{scope}
  \begin{scope}[shift={(0.25,-.1)}]\tkzDrop\end{scope}
}

% 公共别名: pst-labo 中 \tkzColumn 是用户可调用的公共宏 (小写)
% tikz-chemlab 内部用 \tkzColumn (驼峰), 这里注册公共别名保持兼容
\let\tkzColumn\tkzColumn

% réfrigérant (condenseur)
\newcommand{\tkzCondenser}{%
  \draw ([shift={(3,0)}]90:0.5) arc[start angle=90,end angle=270,radius=0.5];
  % pst-labo: \psarc(7,0){0.5}{270}{90} → 逆时针经过 0°/右侧
  \draw ([shift={(7,0)}]270:0.5) arc[start angle=270,end angle=450,radius=0.5];
  \fill[white] (0,-.05) rectangle (9,0.05);
  \begin{scope}[line width=0.4pt]
    \draw (-.075,.05) -- (9,0.05);
    \draw (0.02,-.05) -- (9,-.05);
  \end{scope}
  \draw (3,0.5) -- (3.5,0.5);
  \draw (3,-.5) -- (6.5,-.5);
  \draw (3.6,0.5) -- (7,0.5);
  \begin{scope}[line width=0.4pt]
    \draw (3.5,0.5) -- (3.5,1.5);
    \draw (3.6,0.5) -- (3.6,1.5);
    \draw (6.6,-.5) -- (6.6,-1.5);
    \draw (6.5,-.5) -- (6.5,-1.5);
  \end{scope}
  \draw (6.6,-.5) -- (7,-.5);
  \draw[->] (6.55,-2.5) -- (6.55,-1.5);
  \draw[->] (3.55,1.5)  -- (3.55,2.5);
}

\newcommand{\tkzThermometer}{%
  % pst-labo: \psline[linearc=0.05](-.05,.25)(-.05,2.75)(0.05,2.75)(0.05,0.25)
  % linearc=0.05 × unit, TikZ: rounded corners=0.05 (受 scale 缩放)
  \draw[rounded corners=0.025cm]
    (-.05,.25) -- (-.05,2.75) -- (0.05,2.75) -- (0.05,0.25) -- cycle;
  % pst-labo: \psframe[framearc=0.8,fillstyle=solid,fillcolor=red](-.15,-.25)(0.15,0.25)
  % framearc=0.8: 80% of half-min(0.3,0.5)=0.15 → radius=0.12
  \filldraw[fill=red,draw=black,rounded corners=0.06cm] (-.15,-.25) rectangle (0.15,0.25);
  % pst-labo: \psline[linecolor=red](0,2) — 红色竖线, 使用默认线宽
  % 在 pst-labo 渲染中较粗, 看起来像红色填充
  \draw[red,line width=0.8pt] (0,0) -- (0,2);
}

\newcommand{\tkzHeater}{%
  \draw (1.17,2.5) -- (0,2.5) -- (0,0) -- (5.5,0) -- (5.5,2.5) -- (4.33,2.5);
}

% 内部小量筒轮廓 - 路径版本 (供 clip 使用)
\newcommand{\tkzGraduatedCylinderPath}{%
  (-0.5,6) -- (-0.5,.5)
  -- ([shift={(-1,0.5)}]0:0.5) arc[start angle=0,end angle=-90,radius=0.5]
  -- (-1,0) -- (1,0)
  -- ([shift={(1,0.5)}]270:0.5) arc[start angle=270,end angle=180,radius=0.5]
  -- (0.5,0.5) -- (0.5,6) -- cycle
}

% 内部小量筒 (描边)
% pst-labo: \psset{linewidth=1.5\pslinewidth} — \pslinewidth 默认 0.8pt (ABSOLUTE)
% → 1.5*0.8pt = 1.2pt (绝对值, 不随 unit/scale 缩放)
% 注意: TikZ 的 [line width=N\pgflinewidth] 有循环引用 bug — \pgflinewidth 在
% 乘法前被更新, 导致 line width=3\pgflinewidth 实际给出 3.6pt 而非 1.2pt。
% 必须用显式 pt 值。
\newcommand{\tkzGraduatedCylinderSmall}{%
  \begin{scope}[line width=0.72pt,line cap=round,line join=round]
    \draw (-0.5,6) -- (-0.5,.5)
      -- ([shift={(-1,0.5)}]0:0.5) arc[start angle=0,end angle=-90,radius=0.5]
      -- (-1,0) -- (1,0)
      -- ([shift={(1,0.5)}]270:0.5) arc[start angle=270,end angle=180,radius=0.5]
      -- (0.5,0.5) -- (0.5,6);
  \end{scope}
}

% 带液体的量筒 (内部使用, 重命名以避免与主文件 \pstEproulette 冲突)
\newcommand{\tkzGraduatedCylinderWithLiquid}{%
  \begin{scope}
    \clip \tkzGraduatedCylinderPath;
    % pst-labo: \psframe*[linecolor=\psk@Distillation@DistillateColor]
    \fill[\tkzDistillateColor] (-1.5,-1) rectangle (1.5,1);
  \end{scope}
  \tkzGraduatedCylinderSmall
}

% (注 : 原 PSTricks 中 \pst@roundFlask 已合并到 \tkzRoundFlaskSimple, 此处不再单独定义,
%  避免与 tikz-chemlab.tex 中的用户命令 \tkzRoundFlask 冲突)

% 自定义密集斜线图案 (匹配 pst-labo hlines 的密度)
{\pgfpointorigin}{\pgfpoint{0.012cm}{0.012cm}}{\pgfpoint{0.012cm}{0.012cm}}
{\pgfsetlinewidth{0.0028cm}\pgfpathmoveto{\pgfpoint{0cm}{0cm}}\pgfpathlineto{\pgfpoint{0.012cm}{0.012cm}}\pgfusepath{stroke}}
{\pgfpointorigin}{\pgfpoint{0.012cm}{0.012cm}}{\pgfpoint{0.012cm}{0.012cm}}
{\pgfsetlinewidth{0.004cm}\pgfpathmoveto{\pgfpoint{0cm}{0cm}}\pgfpathlineto{\pgfpoint{0.012cm}{0.012cm}}\pgfusepath{stroke}}
% \pst@DistillationFractionnee 主图
\newcommand{\tkzDistillationFractionated}{%
  % --- roundFlask ---
  \def\flaskLocal{%
    % pst-labo: fillstyle=hlines,hatchwidth=0.007,hatchsep=0.03 → 水平线填充
    \begin{scope}[pattern={Lines[angle=45,distance={1pt/sqrt(2)}]},pattern color=gray]
      \fill (-.5,2.25) rectangle (-.25,3);
      \fill (0.5,2.25)  rectangle (0.25,3);
    \end{scope}
    % pst-labo 线宽: neck=0.053, arcs=0.03 (都是 ×unit 用户单位, 随 scale 缩放)
    % \chemlabUpdateLineWidths 已设置 \chemlabLineWidth = 0.053*scale
    \chemlabUpdateLineWidths
    \draw[line width=\chemlabLineWidth] (0.5,1.5) -- (0.5,3);
    \draw[line width=\chemlabLineWidth] (-.5,1.5) -- (-.5,3);
    % pst-labo: \psarc(0,0){1.581}{108.435}{71.565} 逆时针走长路 (经底部),
    %   左右侧弧可见, 底部被液体填充覆盖, 顶部无可见弧线。此处拆分为
    %   左右侧两段 (底部由下方轮廓弧绘制, 顶部不画)。
    % 右上侧弧: 71° → 0° (连接右颈底部到球体最右点)
    \draw[line width=\chemlabLineWidth] ([shift={(0,0)}]71.565:1.581) arc[start angle=71.565,end angle=0,radius=1.581];
    % 左上侧弧: 180° → 108° (连接球体最左点到左颈底部)
    \draw[line width=\chemlabLineWidth] ([shift={(0,0)}]180:1.581) arc[start angle=180,end angle=108.435,radius=1.581];
    % pst-labo: \psset{style=\psk@Distillation@MixtureStyle}\fill
    % 需要用 \shade (不是 \fill) 来产生渐变效果
    \shade[\tkzMixtureStyle] plot[smooth,tension=0.5] coordinates {%
      (-1.581,0) (-1.25,0.1) (-.75,-.1) (0,.1) (0.75,-.1) (1.25,.1) (1.581,0)}
      arc[start angle=0,end angle=-180,radius=1.581];
    % 底部轮廓弧: pst-labo \psarc[linewidth=0.03] → 0.03*unit, 用 \chemlabLineWidth 0.053
    % (pst-labo 中 0.03 vs 0.053 差异微小, 统一用 0.053*scale 避免额外变量)
    \draw[line width=\chemlabLineWidth] (0,0) ++(180:1.581) arc[start angle=180,end angle=360,radius=1.581];
  }%
  % --- 主图整体平移 ---
  \begin{scope}[shift={(0,-3)}]
    \begin{scope}[shift={(2.75,13.5)}]\tkzThermometer\end{scope}
    \begin{scope}[shift={(2.75,2.5)}]\flaskLocal\end{scope}
    \begin{scope}[shift={(2.75,4)}]\tkzColumn\end{scope}
    \begin{scope}[shift={(3,13.4)},rotate=-45]\tkzCondenser\end{scope}
    \begin{scope}[shift={(0,0)}]\tkzHeater\end{scope}
    \begin{scope}[shift={(9.4,0)}]\tkzGraduatedCylinderWithLiquid
      % 蒸馏装置量筒刻度 : 主刻度 1.5× 细刻度
      \pgfmathsetlengthmacro{\chemlab@distill@major}{1.5*\chemlabGraduationLineWidth}%
      \foreach \Level in {1,2,...,5} {\draw[line width=\chemlab@distill@major] (-.4,\Level) -- (-.1,\Level);}
      \foreach \Level in {0.2,0.4,...,4.8} {\draw[line width=\chemlabGraduationLineWidth] (-.4,\Level) -- (-.2,\Level);}
    \end{scope}
    \begin{scope}[shift={(9.37,7)}]\tkzDrop\end{scope}
    \begin{scope}[shift={(9.37,6)}]\tkzDrop\end{scope}
    \begin{scope}[shift={(9.37,4)}]\tkzDrop\end{scope}
    \begin{scope}[shift={(9.37,1.2)}]\tkzDrop\end{scope}
    % 加热盘 (搅拌器)
    \begin{scope}[shift={(0.75,.75)}]
      \draw (0,0) circle (.5);
      \foreach \i in {0,30,...,180} {\draw (\i:0.55) -- (\i:0.65);}
      \draw[->,line width=0.035cm] (0,0) -- (0,0.45);
    \end{scope}
  \end{scope}
}

% ============================================================
% 9. 本生灯 (Bec Bunsen) - 不带 grille
% ============================================================
\newcommand{\tkzBunsenBurner}{%
  \chemlabUpdateLineWidths\begin{scope}[line width=\chemlabLineWidth,line cap=round,line join=round]
    \draw (-1.25,0) rectangle (1.25,0.25);
    \draw (-.5,1.25) rectangle (0.5,2.25);
    \foreach \n in {-0.3,0.0,0.3} {\draw (\n,1.75) circle (0.1);}
    \draw (-.25,2.25) rectangle (0.25,4.25);
    \draw (0.25,1.25) -- (0.25,0.5) -- (1.25,0.25);
    \draw (-1.25,0.25) -- (-.25,0.5) -- (-0.25,0.75);
    \draw (-2.25,0.75) -- (-.25,0.75);
    \draw (-2.25,1) -- (-.25,1);
    \draw (-.25,0.875) ellipse (0.1 and 0.125);
    \fill[white] (-2.25,0.75) rectangle (-0.25,1);
    \draw (-2.25,0.75) -- (-0.25,0.75);
    \draw (-2.25,1) -- (-0.25,1) -- (-.25,1.25);
    \draw plot[smooth,tension=0.5] coordinates {(-0.25,0.5) (0,0.4) (0.25,0.5)};
    % flamme
    % pst-labo 源码: \rput(0,4.25){\psclip{\psbezier[...gradient...] (12点)}\pspolygon[...gradient...](-0.25,0)(0.25,0)(0,1)\endpsclip}
    % \psbezier 12个点 + BezierNArray padding (0,0) = 13点 = 4段贝塞尔
    % 经像素级验证, 正确路径从 (0,0) 出发, 经4段贝塞尔回到 (0,0), 无闭合直线:
    %   Seg 1: (0,0) → c1=P1=(-0.25,0), c2=P2=(-0.35,0.5) → P3=(-0.4,0.75) [左下→左中]
    %   Seg 2: P3=(-0.4,0.75) → c1=P4=(-0.35,1), c2=P5=(-0.25,1.5) → P6=(0.5,2) [左中→顶]
    %   Seg 3: P6=(0.5,2) → c1=P7=(0.25,1.5), c2=P8=(0.35,1) → P9=(0.4,0.75) [顶→右中]
    %   Seg 4: P9=(0.4,0.75) → c1=P10=(0.35,0.5), c2=P11=(0.25,0) → (0,0) [右中→底]
    % 渐变: gradbegin=PaleOrange (底), gradend=yellow (顶), gradmidpoint=0, gradangle=0 (垂直)
    \begin{scope}[shift={(0,4.25)}]
      % 外焰填充: PaleOrange (底) -> yellow (顶) 线性垂直渐变
      \shade[top color=yellow, bottom color=PaleOrange]
        (0,0) .. controls (-0.25,0) and (-0.35,0.5)
        .. (-0.4,0.75) .. controls (-0.35,1) and (-0.25,1.5)
        .. (0.5,2) .. controls (0.25,1.5) and (0.35,1)
        .. (0.4,0.75) .. controls (0.35,0.5) and (0.25,0)
        .. (0,0);
      % 内焰: 三角形 cyan(底)->white(顶) 渐变, 裁剪到外焰形状内
      \begin{scope}
        \clip
          (0,0) .. controls (-0.25,0) and (-0.35,0.5)
          .. (-0.4,0.75) .. controls (-0.35,1) and (-0.25,1.5)
          .. (0.5,2) .. controls (0.25,1.5) and (0.35,1)
          .. (0.4,0.75) .. controls (0.35,0.5) and (0.25,0)
          .. (0,0);
        \shade[top color=white, bottom color=cyan]
          (-0.25,0) -- (0.25,0) -- (0,1) -- cycle;
      \end{scope}
    \end{scope}
  \end{scope}
}

% 火焰 + 网格 (FlammeGrille 形状) - 路径版本 (供 clip 使用)
\newcommand{\tkzFlameGrillPath}{%
  ([shift={(0,0.75)}]270:0.75) arc[start angle=270,end angle=360,radius=0.75]
  -- ([shift={(1.5,0.75)}]180:0.75) arc[start angle=180,end angle=90,radius=0.75]
  -- (1.5,1.5) -- (-1.5,1.5)
  -- ([shift={(-1.5,0.75)}]90:0.75) arc[start angle=90,end angle=0,radius=0.75]
  -- ([shift={(0,0.75)}]180:0.75) arc[start angle=180,end angle=270,radius=0.75]
  -- cycle
}

\newcommand{\tkzFlameGrill}{%
  \fill[flameAndGrill] \tkzFlameGrillPath;
}

% ============================================================
% 10. 本生灯 - 带 grille
% ============================================================
\newcommand{\tkzBunsenBurnerGrill}{%
  \chemlabUpdateLineWidths\begin{scope}[line width=\chemlabLineWidth,line cap=round,line join=round]
    \draw (-1.25,0) rectangle (1.25,0.25);
    \draw (-.5,1.25) rectangle (0.5,2.25);
    \foreach \n in {-0.3,0.0,0.3} {\draw (\n,1.75) circle (0.1);}
    \draw (-.25,2.25) rectangle (0.25,4.25);
    \draw (0.25,1.25) -- (0.25,0.5) -- (1.25,0.25);
    \draw (-1.25,0.25) -- (-.25,0.5) -- (-0.25,0.75);
    \draw (-2.25,0.75) -- (-.25,0.75);
    \draw (-2.25,1) -- (-.25,1);
    \draw (-.25,0.875) ellipse (0.1 and 0.125);
    \fill[white] (-2.25,0.75) rectangle (-0.25,1);
    \draw (-2.25,0.75) -- (-0.25,0.75);
    \draw (-2.25,1) -- (-0.25,1) -- (-.25,1.25);
    \draw plot[smooth,tension=0.5] coordinates {(-0.25,0.5) (0,0.4) (0.25,0.5)};
    % flamme (clip 到 FlammeGrille)
    % 原 pst-labo 火焰分两层:
    %   外焰: \pst@FlammeGrille 整体填充 flameAndGrill (PaleOrange->yellow 渐变)
    %   内焰: 三角形 cyan->white 渐变 (在内焰区域)
    % 颜色和大小可通过 \chemlabset 调整
    \begin{scope}[shift={(0,4.25)},scale=\chemlabFlameScale]
      \begin{scope}
        \clip \tkzFlameGrillPath;
        % 外焰: PaleOrange (底) -> yellow (顶) 线性垂直渐变
        % 对应 pst-labo: fillstyle=gradient, gradbegin=PaleOrange, gradend=yellow
        % 渐变范围与火焰形状 (y=0..1.5) 一致, 避免底部颜色被稀释
        \shade[top color=\chemlabFlameOuterEndColor, bottom color=\chemlabFlameOuterColor]
          (-1.5,0) rectangle (1.5,1.5);
        % 内焰: cyan (底) -> white (顶) 线性垂直渐变 (三角形)
        \shade[top color=\chemlabFlameInnerEndColor, bottom color=\chemlabFlameInnerColor]
          (-0.25,0) -- (0.25,0) -- (0,1) -- cycle;
      \end{scope}
    \end{scope}
    \draw[dashed,line width=\chemlabEllipseLineWidth] (-2,5.75) -- (2,5.75);
  \end{scope}
}

% ============================================================
% 11. 烧瓶 (Ballon)
% ============================================================
% 路径版本 - 供 clip 使用 (闭合轮廓)
\newcommand{\tkzRoundFlaskSimplePath}{%
  ([shift={(-0.7,3.4)}]90:0.2) arc[start angle=90,end angle=0,radius=0.2]
  -- (-0.5,3.4) -- (-0.5,2.4)
  -- (-0.489,2.401)
  arc[start angle=113,end angle=427,radius=1.25]
  -- (0.5,2.4) -- (0.5,3.4)
  -- ([shift={(0.7,3.4)}]180:0.2) arc[start angle=180,end angle=90,radius=0.2]
  -- (0,3.5) -- cycle
}

\newcommand{\tkzRoundFlaskSimple}{%
  \chemlabUpdateLineWidths\begin{scope}[line width=\chemlabLineWidth,line cap=round,line join=round]
    % 左侧: 瓶口圆弧 + 颈部直线 + 球体大圆弧, 合并为单条路径以平滑过渡
    \draw ([shift={(-0.7,3.4)}]90:0.2) arc[start angle=90,end angle=0,radius=0.2]
          -- (-0.5,3.4) -- (-0.5,2.4)
          -- (-0.489,2.401) arc[start angle=113,end angle=427,radius=1.25];
    % 右侧: 球体终点(0.489,2.401) -> 颈部直线 -> 瓶口圆弧
    \draw (0.489,2.401) -- (0.5,2.4) -- (0.5,3.4)
          -- ([shift={(0.7,3.4)}]180:0.2) arc[start angle=180,end angle=90,radius=0.2];
    % 瓶口椭圆: 白色填充覆盖颈部上端圆头, 避免穿入椭圆
    % pst-labo: ballon椭圆用默认线宽 0.053 (与主轮廓一致, 不是 0.07)
    \filldraw[fill=white,draw=black,line width=\chemlabLineWidth] (0,3.6) ellipse (0.7 and 0.1);
  \end{scope}
}

\newcommand{\tkzRoundFlaskGrill}{%
  \chemlabUpdateLineWidths\begin{scope}[line width=\chemlabLineWidth,line cap=round,line join=round]
    % 左侧: 瓶口圆弧 + 颈部直线 + 球体大圆弧, 合并为单条路径
    \draw ([shift={(-0.7,2.15)}]90:0.2) arc[start angle=90,end angle=0,radius=0.2]
          -- (-0.5,2.15) -- (-0.5,1.4)
          -- (-0.489,1.401) arc[start angle=113,end angle=427,radius=1.25];
    % 右侧: 球体终点 -> 颈部直线 -> 瓶口圆弧
    \draw (0.489,1.401) -- (0.5,1.15) -- (0.5,3.4)
          -- ([shift={(0.7,2.15)}]180:0.2) arc[start angle=180,end angle=90,radius=0.2];
    % 瓶口椭圆: 白色填充覆盖颈部上端圆头, 避免穿入椭圆
    % pst-labo: ballon椭圆用默认线宽 0.053 (与主轮廓一致)
    \filldraw[fill=white,draw=black,line width=\chemlabLineWidth] (0,2.35) ellipse (0.7 and 0.1);
  \end{scope}
}

% ============================================================
% 12. 试管 (TubeEssais)
% ============================================================
\newcommand{\tkzTestTubeStopper}{%
  \chemlabUpdateLineWidths
  \begin{scope}[fill=LightGray,draw=black,line width=\chemlabLineWidth,line cap=round,line join=round]
    \filldraw (0,-0.3) ellipse (0.4 and 0.1);
    \filldraw (-0.4,-0.3) -- (-0.6,0.5) -- (0.6,0.5) -- (0.4,-0.3) -- cycle;
    \fill[draw=none] (0,-0.3) ellipse (0.4 and 0.1);
    \filldraw (0,0.5) ellipse (0.6 and 0.1);
  \end{scope}
  \draw[line width=\chemlabEllipseLineWidth] plot[smooth,tension=0.5]
        coordinates {(-0.5,0) (0,-0.07) (0.5,0)};
}

\newcommand{\tkzTestTubeBody}{%
  \chemlabUpdateLineWidths\begin{scope}[line width=\chemlabLineWidth,line cap=round,line join=round]
    % 合并左侧直线 + 底部圆弧 + 右侧直线为单条路径, 平滑转角
    \draw (-0.5,3.5) -- (-0.5,0.5)
          arc[start angle=180,end angle=360,radius=0.5]
          -- (0.5,3.5);
    % 试管口椭圆 (原 pst-labo 用 0.08cm, 比其他椭圆粗)
    \draw[line width=\chemlabTubeEllipseLineWidth] (0,3.5) ellipse (0.5 and 0.1);
  \end{scope}
}

\newcommand{\tkzTestTubeStraightTube}{%
  \draw[line width=0.5\pgflinewidth] (0.05,-1) -- (0.05,5);
  \draw[line width=0.5\pgflinewidth] (-0.05,-1) -- (-0.05,5);
  \tkzTestTubeStopper
  \draw[fill=white,rounded corners=0.5,line width=0.5\pgflinewidth]
       (-0.05,5) rectangle (0.05,0.5);
  \draw[white,line width=0.14cm] (-0.1,5) -- (0.1,5);
}

\newcommand{\tkzTestTubeElbowTube}{%
  \draw[line width=0.5\pgflinewidth] (0.05,-1) -- (0.05,0);
  \draw[line width=0.5\pgflinewidth] (-0.05,-1) -- (-0.05,0);
  \tkzTestTubeStopper
  % 弯管路径: pst-labo \pspolygon(...)(8,2)(8,1.9)...
  \filldraw[fill=white,rounded corners=0.05cm,line width=0.5\pgflinewidth]
       (-0.05,0.5) -- (-0.05,2) -- (8,2) -- (8,1.9)
       -- (0.05,1.9) -- (0.05,0.5) -- cycle;
  \draw[white,line width=0.14cm] (7.95,1.9) -- (8.05,2);
}

\newcommand{\tkzTestTubeUTube}{%
  \draw[line width=0.5\pgflinewidth] (0.05,-1) -- (0.05,0);
  \draw[line width=0.5\pgflinewidth] (-0.05,-1) -- (-0.05,0);
  \tkzTestTubeStopper
  \filldraw[fill=white,rounded corners=0.05cm,line width=0.5\pgflinewidth]
       (-0.05,0.5) -- (-0.05,2) -- (4.5,2) -- (4.5,-3)
       -- (4.4,-3) -- (4.4,1.9) -- (0.05,1.9) -- (0.05,0.5) -- cycle;
  \draw[white,line width=0.14cm] (4.3,-3) -- (4.6,-3);
}

\newcommand{\tkzTestTubeUTubeLong}{%
  \draw[line width=0.5\pgflinewidth] (0.05,-1) -- (0.05,0);
  \draw[line width=0.5\pgflinewidth] (-0.05,-1) -- (-0.05,0);
  \tkzTestTubeStopper
  \filldraw[fill=white,rounded corners=0.05cm,line width=0.5\pgflinewidth]
       (-0.05,0.5) -- (-0.05,2) -- (4.5,2) -- (4.5,-9)
       -- (4.4,-9) -- (4.4,1.9) -- (0.05,1.9) -- (0.05,0.5) -- cycle;
  \draw[white,line width=0.14cm] (4.3,-9) -- (4.6,-9);
}

\newcommand{\tkzCurvedTube}{%
  \draw[line width=0.5\pgflinewidth] (0.05,-1) -- (0.05,0);
  \draw[line width=0.5\pgflinewidth] (-0.05,-1) -- (-0.05,0);
  \tkzTestTubeStopper
  \begin{scope}[line width=0.4pt]
    \draw (-0.05,0.5) -- (-0.05,1);
    \draw ([shift={(0.5,1)}]180:0.55) arc[start angle=180,end angle=90,radius=0.55];
    \draw (0.5,1.55) -- (5.5,1.55);
    \draw ([shift={(5.5,1)}]90:0.55) arc[start angle=90,end angle=0,radius=0.55];
    \draw (6.05,1) -- (6.05,-8.8);
    \draw ([shift={(6.5,-8.8)}]180:0.45) arc[start angle=180,end angle=360,radius=0.45];
    \draw (6.95,-8.8) -- (6.95,-8.3);
    \draw (7.05,-8.3) -- (7.05,-8.8);
    \draw ([shift={(6.5,-8.8)}]0:0.55) arc[start angle=0,end angle=-180,radius=0.55];
    \draw (5.95,-8.8) -- (5.95,1);
    \draw ([shift={(5.5,1)}]0:0.45) arc[start angle=0,end angle=90,radius=0.45];
    \draw (5.5,1.45) -- (0.5,1.45);
    \draw ([shift={(0.5,1)}]90:0.45) arc[start angle=90,end angle=180,radius=0.45];
    \draw (0.05,1) -- (0.05,0.5);
    \draw (-0.05,0.5) -- (0.05,0.5);
  \end{scope}
  \fill[white] (-0.025,0.52) rectangle (0.020,0.6);
}

\newcommand{\tkzCurvedTubeShort}{%
  \draw[line width=0.5\pgflinewidth] (0.05,-1) -- (0.05,0);
  \draw[line width=0.5\pgflinewidth] (-0.05,-1) -- (-0.05,0);
  \tkzTestTubeStopper
  \draw (-0.05,0.5) -- (-0.05,1);
  \begin{scope}[line width=0.4pt]
    \draw ([shift={(0.5,1)}]180:0.55) arc[start angle=180,end angle=90,radius=0.55];
    \draw (0.5,1.55) -- (3.5,1.55);
    \draw ([shift={(3.5,1)}]90:0.55) arc[start angle=90,end angle=0,radius=0.55];
    \draw (4.05,1) -- (4.05,-2.8);
    \draw ([shift={(4.5,-2.8)}]180:0.45) arc[start angle=180,end angle=360,radius=0.45];
    \draw (4.95,-2.8) -- (4.95,-2.3);
    \draw (5.05,-2.3) -- (5.05,-2.8);
    \draw ([shift={(4.5,-2.8)}]0:0.55) arc[start angle=0,end angle=-180,radius=0.55];
    \draw (3.95,-2.8) -- (3.95,1);
    \draw ([shift={(3.5,1)}]0:0.45) arc[start angle=0,end angle=90,radius=0.45];
    \draw (3.5,1.45) -- (0.5,1.45);
    \draw ([shift={(0.5,1)}]90:0.45) arc[start angle=90,end angle=180,radius=0.45];
    \draw (0.05,1) -- (0.05,0.5);
    \draw (-0.05,0.5) -- (0.05,0.5);
  \end{scope}
  \fill[white] (-0.025,0.52) rectangle (0.020,0.6);
}

\newcommand{\tkzCurvedTubeShortDoubleTube}{%
  \begin{scope}[line width=0.4pt]
    \draw (-0.05,0.5) -- (-0.05,1);
    \draw ([shift={(0.5,1)}]180:0.55) arc[start angle=180,end angle=90,radius=0.55];
    \draw (0.5,1.55) -- (3.5,1.55);
    \draw ([shift={(3.5,1)}]90:0.55) arc[start angle=90,end angle=0,radius=0.55];
    \draw (4.05,1) -- (4.05,-2.8);
    \draw ([shift={(4.5,-2.8)}]180:0.45) arc[start angle=180,end angle=360,radius=0.45];
    \draw (4.95,-2.8) -- (4.95,-2.3);
    \draw (5.05,-2.3) -- (5.05,-2.8);
    \draw ([shift={(4.5,-2.8)}]0:0.55) arc[start angle=0,end angle=-180,radius=0.55];
    \draw (3.95,-2.8) -- (3.95,1);
    \draw ([shift={(3.5,1)}]0:0.45) arc[start angle=0,end angle=90,radius=0.45];
    \draw (3.5,1.45) -- (0.5,1.45);
    \draw ([shift={(0.5,1)}]90:0.45) arc[start angle=90,end angle=180,radius=0.45];
    \draw (0.05,1) -- (0.05,0.5);
    \draw (-0.05,0.5) -- (0.05,0.5);
  \end{scope}
  \fill[white] (-0.025,0.52) rectangle (0.018,0.8);
}

\newcommand{\tkzTestTubeDoubleTube}{%
  \begin{scope}[line width=0.4pt]
    \draw (-0.3,-3) -- (-0.3,0);
    \draw (-.2,-3) -- (-0.2,0);
    \draw (0.25,-1) -- (0.25,0);
    \draw (.15,-1) -- (.15,0);
  \end{scope}
  \tkzTestTubeStopper
  \begin{scope}[shift={(0.2,0)}]\tkzCurvedTubeShortDoubleTube\end{scope}
  \begin{scope}[shift={(0,0.5)}]
    \begin{scope}[line width=0.4pt]
      \filldraw[fill=white,rounded corners=0.05cm]
        (-.3,0) -- (-0.3,2) -- (-0.5,2.2) -- (-0.5,3) -- (0,3)
                -- (0,2.2) -- (-0.2,2) -- (-0.2,0) -- cycle;
      \draw (0.1,1.55) -- (-.5,1.55) -- (-0.5,1.7) -- (-.55,1.7)
            -- (-.55,1.3) -- (-0.5,1.3) -- (-0.5,1.55) -- (-.5,1.45) -- (.1,1.45);
      \filldraw[fill=white,rounded corners=0.2] (-0.4,1.3) rectangle (-0.1,1.7);
    \end{scope}
    \fill[white] (-0.25,3) ellipse (0.3 and 0.1);
  \end{scope}
}

% 试管夹 (Pince) - 几何与原 pst-labo pst-laboObj.tex line 483-508 一致
% 根据原 pst-labo 实际渲染 (5.3x 截图) 校准:
%   - 上边缘凹弧较浅 (原 \pscurve 3点几乎共线, 浅弧)
%   - 顶部小凸起为圆弧形 (\psarcn 圆弧)
%   - 所有转角为圆角 (原 \pscustom 默认 linejoin=0 即 round)
\newcommand{\tkzTestTubeClamp}{%
  % 第1块: 上部夹爪
  % 上边缘(1a, 夹合面): 下凹弧贴合试管顶部弧面
  %   试管半径 0.5, 在 y=3.2 处试管顶部弧面下凹约 0.1, 用贝塞尔曲线模拟
  %   从 (0.5,0) 经 (0,-0.1) 到 (-0.5,0), 控制点 y=-0.13 产生平滑下凹
  \filldraw[fill=PaleOrange,draw=black,line cap=round,line join=round]
    (0.5,0) .. controls (0.25,-0.13) and (-0.25,-0.13) .. (-0.5,0)
    -- (-0.8,0)
    -- ([shift={(-0.5,0)}]180:0.3) arc[start angle=180,end angle=270,radius=0.3]
    -- (-.5,-.3) -- (2.5,-.3) -- (2.4,-.2) -- (1,0) -- (0.5,0) -- cycle;
  % 第2块: 下部夹爪
  % 注: 原 \psarcn(-.5,0){0.3}{-90}{180} 是顺时针从 -90°(=270°) 到 180°, 跨度 90° (四分之一圆弧)
  % TikZ arc 当 end<start 时顺时针, 需用 end=-180 让其走短弧 (避免四分之三圆弧)
  \filldraw[fill=PaleOrange,draw=black,line cap=round,line join=round]
    (-0.8,0) -- (-0.8,-0.3)
    -- ([shift={(-0.5,-0.3)}]180:0.3) arc[start angle=180,end angle=270,radius=0.3]
    -- (-0.5,-0.6) -- (2.5,-.6) -- (2.5,-.3) -- (-.5,-.3)
    -- ([shift={(-.5,0)}]-90:0.3) arc[start angle=-90,end angle=-180,radius=0.3] -- cycle;
  % 第3块: 顶部小凸起 - 原 \psarcn 顺时针 180°->90°, 圆弧形凸起
  % 去掉 (-0.8,0.05)->(-0.75,0.05) 的多余水平线段, 直接从 (-0.5,0.05) 到圆弧起点
  % 使凸起呈纯圆弧形 (与原 pst-labo 渲染一致)
  \filldraw[fill=PaleOrange,draw=black,line cap=round,line join=round]
    (-.5,0.05) -- ([shift={(-0.5,0.05)}]180:0.25)
    arc[start angle=180,end angle=90,radius=0.25]
    -- (-.5,0.3) -- (-0.5,0.05) -- cycle;
  % 第4块: 夹臂上部
  \filldraw[fill=PaleOrange,draw=black,line cap=round,line join=round]
    (0.5,0.3) -- (4.9,0.3) -- (5,0.2) -- (1,0) -- (0.5,0) -- (0.5,0.3) -- cycle;
  % 第5块: 夹臂下部 (Brown 棕色)
  \filldraw[fill=Brown,draw=black,line cap=round,line join=round]
    (0.5,0) -- (5,0.2) -- (5,0) -- (2,-0.2) -- (1,0) -- cycle;
  % 螺丝椭圆
  \draw[line width=1.5\pgflinewidth,line cap=round] (1,0) ellipse (0.3 and 0.1);
  % 螺丝下方的杆
  \draw[line width=1.5\pgflinewidth,line cap=round,line join=round]
        (0.7,0) -- (0.4,-.3) -- (0.4,-.6);
}

% ============================================================
% 13. 容量瓶 (Fiole jaugée)
% ============================================================
\newcommand{\tkzVolumetricFlask}{%
  \chemlabUpdateLineWidths
  \draw[rounded corners=0.1cm,line width=\chemlabLineWidth,line cap=round,line join=round]
       (-0.21,3.5) -- (-0.1,3.5) -- (-0.1,1.5)
       -- (-0.75,0) -- (0.75,0) -- (0.1,1.5) -- (0.1,3.5) -- (0.21,3.5);
  \draw[line width=\chemlabGraduationLineWidth,line cap=round] (-0.1,2.5) -- (0.1,2.5);
}

% ============================================================
% 14. 结晶皿 (Cristallisoir)
% ============================================================
\newcommand{\Cristallisoir}{%
  \chemlabUpdateLineWidths
  \begin{scope}[line width=\chemlabLineWidth,line cap=round,line join=round]
  \draw (-1.9,2) -- (-2,2);
  \draw ([shift={(-2,1.95)}]90:0.05) arc[start angle=90,end angle=270,radius=0.05];
  \draw (-2,1.9) -- (-2,1);
  \draw ([shift={(-1,1)}]180:1) arc[start angle=180,end angle=270,radius=1];
  \draw (-1,0) -- (1,0);
  \draw ([shift={(1,1)}]270:1) arc[start angle=270,end angle=360,radius=1];
  \draw (2,1) -- (2,1.9);
  \draw ([shift={(2,1.95)}]270:0.05) arc[start angle=270,end angle=90,radius=0.05];
  \draw (2,2) -- (1.9,2);
  \draw (1.9,2) -- (1.9,1);
  \draw ([shift={(1,1)}]0:0.9) arc[start angle=0,end angle=-90,radius=0.9];
  \draw (1,0.1) -- (-1,0.1);
  \draw ([shift={(-1,1)}]270:0.9) arc[start angle=270,end angle=180,radius=0.9];
  \draw (-1.9,1) -- (-1.9,2);
  \end{scope}
}

\newcommand{\InterieurCristallisoir}{%
  (1.9,2) -- (1.9,1)
  -- ([shift={(1,1)}]0:0.9) arc[start angle=0,end angle=-90,radius=0.9]
  -- (1,0.1) -- (-1,0.1)
  -- ([shift={(-1,1)}]270:0.9) arc[start angle=270,end angle=180,radius=0.9]
  -- (-1.9,1) -- (-1.9,2) -- cycle
}

\newcommand{\tkzCrystallizingDish}{%
  \begin{scope}
    \clip \InterieurCristallisoir;
    \fill[LightBlue1] (-3,-1) rectangle (3,1.5);
  \end{scope}
  \Cristallisoir
}

% ============================================================
% 15. 倒置试管 (TubeRenverse)
% ============================================================
\newcommand{\tkzInvertedTube}{%
  \chemlabUpdateLineWidths
  \begin{scope}[line width=\chemlabEllipseLineWidth,line cap=round,line join=round]
    \fill[LightBlue1] (-0.5,0) rectangle (0.5,5.5);
    % 合并左侧 + 顶部圆弧 + 右侧为单条路径
    % pst-labo: \psarcn(0,6){0.5}{180}{0} = 从180°顺时针到0°, 经过顶部90°
    % TikZ: arc[start=180, end=0] = 从180°逆时针(减小)到0°, 经过顶部90°
    \draw (-0.5,0) -- (-0.5,6)
          arc[start angle=180,end angle=0,radius=0.5]
          -- (0.5,0);
  \end{scope}
}

% ============================================================
% 16. 试管架 (ReposeTube)
% ============================================================
\newcommand{\tkzTestTubeRack}{%
  \begin{scope}[line width=0.1cm,draw=lightgray]
    \draw ([shift={(0,-0.25)}]100:1) arc[start angle=100,end angle=140,radius=1];
    \draw ([shift={(0,-.25)}]150:1) arc[start angle=150,end angle=160,radius=1];
    \draw ([shift={(0,-0.25)}]20:1) arc[start angle=20,end angle=80,radius=1];
  \end{scope}
}

% ============================================================
% 17. 烧杯边缘 + 刻度 (RebordBecher)
% ============================================================
\newcommand{\tkzBeakerRim}{%
  \chemlabUpdateLineWidths
  \draw[line width=\chemlabLineWidth] (0,2.55) ellipse (1.1 and .1);
  % 主刻度 : 粗细为细刻度线的 1.5 倍 (原 pst-labo: 0.03cm = 1.5 × 0.02cm)
  % 用 \pgfmathsetmacro 预算尺寸, 避免 \dimexpr 在 pgfkeys 中的兼容问题
  \pgfmathsetlengthmacro{\chemlab@majorGrad}{1.5*\chemlabGraduationLineWidth}%
  \foreach \nDiv in {0.5,1.0,...,2.0} {%
    \draw[line width=\chemlab@majorGrad] (-.75,\nDiv) -- (-.4,\nDiv);}
  % 细分刻度
  \foreach \nSubDiv in {0.5,0.6,...,2.0} {%
    \draw[line width=\chemlabGraduationLineWidth] (-0.75,\nSubDiv) -- (-.5,\nSubDiv);}
}

% ============================================================
% 18. 烧杯 (Becher)
% ============================================================
% 注 : 原 PSTricks 使用 \psline[linearc=0.5](-1,2.5)(-1,0)(1,0)(1,2.5)
% TikZ 的 rounded corners 与 PSTricks linearc 语义略有不同,
% 因此改用显式圆弧精确模拟 PSTricks linearc=0.5 的行为 :
% 左下角圆弧 : 圆心 (-0.5,0.5), 半径 0.5, 从 180° 到 270°
% 右下角圆弧 : 圆心 (0.5,0.5),  半径 0.5, 从 270° 到 360°
\newcommand{\tkzBeakerBody}{%
  \chemlabUpdateLineWidths
  \draw[line width=\chemlabLineWidth,line cap=round,line join=round]
       (-1,2.5) -- (-1,0.5)
       -- ([shift={(-0.5,0.5)}]180:0.5) arc[start angle=180,end angle=270,radius=0.5]
       -- (0.5,0)
       -- ([shift={(0.5,0.5)}]270:0.5) arc[start angle=270,end angle=360,radius=0.5]
       -- (1,2.5);
}

% ============================================================
% 19. 锥形瓶 (Erlenmeyer)
% ============================================================
% 注 : 原 PSTricks 使用 \psline[linearc=0.3](-0.5,3.5)(-0.5,2.5)(-1.5,0)(1.5,0)(0.5,2.5)(0.5,3.5)
% TikZ 的 rounded corners 与 PSTricks linearc 语义不同,
% 因此改用显式圆弧精确模拟 linearc=0.3 (由 scripts/erlen_arcs5.py 计算) :
%   肩部 (凹转角, 内侧角 21.8°) : 半径 0.3 太大无法相切 (切点距 1.558 > 颈部线段长 1.0), 保留尖角
%   左下 (凸转角, 内侧角 68.2°) : center=(-1.0569,0.3), r=0.3, 158.20° -> 270° (ccw)
%   右下 (凸转角, 内侧角 68.2°) : center=(1.0569,0.3),  r=0.3, -90° -> 21.80° (ccw)
\newcommand{\tkzErlenmeyerBody}{%
  \chemlabUpdateLineWidths
  \draw[line width=\chemlabLineWidth,line cap=round,line join=round]
       (-0.5,3.5) -- (-0.5,2.5)
       -- (-1.3354,0.4114)
       -- ([shift={(-1.0569,0.3)}]158.20:0.3) arc[start angle=158.20,end angle=270,radius=0.3]
       -- (1.0569,0)
       -- ([shift={(1.0569,0.3)}]-90:0.3) arc[start angle=-90,end angle=21.80,radius=0.3]
       -- (0.5,2.5) -- (0.5,3.5);
}

% ============================================================
% 20. 试剂瓶 (Flacon)
% ============================================================
\newcommand{\tkzReagentBottleBody}{%
  \chemlabUpdateLineWidths\begin{scope}[line width=\chemlabLineWidth,line cap=round,line join=round]
    % 左侧瓶口上角小弧 (psarcn 90->0, 顺时针)
    \draw ([shift={(-0.6,3.4)}]90:0.1) arc[start angle=90,end angle=0,radius=0.1];
    % 左侧颈部直线
    \draw (-0.5,3.4) -- (-0.5,3.1);
    % 左侧瓶口下角凹弧 (psarcn 0->-90, 顺时针; TikZ end<start 自动顺时针)
    \draw ([shift={(-0.6,3.1)}]0:0.1) arc[start angle=0,end angle=-90,radius=0.1];
    % 左侧肩部大凸弧 (psarc 90->180, 逆时针)
    \draw ([shift={(-0.6,2.6)}]90:0.4) arc[start angle=90,end angle=180,radius=0.4];
    % 左侧瓶身直线
    \draw (-1,2.6) -- (-1,0.2);
    % 左侧底部圆角 (psarc 180->270, 逆时针)
    \draw ([shift={(-.8,0.2)}]180:0.2) arc[start angle=180,end angle=270,radius=0.2];
    % 底部
    \draw (-.8,0) -- (.8,0);
    % 右侧底部圆角 (psarc -90->0, 逆时针)
    \draw ([shift={(.8,0.2)}]-90:0.2) arc[start angle=-90,end angle=0,radius=0.2];
    % 右侧瓶身直线
    \draw (1,0.2) -- (1,2.6);
    % 右侧肩部大凸弧 (psarc 0->90, 逆时针)
    \draw ([shift={(0.6,2.6)}]0:0.4) arc[start angle=0,end angle=90,radius=0.4];
    % 右侧瓶口下角凹弧 (psarcn -90->180, 顺时针短弧; TikZ 需 end=-180 避免 270° 长弧)
    \draw ([shift={(0.6,3.1)}]-90:0.1) arc[start angle=-90,end angle=-180,radius=0.1];
    % 右侧颈部直线
    \draw (0.5,3.1) -- (0.5,3.4);
    % 右侧瓶口上角小弧 (psarcn 180->90, 顺时针)
    \draw ([shift={(0.6,3.4)}]180:0.1) arc[start angle=180,end angle=90,radius=0.1];
  \end{scope}
}

% ============================================================
% 21. 玻璃杯 (Verre) - F. Bréal
% ============================================================
\newcommand{\tkzWatchGlassBody}{%
  \chemlabUpdateLineWidths\begin{scope}[line width=\chemlabLineWidth,line cap=round,line join=round]
    % 合并左斜线 + 底部圆弧 + 右斜线为单条路径, 平滑转角
    \draw (-1,3.5) -- (-0.5,1)
          arc[start angle=180,end angle=360,radius=0.5]
          -- (1,3.5);
  \end{scope}
}

\newcommand{\tkzWatchGlassBase}{%
  \chemlabUpdateLineWidths\begin{scope}[line width=\chemlabLineWidth,rounded corners=0.03cm,line cap=round,line join=round]
    \draw (-0.5,1) -- (-0.34,0.2) -- (-1.2,0) -- (1.2,0) -- (0.34,0.2) -- (0.5,1);
  \end{scope}
}

% ============================================================
% 22. 漏斗 (Entonnoir)
% ============================================================
\newcommand{\tkzFunnelOutline}{%
  \begin{scope}[line cap=round,line join=round]
    \draw (-0.07,0) -- (-0.07,1) -- (-1,2.5);
    \draw (1,2.5) -- (0.07,1) -- (0.07,0.1);
  \end{scope}
}

\newcommand{\tkzFunnelBody}{%
  \begin{scope}
    \clip (-0.07,0) -- (-0.07,1) -- (-1,2.5)
          -- (1,2.5) -- (0.07,1) -- (0.07,0.1) -- cycle;
    \fill[white] (-2,-1) rectangle (2,3);
  \end{scope}
  \begin{scope}[rounded corners=0.25cm]
    \draw (-0.07,0) -- (-0.07,1) -- (-1,2.5);
    \draw (1,2.5) -- (0.07,1) -- (0.07,0.1);
  \end{scope}
  \draw (0,2.5) ellipse (1 and 0.1);
}

% ============================================================
% 23. 球形冷凝器 (Réfrigérant à boulles)
% ============================================================
\newcommand{\tkzCondenserStem}{%
  % 整个 TigeBoulles 作为一个连续路径 (匹配 pst-labo \pscustom)
  % pst-labo: \pscustom[linestyle=none]{...\fill} 填充 + 单独描边 (无闭合线)
  % TikZ: 分开 \fill (带 cycle 闭合填充) 和 \draw (无 cycle, 顶部敞口)
  % 顶部是敞口的 — 两端弧之间没有连线 (匹配 pst-labo)
  %
  % 填充: 带闭合线 (fill 需要 cycle)
  \fill[white]
    % 顶部圆弧与颈部 - 左侧
    ({-0.35+0.1*cos(90)},{7.5+0.1*sin(90)}) arc[start angle=90,end angle=0,radius=0.1]
    -- (-.25,7.5) -- (-.25,6.6)
    -- ({-.15+0.1*cos(180)},{6.6+0.1*sin(180)}) arc[start angle=180,end angle=270,radius=0.1]
    -- ({-.15+0.1*cos(90)},{6.4+0.1*sin(90)}) arc[start angle=90,end angle=0,radius=0.1]
    -- (-.05,6.4) -- (-0.05,6)
    % 左侧 7 个球状突起 (中心 x=-0.05, 半径 0.25, arc 90→270 向左凸)
    (-0.05,6.0) arc[start angle=90,end angle=270,radius=0.25] -- (-0.05,5.50) -- (-0.05,5.25)
    arc[start angle=90,end angle=270,radius=0.25] -- (-0.05,4.75) -- (-0.05,4.50)
    arc[start angle=90,end angle=270,radius=0.25] -- (-0.05,4.00) -- (-0.05,3.75)
    arc[start angle=90,end angle=270,radius=0.25] -- (-0.05,3.25) -- (-0.05,3.00)
    arc[start angle=90,end angle=270,radius=0.25] -- (-0.05,2.50) -- (-0.05,2.25)
    arc[start angle=90,end angle=270,radius=0.25] -- (-0.05,1.75) -- (-0.05,1.50)
    arc[start angle=90,end angle=270,radius=0.25] -- (-0.05,1.00) -- (-0.05,0.75)
    % 左侧底部到下管口
    -- (-0.05,-1.5) -- (0.05,-1.5) -- (0.05,1)
    % 右侧 7 个球状突起 (中心 x=0.05, 半径 0.25, arc 270→450 向右凸)
    arc[start angle=270,end angle=450,radius=0.25] -- (0.05,1.50) -- (0.05,1.75)
    arc[start angle=270,end angle=450,radius=0.25] -- (0.05,2.25) -- (0.05,2.50)
    arc[start angle=270,end angle=450,radius=0.25] -- (0.05,3.00) -- (0.05,3.25)
    arc[start angle=270,end angle=450,radius=0.25] -- (0.05,3.75) -- (0.05,4.00)
    arc[start angle=270,end angle=450,radius=0.25] -- (0.05,4.50) -- (0.05,4.75)
    arc[start angle=270,end angle=450,radius=0.25] -- (0.05,5.25) -- (0.05,5.50)
    arc[start angle=270,end angle=450,radius=0.25] -- (0.05,6.00) -- (0.05,6.40)
    % 顶部圆弧与颈部 - 右侧
    -- ({0.15+0.1*cos(180)},{6.4+0.1*sin(180)}) arc[start angle=180,end angle=90,radius=0.1]
    -- ({0.15+0.1*cos(-90)},{6.6+0.1*sin(-90)}) arc[start angle=-90,end angle=0,radius=0.1]
    -- (0.25,6.6) -- (0.25,7.5)
    -- ({0.35+0.1*cos(180)},{7.5+0.1*sin(180)}) arc[start angle=180,end angle=90,radius=0.1]
    -- cycle;
  %
  % 描边: 底部敞口 (匹配 pst-labo: 左右侧分别 \psline, 无底部连线)
  \draw
    % 顶部圆弧与颈部 - 左侧
    ({-0.35+0.1*cos(90)},{7.5+0.1*sin(90)}) arc[start angle=90,end angle=0,radius=0.1]
    -- (-.25,7.5) -- (-.25,6.6)
    -- ({-.15+0.1*cos(180)},{6.6+0.1*sin(180)}) arc[start angle=180,end angle=270,radius=0.1]
    -- ({-.15+0.1*cos(90)},{6.4+0.1*sin(90)}) arc[start angle=90,end angle=0,radius=0.1]
    -- (-.05,6.4) -- (-0.05,6)
    % 左侧 7 个球状突起
    (-0.05,6.0) arc[start angle=90,end angle=270,radius=0.25] -- (-0.05,5.50) -- (-0.05,5.25)
    arc[start angle=90,end angle=270,radius=0.25] -- (-0.05,4.75) -- (-0.05,4.50)
    arc[start angle=90,end angle=270,radius=0.25] -- (-0.05,4.00) -- (-0.05,3.75)
    arc[start angle=90,end angle=270,radius=0.25] -- (-0.05,3.25) -- (-0.05,3.00)
    arc[start angle=90,end angle=270,radius=0.25] -- (-0.05,2.50) -- (-0.05,2.25)
    arc[start angle=90,end angle=270,radius=0.25] -- (-0.05,1.75) -- (-0.05,1.50)
    arc[start angle=90,end angle=270,radius=0.25] -- (-0.05,1.00) -- (-0.05,0.75)
    -- (-0.05,-1.5);
  % 右侧 (单独 \draw, 底部不连接左侧)
  \draw
    (0.05,-1.5) -- (0.05,1)
    % 右侧 7 个球状突起
    arc[start angle=270,end angle=450,radius=0.25] -- (0.05,1.50) -- (0.05,1.75)
    arc[start angle=270,end angle=450,radius=0.25] -- (0.05,2.25) -- (0.05,2.50)
    arc[start angle=270,end angle=450,radius=0.25] -- (0.05,3.00) -- (0.05,3.25)
    arc[start angle=270,end angle=450,radius=0.25] -- (0.05,3.75) -- (0.05,4.00)
    arc[start angle=270,end angle=450,radius=0.25] -- (0.05,4.50) -- (0.05,4.75)
    arc[start angle=270,end angle=450,radius=0.25] -- (0.05,5.25) -- (0.05,5.50)
    arc[start angle=270,end angle=450,radius=0.25] -- (0.05,6.00) -- (0.05,6.40)
    -- ({0.15+0.1*cos(180)},{6.4+0.1*sin(180)}) arc[start angle=180,end angle=90,radius=0.1]
    -- ({0.15+0.1*cos(-90)},{6.6+0.1*sin(-90)}) arc[start angle=-90,end angle=0,radius=0.1]
    -- (0.25,6.6) -- (0.25,7.5)
    -- ({0.35+0.1*cos(180)},{7.5+0.1*sin(180)}) arc[start angle=180,end angle=90,radius=0.1];
  % 底部液滴
  \begin{scope}[shift={(0,-1.6)}]\tkzDrop\end{scope}
}

\newcommand{\tkzCondenserBody}{%
  % 渐变填充: pst-labo gradbegin=cyan (底), gradend=PaleOrange (顶), gradmidpoint=0
  % TikZ: bottom color=cyan (底), top color=PaleOrange (顶)
  \shade[top color=PaleOrange,bottom color=cyan]
    ([shift={(0,5.75)}]0:0.5) arc[start angle=0,end angle=180,radius=0.5]
    -- (-0.5,5.75) -- (-0.5,5.5) -- (-1.5,5.5)
    -- (-1.5,5.3) -- (-.5,5.3) -- (-.5,1.7) -- (-1.5,1.7)
    -- (-1.5,1.5) -- (-.5,1.5) -- (-0.5,1)
    -- ([shift={(0,1)}]180:0.5) arc[start angle=180,end angle=360,radius=0.5]
    -- (0.5,1) -- (0.5,5.75) -- cycle;
  % 描边
  \begin{scope}[line width=0.88pt]
    \draw ([shift={(0,5.75)}]0:0.5) arc[start angle=0,end angle=180,radius=0.5];
    \draw (-0.5,5.75) -- (-0.5,5.5) -- (-1.5,5.5);
    \draw (-0.5,5.75) -- (-0.5,5.5) -- (-1.5,5.5);
    \draw (-1.5,5.3) -- (-.5,5.3) -- (-.5,1.7) -- (-1.5,1.7);
    \draw (-1.5,1.5) -- (-.5,1.5) -- (-0.5,1);
    \draw ([shift={(0,1)}]180:0.5) arc[start angle=180,end angle=360,radius=0.5];
    \draw (0.5,1) -- (0.5,5.75);
  \end{scope}
  \draw[->,cyan] (-2.5,1.6) -- (-1.6,1.6);
  \draw[<-,PaleOrange] (-2.5,5.4) -- (-1.6,5.4);
}

\newcommand{\tkzCondenserBulbs}{%
  \begin{scope}[shift={(0,0.2)}]
    \tkzCondenserBody
    \tkzCondenserStem
  \end{scope}
  \tkzTestTubeStopper
  \fill[white,rounded corners=0.025cm]
       (-0.05,0.6) -- (-0.05,0.5) -- (0.05,0.5) -- (0.05,0.6) -- cycle;
}

% ============================================================
% 24. 加热球 (Chauffe-Ballon)
% ============================================================
\newcommand{\tkzHeaterRoundFlask}{%
  \filldraw[fill=white,draw=black] (-2,-1) rectangle (2,1);
  \draw (-1.25,0) circle (0.4);
  \begin{scope}[shift={(-1.25,0)}]
    \foreach \iA in {0,30,...,180} {\draw (\iA:0.45) -- (\iA:0.6);}
    \draw[->] (0,0) -- (60:0.3);
  \end{scope}
  % 按钮区域: pst-labo \psframe[doubleline=true] 产生双线边框
  % PSTricks doubleline: 外框在指定坐标, 内框向内偏移
  % 实测: 外框到内框偏移约 0.13 units (含线宽+间隔, at scale=0.5)
  % 外框 (在指定坐标处)
  \draw[black,line width=0.8pt] (1,-.5) rectangle (1.5,0.5);
  % 内框 (向内偏移, 填充灰色)
  \filldraw[fill=gray,draw=black,line width=0.8pt] ({1+0.14},{-.5+0.14}) rectangle ({1.5-0.14},{0.5-0.14});
  % 红色指示灯
  \fill[red] (1.25,0.7) circle (0.1);
}

% ============================================================
% 25. 回流烧瓶 (Ballon à reflux)
% ============================================================
\newcommand{\tkzRoundFlaskReflux}{%
  \begin{scope}
    \clip \tkzRoundFlaskSimplePath;   % 使用路径版本形成闭合 clip 区域
    \shade[top color=cyan,bottom color=white,draw=cyan]
      plot[smooth,tension=0.5] coordinates {%
        (-1.5,1.5) (-1.375,1.55) (-1.25,1.5) (-1.125,1.45)
        (-1,1.5) (-.875,1.55) (-.75,1.5) (-.625,1.45)
        (-.5,1.5) (-.375,1.55) (-.25,1.5) (-0.125,1.45)
        (0,1.5) (0.125,1.55) (.25,1.5) (0.375,1.45)
        (0.5,1.5) (0.625,1.55) (0.75,1.5) (0.875,1.45)
        (1,1.5) (1.125,1.55) (1.25,1.5) (1.375,1.45) (1.5,1.5)}
      -- (1.5,1.5) -- (1.5,0) -- (-1.5,0) -- (-1.5,1.5) -- cycle;
  \end{scope}
  \tkzRoundFlaskSimple
}

% ============================================================
% 26. 分液漏斗 (Ampoule à décanter / Separate Funnel)
% ============================================================
% 相 : 水相 / 有机相
\definecolor{AqueousPhaseColor}{cmyk}{0.215686,0.062745,0,0}
\definecolor{OrganicPhaseColor}{cmyk}{0.223529,0.094118,0.501961,0}

% pst-labo: \tkzSeparatoryFunnel 设 linewidth=0.04 (用户单位, 随 unit 缩放)
% TikZ: 用 \chemlabLineWidth (基础 0.053cm × scale) 近似, 匹配其他器皿线宽
\newcommand{\tkzSeparatoryFunnelBody}{%
  \chemlabUpdateLineWidths
  \draw[line width=\chemlabLineWidth,line cap=round,line join=round]
    (-0.65,15) -- (-0.55,13.75)
    .. controls (-1.55,13.54) and ( -1.92,12.7) .. (-1.92,11.65)
    .. controls (-1.92,10.6) and (-0.3,9.55) .. (-0.3,8.5)
    -- (-0.3,2)
    .. controls (-0.3,1.5) and (-0.14,1.5) .. (-0.14,1)
    -- (-0.14,-1)
    -- (0.14,-1)
    -- (0.14,1)
    .. controls (0.14,1.5) and (0.3,1.5) .. (0.3,2)
    -- (0.3,8.5)
    .. controls (0.3,9.55) and (1.92,10.6) .. (1.92,11.65)
    .. controls ( 1.92,12.7) and (1.55,13.54) .. (0.55,13.75)
    -- (0.65,15);
}

% 用于 clip 的轮廓
\newcommand{\tkzSeparatoryFunnelOutlineClip}{%
  (-0.65,15) -- (-0.55,13.75)
  .. controls (-1.55,13.54) and ( -1.92,12.7) .. (-1.92,11.65)
  .. controls (-1.92,10.6) and (-0.3,9.55) .. (-0.3,8.5)
  -- (-0.3,2)
  .. controls (-0.3,1.5) and (-0.14,1.5) .. (-0.14,1)
  -- (-0.14,-1) -- (0.14,-1) -- (0.14,1)
  .. controls (0.14,1.5) and (0.3,1.5) .. (0.3,2)
  -- (0.3,8.5)
  .. controls (0.3,9.55) and (1.92,10.6) .. (1.92,11.65)
  .. controls (1.92,12.7) and (1.55,13.54) .. (0.55,13.75)
  -- (0.65,15) -- cycle
}

% 顶盖 (Bouchon) - 继承 linewidth=0.04
\newcommand{\tkzStopper}{%
  \chemlabUpdateLineWidths
  \filldraw[fill=gray!10,draw=gray,line width=\chemlabLineWidth,line cap=round,line join=round]
    (-0.67,15.25) -- (0.67,15.25) -- (0.57,14) -- (-0.57,14) -- cycle;
  \filldraw[fill=gray!5,draw=gray,line width=\chemlabLineWidth,line cap=round,line join=round]
    ([shift={(0.555,15.365)}]-90:0.1145) arc[start angle=-90,end angle=-180,radius=0.1145]
    -- ([shift={(0,15.92)}]-45:0.67) arc[start angle=-45,end angle=225,radius=0.67]
    -- ([shift={(-0.555,15.365)}]0:0.1145) arc[start angle=0,end angle=-90,radius=0.1145]
    -- cycle;
}

% 开启的旋塞 (TapOpen)
% pst-labo: \psframe[fillstyle=solid] 默认 fillcolor=white (PSTricks 默认)
% TikZ: \fill 默认黑色, 必须显式 fill=white
\newcommand{\tkzTapOpen}{%
  \chemlabUpdateLineWidths
  \filldraw[fill=white,line width=\chemlabLineWidth] (-.2,0.5) rectangle (.2,1);
  \draw[line width=\chemlabLineWidth] (-.2,0.4) rectangle (-.3,1.1);
  \draw[line width=\chemlabLineWidth] (.2,0.4) rectangle (.3,1.1);
  \draw[line width=\chemlabLineWidth] (.3,0.6) -- (.4,0.6) -- (.5,0.3) -- (.9,0.3) -- (.9,1.2) -- (.5,1.2)
        -- (.4,0.9) -- (.3,0.9) -- cycle;
  \fill[cyan!50] (-.6,0.4) -- (-.85,0.6) -- (-.85,0.9) -- (-.6,1.1) -- (-.3,1.1) -- (-.3,0.4) -- cycle;
  \draw[line width=\chemlabLineWidth] (-0.6,0.4) -- (-0.6,1.1);
  \draw[line width=\chemlabLineWidth] (0.5,0.3) -- (0.5,1.2);
}

% 关闭的旋塞 (TapClosed)
% pst-labo: \psframe[fillstyle=solid] 默认 fillcolor=white
\newcommand{\tkzTapClosed}{%
  \chemlabUpdateLineWidths
  \filldraw[fill=white,line width=\chemlabLineWidth] (-.2,0.5) rectangle (.2,1);
  \draw[line width=\chemlabLineWidth] (-.2,0.4) rectangle (-.3,1.1);
  \draw[line width=\chemlabLineWidth] (.2,0.4) rectangle (.3,1.1);
  \draw[line width=\chemlabLineWidth] ([shift={(0.35,0.6)}]180:0.05) arc[start angle=180,end angle=90,radius=0.05]
        -- (0.4,0.65) -- (0.9,0.65) -- (0.9,0.85) -- (0.35,0.85)
        -- ([shift={(0.35,0.9)}]-90:0.05) arc[start angle=-90,end angle=-180,radius=0.05]
        -- cycle;
  \draw[line width=\chemlabLineWidth] (0.5,0.65) -- (0.5,0.85);
  \fill[cyan!50] (-.6,0.4) -- (-.85,0.6) -- (-.85,0.9) -- (-.6,1.1) -- (-.3,1.1) -- (-.3,0.4) -- cycle;
  \draw[line width=\chemlabLineWidth] (-0.6,0.4) -- (-0.6,1.1);
}

% \tkzSeparatoryFunnel  (主调用)
% 可选参数 : stopper=true/false, TapOpen=true/false, liquidLevel1, liquidLevel2
\newcommand{\tkzSeparatoryFunnel}[1][]{%
  \begingroup
  \chemlabUpdateLineWidths
  \pgfkeys{
    /tkzSeparatoryFunnel/.cd,
    stopper/.is if=TkzTestTubeStopper,
    TapOpen/.is if=TkzTapOpen,
    liquidLevel1/.store in=\tkzTestTubeLiquidLevelI,
    liquidLevel2/.store in=\tkzTestTubeLiquidLevelII,
    #1}
  % 默认值
  \ifdefined\tkzTestTubeLiquidLevelI\else\def\tkzTestTubeLiquidLevelI{11}\fi
  \ifdefined\tkzTestTubeLiquidLevelII\else\def\tkzTestTubeLiquidLevelII{12}\fi
  % 水相
  \begin{scope}
    \clip \tkzSeparatoryFunnelOutlineClip;
    \fill[AqueousPhaseColor] (-2,1) rectangle (2,\tkzTestTubeLiquidLevelI);
  \end{scope}
  % 有机相
  \begin{scope}
    \clip \tkzSeparatoryFunnelOutlineClip;
    \fill[OrganicPhaseColor] (-2,\tkzTestTubeLiquidLevelI)
          rectangle (2,\tkzTestTubeLiquidLevelII);
  \end{scope}
  \ifTkzTestTubeStopper\tkzStopper\fi
  \tkzSeparatoryFunnelBody
  \ifTkzTapOpen\tkzTapOpen\else\tkzTapClosed\fi
  % 出口管口 (pst-labo: \psellipse[fillstyle=solid,fillcolor=white] 既有白色填充又有黑色轮廓)
  \filldraw[fill=white,line width=\chemlabLineWidth] (0,-1) ellipse (0.15 and 0.25);
  % 瓶口圆点 (填充, 无需线宽)
  \fill[darkgray] (-0.675,15) circle (0.03);
  \fill[darkgray] (0.675,15) circle (0.03);
  \draw[darkgray,line width=\chemlabLineWidth] (-0.65,15) -- (0.65,15);
  % 颈与瓶身接合
  \draw[gray,line width=\chemlabLineWidth] (-0.55,13.75) -- (0.55,13.75);
  \endgroup
}

% ============================================================
% 27. 支架 (pstSupport)
% ============================================================
% pst-labo: \psset{dimen=m,linewidth=0.04} — linewidth=0.04 是用户单位, 随 unit 缩放
% TikZ: 用 \chemlabLineWidth (基础 0.053cm × scale), 与其他器皿一致
\newcommand{\tkzStand}{%
  \chemlabUpdateLineWidths
  \begin{scope}[line width=\chemlabLineWidth]
    \filldraw[fill=woodColor,draw=black,line join=round]
        (-2,-6) rectangle (5,-5);
    \filldraw[fill=woodColor,draw=black,line join=round]
        (2.75,2) rectangle (3.25,16);
    \filldraw[fill=woodColor,draw=black,line join=round]
        ([shift={(2.75,1.5)}]90:0.5) arc[start angle=90,end angle=180,radius=0.5]
        -- (2.25,1.5) -- (2.25,-5) -- (3.75,-5) -- (3.75,1.5)
        -- ([shift={(3.25,1.5)}]0:0.5) arc[start angle=0,end angle=90,radius=0.5]
        -- cycle;
    \draw (2,11.5) -- (1.55,10.5);
    \filldraw[fill=woodColor,draw=black,line join=round]
        (-2.5,11.5) -- (-2,11.5) -- (-1.55,10.5) -- (-2.5,10.5) -- cycle;
    \filldraw[fill=woodColor,draw=black,line join=round]
        (4,11.5) -- (2,11.5) -- (1.55,10.5) -- (4,10.5) -- cycle;
    \draw (-2.5,10.5) rectangle (4,11.5);
    % 旋钮: pst-labo \pscustom[style=woodColor]{...} 填充会自动闭合, 但描边只画实际路径
    % (不画对角闭合边). TikZ: \fill 自动闭合填充, \draw 画开放路径 (无 -- cycle).
    \fill[woodColor]
        (4,11.2) -- (4.5,11.2)
        -- ([shift={(5,11)}]158.2:0.5385) arc[start angle=158.2,end angle=-158.2,radius=0.5385]
        -- (4.5,10.8) -- (4,10.8) -- cycle;
    \draw
        (4,11.2) -- (4.5,11.2)
        -- ([shift={(5,11)}]158.2:0.5385) arc[start angle=158.2,end angle=-158.2,radius=0.5385]
        -- (4.5,10.8) -- (4,10.8);
    \draw (4.5,10.8) -- (4.5,11.2);
  \end{scope}
}

\endinput
    % 形状宏库 (必须先加载)
%    % ============================================================
%  tikz-chemlab.tex
%  化学实验器材 TikZ 绘图包 - 主实现 (参数系统 + 高级命令)
%  原作者 : D. Girou / C. Jorsen / M. Luque / H. Voss
%  TikZ 改写 : Aihua Wang(王爱华) (原 PSTricks 包 pst-labo)
% ============================================================
%
%  依赖 :
%    \usepackage{tikz}
%    \usetikzlibrary{arrows.meta,calc,patterns,patterns.meta,shadings}
%    % ============================================================
%  tikz-chemlabObj.tex
%  化学实验器材 TikZ 绘图包 - 形状宏库
%  原作者 : D. Girou / C. Jorsen / M. Luque / H. Voss
%  TikZ 改写 : Aihua Wang(王爱华) (原 PSTricks 包 pst-laboObj)
% ============================================================
%
%  使用说明 :
%    \usepackage{tikz}
%    \usetikzlibrary{arrows.meta,calc,patterns,patterns.meta,shadings}
%    \input{tikz-chemlabObj}
%    \begin{tikzpicture}
%       \BuretteCorps  ... 等等
%    \end{tikzpicture}
%
%  映射规则 :
%    \psline(A)(B)(C)              -> \draw (A) -- (B) -- (C);
%    \psline{->}(A)(B)             -> \draw[->] (A) -- (B);
%    \pspolygon(A)(B)(C)           -> \filldraw (A) -- (B) -- (C) -- cycle;
%    \psframe(x1,y1)(x2,y2)        -> \draw (x1,y1) rectangle (x2,y2);
%    \psframe*[linecolor=c]        -> \fill[c] (..) rectangle (..);
%    \psframe[framearc=r]          -> \draw[rounded corners=r] (...) rectangle (...);
%    \pscircle(c){r}               -> \draw (c) circle (r);
%    \pscircle*[linecolor=c](c){r} -> \fill[c] (c) circle (r);
%    \psellipse(c)(a,b)            -> \draw (c) ellipse (a and b);
%    \psarc(c){r}{a1}{a2}          -> \draw ([shift={(c)}]a1:r) arc[start angle=a1,end angle=a2,radius=r];
%    \psarcn(c){r}{a1}{a2}         -> \draw ([shift={(c)}]a1:r) arc[start angle=a1,end angle=a2,radius=r];   (注意 : PSTricks 中 arcn 是逆时针 / 角度递减)
%    \psbezier(A)(B)(C)(D)         -> \draw (A) .. controls (B) and (C) .. (D);
%    \pscurve(P1)(P2)...(Pn)       -> \draw plot[smooth,tension=0.5] coordinates {(P1) (P2) ... (Pn)};
%    \pscustom{ ... \fill}         -> \filldraw / \path[fill] (整段路径) -- cycle;
%    \rput(x,y){...}               -> \begin{scope}[shift={(x,y)}] ... \end{scope}
%    \rput{angle}(x,y){...}        -> \begin{scope}[shift={(x,y)},rotate=angle] ... \end{scope}
%    \multido{\i=a+b}{n}{...}      -> \foreach \i in {a,a+b,...,a+(n-1)*b} {...}
%    \psclip{<path>} ... \endpsclip-> \begin{scope} \clip <path>; ... \end{scope}
%    fillstyle=solid,fillcolor=c   -> fill=c
%    fillstyle=gradient,gradbegin=..,gradend=.. -> \shade[left color=..,right color=..]
%    fillstyle=hlines              -> pattern=horizontal lines
%    doubleline=true               -> double
%    linestyle=dashed              -> dashed
%    framearc=r                    -> rounded corners=r (近似)
%    linearc=r                     -> rounded corners=r
% ============================================================

\NeedsTeXFormat{LaTeX2e}
\ProvidesFile{tikz-chemlabObj.tex}[2026/08/03 v1.0.0]

\RequirePackage{tikz}
\usetikzlibrary{arrows.meta,calc,patterns,patterns.meta,shadings,decorations.pathmorphing}

% ------------------------------------------------------------
% 颜色定义 (对应原包中使用的颜色)
% ------------------------------------------------------------
\definecolor{LightGray}{gray}{0.80}      % 浅灰
\definecolor{VeryLightGray}{cmyk}{0,0,0,0.1}  % 极浅灰 (label 标签背景)
\definecolor{LightBlue}{rgb}{0.68,0.85,0.90}
\definecolor{LightBlue1}{rgb}{0.60,0.80,1.00}
\definecolor{VeryPaleOrange}{cmyk}{0,0.1,0.3,0}  % 滴定管液体默认色
\definecolor{PaleOrange}{rgb}{1.00,0.78,0.55}
\definecolor{Brown}{rgb}{0.50,0.25,0.00}
\definecolor{woodColor}{rgb}{0.957,0.84,0.69}
\definecolor{darkWoodColor}{rgb}{0.50,0.25,0.00}
\definecolor{flameAndGrill}{rgb}{0.95,0.55,0.10}  % 估算颜色
\definecolor{Beige}{rgb}{0.96,0.96,0.86}  % Champagne 液体默认色

% ------------------------------------------------------------
% 公共缺省值 (供调用方按需修改)
% ------------------------------------------------------------
\tikzset{
  Distillation/.cd,
  DistillateColor/.store in=\DistillationDistillateColor,
  DistillateColor=LightBlue1,
  MixtureStyle/.style={fill=LightBlue},
}
\newcommand{\tkzDistillationDistillateColor}{LightBlue1}
\newcommand{\tkzDistillationMixtureStyle}{LightBlue}
\newcommand{\tkzTestTubeLiquidLevelI}{11}     % liquidLevel1
\newcommand{\tkzTestTubeLiquidLevelII}{12}    % liquidLevel2
\newif\ifTkzTapOpen
\newif\ifTkzTestTubeStopper

% ============================================================
% 1. 滴定管 (Burette)
% ============================================================
\newcommand{\tkzBuretteBody}{%
  \draw (-0.5,11.5) -- (-0.25,11) -- (-0.25,3) -- (0,2.75) -- (0,0.5);
  \draw (0.2,0.7) -- (0.2,2.75) -- (0.5,3) -- (0.5,11) -- (0.75,11.5);
}

\newcommand{\tkzBuretteStopcock}{%
  % pst-labo: \psframe[fillstyle=solid] → 填充默认白色 + 黑色边框
  \filldraw[fill=white] (-.1,1.6) rectangle (.3,2.1);
  \draw (-.1,1.5) rectangle (-.2,2.2);
  \draw (.3,1.5) rectangle (.4,2.2);
  \draw (.4,1.7) -- (.5,1.7) -- (.6,1.4) -- (.8,1.4)
        -- (.8,2.3) -- (.6,2.3) -- (.5,2) -- (.4,2) -- cycle;
  % 左侧旋钮: pst-labo \pspolygon + \psframe with fillstyle=solid,fillcolor=LightBlue
  % 需要 LightBlue 填充 + 黑色外轮廓线
  \begin{scope}[fill=LightBlue,draw=black]
    \filldraw (-.5,1.5) -- (-.75,1.7) -- (-.75,2) -- (-.5,2.2)
          -- (-.3,2.2) -- (-.3,1.5) -- cycle;
    \filldraw (-.5,1.5) rectangle (-.2,2.2);
  \end{scope}
}

\newcommand{\tkzBuretteGraduation}{%
  % 滴定管刻度 (原 pst-labo 无显式线宽, 这里用细刻度线宽)
  \foreach \n in {3.0,3.32,...,10.68} { \draw[line width=\chemlabGraduationLineWidth] (-0.25,\n) -- (0,\n); }
}

% ============================================================
% 2. pH 计 (pH-mètre)
% ============================================================
\newcommand{\tkzTitrationPHMeter}{%
  % boitier
  \filldraw[fill=white,draw=black,rounded corners=0.125cm] (-1,0) rectangle (1,3);
  \fill[LightGray] (-0.75,2) rectangle (0.75,2.75);
  \node[white] at (0,2.375) {PH};
  \draw (0,1.25) circle (0.3);
  \draw[->] (0,1.25) -- (0,1.5);
  \node at (0.65,1.5) {$^\circ$C};
  \draw (-0.5,0.5) circle (0.2);
  \draw (0.5,0.5) circle (0.2);
  % électrode — 完全透明填充 (不遮挡烧杯刻度线)
  % pst-labo: \psframe[framearc=0.5](2.25,2.5)(2.75,4.5) 外框圆角
  %   顶部圆角被下方 \psframe[LightGray](sharp) 覆盖, 底部圆角被 \pscircle 覆盖
  %   实际可见的外框是直角矩形, 此处用 sharp corners 匹配 pst-labo 可见效果
  \draw (2.25,2.5) rectangle (2.75,4.5);
  % 顶部 LightGray 填充 + 描边 (sharp corners)
  \filldraw[fill=LightGray,draw=black] (2.25,4) rectangle (2.75,4.5);
  \draw (2.5,2.25) circle (0.25);
  % fil
  \draw[line width=0.05cm,LightBlue,rounded corners=0.2cm]
        (2.5,4.5) -- (2.5,5) -- (1.25,4.75) -- (0,5) -- (0,3);
}

% ============================================================
% 3. 滴定支架 (Support dosage)
% ============================================================
\newcommand{\tkzTitrationStand}{%
  % pst-labo: \psframe(-1.5,0)(1.5,1.5) + {\psset{doubleline=true} ...}
  % 线宽: pst-labo 在 \tkzTitration 中设 linewidth=0.053 (×0.4cm = 0.0212cm ≈ 0.6pt)
  \chemlabUpdateLineWidths
  % 平台: 无填充, 只有边框
  \draw[line width=\chemlabLineWidth] (-1.5,0) rectangle (1.5,1.5);
  % 红色按钮: \pscircle*[linecolor=red] → 红色填充, 双线边框也是红色 (不可见)
  \fill[red] (0,0.75) circle (0.2);
  % 回字型矩形按钮: 无填充, 双线边框
  % PSTricks doubleline: 指定坐标在内外线之间
  % 实测: 外线向外偏移 0.085 units, 内线向内偏移 0.034 units
  \draw[line width=\chemlabLineWidth] ({-1.25-0.085},{0.3-0.085}) rectangle ({-0.75+0.085},{1.2+0.085});
  \draw[line width=\chemlabLineWidth] ({-1.25+0.034},{0.3+0.034}) rectangle ({-0.75-0.034},{1.2-0.034});
  % 环形旋钮: 无填充, 双线边框
  % 外线向外偏移, 内线向内偏移
  \draw[line width=\chemlabLineWidth] (1,0.75) circle ({0.25+0.085});
  \draw[line width=\chemlabLineWidth] (1,0.75) circle ({0.25-0.034});
}

% ============================================================
% 4. 磁力搅拌子 (Aimant)
% ============================================================
\newcommand{\tkzTitrationStirBar}{%
  % pst-labo: \psset{linewidth=0.053} 在外层 \tkzTitration 中设置, 随 scale 缩放
  \chemlabUpdateLineWidths
  \draw[line width=\chemlabLineWidth,rounded corners=0.075cm,fill=white] (-0.5,0.1) rectangle (0.5,0.4);
  \draw[line width=\chemlabLineWidth,rounded corners=0.025cm,fill=white] (-0.05,0) rectangle (0.05,0.5);
}

% ============================================================
% 5. 移液管 (Pipette)
% ============================================================
\newcommand{\tkzPipetteBody}{%
  \chemlabUpdateLineWidths\begin{scope}[line width=\chemlabLineWidth,rounded corners=0.05cm]
    \draw (-0.05,0) -- (-0.15,0.2) -- (-0.15,6)
          -- (0.15,6) -- (0.15,0.2) -- (0.05,0) -- cycle;
    \fill[white] (0,6) ellipse (0.15 and 0.05);
    \foreach \Level in {0.5,1.0,...,5.0} {%
      \draw[line width=\chemlabGraduationLineWidth] (0,\Level) -- (0.1,\Level);}
  \end{scope}
}

\newcommand{\tkzLargePipetteBody}{%
  \chemlabUpdateLineWidths\begin{scope}[line width=\chemlabLineWidth,rounded corners=0.05cm]
    \draw (0,0) -- (0,1) -- (-0.23,2.5) -- (-0.23,22.5)
          arc[start angle=-80,end angle=80,radius=0.45];
  \end{scope}
}

% ============================================================
% 6. 液滴 (Goutte)
% ============================================================
% 原文件中 \pst@goutte 被重复定义了两次, 此处使用后者 (unit=0.6 的版本)
% pst-labo: \psset{unit=0.6}\psbezier[linewidth=0.017] — linewidth=0.017 是
%   0.017 用户单位 = 0.017*0.6 = 0.0102cm (随 goutte 的 LOCAL unit=0.6 缩放, 
%   不随外层 unit 变化, 因为 \psset{unit=0.6} 覆盖了外层设置)
% TikZ: line width=0.0102cm (ABSOLUTE, 匹配 pst-labo 的固定 0.0102cm)
\newcommand{\tkzDrop}{%
  \begin{scope}[scale=0.6]
    \draw[line width=0.0102cm]
      (0,0) .. controls (0.25,-.4) and (-.25,-.4) .. (0,0);
  \end{scope}
}

% ============================================================
% 7. 量筒 (Eprouvette)
% ============================================================
\newcommand{\tkzGraduatedCylinderBody}{%
  \draw (-0.5,5) -- (-0.5,0) -- (0.5,0) -- (0.5,4.9) -- (0.6,4.93);
}

\newcommand{\tkzGraduatedCylinderRim}{%
  \draw plot[smooth,tension=0.5] coordinates {%
    (0.6,-0.07) (0.498,0.008) (0.483,0.026) (0.433,0.05)
    (0.353,0.0707) (0.25,0.0866) (0.13,0.0966) (0,0.1)
    (-0.13,0.0966) (-0.25,0.0866) (-0.353,0.0707) (-0.433,0.05)
    (-0.483,0.026) (-0.5,0) (-0.483,-0.026) (-0.353,-0.0707)
    (-0.25,-0.0866) (-0.13,-0.096) (0,-0.1) (0.13,-0.096)
    (0.25,-0.0866) (0.353,-0.0707) (0.433,-0.05)
    (0.498,-0.026) (0.6,-0.07)};
}

\newcommand{\tkzGraduatedCylinderBase}{%
  \fill[white]
    ([shift={(-0.95,0.5)}]0:0.45) arc[start angle=0,end angle=-90,radius=0.45]
    -- ([shift={(-0.95,0.025)}]90:0.025) arc[start angle=90,end angle=270,radius=0.025]
    -- (-0.95,0) -- (0.95,0)
    -- ([shift={(0.95,0.025)}]-90:0.025) arc[start angle=-90,end angle=90,radius=0.025]
    -- ([shift={(0.95,0.5)}]-90:0.45) arc[start angle=-90,end angle=-180,radius=0.45]
    -- plot[smooth,tension=0.5] coordinates {(0.5,0.4) (0,0.25) (-0.5,0.4)}
    -- cycle;
  \draw
    ([shift={(-0.95,0.5)}]0:0.45) arc[start angle=0,end angle=-90,radius=0.45]
    -- ([shift={(-0.95,0.025)}]90:0.025) arc[start angle=90,end angle=270,radius=0.025]
    -- (-0.95,0) -- (0.95,0)
    -- ([shift={(0.95,0.025)}]-90:0.025) arc[start angle=-90,end angle=90,radius=0.025]
    -- ([shift={(0.95,0.5)}]-90:0.45) arc[start angle=-90,end angle=-180,radius=0.45]
    -- plot[smooth,tension=0.5] coordinates {(0.5,0.4) (0,0.25) (-0.5,0.4)};
}

\newcommand{\tkzGraduatedCylinderGraduation}{%
  % 量筒刻度 : 主刻度 1.5× 细刻度 (与烧杯一致)
  \pgfmathsetlengthmacro{\chemlab@eprouvette@major}{1.5*\chemlabGraduationLineWidth}%
  \foreach \nyD in {0.5,1.0,...,4.5} { \draw[line width=\chemlab@eprouvette@major] (-0.4,\nyD) -- (-0.1,\nyD); }
  \foreach \nyd in {0.5,0.6,...,4.4} { \draw[line width=\chemlabGraduationLineWidth] (-0.4,\nyd) -- (-0.25,\nyd); }
}

% ============================================================
% 8. 分馏装置 (Distillation fractionnée)
% ============================================================
% \pst@goutte 这里又被定义了一次, 在 TikZ 版中沿用同一个 \tkzDrop

\newcommand{\tkzColumn}{%
  % pst-labo: linearc=0.3 (×unit, 不是绝对值) → TikZ: rounded corners=0.12cm (= 0.3 × 0.4cm)
  % 尖刺 (pointeG/D) 先画并用 clip 裁剪到蒸馏柱内部, 再画轮廓线覆盖切边。
  %   pst-labo 中 linearc=0.3 (R=0.3 user units), 尖刺端点恰在圆弧切点附近, 几乎
  %   不可见。TikZ 的 rounded corners 实际 R 比标称值大 ~2×, 圆弧切点更低,
  %   尖刺端点 (在 y_col=8 和 7.9) 会突出到圆弧外侧形成多余线头。用 clip 将
  %   尖刺裁剪到柱内, 再画轮廓线覆盖切边, 消除突出线头。
  \begin{scope}
    \clip[rounded corners=0.12cm]
      (-.25,10.5) -- (-.25,8.5) -- (-.75,8) -- (-.75,2.5) -- (-.25,2) -- (-.25,0)
      -- (0.25,0) -- (0.25,2) -- (.75,2.5) -- (.75,8) -- (0.25,8.5) -- (0.25,9.36)
      -- (0.25,10.5) -- cycle;
    % 左侧尖刺 (pointeG)
    \foreach \i in {3,4,...,8} {%
      \draw[line width=0.4\pgflinewidth] (-.75,\i) -- (-.25,\i-0.5)
          -- (-.75,\i-0.1);}
    % 右侧尖刺 (pointeD)
    \foreach \i in {3,4,...,8} {%
      \draw[line width=0.4\pgflinewidth] (.75,\i) -- (.25,\i-0.5)
          -- (.75,\i-0.1);}
  \end{scope}
  \draw[rounded corners=0.12cm]
    (-.25,10.5) -- (-.25,8.5) -- (-.75,8) -- (-.75,2.5)
    -- (-.25,2) -- (-.25,0);
  \draw[rounded corners=0.12cm]
    (0.25,0) -- (0.25,2) -- (.75,2.5) -- (.75,8)
    -- (0.25,8.5) -- (0.25,9.36);
  % 上部断口 (line 3a/3b):
  %   pst-labo 源码 line 3a: \psline(0.25,8.5)(0.25,9.30) 用于填充 linearc 圆角
  %   留下的间隙。实测发现 TikZ 的 rounded corners 实际圆弧半径比标称值大
  %   ~2× (在 scale=0.6 下 rounded corners=0.3cm 给出 R≈0.62cm), 圆弧已覆盖到
  %   x=0.25 处 y_col=8.93 附近, 故 line 3a 在任何 scale/line width 下均被圆弧
  %   覆盖, 不可见。为避免在编号图等彩色调试场景中残留多余线段, 此处不再
  %   单独绘制 line 3a。
  \draw (0.25,9.5) -- (0.25,10.5);
  % 填充 / 阴影
  \begin{scope}[pattern={Lines[angle=45,distance={1pt/sqrt(2)}]},pattern color=gray]
    \fill (-.25,10) rectangle (-.05,10.5);
    \fill (.25,10)  rectangle (0.05,10.5);
  \end{scope}
  \foreach \n in {2.5,3.5,...,7.5} {%
    \begin{scope}[shift={(-.25,\n)}]\tkzDrop\end{scope}
    \begin{scope}[shift={(.25,\n)}]\tkzDrop\end{scope}}
  \begin{scope}[shift={(-0.25,-.1)}]\tkzDrop\end{scope}
  \begin{scope}[shift={(0.25,-.1)}]\tkzDrop\end{scope}
}

% 公共别名: pst-labo 中 \tkzColumn 是用户可调用的公共宏 (小写)
% tikz-chemlab 内部用 \tkzColumn (驼峰), 这里注册公共别名保持兼容
\let\tkzColumn\tkzColumn

% réfrigérant (condenseur)
\newcommand{\tkzCondenser}{%
  \draw ([shift={(3,0)}]90:0.5) arc[start angle=90,end angle=270,radius=0.5];
  % pst-labo: \psarc(7,0){0.5}{270}{90} → 逆时针经过 0°/右侧
  \draw ([shift={(7,0)}]270:0.5) arc[start angle=270,end angle=450,radius=0.5];
  \fill[white] (0,-.05) rectangle (9,0.05);
  \begin{scope}[line width=0.4pt]
    \draw (-.075,.05) -- (9,0.05);
    \draw (0.02,-.05) -- (9,-.05);
  \end{scope}
  \draw (3,0.5) -- (3.5,0.5);
  \draw (3,-.5) -- (6.5,-.5);
  \draw (3.6,0.5) -- (7,0.5);
  \begin{scope}[line width=0.4pt]
    \draw (3.5,0.5) -- (3.5,1.5);
    \draw (3.6,0.5) -- (3.6,1.5);
    \draw (6.6,-.5) -- (6.6,-1.5);
    \draw (6.5,-.5) -- (6.5,-1.5);
  \end{scope}
  \draw (6.6,-.5) -- (7,-.5);
  \draw[->] (6.55,-2.5) -- (6.55,-1.5);
  \draw[->] (3.55,1.5)  -- (3.55,2.5);
}

\newcommand{\tkzThermometer}{%
  % pst-labo: \psline[linearc=0.05](-.05,.25)(-.05,2.75)(0.05,2.75)(0.05,0.25)
  % linearc=0.05 × unit, TikZ: rounded corners=0.05 (受 scale 缩放)
  \draw[rounded corners=0.025cm]
    (-.05,.25) -- (-.05,2.75) -- (0.05,2.75) -- (0.05,0.25) -- cycle;
  % pst-labo: \psframe[framearc=0.8,fillstyle=solid,fillcolor=red](-.15,-.25)(0.15,0.25)
  % framearc=0.8: 80% of half-min(0.3,0.5)=0.15 → radius=0.12
  \filldraw[fill=red,draw=black,rounded corners=0.06cm] (-.15,-.25) rectangle (0.15,0.25);
  % pst-labo: \psline[linecolor=red](0,2) — 红色竖线, 使用默认线宽
  % 在 pst-labo 渲染中较粗, 看起来像红色填充
  \draw[red,line width=0.8pt] (0,0) -- (0,2);
}

\newcommand{\tkzHeater}{%
  \draw (1.17,2.5) -- (0,2.5) -- (0,0) -- (5.5,0) -- (5.5,2.5) -- (4.33,2.5);
}

% 内部小量筒轮廓 - 路径版本 (供 clip 使用)
\newcommand{\tkzGraduatedCylinderPath}{%
  (-0.5,6) -- (-0.5,.5)
  -- ([shift={(-1,0.5)}]0:0.5) arc[start angle=0,end angle=-90,radius=0.5]
  -- (-1,0) -- (1,0)
  -- ([shift={(1,0.5)}]270:0.5) arc[start angle=270,end angle=180,radius=0.5]
  -- (0.5,0.5) -- (0.5,6) -- cycle
}

% 内部小量筒 (描边)
% pst-labo: \psset{linewidth=1.5\pslinewidth} — \pslinewidth 默认 0.8pt (ABSOLUTE)
% → 1.5*0.8pt = 1.2pt (绝对值, 不随 unit/scale 缩放)
% 注意: TikZ 的 [line width=N\pgflinewidth] 有循环引用 bug — \pgflinewidth 在
% 乘法前被更新, 导致 line width=3\pgflinewidth 实际给出 3.6pt 而非 1.2pt。
% 必须用显式 pt 值。
\newcommand{\tkzGraduatedCylinderSmall}{%
  \begin{scope}[line width=0.72pt,line cap=round,line join=round]
    \draw (-0.5,6) -- (-0.5,.5)
      -- ([shift={(-1,0.5)}]0:0.5) arc[start angle=0,end angle=-90,radius=0.5]
      -- (-1,0) -- (1,0)
      -- ([shift={(1,0.5)}]270:0.5) arc[start angle=270,end angle=180,radius=0.5]
      -- (0.5,0.5) -- (0.5,6);
  \end{scope}
}

% 带液体的量筒 (内部使用, 重命名以避免与主文件 \pstEproulette 冲突)
\newcommand{\tkzGraduatedCylinderWithLiquid}{%
  \begin{scope}
    \clip \tkzGraduatedCylinderPath;
    % pst-labo: \psframe*[linecolor=\psk@Distillation@DistillateColor]
    \fill[\tkzDistillateColor] (-1.5,-1) rectangle (1.5,1);
  \end{scope}
  \tkzGraduatedCylinderSmall
}

% (注 : 原 PSTricks 中 \pst@roundFlask 已合并到 \tkzRoundFlaskSimple, 此处不再单独定义,
%  避免与 tikz-chemlab.tex 中的用户命令 \tkzRoundFlask 冲突)

% 自定义密集斜线图案 (匹配 pst-labo hlines 的密度)
{\pgfpointorigin}{\pgfpoint{0.012cm}{0.012cm}}{\pgfpoint{0.012cm}{0.012cm}}
{\pgfsetlinewidth{0.0028cm}\pgfpathmoveto{\pgfpoint{0cm}{0cm}}\pgfpathlineto{\pgfpoint{0.012cm}{0.012cm}}\pgfusepath{stroke}}
{\pgfpointorigin}{\pgfpoint{0.012cm}{0.012cm}}{\pgfpoint{0.012cm}{0.012cm}}
{\pgfsetlinewidth{0.004cm}\pgfpathmoveto{\pgfpoint{0cm}{0cm}}\pgfpathlineto{\pgfpoint{0.012cm}{0.012cm}}\pgfusepath{stroke}}
% \pst@DistillationFractionnee 主图
\newcommand{\tkzDistillationFractionated}{%
  % --- roundFlask ---
  \def\flaskLocal{%
    % pst-labo: fillstyle=hlines,hatchwidth=0.007,hatchsep=0.03 → 水平线填充
    \begin{scope}[pattern={Lines[angle=45,distance={1pt/sqrt(2)}]},pattern color=gray]
      \fill (-.5,2.25) rectangle (-.25,3);
      \fill (0.5,2.25)  rectangle (0.25,3);
    \end{scope}
    % pst-labo 线宽: neck=0.053, arcs=0.03 (都是 ×unit 用户单位, 随 scale 缩放)
    % \chemlabUpdateLineWidths 已设置 \chemlabLineWidth = 0.053*scale
    \chemlabUpdateLineWidths
    \draw[line width=\chemlabLineWidth] (0.5,1.5) -- (0.5,3);
    \draw[line width=\chemlabLineWidth] (-.5,1.5) -- (-.5,3);
    % pst-labo: \psarc(0,0){1.581}{108.435}{71.565} 逆时针走长路 (经底部),
    %   左右侧弧可见, 底部被液体填充覆盖, 顶部无可见弧线。此处拆分为
    %   左右侧两段 (底部由下方轮廓弧绘制, 顶部不画)。
    % 右上侧弧: 71° → 0° (连接右颈底部到球体最右点)
    \draw[line width=\chemlabLineWidth] ([shift={(0,0)}]71.565:1.581) arc[start angle=71.565,end angle=0,radius=1.581];
    % 左上侧弧: 180° → 108° (连接球体最左点到左颈底部)
    \draw[line width=\chemlabLineWidth] ([shift={(0,0)}]180:1.581) arc[start angle=180,end angle=108.435,radius=1.581];
    % pst-labo: \psset{style=\psk@Distillation@MixtureStyle}\fill
    % 需要用 \shade (不是 \fill) 来产生渐变效果
    \shade[\tkzMixtureStyle] plot[smooth,tension=0.5] coordinates {%
      (-1.581,0) (-1.25,0.1) (-.75,-.1) (0,.1) (0.75,-.1) (1.25,.1) (1.581,0)}
      arc[start angle=0,end angle=-180,radius=1.581];
    % 底部轮廓弧: pst-labo \psarc[linewidth=0.03] → 0.03*unit, 用 \chemlabLineWidth 0.053
    % (pst-labo 中 0.03 vs 0.053 差异微小, 统一用 0.053*scale 避免额外变量)
    \draw[line width=\chemlabLineWidth] (0,0) ++(180:1.581) arc[start angle=180,end angle=360,radius=1.581];
  }%
  % --- 主图整体平移 ---
  \begin{scope}[shift={(0,-3)}]
    \begin{scope}[shift={(2.75,13.5)}]\tkzThermometer\end{scope}
    \begin{scope}[shift={(2.75,2.5)}]\flaskLocal\end{scope}
    \begin{scope}[shift={(2.75,4)}]\tkzColumn\end{scope}
    \begin{scope}[shift={(3,13.4)},rotate=-45]\tkzCondenser\end{scope}
    \begin{scope}[shift={(0,0)}]\tkzHeater\end{scope}
    \begin{scope}[shift={(9.4,0)}]\tkzGraduatedCylinderWithLiquid
      % 蒸馏装置量筒刻度 : 主刻度 1.5× 细刻度
      \pgfmathsetlengthmacro{\chemlab@distill@major}{1.5*\chemlabGraduationLineWidth}%
      \foreach \Level in {1,2,...,5} {\draw[line width=\chemlab@distill@major] (-.4,\Level) -- (-.1,\Level);}
      \foreach \Level in {0.2,0.4,...,4.8} {\draw[line width=\chemlabGraduationLineWidth] (-.4,\Level) -- (-.2,\Level);}
    \end{scope}
    \begin{scope}[shift={(9.37,7)}]\tkzDrop\end{scope}
    \begin{scope}[shift={(9.37,6)}]\tkzDrop\end{scope}
    \begin{scope}[shift={(9.37,4)}]\tkzDrop\end{scope}
    \begin{scope}[shift={(9.37,1.2)}]\tkzDrop\end{scope}
    % 加热盘 (搅拌器)
    \begin{scope}[shift={(0.75,.75)}]
      \draw (0,0) circle (.5);
      \foreach \i in {0,30,...,180} {\draw (\i:0.55) -- (\i:0.65);}
      \draw[->,line width=0.035cm] (0,0) -- (0,0.45);
    \end{scope}
  \end{scope}
}

% ============================================================
% 9. 本生灯 (Bec Bunsen) - 不带 grille
% ============================================================
\newcommand{\tkzBunsenBurner}{%
  \chemlabUpdateLineWidths\begin{scope}[line width=\chemlabLineWidth,line cap=round,line join=round]
    \draw (-1.25,0) rectangle (1.25,0.25);
    \draw (-.5,1.25) rectangle (0.5,2.25);
    \foreach \n in {-0.3,0.0,0.3} {\draw (\n,1.75) circle (0.1);}
    \draw (-.25,2.25) rectangle (0.25,4.25);
    \draw (0.25,1.25) -- (0.25,0.5) -- (1.25,0.25);
    \draw (-1.25,0.25) -- (-.25,0.5) -- (-0.25,0.75);
    \draw (-2.25,0.75) -- (-.25,0.75);
    \draw (-2.25,1) -- (-.25,1);
    \draw (-.25,0.875) ellipse (0.1 and 0.125);
    \fill[white] (-2.25,0.75) rectangle (-0.25,1);
    \draw (-2.25,0.75) -- (-0.25,0.75);
    \draw (-2.25,1) -- (-0.25,1) -- (-.25,1.25);
    \draw plot[smooth,tension=0.5] coordinates {(-0.25,0.5) (0,0.4) (0.25,0.5)};
    % flamme
    % pst-labo 源码: \rput(0,4.25){\psclip{\psbezier[...gradient...] (12点)}\pspolygon[...gradient...](-0.25,0)(0.25,0)(0,1)\endpsclip}
    % \psbezier 12个点 + BezierNArray padding (0,0) = 13点 = 4段贝塞尔
    % 经像素级验证, 正确路径从 (0,0) 出发, 经4段贝塞尔回到 (0,0), 无闭合直线:
    %   Seg 1: (0,0) → c1=P1=(-0.25,0), c2=P2=(-0.35,0.5) → P3=(-0.4,0.75) [左下→左中]
    %   Seg 2: P3=(-0.4,0.75) → c1=P4=(-0.35,1), c2=P5=(-0.25,1.5) → P6=(0.5,2) [左中→顶]
    %   Seg 3: P6=(0.5,2) → c1=P7=(0.25,1.5), c2=P8=(0.35,1) → P9=(0.4,0.75) [顶→右中]
    %   Seg 4: P9=(0.4,0.75) → c1=P10=(0.35,0.5), c2=P11=(0.25,0) → (0,0) [右中→底]
    % 渐变: gradbegin=PaleOrange (底), gradend=yellow (顶), gradmidpoint=0, gradangle=0 (垂直)
    \begin{scope}[shift={(0,4.25)}]
      % 外焰填充: PaleOrange (底) -> yellow (顶) 线性垂直渐变
      \shade[top color=yellow, bottom color=PaleOrange]
        (0,0) .. controls (-0.25,0) and (-0.35,0.5)
        .. (-0.4,0.75) .. controls (-0.35,1) and (-0.25,1.5)
        .. (0.5,2) .. controls (0.25,1.5) and (0.35,1)
        .. (0.4,0.75) .. controls (0.35,0.5) and (0.25,0)
        .. (0,0);
      % 内焰: 三角形 cyan(底)->white(顶) 渐变, 裁剪到外焰形状内
      \begin{scope}
        \clip
          (0,0) .. controls (-0.25,0) and (-0.35,0.5)
          .. (-0.4,0.75) .. controls (-0.35,1) and (-0.25,1.5)
          .. (0.5,2) .. controls (0.25,1.5) and (0.35,1)
          .. (0.4,0.75) .. controls (0.35,0.5) and (0.25,0)
          .. (0,0);
        \shade[top color=white, bottom color=cyan]
          (-0.25,0) -- (0.25,0) -- (0,1) -- cycle;
      \end{scope}
    \end{scope}
  \end{scope}
}

% 火焰 + 网格 (FlammeGrille 形状) - 路径版本 (供 clip 使用)
\newcommand{\tkzFlameGrillPath}{%
  ([shift={(0,0.75)}]270:0.75) arc[start angle=270,end angle=360,radius=0.75]
  -- ([shift={(1.5,0.75)}]180:0.75) arc[start angle=180,end angle=90,radius=0.75]
  -- (1.5,1.5) -- (-1.5,1.5)
  -- ([shift={(-1.5,0.75)}]90:0.75) arc[start angle=90,end angle=0,radius=0.75]
  -- ([shift={(0,0.75)}]180:0.75) arc[start angle=180,end angle=270,radius=0.75]
  -- cycle
}

\newcommand{\tkzFlameGrill}{%
  \fill[flameAndGrill] \tkzFlameGrillPath;
}

% ============================================================
% 10. 本生灯 - 带 grille
% ============================================================
\newcommand{\tkzBunsenBurnerGrill}{%
  \chemlabUpdateLineWidths\begin{scope}[line width=\chemlabLineWidth,line cap=round,line join=round]
    \draw (-1.25,0) rectangle (1.25,0.25);
    \draw (-.5,1.25) rectangle (0.5,2.25);
    \foreach \n in {-0.3,0.0,0.3} {\draw (\n,1.75) circle (0.1);}
    \draw (-.25,2.25) rectangle (0.25,4.25);
    \draw (0.25,1.25) -- (0.25,0.5) -- (1.25,0.25);
    \draw (-1.25,0.25) -- (-.25,0.5) -- (-0.25,0.75);
    \draw (-2.25,0.75) -- (-.25,0.75);
    \draw (-2.25,1) -- (-.25,1);
    \draw (-.25,0.875) ellipse (0.1 and 0.125);
    \fill[white] (-2.25,0.75) rectangle (-0.25,1);
    \draw (-2.25,0.75) -- (-0.25,0.75);
    \draw (-2.25,1) -- (-0.25,1) -- (-.25,1.25);
    \draw plot[smooth,tension=0.5] coordinates {(-0.25,0.5) (0,0.4) (0.25,0.5)};
    % flamme (clip 到 FlammeGrille)
    % 原 pst-labo 火焰分两层:
    %   外焰: \pst@FlammeGrille 整体填充 flameAndGrill (PaleOrange->yellow 渐变)
    %   内焰: 三角形 cyan->white 渐变 (在内焰区域)
    % 颜色和大小可通过 \chemlabset 调整
    \begin{scope}[shift={(0,4.25)},scale=\chemlabFlameScale]
      \begin{scope}
        \clip \tkzFlameGrillPath;
        % 外焰: PaleOrange (底) -> yellow (顶) 线性垂直渐变
        % 对应 pst-labo: fillstyle=gradient, gradbegin=PaleOrange, gradend=yellow
        % 渐变范围与火焰形状 (y=0..1.5) 一致, 避免底部颜色被稀释
        \shade[top color=\chemlabFlameOuterEndColor, bottom color=\chemlabFlameOuterColor]
          (-1.5,0) rectangle (1.5,1.5);
        % 内焰: cyan (底) -> white (顶) 线性垂直渐变 (三角形)
        \shade[top color=\chemlabFlameInnerEndColor, bottom color=\chemlabFlameInnerColor]
          (-0.25,0) -- (0.25,0) -- (0,1) -- cycle;
      \end{scope}
    \end{scope}
    \draw[dashed,line width=\chemlabEllipseLineWidth] (-2,5.75) -- (2,5.75);
  \end{scope}
}

% ============================================================
% 11. 烧瓶 (Ballon)
% ============================================================
% 路径版本 - 供 clip 使用 (闭合轮廓)
\newcommand{\tkzRoundFlaskSimplePath}{%
  ([shift={(-0.7,3.4)}]90:0.2) arc[start angle=90,end angle=0,radius=0.2]
  -- (-0.5,3.4) -- (-0.5,2.4)
  -- (-0.489,2.401)
  arc[start angle=113,end angle=427,radius=1.25]
  -- (0.5,2.4) -- (0.5,3.4)
  -- ([shift={(0.7,3.4)}]180:0.2) arc[start angle=180,end angle=90,radius=0.2]
  -- (0,3.5) -- cycle
}

\newcommand{\tkzRoundFlaskSimple}{%
  \chemlabUpdateLineWidths\begin{scope}[line width=\chemlabLineWidth,line cap=round,line join=round]
    % 左侧: 瓶口圆弧 + 颈部直线 + 球体大圆弧, 合并为单条路径以平滑过渡
    \draw ([shift={(-0.7,3.4)}]90:0.2) arc[start angle=90,end angle=0,radius=0.2]
          -- (-0.5,3.4) -- (-0.5,2.4)
          -- (-0.489,2.401) arc[start angle=113,end angle=427,radius=1.25];
    % 右侧: 球体终点(0.489,2.401) -> 颈部直线 -> 瓶口圆弧
    \draw (0.489,2.401) -- (0.5,2.4) -- (0.5,3.4)
          -- ([shift={(0.7,3.4)}]180:0.2) arc[start angle=180,end angle=90,radius=0.2];
    % 瓶口椭圆: 白色填充覆盖颈部上端圆头, 避免穿入椭圆
    % pst-labo: ballon椭圆用默认线宽 0.053 (与主轮廓一致, 不是 0.07)
    \filldraw[fill=white,draw=black,line width=\chemlabLineWidth] (0,3.6) ellipse (0.7 and 0.1);
  \end{scope}
}

\newcommand{\tkzRoundFlaskGrill}{%
  \chemlabUpdateLineWidths\begin{scope}[line width=\chemlabLineWidth,line cap=round,line join=round]
    % 左侧: 瓶口圆弧 + 颈部直线 + 球体大圆弧, 合并为单条路径
    \draw ([shift={(-0.7,2.15)}]90:0.2) arc[start angle=90,end angle=0,radius=0.2]
          -- (-0.5,2.15) -- (-0.5,1.4)
          -- (-0.489,1.401) arc[start angle=113,end angle=427,radius=1.25];
    % 右侧: 球体终点 -> 颈部直线 -> 瓶口圆弧
    \draw (0.489,1.401) -- (0.5,1.15) -- (0.5,3.4)
          -- ([shift={(0.7,2.15)}]180:0.2) arc[start angle=180,end angle=90,radius=0.2];
    % 瓶口椭圆: 白色填充覆盖颈部上端圆头, 避免穿入椭圆
    % pst-labo: ballon椭圆用默认线宽 0.053 (与主轮廓一致)
    \filldraw[fill=white,draw=black,line width=\chemlabLineWidth] (0,2.35) ellipse (0.7 and 0.1);
  \end{scope}
}

% ============================================================
% 12. 试管 (TubeEssais)
% ============================================================
\newcommand{\tkzTestTubeStopper}{%
  \chemlabUpdateLineWidths
  \begin{scope}[fill=LightGray,draw=black,line width=\chemlabLineWidth,line cap=round,line join=round]
    \filldraw (0,-0.3) ellipse (0.4 and 0.1);
    \filldraw (-0.4,-0.3) -- (-0.6,0.5) -- (0.6,0.5) -- (0.4,-0.3) -- cycle;
    \fill[draw=none] (0,-0.3) ellipse (0.4 and 0.1);
    \filldraw (0,0.5) ellipse (0.6 and 0.1);
  \end{scope}
  \draw[line width=\chemlabEllipseLineWidth] plot[smooth,tension=0.5]
        coordinates {(-0.5,0) (0,-0.07) (0.5,0)};
}

\newcommand{\tkzTestTubeBody}{%
  \chemlabUpdateLineWidths\begin{scope}[line width=\chemlabLineWidth,line cap=round,line join=round]
    % 合并左侧直线 + 底部圆弧 + 右侧直线为单条路径, 平滑转角
    \draw (-0.5,3.5) -- (-0.5,0.5)
          arc[start angle=180,end angle=360,radius=0.5]
          -- (0.5,3.5);
    % 试管口椭圆 (原 pst-labo 用 0.08cm, 比其他椭圆粗)
    \draw[line width=\chemlabTubeEllipseLineWidth] (0,3.5) ellipse (0.5 and 0.1);
  \end{scope}
}

\newcommand{\tkzTestTubeStraightTube}{%
  \draw[line width=0.5\pgflinewidth] (0.05,-1) -- (0.05,5);
  \draw[line width=0.5\pgflinewidth] (-0.05,-1) -- (-0.05,5);
  \tkzTestTubeStopper
  \draw[fill=white,rounded corners=0.5,line width=0.5\pgflinewidth]
       (-0.05,5) rectangle (0.05,0.5);
  \draw[white,line width=0.14cm] (-0.1,5) -- (0.1,5);
}

\newcommand{\tkzTestTubeElbowTube}{%
  \draw[line width=0.5\pgflinewidth] (0.05,-1) -- (0.05,0);
  \draw[line width=0.5\pgflinewidth] (-0.05,-1) -- (-0.05,0);
  \tkzTestTubeStopper
  % 弯管路径: pst-labo \pspolygon(...)(8,2)(8,1.9)...
  \filldraw[fill=white,rounded corners=0.05cm,line width=0.5\pgflinewidth]
       (-0.05,0.5) -- (-0.05,2) -- (8,2) -- (8,1.9)
       -- (0.05,1.9) -- (0.05,0.5) -- cycle;
  \draw[white,line width=0.14cm] (7.95,1.9) -- (8.05,2);
}

\newcommand{\tkzTestTubeUTube}{%
  \draw[line width=0.5\pgflinewidth] (0.05,-1) -- (0.05,0);
  \draw[line width=0.5\pgflinewidth] (-0.05,-1) -- (-0.05,0);
  \tkzTestTubeStopper
  \filldraw[fill=white,rounded corners=0.05cm,line width=0.5\pgflinewidth]
       (-0.05,0.5) -- (-0.05,2) -- (4.5,2) -- (4.5,-3)
       -- (4.4,-3) -- (4.4,1.9) -- (0.05,1.9) -- (0.05,0.5) -- cycle;
  \draw[white,line width=0.14cm] (4.3,-3) -- (4.6,-3);
}

\newcommand{\tkzTestTubeUTubeLong}{%
  \draw[line width=0.5\pgflinewidth] (0.05,-1) -- (0.05,0);
  \draw[line width=0.5\pgflinewidth] (-0.05,-1) -- (-0.05,0);
  \tkzTestTubeStopper
  \filldraw[fill=white,rounded corners=0.05cm,line width=0.5\pgflinewidth]
       (-0.05,0.5) -- (-0.05,2) -- (4.5,2) -- (4.5,-9)
       -- (4.4,-9) -- (4.4,1.9) -- (0.05,1.9) -- (0.05,0.5) -- cycle;
  \draw[white,line width=0.14cm] (4.3,-9) -- (4.6,-9);
}

\newcommand{\tkzCurvedTube}{%
  \draw[line width=0.5\pgflinewidth] (0.05,-1) -- (0.05,0);
  \draw[line width=0.5\pgflinewidth] (-0.05,-1) -- (-0.05,0);
  \tkzTestTubeStopper
  \begin{scope}[line width=0.4pt]
    \draw (-0.05,0.5) -- (-0.05,1);
    \draw ([shift={(0.5,1)}]180:0.55) arc[start angle=180,end angle=90,radius=0.55];
    \draw (0.5,1.55) -- (5.5,1.55);
    \draw ([shift={(5.5,1)}]90:0.55) arc[start angle=90,end angle=0,radius=0.55];
    \draw (6.05,1) -- (6.05,-8.8);
    \draw ([shift={(6.5,-8.8)}]180:0.45) arc[start angle=180,end angle=360,radius=0.45];
    \draw (6.95,-8.8) -- (6.95,-8.3);
    \draw (7.05,-8.3) -- (7.05,-8.8);
    \draw ([shift={(6.5,-8.8)}]0:0.55) arc[start angle=0,end angle=-180,radius=0.55];
    \draw (5.95,-8.8) -- (5.95,1);
    \draw ([shift={(5.5,1)}]0:0.45) arc[start angle=0,end angle=90,radius=0.45];
    \draw (5.5,1.45) -- (0.5,1.45);
    \draw ([shift={(0.5,1)}]90:0.45) arc[start angle=90,end angle=180,radius=0.45];
    \draw (0.05,1) -- (0.05,0.5);
    \draw (-0.05,0.5) -- (0.05,0.5);
  \end{scope}
  \fill[white] (-0.025,0.52) rectangle (0.020,0.6);
}

\newcommand{\tkzCurvedTubeShort}{%
  \draw[line width=0.5\pgflinewidth] (0.05,-1) -- (0.05,0);
  \draw[line width=0.5\pgflinewidth] (-0.05,-1) -- (-0.05,0);
  \tkzTestTubeStopper
  \draw (-0.05,0.5) -- (-0.05,1);
  \begin{scope}[line width=0.4pt]
    \draw ([shift={(0.5,1)}]180:0.55) arc[start angle=180,end angle=90,radius=0.55];
    \draw (0.5,1.55) -- (3.5,1.55);
    \draw ([shift={(3.5,1)}]90:0.55) arc[start angle=90,end angle=0,radius=0.55];
    \draw (4.05,1) -- (4.05,-2.8);
    \draw ([shift={(4.5,-2.8)}]180:0.45) arc[start angle=180,end angle=360,radius=0.45];
    \draw (4.95,-2.8) -- (4.95,-2.3);
    \draw (5.05,-2.3) -- (5.05,-2.8);
    \draw ([shift={(4.5,-2.8)}]0:0.55) arc[start angle=0,end angle=-180,radius=0.55];
    \draw (3.95,-2.8) -- (3.95,1);
    \draw ([shift={(3.5,1)}]0:0.45) arc[start angle=0,end angle=90,radius=0.45];
    \draw (3.5,1.45) -- (0.5,1.45);
    \draw ([shift={(0.5,1)}]90:0.45) arc[start angle=90,end angle=180,radius=0.45];
    \draw (0.05,1) -- (0.05,0.5);
    \draw (-0.05,0.5) -- (0.05,0.5);
  \end{scope}
  \fill[white] (-0.025,0.52) rectangle (0.020,0.6);
}

\newcommand{\tkzCurvedTubeShortDoubleTube}{%
  \begin{scope}[line width=0.4pt]
    \draw (-0.05,0.5) -- (-0.05,1);
    \draw ([shift={(0.5,1)}]180:0.55) arc[start angle=180,end angle=90,radius=0.55];
    \draw (0.5,1.55) -- (3.5,1.55);
    \draw ([shift={(3.5,1)}]90:0.55) arc[start angle=90,end angle=0,radius=0.55];
    \draw (4.05,1) -- (4.05,-2.8);
    \draw ([shift={(4.5,-2.8)}]180:0.45) arc[start angle=180,end angle=360,radius=0.45];
    \draw (4.95,-2.8) -- (4.95,-2.3);
    \draw (5.05,-2.3) -- (5.05,-2.8);
    \draw ([shift={(4.5,-2.8)}]0:0.55) arc[start angle=0,end angle=-180,radius=0.55];
    \draw (3.95,-2.8) -- (3.95,1);
    \draw ([shift={(3.5,1)}]0:0.45) arc[start angle=0,end angle=90,radius=0.45];
    \draw (3.5,1.45) -- (0.5,1.45);
    \draw ([shift={(0.5,1)}]90:0.45) arc[start angle=90,end angle=180,radius=0.45];
    \draw (0.05,1) -- (0.05,0.5);
    \draw (-0.05,0.5) -- (0.05,0.5);
  \end{scope}
  \fill[white] (-0.025,0.52) rectangle (0.018,0.8);
}

\newcommand{\tkzTestTubeDoubleTube}{%
  \begin{scope}[line width=0.4pt]
    \draw (-0.3,-3) -- (-0.3,0);
    \draw (-.2,-3) -- (-0.2,0);
    \draw (0.25,-1) -- (0.25,0);
    \draw (.15,-1) -- (.15,0);
  \end{scope}
  \tkzTestTubeStopper
  \begin{scope}[shift={(0.2,0)}]\tkzCurvedTubeShortDoubleTube\end{scope}
  \begin{scope}[shift={(0,0.5)}]
    \begin{scope}[line width=0.4pt]
      \filldraw[fill=white,rounded corners=0.05cm]
        (-.3,0) -- (-0.3,2) -- (-0.5,2.2) -- (-0.5,3) -- (0,3)
                -- (0,2.2) -- (-0.2,2) -- (-0.2,0) -- cycle;
      \draw (0.1,1.55) -- (-.5,1.55) -- (-0.5,1.7) -- (-.55,1.7)
            -- (-.55,1.3) -- (-0.5,1.3) -- (-0.5,1.55) -- (-.5,1.45) -- (.1,1.45);
      \filldraw[fill=white,rounded corners=0.2] (-0.4,1.3) rectangle (-0.1,1.7);
    \end{scope}
    \fill[white] (-0.25,3) ellipse (0.3 and 0.1);
  \end{scope}
}

% 试管夹 (Pince) - 几何与原 pst-labo pst-laboObj.tex line 483-508 一致
% 根据原 pst-labo 实际渲染 (5.3x 截图) 校准:
%   - 上边缘凹弧较浅 (原 \pscurve 3点几乎共线, 浅弧)
%   - 顶部小凸起为圆弧形 (\psarcn 圆弧)
%   - 所有转角为圆角 (原 \pscustom 默认 linejoin=0 即 round)
\newcommand{\tkzTestTubeClamp}{%
  % 第1块: 上部夹爪
  % 上边缘(1a, 夹合面): 下凹弧贴合试管顶部弧面
  %   试管半径 0.5, 在 y=3.2 处试管顶部弧面下凹约 0.1, 用贝塞尔曲线模拟
  %   从 (0.5,0) 经 (0,-0.1) 到 (-0.5,0), 控制点 y=-0.13 产生平滑下凹
  \filldraw[fill=PaleOrange,draw=black,line cap=round,line join=round]
    (0.5,0) .. controls (0.25,-0.13) and (-0.25,-0.13) .. (-0.5,0)
    -- (-0.8,0)
    -- ([shift={(-0.5,0)}]180:0.3) arc[start angle=180,end angle=270,radius=0.3]
    -- (-.5,-.3) -- (2.5,-.3) -- (2.4,-.2) -- (1,0) -- (0.5,0) -- cycle;
  % 第2块: 下部夹爪
  % 注: 原 \psarcn(-.5,0){0.3}{-90}{180} 是顺时针从 -90°(=270°) 到 180°, 跨度 90° (四分之一圆弧)
  % TikZ arc 当 end<start 时顺时针, 需用 end=-180 让其走短弧 (避免四分之三圆弧)
  \filldraw[fill=PaleOrange,draw=black,line cap=round,line join=round]
    (-0.8,0) -- (-0.8,-0.3)
    -- ([shift={(-0.5,-0.3)}]180:0.3) arc[start angle=180,end angle=270,radius=0.3]
    -- (-0.5,-0.6) -- (2.5,-.6) -- (2.5,-.3) -- (-.5,-.3)
    -- ([shift={(-.5,0)}]-90:0.3) arc[start angle=-90,end angle=-180,radius=0.3] -- cycle;
  % 第3块: 顶部小凸起 - 原 \psarcn 顺时针 180°->90°, 圆弧形凸起
  % 去掉 (-0.8,0.05)->(-0.75,0.05) 的多余水平线段, 直接从 (-0.5,0.05) 到圆弧起点
  % 使凸起呈纯圆弧形 (与原 pst-labo 渲染一致)
  \filldraw[fill=PaleOrange,draw=black,line cap=round,line join=round]
    (-.5,0.05) -- ([shift={(-0.5,0.05)}]180:0.25)
    arc[start angle=180,end angle=90,radius=0.25]
    -- (-.5,0.3) -- (-0.5,0.05) -- cycle;
  % 第4块: 夹臂上部
  \filldraw[fill=PaleOrange,draw=black,line cap=round,line join=round]
    (0.5,0.3) -- (4.9,0.3) -- (5,0.2) -- (1,0) -- (0.5,0) -- (0.5,0.3) -- cycle;
  % 第5块: 夹臂下部 (Brown 棕色)
  \filldraw[fill=Brown,draw=black,line cap=round,line join=round]
    (0.5,0) -- (5,0.2) -- (5,0) -- (2,-0.2) -- (1,0) -- cycle;
  % 螺丝椭圆
  \draw[line width=1.5\pgflinewidth,line cap=round] (1,0) ellipse (0.3 and 0.1);
  % 螺丝下方的杆
  \draw[line width=1.5\pgflinewidth,line cap=round,line join=round]
        (0.7,0) -- (0.4,-.3) -- (0.4,-.6);
}

% ============================================================
% 13. 容量瓶 (Fiole jaugée)
% ============================================================
\newcommand{\tkzVolumetricFlask}{%
  \chemlabUpdateLineWidths
  \draw[rounded corners=0.1cm,line width=\chemlabLineWidth,line cap=round,line join=round]
       (-0.21,3.5) -- (-0.1,3.5) -- (-0.1,1.5)
       -- (-0.75,0) -- (0.75,0) -- (0.1,1.5) -- (0.1,3.5) -- (0.21,3.5);
  \draw[line width=\chemlabGraduationLineWidth,line cap=round] (-0.1,2.5) -- (0.1,2.5);
}

% ============================================================
% 14. 结晶皿 (Cristallisoir)
% ============================================================
\newcommand{\Cristallisoir}{%
  \chemlabUpdateLineWidths
  \begin{scope}[line width=\chemlabLineWidth,line cap=round,line join=round]
  \draw (-1.9,2) -- (-2,2);
  \draw ([shift={(-2,1.95)}]90:0.05) arc[start angle=90,end angle=270,radius=0.05];
  \draw (-2,1.9) -- (-2,1);
  \draw ([shift={(-1,1)}]180:1) arc[start angle=180,end angle=270,radius=1];
  \draw (-1,0) -- (1,0);
  \draw ([shift={(1,1)}]270:1) arc[start angle=270,end angle=360,radius=1];
  \draw (2,1) -- (2,1.9);
  \draw ([shift={(2,1.95)}]270:0.05) arc[start angle=270,end angle=90,radius=0.05];
  \draw (2,2) -- (1.9,2);
  \draw (1.9,2) -- (1.9,1);
  \draw ([shift={(1,1)}]0:0.9) arc[start angle=0,end angle=-90,radius=0.9];
  \draw (1,0.1) -- (-1,0.1);
  \draw ([shift={(-1,1)}]270:0.9) arc[start angle=270,end angle=180,radius=0.9];
  \draw (-1.9,1) -- (-1.9,2);
  \end{scope}
}

\newcommand{\InterieurCristallisoir}{%
  (1.9,2) -- (1.9,1)
  -- ([shift={(1,1)}]0:0.9) arc[start angle=0,end angle=-90,radius=0.9]
  -- (1,0.1) -- (-1,0.1)
  -- ([shift={(-1,1)}]270:0.9) arc[start angle=270,end angle=180,radius=0.9]
  -- (-1.9,1) -- (-1.9,2) -- cycle
}

\newcommand{\tkzCrystallizingDish}{%
  \begin{scope}
    \clip \InterieurCristallisoir;
    \fill[LightBlue1] (-3,-1) rectangle (3,1.5);
  \end{scope}
  \Cristallisoir
}

% ============================================================
% 15. 倒置试管 (TubeRenverse)
% ============================================================
\newcommand{\tkzInvertedTube}{%
  \chemlabUpdateLineWidths
  \begin{scope}[line width=\chemlabEllipseLineWidth,line cap=round,line join=round]
    \fill[LightBlue1] (-0.5,0) rectangle (0.5,5.5);
    % 合并左侧 + 顶部圆弧 + 右侧为单条路径
    % pst-labo: \psarcn(0,6){0.5}{180}{0} = 从180°顺时针到0°, 经过顶部90°
    % TikZ: arc[start=180, end=0] = 从180°逆时针(减小)到0°, 经过顶部90°
    \draw (-0.5,0) -- (-0.5,6)
          arc[start angle=180,end angle=0,radius=0.5]
          -- (0.5,0);
  \end{scope}
}

% ============================================================
% 16. 试管架 (ReposeTube)
% ============================================================
\newcommand{\tkzTestTubeRack}{%
  \begin{scope}[line width=0.1cm,draw=lightgray]
    \draw ([shift={(0,-0.25)}]100:1) arc[start angle=100,end angle=140,radius=1];
    \draw ([shift={(0,-.25)}]150:1) arc[start angle=150,end angle=160,radius=1];
    \draw ([shift={(0,-0.25)}]20:1) arc[start angle=20,end angle=80,radius=1];
  \end{scope}
}

% ============================================================
% 17. 烧杯边缘 + 刻度 (RebordBecher)
% ============================================================
\newcommand{\tkzBeakerRim}{%
  \chemlabUpdateLineWidths
  \draw[line width=\chemlabLineWidth] (0,2.55) ellipse (1.1 and .1);
  % 主刻度 : 粗细为细刻度线的 1.5 倍 (原 pst-labo: 0.03cm = 1.5 × 0.02cm)
  % 用 \pgfmathsetmacro 预算尺寸, 避免 \dimexpr 在 pgfkeys 中的兼容问题
  \pgfmathsetlengthmacro{\chemlab@majorGrad}{1.5*\chemlabGraduationLineWidth}%
  \foreach \nDiv in {0.5,1.0,...,2.0} {%
    \draw[line width=\chemlab@majorGrad] (-.75,\nDiv) -- (-.4,\nDiv);}
  % 细分刻度
  \foreach \nSubDiv in {0.5,0.6,...,2.0} {%
    \draw[line width=\chemlabGraduationLineWidth] (-0.75,\nSubDiv) -- (-.5,\nSubDiv);}
}

% ============================================================
% 18. 烧杯 (Becher)
% ============================================================
% 注 : 原 PSTricks 使用 \psline[linearc=0.5](-1,2.5)(-1,0)(1,0)(1,2.5)
% TikZ 的 rounded corners 与 PSTricks linearc 语义略有不同,
% 因此改用显式圆弧精确模拟 PSTricks linearc=0.5 的行为 :
% 左下角圆弧 : 圆心 (-0.5,0.5), 半径 0.5, 从 180° 到 270°
% 右下角圆弧 : 圆心 (0.5,0.5),  半径 0.5, 从 270° 到 360°
\newcommand{\tkzBeakerBody}{%
  \chemlabUpdateLineWidths
  \draw[line width=\chemlabLineWidth,line cap=round,line join=round]
       (-1,2.5) -- (-1,0.5)
       -- ([shift={(-0.5,0.5)}]180:0.5) arc[start angle=180,end angle=270,radius=0.5]
       -- (0.5,0)
       -- ([shift={(0.5,0.5)}]270:0.5) arc[start angle=270,end angle=360,radius=0.5]
       -- (1,2.5);
}

% ============================================================
% 19. 锥形瓶 (Erlenmeyer)
% ============================================================
% 注 : 原 PSTricks 使用 \psline[linearc=0.3](-0.5,3.5)(-0.5,2.5)(-1.5,0)(1.5,0)(0.5,2.5)(0.5,3.5)
% TikZ 的 rounded corners 与 PSTricks linearc 语义不同,
% 因此改用显式圆弧精确模拟 linearc=0.3 (由 scripts/erlen_arcs5.py 计算) :
%   肩部 (凹转角, 内侧角 21.8°) : 半径 0.3 太大无法相切 (切点距 1.558 > 颈部线段长 1.0), 保留尖角
%   左下 (凸转角, 内侧角 68.2°) : center=(-1.0569,0.3), r=0.3, 158.20° -> 270° (ccw)
%   右下 (凸转角, 内侧角 68.2°) : center=(1.0569,0.3),  r=0.3, -90° -> 21.80° (ccw)
\newcommand{\tkzErlenmeyerBody}{%
  \chemlabUpdateLineWidths
  \draw[line width=\chemlabLineWidth,line cap=round,line join=round]
       (-0.5,3.5) -- (-0.5,2.5)
       -- (-1.3354,0.4114)
       -- ([shift={(-1.0569,0.3)}]158.20:0.3) arc[start angle=158.20,end angle=270,radius=0.3]
       -- (1.0569,0)
       -- ([shift={(1.0569,0.3)}]-90:0.3) arc[start angle=-90,end angle=21.80,radius=0.3]
       -- (0.5,2.5) -- (0.5,3.5);
}

% ============================================================
% 20. 试剂瓶 (Flacon)
% ============================================================
\newcommand{\tkzReagentBottleBody}{%
  \chemlabUpdateLineWidths\begin{scope}[line width=\chemlabLineWidth,line cap=round,line join=round]
    % 左侧瓶口上角小弧 (psarcn 90->0, 顺时针)
    \draw ([shift={(-0.6,3.4)}]90:0.1) arc[start angle=90,end angle=0,radius=0.1];
    % 左侧颈部直线
    \draw (-0.5,3.4) -- (-0.5,3.1);
    % 左侧瓶口下角凹弧 (psarcn 0->-90, 顺时针; TikZ end<start 自动顺时针)
    \draw ([shift={(-0.6,3.1)}]0:0.1) arc[start angle=0,end angle=-90,radius=0.1];
    % 左侧肩部大凸弧 (psarc 90->180, 逆时针)
    \draw ([shift={(-0.6,2.6)}]90:0.4) arc[start angle=90,end angle=180,radius=0.4];
    % 左侧瓶身直线
    \draw (-1,2.6) -- (-1,0.2);
    % 左侧底部圆角 (psarc 180->270, 逆时针)
    \draw ([shift={(-.8,0.2)}]180:0.2) arc[start angle=180,end angle=270,radius=0.2];
    % 底部
    \draw (-.8,0) -- (.8,0);
    % 右侧底部圆角 (psarc -90->0, 逆时针)
    \draw ([shift={(.8,0.2)}]-90:0.2) arc[start angle=-90,end angle=0,radius=0.2];
    % 右侧瓶身直线
    \draw (1,0.2) -- (1,2.6);
    % 右侧肩部大凸弧 (psarc 0->90, 逆时针)
    \draw ([shift={(0.6,2.6)}]0:0.4) arc[start angle=0,end angle=90,radius=0.4];
    % 右侧瓶口下角凹弧 (psarcn -90->180, 顺时针短弧; TikZ 需 end=-180 避免 270° 长弧)
    \draw ([shift={(0.6,3.1)}]-90:0.1) arc[start angle=-90,end angle=-180,radius=0.1];
    % 右侧颈部直线
    \draw (0.5,3.1) -- (0.5,3.4);
    % 右侧瓶口上角小弧 (psarcn 180->90, 顺时针)
    \draw ([shift={(0.6,3.4)}]180:0.1) arc[start angle=180,end angle=90,radius=0.1];
  \end{scope}
}

% ============================================================
% 21. 玻璃杯 (Verre) - F. Bréal
% ============================================================
\newcommand{\tkzWatchGlassBody}{%
  \chemlabUpdateLineWidths\begin{scope}[line width=\chemlabLineWidth,line cap=round,line join=round]
    % 合并左斜线 + 底部圆弧 + 右斜线为单条路径, 平滑转角
    \draw (-1,3.5) -- (-0.5,1)
          arc[start angle=180,end angle=360,radius=0.5]
          -- (1,3.5);
  \end{scope}
}

\newcommand{\tkzWatchGlassBase}{%
  \chemlabUpdateLineWidths\begin{scope}[line width=\chemlabLineWidth,rounded corners=0.03cm,line cap=round,line join=round]
    \draw (-0.5,1) -- (-0.34,0.2) -- (-1.2,0) -- (1.2,0) -- (0.34,0.2) -- (0.5,1);
  \end{scope}
}

% ============================================================
% 22. 漏斗 (Entonnoir)
% ============================================================
\newcommand{\tkzFunnelOutline}{%
  \begin{scope}[line cap=round,line join=round]
    \draw (-0.07,0) -- (-0.07,1) -- (-1,2.5);
    \draw (1,2.5) -- (0.07,1) -- (0.07,0.1);
  \end{scope}
}

\newcommand{\tkzFunnelBody}{%
  \begin{scope}
    \clip (-0.07,0) -- (-0.07,1) -- (-1,2.5)
          -- (1,2.5) -- (0.07,1) -- (0.07,0.1) -- cycle;
    \fill[white] (-2,-1) rectangle (2,3);
  \end{scope}
  \begin{scope}[rounded corners=0.25cm]
    \draw (-0.07,0) -- (-0.07,1) -- (-1,2.5);
    \draw (1,2.5) -- (0.07,1) -- (0.07,0.1);
  \end{scope}
  \draw (0,2.5) ellipse (1 and 0.1);
}

% ============================================================
% 23. 球形冷凝器 (Réfrigérant à boulles)
% ============================================================
\newcommand{\tkzCondenserStem}{%
  % 整个 TigeBoulles 作为一个连续路径 (匹配 pst-labo \pscustom)
  % pst-labo: \pscustom[linestyle=none]{...\fill} 填充 + 单独描边 (无闭合线)
  % TikZ: 分开 \fill (带 cycle 闭合填充) 和 \draw (无 cycle, 顶部敞口)
  % 顶部是敞口的 — 两端弧之间没有连线 (匹配 pst-labo)
  %
  % 填充: 带闭合线 (fill 需要 cycle)
  \fill[white]
    % 顶部圆弧与颈部 - 左侧
    ({-0.35+0.1*cos(90)},{7.5+0.1*sin(90)}) arc[start angle=90,end angle=0,radius=0.1]
    -- (-.25,7.5) -- (-.25,6.6)
    -- ({-.15+0.1*cos(180)},{6.6+0.1*sin(180)}) arc[start angle=180,end angle=270,radius=0.1]
    -- ({-.15+0.1*cos(90)},{6.4+0.1*sin(90)}) arc[start angle=90,end angle=0,radius=0.1]
    -- (-.05,6.4) -- (-0.05,6)
    % 左侧 7 个球状突起 (中心 x=-0.05, 半径 0.25, arc 90→270 向左凸)
    (-0.05,6.0) arc[start angle=90,end angle=270,radius=0.25] -- (-0.05,5.50) -- (-0.05,5.25)
    arc[start angle=90,end angle=270,radius=0.25] -- (-0.05,4.75) -- (-0.05,4.50)
    arc[start angle=90,end angle=270,radius=0.25] -- (-0.05,4.00) -- (-0.05,3.75)
    arc[start angle=90,end angle=270,radius=0.25] -- (-0.05,3.25) -- (-0.05,3.00)
    arc[start angle=90,end angle=270,radius=0.25] -- (-0.05,2.50) -- (-0.05,2.25)
    arc[start angle=90,end angle=270,radius=0.25] -- (-0.05,1.75) -- (-0.05,1.50)
    arc[start angle=90,end angle=270,radius=0.25] -- (-0.05,1.00) -- (-0.05,0.75)
    % 左侧底部到下管口
    -- (-0.05,-1.5) -- (0.05,-1.5) -- (0.05,1)
    % 右侧 7 个球状突起 (中心 x=0.05, 半径 0.25, arc 270→450 向右凸)
    arc[start angle=270,end angle=450,radius=0.25] -- (0.05,1.50) -- (0.05,1.75)
    arc[start angle=270,end angle=450,radius=0.25] -- (0.05,2.25) -- (0.05,2.50)
    arc[start angle=270,end angle=450,radius=0.25] -- (0.05,3.00) -- (0.05,3.25)
    arc[start angle=270,end angle=450,radius=0.25] -- (0.05,3.75) -- (0.05,4.00)
    arc[start angle=270,end angle=450,radius=0.25] -- (0.05,4.50) -- (0.05,4.75)
    arc[start angle=270,end angle=450,radius=0.25] -- (0.05,5.25) -- (0.05,5.50)
    arc[start angle=270,end angle=450,radius=0.25] -- (0.05,6.00) -- (0.05,6.40)
    % 顶部圆弧与颈部 - 右侧
    -- ({0.15+0.1*cos(180)},{6.4+0.1*sin(180)}) arc[start angle=180,end angle=90,radius=0.1]
    -- ({0.15+0.1*cos(-90)},{6.6+0.1*sin(-90)}) arc[start angle=-90,end angle=0,radius=0.1]
    -- (0.25,6.6) -- (0.25,7.5)
    -- ({0.35+0.1*cos(180)},{7.5+0.1*sin(180)}) arc[start angle=180,end angle=90,radius=0.1]
    -- cycle;
  %
  % 描边: 底部敞口 (匹配 pst-labo: 左右侧分别 \psline, 无底部连线)
  \draw
    % 顶部圆弧与颈部 - 左侧
    ({-0.35+0.1*cos(90)},{7.5+0.1*sin(90)}) arc[start angle=90,end angle=0,radius=0.1]
    -- (-.25,7.5) -- (-.25,6.6)
    -- ({-.15+0.1*cos(180)},{6.6+0.1*sin(180)}) arc[start angle=180,end angle=270,radius=0.1]
    -- ({-.15+0.1*cos(90)},{6.4+0.1*sin(90)}) arc[start angle=90,end angle=0,radius=0.1]
    -- (-.05,6.4) -- (-0.05,6)
    % 左侧 7 个球状突起
    (-0.05,6.0) arc[start angle=90,end angle=270,radius=0.25] -- (-0.05,5.50) -- (-0.05,5.25)
    arc[start angle=90,end angle=270,radius=0.25] -- (-0.05,4.75) -- (-0.05,4.50)
    arc[start angle=90,end angle=270,radius=0.25] -- (-0.05,4.00) -- (-0.05,3.75)
    arc[start angle=90,end angle=270,radius=0.25] -- (-0.05,3.25) -- (-0.05,3.00)
    arc[start angle=90,end angle=270,radius=0.25] -- (-0.05,2.50) -- (-0.05,2.25)
    arc[start angle=90,end angle=270,radius=0.25] -- (-0.05,1.75) -- (-0.05,1.50)
    arc[start angle=90,end angle=270,radius=0.25] -- (-0.05,1.00) -- (-0.05,0.75)
    -- (-0.05,-1.5);
  % 右侧 (单独 \draw, 底部不连接左侧)
  \draw
    (0.05,-1.5) -- (0.05,1)
    % 右侧 7 个球状突起
    arc[start angle=270,end angle=450,radius=0.25] -- (0.05,1.50) -- (0.05,1.75)
    arc[start angle=270,end angle=450,radius=0.25] -- (0.05,2.25) -- (0.05,2.50)
    arc[start angle=270,end angle=450,radius=0.25] -- (0.05,3.00) -- (0.05,3.25)
    arc[start angle=270,end angle=450,radius=0.25] -- (0.05,3.75) -- (0.05,4.00)
    arc[start angle=270,end angle=450,radius=0.25] -- (0.05,4.50) -- (0.05,4.75)
    arc[start angle=270,end angle=450,radius=0.25] -- (0.05,5.25) -- (0.05,5.50)
    arc[start angle=270,end angle=450,radius=0.25] -- (0.05,6.00) -- (0.05,6.40)
    -- ({0.15+0.1*cos(180)},{6.4+0.1*sin(180)}) arc[start angle=180,end angle=90,radius=0.1]
    -- ({0.15+0.1*cos(-90)},{6.6+0.1*sin(-90)}) arc[start angle=-90,end angle=0,radius=0.1]
    -- (0.25,6.6) -- (0.25,7.5)
    -- ({0.35+0.1*cos(180)},{7.5+0.1*sin(180)}) arc[start angle=180,end angle=90,radius=0.1];
  % 底部液滴
  \begin{scope}[shift={(0,-1.6)}]\tkzDrop\end{scope}
}

\newcommand{\tkzCondenserBody}{%
  % 渐变填充: pst-labo gradbegin=cyan (底), gradend=PaleOrange (顶), gradmidpoint=0
  % TikZ: bottom color=cyan (底), top color=PaleOrange (顶)
  \shade[top color=PaleOrange,bottom color=cyan]
    ([shift={(0,5.75)}]0:0.5) arc[start angle=0,end angle=180,radius=0.5]
    -- (-0.5,5.75) -- (-0.5,5.5) -- (-1.5,5.5)
    -- (-1.5,5.3) -- (-.5,5.3) -- (-.5,1.7) -- (-1.5,1.7)
    -- (-1.5,1.5) -- (-.5,1.5) -- (-0.5,1)
    -- ([shift={(0,1)}]180:0.5) arc[start angle=180,end angle=360,radius=0.5]
    -- (0.5,1) -- (0.5,5.75) -- cycle;
  % 描边
  \begin{scope}[line width=0.88pt]
    \draw ([shift={(0,5.75)}]0:0.5) arc[start angle=0,end angle=180,radius=0.5];
    \draw (-0.5,5.75) -- (-0.5,5.5) -- (-1.5,5.5);
    \draw (-0.5,5.75) -- (-0.5,5.5) -- (-1.5,5.5);
    \draw (-1.5,5.3) -- (-.5,5.3) -- (-.5,1.7) -- (-1.5,1.7);
    \draw (-1.5,1.5) -- (-.5,1.5) -- (-0.5,1);
    \draw ([shift={(0,1)}]180:0.5) arc[start angle=180,end angle=360,radius=0.5];
    \draw (0.5,1) -- (0.5,5.75);
  \end{scope}
  \draw[->,cyan] (-2.5,1.6) -- (-1.6,1.6);
  \draw[<-,PaleOrange] (-2.5,5.4) -- (-1.6,5.4);
}

\newcommand{\tkzCondenserBulbs}{%
  \begin{scope}[shift={(0,0.2)}]
    \tkzCondenserBody
    \tkzCondenserStem
  \end{scope}
  \tkzTestTubeStopper
  \fill[white,rounded corners=0.025cm]
       (-0.05,0.6) -- (-0.05,0.5) -- (0.05,0.5) -- (0.05,0.6) -- cycle;
}

% ============================================================
% 24. 加热球 (Chauffe-Ballon)
% ============================================================
\newcommand{\tkzHeaterRoundFlask}{%
  \filldraw[fill=white,draw=black] (-2,-1) rectangle (2,1);
  \draw (-1.25,0) circle (0.4);
  \begin{scope}[shift={(-1.25,0)}]
    \foreach \iA in {0,30,...,180} {\draw (\iA:0.45) -- (\iA:0.6);}
    \draw[->] (0,0) -- (60:0.3);
  \end{scope}
  % 按钮区域: pst-labo \psframe[doubleline=true] 产生双线边框
  % PSTricks doubleline: 外框在指定坐标, 内框向内偏移
  % 实测: 外框到内框偏移约 0.13 units (含线宽+间隔, at scale=0.5)
  % 外框 (在指定坐标处)
  \draw[black,line width=0.8pt] (1,-.5) rectangle (1.5,0.5);
  % 内框 (向内偏移, 填充灰色)
  \filldraw[fill=gray,draw=black,line width=0.8pt] ({1+0.14},{-.5+0.14}) rectangle ({1.5-0.14},{0.5-0.14});
  % 红色指示灯
  \fill[red] (1.25,0.7) circle (0.1);
}

% ============================================================
% 25. 回流烧瓶 (Ballon à reflux)
% ============================================================
\newcommand{\tkzRoundFlaskReflux}{%
  \begin{scope}
    \clip \tkzRoundFlaskSimplePath;   % 使用路径版本形成闭合 clip 区域
    \shade[top color=cyan,bottom color=white,draw=cyan]
      plot[smooth,tension=0.5] coordinates {%
        (-1.5,1.5) (-1.375,1.55) (-1.25,1.5) (-1.125,1.45)
        (-1,1.5) (-.875,1.55) (-.75,1.5) (-.625,1.45)
        (-.5,1.5) (-.375,1.55) (-.25,1.5) (-0.125,1.45)
        (0,1.5) (0.125,1.55) (.25,1.5) (0.375,1.45)
        (0.5,1.5) (0.625,1.55) (0.75,1.5) (0.875,1.45)
        (1,1.5) (1.125,1.55) (1.25,1.5) (1.375,1.45) (1.5,1.5)}
      -- (1.5,1.5) -- (1.5,0) -- (-1.5,0) -- (-1.5,1.5) -- cycle;
  \end{scope}
  \tkzRoundFlaskSimple
}

% ============================================================
% 26. 分液漏斗 (Ampoule à décanter / Separate Funnel)
% ============================================================
% 相 : 水相 / 有机相
\definecolor{AqueousPhaseColor}{cmyk}{0.215686,0.062745,0,0}
\definecolor{OrganicPhaseColor}{cmyk}{0.223529,0.094118,0.501961,0}

% pst-labo: \tkzSeparatoryFunnel 设 linewidth=0.04 (用户单位, 随 unit 缩放)
% TikZ: 用 \chemlabLineWidth (基础 0.053cm × scale) 近似, 匹配其他器皿线宽
\newcommand{\tkzSeparatoryFunnelBody}{%
  \chemlabUpdateLineWidths
  \draw[line width=\chemlabLineWidth,line cap=round,line join=round]
    (-0.65,15) -- (-0.55,13.75)
    .. controls (-1.55,13.54) and ( -1.92,12.7) .. (-1.92,11.65)
    .. controls (-1.92,10.6) and (-0.3,9.55) .. (-0.3,8.5)
    -- (-0.3,2)
    .. controls (-0.3,1.5) and (-0.14,1.5) .. (-0.14,1)
    -- (-0.14,-1)
    -- (0.14,-1)
    -- (0.14,1)
    .. controls (0.14,1.5) and (0.3,1.5) .. (0.3,2)
    -- (0.3,8.5)
    .. controls (0.3,9.55) and (1.92,10.6) .. (1.92,11.65)
    .. controls ( 1.92,12.7) and (1.55,13.54) .. (0.55,13.75)
    -- (0.65,15);
}

% 用于 clip 的轮廓
\newcommand{\tkzSeparatoryFunnelOutlineClip}{%
  (-0.65,15) -- (-0.55,13.75)
  .. controls (-1.55,13.54) and ( -1.92,12.7) .. (-1.92,11.65)
  .. controls (-1.92,10.6) and (-0.3,9.55) .. (-0.3,8.5)
  -- (-0.3,2)
  .. controls (-0.3,1.5) and (-0.14,1.5) .. (-0.14,1)
  -- (-0.14,-1) -- (0.14,-1) -- (0.14,1)
  .. controls (0.14,1.5) and (0.3,1.5) .. (0.3,2)
  -- (0.3,8.5)
  .. controls (0.3,9.55) and (1.92,10.6) .. (1.92,11.65)
  .. controls (1.92,12.7) and (1.55,13.54) .. (0.55,13.75)
  -- (0.65,15) -- cycle
}

% 顶盖 (Bouchon) - 继承 linewidth=0.04
\newcommand{\tkzStopper}{%
  \chemlabUpdateLineWidths
  \filldraw[fill=gray!10,draw=gray,line width=\chemlabLineWidth,line cap=round,line join=round]
    (-0.67,15.25) -- (0.67,15.25) -- (0.57,14) -- (-0.57,14) -- cycle;
  \filldraw[fill=gray!5,draw=gray,line width=\chemlabLineWidth,line cap=round,line join=round]
    ([shift={(0.555,15.365)}]-90:0.1145) arc[start angle=-90,end angle=-180,radius=0.1145]
    -- ([shift={(0,15.92)}]-45:0.67) arc[start angle=-45,end angle=225,radius=0.67]
    -- ([shift={(-0.555,15.365)}]0:0.1145) arc[start angle=0,end angle=-90,radius=0.1145]
    -- cycle;
}

% 开启的旋塞 (TapOpen)
% pst-labo: \psframe[fillstyle=solid] 默认 fillcolor=white (PSTricks 默认)
% TikZ: \fill 默认黑色, 必须显式 fill=white
\newcommand{\tkzTapOpen}{%
  \chemlabUpdateLineWidths
  \filldraw[fill=white,line width=\chemlabLineWidth] (-.2,0.5) rectangle (.2,1);
  \draw[line width=\chemlabLineWidth] (-.2,0.4) rectangle (-.3,1.1);
  \draw[line width=\chemlabLineWidth] (.2,0.4) rectangle (.3,1.1);
  \draw[line width=\chemlabLineWidth] (.3,0.6) -- (.4,0.6) -- (.5,0.3) -- (.9,0.3) -- (.9,1.2) -- (.5,1.2)
        -- (.4,0.9) -- (.3,0.9) -- cycle;
  \fill[cyan!50] (-.6,0.4) -- (-.85,0.6) -- (-.85,0.9) -- (-.6,1.1) -- (-.3,1.1) -- (-.3,0.4) -- cycle;
  \draw[line width=\chemlabLineWidth] (-0.6,0.4) -- (-0.6,1.1);
  \draw[line width=\chemlabLineWidth] (0.5,0.3) -- (0.5,1.2);
}

% 关闭的旋塞 (TapClosed)
% pst-labo: \psframe[fillstyle=solid] 默认 fillcolor=white
\newcommand{\tkzTapClosed}{%
  \chemlabUpdateLineWidths
  \filldraw[fill=white,line width=\chemlabLineWidth] (-.2,0.5) rectangle (.2,1);
  \draw[line width=\chemlabLineWidth] (-.2,0.4) rectangle (-.3,1.1);
  \draw[line width=\chemlabLineWidth] (.2,0.4) rectangle (.3,1.1);
  \draw[line width=\chemlabLineWidth] ([shift={(0.35,0.6)}]180:0.05) arc[start angle=180,end angle=90,radius=0.05]
        -- (0.4,0.65) -- (0.9,0.65) -- (0.9,0.85) -- (0.35,0.85)
        -- ([shift={(0.35,0.9)}]-90:0.05) arc[start angle=-90,end angle=-180,radius=0.05]
        -- cycle;
  \draw[line width=\chemlabLineWidth] (0.5,0.65) -- (0.5,0.85);
  \fill[cyan!50] (-.6,0.4) -- (-.85,0.6) -- (-.85,0.9) -- (-.6,1.1) -- (-.3,1.1) -- (-.3,0.4) -- cycle;
  \draw[line width=\chemlabLineWidth] (-0.6,0.4) -- (-0.6,1.1);
}

% \tkzSeparatoryFunnel  (主调用)
% 可选参数 : stopper=true/false, TapOpen=true/false, liquidLevel1, liquidLevel2
\newcommand{\tkzSeparatoryFunnel}[1][]{%
  \begingroup
  \chemlabUpdateLineWidths
  \pgfkeys{
    /tkzSeparatoryFunnel/.cd,
    stopper/.is if=TkzTestTubeStopper,
    TapOpen/.is if=TkzTapOpen,
    liquidLevel1/.store in=\tkzTestTubeLiquidLevelI,
    liquidLevel2/.store in=\tkzTestTubeLiquidLevelII,
    #1}
  % 默认值
  \ifdefined\tkzTestTubeLiquidLevelI\else\def\tkzTestTubeLiquidLevelI{11}\fi
  \ifdefined\tkzTestTubeLiquidLevelII\else\def\tkzTestTubeLiquidLevelII{12}\fi
  % 水相
  \begin{scope}
    \clip \tkzSeparatoryFunnelOutlineClip;
    \fill[AqueousPhaseColor] (-2,1) rectangle (2,\tkzTestTubeLiquidLevelI);
  \end{scope}
  % 有机相
  \begin{scope}
    \clip \tkzSeparatoryFunnelOutlineClip;
    \fill[OrganicPhaseColor] (-2,\tkzTestTubeLiquidLevelI)
          rectangle (2,\tkzTestTubeLiquidLevelII);
  \end{scope}
  \ifTkzTestTubeStopper\tkzStopper\fi
  \tkzSeparatoryFunnelBody
  \ifTkzTapOpen\tkzTapOpen\else\tkzTapClosed\fi
  % 出口管口 (pst-labo: \psellipse[fillstyle=solid,fillcolor=white] 既有白色填充又有黑色轮廓)
  \filldraw[fill=white,line width=\chemlabLineWidth] (0,-1) ellipse (0.15 and 0.25);
  % 瓶口圆点 (填充, 无需线宽)
  \fill[darkgray] (-0.675,15) circle (0.03);
  \fill[darkgray] (0.675,15) circle (0.03);
  \draw[darkgray,line width=\chemlabLineWidth] (-0.65,15) -- (0.65,15);
  % 颈与瓶身接合
  \draw[gray,line width=\chemlabLineWidth] (-0.55,13.75) -- (0.55,13.75);
  \endgroup
}

% ============================================================
% 27. 支架 (pstSupport)
% ============================================================
% pst-labo: \psset{dimen=m,linewidth=0.04} — linewidth=0.04 是用户单位, 随 unit 缩放
% TikZ: 用 \chemlabLineWidth (基础 0.053cm × scale), 与其他器皿一致
\newcommand{\tkzStand}{%
  \chemlabUpdateLineWidths
  \begin{scope}[line width=\chemlabLineWidth]
    \filldraw[fill=woodColor,draw=black,line join=round]
        (-2,-6) rectangle (5,-5);
    \filldraw[fill=woodColor,draw=black,line join=round]
        (2.75,2) rectangle (3.25,16);
    \filldraw[fill=woodColor,draw=black,line join=round]
        ([shift={(2.75,1.5)}]90:0.5) arc[start angle=90,end angle=180,radius=0.5]
        -- (2.25,1.5) -- (2.25,-5) -- (3.75,-5) -- (3.75,1.5)
        -- ([shift={(3.25,1.5)}]0:0.5) arc[start angle=0,end angle=90,radius=0.5]
        -- cycle;
    \draw (2,11.5) -- (1.55,10.5);
    \filldraw[fill=woodColor,draw=black,line join=round]
        (-2.5,11.5) -- (-2,11.5) -- (-1.55,10.5) -- (-2.5,10.5) -- cycle;
    \filldraw[fill=woodColor,draw=black,line join=round]
        (4,11.5) -- (2,11.5) -- (1.55,10.5) -- (4,10.5) -- cycle;
    \draw (-2.5,10.5) rectangle (4,11.5);
    % 旋钮: pst-labo \pscustom[style=woodColor]{...} 填充会自动闭合, 但描边只画实际路径
    % (不画对角闭合边). TikZ: \fill 自动闭合填充, \draw 画开放路径 (无 -- cycle).
    \fill[woodColor]
        (4,11.2) -- (4.5,11.2)
        -- ([shift={(5,11)}]158.2:0.5385) arc[start angle=158.2,end angle=-158.2,radius=0.5385]
        -- (4.5,10.8) -- (4,10.8) -- cycle;
    \draw
        (4,11.2) -- (4.5,11.2)
        -- ([shift={(5,11)}]158.2:0.5385) arc[start angle=158.2,end angle=-158.2,radius=0.5385]
        -- (4.5,10.8) -- (4,10.8);
    \draw (4.5,10.8) -- (4.5,11.2);
  \end{scope}
}

\endinput
    % 形状宏库 (必须先加载)
%    % ============================================================
%  tikz-chemlab.tex
%  化学实验器材 TikZ 绘图包 - 主实现 (参数系统 + 高级命令)
%  原作者 : D. Girou / C. Jorsen / M. Luque / H. Voss
%  TikZ 改写 : Aihua Wang(王爱华) (原 PSTricks 包 pst-labo)
% ============================================================
%
%  依赖 :
%    \usepackage{tikz}
%    \usetikzlibrary{arrows.meta,calc,patterns,patterns.meta,shadings}
%    \input{tikz-chemlabObj}    % 形状宏库 (必须先加载)
%    \input{tikz-chemlab}       % 本文件
%
%  使用示例 :
%    \begin{tikzpicture}
%       \tkzTestTube[glassType=beaker,liquidLevel1=50]
%    \end{tikzpicture}
%
%  参数键 (与原 pst-labo 兼容, 通过 pgfkeys 传入) :
%    burette, pHMeter, buretteColor, buretteLevel,
%    buretteReagent, beakerReagent, reagentColor, reagentLevel,
%    DistillateColor, MixtureStyle,
%    stopper, straightTube, elbowTube, uTube, uTubeLong,
%    curvedTube, curvedTubeShort, DoubleTube, condenserBulbs,
%    clamp, bubbling, gasCollection, tubeOnly, gasTubeStraight,
%    bunsenBurner, Number, label, magneticStirrer,
%    tiltAngle, liquidLevel1, liquidLevel2, liquidLevel3,
%    liquidStyle1, liquidStyle2, liquidStyle3,
%    substance, solid, glassType, TapOpen
%
%  glassType 取值 : tube / roundFlask / beaker / erlenmeyer / reagentBottle / volumetricFlask / watchGlass
% ============================================================

\NeedsTeXFormat{LaTeX2e}
\ProvidesFile{tikz-chemlab.tex}[2026/08/03 v1.0.0]

% --- 文档类检测 ----------------------------------------------
% standalone 会裁剪到 \useasboundingbox, article/book 不会
% 用于在 \tkzHeatFlask 中选择 bbox 策略
\@ifclassloaded{standalone}{\def\chemlab@isstandalone{T}}{\def\chemlab@isstandalone{F}}

% --- 全局 xshift 累加器 (用于多个 \tkzHeatFlask/\tkzRoundFlask 自动推进原点) ---
\newdimen\chemlab@xshift
\chemlab@xshift=0pt
% --- 首次调用标志 (用于在首个 \tkzHeatFlask 设置全局 clip) ---
\newif\ifchemlab@firstcall
\chemlab@firstcalltrue

% --- 依赖加载 ------------------------------------------------
\RequirePackage{tikz}
\usetikzlibrary{arrows.meta,calc,patterns,patterns.meta,shadings,decorations.pathmorphing,decorations.pathreplacing}
\@input{tikz-chemlabObj.tex}

% --- 每个 tikzpicture 开始时自动重置 xshift 和 firstcall ---
% 这样用户无需手动调用 \chemlabresetxshift
\tikzset{every picture/.append style={execute at begin picture={%
  \chemlab@xshift=0pt\chemlab@firstcalltrue%
}}}

% --- 颜色 (原 pst-labo.tex 中定义, 与 Obj 文件中的部分重复此处保留以覆盖) ---
\definecolor{Beige}{cmyk}{0,0,0.1,0.04}
\definecolor{LightGray}{cmyk}{0,0,0,0.2}
\definecolor{VeryLightGray}{cmyk}{0,0,0,0.1}
\definecolor{VeryPaleOrange}{cmyk}{0,0.1,0.3,0}
\definecolor{PaleOrange}{cmyk}{0,0.2,0.4,0}
\definecolor{LightBlue1}{cmyk}{0.2,0,0,0}
\definecolor{LightBlue}{cmyk}{0.22,0.05,0,0.1}
\definecolor{Copper}{cmyk}{0,0.9,0.9,0.2}
\definecolor{Brown}{cmyk}{0,0.3,0.5,.3}

% --- 样式 (对应原 \newpsstyle) --------------------------------
% 使用 \tikzset 定义命名样式; 调用 : \fill[liquidStyle1] ...
\tikzset{
  liquidStyle1/.style={fill=cyan,draw=none},
  liquidStyle2/.style={fill=yellow,draw=none},
  liquidStyle3/.style={fill=magenta,draw=none},
  Ball3D/.style={shading=ball,ball color=cyan,draw=none},
  Champagne/.style={fill=Beige,draw=none},
  Blood/.style={fill=red,draw=none},
  Cobalt/.style={fill=blue,line width=0.2cm},
  Oil/.style={fill=yellow,draw=none},
  Vinegar/.style={fill=magenta,draw=none},
  % pst-labo Diffusion: gradbegin=rgb(0,0.1,0.95), gradend=cyan, gradmidpoint=0
  % gradmidpoint=0: 渐变中点在底部, 大部分区域为 cyan, 底部为 blue
  Diffusion/.style={bottom color={rgb,255:red,0;green,25;blue,242},top color=cyan,draw=none},
  MixtureDiffusion2/.style={bottom color=white,top color=red,draw=red},
  flameAndGrill/.style={bottom color=PaleOrange,top color=yellow,draw=none},
  YellowStripes/.style={pattern=north east lines,pattern color=yellow},
  BlueDiffusion/.style={bottom color=green,top color=cyan,draw=none},
  % 分液漏斗两相液体 (对应 pst-labo \newpsstyle{AqueousPhase}/{OrganicPhase})
  % 颜色定义在 tikz-chemlabObj.tex: AqueousPhaseColor / OrganicPhaseColor
  AqueousPhase/.style={fill=AqueousPhaseColor,draw=none},
  OrganicPhase/.style={fill=OrganicPhaseColor,draw=none},
  % 额外 : 液体通用填充
  tkz liquid/.style={draw=none},
}

% --- 布尔标志 (\newif) ----------------------------------------
\newif\ifTkzBurette
\newif\ifTkzTitrationPHMeter
\newif\ifTkzTestTubeStopper
\newif\ifTkzTestTubeStraightTube
\newif\ifTkzTestTubeElbowTube
\newif\ifTkzTestTubeUTube
\newif\ifTkzTestTubeUTubeLong
\newif\ifTkzTestTubeCurvedTube
\newif\ifTkzTestTubeCurvedTubeShort
\newif\ifTkzTestTubeDoubleTube
\newif\ifTkzTestTubeCondenserBulbs
\newif\ifTkzTestTubeClamp
\newif\ifTkzBubbling
\newif\ifTkzGasCollection
\newif\ifTkzTubeOnly
\newif\ifTkzGasTubeStraight
\newif\ifTkzBunsenBurner
\newif\ifTkzLabel
\newif\ifTkzMagneticStirrer
\newif\ifTkzTapOpen            % (用于 \tkzSeparatoryFunnel)

% --- 参数存储宏 ----------------------------------------------
\def\tkzTiltAngle{0}
\def\tkzLiquidLevelI{50}
\def\tkzLiquidLevelII{0}
\def\tkzLiquidLevelIII{0}
\def\tkzLiquidStyleI{liquidStyle1}
\def\tkzLiquidStyleII{liquidStyle2}
\def\tkzLiquidStyleIII{liquidStyle3}
\def\tkzSubstance{}
\def\tkzSolid{}
\def\tkzNumber{}
\def\tkzGlassType{0}          % 0..6
\def\tkzGlassTypeName{tube}   % 字符串名称 (供嵌套调用传递)
\def\tkzBuretteColor{VeryPaleOrange}
\def\tkzBuretteLevel{20}
\def\tkzBuretteReagent{}
\def\tkzBeakerReagent{}
\def\tkzReagentColor{PaleOrange}
\def\tkzReagentLevel{2}
\def\tkzDistillateColor{yellow}
\def\tkzMixtureStyle{BlueDiffusion}

% --- glassType 字符串 -> 数字 --------------------------------
\newcommand{\tkzGlassTypeLookup}[1]{%
  \ifnum\pdfstrcmp{#1}{tube}=0 0\else
  \ifnum\pdfstrcmp{#1}{roundFlask}=0 1\else
  \ifnum\pdfstrcmp{#1}{beaker}=0 2\else
  \ifnum\pdfstrcmp{#1}{erlenmeyer}=0 3\else
  \ifnum\pdfstrcmp{#1}{reagentBottle}=0 4\else
  \ifnum\pdfstrcmp{#1}{volumetricFlask}=0 5\else
  \ifnum\pdfstrcmp{#1}{watchGlass}=0 6\else
  0\fi\fi\fi\fi\fi\fi\fi
}

% --- 线条粗细变量 (全局, 可由用户通过 \chemlabset 修改) ---
% 四类线宽可独立控制:
%   \chemlabLineWidth            - 主轮廓线 (玻璃器皿外壁)
%   \chemlabEllipseLineWidth     - 瓶口椭圆 (erlenmeyer/reagentBottle/倒置试管, 略粗)
%   \chemlabTubeEllipseLineWidth - 试管口椭圆 (更粗, 原 pst-labo 用 0.08cm)
%   \chemlabGraduationLineWidth  - 刻度线 (较细)
% 默认值与原 pst-labo 完全一致 (在 unit=1cm 时)
% 注: 与 pst-labo 一致, 线宽会自动按当前 tikz scale 缩放
%     (例如 scale=0.5 时, 实际线宽 = 基础线宽 × 0.5)
\def\chemlabBaseLineWidth{0.053cm}
\def\chemlabBaseEllipseLineWidth{0.07cm}
\def\chemlabBaseTubeEllipseLineWidth{0.08cm}
\def\chemlabBaseGraduationLineWidth{0.02cm}

% \chemlab@scaleLength{<基础长度>} - 返回按当前 tikz 缩放因子缩放后的长度
% 用于线宽跟随 scale=N 自动缩放, 与 pst-labo 的 unit=Ncm 行为一致
% 实现: 在每条 \draw 之前由 \chemlabUpdateLineWidths 把缩放后的长度
% 存入 dimen 寄存器, 绘制时直接使用寄存器
\newdimen\chemlabLineWidth
\newdimen\chemlabEllipseLineWidth
\newdimen\chemlabTubeEllipseLineWidth
\newdimen\chemlabGraduationLineWidth
\edef\chemlab@aa{1}\edef\chemlab@bb{0}\edef\chemlab@cc{0}\edef\chemlab@dd{1}\edef\chemlab@xx{0}\edef\chemlab@yy{0}
\newcommand{\chemlabUpdateLineWidths}{%
  \pgfgettransformentries{\chemlab@aa}{\chemlab@bb}{\chemlab@cc}{\chemlab@dd}{\chemlab@xx}{\chemlab@yy}%
  \pgfmathsetlength{\chemlabLineWidth}{\chemlabBaseLineWidth * (\chemlab@aa + \chemlab@dd) / 2}%
  \pgfmathsetlength{\chemlabEllipseLineWidth}{\chemlabBaseEllipseLineWidth * (\chemlab@aa + \chemlab@dd) / 2}%
  \pgfmathsetlength{\chemlabTubeEllipseLineWidth}{\chemlabBaseTubeEllipseLineWidth * (\chemlab@aa + \chemlab@dd) / 2}%
  \pgfmathsetlength{\chemlabGraduationLineWidth}{\chemlabBaseGraduationLineWidth * (\chemlab@aa + \chemlab@dd) / 2}%
}
% 默认值 (未调用 UpdateLineWidths 时的回退)
\chemlabLineWidth=0.053cm
\chemlabEllipseLineWidth=0.07cm
\chemlabTubeEllipseLineWidth=0.08cm
\chemlabGraduationLineWidth=0.02cm

% --- 火焰参数 (全局, 可由用户通过 \chemlabset 修改) ---
% 颜色 (默认与原 pst-labo 一致)
\def\chemlabFlameOuterColor{PaleOrange}   % 外焰起始色 (inner color)
\def\chemlabFlameOuterEndColor{yellow}    % 外焰结束色 (outer color)
\def\chemlabFlameInnerColor{cyan}         % 内焰起始色 (inner color)
\def\chemlabFlameInnerEndColor{white}     % 内焰结束色 (outer color)
% 大小 (默认与原 pst-labo 一致)
% 火焰缩放因子: 1.0=原版大小, 0.5=半高, 2.0=双高
\def\chemlabFlameScale{1.0}

% 线条粗细关键字 -> 数值转换 (对应 TikZ 标准)
% thin=0.4pt, semithick=0.6pt, thick=0.8pt, very thick=1.2pt, ultra thick=1.6pt
% 用法: \chemlab@resolveLineWidth{<关键字或数值>}{<目标宏>}
% 结果: 将解析后的数值存入 <目标宏>
\def\chemlab@thin{thin}
\def\chemlab@semithick{semithick}
\def\chemlab@thick{thick}
\def\chemlab@verythick{very thick}
\def\chemlab@ultrathick{ultra thick}
\newcommand{\chemlab@resolveLineWidth}[2]{%
  \def\chemlab@tmp{#1}%
  \ifx\chemlab@tmp\chemlab@thin          \def#2{0.4pt}%
  \else\ifx\chemlab@tmp\chemlab@semithick \def#2{0.6pt}%
  \else\ifx\chemlab@tmp\chemlab@thick     \def#2{0.8pt}%
  \else\ifx\chemlab@tmp\chemlab@verythick \def#2{1.2pt}%
  \else\ifx\chemlab@tmp\chemlab@ultrathick\def#2{1.6pt}%
  \else  % 直接当作数值 (如 0.04cm, 1pt)
    \def#2{#1}%
  \fi\fi\fi\fi\fi
}

% 三个独立设置器 (供 pgfkeys 调用)
\newcommand{\chemlabSetOutlineWidth}[1]{\chemlab@resolveLineWidth{#1}{\chemlabBaseLineWidth}}
\newcommand{\chemlabSetEllipseWidth}[1]{\chemlab@resolveLineWidth{#1}{\chemlabBaseEllipseLineWidth}}
\newcommand{\chemlabSetGraduationWidth}[1]{\chemlab@resolveLineWidth{#1}{\chemlabBaseGraduationLineWidth}}
\newcommand{\chemlabSetTubeEllipseWidth}[1]{\chemlab@resolveLineWidth{#1}{\chemlabBaseTubeEllipseLineWidth}}

% 同时设置四类线宽的快捷方式
\newcommand{\chemlabSetAllLineWidth}[1]{%
  \chemlab@resolveLineWidth{#1}{\chemlabBaseLineWidth}%
  \chemlab@resolveLineWidth{#1}{\chemlabBaseEllipseLineWidth}%
  \chemlab@resolveLineWidth{#1}{\chemlabBaseTubeEllipseLineWidth}%
  \chemlab@resolveLineWidth{#1}{\chemlabBaseGraduationLineWidth}%
}

% ===== 预设样式 (一键设置四类线宽) =====
% 通过 \chemlabset{style=publication} 等使用
% 预设说明 (数值单位 pt, 1pt ≈ 0.0353cm) :
%   default      - 原 pst-labo 默认 (轮廓 1.508pt, 椭圆 1.984pt, 试管椭圆 2.268pt, 刻度 0.567pt)
%   publication  - 出版级 (细线, 适合印刷) : 轮廓 thin, 椭圆 semithick, 试管椭圆 thick, 刻度 thin
%   presentation- 演示级 (粗线, 适合投影) : 轮廓 thick, 椭圆 very thick, 试管椭圆 ultra thick, 刻度 semithick
%   schematic    - 示意图 (中等, 平衡)   : 轮廓 semithick, 椭圆 thick, 试管椭圆 very thick, 刻度 thin
%   draft        - 草稿 (极细, 快速预览) : 全部 thin
%   bold         - 加粗 (高对比)         : 轮廓 very thick, 椭圆 ultra thick, 试管椭圆 ultra thick, 刻度 thick
\def\chemlab@style@default{default}
\def\chemlab@style@publication{publication}
\def\chemlab@style@presentation{presentation}
\def\chemlab@style@schematic{schematic}
\def\chemlab@style@draft{draft}
\def\chemlab@style@bold{bold}
\newcommand{\chemlabSetStyle}[1]{%
  \def\chemlab@styletmp{#1}%
  \ifx\chemlab@styletmp\chemlab@style@default
    \def\chemlabBaseLineWidth{0.053cm}\def\chemlabBaseEllipseLineWidth{0.07cm}\def\chemlabBaseTubeEllipseLineWidth{0.08cm}\def\chemlabBaseGraduationLineWidth{0.02cm}%
  \else\ifx\chemlab@styletmp\chemlab@style@publication
    \def\chemlabBaseLineWidth{0.4pt}\def\chemlabBaseEllipseLineWidth{0.6pt}\def\chemlabBaseTubeEllipseLineWidth{0.8pt}\def\chemlabBaseGraduationLineWidth{0.4pt}%
  \else\ifx\chemlab@styletmp\chemlab@style@presentation
    \def\chemlabBaseLineWidth{0.8pt}\def\chemlabBaseEllipseLineWidth{1.2pt}\def\chemlabBaseTubeEllipseLineWidth{1.6pt}\def\chemlabBaseGraduationLineWidth{0.6pt}%
  \else\ifx\chemlab@styletmp\chemlab@style@schematic
    \def\chemlabBaseLineWidth{0.6pt}\def\chemlabBaseEllipseLineWidth{0.8pt}\def\chemlabBaseTubeEllipseLineWidth{1.2pt}\def\chemlabBaseGraduationLineWidth{0.4pt}%
  \else\ifx\chemlab@styletmp\chemlab@style@draft
    \def\chemlabBaseLineWidth{0.4pt}\def\chemlabBaseEllipseLineWidth{0.4pt}\def\chemlabBaseTubeEllipseLineWidth{0.4pt}\def\chemlabBaseGraduationLineWidth{0.4pt}%
  \else\ifx\chemlab@styletmp\chemlab@style@bold
    \def\chemlabBaseLineWidth{1.2pt}\def\chemlabBaseEllipseLineWidth{1.6pt}\def\chemlabBaseTubeEllipseLineWidth{1.6pt}\def\chemlabBaseGraduationLineWidth{0.8pt}%
  \else
    \typeout{tikz-chemlab: unknown style `#1'. Available: default, publication, presentation, schematic, draft, bold.}%
  \fi\fi\fi\fi\fi\fi
}

% 用户接口: \chemlabset{line width=thick} 或 \chemlabset{outline line width=thick, ellipse line width=0.1cm}
% 或 \chemlabset{style=publication}
\newcommand{\chemlabset}[1]{\pgfkeys{/tikz-chemlab/.cd, #1}}

% --- pgfkeys 参数族 ------------------------------------------
\pgfkeys{
  /tikz-chemlab/.cd,
  % ===== 线条粗细控制 (全局, 三类可独立设置) =====
  % line width / linewidth : 同时设置三类线宽 (轮廓+椭圆+刻度)
  % 可设为: thin / semithick / thick / very thick / ultra thick
  % 或直接给数值 (如 0.04cm, 1pt 等)
  line width/.code={\chemlabSetAllLineWidth{#1}},
  linewidth/.code={\chemlabSetAllLineWidth{#1}},
  % outline line width : 仅设置主轮廓线
  outline line width/.code={\chemlabSetOutlineWidth{#1}},
  outline width/.code={\chemlabSetOutlineWidth{#1}},
  % ellipse line width : 仅设置瓶口椭圆 (erlenmeyer/reagentBottle/倒置试管)
  ellipse line width/.code={\chemlabSetEllipseWidth{#1}},
  ellipse width/.code={\chemlabSetEllipseWidth{#1}},
  % tube ellipse line width : 仅设置试管口椭圆 (原 pst-labo 用 0.08cm, 比其他椭圆粗)
  tube ellipse line width/.code={\chemlabSetTubeEllipseWidth{#1}},
  tube ellipse width/.code={\chemlabSetTubeEllipseWidth{#1}},
  % graduation line width : 仅设置刻度线
  graduation line width/.code={\chemlabSetGraduationWidth{#1}},
  graduation width/.code={\chemlabSetGraduationWidth{#1}},
  % style : 预设样式 (一键设置三类线宽)
  % 可选: default / publication / presentation / schematic / draft / bold
  style/.code={\chemlabSetStyle{#1}},
  % ===== 火焰参数 (本生灯) =====
  % 颜色: 可用任意 TikZ 颜色名 (如 cyan, yellow, red, PaleOrange 等)
  flame outer color/.store in=\chemlabFlameOuterColor,
  flame outer end color/.store in=\chemlabFlameOuterEndColor,
  flame inner color/.store in=\chemlabFlameInnerColor,
  flame inner end color/.store in=\chemlabFlameInnerEndColor,
  % 大小: 缩放因子 (1.0=原版, 0.5=半高, 2.0=双高)
  flame scale/.store in=\chemlabFlameScale,
  burette/.is if=TkzBurette,
  pHMeter/.is if=TkzTitrationPHMeter,
  stopper/.is if=TkzTestTubeStopper,
  straightTube/.is if=TkzTestTubeStraightTube,
  elbowTube/.is if=TkzTestTubeElbowTube,
  uTube/.is if=TkzTestTubeUTube,
  uTubeLong/.is if=TkzTestTubeUTubeLong,
  curvedTube/.is if=TkzTestTubeCurvedTube,
  curvedTubeShort/.is if=TkzTestTubeCurvedTubeShort,
  doubleTube/.is if=TkzTestTubeDoubleTube,
  condenserBulbs/.is if=TkzTestTubeCondenserBulbs,
  clamp/.is if=TkzTestTubeClamp,
  noDefaults/.is if=chemlab@noDefaults,
  bubbling/.is if=TkzBubbling,
  gasCollection/.is if=TkzGasCollection,
  tubeOnly/.is if=TkzTubeOnly,
  gasTubeStraight/.is if=TkzGasTubeStraight,
  bunsenBurner/.is if=TkzBunsenBurner,
  label/.is if=TkzLabel,
  magneticStirrer/.is if=TkzMagneticStirrer,
  TapOpen/.is if=TkzTapOpen,
  tiltAngle/.store in=\tkzTiltAngle,
  liquidLevel1/.store in=\tkzLiquidLevelI,
  liquidLevel2/.store in=\tkzLiquidLevelII,
  liquidLevel3/.store in=\tkzLiquidLevelIII,
  liquidStyle1/.store in=\tkzLiquidStyleI,
  liquidStyle2/.store in=\tkzLiquidStyleII,
  liquidStyle3/.store in=\tkzLiquidStyleIII,
  substance/.store in=\tkzSubstance,
  solid/.store in=\tkzSolid,
  Number/.store in=\tkzNumber,
  glassType/.code={\edef\tkzGlassType{\tkzGlassTypeLookup{#1}}\def\tkzGlassTypeName{#1}},
  buretteColor/.store in=\tkzBuretteColor,
  buretteLevel/.store in=\tkzBuretteLevel,
  buretteReagent/.store in=\tkzBuretteReagent,
  beakerReagent/.store in=\tkzBeakerReagent,
  reagentColor/.store in=\tkzReagentColor,
  reagentLevel/.store in=\tkzReagentLevel,
  DistillateColor/.store in=\tkzDistillateColor,
  MixtureStyle/.store in=\tkzMixtureStyle,
}

% --- 默认值 (对应原 \psset[pst-labo]{...}) -------------------
\newcommand{\tikzchemlabdefaults}{%
  \TkzBurettetrue
  \TkzTitrationPHMeterfalse
  \TkzTestTubeStopperfalse
  \TkzTestTubeStraightTubefalse
  \TkzTestTubeElbowTubefalse
  \TkzTestTubeUTubefalse
  \TkzTestTubeUTubeLongfalse
  \TkzTestTubeCurvedTubefalse
  \TkzTestTubeCurvedTubeShortfalse
  \TkzTestTubeDoubleTubefalse
  \TkzTestTubeCondenserBulbsfalse
  \TkzTestTubeClampfalse
  \TkzBubblingfalse
  \TkzGasCollectionfalse
  \TkzTubeOnlyfalse
  \TkzGasTubeStraightfalse
  \TkzBunsenBurnertrue
  \TkzLabelfalse
  \TkzMagneticStirrertrue
  \TkzTapOpenfalse
  \def\tkzTiltAngle{0}%
  \def\tkzLiquidLevelI{50}%
  \def\tkzLiquidLevelII{0}%
  \def\tkzLiquidLevelIII{0}%
  \def\tkzLiquidStyleI{liquidStyle1}%
  \def\tkzLiquidStyleII{liquidStyle2}%
  \def\tkzLiquidStyleIII{liquidStyle3}%
  \def\tkzSubstance{}%
  \def\tkzSolid{}%
  \def\tkzNumber{}%
  \def\tkzGlassType{0}%
  \def\tkzGlassTypeName{tube}%
  \def\tkzBuretteColor{VeryPaleOrange}%
  \def\tkzBuretteLevel{20}%
  \def\tkzBuretteReagent{}%
  \def\tkzBeakerReagent{}%
  \def\tkzReagentColor{PaleOrange}%
  \def\tkzReagentLevel{2}%
  \def\tkzDistillateColor{yellow}%
  \def\tkzMixtureStyle{BlueDiffusion}%
}

% ============================================================
% 玻璃器皿路径版本 (供 \clip 使用, 不带 \draw/\fill)
% ============================================================

% 滴定管路径 (burette corps)
\newcommand{\tkzBuretteBodyPath}{%
  (-0.5,11.5) -- (-0.25,11) -- (-0.25,3) -- (0,2.75) -- (0,0.5)
  -- (0.2,0.7) -- (0.2,2.75) -- (0.5,3) -- (0.5,11) -- (0.75,11.5) -- cycle
}

% 移液管路径 (pipette corps + 液滴)
\newcommand{\tkzPipetteBodyPath}{%
  (-0.05,0) -- (-0.15,0.2) -- (-0.15,6)
  -- (0.15,6) -- (0.15,0.2) -- (0.05,0) -- cycle
  (0,0) .. controls (0.25,-.4) and (-.25,-.4) .. (0,0)
}

% 试管路径 (TubeEssais corps)
\newcommand{\tkzTestTubeBodyPath}{%
  (-0.5,3.5) -- (-0.5,0.5)
  arc[start angle=180,end angle=360,radius=0.5]
  -- (0.5,3.5) -- cycle
}

% 烧杯路径 (Becher corps) - 显式圆弧, 与 \tkzBeakerBody 一致
\newcommand{\tkzBeakerBodyPath}{%
  (-1,2.5) -- (-1,0.5)
  -- ([shift={(-0.5,0.5)}]180:0.5) arc[start angle=180,end angle=270,radius=0.5]
  -- (0.5,0)
  -- ([shift={(0.5,0.5)}]270:0.5) arc[start angle=270,end angle=360,radius=0.5]
  -- (1,2.5) -- cycle
}

% 锥形瓶路径 (Erlenmeyer corps) - 显式圆弧, 与 \tkzErlenmeyerBody 一致
\newcommand{\tkzErlenmeyerBodyPath}{%
  (-0.5,3.5) -- (-0.5,2.5)
  -- (-1.3354,0.4114)
  -- ([shift={(-1.0569,0.3)}]158.20:0.3) arc[start angle=158.20,end angle=270,radius=0.3]
  -- (1.0569,0)
  -- ([shift={(1.0569,0.3)}]-90:0.3) arc[start angle=-90,end angle=21.80,radius=0.3]
  -- (0.5,2.5) -- (0.5,3.5) -- cycle
}

% 试剂瓶路径 (Flacon corps)
\newcommand{\tkzReagentBottleBodyPath}{%
  ([shift={(-0.6,3.4)}]90:0.1) arc[start angle=90,end angle=0,radius=0.1]
  -- (-0.5,3.4) -- (-0.5,3.1)
  -- ([shift={(-0.6,3.1)}]0:0.1) arc[start angle=0,end angle=-90,radius=0.1]
  -- ([shift={(-0.6,2.6)}]90:0.4) arc[start angle=90,end angle=180,radius=0.4]
  -- (-1,2.6) -- (-1,0.2)
  -- ([shift={(-.8,0.2)}]180:0.2) arc[start angle=180,end angle=270,radius=0.2]
  -- (-.8,0) -- (.8,0)
  -- ([shift={(.8,0.2)}]-90:0.2) arc[start angle=-90,end angle=0,radius=0.2]
  -- (1,0.2) -- (1,2.6)
  -- ([shift={(0.6,2.6)}]0:0.4) arc[start angle=0,end angle=90,radius=0.4]
  -- ([shift={(0.6,3.1)}]-90:0.1) arc[start angle=-90,end angle=-180,radius=0.1]
  -- (0.5,3.1) -- (0.5,3.4)
  -- ([shift={(0.6,3.4)}]180:0.1) arc[start angle=180,end angle=90,radius=0.1]
  -- cycle
}

% 容量瓶路径 (FioleJauge)
\newcommand{\tkzVolumetricFlaskPath}{%
  (-0.21,3.5) -- (-0.1,3.5) -- (-0.1,1.5)
  -- (-0.75,0) -- (0.75,0) -- (0.1,1.5) -- (0.1,3.5) -- (0.21,3.5) -- cycle
}

% 玻璃杯路径 (Verre corps)
\newcommand{\tkzWatchGlassBodyPath}{%
  (-1,3.5) -- (-0.5,1)
  arc[start angle=180,end angle=360,radius=0.5]
  -- (1,3.5) -- cycle
}

% 量筒路径 (Eprouvette corps)
\newcommand{\tkzGraduatedCylinderBodyPath}{%
  (-0.5,5) -- (-0.5,0) -- (0.5,0) -- (0.5,4.9) -- (0.6,4.93) -- cycle
}

% 统一调度 : 根据 \tkzGlassType 选择 clip 路径
\newcommand{\tkzGlassClipPath}{%
  \ifcase\tkzGlassType
    \tkzTestTubeBodyPath%        0 tube
  \or
    \tkzRoundFlaskSimplePath%          1 roundFlask
  \or
    \tkzBeakerBodyPath%           2 beaker
  \or
    \tkzErlenmeyerBodyPath%            3 erlenmeyer
  \or
    \tkzReagentBottleBodyPath%           4 reagentBottle
  \or
    \tkzVolumetricFlaskPath%            5 volumetricFlask
  \or
    \tkzWatchGlassBodyPath%            6 watchGlass
  \else
    \tkzTestTubeBodyPath%
  \fi
}

% 各 glassType 对应的 clip 圆角值 (与 \tkzGlassDraw 中各器皿的圆角保持一致,
% 否则液体会溢出实际绘制的轮廓外)
% 通过 \edef 预先展开为纯数字, 避免在 \clip[...] 选项中嵌套 \ifcase
% 注 : beaker(2) 和 erlenmeyer(3) 的路径已用显式圆弧, 故此处用 0pt
\newcommand{\tkzGlassClipOptions}{0pt}
\newcommand{\tkzGlassClipOptionsSet}{%
  \ifcase\tkzGlassType
    \edef\tkzGlassClipOptions{0pt}%       0 tube
  \or
    \edef\tkzGlassClipOptions{0pt}%       1 roundFlask
  \or
    \edef\tkzGlassClipOptions{0pt}%       2 beaker (路径自带圆弧)
  \or
    \edef\tkzGlassClipOptions{0pt}%       3 erlenmeyer  (路径自带圆弧)
  \or
    \edef\tkzGlassClipOptions{0pt}%       4 reagentBottle
  \or
    \edef\tkzGlassClipOptions{0.1cm}%     5 volumetricFlask
  \or
    \edef\tkzGlassClipOptions{0pt}%       6 watchGlass
  \else
    \edef\tkzGlassClipOptions{0pt}%
  \fi
}

% 统一调度 : 根据 \tkzGlassType 绘制对应玻璃器皿轮廓
\newcommand{\tkzGlassDraw}{%
  \ifcase\tkzGlassType
    \tkzTestTubeBody%                                0 tube
  \or
    \tkzRoundFlaskSimple%                                   1 roundFlask
  \or
    \tkzBeakerBody\relax\tkzBeakerRim%              2 beaker
  \or
    \tkzErlenmeyerBody\relax\filldraw[fill=white,draw=black,line width=\chemlabEllipseLineWidth,line cap=round] (0,3.5) ellipse (0.6 and 0.1);% 3 erlenmeyer
  \or
    \tkzReagentBottleBody\relax\filldraw[fill=white,draw=black,line width=\chemlabEllipseLineWidth,line cap=round] (0,3.5) ellipse (0.6 and 0.1);% 4 reagentBottle
  \or
    \tkzVolumetricFlask%                                     5 volumetricFlask
  \or
    \tkzWatchGlassBody\relax\tkzWatchGlassBase%                  6 watchGlass
  \fi
}

% ============================================================
% 液体多边形 helper (对应 \pst@Eprouvette@Liquide / \pst@TubeEssais@Liquide)
% ============================================================
% 参数 : #1=niveau(0-100), #2=style, #3=angle, #4=液面比例高度(3 或 4.5)
\newcommand{\tkzLiquidFill}[4]{%
  \pgfmathsetmacro{\liquidY}{#4 * #1 / 100 * cos(#3)}%
  \fill[#2] (-4,\liquidY) -- (-4,-1) -- (5,-1) -- (5,\liquidY) -- cycle;
}

% 试管液体 (高度 3)
\newcommand{\tkzTestTubeLiquid}[3]{\tkzLiquidFill{#1}{#2}{#3}{3}}
% 量筒液体 (高度 4.5)
\newcommand{\tkzGraduatedCylinderLiquid}[3]{\tkzLiquidFill{#1}{#2}{#3}{4.5}}

% ============================================================
% 1. \tkzTitration  滴定装置
% ============================================================
\newcommand{\tkzTitration}[1][]{%
  % 保存全局参数 (匹配 pst-labo \psset 的全局继承行为)
  \edef\chemlab@savedGlassType{\tkzGlassType}%
  \edef\chemlab@savedGlassTypeName{\tkzGlassTypeName}%
  \edef\chemlab@savedLiquidLevelI{\tkzLiquidLevelI}%
  \tikzchemlabdefaults
  \pgfkeys{/tikz-chemlab/.cd, #1}%
  % 恢复全局参数 (如果命令选项未指定)
  \ifnum\tkzGlassType=0 \edef\tkzGlassType{\chemlab@savedGlassType}\edef\tkzGlassTypeName{\chemlab@savedGlassTypeName}\fi
  \ifnum\tkzLiquidLevelI=50\edef\tkzLiquidLevelI{\chemlab@savedLiquidLevelI}\fi
  % --- 方案 B 局部修复: 忠实复刻 pst-labo \tkzRoundFlask 的 fallback 行为 ---
  % pst-labo.tex L478-522: \tkzTitration 内部调用 \tkzRoundFlask, 其 \ifcase\psk@glassType
  %   0(tube)/1(roundFlask) -> \pst@roundFlask
  %   2(beaker)         -> \pst@Becher@Corps
  %   3(erlenmeyer)          -> \pst@Erlen@Corps
  %   4(reagentBottle)         -> \pst@Flacon@Corps
  %   5(volumetricFlask)/6(Verre)/未知 -> \pst@roundFlask (fallback, L485 注释 "roundFlask is default")
  % tikz-chemlab 的 \tkzGlassDraw 对 glassType=0 画的是 \tkzTestTubeBody (试管),
  % 与原版不一致. 此处仅在 \tkzTitration 内部把 tube/volumetricFlask/Verre/未知 重映射为 roundFlask,
  % 使 \tkzGlassDraw/\tkzGlassClipPath 走 roundFlask 分支. 不影响其他命令.
  \ifnum\tkzGlassType=0 \def\tkzGlassType{1}\def\tkzGlassTypeName{roundFlask}\fi
  \ifnum\tkzGlassType=5 \def\tkzGlassType{1}\def\tkzGlassTypeName{roundFlask}\fi
  \ifnum\tkzGlassType=6 \def\tkzGlassType{1}\def\tkzGlassTypeName{roundFlask}\fi
  \ifnum\tkzGlassType>6 \def\tkzGlassType{1}\def\tkzGlassTypeName{roundFlask}\fi
  \begin{scope}[xshift=\chemlab@xshift]
    % 边界框
    \ifTkzBurette
      \ifTkzTitrationPHMeter
        \useasboundingbox (-5,0) rectangle (3,16);
      \else
        \useasboundingbox (-2,0) rectangle (2,16);
      \fi
    \else
      \ifTkzTitrationPHMeter
        \useasboundingbox (-5,0) rectangle (3,5);
      \else
        \useasboundingbox (-2,0) rectangle (2,5);
      \fi
    \fi
    % 烧瓶 + 液体 (匹配 pst-labo \rput(0,3.5){\tkzRoundFlask})
    % pst-labo: \tkzRoundFlask 的 pspicture bbox = (-2,0)(2,4), 中心 (0,2)
    % \rput(0,3.5) 使用 ref=c (中心对齐), 所以局部原点 (瓶底) 位于世界 (0, 3.5-2) = (0, 1.5)
    % 锥形瓶底部 (y=0 local) 正好放在搅拌器台面上 (y=1.5 world)
    \begin{scope}[shift={(0,1.5)}]
      % 液体 (clip 到玻璃内部 ∩ 液体矩形)
      \begin{scope}
        \tkzGlassClipOptionsSet
        \clip[rounded corners=\tkzGlassClipOptions] \tkzGlassClipPath;
        \ifnum\tkzLiquidLevelI>0
          \pgfmathsetmacro{\chemlab@liquidClipY}{3 * \tkzLiquidLevelI / 100 * cos(\tkzTiltAngle)}%
          \clip (-4,-1.5) rectangle (5,\chemlab@liquidClipY);
        \fi
        \ifnum\tkzLiquidLevelI>0
          \tkzTestTubeLiquid{\tkzLiquidLevelI}{\tkzLiquidStyleI}{\tkzTiltAngle}
        \fi
        \tkzSubstance
        \tkzSolid
      \end{scope}
      % 玻璃器皿轮廓 (根据 glassType 选择: erlenmeyer/beaker/roundFlask 等)
      \tkzGlassDraw
    \end{scope}
    % 滴定管 + 液体
    \ifTkzBurette
      \begin{scope}[shift={(0,4.5)}]
        % 液体填充 (clip 到 burette 内部)
        \begin{scope}
          \clip \tkzBuretteBodyPath;
          \pgfmathsetmacro{\buretteY}{\tkzBuretteLevel * 0.32 + 2.68}
          \fill[\tkzBuretteColor] (-1,-1) rectangle (1,\buretteY);
        \end{scope}
        \tkzBuretteBody
        \tkzBuretteStopcock
        \tkzBuretteGraduation
      \end{scope}
    \fi
    % pH 计
    \ifTkzTitrationPHMeter
      \begin{scope}[shift={(-3,0)}]\tkzTitrationPHMeter\end{scope}
    \fi
    % 搅拌子 + 支架
    \ifTkzMagneticStirrer
      \begin{scope}[shift={(0,1.5)}]\tkzTitrationStirBar\end{scope}
      \begin{scope}[shift={(0,0)}]\tkzTitrationStand\end{scope}
    \else
      % pst-labo: \psset{linewidth=0.053}\psframe(-1.5,0)(1.5,1.5)
      % linewidth=0.053 是用户单位 (×unit), 随 scale 缩放
      \chemlabUpdateLineWidths
      \draw[line width=\chemlabLineWidth] (-1.5,0) rectangle (1.5,1.5);
    \fi
    % 反应物标签
    \ifTkzBurette
      \node[anchor=west] at (2,11) {\tkzBuretteReagent};
    \fi
    \node[anchor=west] at (2,3.5) {\tkzBeakerReagent};
  \end{scope}
  % 自动推进原点 (匹配 pst-labo \pspicture 的布局宽度)
  \ifTkzTitrationPHMeter
    \advance\chemlab@xshift by 8cm %  width = 8
  \else
    \advance\chemlab@xshift by 4cm %  width = 4
  \fi
}

% ============================================================
% 2. \tkzPipette  移液管
% ============================================================
\newcommand{\tkzPipette}[1][]{%
  \begingroup
  \tikzchemlabdefaults
  \pgfkeys{/tikz-chemlab/.cd, #1}%
  \begin{scope}
    \useasboundingbox (-1,-1) rectangle (1,6);
    % 液体
    \begin{scope}
      \clip \tkzPipetteBodyPath;
      \pgfmathsetmacro{\pipetteY}{\tkzReagentLevel / 2}
      \fill[\tkzReagentColor] (-1,-2) rectangle (1,\pipetteY);
    \end{scope}
    \tkzPipetteBody
    \tkzDrop
  \end{scope}
  \endgroup
}

% ============================================================
% 3. \tkzLargePipette  大移液管
% ============================================================
\newcommand{\tkzLargePipette}[1][]{%
  \begingroup
  \tikzchemlabdefaults
  \pgfkeys{/tikz-chemlab/.cd, #1}%
  \begin{scope}
    \useasboundingbox (-1,-1) rectangle (1,6);
    % 注 : 原文 \pst@pipette 是已注释的代码, 此处使用 bigpipette 形状
    \begin{scope}
      % bigpipette 形状比较简单, 直接 clip 一个矩形区域近似
      \clip (0,0) -- (0,1) -- (-0.23,2.5) -- (-0.23,22.5)
            arc[start angle=-80,end angle=80,radius=0.45] -- cycle;
      \pgfmathsetmacro{\pipetteY}{\tkzReagentLevel / 2}
      \fill[\tkzReagentColor] (-1,-2) rectangle (1,\pipetteY);
    \end{scope}
    \tkzLargePipetteBody
  \end{scope}
  \endgroup
}

% ============================================================
% 4. \tkzGraduatedCylinder  量筒
% ============================================================
\newcommand{\tkzGraduatedCylinder}[1][]{%
  \begingroup
  \tikzchemlabdefaults
  \pgfkeys{/tikz-chemlab/.cd, #1}%
  \begin{scope}
    % 注: pst-labo 用 \pspicture(-2,0)(2,6) 设置 bbox 用于 \rput 放置, 但内容可超出 bbox
    % 在 TikZ 中, useasboundingbox 会覆盖整个 picture 的 bbox, 不适合嵌套使用
    % 因此这里不设 useasboundingbox, 让外层命令 (如 Figure 中的 pspicture) 控制 bbox
    % standalone 编译单独的 \tkzGraduatedCylinder 时, 自动检测内容 bbox
    % 液体 (clip 到量筒内部, 旋转 tiltAngle)
    \begin{scope}[rotate=\tkzTiltAngle]
      \begin{scope}
        \clip \tkzGraduatedCylinderBodyPath;
        \ifnum\tkzLiquidLevelI>0
          \tkzGraduatedCylinderLiquid{\tkzLiquidLevelI}{\tkzLiquidStyleI}{\tkzTiltAngle}
        \fi
        \ifnum\tkzLiquidLevelII>0
          \tkzGraduatedCylinderLiquid{\tkzLiquidLevelII}{\tkzLiquidStyleII}{\tkzTiltAngle}
        \fi
        \ifnum\tkzLiquidLevelIII>0
          \tkzGraduatedCylinderLiquid{\tkzLiquidLevelIII}{\tkzLiquidStyleIII}{\tkzTiltAngle}
        \fi
        \tkzSubstance
      \end{scope}
      \tkzGraduatedCylinderBody
      \tkzGraduatedCylinderBase
      \tkzGraduatedCylinderGraduation
      \begin{scope}[shift={(0,5)}]\tkzGraduatedCylinderRim\end{scope}
    \end{scope}
  \end{scope}
  \endgroup
}

% ============================================================
% 5. \tkzDistillation  蒸馏装置
% ============================================================
\newcommand{\tkzDistillation}[1][]{%
  \begingroup
  \tikzchemlabdefaults
  \pgfkeys{/tikz-chemlab/.cd, #1}%
  \begin{scope}
    \useasboundingbox (-4,-10) rectangle (8,7);
    \begin{scope}[shift={(-3,-7)}]\tkzDistillationFractionated\end{scope}
  \end{scope}
  \endgroup
}

% ============================================================
% 6. \tkzFunnel  漏斗 + 试管
% ============================================================
\newcommand{\tkzFunnel}[1][]{%
  \begingroup
  \tikzchemlabdefaults
  \pgfkeys{/tikz-chemlab/.cd, #1}%
  % 设置 noDefaults=true, 使 \tkzTestTube 继承用户参数
  \chemlab@noDefaultstrue
  \begin{scope}
    % pst-labo: \rput(0,0){\tkzTestTube} 用 pspicture bbox CENTER 对齐
    % tikz: shift = rput_pos - bbox_center (in local coords)
    % pspicture bbox 取决于 glassType 和 tiltAngle:
    %   beaker (glassType=2), tiltAngle=0: (-0.5,0)(1,2.8), center=(0.25,1.4)
    %   tiltAngle<0: (-1,0)(3,4), center=(1,2)
    %   tiltAngle>0: (-3,0)(1,4), center=(-1,2)
    %   其他 (tiltAngle=0, glassType!=2): (-1.5,0)(1.5,4), center=(0,2)
    \ifnum\tkzTiltAngle=0
      \ifnum\tkzGlassType=2
        % beaker, 直立: center=(0.25,1.4)
        \begin{scope}[shift={(-0.25,-1.4)}]\tkzTestTube[noDefaults=true]\end{scope}
      \else
        % roundFlask/reagentBottle/erlenmeyer, 直立: center=(0,2)
        \begin{scope}[shift={(0,-2)}]\tkzTestTube[noDefaults=true]\end{scope}
      \fi
    \else
      \ifnum\tkzTiltAngle>0
        % 倾斜>0: center=(-1,2)
        \begin{scope}[shift={(1,-2)}]\tkzTestTube[noDefaults=true]\end{scope}
      \else
        % 倾斜<0: center=(1,2)
        \begin{scope}[shift={(-1,-2)}]\tkzTestTube[noDefaults=true]\end{scope}
      \fi
    \fi
    % 漏斗: pst-labo \rput(0,1.5){\pst@Entonnoir@Corps}
    % \pst@Entonnoir@Corps 路径 bbox: x=[-1,1], y=[0,2.6] (含椭圆顶 2.5+0.1)
    % bbox center = (0, 1.3)
    % \rput(0,1.5) ref=c: center 置于 (0,1.5)
    % shift = rput_pos - center = (0, 1.5-1.3) = (0, 0.2)
    % 即管尖 (路径原点 0,0) 位于 (0, 0.2)
    \begin{scope}[shift={(0,0.2)}]\tkzFunnelBody\end{scope}
  \end{scope}
  \chemlab@noDefaultsfalse
  \endgroup
}

% ============================================================
% 7. \tkzTestTube  试管 (核心命令)
% ============================================================
\newcommand{\tkzTestTube}[1][]{%
  \begingroup
  \ifchemlab@noDefaults\else\tikzchemlabdefaults\fi
  \pgfkeys{/tikz-chemlab/.cd, #1}%
  \begin{scope}
    % 注: pst-labo 用 \pspicture(bbox) 设置 bbox 用于 \rput 放置, 但内容可超出 bbox
    % 在 TikZ 中, useasboundingbox 会覆盖整个 picture 的 bbox, 不适合嵌套使用
    % 因此这里不设 useasboundingbox, 让外层命令 (如 \tkzHeatTube) 控制 bbox
    % standalone 编译单独的 \tkzTestTube 时, 自动检测内容 bbox
    % 液体 (clip 到旋转的玻璃器皿内部, 但液面保持水平)
    % 原版 pst-labo: \psclip{ \rput{\angle}{<轮廓>} \psframe<液面不旋转> }
    % pst-labo 的 \psclip 产生两次裁剪 (PostScript clip 是累积的=交集):
    %   1. 裁剪到玻璃形状 (来自 \pscustom)
    %   2. 裁剪到液体矩形 (来自 \psframe(-3,-1)(5, liquid_level))
    % → 物质被裁剪到 玻璃形状 ∩ 液体矩形 = 液面以下
    % TikZ: 先旋转 clip 玻璃路径, 再反向旋转 clip 液体矩形, 然后画液体和物质
    \begin{scope}[rotate=\tkzTiltAngle]
      \begin{scope}
        \tkzGlassClipOptionsSet
        \clip[rounded corners=\tkzGlassClipOptions] \tkzGlassClipPath;
        % 液面: 在旋转的 clip 内, 但反向旋转使液面水平
        % 液面高度需乘 cos(tiltAngle) (原版 pst-labo 的行为)
        \begin{scope}[rotate=-\tkzTiltAngle]
          % 第二次裁剪: 液体矩形 (匹配 pst-labo \psframe 裁剪)
          % 只在有液体时才裁剪; 无液体时不裁剪 (物质可见到玻璃顶部)
          \ifnum\tkzLiquidLevelI>0
            \pgfmathsetmacro{\chemlab@liquidClipY}{3 * \tkzLiquidLevelI / 100 * cos(\tkzTiltAngle)}%
            \clip (-4,-1.5) rectangle (5,\chemlab@liquidClipY);
          \fi
          \ifnum\tkzLiquidLevelI>0
            \tkzTestTubeLiquid{\tkzLiquidLevelI}{\tkzLiquidStyleI}{\tkzTiltAngle}
          \fi
          \ifnum\tkzLiquidLevelII>0
            \tkzTestTubeLiquid{\tkzLiquidLevelII}{\tkzLiquidStyleII}{\tkzTiltAngle}
          \fi
          \ifnum\tkzLiquidLevelIII>0
            \tkzTestTubeLiquid{\tkzLiquidLevelIII}{\tkzLiquidStyleIII}{\tkzTiltAngle}
          \fi
          \tkzSubstance
          \tkzSolid
        \end{scope}
      \end{scope}
      % 标签
      \ifTkzLabel
        \node[fill=VeryLightGray,draw=none,inner sep=2pt] at (0,2) {\tkzNumber};
      \fi
      % 玻璃器皿轮廓
      \tkzGlassDraw
      % 顶部附件 (straightTube / elbowTube / uTube / stopper)
      \ifTkzTestTubeStraightTube
        \begin{scope}[shift={(0,3.5)}]\tkzTestTubeStraightTube\end{scope}
      \else
        \ifTkzTestTubeElbowTube
          \begin{scope}[shift={(0,3.5)}]\tkzTestTubeElbowTube\end{scope}
        \else
          \ifTkzTestTubeUTube
            \begin{scope}[shift={(0,3.5)}]\tkzTestTubeUTube\end{scope}
          \else
            \ifTkzTestTubeStopper
              \begin{scope}[shift={(0,3.5)}]\tkzTestTubeStopper\end{scope}
            \fi
          \fi
        \fi
      \fi
      % 试管夹
      \ifTkzTestTubeClamp
        \begin{scope}[shift={(0,3.2)}]\tkzTestTubeClamp\end{scope}
      \fi
    \end{scope}
  \end{scope}
  \endgroup
}

% ============================================================
% 8. \tkzHeatTube  加热试管
% ============================================================
% pstChauffageTube: 保存用户参数, 传递给 \tkzTestTube
% pst-labo 用 \psset 全局继承; tikz-chemlab 用 \tkzTestTube[noDefaults] 跳过重置
\newif\ifchemlab@noDefaults
\newcommand{\tkzHeatTube}[1][]{%
  \begingroup
  \tikzchemlabdefaults
  \pgfkeys{/tikz-chemlab/.cd, #1}%
  \chemlab@noDefaultstrue
  \begin{scope}
    % pst-labo bbox: tubeOnly=(-2,-5)(5,4), 其他=(-3,-5)(10,4)
    % 不使用 \useasboundingbox — pst-labo 的 pspicture 不裁剪内容,
    % 内容可以超出 bbox (如 elbowTube 弯管延伸到 x>5)
    % 让 TikZ 自动检测完整内容 bbox
    \ifTkzBubbling
      % 上方试管/烧瓶: 继承用户设置的 glassType, 倾斜-60°, 有弯管
      \edef\chemlab@tmpGlassType{\tkzGlassTypeName}%
      % pst-labo: \rput(2.5,2.5){\tkzTestTube[elbowTube=true,tiltAngle=-60]}
      % elbowTube=true → bbox (-1,0)(6,4) center (2.5,2)
      % → 局部原点 (瓶底) 世界 (2.5-2.5, 2.5-2) = (0, 0.5)
      \begin{scope}[shift={(0,0.5)}]
        \tkzTestTube[elbowTube=true,tiltAngle=-60,noDefaults=true]
      \end{scope}
      % 下方试管: pst-labo \rput(8.2,-3){\tkzTestTube[tiltAngle=30,...]}
      % tiltAngle>0 → bbox (-3,0)(1,4) center (-1,2)
      % → 局部原点 (试管底) 世界 (8.2-(-1), -3-2) = (9.2, -5)
      \begin{scope}[shift={(9.2,-5)}]
        \tkzTestTube[liquidLevel1=70,glassType=tube,elbowTube=false,
                       tiltAngle=30,
                       liquidStyle1=Champagne,
                       substance={\tkzChampagneBubbles[10]},clamp=false]
      \end{scope}
    \else
      % 非barbotage: 上方器皿继承 glassType
      \edef\chemlab@tmpGlassType{\tkzGlassTypeName}%
      \ifcase\tkzGlassType
        % glassType=tube
        \ifTkzTubeOnly
          % tubeOnly + tube: pst-labo \rput(1.5,2.5)
          % tiltAngle<0 → bbox (-1,0)(3,4) center (1,2) → 世界 (0.5, 0.5)
          \begin{scope}[shift={(0.5,0.5)}]\tkzTestTube[tiltAngle=-60,noDefaults=true]\end{scope}
        \else
          \ifTkzTestTubeElbowTube
            % elbowTube: pst-labo \rput(2,2.5)
            % elbowTube=true → bbox (-1,0)(6,4) center (2.5,2) → 世界 (-0.5, 0.5)
            \begin{scope}[shift={(-0.5,0.5)}]\tkzTestTube[tiltAngle=-60,noDefaults=true]\end{scope}
          \else
            % 无tubeCoude: pst-labo \rput(1,2.5)
            % tiltAngle<0 → bbox (-1,0)(3,4) center (1,2) → 世界 (0, 0.5)
            \begin{scope}[shift={(0,0.5)}]\tkzTestTube[tiltAngle=-60,noDefaults=true]\end{scope}
          \fi
        \fi
      \or
        % glassType=roundFlask (无论 tubeOnly 与否, pst-labo 都是 \rput(2,2.5))
        % tiltAngle<0 → bbox (-1,0)(3,4) center (1,2) → 局部原点世界 (2-1, 2.5-2) = (1, 0.5)
        \begin{scope}[shift={(1,0.5)}]\tkzTestTube[tiltAngle=-60,noDefaults=true]\end{scope}
      \else
        % 其他 glassType: pst-labo \rput(1,2.5) → 世界 (0, 0.5)
        \begin{scope}[shift={(0,0.5)}]\tkzTestTube[tiltAngle=-60,noDefaults=true]\end{scope}
      \fi
    \fi
    \ifTkzBunsenBurner
      % pst-labo: \rput(1,-5){\pst@BecBunsen} (无 pspicture, 实测原点对齐)
      % 灯底位于世界 (1, -5)
      \begin{scope}[shift={(1,-5)}]\tkzBunsenBurner\end{scope}
    \fi
  \end{scope}
  \chemlab@noDefaultsfalse
  \endgroup
}

% ============================================================
% 9. \tkzHeatFlask  加热烧瓶
% ============================================================
\newcommand{\tkzHeatFlask}[1][]{%
  \tikzchemlabdefaults
  \pgfkeys{/tikz-chemlab/.cd, #1}%
  % pst-labo: 每个 \tkzHeatFlask 有独立 pspicture, 不跨实例裁剪
  % tikz-chemlab: 不设全局 clip, 让所有内容自由绘制
  % pst-labo: \tkzRoundFlask 中 glassType=0 (tube) 也画 \pst@roundFlask
  % tikz-chemlab: 强制 glassType=0 → 1 (roundFlask)
  \begin{scope}[xshift=\chemlab@xshift]
    % 不设 \useasboundingbox, 让 TikZ 自动检测完整内容 bbox
    % 物体间相对位置由 shift 控制, 与 pst-labo 一致
    \ifTkzBunsenBurner\def\cnth{-5}\else\def\cnth{1}\fi
    % 气体回收装置
    \ifTkzGasCollection
      \ifTkzTestTubeCurvedTube
        \begin{scope}[shift={(8,-5)}]
          \tkzCrystallizingDish\relax\tkzTestTubeRack
          \begin{scope}[shift={(0,0.7)}]\tkzInvertedTube\relax\tkzBubbles\end{scope}
        \end{scope}
      \else\ifTkzTestTubeCurvedTubeShort
        \TkzBunsenBurnerfalse
        \begin{scope}[shift={(6,1)}]
          \tkzCrystallizingDish\relax\tkzTestTubeRack
          \begin{scope}[shift={(0,0.7)}]\tkzInvertedTube\relax\tkzBubbles\end{scope}
        \end{scope}
      \else\ifTkzTestTubeDoubleTube
        \TkzBunsenBurnerfalse
        \begin{scope}[shift={(6.2,1)}]
          \tkzCrystallizingDish\relax\tkzTestTubeRack
          \begin{scope}[shift={(0,0.7)}]\tkzInvertedTube\relax\tkzBubbles\end{scope}
        \end{scope}
      \fi\fi\fi
    \fi
    % 烧瓶 + 液体 + 导管
    % pst-labo: \rput(1,3){\tkzRoundFlask} 默认使用 bbox 中心对齐 (ref=c)
    % \tkzRoundFlask 的 pspicture bbox = (-2,0)(2,4), 中心 (0,2)
    % 所以局部原点 (瓶底) 实际位于世界 (1, 3-2) = (1, 1)
    \begin{scope}[shift={(1,1)}]
      % 液体 (clip 到烧瓶内部 ∩ 液体矩形, 匹配 pst-labo \psclip 双重裁剪)
      \begin{scope}
        \tkzGlassClipOptionsSet
        \clip[rounded corners=\tkzGlassClipOptions] \tkzGlassClipPath;
        % 第二次裁剪: 液体矩形 (匹配 pst-labo \psframe 裁剪)
        % 只在有液体时才裁剪; 无液体时不裁剪 (物质可见到玻璃顶部)
        \ifnum\tkzLiquidLevelI>0
          \pgfmathsetmacro{\chemlab@liquidClipY}{3 * \tkzLiquidLevelI / 100 * cos(\tkzTiltAngle)}%
          \clip (-4,-1.5) rectangle (5,\chemlab@liquidClipY);
        \fi
        \ifnum\tkzLiquidLevelI>0
          \tkzTestTubeLiquid{\tkzLiquidLevelI}{\tkzLiquidStyleI}{\tkzTiltAngle}
        \fi
        \tkzSubstance
        \tkzSolid
      \end{scope}
      % 玻璃器皿轮廓 (根据 glassType 选择: roundFlask/erlenmeyer/beaker 等)
      \tkzGlassDraw
      % 导管 (在瓶口 y=3.5 处) — 匹配 pst-labo \tkzRoundFlask 的全部导管类型
      \ifTkzTestTubeStraightTube
        \begin{scope}[shift={(0,3.5)}]\tkzTestTubeStraightTube\end{scope}
      \else\ifTkzTestTubeElbowTube
        \begin{scope}[shift={(0,3.5)}]\tkzTestTubeElbowTube\end{scope}
      \else\ifTkzTestTubeUTube
        \begin{scope}[shift={(0,3.5)}]\tkzTestTubeUTube\end{scope}
      \else\ifTkzTestTubeUTubeLong
        \begin{scope}[shift={(0,3.5)}]\tkzTestTubeUTubeLong\end{scope}
      \else\ifTkzTestTubeCurvedTube
        \begin{scope}[shift={(0,3.5)}]\tkzCurvedTube\end{scope}
      \else\ifTkzTestTubeCurvedTubeShort
        \begin{scope}[shift={(0,3.5)}]\tkzCurvedTubeShort\end{scope}
      \else\ifTkzTestTubeStopper
        \begin{scope}[shift={(0,3.5)}]\tkzTestTubeStopper\end{scope}
      \else\ifTkzTestTubeDoubleTube
        \begin{scope}[shift={(0,3.5)}]\tkzTestTubeDoubleTube\end{scope}
      \fi\fi\fi\fi\fi\fi\fi\fi
    \end{scope}
    \ifTkzBubbling
      % pst-labo: \rput(5.3,-3){\tkzTestTube[...]}, tiltAngle=0 默认, glassType=tube
      % bbox (-1.5,0)(1.5,4) center (0,2) → 局部原点 (试管底) 世界 (5.3, -5)
      \begin{scope}[shift={(5.3,-5)}]
        \tkzTestTube[liquidLevel1=70,glassType=tube]
      \end{scope}
    \fi
    \ifTkzBunsenBurner
      % pst-labo: \rput(1,-5){\pst@BecBunsenGrille} (无 pspicture, 实测原点对齐)
      \begin{scope}[shift={(1,-5)}]\tkzBunsenBurnerGrill\end{scope}
    \fi
  \end{scope}
  % 自动推进原点 (匹配 pst-labo \pspicture 的布局宽度)
  % 使用全局 dimen 累加, scope 的 xshift 会正确应用
  % pst-labo bbox 宽度: curvedTubeShort/DoubleTube = 11, 其他 = 6
  \ifTkzTestTubeUTubeLong
    % uTubeLong: 不推进
  \else\ifTkzTestTubeUTube
    % uTube: 不推进
  \else\ifTkzTestTubeCurvedTubeShort
    \advance\chemlab@xshift by 11cm %  bbox (-2,0)(9,8), width = 11
  \else\ifTkzTestTubeDoubleTube
    \advance\chemlab@xshift by 11cm %  bbox (-2,0)(9,8), width = 11
  \else
    \advance\chemlab@xshift by 6cm %  bbox (-2,\cnth)(4,6), width = 6
  \fi\fi\fi\fi
}

% ============================================================
% 10. \tkzRoundFlask  烧瓶 (高级命令)
% ============================================================
\newcommand{\tkzRoundFlask}[1][]{%
  \tikzchemlabdefaults
  \pgfkeys{/tikz-chemlab/.cd, #1}%
  % pst-labo: \tkzRoundFlask 中 glassType=0 (tube) 也画 \pst@roundFlask (case 0 = roundFlask)
  % tikz-chemlab: 强制 glassType=0 → 1 (roundFlask), 避免画试管轮廓
  \begin{scope}[xshift=\chemlab@xshift]
    % 不设 \useasboundingbox, 让 TikZ 自动检测完整内容 bbox
    % bbox 坐标使用局部坐标系, 会随 scale 自动缩放
    % 物体间相对位置由 shift 控制, 与 pst-labo 一致
    % 液体
    \begin{scope}[rotate=\tkzTiltAngle]
      \begin{scope}
        \tkzGlassClipOptionsSet
        \clip[rounded corners=\tkzGlassClipOptions] \tkzGlassClipPath;
        % 第二次裁剪: 液体矩形 (匹配 pst-labo \psframe 裁剪)
        % 只在有液体时才裁剪; 无液体时不裁剪 (物质可见到玻璃顶部)
        \ifnum\tkzLiquidLevelI>0
          \pgfmathsetmacro{\chemlab@liquidClipY}{3 * \tkzLiquidLevelI / 100 * cos(\tkzTiltAngle)}%
          \clip (-4,-1.5) rectangle (5,\chemlab@liquidClipY);
        \fi
        \ifnum\tkzLiquidLevelI>0
          \tkzTestTubeLiquid{\tkzLiquidLevelI}{\tkzLiquidStyleI}{\tkzTiltAngle}
        \fi
        \ifnum\tkzLiquidLevelII>0
          \tkzTestTubeLiquid{\tkzLiquidLevelII}{\tkzLiquidStyleII}{\tkzTiltAngle}
        \fi
        \ifnum\tkzLiquidLevelIII>0
          \tkzTestTubeLiquid{\tkzLiquidLevelIII}{\tkzLiquidStyleIII}{\tkzTiltAngle}
        \fi
        \tkzSubstance
        \tkzSolid
      \end{scope}
      % 标签
      \ifTkzLabel
        \node[fill=VeryLightGray,draw=none,inner sep=2pt] at (0,1.5) {\tkzNumber};
      \fi
      % 玻璃器皿轮廓
      \tkzGlassDraw
      % 顶部附件 (按优先级)
      \ifTkzTestTubeStraightTube
        \begin{scope}[shift={(0,3.5)}]\tkzTestTubeStraightTube\end{scope}
      \else\ifTkzTestTubeElbowTube
        \begin{scope}[shift={(0,3.5)}]\tkzTestTubeElbowTube\end{scope}
      \else\ifTkzTestTubeUTube
        \begin{scope}[shift={(0,3.5)}]\tkzTestTubeUTube\end{scope}
      \else\ifTkzTestTubeUTubeLong
        \begin{scope}[shift={(0,3.5)}]\tkzTestTubeUTubeLong\end{scope}
      \else\ifTkzTestTubeCurvedTube
        \begin{scope}[shift={(0,3.5)}]\tkzCurvedTube\end{scope}
      \else\ifTkzTestTubeCurvedTubeShort
        \begin{scope}[shift={(0,3.5)}]\tkzCurvedTubeShort\end{scope}
      \else\ifTkzTestTubeStopper
        \begin{scope}[shift={(0,3.5)}]\tkzTestTubeStopper\end{scope}
      \else\ifTkzTestTubeDoubleTube
        \begin{scope}[shift={(0,3.5)}]\tkzTestTubeDoubleTube\end{scope}
      \else\ifTkzTestTubeCondenserBulbs
        \tkzRoundFlaskReflux
        \begin{scope}[shift={(0,3.5)}]\tkzCondenserBulbs\end{scope}
        \tkzHeaterRoundFlask
      \fi\fi\fi\fi\fi\fi\fi\fi\fi
    \end{scope}
  \end{scope}
  % 自动推进原点 (匹配 pst-labo \pspicture 的布局宽度)
  % pst-labo bbox: condenserBulbs = (-2,0)(2,12), 其他 = (-2,0)(2,4), width = 4
  \advance\chemlab@xshift by 4cm %  bbox width = 4
}

% ============================================================
% 11. \tkzChampagneBubbles  香槟气泡
% ============================================================
% 原版 pst-labo 源码:
%   \setrandim{\BubbleX}{-1pt}{1pt}    -- BulleX ∈ [-1, 1] (纯数字, cm)
%   \setrandim{\BubbleY}{0pt}{2pt}     -- BulleY ∈ [0, 2] (纯数字)
%   Y 坐标 = \pointless\BubbleY * liquidLevel1/100  (乘以液位百分比)
%   半径 = {\BubbleY} (带 pt 单位, 1pt = 0.0353 cm)
%   \pst@dimA 分级放大: BulleY<1pt → ×3, 1pt≤BulleY<2pt → ×2
%
% 实测原版渲染图:
%   气泡为黑色边空心圆 (\pscircle, fillstyle=none)
%   半径 0.05~0.175 单位 (比源码计算 0~0.071 大, 可能是渲染缩放)
%   气泡团中心 (世界): (7.89, -0.975)
\newcommand{\tkzChampagneBubbles}[1][25]{%
  \foreach \iBubble in {1,...,#1}{%
    % BulleX ∈ [-1, 1] (纯数字, cm)
    \pgfmathsetmacro{\BubbleX}{-1 + 2*rnd}%
    % BulleY ∈ [0, 2] (纯数字)
    \pgfmathsetmacro{\BubbleY}{2*rnd}%
    % Y 坐标 = BulleY * liquidLevel1/100 (匹配源码逻辑)
    \pgfmathsetmacro{\BubbleYpos}{\BubbleY * \tkzLiquidLevelI / 100}%
    % dimA 分级放大 (匹配源码逻辑)
    \pgfmathsetmacro{\BubbleRand}{rnd}%
    \ifdim\BubbleRand pt<0.5pt
      % BulleY < 1pt (50%): dimA = BulleY * 3
      \pgfmathsetmacro{\BubbleR}{\BubbleY * 0.0875}%
    \else
      % 1pt <= BulleY < 2pt (50%): dimA = BulleY * 2
      \pgfmathsetmacro{\BubbleR}{\BubbleY * 0.0875}%
    \fi
    \begin{scope}[rotate=\tkzTiltAngle]
      \draw[draw=black, line width=0.6pt] (\BubbleX,\BubbleYpos) circle (\BubbleR);
    \end{scope}
  }%
}

% ============================================================
% 12. \tkzBubbles  倒置管中气泡
% ============================================================
\newcommand{\tkzBubbles}[1][25]{%
  \pgfgettransformentries{\clab@aa}{\clab@bb}{\clab@cc}{\clab@dd}{\clab@xx}{\clab@yy}%
  \pgfmathsetmacro{\clab@scale}{(\clab@aa + \clab@dd) / 2}%
  \foreach \iBubble in {1,...,#1}{%
    % pst-labo: \setrandim{BulleX}{-0.1pt}{0.1pt} → \pointless 去掉pt, 数值直接用作局部坐标
    \pgfmathsetmacro{\BubbleX}{-0.1 + 0.2*rnd}%
    \pgfmathsetmacro{\BubbleY}{0.15 + 2.85*rnd}%
    % pst-labo: dimh = BulleY * (3 if <1pt, 2 if <2pt, 1 if >=2pt)
    % dotscale = dimh, dot radius = dotscale * dotsize/2
    % dotsize default = 2pt + 2*linewidth = 2pt + 2*0.8pt = 3.6pt → radius = dimh * 1.8pt
    % 转换为局部坐标: local = pt / (72 * scale)
    \pgfmathparse{\BubbleY<1?3:(\BubbleY<2?2:1)}%
    \let\BubbleScale\pgfmathresult
    \pgfmathsetmacro{\BubbleDimh}{\BubbleY * \BubbleScale}%
    \pgfmathsetmacro{\BubbleR}{1.8 * \BubbleDimh / (72 * \clab@scale)}%
    % Y position = BulleY * 1.5 (局部坐标, 匹配 pst-labo \pointless\BubbleY 1.5 mul)
    \pgfmathsetmacro{\BubbleYpos}{\BubbleY * 1.5}%
    \begin{scope}[rotate=\tkzTiltAngle]
      % pst-labo: dotstyle=o (open circle = hollow, black outline)
      \draw[draw=black, line width=0.4pt, fill=white] (\BubbleX,\BubbleYpos) circle (\BubbleR);
    \end{scope}
  }%
}

% ============================================================
% 13. \tkzFilaments  絮状物
% ============================================================
\newcommand{\tkzFilaments}[2][5]{%
  \pgfgettransformentries{\fil@aa}{\fil@bb}{\fil@cc}{\fil@dd}{\fil@xx}{\fil@yy}%
  \pgfmathsetmacro{\fil@scale}{(\fil@aa + \fil@dd) / 2}%
  \foreach \iFil in {1,...,#1}{%
    % pst-labo: \setrandim{\pst@FilamentWidth}{0pt}{0.05pt} → \pointless = [0, 0.05)
    % linewidth=\pointless\pst@FilamentWidth 解释为 N × unit (PSTricks 规则)
    % 实际线宽 = [0, 0.05) × scale cm (scale 对应 unit)
    \pgfmathsetmacro{\fw}{0.05 * \fil@scale * rnd}%
    \pgfmathsetmacro{\myRand}{rnd}%  ONE random value per filament (匹配 pst-labo MyRand)
    \begin{scope}[rotate=\tkzTiltAngle]
      % pst-labo \pscurve 通过4个点的平滑曲线 (不是贝塞尔):
      %   P1: (MyRand-1, 2*MyRand-1)     = ([-1,0), [-1,1))
      %   P2: (MyRand+0.2, MyRand+2)     = ([0.2,1.2), [2,3))
      %   P3: (MyRand, MyRand-1)         = ([0,1), [-1,0))
      %   P4: (MyRand+0.5, MyRand+4)     = ([0.5,1.5), [4,5))
      \draw[line width=\fw cm,#2] plot[smooth,tension=0.5] coordinates {
        ({\myRand-1},{2*\myRand-1})
        ({\myRand+0.2},{\myRand+2})
        ({\myRand},{\myRand-1})
        ({\myRand+0.5},{\myRand+4})
      };
    \end{scope}
  }%
}

% ============================================================
% 14. \tkzIronNail  铁钉
% ============================================================
% 单个铁钉形状 (单位 0.5)
\newcommand{\tkzNails}{%
  % pst-labo 实际 SLW=0.71135pt=0.025cm (从 xdv 提取)
  % \psarc[unit=0.5,linewidth=0.1] 的 \psline 继承同一 linewidth
  \draw[scale=0.5,line width=0.025cm,blue]
    ([shift={(0,0.5)}]190:0.4) arc[start angle=190,end angle=260,radius=0.4];
  \draw[scale=0.5,line width=0.025cm,blue] (-0.28,0.2) -- (0.5,1);
}

\newcommand{\tkzIronNail}[1][60]{%
  \pgfgettransformentries{\clab@aa}{\clab@bb}{\clab@cc}{\clab@dd}{\clab@xx}{\clab@yy}%
  \pgfmathsetmacro{\clab@scale}{(\clab@aa + \clab@dd) / 2}%
  \foreach \iNail in {1,...,#1}{%
    % pst-labo: \setrandim{ClouX}{-1cm}{1cm} → 物理坐标, 需除以 scale 转换为局部坐标
    \pgfmathsetmacro{\NailX}{(-1 + 2*rnd) / \clab@scale}%
    \pgfmathsetmacro{\NailY}{(-0.2 + 0.6*rnd) / \clab@scale}%
    \pgfmathsetmacro{\NailAngle}{-90 + 180*rnd}%
    \begin{scope}[rotate=\tkzTiltAngle]
      \begin{scope}[shift={(\NailX,\NailY)},rotate=\NailAngle]
        \tkzNails
      \end{scope}
    \end{scope}
  }%
}

% ============================================================
% 15. \tkzCopperTurnings  铜屑
% ============================================================
% 单个铜屑形状 (单位 0.5)
\newcommand{\tkzCopper}{%
  % pst-labo 实际 SLW=0.28458pt=0.01cm (从 xdv 提取)
  \draw[scale=0.5,line width=0.01cm,Copper]
    (0.25,0.25) .. controls (0,0) and (-0.25,0.25) .. (0.25,0.75)
    .. controls (0,0.5) and (-0.25,0.75) .. (0.25,1.25)
    .. controls (0,1) and (-0.25,1.25) .. (0.25,1.75);
}

\newcommand{\tkzCopperTurnings}[1][30]{%
  \pgfgettransformentries{\clab@aa}{\clab@bb}{\clab@cc}{\clab@dd}{\clab@xx}{\clab@yy}%
  \pgfmathsetmacro{\clab@scale}{(\clab@aa + \clab@dd) / 2}%
  \foreach \iT in {1,...,#1}{%
    \pgfmathsetmacro{\TX}{(-1 + 2*rnd) / \clab@scale}%
    \pgfmathsetmacro{\TY}{(-0.2 + 0.5*rnd) / \clab@scale}%
    \pgfmathsetmacro{\TAngle}{-180 + 360*rnd}%
    \begin{scope}[rotate=\tkzTiltAngle]
      \begin{scope}[shift={(\TX,\TY)},rotate=\TAngle]
        \tkzCopper
      \end{scope}
    \end{scope}
  }%
}

% ============================================================
% 16. \tkzZincGranules  锌粒
% ============================================================
% 单个锌粒形状 (单位 0.5)
\newcommand{\tkzZinc}{%
  \filldraw[scale=0.5,fill=LightGray,draw=black]
    plot[smooth,tension=0.5] coordinates {(0,0)(-0.25,0.25)(0,0.5)
      (0.25,0.75)(0.5,0.5)(0.15,0.4)(0,0)} -- cycle;
}

\newcommand{\tkzZincGranules}[1][25]{%
  \pgfgettransformentries{\clab@aa}{\clab@bb}{\clab@cc}{\clab@dd}{\clab@xx}{\clab@yy}%
  \pgfmathsetmacro{\clab@scale}{(\clab@aa + \clab@dd) / 2}%
  \foreach \iG in {1,...,#1}{%
    \pgfmathsetmacro{\GX}{(-1 + 2*rnd) / \clab@scale}%
    \pgfmathsetmacro{\GY}{(-0.2 + 0.5*rnd) / \clab@scale}%
    \pgfmathsetmacro{\GAngle}{-180 + 360*rnd}%
    \begin{scope}[rotate=\tkzTiltAngle]
      \begin{scope}[shift={(\GX,\GY)},rotate=\GAngle]
        \tkzZinc
      \end{scope}
    \end{scope}
  }%
}

% ============================================================
% 17. \tkzBalls  3D 球状物
% ============================================================
\newcommand{\tkzBalls}[1][50]{%
  \foreach \iB in {1,...,#1}{%
    \pgfmathsetmacro{\BX}{-1.5 + 3*rnd}%
    \pgfmathsetmacro{\BY}{3*rnd}%
    \pgfmathsetmacro{\BR}{0.01 + 0.1*rnd}%
    % 按 Y 位置调整半径
    \pgfmathparse{\BY>2?3:(\BY>1?2:1)}%
    \let\Bscale\pgfmathresult
    \pgfmathsetmacro{\BRscaled}{\BR * \Bscale}%
    \begin{scope}[rotate=\tkzTiltAngle]
      % pst-labo Ball3D: GradientCircle=true
      % gradbegin=CMYK(0,0.1,0.95,0.1) → 实际渲染为 RGB(0,0.1,0.95)=蓝色 (固定颜色)
      % gradend=white
      % 径向渐变: 中心=white, 边缘=蓝色
      \definecolor{billeBlue}{rgb}{0,0.1,0.95}%
      \shade[inner color=white, outer color=billeBlue] (\BX,\BY) circle (\BRscaled);
    \end{scope}
  }%
}

% ============================================================
% 18. \tkzBubblesThreeD  2D 气泡
% ============================================================
\newcommand{\tkzBubblesThreeD}[2][50]{%
  \foreach \iB in {1,...,#1}{%
    \pgfmathsetmacro{\BX}{-1.5 + 3*rnd}%
    \pgfmathsetmacro{\BY}{3*rnd}%
    \pgfmathsetmacro{\BR}{0.01 + 0.1*rnd}%
    \pgfmathparse{\BY>2?3:(\BY>1?2:1)}%
    \let\Bscale\pgfmathresult
    \pgfmathsetmacro{\BRscaled}{\BR * \Bscale}%
    \begin{scope}[rotate=\tkzTiltAngle]
      \fill[#2,line width=0.2pt] (\BX,\BY) circle (\BRscaled);
    \end{scope}
  }%
}

\endinput
       % 本文件
%
%  使用示例 :
%    \begin{tikzpicture}
%       \tkzTestTube[glassType=beaker,liquidLevel1=50]
%    \end{tikzpicture}
%
%  参数键 (与原 pst-labo 兼容, 通过 pgfkeys 传入) :
%    burette, pHMeter, buretteColor, buretteLevel,
%    buretteReagent, beakerReagent, reagentColor, reagentLevel,
%    DistillateColor, MixtureStyle,
%    stopper, straightTube, elbowTube, uTube, uTubeLong,
%    curvedTube, curvedTubeShort, DoubleTube, condenserBulbs,
%    clamp, bubbling, gasCollection, tubeOnly, gasTubeStraight,
%    bunsenBurner, Number, label, magneticStirrer,
%    tiltAngle, liquidLevel1, liquidLevel2, liquidLevel3,
%    liquidStyle1, liquidStyle2, liquidStyle3,
%    substance, solid, glassType, TapOpen
%
%  glassType 取值 : tube / roundFlask / beaker / erlenmeyer / reagentBottle / volumetricFlask / watchGlass
% ============================================================

\NeedsTeXFormat{LaTeX2e}
\ProvidesFile{tikz-chemlab.tex}[2026/08/03 v1.0.0]

% --- 文档类检测 ----------------------------------------------
% standalone 会裁剪到 \useasboundingbox, article/book 不会
% 用于在 \tkzHeatFlask 中选择 bbox 策略
\@ifclassloaded{standalone}{\def\chemlab@isstandalone{T}}{\def\chemlab@isstandalone{F}}

% --- 全局 xshift 累加器 (用于多个 \tkzHeatFlask/\tkzRoundFlask 自动推进原点) ---
\newdimen\chemlab@xshift
\chemlab@xshift=0pt
% --- 首次调用标志 (用于在首个 \tkzHeatFlask 设置全局 clip) ---
\newif\ifchemlab@firstcall
\chemlab@firstcalltrue

% --- 依赖加载 ------------------------------------------------
\RequirePackage{tikz}
\usetikzlibrary{arrows.meta,calc,patterns,patterns.meta,shadings,decorations.pathmorphing,decorations.pathreplacing}
\@input{tikz-chemlabObj.tex}

% --- 每个 tikzpicture 开始时自动重置 xshift 和 firstcall ---
% 这样用户无需手动调用 \chemlabresetxshift
\tikzset{every picture/.append style={execute at begin picture={%
  \chemlab@xshift=0pt\chemlab@firstcalltrue%
}}}

% --- 颜色 (原 pst-labo.tex 中定义, 与 Obj 文件中的部分重复此处保留以覆盖) ---
\definecolor{Beige}{cmyk}{0,0,0.1,0.04}
\definecolor{LightGray}{cmyk}{0,0,0,0.2}
\definecolor{VeryLightGray}{cmyk}{0,0,0,0.1}
\definecolor{VeryPaleOrange}{cmyk}{0,0.1,0.3,0}
\definecolor{PaleOrange}{cmyk}{0,0.2,0.4,0}
\definecolor{LightBlue1}{cmyk}{0.2,0,0,0}
\definecolor{LightBlue}{cmyk}{0.22,0.05,0,0.1}
\definecolor{Copper}{cmyk}{0,0.9,0.9,0.2}
\definecolor{Brown}{cmyk}{0,0.3,0.5,.3}

% --- 样式 (对应原 \newpsstyle) --------------------------------
% 使用 \tikzset 定义命名样式; 调用 : \fill[liquidStyle1] ...
\tikzset{
  liquidStyle1/.style={fill=cyan,draw=none},
  liquidStyle2/.style={fill=yellow,draw=none},
  liquidStyle3/.style={fill=magenta,draw=none},
  Ball3D/.style={shading=ball,ball color=cyan,draw=none},
  Champagne/.style={fill=Beige,draw=none},
  Blood/.style={fill=red,draw=none},
  Cobalt/.style={fill=blue,line width=0.2cm},
  Oil/.style={fill=yellow,draw=none},
  Vinegar/.style={fill=magenta,draw=none},
  % pst-labo Diffusion: gradbegin=rgb(0,0.1,0.95), gradend=cyan, gradmidpoint=0
  % gradmidpoint=0: 渐变中点在底部, 大部分区域为 cyan, 底部为 blue
  Diffusion/.style={bottom color={rgb,255:red,0;green,25;blue,242},top color=cyan,draw=none},
  MixtureDiffusion2/.style={bottom color=white,top color=red,draw=red},
  flameAndGrill/.style={bottom color=PaleOrange,top color=yellow,draw=none},
  YellowStripes/.style={pattern=north east lines,pattern color=yellow},
  BlueDiffusion/.style={bottom color=green,top color=cyan,draw=none},
  % 分液漏斗两相液体 (对应 pst-labo \newpsstyle{AqueousPhase}/{OrganicPhase})
  % 颜色定义在 tikz-chemlabObj.tex: AqueousPhaseColor / OrganicPhaseColor
  AqueousPhase/.style={fill=AqueousPhaseColor,draw=none},
  OrganicPhase/.style={fill=OrganicPhaseColor,draw=none},
  % 额外 : 液体通用填充
  tkz liquid/.style={draw=none},
}

% --- 布尔标志 (\newif) ----------------------------------------
\newif\ifTkzBurette
\newif\ifTkzTitrationPHMeter
\newif\ifTkzTestTubeStopper
\newif\ifTkzTestTubeStraightTube
\newif\ifTkzTestTubeElbowTube
\newif\ifTkzTestTubeUTube
\newif\ifTkzTestTubeUTubeLong
\newif\ifTkzTestTubeCurvedTube
\newif\ifTkzTestTubeCurvedTubeShort
\newif\ifTkzTestTubeDoubleTube
\newif\ifTkzTestTubeCondenserBulbs
\newif\ifTkzTestTubeClamp
\newif\ifTkzBubbling
\newif\ifTkzGasCollection
\newif\ifTkzTubeOnly
\newif\ifTkzGasTubeStraight
\newif\ifTkzBunsenBurner
\newif\ifTkzLabel
\newif\ifTkzMagneticStirrer
\newif\ifTkzTapOpen            % (用于 \tkzSeparatoryFunnel)

% --- 参数存储宏 ----------------------------------------------
\def\tkzTiltAngle{0}
\def\tkzLiquidLevelI{50}
\def\tkzLiquidLevelII{0}
\def\tkzLiquidLevelIII{0}
\def\tkzLiquidStyleI{liquidStyle1}
\def\tkzLiquidStyleII{liquidStyle2}
\def\tkzLiquidStyleIII{liquidStyle3}
\def\tkzSubstance{}
\def\tkzSolid{}
\def\tkzNumber{}
\def\tkzGlassType{0}          % 0..6
\def\tkzGlassTypeName{tube}   % 字符串名称 (供嵌套调用传递)
\def\tkzBuretteColor{VeryPaleOrange}
\def\tkzBuretteLevel{20}
\def\tkzBuretteReagent{}
\def\tkzBeakerReagent{}
\def\tkzReagentColor{PaleOrange}
\def\tkzReagentLevel{2}
\def\tkzDistillateColor{yellow}
\def\tkzMixtureStyle{BlueDiffusion}

% --- glassType 字符串 -> 数字 --------------------------------
\newcommand{\tkzGlassTypeLookup}[1]{%
  \ifnum\pdfstrcmp{#1}{tube}=0 0\else
  \ifnum\pdfstrcmp{#1}{roundFlask}=0 1\else
  \ifnum\pdfstrcmp{#1}{beaker}=0 2\else
  \ifnum\pdfstrcmp{#1}{erlenmeyer}=0 3\else
  \ifnum\pdfstrcmp{#1}{reagentBottle}=0 4\else
  \ifnum\pdfstrcmp{#1}{volumetricFlask}=0 5\else
  \ifnum\pdfstrcmp{#1}{watchGlass}=0 6\else
  0\fi\fi\fi\fi\fi\fi\fi
}

% --- 线条粗细变量 (全局, 可由用户通过 \chemlabset 修改) ---
% 四类线宽可独立控制:
%   \chemlabLineWidth            - 主轮廓线 (玻璃器皿外壁)
%   \chemlabEllipseLineWidth     - 瓶口椭圆 (erlenmeyer/reagentBottle/倒置试管, 略粗)
%   \chemlabTubeEllipseLineWidth - 试管口椭圆 (更粗, 原 pst-labo 用 0.08cm)
%   \chemlabGraduationLineWidth  - 刻度线 (较细)
% 默认值与原 pst-labo 完全一致 (在 unit=1cm 时)
% 注: 与 pst-labo 一致, 线宽会自动按当前 tikz scale 缩放
%     (例如 scale=0.5 时, 实际线宽 = 基础线宽 × 0.5)
\def\chemlabBaseLineWidth{0.053cm}
\def\chemlabBaseEllipseLineWidth{0.07cm}
\def\chemlabBaseTubeEllipseLineWidth{0.08cm}
\def\chemlabBaseGraduationLineWidth{0.02cm}

% \chemlab@scaleLength{<基础长度>} - 返回按当前 tikz 缩放因子缩放后的长度
% 用于线宽跟随 scale=N 自动缩放, 与 pst-labo 的 unit=Ncm 行为一致
% 实现: 在每条 \draw 之前由 \chemlabUpdateLineWidths 把缩放后的长度
% 存入 dimen 寄存器, 绘制时直接使用寄存器
\newdimen\chemlabLineWidth
\newdimen\chemlabEllipseLineWidth
\newdimen\chemlabTubeEllipseLineWidth
\newdimen\chemlabGraduationLineWidth
\edef\chemlab@aa{1}\edef\chemlab@bb{0}\edef\chemlab@cc{0}\edef\chemlab@dd{1}\edef\chemlab@xx{0}\edef\chemlab@yy{0}
\newcommand{\chemlabUpdateLineWidths}{%
  \pgfgettransformentries{\chemlab@aa}{\chemlab@bb}{\chemlab@cc}{\chemlab@dd}{\chemlab@xx}{\chemlab@yy}%
  \pgfmathsetlength{\chemlabLineWidth}{\chemlabBaseLineWidth * (\chemlab@aa + \chemlab@dd) / 2}%
  \pgfmathsetlength{\chemlabEllipseLineWidth}{\chemlabBaseEllipseLineWidth * (\chemlab@aa + \chemlab@dd) / 2}%
  \pgfmathsetlength{\chemlabTubeEllipseLineWidth}{\chemlabBaseTubeEllipseLineWidth * (\chemlab@aa + \chemlab@dd) / 2}%
  \pgfmathsetlength{\chemlabGraduationLineWidth}{\chemlabBaseGraduationLineWidth * (\chemlab@aa + \chemlab@dd) / 2}%
}
% 默认值 (未调用 UpdateLineWidths 时的回退)
\chemlabLineWidth=0.053cm
\chemlabEllipseLineWidth=0.07cm
\chemlabTubeEllipseLineWidth=0.08cm
\chemlabGraduationLineWidth=0.02cm

% --- 火焰参数 (全局, 可由用户通过 \chemlabset 修改) ---
% 颜色 (默认与原 pst-labo 一致)
\def\chemlabFlameOuterColor{PaleOrange}   % 外焰起始色 (inner color)
\def\chemlabFlameOuterEndColor{yellow}    % 外焰结束色 (outer color)
\def\chemlabFlameInnerColor{cyan}         % 内焰起始色 (inner color)
\def\chemlabFlameInnerEndColor{white}     % 内焰结束色 (outer color)
% 大小 (默认与原 pst-labo 一致)
% 火焰缩放因子: 1.0=原版大小, 0.5=半高, 2.0=双高
\def\chemlabFlameScale{1.0}

% 线条粗细关键字 -> 数值转换 (对应 TikZ 标准)
% thin=0.4pt, semithick=0.6pt, thick=0.8pt, very thick=1.2pt, ultra thick=1.6pt
% 用法: \chemlab@resolveLineWidth{<关键字或数值>}{<目标宏>}
% 结果: 将解析后的数值存入 <目标宏>
\def\chemlab@thin{thin}
\def\chemlab@semithick{semithick}
\def\chemlab@thick{thick}
\def\chemlab@verythick{very thick}
\def\chemlab@ultrathick{ultra thick}
\newcommand{\chemlab@resolveLineWidth}[2]{%
  \def\chemlab@tmp{#1}%
  \ifx\chemlab@tmp\chemlab@thin          \def#2{0.4pt}%
  \else\ifx\chemlab@tmp\chemlab@semithick \def#2{0.6pt}%
  \else\ifx\chemlab@tmp\chemlab@thick     \def#2{0.8pt}%
  \else\ifx\chemlab@tmp\chemlab@verythick \def#2{1.2pt}%
  \else\ifx\chemlab@tmp\chemlab@ultrathick\def#2{1.6pt}%
  \else  % 直接当作数值 (如 0.04cm, 1pt)
    \def#2{#1}%
  \fi\fi\fi\fi\fi
}

% 三个独立设置器 (供 pgfkeys 调用)
\newcommand{\chemlabSetOutlineWidth}[1]{\chemlab@resolveLineWidth{#1}{\chemlabBaseLineWidth}}
\newcommand{\chemlabSetEllipseWidth}[1]{\chemlab@resolveLineWidth{#1}{\chemlabBaseEllipseLineWidth}}
\newcommand{\chemlabSetGraduationWidth}[1]{\chemlab@resolveLineWidth{#1}{\chemlabBaseGraduationLineWidth}}
\newcommand{\chemlabSetTubeEllipseWidth}[1]{\chemlab@resolveLineWidth{#1}{\chemlabBaseTubeEllipseLineWidth}}

% 同时设置四类线宽的快捷方式
\newcommand{\chemlabSetAllLineWidth}[1]{%
  \chemlab@resolveLineWidth{#1}{\chemlabBaseLineWidth}%
  \chemlab@resolveLineWidth{#1}{\chemlabBaseEllipseLineWidth}%
  \chemlab@resolveLineWidth{#1}{\chemlabBaseTubeEllipseLineWidth}%
  \chemlab@resolveLineWidth{#1}{\chemlabBaseGraduationLineWidth}%
}

% ===== 预设样式 (一键设置四类线宽) =====
% 通过 \chemlabset{style=publication} 等使用
% 预设说明 (数值单位 pt, 1pt ≈ 0.0353cm) :
%   default      - 原 pst-labo 默认 (轮廓 1.508pt, 椭圆 1.984pt, 试管椭圆 2.268pt, 刻度 0.567pt)
%   publication  - 出版级 (细线, 适合印刷) : 轮廓 thin, 椭圆 semithick, 试管椭圆 thick, 刻度 thin
%   presentation- 演示级 (粗线, 适合投影) : 轮廓 thick, 椭圆 very thick, 试管椭圆 ultra thick, 刻度 semithick
%   schematic    - 示意图 (中等, 平衡)   : 轮廓 semithick, 椭圆 thick, 试管椭圆 very thick, 刻度 thin
%   draft        - 草稿 (极细, 快速预览) : 全部 thin
%   bold         - 加粗 (高对比)         : 轮廓 very thick, 椭圆 ultra thick, 试管椭圆 ultra thick, 刻度 thick
\def\chemlab@style@default{default}
\def\chemlab@style@publication{publication}
\def\chemlab@style@presentation{presentation}
\def\chemlab@style@schematic{schematic}
\def\chemlab@style@draft{draft}
\def\chemlab@style@bold{bold}
\newcommand{\chemlabSetStyle}[1]{%
  \def\chemlab@styletmp{#1}%
  \ifx\chemlab@styletmp\chemlab@style@default
    \def\chemlabBaseLineWidth{0.053cm}\def\chemlabBaseEllipseLineWidth{0.07cm}\def\chemlabBaseTubeEllipseLineWidth{0.08cm}\def\chemlabBaseGraduationLineWidth{0.02cm}%
  \else\ifx\chemlab@styletmp\chemlab@style@publication
    \def\chemlabBaseLineWidth{0.4pt}\def\chemlabBaseEllipseLineWidth{0.6pt}\def\chemlabBaseTubeEllipseLineWidth{0.8pt}\def\chemlabBaseGraduationLineWidth{0.4pt}%
  \else\ifx\chemlab@styletmp\chemlab@style@presentation
    \def\chemlabBaseLineWidth{0.8pt}\def\chemlabBaseEllipseLineWidth{1.2pt}\def\chemlabBaseTubeEllipseLineWidth{1.6pt}\def\chemlabBaseGraduationLineWidth{0.6pt}%
  \else\ifx\chemlab@styletmp\chemlab@style@schematic
    \def\chemlabBaseLineWidth{0.6pt}\def\chemlabBaseEllipseLineWidth{0.8pt}\def\chemlabBaseTubeEllipseLineWidth{1.2pt}\def\chemlabBaseGraduationLineWidth{0.4pt}%
  \else\ifx\chemlab@styletmp\chemlab@style@draft
    \def\chemlabBaseLineWidth{0.4pt}\def\chemlabBaseEllipseLineWidth{0.4pt}\def\chemlabBaseTubeEllipseLineWidth{0.4pt}\def\chemlabBaseGraduationLineWidth{0.4pt}%
  \else\ifx\chemlab@styletmp\chemlab@style@bold
    \def\chemlabBaseLineWidth{1.2pt}\def\chemlabBaseEllipseLineWidth{1.6pt}\def\chemlabBaseTubeEllipseLineWidth{1.6pt}\def\chemlabBaseGraduationLineWidth{0.8pt}%
  \else
    \typeout{tikz-chemlab: unknown style `#1'. Available: default, publication, presentation, schematic, draft, bold.}%
  \fi\fi\fi\fi\fi\fi
}

% 用户接口: \chemlabset{line width=thick} 或 \chemlabset{outline line width=thick, ellipse line width=0.1cm}
% 或 \chemlabset{style=publication}
\newcommand{\chemlabset}[1]{\pgfkeys{/tikz-chemlab/.cd, #1}}

% --- pgfkeys 参数族 ------------------------------------------
\pgfkeys{
  /tikz-chemlab/.cd,
  % ===== 线条粗细控制 (全局, 三类可独立设置) =====
  % line width / linewidth : 同时设置三类线宽 (轮廓+椭圆+刻度)
  % 可设为: thin / semithick / thick / very thick / ultra thick
  % 或直接给数值 (如 0.04cm, 1pt 等)
  line width/.code={\chemlabSetAllLineWidth{#1}},
  linewidth/.code={\chemlabSetAllLineWidth{#1}},
  % outline line width : 仅设置主轮廓线
  outline line width/.code={\chemlabSetOutlineWidth{#1}},
  outline width/.code={\chemlabSetOutlineWidth{#1}},
  % ellipse line width : 仅设置瓶口椭圆 (erlenmeyer/reagentBottle/倒置试管)
  ellipse line width/.code={\chemlabSetEllipseWidth{#1}},
  ellipse width/.code={\chemlabSetEllipseWidth{#1}},
  % tube ellipse line width : 仅设置试管口椭圆 (原 pst-labo 用 0.08cm, 比其他椭圆粗)
  tube ellipse line width/.code={\chemlabSetTubeEllipseWidth{#1}},
  tube ellipse width/.code={\chemlabSetTubeEllipseWidth{#1}},
  % graduation line width : 仅设置刻度线
  graduation line width/.code={\chemlabSetGraduationWidth{#1}},
  graduation width/.code={\chemlabSetGraduationWidth{#1}},
  % style : 预设样式 (一键设置三类线宽)
  % 可选: default / publication / presentation / schematic / draft / bold
  style/.code={\chemlabSetStyle{#1}},
  % ===== 火焰参数 (本生灯) =====
  % 颜色: 可用任意 TikZ 颜色名 (如 cyan, yellow, red, PaleOrange 等)
  flame outer color/.store in=\chemlabFlameOuterColor,
  flame outer end color/.store in=\chemlabFlameOuterEndColor,
  flame inner color/.store in=\chemlabFlameInnerColor,
  flame inner end color/.store in=\chemlabFlameInnerEndColor,
  % 大小: 缩放因子 (1.0=原版, 0.5=半高, 2.0=双高)
  flame scale/.store in=\chemlabFlameScale,
  burette/.is if=TkzBurette,
  pHMeter/.is if=TkzTitrationPHMeter,
  stopper/.is if=TkzTestTubeStopper,
  straightTube/.is if=TkzTestTubeStraightTube,
  elbowTube/.is if=TkzTestTubeElbowTube,
  uTube/.is if=TkzTestTubeUTube,
  uTubeLong/.is if=TkzTestTubeUTubeLong,
  curvedTube/.is if=TkzTestTubeCurvedTube,
  curvedTubeShort/.is if=TkzTestTubeCurvedTubeShort,
  doubleTube/.is if=TkzTestTubeDoubleTube,
  condenserBulbs/.is if=TkzTestTubeCondenserBulbs,
  clamp/.is if=TkzTestTubeClamp,
  noDefaults/.is if=chemlab@noDefaults,
  bubbling/.is if=TkzBubbling,
  gasCollection/.is if=TkzGasCollection,
  tubeOnly/.is if=TkzTubeOnly,
  gasTubeStraight/.is if=TkzGasTubeStraight,
  bunsenBurner/.is if=TkzBunsenBurner,
  label/.is if=TkzLabel,
  magneticStirrer/.is if=TkzMagneticStirrer,
  TapOpen/.is if=TkzTapOpen,
  tiltAngle/.store in=\tkzTiltAngle,
  liquidLevel1/.store in=\tkzLiquidLevelI,
  liquidLevel2/.store in=\tkzLiquidLevelII,
  liquidLevel3/.store in=\tkzLiquidLevelIII,
  liquidStyle1/.store in=\tkzLiquidStyleI,
  liquidStyle2/.store in=\tkzLiquidStyleII,
  liquidStyle3/.store in=\tkzLiquidStyleIII,
  substance/.store in=\tkzSubstance,
  solid/.store in=\tkzSolid,
  Number/.store in=\tkzNumber,
  glassType/.code={\edef\tkzGlassType{\tkzGlassTypeLookup{#1}}\def\tkzGlassTypeName{#1}},
  buretteColor/.store in=\tkzBuretteColor,
  buretteLevel/.store in=\tkzBuretteLevel,
  buretteReagent/.store in=\tkzBuretteReagent,
  beakerReagent/.store in=\tkzBeakerReagent,
  reagentColor/.store in=\tkzReagentColor,
  reagentLevel/.store in=\tkzReagentLevel,
  DistillateColor/.store in=\tkzDistillateColor,
  MixtureStyle/.store in=\tkzMixtureStyle,
}

% --- 默认值 (对应原 \psset[pst-labo]{...}) -------------------
\newcommand{\tikzchemlabdefaults}{%
  \TkzBurettetrue
  \TkzTitrationPHMeterfalse
  \TkzTestTubeStopperfalse
  \TkzTestTubeStraightTubefalse
  \TkzTestTubeElbowTubefalse
  \TkzTestTubeUTubefalse
  \TkzTestTubeUTubeLongfalse
  \TkzTestTubeCurvedTubefalse
  \TkzTestTubeCurvedTubeShortfalse
  \TkzTestTubeDoubleTubefalse
  \TkzTestTubeCondenserBulbsfalse
  \TkzTestTubeClampfalse
  \TkzBubblingfalse
  \TkzGasCollectionfalse
  \TkzTubeOnlyfalse
  \TkzGasTubeStraightfalse
  \TkzBunsenBurnertrue
  \TkzLabelfalse
  \TkzMagneticStirrertrue
  \TkzTapOpenfalse
  \def\tkzTiltAngle{0}%
  \def\tkzLiquidLevelI{50}%
  \def\tkzLiquidLevelII{0}%
  \def\tkzLiquidLevelIII{0}%
  \def\tkzLiquidStyleI{liquidStyle1}%
  \def\tkzLiquidStyleII{liquidStyle2}%
  \def\tkzLiquidStyleIII{liquidStyle3}%
  \def\tkzSubstance{}%
  \def\tkzSolid{}%
  \def\tkzNumber{}%
  \def\tkzGlassType{0}%
  \def\tkzGlassTypeName{tube}%
  \def\tkzBuretteColor{VeryPaleOrange}%
  \def\tkzBuretteLevel{20}%
  \def\tkzBuretteReagent{}%
  \def\tkzBeakerReagent{}%
  \def\tkzReagentColor{PaleOrange}%
  \def\tkzReagentLevel{2}%
  \def\tkzDistillateColor{yellow}%
  \def\tkzMixtureStyle{BlueDiffusion}%
}

% ============================================================
% 玻璃器皿路径版本 (供 \clip 使用, 不带 \draw/\fill)
% ============================================================

% 滴定管路径 (burette corps)
\newcommand{\tkzBuretteBodyPath}{%
  (-0.5,11.5) -- (-0.25,11) -- (-0.25,3) -- (0,2.75) -- (0,0.5)
  -- (0.2,0.7) -- (0.2,2.75) -- (0.5,3) -- (0.5,11) -- (0.75,11.5) -- cycle
}

% 移液管路径 (pipette corps + 液滴)
\newcommand{\tkzPipetteBodyPath}{%
  (-0.05,0) -- (-0.15,0.2) -- (-0.15,6)
  -- (0.15,6) -- (0.15,0.2) -- (0.05,0) -- cycle
  (0,0) .. controls (0.25,-.4) and (-.25,-.4) .. (0,0)
}

% 试管路径 (TubeEssais corps)
\newcommand{\tkzTestTubeBodyPath}{%
  (-0.5,3.5) -- (-0.5,0.5)
  arc[start angle=180,end angle=360,radius=0.5]
  -- (0.5,3.5) -- cycle
}

% 烧杯路径 (Becher corps) - 显式圆弧, 与 \tkzBeakerBody 一致
\newcommand{\tkzBeakerBodyPath}{%
  (-1,2.5) -- (-1,0.5)
  -- ([shift={(-0.5,0.5)}]180:0.5) arc[start angle=180,end angle=270,radius=0.5]
  -- (0.5,0)
  -- ([shift={(0.5,0.5)}]270:0.5) arc[start angle=270,end angle=360,radius=0.5]
  -- (1,2.5) -- cycle
}

% 锥形瓶路径 (Erlenmeyer corps) - 显式圆弧, 与 \tkzErlenmeyerBody 一致
\newcommand{\tkzErlenmeyerBodyPath}{%
  (-0.5,3.5) -- (-0.5,2.5)
  -- (-1.3354,0.4114)
  -- ([shift={(-1.0569,0.3)}]158.20:0.3) arc[start angle=158.20,end angle=270,radius=0.3]
  -- (1.0569,0)
  -- ([shift={(1.0569,0.3)}]-90:0.3) arc[start angle=-90,end angle=21.80,radius=0.3]
  -- (0.5,2.5) -- (0.5,3.5) -- cycle
}

% 试剂瓶路径 (Flacon corps)
\newcommand{\tkzReagentBottleBodyPath}{%
  ([shift={(-0.6,3.4)}]90:0.1) arc[start angle=90,end angle=0,radius=0.1]
  -- (-0.5,3.4) -- (-0.5,3.1)
  -- ([shift={(-0.6,3.1)}]0:0.1) arc[start angle=0,end angle=-90,radius=0.1]
  -- ([shift={(-0.6,2.6)}]90:0.4) arc[start angle=90,end angle=180,radius=0.4]
  -- (-1,2.6) -- (-1,0.2)
  -- ([shift={(-.8,0.2)}]180:0.2) arc[start angle=180,end angle=270,radius=0.2]
  -- (-.8,0) -- (.8,0)
  -- ([shift={(.8,0.2)}]-90:0.2) arc[start angle=-90,end angle=0,radius=0.2]
  -- (1,0.2) -- (1,2.6)
  -- ([shift={(0.6,2.6)}]0:0.4) arc[start angle=0,end angle=90,radius=0.4]
  -- ([shift={(0.6,3.1)}]-90:0.1) arc[start angle=-90,end angle=-180,radius=0.1]
  -- (0.5,3.1) -- (0.5,3.4)
  -- ([shift={(0.6,3.4)}]180:0.1) arc[start angle=180,end angle=90,radius=0.1]
  -- cycle
}

% 容量瓶路径 (FioleJauge)
\newcommand{\tkzVolumetricFlaskPath}{%
  (-0.21,3.5) -- (-0.1,3.5) -- (-0.1,1.5)
  -- (-0.75,0) -- (0.75,0) -- (0.1,1.5) -- (0.1,3.5) -- (0.21,3.5) -- cycle
}

% 玻璃杯路径 (Verre corps)
\newcommand{\tkzWatchGlassBodyPath}{%
  (-1,3.5) -- (-0.5,1)
  arc[start angle=180,end angle=360,radius=0.5]
  -- (1,3.5) -- cycle
}

% 量筒路径 (Eprouvette corps)
\newcommand{\tkzGraduatedCylinderBodyPath}{%
  (-0.5,5) -- (-0.5,0) -- (0.5,0) -- (0.5,4.9) -- (0.6,4.93) -- cycle
}

% 统一调度 : 根据 \tkzGlassType 选择 clip 路径
\newcommand{\tkzGlassClipPath}{%
  \ifcase\tkzGlassType
    \tkzTestTubeBodyPath%        0 tube
  \or
    \tkzRoundFlaskSimplePath%          1 roundFlask
  \or
    \tkzBeakerBodyPath%           2 beaker
  \or
    \tkzErlenmeyerBodyPath%            3 erlenmeyer
  \or
    \tkzReagentBottleBodyPath%           4 reagentBottle
  \or
    \tkzVolumetricFlaskPath%            5 volumetricFlask
  \or
    \tkzWatchGlassBodyPath%            6 watchGlass
  \else
    \tkzTestTubeBodyPath%
  \fi
}

% 各 glassType 对应的 clip 圆角值 (与 \tkzGlassDraw 中各器皿的圆角保持一致,
% 否则液体会溢出实际绘制的轮廓外)
% 通过 \edef 预先展开为纯数字, 避免在 \clip[...] 选项中嵌套 \ifcase
% 注 : beaker(2) 和 erlenmeyer(3) 的路径已用显式圆弧, 故此处用 0pt
\newcommand{\tkzGlassClipOptions}{0pt}
\newcommand{\tkzGlassClipOptionsSet}{%
  \ifcase\tkzGlassType
    \edef\tkzGlassClipOptions{0pt}%       0 tube
  \or
    \edef\tkzGlassClipOptions{0pt}%       1 roundFlask
  \or
    \edef\tkzGlassClipOptions{0pt}%       2 beaker (路径自带圆弧)
  \or
    \edef\tkzGlassClipOptions{0pt}%       3 erlenmeyer  (路径自带圆弧)
  \or
    \edef\tkzGlassClipOptions{0pt}%       4 reagentBottle
  \or
    \edef\tkzGlassClipOptions{0.1cm}%     5 volumetricFlask
  \or
    \edef\tkzGlassClipOptions{0pt}%       6 watchGlass
  \else
    \edef\tkzGlassClipOptions{0pt}%
  \fi
}

% 统一调度 : 根据 \tkzGlassType 绘制对应玻璃器皿轮廓
\newcommand{\tkzGlassDraw}{%
  \ifcase\tkzGlassType
    \tkzTestTubeBody%                                0 tube
  \or
    \tkzRoundFlaskSimple%                                   1 roundFlask
  \or
    \tkzBeakerBody\relax\tkzBeakerRim%              2 beaker
  \or
    \tkzErlenmeyerBody\relax\filldraw[fill=white,draw=black,line width=\chemlabEllipseLineWidth,line cap=round] (0,3.5) ellipse (0.6 and 0.1);% 3 erlenmeyer
  \or
    \tkzReagentBottleBody\relax\filldraw[fill=white,draw=black,line width=\chemlabEllipseLineWidth,line cap=round] (0,3.5) ellipse (0.6 and 0.1);% 4 reagentBottle
  \or
    \tkzVolumetricFlask%                                     5 volumetricFlask
  \or
    \tkzWatchGlassBody\relax\tkzWatchGlassBase%                  6 watchGlass
  \fi
}

% ============================================================
% 液体多边形 helper (对应 \pst@Eprouvette@Liquide / \pst@TubeEssais@Liquide)
% ============================================================
% 参数 : #1=niveau(0-100), #2=style, #3=angle, #4=液面比例高度(3 或 4.5)
\newcommand{\tkzLiquidFill}[4]{%
  \pgfmathsetmacro{\liquidY}{#4 * #1 / 100 * cos(#3)}%
  \fill[#2] (-4,\liquidY) -- (-4,-1) -- (5,-1) -- (5,\liquidY) -- cycle;
}

% 试管液体 (高度 3)
\newcommand{\tkzTestTubeLiquid}[3]{\tkzLiquidFill{#1}{#2}{#3}{3}}
% 量筒液体 (高度 4.5)
\newcommand{\tkzGraduatedCylinderLiquid}[3]{\tkzLiquidFill{#1}{#2}{#3}{4.5}}

% ============================================================
% 1. \tkzTitration  滴定装置
% ============================================================
\newcommand{\tkzTitration}[1][]{%
  % 保存全局参数 (匹配 pst-labo \psset 的全局继承行为)
  \edef\chemlab@savedGlassType{\tkzGlassType}%
  \edef\chemlab@savedGlassTypeName{\tkzGlassTypeName}%
  \edef\chemlab@savedLiquidLevelI{\tkzLiquidLevelI}%
  \tikzchemlabdefaults
  \pgfkeys{/tikz-chemlab/.cd, #1}%
  % 恢复全局参数 (如果命令选项未指定)
  \ifnum\tkzGlassType=0 \edef\tkzGlassType{\chemlab@savedGlassType}\edef\tkzGlassTypeName{\chemlab@savedGlassTypeName}\fi
  \ifnum\tkzLiquidLevelI=50\edef\tkzLiquidLevelI{\chemlab@savedLiquidLevelI}\fi
  % --- 方案 B 局部修复: 忠实复刻 pst-labo \tkzRoundFlask 的 fallback 行为 ---
  % pst-labo.tex L478-522: \tkzTitration 内部调用 \tkzRoundFlask, 其 \ifcase\psk@glassType
  %   0(tube)/1(roundFlask) -> \pst@roundFlask
  %   2(beaker)         -> \pst@Becher@Corps
  %   3(erlenmeyer)          -> \pst@Erlen@Corps
  %   4(reagentBottle)         -> \pst@Flacon@Corps
  %   5(volumetricFlask)/6(Verre)/未知 -> \pst@roundFlask (fallback, L485 注释 "roundFlask is default")
  % tikz-chemlab 的 \tkzGlassDraw 对 glassType=0 画的是 \tkzTestTubeBody (试管),
  % 与原版不一致. 此处仅在 \tkzTitration 内部把 tube/volumetricFlask/Verre/未知 重映射为 roundFlask,
  % 使 \tkzGlassDraw/\tkzGlassClipPath 走 roundFlask 分支. 不影响其他命令.
  \ifnum\tkzGlassType=0 \def\tkzGlassType{1}\def\tkzGlassTypeName{roundFlask}\fi
  \ifnum\tkzGlassType=5 \def\tkzGlassType{1}\def\tkzGlassTypeName{roundFlask}\fi
  \ifnum\tkzGlassType=6 \def\tkzGlassType{1}\def\tkzGlassTypeName{roundFlask}\fi
  \ifnum\tkzGlassType>6 \def\tkzGlassType{1}\def\tkzGlassTypeName{roundFlask}\fi
  \begin{scope}[xshift=\chemlab@xshift]
    % 边界框
    \ifTkzBurette
      \ifTkzTitrationPHMeter
        \useasboundingbox (-5,0) rectangle (3,16);
      \else
        \useasboundingbox (-2,0) rectangle (2,16);
      \fi
    \else
      \ifTkzTitrationPHMeter
        \useasboundingbox (-5,0) rectangle (3,5);
      \else
        \useasboundingbox (-2,0) rectangle (2,5);
      \fi
    \fi
    % 烧瓶 + 液体 (匹配 pst-labo \rput(0,3.5){\tkzRoundFlask})
    % pst-labo: \tkzRoundFlask 的 pspicture bbox = (-2,0)(2,4), 中心 (0,2)
    % \rput(0,3.5) 使用 ref=c (中心对齐), 所以局部原点 (瓶底) 位于世界 (0, 3.5-2) = (0, 1.5)
    % 锥形瓶底部 (y=0 local) 正好放在搅拌器台面上 (y=1.5 world)
    \begin{scope}[shift={(0,1.5)}]
      % 液体 (clip 到玻璃内部 ∩ 液体矩形)
      \begin{scope}
        \tkzGlassClipOptionsSet
        \clip[rounded corners=\tkzGlassClipOptions] \tkzGlassClipPath;
        \ifnum\tkzLiquidLevelI>0
          \pgfmathsetmacro{\chemlab@liquidClipY}{3 * \tkzLiquidLevelI / 100 * cos(\tkzTiltAngle)}%
          \clip (-4,-1.5) rectangle (5,\chemlab@liquidClipY);
        \fi
        \ifnum\tkzLiquidLevelI>0
          \tkzTestTubeLiquid{\tkzLiquidLevelI}{\tkzLiquidStyleI}{\tkzTiltAngle}
        \fi
        \tkzSubstance
        \tkzSolid
      \end{scope}
      % 玻璃器皿轮廓 (根据 glassType 选择: erlenmeyer/beaker/roundFlask 等)
      \tkzGlassDraw
    \end{scope}
    % 滴定管 + 液体
    \ifTkzBurette
      \begin{scope}[shift={(0,4.5)}]
        % 液体填充 (clip 到 burette 内部)
        \begin{scope}
          \clip \tkzBuretteBodyPath;
          \pgfmathsetmacro{\buretteY}{\tkzBuretteLevel * 0.32 + 2.68}
          \fill[\tkzBuretteColor] (-1,-1) rectangle (1,\buretteY);
        \end{scope}
        \tkzBuretteBody
        \tkzBuretteStopcock
        \tkzBuretteGraduation
      \end{scope}
    \fi
    % pH 计
    \ifTkzTitrationPHMeter
      \begin{scope}[shift={(-3,0)}]\tkzTitrationPHMeter\end{scope}
    \fi
    % 搅拌子 + 支架
    \ifTkzMagneticStirrer
      \begin{scope}[shift={(0,1.5)}]\tkzTitrationStirBar\end{scope}
      \begin{scope}[shift={(0,0)}]\tkzTitrationStand\end{scope}
    \else
      % pst-labo: \psset{linewidth=0.053}\psframe(-1.5,0)(1.5,1.5)
      % linewidth=0.053 是用户单位 (×unit), 随 scale 缩放
      \chemlabUpdateLineWidths
      \draw[line width=\chemlabLineWidth] (-1.5,0) rectangle (1.5,1.5);
    \fi
    % 反应物标签
    \ifTkzBurette
      \node[anchor=west] at (2,11) {\tkzBuretteReagent};
    \fi
    \node[anchor=west] at (2,3.5) {\tkzBeakerReagent};
  \end{scope}
  % 自动推进原点 (匹配 pst-labo \pspicture 的布局宽度)
  \ifTkzTitrationPHMeter
    \advance\chemlab@xshift by 8cm %  width = 8
  \else
    \advance\chemlab@xshift by 4cm %  width = 4
  \fi
}

% ============================================================
% 2. \tkzPipette  移液管
% ============================================================
\newcommand{\tkzPipette}[1][]{%
  \begingroup
  \tikzchemlabdefaults
  \pgfkeys{/tikz-chemlab/.cd, #1}%
  \begin{scope}
    \useasboundingbox (-1,-1) rectangle (1,6);
    % 液体
    \begin{scope}
      \clip \tkzPipetteBodyPath;
      \pgfmathsetmacro{\pipetteY}{\tkzReagentLevel / 2}
      \fill[\tkzReagentColor] (-1,-2) rectangle (1,\pipetteY);
    \end{scope}
    \tkzPipetteBody
    \tkzDrop
  \end{scope}
  \endgroup
}

% ============================================================
% 3. \tkzLargePipette  大移液管
% ============================================================
\newcommand{\tkzLargePipette}[1][]{%
  \begingroup
  \tikzchemlabdefaults
  \pgfkeys{/tikz-chemlab/.cd, #1}%
  \begin{scope}
    \useasboundingbox (-1,-1) rectangle (1,6);
    % 注 : 原文 \pst@pipette 是已注释的代码, 此处使用 bigpipette 形状
    \begin{scope}
      % bigpipette 形状比较简单, 直接 clip 一个矩形区域近似
      \clip (0,0) -- (0,1) -- (-0.23,2.5) -- (-0.23,22.5)
            arc[start angle=-80,end angle=80,radius=0.45] -- cycle;
      \pgfmathsetmacro{\pipetteY}{\tkzReagentLevel / 2}
      \fill[\tkzReagentColor] (-1,-2) rectangle (1,\pipetteY);
    \end{scope}
    \tkzLargePipetteBody
  \end{scope}
  \endgroup
}

% ============================================================
% 4. \tkzGraduatedCylinder  量筒
% ============================================================
\newcommand{\tkzGraduatedCylinder}[1][]{%
  \begingroup
  \tikzchemlabdefaults
  \pgfkeys{/tikz-chemlab/.cd, #1}%
  \begin{scope}
    % 注: pst-labo 用 \pspicture(-2,0)(2,6) 设置 bbox 用于 \rput 放置, 但内容可超出 bbox
    % 在 TikZ 中, useasboundingbox 会覆盖整个 picture 的 bbox, 不适合嵌套使用
    % 因此这里不设 useasboundingbox, 让外层命令 (如 Figure 中的 pspicture) 控制 bbox
    % standalone 编译单独的 \tkzGraduatedCylinder 时, 自动检测内容 bbox
    % 液体 (clip 到量筒内部, 旋转 tiltAngle)
    \begin{scope}[rotate=\tkzTiltAngle]
      \begin{scope}
        \clip \tkzGraduatedCylinderBodyPath;
        \ifnum\tkzLiquidLevelI>0
          \tkzGraduatedCylinderLiquid{\tkzLiquidLevelI}{\tkzLiquidStyleI}{\tkzTiltAngle}
        \fi
        \ifnum\tkzLiquidLevelII>0
          \tkzGraduatedCylinderLiquid{\tkzLiquidLevelII}{\tkzLiquidStyleII}{\tkzTiltAngle}
        \fi
        \ifnum\tkzLiquidLevelIII>0
          \tkzGraduatedCylinderLiquid{\tkzLiquidLevelIII}{\tkzLiquidStyleIII}{\tkzTiltAngle}
        \fi
        \tkzSubstance
      \end{scope}
      \tkzGraduatedCylinderBody
      \tkzGraduatedCylinderBase
      \tkzGraduatedCylinderGraduation
      \begin{scope}[shift={(0,5)}]\tkzGraduatedCylinderRim\end{scope}
    \end{scope}
  \end{scope}
  \endgroup
}

% ============================================================
% 5. \tkzDistillation  蒸馏装置
% ============================================================
\newcommand{\tkzDistillation}[1][]{%
  \begingroup
  \tikzchemlabdefaults
  \pgfkeys{/tikz-chemlab/.cd, #1}%
  \begin{scope}
    \useasboundingbox (-4,-10) rectangle (8,7);
    \begin{scope}[shift={(-3,-7)}]\tkzDistillationFractionated\end{scope}
  \end{scope}
  \endgroup
}

% ============================================================
% 6. \tkzFunnel  漏斗 + 试管
% ============================================================
\newcommand{\tkzFunnel}[1][]{%
  \begingroup
  \tikzchemlabdefaults
  \pgfkeys{/tikz-chemlab/.cd, #1}%
  % 设置 noDefaults=true, 使 \tkzTestTube 继承用户参数
  \chemlab@noDefaultstrue
  \begin{scope}
    % pst-labo: \rput(0,0){\tkzTestTube} 用 pspicture bbox CENTER 对齐
    % tikz: shift = rput_pos - bbox_center (in local coords)
    % pspicture bbox 取决于 glassType 和 tiltAngle:
    %   beaker (glassType=2), tiltAngle=0: (-0.5,0)(1,2.8), center=(0.25,1.4)
    %   tiltAngle<0: (-1,0)(3,4), center=(1,2)
    %   tiltAngle>0: (-3,0)(1,4), center=(-1,2)
    %   其他 (tiltAngle=0, glassType!=2): (-1.5,0)(1.5,4), center=(0,2)
    \ifnum\tkzTiltAngle=0
      \ifnum\tkzGlassType=2
        % beaker, 直立: center=(0.25,1.4)
        \begin{scope}[shift={(-0.25,-1.4)}]\tkzTestTube[noDefaults=true]\end{scope}
      \else
        % roundFlask/reagentBottle/erlenmeyer, 直立: center=(0,2)
        \begin{scope}[shift={(0,-2)}]\tkzTestTube[noDefaults=true]\end{scope}
      \fi
    \else
      \ifnum\tkzTiltAngle>0
        % 倾斜>0: center=(-1,2)
        \begin{scope}[shift={(1,-2)}]\tkzTestTube[noDefaults=true]\end{scope}
      \else
        % 倾斜<0: center=(1,2)
        \begin{scope}[shift={(-1,-2)}]\tkzTestTube[noDefaults=true]\end{scope}
      \fi
    \fi
    % 漏斗: pst-labo \rput(0,1.5){\pst@Entonnoir@Corps}
    % \pst@Entonnoir@Corps 路径 bbox: x=[-1,1], y=[0,2.6] (含椭圆顶 2.5+0.1)
    % bbox center = (0, 1.3)
    % \rput(0,1.5) ref=c: center 置于 (0,1.5)
    % shift = rput_pos - center = (0, 1.5-1.3) = (0, 0.2)
    % 即管尖 (路径原点 0,0) 位于 (0, 0.2)
    \begin{scope}[shift={(0,0.2)}]\tkzFunnelBody\end{scope}
  \end{scope}
  \chemlab@noDefaultsfalse
  \endgroup
}

% ============================================================
% 7. \tkzTestTube  试管 (核心命令)
% ============================================================
\newcommand{\tkzTestTube}[1][]{%
  \begingroup
  \ifchemlab@noDefaults\else\tikzchemlabdefaults\fi
  \pgfkeys{/tikz-chemlab/.cd, #1}%
  \begin{scope}
    % 注: pst-labo 用 \pspicture(bbox) 设置 bbox 用于 \rput 放置, 但内容可超出 bbox
    % 在 TikZ 中, useasboundingbox 会覆盖整个 picture 的 bbox, 不适合嵌套使用
    % 因此这里不设 useasboundingbox, 让外层命令 (如 \tkzHeatTube) 控制 bbox
    % standalone 编译单独的 \tkzTestTube 时, 自动检测内容 bbox
    % 液体 (clip 到旋转的玻璃器皿内部, 但液面保持水平)
    % 原版 pst-labo: \psclip{ \rput{\angle}{<轮廓>} \psframe<液面不旋转> }
    % pst-labo 的 \psclip 产生两次裁剪 (PostScript clip 是累积的=交集):
    %   1. 裁剪到玻璃形状 (来自 \pscustom)
    %   2. 裁剪到液体矩形 (来自 \psframe(-3,-1)(5, liquid_level))
    % → 物质被裁剪到 玻璃形状 ∩ 液体矩形 = 液面以下
    % TikZ: 先旋转 clip 玻璃路径, 再反向旋转 clip 液体矩形, 然后画液体和物质
    \begin{scope}[rotate=\tkzTiltAngle]
      \begin{scope}
        \tkzGlassClipOptionsSet
        \clip[rounded corners=\tkzGlassClipOptions] \tkzGlassClipPath;
        % 液面: 在旋转的 clip 内, 但反向旋转使液面水平
        % 液面高度需乘 cos(tiltAngle) (原版 pst-labo 的行为)
        \begin{scope}[rotate=-\tkzTiltAngle]
          % 第二次裁剪: 液体矩形 (匹配 pst-labo \psframe 裁剪)
          % 只在有液体时才裁剪; 无液体时不裁剪 (物质可见到玻璃顶部)
          \ifnum\tkzLiquidLevelI>0
            \pgfmathsetmacro{\chemlab@liquidClipY}{3 * \tkzLiquidLevelI / 100 * cos(\tkzTiltAngle)}%
            \clip (-4,-1.5) rectangle (5,\chemlab@liquidClipY);
          \fi
          \ifnum\tkzLiquidLevelI>0
            \tkzTestTubeLiquid{\tkzLiquidLevelI}{\tkzLiquidStyleI}{\tkzTiltAngle}
          \fi
          \ifnum\tkzLiquidLevelII>0
            \tkzTestTubeLiquid{\tkzLiquidLevelII}{\tkzLiquidStyleII}{\tkzTiltAngle}
          \fi
          \ifnum\tkzLiquidLevelIII>0
            \tkzTestTubeLiquid{\tkzLiquidLevelIII}{\tkzLiquidStyleIII}{\tkzTiltAngle}
          \fi
          \tkzSubstance
          \tkzSolid
        \end{scope}
      \end{scope}
      % 标签
      \ifTkzLabel
        \node[fill=VeryLightGray,draw=none,inner sep=2pt] at (0,2) {\tkzNumber};
      \fi
      % 玻璃器皿轮廓
      \tkzGlassDraw
      % 顶部附件 (straightTube / elbowTube / uTube / stopper)
      \ifTkzTestTubeStraightTube
        \begin{scope}[shift={(0,3.5)}]\tkzTestTubeStraightTube\end{scope}
      \else
        \ifTkzTestTubeElbowTube
          \begin{scope}[shift={(0,3.5)}]\tkzTestTubeElbowTube\end{scope}
        \else
          \ifTkzTestTubeUTube
            \begin{scope}[shift={(0,3.5)}]\tkzTestTubeUTube\end{scope}
          \else
            \ifTkzTestTubeStopper
              \begin{scope}[shift={(0,3.5)}]\tkzTestTubeStopper\end{scope}
            \fi
          \fi
        \fi
      \fi
      % 试管夹
      \ifTkzTestTubeClamp
        \begin{scope}[shift={(0,3.2)}]\tkzTestTubeClamp\end{scope}
      \fi
    \end{scope}
  \end{scope}
  \endgroup
}

% ============================================================
% 8. \tkzHeatTube  加热试管
% ============================================================
% pstChauffageTube: 保存用户参数, 传递给 \tkzTestTube
% pst-labo 用 \psset 全局继承; tikz-chemlab 用 \tkzTestTube[noDefaults] 跳过重置
\newif\ifchemlab@noDefaults
\newcommand{\tkzHeatTube}[1][]{%
  \begingroup
  \tikzchemlabdefaults
  \pgfkeys{/tikz-chemlab/.cd, #1}%
  \chemlab@noDefaultstrue
  \begin{scope}
    % pst-labo bbox: tubeOnly=(-2,-5)(5,4), 其他=(-3,-5)(10,4)
    % 不使用 \useasboundingbox — pst-labo 的 pspicture 不裁剪内容,
    % 内容可以超出 bbox (如 elbowTube 弯管延伸到 x>5)
    % 让 TikZ 自动检测完整内容 bbox
    \ifTkzBubbling
      % 上方试管/烧瓶: 继承用户设置的 glassType, 倾斜-60°, 有弯管
      \edef\chemlab@tmpGlassType{\tkzGlassTypeName}%
      % pst-labo: \rput(2.5,2.5){\tkzTestTube[elbowTube=true,tiltAngle=-60]}
      % elbowTube=true → bbox (-1,0)(6,4) center (2.5,2)
      % → 局部原点 (瓶底) 世界 (2.5-2.5, 2.5-2) = (0, 0.5)
      \begin{scope}[shift={(0,0.5)}]
        \tkzTestTube[elbowTube=true,tiltAngle=-60,noDefaults=true]
      \end{scope}
      % 下方试管: pst-labo \rput(8.2,-3){\tkzTestTube[tiltAngle=30,...]}
      % tiltAngle>0 → bbox (-3,0)(1,4) center (-1,2)
      % → 局部原点 (试管底) 世界 (8.2-(-1), -3-2) = (9.2, -5)
      \begin{scope}[shift={(9.2,-5)}]
        \tkzTestTube[liquidLevel1=70,glassType=tube,elbowTube=false,
                       tiltAngle=30,
                       liquidStyle1=Champagne,
                       substance={\tkzChampagneBubbles[10]},clamp=false]
      \end{scope}
    \else
      % 非barbotage: 上方器皿继承 glassType
      \edef\chemlab@tmpGlassType{\tkzGlassTypeName}%
      \ifcase\tkzGlassType
        % glassType=tube
        \ifTkzTubeOnly
          % tubeOnly + tube: pst-labo \rput(1.5,2.5)
          % tiltAngle<0 → bbox (-1,0)(3,4) center (1,2) → 世界 (0.5, 0.5)
          \begin{scope}[shift={(0.5,0.5)}]\tkzTestTube[tiltAngle=-60,noDefaults=true]\end{scope}
        \else
          \ifTkzTestTubeElbowTube
            % elbowTube: pst-labo \rput(2,2.5)
            % elbowTube=true → bbox (-1,0)(6,4) center (2.5,2) → 世界 (-0.5, 0.5)
            \begin{scope}[shift={(-0.5,0.5)}]\tkzTestTube[tiltAngle=-60,noDefaults=true]\end{scope}
          \else
            % 无tubeCoude: pst-labo \rput(1,2.5)
            % tiltAngle<0 → bbox (-1,0)(3,4) center (1,2) → 世界 (0, 0.5)
            \begin{scope}[shift={(0,0.5)}]\tkzTestTube[tiltAngle=-60,noDefaults=true]\end{scope}
          \fi
        \fi
      \or
        % glassType=roundFlask (无论 tubeOnly 与否, pst-labo 都是 \rput(2,2.5))
        % tiltAngle<0 → bbox (-1,0)(3,4) center (1,2) → 局部原点世界 (2-1, 2.5-2) = (1, 0.5)
        \begin{scope}[shift={(1,0.5)}]\tkzTestTube[tiltAngle=-60,noDefaults=true]\end{scope}
      \else
        % 其他 glassType: pst-labo \rput(1,2.5) → 世界 (0, 0.5)
        \begin{scope}[shift={(0,0.5)}]\tkzTestTube[tiltAngle=-60,noDefaults=true]\end{scope}
      \fi
    \fi
    \ifTkzBunsenBurner
      % pst-labo: \rput(1,-5){\pst@BecBunsen} (无 pspicture, 实测原点对齐)
      % 灯底位于世界 (1, -5)
      \begin{scope}[shift={(1,-5)}]\tkzBunsenBurner\end{scope}
    \fi
  \end{scope}
  \chemlab@noDefaultsfalse
  \endgroup
}

% ============================================================
% 9. \tkzHeatFlask  加热烧瓶
% ============================================================
\newcommand{\tkzHeatFlask}[1][]{%
  \tikzchemlabdefaults
  \pgfkeys{/tikz-chemlab/.cd, #1}%
  % pst-labo: 每个 \tkzHeatFlask 有独立 pspicture, 不跨实例裁剪
  % tikz-chemlab: 不设全局 clip, 让所有内容自由绘制
  % pst-labo: \tkzRoundFlask 中 glassType=0 (tube) 也画 \pst@roundFlask
  % tikz-chemlab: 强制 glassType=0 → 1 (roundFlask)
  \begin{scope}[xshift=\chemlab@xshift]
    % 不设 \useasboundingbox, 让 TikZ 自动检测完整内容 bbox
    % 物体间相对位置由 shift 控制, 与 pst-labo 一致
    \ifTkzBunsenBurner\def\cnth{-5}\else\def\cnth{1}\fi
    % 气体回收装置
    \ifTkzGasCollection
      \ifTkzTestTubeCurvedTube
        \begin{scope}[shift={(8,-5)}]
          \tkzCrystallizingDish\relax\tkzTestTubeRack
          \begin{scope}[shift={(0,0.7)}]\tkzInvertedTube\relax\tkzBubbles\end{scope}
        \end{scope}
      \else\ifTkzTestTubeCurvedTubeShort
        \TkzBunsenBurnerfalse
        \begin{scope}[shift={(6,1)}]
          \tkzCrystallizingDish\relax\tkzTestTubeRack
          \begin{scope}[shift={(0,0.7)}]\tkzInvertedTube\relax\tkzBubbles\end{scope}
        \end{scope}
      \else\ifTkzTestTubeDoubleTube
        \TkzBunsenBurnerfalse
        \begin{scope}[shift={(6.2,1)}]
          \tkzCrystallizingDish\relax\tkzTestTubeRack
          \begin{scope}[shift={(0,0.7)}]\tkzInvertedTube\relax\tkzBubbles\end{scope}
        \end{scope}
      \fi\fi\fi
    \fi
    % 烧瓶 + 液体 + 导管
    % pst-labo: \rput(1,3){\tkzRoundFlask} 默认使用 bbox 中心对齐 (ref=c)
    % \tkzRoundFlask 的 pspicture bbox = (-2,0)(2,4), 中心 (0,2)
    % 所以局部原点 (瓶底) 实际位于世界 (1, 3-2) = (1, 1)
    \begin{scope}[shift={(1,1)}]
      % 液体 (clip 到烧瓶内部 ∩ 液体矩形, 匹配 pst-labo \psclip 双重裁剪)
      \begin{scope}
        \tkzGlassClipOptionsSet
        \clip[rounded corners=\tkzGlassClipOptions] \tkzGlassClipPath;
        % 第二次裁剪: 液体矩形 (匹配 pst-labo \psframe 裁剪)
        % 只在有液体时才裁剪; 无液体时不裁剪 (物质可见到玻璃顶部)
        \ifnum\tkzLiquidLevelI>0
          \pgfmathsetmacro{\chemlab@liquidClipY}{3 * \tkzLiquidLevelI / 100 * cos(\tkzTiltAngle)}%
          \clip (-4,-1.5) rectangle (5,\chemlab@liquidClipY);
        \fi
        \ifnum\tkzLiquidLevelI>0
          \tkzTestTubeLiquid{\tkzLiquidLevelI}{\tkzLiquidStyleI}{\tkzTiltAngle}
        \fi
        \tkzSubstance
        \tkzSolid
      \end{scope}
      % 玻璃器皿轮廓 (根据 glassType 选择: roundFlask/erlenmeyer/beaker 等)
      \tkzGlassDraw
      % 导管 (在瓶口 y=3.5 处) — 匹配 pst-labo \tkzRoundFlask 的全部导管类型
      \ifTkzTestTubeStraightTube
        \begin{scope}[shift={(0,3.5)}]\tkzTestTubeStraightTube\end{scope}
      \else\ifTkzTestTubeElbowTube
        \begin{scope}[shift={(0,3.5)}]\tkzTestTubeElbowTube\end{scope}
      \else\ifTkzTestTubeUTube
        \begin{scope}[shift={(0,3.5)}]\tkzTestTubeUTube\end{scope}
      \else\ifTkzTestTubeUTubeLong
        \begin{scope}[shift={(0,3.5)}]\tkzTestTubeUTubeLong\end{scope}
      \else\ifTkzTestTubeCurvedTube
        \begin{scope}[shift={(0,3.5)}]\tkzCurvedTube\end{scope}
      \else\ifTkzTestTubeCurvedTubeShort
        \begin{scope}[shift={(0,3.5)}]\tkzCurvedTubeShort\end{scope}
      \else\ifTkzTestTubeStopper
        \begin{scope}[shift={(0,3.5)}]\tkzTestTubeStopper\end{scope}
      \else\ifTkzTestTubeDoubleTube
        \begin{scope}[shift={(0,3.5)}]\tkzTestTubeDoubleTube\end{scope}
      \fi\fi\fi\fi\fi\fi\fi\fi
    \end{scope}
    \ifTkzBubbling
      % pst-labo: \rput(5.3,-3){\tkzTestTube[...]}, tiltAngle=0 默认, glassType=tube
      % bbox (-1.5,0)(1.5,4) center (0,2) → 局部原点 (试管底) 世界 (5.3, -5)
      \begin{scope}[shift={(5.3,-5)}]
        \tkzTestTube[liquidLevel1=70,glassType=tube]
      \end{scope}
    \fi
    \ifTkzBunsenBurner
      % pst-labo: \rput(1,-5){\pst@BecBunsenGrille} (无 pspicture, 实测原点对齐)
      \begin{scope}[shift={(1,-5)}]\tkzBunsenBurnerGrill\end{scope}
    \fi
  \end{scope}
  % 自动推进原点 (匹配 pst-labo \pspicture 的布局宽度)
  % 使用全局 dimen 累加, scope 的 xshift 会正确应用
  % pst-labo bbox 宽度: curvedTubeShort/DoubleTube = 11, 其他 = 6
  \ifTkzTestTubeUTubeLong
    % uTubeLong: 不推进
  \else\ifTkzTestTubeUTube
    % uTube: 不推进
  \else\ifTkzTestTubeCurvedTubeShort
    \advance\chemlab@xshift by 11cm %  bbox (-2,0)(9,8), width = 11
  \else\ifTkzTestTubeDoubleTube
    \advance\chemlab@xshift by 11cm %  bbox (-2,0)(9,8), width = 11
  \else
    \advance\chemlab@xshift by 6cm %  bbox (-2,\cnth)(4,6), width = 6
  \fi\fi\fi\fi
}

% ============================================================
% 10. \tkzRoundFlask  烧瓶 (高级命令)
% ============================================================
\newcommand{\tkzRoundFlask}[1][]{%
  \tikzchemlabdefaults
  \pgfkeys{/tikz-chemlab/.cd, #1}%
  % pst-labo: \tkzRoundFlask 中 glassType=0 (tube) 也画 \pst@roundFlask (case 0 = roundFlask)
  % tikz-chemlab: 强制 glassType=0 → 1 (roundFlask), 避免画试管轮廓
  \begin{scope}[xshift=\chemlab@xshift]
    % 不设 \useasboundingbox, 让 TikZ 自动检测完整内容 bbox
    % bbox 坐标使用局部坐标系, 会随 scale 自动缩放
    % 物体间相对位置由 shift 控制, 与 pst-labo 一致
    % 液体
    \begin{scope}[rotate=\tkzTiltAngle]
      \begin{scope}
        \tkzGlassClipOptionsSet
        \clip[rounded corners=\tkzGlassClipOptions] \tkzGlassClipPath;
        % 第二次裁剪: 液体矩形 (匹配 pst-labo \psframe 裁剪)
        % 只在有液体时才裁剪; 无液体时不裁剪 (物质可见到玻璃顶部)
        \ifnum\tkzLiquidLevelI>0
          \pgfmathsetmacro{\chemlab@liquidClipY}{3 * \tkzLiquidLevelI / 100 * cos(\tkzTiltAngle)}%
          \clip (-4,-1.5) rectangle (5,\chemlab@liquidClipY);
        \fi
        \ifnum\tkzLiquidLevelI>0
          \tkzTestTubeLiquid{\tkzLiquidLevelI}{\tkzLiquidStyleI}{\tkzTiltAngle}
        \fi
        \ifnum\tkzLiquidLevelII>0
          \tkzTestTubeLiquid{\tkzLiquidLevelII}{\tkzLiquidStyleII}{\tkzTiltAngle}
        \fi
        \ifnum\tkzLiquidLevelIII>0
          \tkzTestTubeLiquid{\tkzLiquidLevelIII}{\tkzLiquidStyleIII}{\tkzTiltAngle}
        \fi
        \tkzSubstance
        \tkzSolid
      \end{scope}
      % 标签
      \ifTkzLabel
        \node[fill=VeryLightGray,draw=none,inner sep=2pt] at (0,1.5) {\tkzNumber};
      \fi
      % 玻璃器皿轮廓
      \tkzGlassDraw
      % 顶部附件 (按优先级)
      \ifTkzTestTubeStraightTube
        \begin{scope}[shift={(0,3.5)}]\tkzTestTubeStraightTube\end{scope}
      \else\ifTkzTestTubeElbowTube
        \begin{scope}[shift={(0,3.5)}]\tkzTestTubeElbowTube\end{scope}
      \else\ifTkzTestTubeUTube
        \begin{scope}[shift={(0,3.5)}]\tkzTestTubeUTube\end{scope}
      \else\ifTkzTestTubeUTubeLong
        \begin{scope}[shift={(0,3.5)}]\tkzTestTubeUTubeLong\end{scope}
      \else\ifTkzTestTubeCurvedTube
        \begin{scope}[shift={(0,3.5)}]\tkzCurvedTube\end{scope}
      \else\ifTkzTestTubeCurvedTubeShort
        \begin{scope}[shift={(0,3.5)}]\tkzCurvedTubeShort\end{scope}
      \else\ifTkzTestTubeStopper
        \begin{scope}[shift={(0,3.5)}]\tkzTestTubeStopper\end{scope}
      \else\ifTkzTestTubeDoubleTube
        \begin{scope}[shift={(0,3.5)}]\tkzTestTubeDoubleTube\end{scope}
      \else\ifTkzTestTubeCondenserBulbs
        \tkzRoundFlaskReflux
        \begin{scope}[shift={(0,3.5)}]\tkzCondenserBulbs\end{scope}
        \tkzHeaterRoundFlask
      \fi\fi\fi\fi\fi\fi\fi\fi\fi
    \end{scope}
  \end{scope}
  % 自动推进原点 (匹配 pst-labo \pspicture 的布局宽度)
  % pst-labo bbox: condenserBulbs = (-2,0)(2,12), 其他 = (-2,0)(2,4), width = 4
  \advance\chemlab@xshift by 4cm %  bbox width = 4
}

% ============================================================
% 11. \tkzChampagneBubbles  香槟气泡
% ============================================================
% 原版 pst-labo 源码:
%   \setrandim{\BubbleX}{-1pt}{1pt}    -- BulleX ∈ [-1, 1] (纯数字, cm)
%   \setrandim{\BubbleY}{0pt}{2pt}     -- BulleY ∈ [0, 2] (纯数字)
%   Y 坐标 = \pointless\BubbleY * liquidLevel1/100  (乘以液位百分比)
%   半径 = {\BubbleY} (带 pt 单位, 1pt = 0.0353 cm)
%   \pst@dimA 分级放大: BulleY<1pt → ×3, 1pt≤BulleY<2pt → ×2
%
% 实测原版渲染图:
%   气泡为黑色边空心圆 (\pscircle, fillstyle=none)
%   半径 0.05~0.175 单位 (比源码计算 0~0.071 大, 可能是渲染缩放)
%   气泡团中心 (世界): (7.89, -0.975)
\newcommand{\tkzChampagneBubbles}[1][25]{%
  \foreach \iBubble in {1,...,#1}{%
    % BulleX ∈ [-1, 1] (纯数字, cm)
    \pgfmathsetmacro{\BubbleX}{-1 + 2*rnd}%
    % BulleY ∈ [0, 2] (纯数字)
    \pgfmathsetmacro{\BubbleY}{2*rnd}%
    % Y 坐标 = BulleY * liquidLevel1/100 (匹配源码逻辑)
    \pgfmathsetmacro{\BubbleYpos}{\BubbleY * \tkzLiquidLevelI / 100}%
    % dimA 分级放大 (匹配源码逻辑)
    \pgfmathsetmacro{\BubbleRand}{rnd}%
    \ifdim\BubbleRand pt<0.5pt
      % BulleY < 1pt (50%): dimA = BulleY * 3
      \pgfmathsetmacro{\BubbleR}{\BubbleY * 0.0875}%
    \else
      % 1pt <= BulleY < 2pt (50%): dimA = BulleY * 2
      \pgfmathsetmacro{\BubbleR}{\BubbleY * 0.0875}%
    \fi
    \begin{scope}[rotate=\tkzTiltAngle]
      \draw[draw=black, line width=0.6pt] (\BubbleX,\BubbleYpos) circle (\BubbleR);
    \end{scope}
  }%
}

% ============================================================
% 12. \tkzBubbles  倒置管中气泡
% ============================================================
\newcommand{\tkzBubbles}[1][25]{%
  \pgfgettransformentries{\clab@aa}{\clab@bb}{\clab@cc}{\clab@dd}{\clab@xx}{\clab@yy}%
  \pgfmathsetmacro{\clab@scale}{(\clab@aa + \clab@dd) / 2}%
  \foreach \iBubble in {1,...,#1}{%
    % pst-labo: \setrandim{BulleX}{-0.1pt}{0.1pt} → \pointless 去掉pt, 数值直接用作局部坐标
    \pgfmathsetmacro{\BubbleX}{-0.1 + 0.2*rnd}%
    \pgfmathsetmacro{\BubbleY}{0.15 + 2.85*rnd}%
    % pst-labo: dimh = BulleY * (3 if <1pt, 2 if <2pt, 1 if >=2pt)
    % dotscale = dimh, dot radius = dotscale * dotsize/2
    % dotsize default = 2pt + 2*linewidth = 2pt + 2*0.8pt = 3.6pt → radius = dimh * 1.8pt
    % 转换为局部坐标: local = pt / (72 * scale)
    \pgfmathparse{\BubbleY<1?3:(\BubbleY<2?2:1)}%
    \let\BubbleScale\pgfmathresult
    \pgfmathsetmacro{\BubbleDimh}{\BubbleY * \BubbleScale}%
    \pgfmathsetmacro{\BubbleR}{1.8 * \BubbleDimh / (72 * \clab@scale)}%
    % Y position = BulleY * 1.5 (局部坐标, 匹配 pst-labo \pointless\BubbleY 1.5 mul)
    \pgfmathsetmacro{\BubbleYpos}{\BubbleY * 1.5}%
    \begin{scope}[rotate=\tkzTiltAngle]
      % pst-labo: dotstyle=o (open circle = hollow, black outline)
      \draw[draw=black, line width=0.4pt, fill=white] (\BubbleX,\BubbleYpos) circle (\BubbleR);
    \end{scope}
  }%
}

% ============================================================
% 13. \tkzFilaments  絮状物
% ============================================================
\newcommand{\tkzFilaments}[2][5]{%
  \pgfgettransformentries{\fil@aa}{\fil@bb}{\fil@cc}{\fil@dd}{\fil@xx}{\fil@yy}%
  \pgfmathsetmacro{\fil@scale}{(\fil@aa + \fil@dd) / 2}%
  \foreach \iFil in {1,...,#1}{%
    % pst-labo: \setrandim{\pst@FilamentWidth}{0pt}{0.05pt} → \pointless = [0, 0.05)
    % linewidth=\pointless\pst@FilamentWidth 解释为 N × unit (PSTricks 规则)
    % 实际线宽 = [0, 0.05) × scale cm (scale 对应 unit)
    \pgfmathsetmacro{\fw}{0.05 * \fil@scale * rnd}%
    \pgfmathsetmacro{\myRand}{rnd}%  ONE random value per filament (匹配 pst-labo MyRand)
    \begin{scope}[rotate=\tkzTiltAngle]
      % pst-labo \pscurve 通过4个点的平滑曲线 (不是贝塞尔):
      %   P1: (MyRand-1, 2*MyRand-1)     = ([-1,0), [-1,1))
      %   P2: (MyRand+0.2, MyRand+2)     = ([0.2,1.2), [2,3))
      %   P3: (MyRand, MyRand-1)         = ([0,1), [-1,0))
      %   P4: (MyRand+0.5, MyRand+4)     = ([0.5,1.5), [4,5))
      \draw[line width=\fw cm,#2] plot[smooth,tension=0.5] coordinates {
        ({\myRand-1},{2*\myRand-1})
        ({\myRand+0.2},{\myRand+2})
        ({\myRand},{\myRand-1})
        ({\myRand+0.5},{\myRand+4})
      };
    \end{scope}
  }%
}

% ============================================================
% 14. \tkzIronNail  铁钉
% ============================================================
% 单个铁钉形状 (单位 0.5)
\newcommand{\tkzNails}{%
  % pst-labo 实际 SLW=0.71135pt=0.025cm (从 xdv 提取)
  % \psarc[unit=0.5,linewidth=0.1] 的 \psline 继承同一 linewidth
  \draw[scale=0.5,line width=0.025cm,blue]
    ([shift={(0,0.5)}]190:0.4) arc[start angle=190,end angle=260,radius=0.4];
  \draw[scale=0.5,line width=0.025cm,blue] (-0.28,0.2) -- (0.5,1);
}

\newcommand{\tkzIronNail}[1][60]{%
  \pgfgettransformentries{\clab@aa}{\clab@bb}{\clab@cc}{\clab@dd}{\clab@xx}{\clab@yy}%
  \pgfmathsetmacro{\clab@scale}{(\clab@aa + \clab@dd) / 2}%
  \foreach \iNail in {1,...,#1}{%
    % pst-labo: \setrandim{ClouX}{-1cm}{1cm} → 物理坐标, 需除以 scale 转换为局部坐标
    \pgfmathsetmacro{\NailX}{(-1 + 2*rnd) / \clab@scale}%
    \pgfmathsetmacro{\NailY}{(-0.2 + 0.6*rnd) / \clab@scale}%
    \pgfmathsetmacro{\NailAngle}{-90 + 180*rnd}%
    \begin{scope}[rotate=\tkzTiltAngle]
      \begin{scope}[shift={(\NailX,\NailY)},rotate=\NailAngle]
        \tkzNails
      \end{scope}
    \end{scope}
  }%
}

% ============================================================
% 15. \tkzCopperTurnings  铜屑
% ============================================================
% 单个铜屑形状 (单位 0.5)
\newcommand{\tkzCopper}{%
  % pst-labo 实际 SLW=0.28458pt=0.01cm (从 xdv 提取)
  \draw[scale=0.5,line width=0.01cm,Copper]
    (0.25,0.25) .. controls (0,0) and (-0.25,0.25) .. (0.25,0.75)
    .. controls (0,0.5) and (-0.25,0.75) .. (0.25,1.25)
    .. controls (0,1) and (-0.25,1.25) .. (0.25,1.75);
}

\newcommand{\tkzCopperTurnings}[1][30]{%
  \pgfgettransformentries{\clab@aa}{\clab@bb}{\clab@cc}{\clab@dd}{\clab@xx}{\clab@yy}%
  \pgfmathsetmacro{\clab@scale}{(\clab@aa + \clab@dd) / 2}%
  \foreach \iT in {1,...,#1}{%
    \pgfmathsetmacro{\TX}{(-1 + 2*rnd) / \clab@scale}%
    \pgfmathsetmacro{\TY}{(-0.2 + 0.5*rnd) / \clab@scale}%
    \pgfmathsetmacro{\TAngle}{-180 + 360*rnd}%
    \begin{scope}[rotate=\tkzTiltAngle]
      \begin{scope}[shift={(\TX,\TY)},rotate=\TAngle]
        \tkzCopper
      \end{scope}
    \end{scope}
  }%
}

% ============================================================
% 16. \tkzZincGranules  锌粒
% ============================================================
% 单个锌粒形状 (单位 0.5)
\newcommand{\tkzZinc}{%
  \filldraw[scale=0.5,fill=LightGray,draw=black]
    plot[smooth,tension=0.5] coordinates {(0,0)(-0.25,0.25)(0,0.5)
      (0.25,0.75)(0.5,0.5)(0.15,0.4)(0,0)} -- cycle;
}

\newcommand{\tkzZincGranules}[1][25]{%
  \pgfgettransformentries{\clab@aa}{\clab@bb}{\clab@cc}{\clab@dd}{\clab@xx}{\clab@yy}%
  \pgfmathsetmacro{\clab@scale}{(\clab@aa + \clab@dd) / 2}%
  \foreach \iG in {1,...,#1}{%
    \pgfmathsetmacro{\GX}{(-1 + 2*rnd) / \clab@scale}%
    \pgfmathsetmacro{\GY}{(-0.2 + 0.5*rnd) / \clab@scale}%
    \pgfmathsetmacro{\GAngle}{-180 + 360*rnd}%
    \begin{scope}[rotate=\tkzTiltAngle]
      \begin{scope}[shift={(\GX,\GY)},rotate=\GAngle]
        \tkzZinc
      \end{scope}
    \end{scope}
  }%
}

% ============================================================
% 17. \tkzBalls  3D 球状物
% ============================================================
\newcommand{\tkzBalls}[1][50]{%
  \foreach \iB in {1,...,#1}{%
    \pgfmathsetmacro{\BX}{-1.5 + 3*rnd}%
    \pgfmathsetmacro{\BY}{3*rnd}%
    \pgfmathsetmacro{\BR}{0.01 + 0.1*rnd}%
    % 按 Y 位置调整半径
    \pgfmathparse{\BY>2?3:(\BY>1?2:1)}%
    \let\Bscale\pgfmathresult
    \pgfmathsetmacro{\BRscaled}{\BR * \Bscale}%
    \begin{scope}[rotate=\tkzTiltAngle]
      % pst-labo Ball3D: GradientCircle=true
      % gradbegin=CMYK(0,0.1,0.95,0.1) → 实际渲染为 RGB(0,0.1,0.95)=蓝色 (固定颜色)
      % gradend=white
      % 径向渐变: 中心=white, 边缘=蓝色
      \definecolor{billeBlue}{rgb}{0,0.1,0.95}%
      \shade[inner color=white, outer color=billeBlue] (\BX,\BY) circle (\BRscaled);
    \end{scope}
  }%
}

% ============================================================
% 18. \tkzBubblesThreeD  2D 气泡
% ============================================================
\newcommand{\tkzBubblesThreeD}[2][50]{%
  \foreach \iB in {1,...,#1}{%
    \pgfmathsetmacro{\BX}{-1.5 + 3*rnd}%
    \pgfmathsetmacro{\BY}{3*rnd}%
    \pgfmathsetmacro{\BR}{0.01 + 0.1*rnd}%
    \pgfmathparse{\BY>2?3:(\BY>1?2:1)}%
    \let\Bscale\pgfmathresult
    \pgfmathsetmacro{\BRscaled}{\BR * \Bscale}%
    \begin{scope}[rotate=\tkzTiltAngle]
      \fill[#2,line width=0.2pt] (\BX,\BY) circle (\BRscaled);
    \end{scope}
  }%
}

\endinput
       % 本文件
%
%  使用示例 :
%    \begin{tikzpicture}
%       \tkzTestTube[glassType=beaker,liquidLevel1=50]
%    \end{tikzpicture}
%
%  参数键 (与原 pst-labo 兼容, 通过 pgfkeys 传入) :
%    burette, pHMeter, buretteColor, buretteLevel,
%    buretteReagent, beakerReagent, reagentColor, reagentLevel,
%    DistillateColor, MixtureStyle,
%    stopper, straightTube, elbowTube, uTube, uTubeLong,
%    curvedTube, curvedTubeShort, DoubleTube, condenserBulbs,
%    clamp, bubbling, gasCollection, tubeOnly, gasTubeStraight,
%    bunsenBurner, Number, label, magneticStirrer,
%    tiltAngle, liquidLevel1, liquidLevel2, liquidLevel3,
%    liquidStyle1, liquidStyle2, liquidStyle3,
%    substance, solid, glassType, TapOpen
%
%  glassType 取值 : tube / roundFlask / beaker / erlenmeyer / reagentBottle / volumetricFlask / watchGlass
% ============================================================

\NeedsTeXFormat{LaTeX2e}
\ProvidesFile{tikz-chemlab.tex}[2026/08/03 v1.0.0]

% --- 文档类检测 ----------------------------------------------
% standalone 会裁剪到 \useasboundingbox, article/book 不会
% 用于在 \tkzHeatFlask 中选择 bbox 策略
\@ifclassloaded{standalone}{\def\chemlab@isstandalone{T}}{\def\chemlab@isstandalone{F}}

% --- 全局 xshift 累加器 (用于多个 \tkzHeatFlask/\tkzRoundFlask 自动推进原点) ---
\newdimen\chemlab@xshift
\chemlab@xshift=0pt
% --- 首次调用标志 (用于在首个 \tkzHeatFlask 设置全局 clip) ---
\newif\ifchemlab@firstcall
\chemlab@firstcalltrue

% --- 依赖加载 ------------------------------------------------
\RequirePackage{tikz}
\usetikzlibrary{arrows.meta,calc,patterns,patterns.meta,shadings,decorations.pathmorphing,decorations.pathreplacing}
\@input{tikz-chemlabObj.tex}

% --- 每个 tikzpicture 开始时自动重置 xshift 和 firstcall ---
% 这样用户无需手动调用 \chemlabresetxshift
\tikzset{every picture/.append style={execute at begin picture={%
  \chemlab@xshift=0pt\chemlab@firstcalltrue%
}}}

% --- 颜色 (原 pst-labo.tex 中定义, 与 Obj 文件中的部分重复此处保留以覆盖) ---
\definecolor{Beige}{cmyk}{0,0,0.1,0.04}
\definecolor{LightGray}{cmyk}{0,0,0,0.2}
\definecolor{VeryLightGray}{cmyk}{0,0,0,0.1}
\definecolor{VeryPaleOrange}{cmyk}{0,0.1,0.3,0}
\definecolor{PaleOrange}{cmyk}{0,0.2,0.4,0}
\definecolor{LightBlue1}{cmyk}{0.2,0,0,0}
\definecolor{LightBlue}{cmyk}{0.22,0.05,0,0.1}
\definecolor{Copper}{cmyk}{0,0.9,0.9,0.2}
\definecolor{Brown}{cmyk}{0,0.3,0.5,.3}

% --- 样式 (对应原 \newpsstyle) --------------------------------
% 使用 \tikzset 定义命名样式; 调用 : \fill[liquidStyle1] ...
\tikzset{
  liquidStyle1/.style={fill=cyan,draw=none},
  liquidStyle2/.style={fill=yellow,draw=none},
  liquidStyle3/.style={fill=magenta,draw=none},
  Ball3D/.style={shading=ball,ball color=cyan,draw=none},
  Champagne/.style={fill=Beige,draw=none},
  Blood/.style={fill=red,draw=none},
  Cobalt/.style={fill=blue,line width=0.2cm},
  Oil/.style={fill=yellow,draw=none},
  Vinegar/.style={fill=magenta,draw=none},
  % pst-labo Diffusion: gradbegin=rgb(0,0.1,0.95), gradend=cyan, gradmidpoint=0
  % gradmidpoint=0: 渐变中点在底部, 大部分区域为 cyan, 底部为 blue
  Diffusion/.style={bottom color={rgb,255:red,0;green,25;blue,242},top color=cyan,draw=none},
  MixtureDiffusion2/.style={bottom color=white,top color=red,draw=red},
  flameAndGrill/.style={bottom color=PaleOrange,top color=yellow,draw=none},
  YellowStripes/.style={pattern=north east lines,pattern color=yellow},
  BlueDiffusion/.style={bottom color=green,top color=cyan,draw=none},
  % 分液漏斗两相液体 (对应 pst-labo \newpsstyle{AqueousPhase}/{OrganicPhase})
  % 颜色定义在 tikz-chemlabObj.tex: AqueousPhaseColor / OrganicPhaseColor
  AqueousPhase/.style={fill=AqueousPhaseColor,draw=none},
  OrganicPhase/.style={fill=OrganicPhaseColor,draw=none},
  % 额外 : 液体通用填充
  tkz liquid/.style={draw=none},
}

% --- 布尔标志 (\newif) ----------------------------------------
\newif\ifTkzBurette
\newif\ifTkzTitrationPHMeter
\newif\ifTkzTestTubeStopper
\newif\ifTkzTestTubeStraightTube
\newif\ifTkzTestTubeElbowTube
\newif\ifTkzTestTubeUTube
\newif\ifTkzTestTubeUTubeLong
\newif\ifTkzTestTubeCurvedTube
\newif\ifTkzTestTubeCurvedTubeShort
\newif\ifTkzTestTubeDoubleTube
\newif\ifTkzTestTubeCondenserBulbs
\newif\ifTkzTestTubeClamp
\newif\ifTkzBubbling
\newif\ifTkzGasCollection
\newif\ifTkzTubeOnly
\newif\ifTkzGasTubeStraight
\newif\ifTkzBunsenBurner
\newif\ifTkzLabel
\newif\ifTkzMagneticStirrer
\newif\ifTkzTapOpen            % (用于 \tkzSeparatoryFunnel)

% --- 参数存储宏 ----------------------------------------------
\def\tkzTiltAngle{0}
\def\tkzLiquidLevelI{50}
\def\tkzLiquidLevelII{0}
\def\tkzLiquidLevelIII{0}
\def\tkzLiquidStyleI{liquidStyle1}
\def\tkzLiquidStyleII{liquidStyle2}
\def\tkzLiquidStyleIII{liquidStyle3}
\def\tkzSubstance{}
\def\tkzSolid{}
\def\tkzNumber{}
\def\tkzGlassType{0}          % 0..6
\def\tkzGlassTypeName{tube}   % 字符串名称 (供嵌套调用传递)
\def\tkzBuretteColor{VeryPaleOrange}
\def\tkzBuretteLevel{20}
\def\tkzBuretteReagent{}
\def\tkzBeakerReagent{}
\def\tkzReagentColor{PaleOrange}
\def\tkzReagentLevel{2}
\def\tkzDistillateColor{yellow}
\def\tkzMixtureStyle{BlueDiffusion}

% --- glassType 字符串 -> 数字 --------------------------------
\newcommand{\tkzGlassTypeLookup}[1]{%
  \ifnum\pdfstrcmp{#1}{tube}=0 0\else
  \ifnum\pdfstrcmp{#1}{roundFlask}=0 1\else
  \ifnum\pdfstrcmp{#1}{beaker}=0 2\else
  \ifnum\pdfstrcmp{#1}{erlenmeyer}=0 3\else
  \ifnum\pdfstrcmp{#1}{reagentBottle}=0 4\else
  \ifnum\pdfstrcmp{#1}{volumetricFlask}=0 5\else
  \ifnum\pdfstrcmp{#1}{watchGlass}=0 6\else
  0\fi\fi\fi\fi\fi\fi\fi
}

% --- 线条粗细变量 (全局, 可由用户通过 \chemlabset 修改) ---
% 四类线宽可独立控制:
%   \chemlabLineWidth            - 主轮廓线 (玻璃器皿外壁)
%   \chemlabEllipseLineWidth     - 瓶口椭圆 (erlenmeyer/reagentBottle/倒置试管, 略粗)
%   \chemlabTubeEllipseLineWidth - 试管口椭圆 (更粗, 原 pst-labo 用 0.08cm)
%   \chemlabGraduationLineWidth  - 刻度线 (较细)
% 默认值与原 pst-labo 完全一致 (在 unit=1cm 时)
% 注: 与 pst-labo 一致, 线宽会自动按当前 tikz scale 缩放
%     (例如 scale=0.5 时, 实际线宽 = 基础线宽 × 0.5)
\def\chemlabBaseLineWidth{0.053cm}
\def\chemlabBaseEllipseLineWidth{0.07cm}
\def\chemlabBaseTubeEllipseLineWidth{0.08cm}
\def\chemlabBaseGraduationLineWidth{0.02cm}

% \chemlab@scaleLength{<基础长度>} - 返回按当前 tikz 缩放因子缩放后的长度
% 用于线宽跟随 scale=N 自动缩放, 与 pst-labo 的 unit=Ncm 行为一致
% 实现: 在每条 \draw 之前由 \chemlabUpdateLineWidths 把缩放后的长度
% 存入 dimen 寄存器, 绘制时直接使用寄存器
\newdimen\chemlabLineWidth
\newdimen\chemlabEllipseLineWidth
\newdimen\chemlabTubeEllipseLineWidth
\newdimen\chemlabGraduationLineWidth
\edef\chemlab@aa{1}\edef\chemlab@bb{0}\edef\chemlab@cc{0}\edef\chemlab@dd{1}\edef\chemlab@xx{0}\edef\chemlab@yy{0}
\newcommand{\chemlabUpdateLineWidths}{%
  \pgfgettransformentries{\chemlab@aa}{\chemlab@bb}{\chemlab@cc}{\chemlab@dd}{\chemlab@xx}{\chemlab@yy}%
  \pgfmathsetlength{\chemlabLineWidth}{\chemlabBaseLineWidth * (\chemlab@aa + \chemlab@dd) / 2}%
  \pgfmathsetlength{\chemlabEllipseLineWidth}{\chemlabBaseEllipseLineWidth * (\chemlab@aa + \chemlab@dd) / 2}%
  \pgfmathsetlength{\chemlabTubeEllipseLineWidth}{\chemlabBaseTubeEllipseLineWidth * (\chemlab@aa + \chemlab@dd) / 2}%
  \pgfmathsetlength{\chemlabGraduationLineWidth}{\chemlabBaseGraduationLineWidth * (\chemlab@aa + \chemlab@dd) / 2}%
}
% 默认值 (未调用 UpdateLineWidths 时的回退)
\chemlabLineWidth=0.053cm
\chemlabEllipseLineWidth=0.07cm
\chemlabTubeEllipseLineWidth=0.08cm
\chemlabGraduationLineWidth=0.02cm

% --- 火焰参数 (全局, 可由用户通过 \chemlabset 修改) ---
% 颜色 (默认与原 pst-labo 一致)
\def\chemlabFlameOuterColor{PaleOrange}   % 外焰起始色 (inner color)
\def\chemlabFlameOuterEndColor{yellow}    % 外焰结束色 (outer color)
\def\chemlabFlameInnerColor{cyan}         % 内焰起始色 (inner color)
\def\chemlabFlameInnerEndColor{white}     % 内焰结束色 (outer color)
% 大小 (默认与原 pst-labo 一致)
% 火焰缩放因子: 1.0=原版大小, 0.5=半高, 2.0=双高
\def\chemlabFlameScale{1.0}

% 线条粗细关键字 -> 数值转换 (对应 TikZ 标准)
% thin=0.4pt, semithick=0.6pt, thick=0.8pt, very thick=1.2pt, ultra thick=1.6pt
% 用法: \chemlab@resolveLineWidth{<关键字或数值>}{<目标宏>}
% 结果: 将解析后的数值存入 <目标宏>
\def\chemlab@thin{thin}
\def\chemlab@semithick{semithick}
\def\chemlab@thick{thick}
\def\chemlab@verythick{very thick}
\def\chemlab@ultrathick{ultra thick}
\newcommand{\chemlab@resolveLineWidth}[2]{%
  \def\chemlab@tmp{#1}%
  \ifx\chemlab@tmp\chemlab@thin          \def#2{0.4pt}%
  \else\ifx\chemlab@tmp\chemlab@semithick \def#2{0.6pt}%
  \else\ifx\chemlab@tmp\chemlab@thick     \def#2{0.8pt}%
  \else\ifx\chemlab@tmp\chemlab@verythick \def#2{1.2pt}%
  \else\ifx\chemlab@tmp\chemlab@ultrathick\def#2{1.6pt}%
  \else  % 直接当作数值 (如 0.04cm, 1pt)
    \def#2{#1}%
  \fi\fi\fi\fi\fi
}

% 三个独立设置器 (供 pgfkeys 调用)
\newcommand{\chemlabSetOutlineWidth}[1]{\chemlab@resolveLineWidth{#1}{\chemlabBaseLineWidth}}
\newcommand{\chemlabSetEllipseWidth}[1]{\chemlab@resolveLineWidth{#1}{\chemlabBaseEllipseLineWidth}}
\newcommand{\chemlabSetGraduationWidth}[1]{\chemlab@resolveLineWidth{#1}{\chemlabBaseGraduationLineWidth}}
\newcommand{\chemlabSetTubeEllipseWidth}[1]{\chemlab@resolveLineWidth{#1}{\chemlabBaseTubeEllipseLineWidth}}

% 同时设置四类线宽的快捷方式
\newcommand{\chemlabSetAllLineWidth}[1]{%
  \chemlab@resolveLineWidth{#1}{\chemlabBaseLineWidth}%
  \chemlab@resolveLineWidth{#1}{\chemlabBaseEllipseLineWidth}%
  \chemlab@resolveLineWidth{#1}{\chemlabBaseTubeEllipseLineWidth}%
  \chemlab@resolveLineWidth{#1}{\chemlabBaseGraduationLineWidth}%
}

% ===== 预设样式 (一键设置四类线宽) =====
% 通过 \chemlabset{style=publication} 等使用
% 预设说明 (数值单位 pt, 1pt ≈ 0.0353cm) :
%   default      - 原 pst-labo 默认 (轮廓 1.508pt, 椭圆 1.984pt, 试管椭圆 2.268pt, 刻度 0.567pt)
%   publication  - 出版级 (细线, 适合印刷) : 轮廓 thin, 椭圆 semithick, 试管椭圆 thick, 刻度 thin
%   presentation- 演示级 (粗线, 适合投影) : 轮廓 thick, 椭圆 very thick, 试管椭圆 ultra thick, 刻度 semithick
%   schematic    - 示意图 (中等, 平衡)   : 轮廓 semithick, 椭圆 thick, 试管椭圆 very thick, 刻度 thin
%   draft        - 草稿 (极细, 快速预览) : 全部 thin
%   bold         - 加粗 (高对比)         : 轮廓 very thick, 椭圆 ultra thick, 试管椭圆 ultra thick, 刻度 thick
\def\chemlab@style@default{default}
\def\chemlab@style@publication{publication}
\def\chemlab@style@presentation{presentation}
\def\chemlab@style@schematic{schematic}
\def\chemlab@style@draft{draft}
\def\chemlab@style@bold{bold}
\newcommand{\chemlabSetStyle}[1]{%
  \def\chemlab@styletmp{#1}%
  \ifx\chemlab@styletmp\chemlab@style@default
    \def\chemlabBaseLineWidth{0.053cm}\def\chemlabBaseEllipseLineWidth{0.07cm}\def\chemlabBaseTubeEllipseLineWidth{0.08cm}\def\chemlabBaseGraduationLineWidth{0.02cm}%
  \else\ifx\chemlab@styletmp\chemlab@style@publication
    \def\chemlabBaseLineWidth{0.4pt}\def\chemlabBaseEllipseLineWidth{0.6pt}\def\chemlabBaseTubeEllipseLineWidth{0.8pt}\def\chemlabBaseGraduationLineWidth{0.4pt}%
  \else\ifx\chemlab@styletmp\chemlab@style@presentation
    \def\chemlabBaseLineWidth{0.8pt}\def\chemlabBaseEllipseLineWidth{1.2pt}\def\chemlabBaseTubeEllipseLineWidth{1.6pt}\def\chemlabBaseGraduationLineWidth{0.6pt}%
  \else\ifx\chemlab@styletmp\chemlab@style@schematic
    \def\chemlabBaseLineWidth{0.6pt}\def\chemlabBaseEllipseLineWidth{0.8pt}\def\chemlabBaseTubeEllipseLineWidth{1.2pt}\def\chemlabBaseGraduationLineWidth{0.4pt}%
  \else\ifx\chemlab@styletmp\chemlab@style@draft
    \def\chemlabBaseLineWidth{0.4pt}\def\chemlabBaseEllipseLineWidth{0.4pt}\def\chemlabBaseTubeEllipseLineWidth{0.4pt}\def\chemlabBaseGraduationLineWidth{0.4pt}%
  \else\ifx\chemlab@styletmp\chemlab@style@bold
    \def\chemlabBaseLineWidth{1.2pt}\def\chemlabBaseEllipseLineWidth{1.6pt}\def\chemlabBaseTubeEllipseLineWidth{1.6pt}\def\chemlabBaseGraduationLineWidth{0.8pt}%
  \else
    \typeout{tikz-chemlab: unknown style `#1'. Available: default, publication, presentation, schematic, draft, bold.}%
  \fi\fi\fi\fi\fi\fi
}

% 用户接口: \chemlabset{line width=thick} 或 \chemlabset{outline line width=thick, ellipse line width=0.1cm}
% 或 \chemlabset{style=publication}
\newcommand{\chemlabset}[1]{\pgfkeys{/tikz-chemlab/.cd, #1}}

% --- pgfkeys 参数族 ------------------------------------------
\pgfkeys{
  /tikz-chemlab/.cd,
  % ===== 线条粗细控制 (全局, 三类可独立设置) =====
  % line width / linewidth : 同时设置三类线宽 (轮廓+椭圆+刻度)
  % 可设为: thin / semithick / thick / very thick / ultra thick
  % 或直接给数值 (如 0.04cm, 1pt 等)
  line width/.code={\chemlabSetAllLineWidth{#1}},
  linewidth/.code={\chemlabSetAllLineWidth{#1}},
  % outline line width : 仅设置主轮廓线
  outline line width/.code={\chemlabSetOutlineWidth{#1}},
  outline width/.code={\chemlabSetOutlineWidth{#1}},
  % ellipse line width : 仅设置瓶口椭圆 (erlenmeyer/reagentBottle/倒置试管)
  ellipse line width/.code={\chemlabSetEllipseWidth{#1}},
  ellipse width/.code={\chemlabSetEllipseWidth{#1}},
  % tube ellipse line width : 仅设置试管口椭圆 (原 pst-labo 用 0.08cm, 比其他椭圆粗)
  tube ellipse line width/.code={\chemlabSetTubeEllipseWidth{#1}},
  tube ellipse width/.code={\chemlabSetTubeEllipseWidth{#1}},
  % graduation line width : 仅设置刻度线
  graduation line width/.code={\chemlabSetGraduationWidth{#1}},
  graduation width/.code={\chemlabSetGraduationWidth{#1}},
  % style : 预设样式 (一键设置三类线宽)
  % 可选: default / publication / presentation / schematic / draft / bold
  style/.code={\chemlabSetStyle{#1}},
  % ===== 火焰参数 (本生灯) =====
  % 颜色: 可用任意 TikZ 颜色名 (如 cyan, yellow, red, PaleOrange 等)
  flame outer color/.store in=\chemlabFlameOuterColor,
  flame outer end color/.store in=\chemlabFlameOuterEndColor,
  flame inner color/.store in=\chemlabFlameInnerColor,
  flame inner end color/.store in=\chemlabFlameInnerEndColor,
  % 大小: 缩放因子 (1.0=原版, 0.5=半高, 2.0=双高)
  flame scale/.store in=\chemlabFlameScale,
  burette/.is if=TkzBurette,
  pHMeter/.is if=TkzTitrationPHMeter,
  stopper/.is if=TkzTestTubeStopper,
  straightTube/.is if=TkzTestTubeStraightTube,
  elbowTube/.is if=TkzTestTubeElbowTube,
  uTube/.is if=TkzTestTubeUTube,
  uTubeLong/.is if=TkzTestTubeUTubeLong,
  curvedTube/.is if=TkzTestTubeCurvedTube,
  curvedTubeShort/.is if=TkzTestTubeCurvedTubeShort,
  doubleTube/.is if=TkzTestTubeDoubleTube,
  condenserBulbs/.is if=TkzTestTubeCondenserBulbs,
  clamp/.is if=TkzTestTubeClamp,
  noDefaults/.is if=chemlab@noDefaults,
  bubbling/.is if=TkzBubbling,
  gasCollection/.is if=TkzGasCollection,
  tubeOnly/.is if=TkzTubeOnly,
  gasTubeStraight/.is if=TkzGasTubeStraight,
  bunsenBurner/.is if=TkzBunsenBurner,
  label/.is if=TkzLabel,
  magneticStirrer/.is if=TkzMagneticStirrer,
  TapOpen/.is if=TkzTapOpen,
  tiltAngle/.store in=\tkzTiltAngle,
  liquidLevel1/.store in=\tkzLiquidLevelI,
  liquidLevel2/.store in=\tkzLiquidLevelII,
  liquidLevel3/.store in=\tkzLiquidLevelIII,
  liquidStyle1/.store in=\tkzLiquidStyleI,
  liquidStyle2/.store in=\tkzLiquidStyleII,
  liquidStyle3/.store in=\tkzLiquidStyleIII,
  substance/.store in=\tkzSubstance,
  solid/.store in=\tkzSolid,
  Number/.store in=\tkzNumber,
  glassType/.code={\edef\tkzGlassType{\tkzGlassTypeLookup{#1}}\def\tkzGlassTypeName{#1}},
  buretteColor/.store in=\tkzBuretteColor,
  buretteLevel/.store in=\tkzBuretteLevel,
  buretteReagent/.store in=\tkzBuretteReagent,
  beakerReagent/.store in=\tkzBeakerReagent,
  reagentColor/.store in=\tkzReagentColor,
  reagentLevel/.store in=\tkzReagentLevel,
  DistillateColor/.store in=\tkzDistillateColor,
  MixtureStyle/.store in=\tkzMixtureStyle,
}

% --- 默认值 (对应原 \psset[pst-labo]{...}) -------------------
\newcommand{\tikzchemlabdefaults}{%
  \TkzBurettetrue
  \TkzTitrationPHMeterfalse
  \TkzTestTubeStopperfalse
  \TkzTestTubeStraightTubefalse
  \TkzTestTubeElbowTubefalse
  \TkzTestTubeUTubefalse
  \TkzTestTubeUTubeLongfalse
  \TkzTestTubeCurvedTubefalse
  \TkzTestTubeCurvedTubeShortfalse
  \TkzTestTubeDoubleTubefalse
  \TkzTestTubeCondenserBulbsfalse
  \TkzTestTubeClampfalse
  \TkzBubblingfalse
  \TkzGasCollectionfalse
  \TkzTubeOnlyfalse
  \TkzGasTubeStraightfalse
  \TkzBunsenBurnertrue
  \TkzLabelfalse
  \TkzMagneticStirrertrue
  \TkzTapOpenfalse
  \def\tkzTiltAngle{0}%
  \def\tkzLiquidLevelI{50}%
  \def\tkzLiquidLevelII{0}%
  \def\tkzLiquidLevelIII{0}%
  \def\tkzLiquidStyleI{liquidStyle1}%
  \def\tkzLiquidStyleII{liquidStyle2}%
  \def\tkzLiquidStyleIII{liquidStyle3}%
  \def\tkzSubstance{}%
  \def\tkzSolid{}%
  \def\tkzNumber{}%
  \def\tkzGlassType{0}%
  \def\tkzGlassTypeName{tube}%
  \def\tkzBuretteColor{VeryPaleOrange}%
  \def\tkzBuretteLevel{20}%
  \def\tkzBuretteReagent{}%
  \def\tkzBeakerReagent{}%
  \def\tkzReagentColor{PaleOrange}%
  \def\tkzReagentLevel{2}%
  \def\tkzDistillateColor{yellow}%
  \def\tkzMixtureStyle{BlueDiffusion}%
}

% ============================================================
% 玻璃器皿路径版本 (供 \clip 使用, 不带 \draw/\fill)
% ============================================================

% 滴定管路径 (burette corps)
\newcommand{\tkzBuretteBodyPath}{%
  (-0.5,11.5) -- (-0.25,11) -- (-0.25,3) -- (0,2.75) -- (0,0.5)
  -- (0.2,0.7) -- (0.2,2.75) -- (0.5,3) -- (0.5,11) -- (0.75,11.5) -- cycle
}

% 移液管路径 (pipette corps + 液滴)
\newcommand{\tkzPipetteBodyPath}{%
  (-0.05,0) -- (-0.15,0.2) -- (-0.15,6)
  -- (0.15,6) -- (0.15,0.2) -- (0.05,0) -- cycle
  (0,0) .. controls (0.25,-.4) and (-.25,-.4) .. (0,0)
}

% 试管路径 (TubeEssais corps)
\newcommand{\tkzTestTubeBodyPath}{%
  (-0.5,3.5) -- (-0.5,0.5)
  arc[start angle=180,end angle=360,radius=0.5]
  -- (0.5,3.5) -- cycle
}

% 烧杯路径 (Becher corps) - 显式圆弧, 与 \tkzBeakerBody 一致
\newcommand{\tkzBeakerBodyPath}{%
  (-1,2.5) -- (-1,0.5)
  -- ([shift={(-0.5,0.5)}]180:0.5) arc[start angle=180,end angle=270,radius=0.5]
  -- (0.5,0)
  -- ([shift={(0.5,0.5)}]270:0.5) arc[start angle=270,end angle=360,radius=0.5]
  -- (1,2.5) -- cycle
}

% 锥形瓶路径 (Erlenmeyer corps) - 显式圆弧, 与 \tkzErlenmeyerBody 一致
\newcommand{\tkzErlenmeyerBodyPath}{%
  (-0.5,3.5) -- (-0.5,2.5)
  -- (-1.3354,0.4114)
  -- ([shift={(-1.0569,0.3)}]158.20:0.3) arc[start angle=158.20,end angle=270,radius=0.3]
  -- (1.0569,0)
  -- ([shift={(1.0569,0.3)}]-90:0.3) arc[start angle=-90,end angle=21.80,radius=0.3]
  -- (0.5,2.5) -- (0.5,3.5) -- cycle
}

% 试剂瓶路径 (Flacon corps)
\newcommand{\tkzReagentBottleBodyPath}{%
  ([shift={(-0.6,3.4)}]90:0.1) arc[start angle=90,end angle=0,radius=0.1]
  -- (-0.5,3.4) -- (-0.5,3.1)
  -- ([shift={(-0.6,3.1)}]0:0.1) arc[start angle=0,end angle=-90,radius=0.1]
  -- ([shift={(-0.6,2.6)}]90:0.4) arc[start angle=90,end angle=180,radius=0.4]
  -- (-1,2.6) -- (-1,0.2)
  -- ([shift={(-.8,0.2)}]180:0.2) arc[start angle=180,end angle=270,radius=0.2]
  -- (-.8,0) -- (.8,0)
  -- ([shift={(.8,0.2)}]-90:0.2) arc[start angle=-90,end angle=0,radius=0.2]
  -- (1,0.2) -- (1,2.6)
  -- ([shift={(0.6,2.6)}]0:0.4) arc[start angle=0,end angle=90,radius=0.4]
  -- ([shift={(0.6,3.1)}]-90:0.1) arc[start angle=-90,end angle=-180,radius=0.1]
  -- (0.5,3.1) -- (0.5,3.4)
  -- ([shift={(0.6,3.4)}]180:0.1) arc[start angle=180,end angle=90,radius=0.1]
  -- cycle
}

% 容量瓶路径 (FioleJauge)
\newcommand{\tkzVolumetricFlaskPath}{%
  (-0.21,3.5) -- (-0.1,3.5) -- (-0.1,1.5)
  -- (-0.75,0) -- (0.75,0) -- (0.1,1.5) -- (0.1,3.5) -- (0.21,3.5) -- cycle
}

% 玻璃杯路径 (Verre corps)
\newcommand{\tkzWatchGlassBodyPath}{%
  (-1,3.5) -- (-0.5,1)
  arc[start angle=180,end angle=360,radius=0.5]
  -- (1,3.5) -- cycle
}

% 量筒路径 (Eprouvette corps)
\newcommand{\tkzGraduatedCylinderBodyPath}{%
  (-0.5,5) -- (-0.5,0) -- (0.5,0) -- (0.5,4.9) -- (0.6,4.93) -- cycle
}

% 统一调度 : 根据 \tkzGlassType 选择 clip 路径
\newcommand{\tkzGlassClipPath}{%
  \ifcase\tkzGlassType
    \tkzTestTubeBodyPath%        0 tube
  \or
    \tkzRoundFlaskSimplePath%          1 roundFlask
  \or
    \tkzBeakerBodyPath%           2 beaker
  \or
    \tkzErlenmeyerBodyPath%            3 erlenmeyer
  \or
    \tkzReagentBottleBodyPath%           4 reagentBottle
  \or
    \tkzVolumetricFlaskPath%            5 volumetricFlask
  \or
    \tkzWatchGlassBodyPath%            6 watchGlass
  \else
    \tkzTestTubeBodyPath%
  \fi
}

% 各 glassType 对应的 clip 圆角值 (与 \tkzGlassDraw 中各器皿的圆角保持一致,
% 否则液体会溢出实际绘制的轮廓外)
% 通过 \edef 预先展开为纯数字, 避免在 \clip[...] 选项中嵌套 \ifcase
% 注 : beaker(2) 和 erlenmeyer(3) 的路径已用显式圆弧, 故此处用 0pt
\newcommand{\tkzGlassClipOptions}{0pt}
\newcommand{\tkzGlassClipOptionsSet}{%
  \ifcase\tkzGlassType
    \edef\tkzGlassClipOptions{0pt}%       0 tube
  \or
    \edef\tkzGlassClipOptions{0pt}%       1 roundFlask
  \or
    \edef\tkzGlassClipOptions{0pt}%       2 beaker (路径自带圆弧)
  \or
    \edef\tkzGlassClipOptions{0pt}%       3 erlenmeyer  (路径自带圆弧)
  \or
    \edef\tkzGlassClipOptions{0pt}%       4 reagentBottle
  \or
    \edef\tkzGlassClipOptions{0.1cm}%     5 volumetricFlask
  \or
    \edef\tkzGlassClipOptions{0pt}%       6 watchGlass
  \else
    \edef\tkzGlassClipOptions{0pt}%
  \fi
}

% 统一调度 : 根据 \tkzGlassType 绘制对应玻璃器皿轮廓
\newcommand{\tkzGlassDraw}{%
  \ifcase\tkzGlassType
    \tkzTestTubeBody%                                0 tube
  \or
    \tkzRoundFlaskSimple%                                   1 roundFlask
  \or
    \tkzBeakerBody\relax\tkzBeakerRim%              2 beaker
  \or
    \tkzErlenmeyerBody\relax\filldraw[fill=white,draw=black,line width=\chemlabEllipseLineWidth,line cap=round] (0,3.5) ellipse (0.6 and 0.1);% 3 erlenmeyer
  \or
    \tkzReagentBottleBody\relax\filldraw[fill=white,draw=black,line width=\chemlabEllipseLineWidth,line cap=round] (0,3.5) ellipse (0.6 and 0.1);% 4 reagentBottle
  \or
    \tkzVolumetricFlask%                                     5 volumetricFlask
  \or
    \tkzWatchGlassBody\relax\tkzWatchGlassBase%                  6 watchGlass
  \fi
}

% ============================================================
% 液体多边形 helper (对应 \pst@Eprouvette@Liquide / \pst@TubeEssais@Liquide)
% ============================================================
% 参数 : #1=niveau(0-100), #2=style, #3=angle, #4=液面比例高度(3 或 4.5)
\newcommand{\tkzLiquidFill}[4]{%
  \pgfmathsetmacro{\liquidY}{#4 * #1 / 100 * cos(#3)}%
  \fill[#2] (-4,\liquidY) -- (-4,-1) -- (5,-1) -- (5,\liquidY) -- cycle;
}

% 试管液体 (高度 3)
\newcommand{\tkzTestTubeLiquid}[3]{\tkzLiquidFill{#1}{#2}{#3}{3}}
% 量筒液体 (高度 4.5)
\newcommand{\tkzGraduatedCylinderLiquid}[3]{\tkzLiquidFill{#1}{#2}{#3}{4.5}}

% ============================================================
% 1. \tkzTitration  滴定装置
% ============================================================
\newcommand{\tkzTitration}[1][]{%
  % 保存全局参数 (匹配 pst-labo \psset 的全局继承行为)
  \edef\chemlab@savedGlassType{\tkzGlassType}%
  \edef\chemlab@savedGlassTypeName{\tkzGlassTypeName}%
  \edef\chemlab@savedLiquidLevelI{\tkzLiquidLevelI}%
  \tikzchemlabdefaults
  \pgfkeys{/tikz-chemlab/.cd, #1}%
  % 恢复全局参数 (如果命令选项未指定)
  \ifnum\tkzGlassType=0 \edef\tkzGlassType{\chemlab@savedGlassType}\edef\tkzGlassTypeName{\chemlab@savedGlassTypeName}\fi
  \ifnum\tkzLiquidLevelI=50\edef\tkzLiquidLevelI{\chemlab@savedLiquidLevelI}\fi
  % --- 方案 B 局部修复: 忠实复刻 pst-labo \tkzRoundFlask 的 fallback 行为 ---
  % pst-labo.tex L478-522: \tkzTitration 内部调用 \tkzRoundFlask, 其 \ifcase\psk@glassType
  %   0(tube)/1(roundFlask) -> \pst@roundFlask
  %   2(beaker)         -> \pst@Becher@Corps
  %   3(erlenmeyer)          -> \pst@Erlen@Corps
  %   4(reagentBottle)         -> \pst@Flacon@Corps
  %   5(volumetricFlask)/6(Verre)/未知 -> \pst@roundFlask (fallback, L485 注释 "roundFlask is default")
  % tikz-chemlab 的 \tkzGlassDraw 对 glassType=0 画的是 \tkzTestTubeBody (试管),
  % 与原版不一致. 此处仅在 \tkzTitration 内部把 tube/volumetricFlask/Verre/未知 重映射为 roundFlask,
  % 使 \tkzGlassDraw/\tkzGlassClipPath 走 roundFlask 分支. 不影响其他命令.
  \ifnum\tkzGlassType=0 \def\tkzGlassType{1}\def\tkzGlassTypeName{roundFlask}\fi
  \ifnum\tkzGlassType=5 \def\tkzGlassType{1}\def\tkzGlassTypeName{roundFlask}\fi
  \ifnum\tkzGlassType=6 \def\tkzGlassType{1}\def\tkzGlassTypeName{roundFlask}\fi
  \ifnum\tkzGlassType>6 \def\tkzGlassType{1}\def\tkzGlassTypeName{roundFlask}\fi
  \begin{scope}[xshift=\chemlab@xshift]
    % 边界框
    \ifTkzBurette
      \ifTkzTitrationPHMeter
        \useasboundingbox (-5,0) rectangle (3,16);
      \else
        \useasboundingbox (-2,0) rectangle (2,16);
      \fi
    \else
      \ifTkzTitrationPHMeter
        \useasboundingbox (-5,0) rectangle (3,5);
      \else
        \useasboundingbox (-2,0) rectangle (2,5);
      \fi
    \fi
    % 烧瓶 + 液体 (匹配 pst-labo \rput(0,3.5){\tkzRoundFlask})
    % pst-labo: \tkzRoundFlask 的 pspicture bbox = (-2,0)(2,4), 中心 (0,2)
    % \rput(0,3.5) 使用 ref=c (中心对齐), 所以局部原点 (瓶底) 位于世界 (0, 3.5-2) = (0, 1.5)
    % 锥形瓶底部 (y=0 local) 正好放在搅拌器台面上 (y=1.5 world)
    \begin{scope}[shift={(0,1.5)}]
      % 液体 (clip 到玻璃内部 ∩ 液体矩形)
      \begin{scope}
        \tkzGlassClipOptionsSet
        \clip[rounded corners=\tkzGlassClipOptions] \tkzGlassClipPath;
        \ifnum\tkzLiquidLevelI>0
          \pgfmathsetmacro{\chemlab@liquidClipY}{3 * \tkzLiquidLevelI / 100 * cos(\tkzTiltAngle)}%
          \clip (-4,-1.5) rectangle (5,\chemlab@liquidClipY);
        \fi
        \ifnum\tkzLiquidLevelI>0
          \tkzTestTubeLiquid{\tkzLiquidLevelI}{\tkzLiquidStyleI}{\tkzTiltAngle}
        \fi
        \tkzSubstance
        \tkzSolid
      \end{scope}
      % 玻璃器皿轮廓 (根据 glassType 选择: erlenmeyer/beaker/roundFlask 等)
      \tkzGlassDraw
    \end{scope}
    % 滴定管 + 液体
    \ifTkzBurette
      \begin{scope}[shift={(0,4.5)}]
        % 液体填充 (clip 到 burette 内部)
        \begin{scope}
          \clip \tkzBuretteBodyPath;
          \pgfmathsetmacro{\buretteY}{\tkzBuretteLevel * 0.32 + 2.68}
          \fill[\tkzBuretteColor] (-1,-1) rectangle (1,\buretteY);
        \end{scope}
        \tkzBuretteBody
        \tkzBuretteStopcock
        \tkzBuretteGraduation
      \end{scope}
    \fi
    % pH 计
    \ifTkzTitrationPHMeter
      \begin{scope}[shift={(-3,0)}]\tkzTitrationPHMeter\end{scope}
    \fi
    % 搅拌子 + 支架
    \ifTkzMagneticStirrer
      \begin{scope}[shift={(0,1.5)}]\tkzTitrationStirBar\end{scope}
      \begin{scope}[shift={(0,0)}]\tkzTitrationStand\end{scope}
    \else
      % pst-labo: \psset{linewidth=0.053}\psframe(-1.5,0)(1.5,1.5)
      % linewidth=0.053 是用户单位 (×unit), 随 scale 缩放
      \chemlabUpdateLineWidths
      \draw[line width=\chemlabLineWidth] (-1.5,0) rectangle (1.5,1.5);
    \fi
    % 反应物标签
    \ifTkzBurette
      \node[anchor=west] at (2,11) {\tkzBuretteReagent};
    \fi
    \node[anchor=west] at (2,3.5) {\tkzBeakerReagent};
  \end{scope}
  % 自动推进原点 (匹配 pst-labo \pspicture 的布局宽度)
  \ifTkzTitrationPHMeter
    \advance\chemlab@xshift by 8cm %  width = 8
  \else
    \advance\chemlab@xshift by 4cm %  width = 4
  \fi
}

% ============================================================
% 2. \tkzPipette  移液管
% ============================================================
\newcommand{\tkzPipette}[1][]{%
  \begingroup
  \tikzchemlabdefaults
  \pgfkeys{/tikz-chemlab/.cd, #1}%
  \begin{scope}
    \useasboundingbox (-1,-1) rectangle (1,6);
    % 液体
    \begin{scope}
      \clip \tkzPipetteBodyPath;
      \pgfmathsetmacro{\pipetteY}{\tkzReagentLevel / 2}
      \fill[\tkzReagentColor] (-1,-2) rectangle (1,\pipetteY);
    \end{scope}
    \tkzPipetteBody
    \tkzDrop
  \end{scope}
  \endgroup
}

% ============================================================
% 3. \tkzLargePipette  大移液管
% ============================================================
\newcommand{\tkzLargePipette}[1][]{%
  \begingroup
  \tikzchemlabdefaults
  \pgfkeys{/tikz-chemlab/.cd, #1}%
  \begin{scope}
    \useasboundingbox (-1,-1) rectangle (1,6);
    % 注 : 原文 \pst@pipette 是已注释的代码, 此处使用 bigpipette 形状
    \begin{scope}
      % bigpipette 形状比较简单, 直接 clip 一个矩形区域近似
      \clip (0,0) -- (0,1) -- (-0.23,2.5) -- (-0.23,22.5)
            arc[start angle=-80,end angle=80,radius=0.45] -- cycle;
      \pgfmathsetmacro{\pipetteY}{\tkzReagentLevel / 2}
      \fill[\tkzReagentColor] (-1,-2) rectangle (1,\pipetteY);
    \end{scope}
    \tkzLargePipetteBody
  \end{scope}
  \endgroup
}

% ============================================================
% 4. \tkzGraduatedCylinder  量筒
% ============================================================
\newcommand{\tkzGraduatedCylinder}[1][]{%
  \begingroup
  \tikzchemlabdefaults
  \pgfkeys{/tikz-chemlab/.cd, #1}%
  \begin{scope}
    % 注: pst-labo 用 \pspicture(-2,0)(2,6) 设置 bbox 用于 \rput 放置, 但内容可超出 bbox
    % 在 TikZ 中, useasboundingbox 会覆盖整个 picture 的 bbox, 不适合嵌套使用
    % 因此这里不设 useasboundingbox, 让外层命令 (如 Figure 中的 pspicture) 控制 bbox
    % standalone 编译单独的 \tkzGraduatedCylinder 时, 自动检测内容 bbox
    % 液体 (clip 到量筒内部, 旋转 tiltAngle)
    \begin{scope}[rotate=\tkzTiltAngle]
      \begin{scope}
        \clip \tkzGraduatedCylinderBodyPath;
        \ifnum\tkzLiquidLevelI>0
          \tkzGraduatedCylinderLiquid{\tkzLiquidLevelI}{\tkzLiquidStyleI}{\tkzTiltAngle}
        \fi
        \ifnum\tkzLiquidLevelII>0
          \tkzGraduatedCylinderLiquid{\tkzLiquidLevelII}{\tkzLiquidStyleII}{\tkzTiltAngle}
        \fi
        \ifnum\tkzLiquidLevelIII>0
          \tkzGraduatedCylinderLiquid{\tkzLiquidLevelIII}{\tkzLiquidStyleIII}{\tkzTiltAngle}
        \fi
        \tkzSubstance
      \end{scope}
      \tkzGraduatedCylinderBody
      \tkzGraduatedCylinderBase
      \tkzGraduatedCylinderGraduation
      \begin{scope}[shift={(0,5)}]\tkzGraduatedCylinderRim\end{scope}
    \end{scope}
  \end{scope}
  \endgroup
}

% ============================================================
% 5. \tkzDistillation  蒸馏装置
% ============================================================
\newcommand{\tkzDistillation}[1][]{%
  \begingroup
  \tikzchemlabdefaults
  \pgfkeys{/tikz-chemlab/.cd, #1}%
  \begin{scope}
    \useasboundingbox (-4,-10) rectangle (8,7);
    \begin{scope}[shift={(-3,-7)}]\tkzDistillationFractionated\end{scope}
  \end{scope}
  \endgroup
}

% ============================================================
% 6. \tkzFunnel  漏斗 + 试管
% ============================================================
\newcommand{\tkzFunnel}[1][]{%
  \begingroup
  \tikzchemlabdefaults
  \pgfkeys{/tikz-chemlab/.cd, #1}%
  % 设置 noDefaults=true, 使 \tkzTestTube 继承用户参数
  \chemlab@noDefaultstrue
  \begin{scope}
    % pst-labo: \rput(0,0){\tkzTestTube} 用 pspicture bbox CENTER 对齐
    % tikz: shift = rput_pos - bbox_center (in local coords)
    % pspicture bbox 取决于 glassType 和 tiltAngle:
    %   beaker (glassType=2), tiltAngle=0: (-0.5,0)(1,2.8), center=(0.25,1.4)
    %   tiltAngle<0: (-1,0)(3,4), center=(1,2)
    %   tiltAngle>0: (-3,0)(1,4), center=(-1,2)
    %   其他 (tiltAngle=0, glassType!=2): (-1.5,0)(1.5,4), center=(0,2)
    \ifnum\tkzTiltAngle=0
      \ifnum\tkzGlassType=2
        % beaker, 直立: center=(0.25,1.4)
        \begin{scope}[shift={(-0.25,-1.4)}]\tkzTestTube[noDefaults=true]\end{scope}
      \else
        % roundFlask/reagentBottle/erlenmeyer, 直立: center=(0,2)
        \begin{scope}[shift={(0,-2)}]\tkzTestTube[noDefaults=true]\end{scope}
      \fi
    \else
      \ifnum\tkzTiltAngle>0
        % 倾斜>0: center=(-1,2)
        \begin{scope}[shift={(1,-2)}]\tkzTestTube[noDefaults=true]\end{scope}
      \else
        % 倾斜<0: center=(1,2)
        \begin{scope}[shift={(-1,-2)}]\tkzTestTube[noDefaults=true]\end{scope}
      \fi
    \fi
    % 漏斗: pst-labo \rput(0,1.5){\pst@Entonnoir@Corps}
    % \pst@Entonnoir@Corps 路径 bbox: x=[-1,1], y=[0,2.6] (含椭圆顶 2.5+0.1)
    % bbox center = (0, 1.3)
    % \rput(0,1.5) ref=c: center 置于 (0,1.5)
    % shift = rput_pos - center = (0, 1.5-1.3) = (0, 0.2)
    % 即管尖 (路径原点 0,0) 位于 (0, 0.2)
    \begin{scope}[shift={(0,0.2)}]\tkzFunnelBody\end{scope}
  \end{scope}
  \chemlab@noDefaultsfalse
  \endgroup
}

% ============================================================
% 7. \tkzTestTube  试管 (核心命令)
% ============================================================
\newcommand{\tkzTestTube}[1][]{%
  \begingroup
  \ifchemlab@noDefaults\else\tikzchemlabdefaults\fi
  \pgfkeys{/tikz-chemlab/.cd, #1}%
  \begin{scope}
    % 注: pst-labo 用 \pspicture(bbox) 设置 bbox 用于 \rput 放置, 但内容可超出 bbox
    % 在 TikZ 中, useasboundingbox 会覆盖整个 picture 的 bbox, 不适合嵌套使用
    % 因此这里不设 useasboundingbox, 让外层命令 (如 \tkzHeatTube) 控制 bbox
    % standalone 编译单独的 \tkzTestTube 时, 自动检测内容 bbox
    % 液体 (clip 到旋转的玻璃器皿内部, 但液面保持水平)
    % 原版 pst-labo: \psclip{ \rput{\angle}{<轮廓>} \psframe<液面不旋转> }
    % pst-labo 的 \psclip 产生两次裁剪 (PostScript clip 是累积的=交集):
    %   1. 裁剪到玻璃形状 (来自 \pscustom)
    %   2. 裁剪到液体矩形 (来自 \psframe(-3,-1)(5, liquid_level))
    % → 物质被裁剪到 玻璃形状 ∩ 液体矩形 = 液面以下
    % TikZ: 先旋转 clip 玻璃路径, 再反向旋转 clip 液体矩形, 然后画液体和物质
    \begin{scope}[rotate=\tkzTiltAngle]
      \begin{scope}
        \tkzGlassClipOptionsSet
        \clip[rounded corners=\tkzGlassClipOptions] \tkzGlassClipPath;
        % 液面: 在旋转的 clip 内, 但反向旋转使液面水平
        % 液面高度需乘 cos(tiltAngle) (原版 pst-labo 的行为)
        \begin{scope}[rotate=-\tkzTiltAngle]
          % 第二次裁剪: 液体矩形 (匹配 pst-labo \psframe 裁剪)
          % 只在有液体时才裁剪; 无液体时不裁剪 (物质可见到玻璃顶部)
          \ifnum\tkzLiquidLevelI>0
            \pgfmathsetmacro{\chemlab@liquidClipY}{3 * \tkzLiquidLevelI / 100 * cos(\tkzTiltAngle)}%
            \clip (-4,-1.5) rectangle (5,\chemlab@liquidClipY);
          \fi
          \ifnum\tkzLiquidLevelI>0
            \tkzTestTubeLiquid{\tkzLiquidLevelI}{\tkzLiquidStyleI}{\tkzTiltAngle}
          \fi
          \ifnum\tkzLiquidLevelII>0
            \tkzTestTubeLiquid{\tkzLiquidLevelII}{\tkzLiquidStyleII}{\tkzTiltAngle}
          \fi
          \ifnum\tkzLiquidLevelIII>0
            \tkzTestTubeLiquid{\tkzLiquidLevelIII}{\tkzLiquidStyleIII}{\tkzTiltAngle}
          \fi
          \tkzSubstance
          \tkzSolid
        \end{scope}
      \end{scope}
      % 标签
      \ifTkzLabel
        \node[fill=VeryLightGray,draw=none,inner sep=2pt] at (0,2) {\tkzNumber};
      \fi
      % 玻璃器皿轮廓
      \tkzGlassDraw
      % 顶部附件 (straightTube / elbowTube / uTube / stopper)
      \ifTkzTestTubeStraightTube
        \begin{scope}[shift={(0,3.5)}]\tkzTestTubeStraightTube\end{scope}
      \else
        \ifTkzTestTubeElbowTube
          \begin{scope}[shift={(0,3.5)}]\tkzTestTubeElbowTube\end{scope}
        \else
          \ifTkzTestTubeUTube
            \begin{scope}[shift={(0,3.5)}]\tkzTestTubeUTube\end{scope}
          \else
            \ifTkzTestTubeStopper
              \begin{scope}[shift={(0,3.5)}]\tkzTestTubeStopper\end{scope}
            \fi
          \fi
        \fi
      \fi
      % 试管夹
      \ifTkzTestTubeClamp
        \begin{scope}[shift={(0,3.2)}]\tkzTestTubeClamp\end{scope}
      \fi
    \end{scope}
  \end{scope}
  \endgroup
}

% ============================================================
% 8. \tkzHeatTube  加热试管
% ============================================================
% pstChauffageTube: 保存用户参数, 传递给 \tkzTestTube
% pst-labo 用 \psset 全局继承; tikz-chemlab 用 \tkzTestTube[noDefaults] 跳过重置
\newif\ifchemlab@noDefaults
\newcommand{\tkzHeatTube}[1][]{%
  \begingroup
  \tikzchemlabdefaults
  \pgfkeys{/tikz-chemlab/.cd, #1}%
  \chemlab@noDefaultstrue
  \begin{scope}
    % pst-labo bbox: tubeOnly=(-2,-5)(5,4), 其他=(-3,-5)(10,4)
    % 不使用 \useasboundingbox — pst-labo 的 pspicture 不裁剪内容,
    % 内容可以超出 bbox (如 elbowTube 弯管延伸到 x>5)
    % 让 TikZ 自动检测完整内容 bbox
    \ifTkzBubbling
      % 上方试管/烧瓶: 继承用户设置的 glassType, 倾斜-60°, 有弯管
      \edef\chemlab@tmpGlassType{\tkzGlassTypeName}%
      % pst-labo: \rput(2.5,2.5){\tkzTestTube[elbowTube=true,tiltAngle=-60]}
      % elbowTube=true → bbox (-1,0)(6,4) center (2.5,2)
      % → 局部原点 (瓶底) 世界 (2.5-2.5, 2.5-2) = (0, 0.5)
      \begin{scope}[shift={(0,0.5)}]
        \tkzTestTube[elbowTube=true,tiltAngle=-60,noDefaults=true]
      \end{scope}
      % 下方试管: pst-labo \rput(8.2,-3){\tkzTestTube[tiltAngle=30,...]}
      % tiltAngle>0 → bbox (-3,0)(1,4) center (-1,2)
      % → 局部原点 (试管底) 世界 (8.2-(-1), -3-2) = (9.2, -5)
      \begin{scope}[shift={(9.2,-5)}]
        \tkzTestTube[liquidLevel1=70,glassType=tube,elbowTube=false,
                       tiltAngle=30,
                       liquidStyle1=Champagne,
                       substance={\tkzChampagneBubbles[10]},clamp=false]
      \end{scope}
    \else
      % 非barbotage: 上方器皿继承 glassType
      \edef\chemlab@tmpGlassType{\tkzGlassTypeName}%
      \ifcase\tkzGlassType
        % glassType=tube
        \ifTkzTubeOnly
          % tubeOnly + tube: pst-labo \rput(1.5,2.5)
          % tiltAngle<0 → bbox (-1,0)(3,4) center (1,2) → 世界 (0.5, 0.5)
          \begin{scope}[shift={(0.5,0.5)}]\tkzTestTube[tiltAngle=-60,noDefaults=true]\end{scope}
        \else
          \ifTkzTestTubeElbowTube
            % elbowTube: pst-labo \rput(2,2.5)
            % elbowTube=true → bbox (-1,0)(6,4) center (2.5,2) → 世界 (-0.5, 0.5)
            \begin{scope}[shift={(-0.5,0.5)}]\tkzTestTube[tiltAngle=-60,noDefaults=true]\end{scope}
          \else
            % 无tubeCoude: pst-labo \rput(1,2.5)
            % tiltAngle<0 → bbox (-1,0)(3,4) center (1,2) → 世界 (0, 0.5)
            \begin{scope}[shift={(0,0.5)}]\tkzTestTube[tiltAngle=-60,noDefaults=true]\end{scope}
          \fi
        \fi
      \or
        % glassType=roundFlask (无论 tubeOnly 与否, pst-labo 都是 \rput(2,2.5))
        % tiltAngle<0 → bbox (-1,0)(3,4) center (1,2) → 局部原点世界 (2-1, 2.5-2) = (1, 0.5)
        \begin{scope}[shift={(1,0.5)}]\tkzTestTube[tiltAngle=-60,noDefaults=true]\end{scope}
      \else
        % 其他 glassType: pst-labo \rput(1,2.5) → 世界 (0, 0.5)
        \begin{scope}[shift={(0,0.5)}]\tkzTestTube[tiltAngle=-60,noDefaults=true]\end{scope}
      \fi
    \fi
    \ifTkzBunsenBurner
      % pst-labo: \rput(1,-5){\pst@BecBunsen} (无 pspicture, 实测原点对齐)
      % 灯底位于世界 (1, -5)
      \begin{scope}[shift={(1,-5)}]\tkzBunsenBurner\end{scope}
    \fi
  \end{scope}
  \chemlab@noDefaultsfalse
  \endgroup
}

% ============================================================
% 9. \tkzHeatFlask  加热烧瓶
% ============================================================
\newcommand{\tkzHeatFlask}[1][]{%
  \tikzchemlabdefaults
  \pgfkeys{/tikz-chemlab/.cd, #1}%
  % pst-labo: 每个 \tkzHeatFlask 有独立 pspicture, 不跨实例裁剪
  % tikz-chemlab: 不设全局 clip, 让所有内容自由绘制
  % pst-labo: \tkzRoundFlask 中 glassType=0 (tube) 也画 \pst@roundFlask
  % tikz-chemlab: 强制 glassType=0 → 1 (roundFlask)
  \begin{scope}[xshift=\chemlab@xshift]
    % 不设 \useasboundingbox, 让 TikZ 自动检测完整内容 bbox
    % 物体间相对位置由 shift 控制, 与 pst-labo 一致
    \ifTkzBunsenBurner\def\cnth{-5}\else\def\cnth{1}\fi
    % 气体回收装置
    \ifTkzGasCollection
      \ifTkzTestTubeCurvedTube
        \begin{scope}[shift={(8,-5)}]
          \tkzCrystallizingDish\relax\tkzTestTubeRack
          \begin{scope}[shift={(0,0.7)}]\tkzInvertedTube\relax\tkzBubbles\end{scope}
        \end{scope}
      \else\ifTkzTestTubeCurvedTubeShort
        \TkzBunsenBurnerfalse
        \begin{scope}[shift={(6,1)}]
          \tkzCrystallizingDish\relax\tkzTestTubeRack
          \begin{scope}[shift={(0,0.7)}]\tkzInvertedTube\relax\tkzBubbles\end{scope}
        \end{scope}
      \else\ifTkzTestTubeDoubleTube
        \TkzBunsenBurnerfalse
        \begin{scope}[shift={(6.2,1)}]
          \tkzCrystallizingDish\relax\tkzTestTubeRack
          \begin{scope}[shift={(0,0.7)}]\tkzInvertedTube\relax\tkzBubbles\end{scope}
        \end{scope}
      \fi\fi\fi
    \fi
    % 烧瓶 + 液体 + 导管
    % pst-labo: \rput(1,3){\tkzRoundFlask} 默认使用 bbox 中心对齐 (ref=c)
    % \tkzRoundFlask 的 pspicture bbox = (-2,0)(2,4), 中心 (0,2)
    % 所以局部原点 (瓶底) 实际位于世界 (1, 3-2) = (1, 1)
    \begin{scope}[shift={(1,1)}]
      % 液体 (clip 到烧瓶内部 ∩ 液体矩形, 匹配 pst-labo \psclip 双重裁剪)
      \begin{scope}
        \tkzGlassClipOptionsSet
        \clip[rounded corners=\tkzGlassClipOptions] \tkzGlassClipPath;
        % 第二次裁剪: 液体矩形 (匹配 pst-labo \psframe 裁剪)
        % 只在有液体时才裁剪; 无液体时不裁剪 (物质可见到玻璃顶部)
        \ifnum\tkzLiquidLevelI>0
          \pgfmathsetmacro{\chemlab@liquidClipY}{3 * \tkzLiquidLevelI / 100 * cos(\tkzTiltAngle)}%
          \clip (-4,-1.5) rectangle (5,\chemlab@liquidClipY);
        \fi
        \ifnum\tkzLiquidLevelI>0
          \tkzTestTubeLiquid{\tkzLiquidLevelI}{\tkzLiquidStyleI}{\tkzTiltAngle}
        \fi
        \tkzSubstance
        \tkzSolid
      \end{scope}
      % 玻璃器皿轮廓 (根据 glassType 选择: roundFlask/erlenmeyer/beaker 等)
      \tkzGlassDraw
      % 导管 (在瓶口 y=3.5 处) — 匹配 pst-labo \tkzRoundFlask 的全部导管类型
      \ifTkzTestTubeStraightTube
        \begin{scope}[shift={(0,3.5)}]\tkzTestTubeStraightTube\end{scope}
      \else\ifTkzTestTubeElbowTube
        \begin{scope}[shift={(0,3.5)}]\tkzTestTubeElbowTube\end{scope}
      \else\ifTkzTestTubeUTube
        \begin{scope}[shift={(0,3.5)}]\tkzTestTubeUTube\end{scope}
      \else\ifTkzTestTubeUTubeLong
        \begin{scope}[shift={(0,3.5)}]\tkzTestTubeUTubeLong\end{scope}
      \else\ifTkzTestTubeCurvedTube
        \begin{scope}[shift={(0,3.5)}]\tkzCurvedTube\end{scope}
      \else\ifTkzTestTubeCurvedTubeShort
        \begin{scope}[shift={(0,3.5)}]\tkzCurvedTubeShort\end{scope}
      \else\ifTkzTestTubeStopper
        \begin{scope}[shift={(0,3.5)}]\tkzTestTubeStopper\end{scope}
      \else\ifTkzTestTubeDoubleTube
        \begin{scope}[shift={(0,3.5)}]\tkzTestTubeDoubleTube\end{scope}
      \fi\fi\fi\fi\fi\fi\fi\fi
    \end{scope}
    \ifTkzBubbling
      % pst-labo: \rput(5.3,-3){\tkzTestTube[...]}, tiltAngle=0 默认, glassType=tube
      % bbox (-1.5,0)(1.5,4) center (0,2) → 局部原点 (试管底) 世界 (5.3, -5)
      \begin{scope}[shift={(5.3,-5)}]
        \tkzTestTube[liquidLevel1=70,glassType=tube]
      \end{scope}
    \fi
    \ifTkzBunsenBurner
      % pst-labo: \rput(1,-5){\pst@BecBunsenGrille} (无 pspicture, 实测原点对齐)
      \begin{scope}[shift={(1,-5)}]\tkzBunsenBurnerGrill\end{scope}
    \fi
  \end{scope}
  % 自动推进原点 (匹配 pst-labo \pspicture 的布局宽度)
  % 使用全局 dimen 累加, scope 的 xshift 会正确应用
  % pst-labo bbox 宽度: curvedTubeShort/DoubleTube = 11, 其他 = 6
  \ifTkzTestTubeUTubeLong
    % uTubeLong: 不推进
  \else\ifTkzTestTubeUTube
    % uTube: 不推进
  \else\ifTkzTestTubeCurvedTubeShort
    \advance\chemlab@xshift by 11cm %  bbox (-2,0)(9,8), width = 11
  \else\ifTkzTestTubeDoubleTube
    \advance\chemlab@xshift by 11cm %  bbox (-2,0)(9,8), width = 11
  \else
    \advance\chemlab@xshift by 6cm %  bbox (-2,\cnth)(4,6), width = 6
  \fi\fi\fi\fi
}

% ============================================================
% 10. \tkzRoundFlask  烧瓶 (高级命令)
% ============================================================
\newcommand{\tkzRoundFlask}[1][]{%
  \tikzchemlabdefaults
  \pgfkeys{/tikz-chemlab/.cd, #1}%
  % pst-labo: \tkzRoundFlask 中 glassType=0 (tube) 也画 \pst@roundFlask (case 0 = roundFlask)
  % tikz-chemlab: 强制 glassType=0 → 1 (roundFlask), 避免画试管轮廓
  \begin{scope}[xshift=\chemlab@xshift]
    % 不设 \useasboundingbox, 让 TikZ 自动检测完整内容 bbox
    % bbox 坐标使用局部坐标系, 会随 scale 自动缩放
    % 物体间相对位置由 shift 控制, 与 pst-labo 一致
    % 液体
    \begin{scope}[rotate=\tkzTiltAngle]
      \begin{scope}
        \tkzGlassClipOptionsSet
        \clip[rounded corners=\tkzGlassClipOptions] \tkzGlassClipPath;
        % 第二次裁剪: 液体矩形 (匹配 pst-labo \psframe 裁剪)
        % 只在有液体时才裁剪; 无液体时不裁剪 (物质可见到玻璃顶部)
        \ifnum\tkzLiquidLevelI>0
          \pgfmathsetmacro{\chemlab@liquidClipY}{3 * \tkzLiquidLevelI / 100 * cos(\tkzTiltAngle)}%
          \clip (-4,-1.5) rectangle (5,\chemlab@liquidClipY);
        \fi
        \ifnum\tkzLiquidLevelI>0
          \tkzTestTubeLiquid{\tkzLiquidLevelI}{\tkzLiquidStyleI}{\tkzTiltAngle}
        \fi
        \ifnum\tkzLiquidLevelII>0
          \tkzTestTubeLiquid{\tkzLiquidLevelII}{\tkzLiquidStyleII}{\tkzTiltAngle}
        \fi
        \ifnum\tkzLiquidLevelIII>0
          \tkzTestTubeLiquid{\tkzLiquidLevelIII}{\tkzLiquidStyleIII}{\tkzTiltAngle}
        \fi
        \tkzSubstance
        \tkzSolid
      \end{scope}
      % 标签
      \ifTkzLabel
        \node[fill=VeryLightGray,draw=none,inner sep=2pt] at (0,1.5) {\tkzNumber};
      \fi
      % 玻璃器皿轮廓
      \tkzGlassDraw
      % 顶部附件 (按优先级)
      \ifTkzTestTubeStraightTube
        \begin{scope}[shift={(0,3.5)}]\tkzTestTubeStraightTube\end{scope}
      \else\ifTkzTestTubeElbowTube
        \begin{scope}[shift={(0,3.5)}]\tkzTestTubeElbowTube\end{scope}
      \else\ifTkzTestTubeUTube
        \begin{scope}[shift={(0,3.5)}]\tkzTestTubeUTube\end{scope}
      \else\ifTkzTestTubeUTubeLong
        \begin{scope}[shift={(0,3.5)}]\tkzTestTubeUTubeLong\end{scope}
      \else\ifTkzTestTubeCurvedTube
        \begin{scope}[shift={(0,3.5)}]\tkzCurvedTube\end{scope}
      \else\ifTkzTestTubeCurvedTubeShort
        \begin{scope}[shift={(0,3.5)}]\tkzCurvedTubeShort\end{scope}
      \else\ifTkzTestTubeStopper
        \begin{scope}[shift={(0,3.5)}]\tkzTestTubeStopper\end{scope}
      \else\ifTkzTestTubeDoubleTube
        \begin{scope}[shift={(0,3.5)}]\tkzTestTubeDoubleTube\end{scope}
      \else\ifTkzTestTubeCondenserBulbs
        \tkzRoundFlaskReflux
        \begin{scope}[shift={(0,3.5)}]\tkzCondenserBulbs\end{scope}
        \tkzHeaterRoundFlask
      \fi\fi\fi\fi\fi\fi\fi\fi\fi
    \end{scope}
  \end{scope}
  % 自动推进原点 (匹配 pst-labo \pspicture 的布局宽度)
  % pst-labo bbox: condenserBulbs = (-2,0)(2,12), 其他 = (-2,0)(2,4), width = 4
  \advance\chemlab@xshift by 4cm %  bbox width = 4
}

% ============================================================
% 11. \tkzChampagneBubbles  香槟气泡
% ============================================================
% 原版 pst-labo 源码:
%   \setrandim{\BubbleX}{-1pt}{1pt}    -- BulleX ∈ [-1, 1] (纯数字, cm)
%   \setrandim{\BubbleY}{0pt}{2pt}     -- BulleY ∈ [0, 2] (纯数字)
%   Y 坐标 = \pointless\BubbleY * liquidLevel1/100  (乘以液位百分比)
%   半径 = {\BubbleY} (带 pt 单位, 1pt = 0.0353 cm)
%   \pst@dimA 分级放大: BulleY<1pt → ×3, 1pt≤BulleY<2pt → ×2
%
% 实测原版渲染图:
%   气泡为黑色边空心圆 (\pscircle, fillstyle=none)
%   半径 0.05~0.175 单位 (比源码计算 0~0.071 大, 可能是渲染缩放)
%   气泡团中心 (世界): (7.89, -0.975)
\newcommand{\tkzChampagneBubbles}[1][25]{%
  \foreach \iBubble in {1,...,#1}{%
    % BulleX ∈ [-1, 1] (纯数字, cm)
    \pgfmathsetmacro{\BubbleX}{-1 + 2*rnd}%
    % BulleY ∈ [0, 2] (纯数字)
    \pgfmathsetmacro{\BubbleY}{2*rnd}%
    % Y 坐标 = BulleY * liquidLevel1/100 (匹配源码逻辑)
    \pgfmathsetmacro{\BubbleYpos}{\BubbleY * \tkzLiquidLevelI / 100}%
    % dimA 分级放大 (匹配源码逻辑)
    \pgfmathsetmacro{\BubbleRand}{rnd}%
    \ifdim\BubbleRand pt<0.5pt
      % BulleY < 1pt (50%): dimA = BulleY * 3
      \pgfmathsetmacro{\BubbleR}{\BubbleY * 0.0875}%
    \else
      % 1pt <= BulleY < 2pt (50%): dimA = BulleY * 2
      \pgfmathsetmacro{\BubbleR}{\BubbleY * 0.0875}%
    \fi
    \begin{scope}[rotate=\tkzTiltAngle]
      \draw[draw=black, line width=0.6pt] (\BubbleX,\BubbleYpos) circle (\BubbleR);
    \end{scope}
  }%
}

% ============================================================
% 12. \tkzBubbles  倒置管中气泡
% ============================================================
\newcommand{\tkzBubbles}[1][25]{%
  \pgfgettransformentries{\clab@aa}{\clab@bb}{\clab@cc}{\clab@dd}{\clab@xx}{\clab@yy}%
  \pgfmathsetmacro{\clab@scale}{(\clab@aa + \clab@dd) / 2}%
  \foreach \iBubble in {1,...,#1}{%
    % pst-labo: \setrandim{BulleX}{-0.1pt}{0.1pt} → \pointless 去掉pt, 数值直接用作局部坐标
    \pgfmathsetmacro{\BubbleX}{-0.1 + 0.2*rnd}%
    \pgfmathsetmacro{\BubbleY}{0.15 + 2.85*rnd}%
    % pst-labo: dimh = BulleY * (3 if <1pt, 2 if <2pt, 1 if >=2pt)
    % dotscale = dimh, dot radius = dotscale * dotsize/2
    % dotsize default = 2pt + 2*linewidth = 2pt + 2*0.8pt = 3.6pt → radius = dimh * 1.8pt
    % 转换为局部坐标: local = pt / (72 * scale)
    \pgfmathparse{\BubbleY<1?3:(\BubbleY<2?2:1)}%
    \let\BubbleScale\pgfmathresult
    \pgfmathsetmacro{\BubbleDimh}{\BubbleY * \BubbleScale}%
    \pgfmathsetmacro{\BubbleR}{1.8 * \BubbleDimh / (72 * \clab@scale)}%
    % Y position = BulleY * 1.5 (局部坐标, 匹配 pst-labo \pointless\BubbleY 1.5 mul)
    \pgfmathsetmacro{\BubbleYpos}{\BubbleY * 1.5}%
    \begin{scope}[rotate=\tkzTiltAngle]
      % pst-labo: dotstyle=o (open circle = hollow, black outline)
      \draw[draw=black, line width=0.4pt, fill=white] (\BubbleX,\BubbleYpos) circle (\BubbleR);
    \end{scope}
  }%
}

% ============================================================
% 13. \tkzFilaments  絮状物
% ============================================================
\newcommand{\tkzFilaments}[2][5]{%
  \pgfgettransformentries{\fil@aa}{\fil@bb}{\fil@cc}{\fil@dd}{\fil@xx}{\fil@yy}%
  \pgfmathsetmacro{\fil@scale}{(\fil@aa + \fil@dd) / 2}%
  \foreach \iFil in {1,...,#1}{%
    % pst-labo: \setrandim{\pst@FilamentWidth}{0pt}{0.05pt} → \pointless = [0, 0.05)
    % linewidth=\pointless\pst@FilamentWidth 解释为 N × unit (PSTricks 规则)
    % 实际线宽 = [0, 0.05) × scale cm (scale 对应 unit)
    \pgfmathsetmacro{\fw}{0.05 * \fil@scale * rnd}%
    \pgfmathsetmacro{\myRand}{rnd}%  ONE random value per filament (匹配 pst-labo MyRand)
    \begin{scope}[rotate=\tkzTiltAngle]
      % pst-labo \pscurve 通过4个点的平滑曲线 (不是贝塞尔):
      %   P1: (MyRand-1, 2*MyRand-1)     = ([-1,0), [-1,1))
      %   P2: (MyRand+0.2, MyRand+2)     = ([0.2,1.2), [2,3))
      %   P3: (MyRand, MyRand-1)         = ([0,1), [-1,0))
      %   P4: (MyRand+0.5, MyRand+4)     = ([0.5,1.5), [4,5))
      \draw[line width=\fw cm,#2] plot[smooth,tension=0.5] coordinates {
        ({\myRand-1},{2*\myRand-1})
        ({\myRand+0.2},{\myRand+2})
        ({\myRand},{\myRand-1})
        ({\myRand+0.5},{\myRand+4})
      };
    \end{scope}
  }%
}

% ============================================================
% 14. \tkzIronNail  铁钉
% ============================================================
% 单个铁钉形状 (单位 0.5)
\newcommand{\tkzNails}{%
  % pst-labo 实际 SLW=0.71135pt=0.025cm (从 xdv 提取)
  % \psarc[unit=0.5,linewidth=0.1] 的 \psline 继承同一 linewidth
  \draw[scale=0.5,line width=0.025cm,blue]
    ([shift={(0,0.5)}]190:0.4) arc[start angle=190,end angle=260,radius=0.4];
  \draw[scale=0.5,line width=0.025cm,blue] (-0.28,0.2) -- (0.5,1);
}

\newcommand{\tkzIronNail}[1][60]{%
  \pgfgettransformentries{\clab@aa}{\clab@bb}{\clab@cc}{\clab@dd}{\clab@xx}{\clab@yy}%
  \pgfmathsetmacro{\clab@scale}{(\clab@aa + \clab@dd) / 2}%
  \foreach \iNail in {1,...,#1}{%
    % pst-labo: \setrandim{ClouX}{-1cm}{1cm} → 物理坐标, 需除以 scale 转换为局部坐标
    \pgfmathsetmacro{\NailX}{(-1 + 2*rnd) / \clab@scale}%
    \pgfmathsetmacro{\NailY}{(-0.2 + 0.6*rnd) / \clab@scale}%
    \pgfmathsetmacro{\NailAngle}{-90 + 180*rnd}%
    \begin{scope}[rotate=\tkzTiltAngle]
      \begin{scope}[shift={(\NailX,\NailY)},rotate=\NailAngle]
        \tkzNails
      \end{scope}
    \end{scope}
  }%
}

% ============================================================
% 15. \tkzCopperTurnings  铜屑
% ============================================================
% 单个铜屑形状 (单位 0.5)
\newcommand{\tkzCopper}{%
  % pst-labo 实际 SLW=0.28458pt=0.01cm (从 xdv 提取)
  \draw[scale=0.5,line width=0.01cm,Copper]
    (0.25,0.25) .. controls (0,0) and (-0.25,0.25) .. (0.25,0.75)
    .. controls (0,0.5) and (-0.25,0.75) .. (0.25,1.25)
    .. controls (0,1) and (-0.25,1.25) .. (0.25,1.75);
}

\newcommand{\tkzCopperTurnings}[1][30]{%
  \pgfgettransformentries{\clab@aa}{\clab@bb}{\clab@cc}{\clab@dd}{\clab@xx}{\clab@yy}%
  \pgfmathsetmacro{\clab@scale}{(\clab@aa + \clab@dd) / 2}%
  \foreach \iT in {1,...,#1}{%
    \pgfmathsetmacro{\TX}{(-1 + 2*rnd) / \clab@scale}%
    \pgfmathsetmacro{\TY}{(-0.2 + 0.5*rnd) / \clab@scale}%
    \pgfmathsetmacro{\TAngle}{-180 + 360*rnd}%
    \begin{scope}[rotate=\tkzTiltAngle]
      \begin{scope}[shift={(\TX,\TY)},rotate=\TAngle]
        \tkzCopper
      \end{scope}
    \end{scope}
  }%
}

% ============================================================
% 16. \tkzZincGranules  锌粒
% ============================================================
% 单个锌粒形状 (单位 0.5)
\newcommand{\tkzZinc}{%
  \filldraw[scale=0.5,fill=LightGray,draw=black]
    plot[smooth,tension=0.5] coordinates {(0,0)(-0.25,0.25)(0,0.5)
      (0.25,0.75)(0.5,0.5)(0.15,0.4)(0,0)} -- cycle;
}

\newcommand{\tkzZincGranules}[1][25]{%
  \pgfgettransformentries{\clab@aa}{\clab@bb}{\clab@cc}{\clab@dd}{\clab@xx}{\clab@yy}%
  \pgfmathsetmacro{\clab@scale}{(\clab@aa + \clab@dd) / 2}%
  \foreach \iG in {1,...,#1}{%
    \pgfmathsetmacro{\GX}{(-1 + 2*rnd) / \clab@scale}%
    \pgfmathsetmacro{\GY}{(-0.2 + 0.5*rnd) / \clab@scale}%
    \pgfmathsetmacro{\GAngle}{-180 + 360*rnd}%
    \begin{scope}[rotate=\tkzTiltAngle]
      \begin{scope}[shift={(\GX,\GY)},rotate=\GAngle]
        \tkzZinc
      \end{scope}
    \end{scope}
  }%
}

% ============================================================
% 17. \tkzBalls  3D 球状物
% ============================================================
\newcommand{\tkzBalls}[1][50]{%
  \foreach \iB in {1,...,#1}{%
    \pgfmathsetmacro{\BX}{-1.5 + 3*rnd}%
    \pgfmathsetmacro{\BY}{3*rnd}%
    \pgfmathsetmacro{\BR}{0.01 + 0.1*rnd}%
    % 按 Y 位置调整半径
    \pgfmathparse{\BY>2?3:(\BY>1?2:1)}%
    \let\Bscale\pgfmathresult
    \pgfmathsetmacro{\BRscaled}{\BR * \Bscale}%
    \begin{scope}[rotate=\tkzTiltAngle]
      % pst-labo Ball3D: GradientCircle=true
      % gradbegin=CMYK(0,0.1,0.95,0.1) → 实际渲染为 RGB(0,0.1,0.95)=蓝色 (固定颜色)
      % gradend=white
      % 径向渐变: 中心=white, 边缘=蓝色
      \definecolor{billeBlue}{rgb}{0,0.1,0.95}%
      \shade[inner color=white, outer color=billeBlue] (\BX,\BY) circle (\BRscaled);
    \end{scope}
  }%
}

% ============================================================
% 18. \tkzBubblesThreeD  2D 气泡
% ============================================================
\newcommand{\tkzBubblesThreeD}[2][50]{%
  \foreach \iB in {1,...,#1}{%
    \pgfmathsetmacro{\BX}{-1.5 + 3*rnd}%
    \pgfmathsetmacro{\BY}{3*rnd}%
    \pgfmathsetmacro{\BR}{0.01 + 0.1*rnd}%
    \pgfmathparse{\BY>2?3:(\BY>1?2:1)}%
    \let\Bscale\pgfmathresult
    \pgfmathsetmacro{\BRscaled}{\BR * \Bscale}%
    \begin{scope}[rotate=\tkzTiltAngle]
      \fill[#2,line width=0.2pt] (\BX,\BY) circle (\BRscaled);
    \end{scope}
  }%
}

\endinput
       % 本文件
%
%  使用示例 :
%    \begin{tikzpicture}
%       \tkzTestTube[glassType=beaker,liquidLevel1=50]
%    \end{tikzpicture}
%
%  参数键 (与原 pst-labo 兼容, 通过 pgfkeys 传入) :
%    burette, pHMeter, buretteColor, buretteLevel,
%    buretteReagent, beakerReagent, reagentColor, reagentLevel,
%    DistillateColor, MixtureStyle,
%    stopper, straightTube, elbowTube, uTube, uTubeLong,
%    curvedTube, curvedTubeShort, DoubleTube, condenserBulbs,
%    clamp, bubbling, gasCollection, tubeOnly, gasTubeStraight,
%    bunsenBurner, Number, label, magneticStirrer,
%    tiltAngle, liquidLevel1, liquidLevel2, liquidLevel3,
%    liquidStyle1, liquidStyle2, liquidStyle3,
%    substance, solid, glassType, TapOpen
%
%  glassType 取值 : tube / roundFlask / beaker / erlenmeyer / reagentBottle / volumetricFlask / watchGlass
% ============================================================

\NeedsTeXFormat{LaTeX2e}
\ProvidesFile{tikz-chemlab.tex}[2026/08/03 v1.0.0]

% --- 文档类检测 ----------------------------------------------
% standalone 会裁剪到 \useasboundingbox, article/book 不会
% 用于在 \tkzHeatFlask 中选择 bbox 策略
\@ifclassloaded{standalone}{\def\chemlab@isstandalone{T}}{\def\chemlab@isstandalone{F}}

% --- 全局 xshift 累加器 (用于多个 \tkzHeatFlask/\tkzRoundFlask 自动推进原点) ---
\newdimen\chemlab@xshift
\chemlab@xshift=0pt
% --- 首次调用标志 (用于在首个 \tkzHeatFlask 设置全局 clip) ---
\newif\ifchemlab@firstcall
\chemlab@firstcalltrue

% --- 依赖加载 ------------------------------------------------
\RequirePackage{tikz}
\usetikzlibrary{arrows.meta,calc,patterns,patterns.meta,shadings,decorations.pathmorphing,decorations.pathreplacing}
\@input{tikz-chemlabObj.tex}

% --- 每个 tikzpicture 开始时自动重置 xshift 和 firstcall ---
% 这样用户无需手动调用 \chemlabresetxshift
\tikzset{every picture/.append style={execute at begin picture={%
  \chemlab@xshift=0pt\chemlab@firstcalltrue%
}}}

% --- 颜色 (原 pst-labo.tex 中定义, 与 Obj 文件中的部分重复此处保留以覆盖) ---
\definecolor{Beige}{cmyk}{0,0,0.1,0.04}
\definecolor{LightGray}{cmyk}{0,0,0,0.2}
\definecolor{VeryLightGray}{cmyk}{0,0,0,0.1}
\definecolor{VeryPaleOrange}{cmyk}{0,0.1,0.3,0}
\definecolor{PaleOrange}{cmyk}{0,0.2,0.4,0}
\definecolor{LightBlue1}{cmyk}{0.2,0,0,0}
\definecolor{LightBlue}{cmyk}{0.22,0.05,0,0.1}
\definecolor{Copper}{cmyk}{0,0.9,0.9,0.2}
\definecolor{Brown}{cmyk}{0,0.3,0.5,.3}

% --- 样式 (对应原 \newpsstyle) --------------------------------
% 使用 \tikzset 定义命名样式; 调用 : \fill[liquidStyle1] ...
\tikzset{
  liquidStyle1/.style={fill=cyan,draw=none},
  liquidStyle2/.style={fill=yellow,draw=none},
  liquidStyle3/.style={fill=magenta,draw=none},
  Ball3D/.style={shading=ball,ball color=cyan,draw=none},
  Champagne/.style={fill=Beige,draw=none},
  Blood/.style={fill=red,draw=none},
  Cobalt/.style={fill=blue,line width=0.2cm},
  Oil/.style={fill=yellow,draw=none},
  Vinegar/.style={fill=magenta,draw=none},
  % pst-labo Diffusion: gradbegin=rgb(0,0.1,0.95), gradend=cyan, gradmidpoint=0
  % gradmidpoint=0: 渐变中点在底部, 大部分区域为 cyan, 底部为 blue
  Diffusion/.style={bottom color={rgb,255:red,0;green,25;blue,242},top color=cyan,draw=none},
  MixtureDiffusion2/.style={bottom color=white,top color=red,draw=red},
  flameAndGrill/.style={bottom color=PaleOrange,top color=yellow,draw=none},
  YellowStripes/.style={pattern=north east lines,pattern color=yellow},
  BlueDiffusion/.style={bottom color=green,top color=cyan,draw=none},
  % 分液漏斗两相液体 (对应 pst-labo \newpsstyle{AqueousPhase}/{OrganicPhase})
  % 颜色定义在 tikz-chemlabObj.tex: AqueousPhaseColor / OrganicPhaseColor
  AqueousPhase/.style={fill=AqueousPhaseColor,draw=none},
  OrganicPhase/.style={fill=OrganicPhaseColor,draw=none},
  % 额外 : 液体通用填充
  tkz liquid/.style={draw=none},
}

% --- 布尔标志 (\newif) ----------------------------------------
\newif\ifTkzBurette
\newif\ifTkzTitrationPHMeter
\newif\ifTkzTestTubeStopper
\newif\ifTkzTestTubeStraightTube
\newif\ifTkzTestTubeElbowTube
\newif\ifTkzTestTubeUTube
\newif\ifTkzTestTubeUTubeLong
\newif\ifTkzTestTubeCurvedTube
\newif\ifTkzTestTubeCurvedTubeShort
\newif\ifTkzTestTubeDoubleTube
\newif\ifTkzTestTubeCondenserBulbs
\newif\ifTkzTestTubeClamp
\newif\ifTkzBubbling
\newif\ifTkzGasCollection
\newif\ifTkzTubeOnly
\newif\ifTkzGasTubeStraight
\newif\ifTkzBunsenBurner
\newif\ifTkzLabel
\newif\ifTkzMagneticStirrer
\newif\ifTkzTapOpen            % (用于 \tkzSeparatoryFunnel)

% --- 参数存储宏 ----------------------------------------------
\def\tkzTiltAngle{0}
\def\tkzLiquidLevelI{50}
\def\tkzLiquidLevelII{0}
\def\tkzLiquidLevelIII{0}
\def\tkzLiquidStyleI{liquidStyle1}
\def\tkzLiquidStyleII{liquidStyle2}
\def\tkzLiquidStyleIII{liquidStyle3}
\def\tkzSubstance{}
\def\tkzSolid{}
\def\tkzNumber{}
\def\tkzGlassType{0}          % 0..6
\def\tkzGlassTypeName{tube}   % 字符串名称 (供嵌套调用传递)
\def\tkzBuretteColor{VeryPaleOrange}
\def\tkzBuretteLevel{20}
\def\tkzBuretteReagent{}
\def\tkzBeakerReagent{}
\def\tkzReagentColor{PaleOrange}
\def\tkzReagentLevel{2}
\def\tkzDistillateColor{yellow}
\def\tkzMixtureStyle{BlueDiffusion}

% --- glassType 字符串 -> 数字 --------------------------------
\newcommand{\tkzGlassTypeLookup}[1]{%
  \ifnum\pdfstrcmp{#1}{tube}=0 0\else
  \ifnum\pdfstrcmp{#1}{roundFlask}=0 1\else
  \ifnum\pdfstrcmp{#1}{beaker}=0 2\else
  \ifnum\pdfstrcmp{#1}{erlenmeyer}=0 3\else
  \ifnum\pdfstrcmp{#1}{reagentBottle}=0 4\else
  \ifnum\pdfstrcmp{#1}{volumetricFlask}=0 5\else
  \ifnum\pdfstrcmp{#1}{watchGlass}=0 6\else
  0\fi\fi\fi\fi\fi\fi\fi
}

% --- 线条粗细变量 (全局, 可由用户通过 \chemlabset 修改) ---
% 四类线宽可独立控制:
%   \chemlabLineWidth            - 主轮廓线 (玻璃器皿外壁)
%   \chemlabEllipseLineWidth     - 瓶口椭圆 (erlenmeyer/reagentBottle/倒置试管, 略粗)
%   \chemlabTubeEllipseLineWidth - 试管口椭圆 (更粗, 原 pst-labo 用 0.08cm)
%   \chemlabGraduationLineWidth  - 刻度线 (较细)
% 默认值与原 pst-labo 完全一致 (在 unit=1cm 时)
% 注: 与 pst-labo 一致, 线宽会自动按当前 tikz scale 缩放
%     (例如 scale=0.5 时, 实际线宽 = 基础线宽 × 0.5)
\def\chemlabBaseLineWidth{0.053cm}
\def\chemlabBaseEllipseLineWidth{0.07cm}
\def\chemlabBaseTubeEllipseLineWidth{0.08cm}
\def\chemlabBaseGraduationLineWidth{0.02cm}

% \chemlab@scaleLength{<基础长度>} - 返回按当前 tikz 缩放因子缩放后的长度
% 用于线宽跟随 scale=N 自动缩放, 与 pst-labo 的 unit=Ncm 行为一致
% 实现: 在每条 \draw 之前由 \chemlabUpdateLineWidths 把缩放后的长度
% 存入 dimen 寄存器, 绘制时直接使用寄存器
\newdimen\chemlabLineWidth
\newdimen\chemlabEllipseLineWidth
\newdimen\chemlabTubeEllipseLineWidth
\newdimen\chemlabGraduationLineWidth
\edef\chemlab@aa{1}\edef\chemlab@bb{0}\edef\chemlab@cc{0}\edef\chemlab@dd{1}\edef\chemlab@xx{0}\edef\chemlab@yy{0}
\newcommand{\chemlabUpdateLineWidths}{%
  \pgfgettransformentries{\chemlab@aa}{\chemlab@bb}{\chemlab@cc}{\chemlab@dd}{\chemlab@xx}{\chemlab@yy}%
  \pgfmathsetlength{\chemlabLineWidth}{\chemlabBaseLineWidth * (\chemlab@aa + \chemlab@dd) / 2}%
  \pgfmathsetlength{\chemlabEllipseLineWidth}{\chemlabBaseEllipseLineWidth * (\chemlab@aa + \chemlab@dd) / 2}%
  \pgfmathsetlength{\chemlabTubeEllipseLineWidth}{\chemlabBaseTubeEllipseLineWidth * (\chemlab@aa + \chemlab@dd) / 2}%
  \pgfmathsetlength{\chemlabGraduationLineWidth}{\chemlabBaseGraduationLineWidth * (\chemlab@aa + \chemlab@dd) / 2}%
}
% 默认值 (未调用 UpdateLineWidths 时的回退)
\chemlabLineWidth=0.053cm
\chemlabEllipseLineWidth=0.07cm
\chemlabTubeEllipseLineWidth=0.08cm
\chemlabGraduationLineWidth=0.02cm

% --- 火焰参数 (全局, 可由用户通过 \chemlabset 修改) ---
% 颜色 (默认与原 pst-labo 一致)
\def\chemlabFlameOuterColor{PaleOrange}   % 外焰起始色 (inner color)
\def\chemlabFlameOuterEndColor{yellow}    % 外焰结束色 (outer color)
\def\chemlabFlameInnerColor{cyan}         % 内焰起始色 (inner color)
\def\chemlabFlameInnerEndColor{white}     % 内焰结束色 (outer color)
% 大小 (默认与原 pst-labo 一致)
% 火焰缩放因子: 1.0=原版大小, 0.5=半高, 2.0=双高
\def\chemlabFlameScale{1.0}

% 线条粗细关键字 -> 数值转换 (对应 TikZ 标准)
% thin=0.4pt, semithick=0.6pt, thick=0.8pt, very thick=1.2pt, ultra thick=1.6pt
% 用法: \chemlab@resolveLineWidth{<关键字或数值>}{<目标宏>}
% 结果: 将解析后的数值存入 <目标宏>
\def\chemlab@thin{thin}
\def\chemlab@semithick{semithick}
\def\chemlab@thick{thick}
\def\chemlab@verythick{very thick}
\def\chemlab@ultrathick{ultra thick}
\newcommand{\chemlab@resolveLineWidth}[2]{%
  \def\chemlab@tmp{#1}%
  \ifx\chemlab@tmp\chemlab@thin          \def#2{0.4pt}%
  \else\ifx\chemlab@tmp\chemlab@semithick \def#2{0.6pt}%
  \else\ifx\chemlab@tmp\chemlab@thick     \def#2{0.8pt}%
  \else\ifx\chemlab@tmp\chemlab@verythick \def#2{1.2pt}%
  \else\ifx\chemlab@tmp\chemlab@ultrathick\def#2{1.6pt}%
  \else  % 直接当作数值 (如 0.04cm, 1pt)
    \def#2{#1}%
  \fi\fi\fi\fi\fi
}

% 三个独立设置器 (供 pgfkeys 调用)
\newcommand{\chemlabSetOutlineWidth}[1]{\chemlab@resolveLineWidth{#1}{\chemlabBaseLineWidth}}
\newcommand{\chemlabSetEllipseWidth}[1]{\chemlab@resolveLineWidth{#1}{\chemlabBaseEllipseLineWidth}}
\newcommand{\chemlabSetGraduationWidth}[1]{\chemlab@resolveLineWidth{#1}{\chemlabBaseGraduationLineWidth}}
\newcommand{\chemlabSetTubeEllipseWidth}[1]{\chemlab@resolveLineWidth{#1}{\chemlabBaseTubeEllipseLineWidth}}

% 同时设置四类线宽的快捷方式
\newcommand{\chemlabSetAllLineWidth}[1]{%
  \chemlab@resolveLineWidth{#1}{\chemlabBaseLineWidth}%
  \chemlab@resolveLineWidth{#1}{\chemlabBaseEllipseLineWidth}%
  \chemlab@resolveLineWidth{#1}{\chemlabBaseTubeEllipseLineWidth}%
  \chemlab@resolveLineWidth{#1}{\chemlabBaseGraduationLineWidth}%
}

% ===== 预设样式 (一键设置四类线宽) =====
% 通过 \chemlabset{style=publication} 等使用
% 预设说明 (数值单位 pt, 1pt ≈ 0.0353cm) :
%   default      - 原 pst-labo 默认 (轮廓 1.508pt, 椭圆 1.984pt, 试管椭圆 2.268pt, 刻度 0.567pt)
%   publication  - 出版级 (细线, 适合印刷) : 轮廓 thin, 椭圆 semithick, 试管椭圆 thick, 刻度 thin
%   presentation- 演示级 (粗线, 适合投影) : 轮廓 thick, 椭圆 very thick, 试管椭圆 ultra thick, 刻度 semithick
%   schematic    - 示意图 (中等, 平衡)   : 轮廓 semithick, 椭圆 thick, 试管椭圆 very thick, 刻度 thin
%   draft        - 草稿 (极细, 快速预览) : 全部 thin
%   bold         - 加粗 (高对比)         : 轮廓 very thick, 椭圆 ultra thick, 试管椭圆 ultra thick, 刻度 thick
\def\chemlab@style@default{default}
\def\chemlab@style@publication{publication}
\def\chemlab@style@presentation{presentation}
\def\chemlab@style@schematic{schematic}
\def\chemlab@style@draft{draft}
\def\chemlab@style@bold{bold}
\newcommand{\chemlabSetStyle}[1]{%
  \def\chemlab@styletmp{#1}%
  \ifx\chemlab@styletmp\chemlab@style@default
    \def\chemlabBaseLineWidth{0.053cm}\def\chemlabBaseEllipseLineWidth{0.07cm}\def\chemlabBaseTubeEllipseLineWidth{0.08cm}\def\chemlabBaseGraduationLineWidth{0.02cm}%
  \else\ifx\chemlab@styletmp\chemlab@style@publication
    \def\chemlabBaseLineWidth{0.4pt}\def\chemlabBaseEllipseLineWidth{0.6pt}\def\chemlabBaseTubeEllipseLineWidth{0.8pt}\def\chemlabBaseGraduationLineWidth{0.4pt}%
  \else\ifx\chemlab@styletmp\chemlab@style@presentation
    \def\chemlabBaseLineWidth{0.8pt}\def\chemlabBaseEllipseLineWidth{1.2pt}\def\chemlabBaseTubeEllipseLineWidth{1.6pt}\def\chemlabBaseGraduationLineWidth{0.6pt}%
  \else\ifx\chemlab@styletmp\chemlab@style@schematic
    \def\chemlabBaseLineWidth{0.6pt}\def\chemlabBaseEllipseLineWidth{0.8pt}\def\chemlabBaseTubeEllipseLineWidth{1.2pt}\def\chemlabBaseGraduationLineWidth{0.4pt}%
  \else\ifx\chemlab@styletmp\chemlab@style@draft
    \def\chemlabBaseLineWidth{0.4pt}\def\chemlabBaseEllipseLineWidth{0.4pt}\def\chemlabBaseTubeEllipseLineWidth{0.4pt}\def\chemlabBaseGraduationLineWidth{0.4pt}%
  \else\ifx\chemlab@styletmp\chemlab@style@bold
    \def\chemlabBaseLineWidth{1.2pt}\def\chemlabBaseEllipseLineWidth{1.6pt}\def\chemlabBaseTubeEllipseLineWidth{1.6pt}\def\chemlabBaseGraduationLineWidth{0.8pt}%
  \else
    \typeout{tikz-chemlab: unknown style `#1'. Available: default, publication, presentation, schematic, draft, bold.}%
  \fi\fi\fi\fi\fi\fi
}

% 用户接口: \chemlabset{line width=thick} 或 \chemlabset{outline line width=thick, ellipse line width=0.1cm}
% 或 \chemlabset{style=publication}
\newcommand{\chemlabset}[1]{\pgfkeys{/tikz-chemlab/.cd, #1}}

% --- pgfkeys 参数族 ------------------------------------------
\pgfkeys{
  /tikz-chemlab/.cd,
  % ===== 线条粗细控制 (全局, 三类可独立设置) =====
  % line width / linewidth : 同时设置三类线宽 (轮廓+椭圆+刻度)
  % 可设为: thin / semithick / thick / very thick / ultra thick
  % 或直接给数值 (如 0.04cm, 1pt 等)
  line width/.code={\chemlabSetAllLineWidth{#1}},
  linewidth/.code={\chemlabSetAllLineWidth{#1}},
  % outline line width : 仅设置主轮廓线
  outline line width/.code={\chemlabSetOutlineWidth{#1}},
  outline width/.code={\chemlabSetOutlineWidth{#1}},
  % ellipse line width : 仅设置瓶口椭圆 (erlenmeyer/reagentBottle/倒置试管)
  ellipse line width/.code={\chemlabSetEllipseWidth{#1}},
  ellipse width/.code={\chemlabSetEllipseWidth{#1}},
  % tube ellipse line width : 仅设置试管口椭圆 (原 pst-labo 用 0.08cm, 比其他椭圆粗)
  tube ellipse line width/.code={\chemlabSetTubeEllipseWidth{#1}},
  tube ellipse width/.code={\chemlabSetTubeEllipseWidth{#1}},
  % graduation line width : 仅设置刻度线
  graduation line width/.code={\chemlabSetGraduationWidth{#1}},
  graduation width/.code={\chemlabSetGraduationWidth{#1}},
  % style : 预设样式 (一键设置三类线宽)
  % 可选: default / publication / presentation / schematic / draft / bold
  style/.code={\chemlabSetStyle{#1}},
  % ===== 火焰参数 (本生灯) =====
  % 颜色: 可用任意 TikZ 颜色名 (如 cyan, yellow, red, PaleOrange 等)
  flame outer color/.store in=\chemlabFlameOuterColor,
  flame outer end color/.store in=\chemlabFlameOuterEndColor,
  flame inner color/.store in=\chemlabFlameInnerColor,
  flame inner end color/.store in=\chemlabFlameInnerEndColor,
  % 大小: 缩放因子 (1.0=原版, 0.5=半高, 2.0=双高)
  flame scale/.store in=\chemlabFlameScale,
  burette/.is if=TkzBurette,
  pHMeter/.is if=TkzTitrationPHMeter,
  stopper/.is if=TkzTestTubeStopper,
  straightTube/.is if=TkzTestTubeStraightTube,
  elbowTube/.is if=TkzTestTubeElbowTube,
  uTube/.is if=TkzTestTubeUTube,
  uTubeLong/.is if=TkzTestTubeUTubeLong,
  curvedTube/.is if=TkzTestTubeCurvedTube,
  curvedTubeShort/.is if=TkzTestTubeCurvedTubeShort,
  doubleTube/.is if=TkzTestTubeDoubleTube,
  condenserBulbs/.is if=TkzTestTubeCondenserBulbs,
  clamp/.is if=TkzTestTubeClamp,
  noDefaults/.is if=chemlab@noDefaults,
  bubbling/.is if=TkzBubbling,
  gasCollection/.is if=TkzGasCollection,
  tubeOnly/.is if=TkzTubeOnly,
  gasTubeStraight/.is if=TkzGasTubeStraight,
  bunsenBurner/.is if=TkzBunsenBurner,
  label/.is if=TkzLabel,
  magneticStirrer/.is if=TkzMagneticStirrer,
  TapOpen/.is if=TkzTapOpen,
  tiltAngle/.store in=\tkzTiltAngle,
  liquidLevel1/.store in=\tkzLiquidLevelI,
  liquidLevel2/.store in=\tkzLiquidLevelII,
  liquidLevel3/.store in=\tkzLiquidLevelIII,
  liquidStyle1/.store in=\tkzLiquidStyleI,
  liquidStyle2/.store in=\tkzLiquidStyleII,
  liquidStyle3/.store in=\tkzLiquidStyleIII,
  substance/.store in=\tkzSubstance,
  solid/.store in=\tkzSolid,
  Number/.store in=\tkzNumber,
  glassType/.code={\edef\tkzGlassType{\tkzGlassTypeLookup{#1}}\def\tkzGlassTypeName{#1}},
  buretteColor/.store in=\tkzBuretteColor,
  buretteLevel/.store in=\tkzBuretteLevel,
  buretteReagent/.store in=\tkzBuretteReagent,
  beakerReagent/.store in=\tkzBeakerReagent,
  reagentColor/.store in=\tkzReagentColor,
  reagentLevel/.store in=\tkzReagentLevel,
  DistillateColor/.store in=\tkzDistillateColor,
  MixtureStyle/.store in=\tkzMixtureStyle,
}

% --- 默认值 (对应原 \psset[pst-labo]{...}) -------------------
\newcommand{\tikzchemlabdefaults}{%
  \TkzBurettetrue
  \TkzTitrationPHMeterfalse
  \TkzTestTubeStopperfalse
  \TkzTestTubeStraightTubefalse
  \TkzTestTubeElbowTubefalse
  \TkzTestTubeUTubefalse
  \TkzTestTubeUTubeLongfalse
  \TkzTestTubeCurvedTubefalse
  \TkzTestTubeCurvedTubeShortfalse
  \TkzTestTubeDoubleTubefalse
  \TkzTestTubeCondenserBulbsfalse
  \TkzTestTubeClampfalse
  \TkzBubblingfalse
  \TkzGasCollectionfalse
  \TkzTubeOnlyfalse
  \TkzGasTubeStraightfalse
  \TkzBunsenBurnertrue
  \TkzLabelfalse
  \TkzMagneticStirrertrue
  \TkzTapOpenfalse
  \def\tkzTiltAngle{0}%
  \def\tkzLiquidLevelI{50}%
  \def\tkzLiquidLevelII{0}%
  \def\tkzLiquidLevelIII{0}%
  \def\tkzLiquidStyleI{liquidStyle1}%
  \def\tkzLiquidStyleII{liquidStyle2}%
  \def\tkzLiquidStyleIII{liquidStyle3}%
  \def\tkzSubstance{}%
  \def\tkzSolid{}%
  \def\tkzNumber{}%
  \def\tkzGlassType{0}%
  \def\tkzGlassTypeName{tube}%
  \def\tkzBuretteColor{VeryPaleOrange}%
  \def\tkzBuretteLevel{20}%
  \def\tkzBuretteReagent{}%
  \def\tkzBeakerReagent{}%
  \def\tkzReagentColor{PaleOrange}%
  \def\tkzReagentLevel{2}%
  \def\tkzDistillateColor{yellow}%
  \def\tkzMixtureStyle{BlueDiffusion}%
}

% ============================================================
% 玻璃器皿路径版本 (供 \clip 使用, 不带 \draw/\fill)
% ============================================================

% 滴定管路径 (burette corps)
\newcommand{\tkzBuretteBodyPath}{%
  (-0.5,11.5) -- (-0.25,11) -- (-0.25,3) -- (0,2.75) -- (0,0.5)
  -- (0.2,0.7) -- (0.2,2.75) -- (0.5,3) -- (0.5,11) -- (0.75,11.5) -- cycle
}

% 移液管路径 (pipette corps + 液滴)
\newcommand{\tkzPipetteBodyPath}{%
  (-0.05,0) -- (-0.15,0.2) -- (-0.15,6)
  -- (0.15,6) -- (0.15,0.2) -- (0.05,0) -- cycle
  (0,0) .. controls (0.25,-.4) and (-.25,-.4) .. (0,0)
}

% 试管路径 (TubeEssais corps)
\newcommand{\tkzTestTubeBodyPath}{%
  (-0.5,3.5) -- (-0.5,0.5)
  arc[start angle=180,end angle=360,radius=0.5]
  -- (0.5,3.5) -- cycle
}

% 烧杯路径 (Becher corps) - 显式圆弧, 与 \tkzBeakerBody 一致
\newcommand{\tkzBeakerBodyPath}{%
  (-1,2.5) -- (-1,0.5)
  -- ([shift={(-0.5,0.5)}]180:0.5) arc[start angle=180,end angle=270,radius=0.5]
  -- (0.5,0)
  -- ([shift={(0.5,0.5)}]270:0.5) arc[start angle=270,end angle=360,radius=0.5]
  -- (1,2.5) -- cycle
}

% 锥形瓶路径 (Erlenmeyer corps) - 显式圆弧, 与 \tkzErlenmeyerBody 一致
\newcommand{\tkzErlenmeyerBodyPath}{%
  (-0.5,3.5) -- (-0.5,2.5)
  -- (-1.3354,0.4114)
  -- ([shift={(-1.0569,0.3)}]158.20:0.3) arc[start angle=158.20,end angle=270,radius=0.3]
  -- (1.0569,0)
  -- ([shift={(1.0569,0.3)}]-90:0.3) arc[start angle=-90,end angle=21.80,radius=0.3]
  -- (0.5,2.5) -- (0.5,3.5) -- cycle
}

% 试剂瓶路径 (Flacon corps)
\newcommand{\tkzReagentBottleBodyPath}{%
  ([shift={(-0.6,3.4)}]90:0.1) arc[start angle=90,end angle=0,radius=0.1]
  -- (-0.5,3.4) -- (-0.5,3.1)
  -- ([shift={(-0.6,3.1)}]0:0.1) arc[start angle=0,end angle=-90,radius=0.1]
  -- ([shift={(-0.6,2.6)}]90:0.4) arc[start angle=90,end angle=180,radius=0.4]
  -- (-1,2.6) -- (-1,0.2)
  -- ([shift={(-.8,0.2)}]180:0.2) arc[start angle=180,end angle=270,radius=0.2]
  -- (-.8,0) -- (.8,0)
  -- ([shift={(.8,0.2)}]-90:0.2) arc[start angle=-90,end angle=0,radius=0.2]
  -- (1,0.2) -- (1,2.6)
  -- ([shift={(0.6,2.6)}]0:0.4) arc[start angle=0,end angle=90,radius=0.4]
  -- ([shift={(0.6,3.1)}]-90:0.1) arc[start angle=-90,end angle=-180,radius=0.1]
  -- (0.5,3.1) -- (0.5,3.4)
  -- ([shift={(0.6,3.4)}]180:0.1) arc[start angle=180,end angle=90,radius=0.1]
  -- cycle
}

% 容量瓶路径 (FioleJauge)
\newcommand{\tkzVolumetricFlaskPath}{%
  (-0.21,3.5) -- (-0.1,3.5) -- (-0.1,1.5)
  -- (-0.75,0) -- (0.75,0) -- (0.1,1.5) -- (0.1,3.5) -- (0.21,3.5) -- cycle
}

% 玻璃杯路径 (Verre corps)
\newcommand{\tkzWatchGlassBodyPath}{%
  (-1,3.5) -- (-0.5,1)
  arc[start angle=180,end angle=360,radius=0.5]
  -- (1,3.5) -- cycle
}

% 量筒路径 (Eprouvette corps)
\newcommand{\tkzGraduatedCylinderBodyPath}{%
  (-0.5,5) -- (-0.5,0) -- (0.5,0) -- (0.5,4.9) -- (0.6,4.93) -- cycle
}

% 统一调度 : 根据 \tkzGlassType 选择 clip 路径
\newcommand{\tkzGlassClipPath}{%
  \ifcase\tkzGlassType
    \tkzTestTubeBodyPath%        0 tube
  \or
    \tkzRoundFlaskSimplePath%          1 roundFlask
  \or
    \tkzBeakerBodyPath%           2 beaker
  \or
    \tkzErlenmeyerBodyPath%            3 erlenmeyer
  \or
    \tkzReagentBottleBodyPath%           4 reagentBottle
  \or
    \tkzVolumetricFlaskPath%            5 volumetricFlask
  \or
    \tkzWatchGlassBodyPath%            6 watchGlass
  \else
    \tkzTestTubeBodyPath%
  \fi
}

% 各 glassType 对应的 clip 圆角值 (与 \tkzGlassDraw 中各器皿的圆角保持一致,
% 否则液体会溢出实际绘制的轮廓外)
% 通过 \edef 预先展开为纯数字, 避免在 \clip[...] 选项中嵌套 \ifcase
% 注 : beaker(2) 和 erlenmeyer(3) 的路径已用显式圆弧, 故此处用 0pt
\newcommand{\tkzGlassClipOptions}{0pt}
\newcommand{\tkzGlassClipOptionsSet}{%
  \ifcase\tkzGlassType
    \edef\tkzGlassClipOptions{0pt}%       0 tube
  \or
    \edef\tkzGlassClipOptions{0pt}%       1 roundFlask
  \or
    \edef\tkzGlassClipOptions{0pt}%       2 beaker (路径自带圆弧)
  \or
    \edef\tkzGlassClipOptions{0pt}%       3 erlenmeyer  (路径自带圆弧)
  \or
    \edef\tkzGlassClipOptions{0pt}%       4 reagentBottle
  \or
    \edef\tkzGlassClipOptions{0.1cm}%     5 volumetricFlask
  \or
    \edef\tkzGlassClipOptions{0pt}%       6 watchGlass
  \else
    \edef\tkzGlassClipOptions{0pt}%
  \fi
}

% 统一调度 : 根据 \tkzGlassType 绘制对应玻璃器皿轮廓
\newcommand{\tkzGlassDraw}{%
  \ifcase\tkzGlassType
    \tkzTestTubeBody%                                0 tube
  \or
    \tkzRoundFlaskSimple%                                   1 roundFlask
  \or
    \tkzBeakerBody\relax\tkzBeakerRim%              2 beaker
  \or
    \tkzErlenmeyerBody\relax\filldraw[fill=white,draw=black,line width=\chemlabEllipseLineWidth,line cap=round] (0,3.5) ellipse (0.6 and 0.1);% 3 erlenmeyer
  \or
    \tkzReagentBottleBody\relax\filldraw[fill=white,draw=black,line width=\chemlabEllipseLineWidth,line cap=round] (0,3.5) ellipse (0.6 and 0.1);% 4 reagentBottle
  \or
    \tkzVolumetricFlask%                                     5 volumetricFlask
  \or
    \tkzWatchGlassBody\relax\tkzWatchGlassBase%                  6 watchGlass
  \fi
}

% ============================================================
% 液体多边形 helper (对应 \pst@Eprouvette@Liquide / \pst@TubeEssais@Liquide)
% ============================================================
% 参数 : #1=niveau(0-100), #2=style, #3=angle, #4=液面比例高度(3 或 4.5)
\newcommand{\tkzLiquidFill}[4]{%
  \pgfmathsetmacro{\liquidY}{#4 * #1 / 100 * cos(#3)}%
  \fill[#2] (-4,\liquidY) -- (-4,-1) -- (5,-1) -- (5,\liquidY) -- cycle;
}

% 试管液体 (高度 3)
\newcommand{\tkzTestTubeLiquid}[3]{\tkzLiquidFill{#1}{#2}{#3}{3}}
% 量筒液体 (高度 4.5)
\newcommand{\tkzGraduatedCylinderLiquid}[3]{\tkzLiquidFill{#1}{#2}{#3}{4.5}}

% ============================================================
% 1. \tkzTitration  滴定装置
% ============================================================
\newcommand{\tkzTitration}[1][]{%
  % 保存全局参数 (匹配 pst-labo \psset 的全局继承行为)
  \edef\chemlab@savedGlassType{\tkzGlassType}%
  \edef\chemlab@savedGlassTypeName{\tkzGlassTypeName}%
  \edef\chemlab@savedLiquidLevelI{\tkzLiquidLevelI}%
  \tikzchemlabdefaults
  \pgfkeys{/tikz-chemlab/.cd, #1}%
  % 恢复全局参数 (如果命令选项未指定)
  \ifnum\tkzGlassType=0 \edef\tkzGlassType{\chemlab@savedGlassType}\edef\tkzGlassTypeName{\chemlab@savedGlassTypeName}\fi
  \ifnum\tkzLiquidLevelI=50\edef\tkzLiquidLevelI{\chemlab@savedLiquidLevelI}\fi
  % --- 方案 B 局部修复: 忠实复刻 pst-labo \tkzRoundFlask 的 fallback 行为 ---
  % pst-labo.tex L478-522: \tkzTitration 内部调用 \tkzRoundFlask, 其 \ifcase\psk@glassType
  %   0(tube)/1(roundFlask) -> \pst@roundFlask
  %   2(beaker)         -> \pst@Becher@Corps
  %   3(erlenmeyer)          -> \pst@Erlen@Corps
  %   4(reagentBottle)         -> \pst@Flacon@Corps
  %   5(volumetricFlask)/6(Verre)/未知 -> \pst@roundFlask (fallback, L485 注释 "roundFlask is default")
  % tikz-chemlab 的 \tkzGlassDraw 对 glassType=0 画的是 \tkzTestTubeBody (试管),
  % 与原版不一致. 此处仅在 \tkzTitration 内部把 tube/volumetricFlask/Verre/未知 重映射为 roundFlask,
  % 使 \tkzGlassDraw/\tkzGlassClipPath 走 roundFlask 分支. 不影响其他命令.
  \ifnum\tkzGlassType=0 \def\tkzGlassType{1}\def\tkzGlassTypeName{roundFlask}\fi
  \ifnum\tkzGlassType=5 \def\tkzGlassType{1}\def\tkzGlassTypeName{roundFlask}\fi
  \ifnum\tkzGlassType=6 \def\tkzGlassType{1}\def\tkzGlassTypeName{roundFlask}\fi
  \ifnum\tkzGlassType>6 \def\tkzGlassType{1}\def\tkzGlassTypeName{roundFlask}\fi
  \begin{scope}[xshift=\chemlab@xshift]
    % 边界框
    \ifTkzBurette
      \ifTkzTitrationPHMeter
        \useasboundingbox (-5,0) rectangle (3,16);
      \else
        \useasboundingbox (-2,0) rectangle (2,16);
      \fi
    \else
      \ifTkzTitrationPHMeter
        \useasboundingbox (-5,0) rectangle (3,5);
      \else
        \useasboundingbox (-2,0) rectangle (2,5);
      \fi
    \fi
    % 烧瓶 + 液体 (匹配 pst-labo \rput(0,3.5){\tkzRoundFlask})
    % pst-labo: \tkzRoundFlask 的 pspicture bbox = (-2,0)(2,4), 中心 (0,2)
    % \rput(0,3.5) 使用 ref=c (中心对齐), 所以局部原点 (瓶底) 位于世界 (0, 3.5-2) = (0, 1.5)
    % 锥形瓶底部 (y=0 local) 正好放在搅拌器台面上 (y=1.5 world)
    \begin{scope}[shift={(0,1.5)}]
      % 液体 (clip 到玻璃内部 ∩ 液体矩形)
      \begin{scope}
        \tkzGlassClipOptionsSet
        \clip[rounded corners=\tkzGlassClipOptions] \tkzGlassClipPath;
        \ifnum\tkzLiquidLevelI>0
          \pgfmathsetmacro{\chemlab@liquidClipY}{3 * \tkzLiquidLevelI / 100 * cos(\tkzTiltAngle)}%
          \clip (-4,-1.5) rectangle (5,\chemlab@liquidClipY);
        \fi
        \ifnum\tkzLiquidLevelI>0
          \tkzTestTubeLiquid{\tkzLiquidLevelI}{\tkzLiquidStyleI}{\tkzTiltAngle}
        \fi
        \tkzSubstance
        \tkzSolid
      \end{scope}
      % 玻璃器皿轮廓 (根据 glassType 选择: erlenmeyer/beaker/roundFlask 等)
      \tkzGlassDraw
    \end{scope}
    % 滴定管 + 液体
    \ifTkzBurette
      \begin{scope}[shift={(0,4.5)}]
        % 液体填充 (clip 到 burette 内部)
        \begin{scope}
          \clip \tkzBuretteBodyPath;
          \pgfmathsetmacro{\buretteY}{\tkzBuretteLevel * 0.32 + 2.68}
          \fill[\tkzBuretteColor] (-1,-1) rectangle (1,\buretteY);
        \end{scope}
        \tkzBuretteBody
        \tkzBuretteStopcock
        \tkzBuretteGraduation
      \end{scope}
    \fi
    % pH 计
    \ifTkzTitrationPHMeter
      \begin{scope}[shift={(-3,0)}]\tkzTitrationPHMeter\end{scope}
    \fi
    % 搅拌子 + 支架
    \ifTkzMagneticStirrer
      \begin{scope}[shift={(0,1.5)}]\tkzTitrationStirBar\end{scope}
      \begin{scope}[shift={(0,0)}]\tkzTitrationStand\end{scope}
    \else
      % pst-labo: \psset{linewidth=0.053}\psframe(-1.5,0)(1.5,1.5)
      % linewidth=0.053 是用户单位 (×unit), 随 scale 缩放
      \chemlabUpdateLineWidths
      \draw[line width=\chemlabLineWidth] (-1.5,0) rectangle (1.5,1.5);
    \fi
    % 反应物标签
    \ifTkzBurette
      \node[anchor=west] at (2,11) {\tkzBuretteReagent};
    \fi
    \node[anchor=west] at (2,3.5) {\tkzBeakerReagent};
  \end{scope}
  % 自动推进原点 (匹配 pst-labo \pspicture 的布局宽度)
  \ifTkzTitrationPHMeter
    \advance\chemlab@xshift by 8cm %  width = 8
  \else
    \advance\chemlab@xshift by 4cm %  width = 4
  \fi
}

% ============================================================
% 2. \tkzPipette  移液管
% ============================================================
\newcommand{\tkzPipette}[1][]{%
  \begingroup
  \tikzchemlabdefaults
  \pgfkeys{/tikz-chemlab/.cd, #1}%
  \begin{scope}
    \useasboundingbox (-1,-1) rectangle (1,6);
    % 液体
    \begin{scope}
      \clip \tkzPipetteBodyPath;
      \pgfmathsetmacro{\pipetteY}{\tkzReagentLevel / 2}
      \fill[\tkzReagentColor] (-1,-2) rectangle (1,\pipetteY);
    \end{scope}
    \tkzPipetteBody
    \tkzDrop
  \end{scope}
  \endgroup
}

% ============================================================
% 3. \tkzLargePipette  大移液管
% ============================================================
\newcommand{\tkzLargePipette}[1][]{%
  \begingroup
  \tikzchemlabdefaults
  \pgfkeys{/tikz-chemlab/.cd, #1}%
  \begin{scope}
    \useasboundingbox (-1,-1) rectangle (1,6);
    % 注 : 原文 \pst@pipette 是已注释的代码, 此处使用 bigpipette 形状
    \begin{scope}
      % bigpipette 形状比较简单, 直接 clip 一个矩形区域近似
      \clip (0,0) -- (0,1) -- (-0.23,2.5) -- (-0.23,22.5)
            arc[start angle=-80,end angle=80,radius=0.45] -- cycle;
      \pgfmathsetmacro{\pipetteY}{\tkzReagentLevel / 2}
      \fill[\tkzReagentColor] (-1,-2) rectangle (1,\pipetteY);
    \end{scope}
    \tkzLargePipetteBody
  \end{scope}
  \endgroup
}

% ============================================================
% 4. \tkzGraduatedCylinder  量筒
% ============================================================
\newcommand{\tkzGraduatedCylinder}[1][]{%
  \begingroup
  \tikzchemlabdefaults
  \pgfkeys{/tikz-chemlab/.cd, #1}%
  \begin{scope}
    % 注: pst-labo 用 \pspicture(-2,0)(2,6) 设置 bbox 用于 \rput 放置, 但内容可超出 bbox
    % 在 TikZ 中, useasboundingbox 会覆盖整个 picture 的 bbox, 不适合嵌套使用
    % 因此这里不设 useasboundingbox, 让外层命令 (如 Figure 中的 pspicture) 控制 bbox
    % standalone 编译单独的 \tkzGraduatedCylinder 时, 自动检测内容 bbox
    % 液体 (clip 到量筒内部, 旋转 tiltAngle)
    \begin{scope}[rotate=\tkzTiltAngle]
      \begin{scope}
        \clip \tkzGraduatedCylinderBodyPath;
        \ifnum\tkzLiquidLevelI>0
          \tkzGraduatedCylinderLiquid{\tkzLiquidLevelI}{\tkzLiquidStyleI}{\tkzTiltAngle}
        \fi
        \ifnum\tkzLiquidLevelII>0
          \tkzGraduatedCylinderLiquid{\tkzLiquidLevelII}{\tkzLiquidStyleII}{\tkzTiltAngle}
        \fi
        \ifnum\tkzLiquidLevelIII>0
          \tkzGraduatedCylinderLiquid{\tkzLiquidLevelIII}{\tkzLiquidStyleIII}{\tkzTiltAngle}
        \fi
        \tkzSubstance
      \end{scope}
      \tkzGraduatedCylinderBody
      \tkzGraduatedCylinderBase
      \tkzGraduatedCylinderGraduation
      \begin{scope}[shift={(0,5)}]\tkzGraduatedCylinderRim\end{scope}
    \end{scope}
  \end{scope}
  \endgroup
}

% ============================================================
% 5. \tkzDistillation  蒸馏装置
% ============================================================
\newcommand{\tkzDistillation}[1][]{%
  \begingroup
  \tikzchemlabdefaults
  \pgfkeys{/tikz-chemlab/.cd, #1}%
  \begin{scope}
    \useasboundingbox (-4,-10) rectangle (8,7);
    \begin{scope}[shift={(-3,-7)}]\tkzDistillationFractionated\end{scope}
  \end{scope}
  \endgroup
}

% ============================================================
% 6. \tkzFunnel  漏斗 + 试管
% ============================================================
\newcommand{\tkzFunnel}[1][]{%
  \begingroup
  \tikzchemlabdefaults
  \pgfkeys{/tikz-chemlab/.cd, #1}%
  % 设置 noDefaults=true, 使 \tkzTestTube 继承用户参数
  \chemlab@noDefaultstrue
  \begin{scope}
    % pst-labo: \rput(0,0){\tkzTestTube} 用 pspicture bbox CENTER 对齐
    % tikz: shift = rput_pos - bbox_center (in local coords)
    % pspicture bbox 取决于 glassType 和 tiltAngle:
    %   beaker (glassType=2), tiltAngle=0: (-0.5,0)(1,2.8), center=(0.25,1.4)
    %   tiltAngle<0: (-1,0)(3,4), center=(1,2)
    %   tiltAngle>0: (-3,0)(1,4), center=(-1,2)
    %   其他 (tiltAngle=0, glassType!=2): (-1.5,0)(1.5,4), center=(0,2)
    \ifnum\tkzTiltAngle=0
      \ifnum\tkzGlassType=2
        % beaker, 直立: center=(0.25,1.4)
        \begin{scope}[shift={(-0.25,-1.4)}]\tkzTestTube[noDefaults=true]\end{scope}
      \else
        % roundFlask/reagentBottle/erlenmeyer, 直立: center=(0,2)
        \begin{scope}[shift={(0,-2)}]\tkzTestTube[noDefaults=true]\end{scope}
      \fi
    \else
      \ifnum\tkzTiltAngle>0
        % 倾斜>0: center=(-1,2)
        \begin{scope}[shift={(1,-2)}]\tkzTestTube[noDefaults=true]\end{scope}
      \else
        % 倾斜<0: center=(1,2)
        \begin{scope}[shift={(-1,-2)}]\tkzTestTube[noDefaults=true]\end{scope}
      \fi
    \fi
    % 漏斗: pst-labo \rput(0,1.5){\pst@Entonnoir@Corps}
    % \pst@Entonnoir@Corps 路径 bbox: x=[-1,1], y=[0,2.6] (含椭圆顶 2.5+0.1)
    % bbox center = (0, 1.3)
    % \rput(0,1.5) ref=c: center 置于 (0,1.5)
    % shift = rput_pos - center = (0, 1.5-1.3) = (0, 0.2)
    % 即管尖 (路径原点 0,0) 位于 (0, 0.2)
    \begin{scope}[shift={(0,0.2)}]\tkzFunnelBody\end{scope}
  \end{scope}
  \chemlab@noDefaultsfalse
  \endgroup
}

% ============================================================
% 7. \tkzTestTube  试管 (核心命令)
% ============================================================
\newcommand{\tkzTestTube}[1][]{%
  \begingroup
  \ifchemlab@noDefaults\else\tikzchemlabdefaults\fi
  \pgfkeys{/tikz-chemlab/.cd, #1}%
  \begin{scope}
    % 注: pst-labo 用 \pspicture(bbox) 设置 bbox 用于 \rput 放置, 但内容可超出 bbox
    % 在 TikZ 中, useasboundingbox 会覆盖整个 picture 的 bbox, 不适合嵌套使用
    % 因此这里不设 useasboundingbox, 让外层命令 (如 \tkzHeatTube) 控制 bbox
    % standalone 编译单独的 \tkzTestTube 时, 自动检测内容 bbox
    % 液体 (clip 到旋转的玻璃器皿内部, 但液面保持水平)
    % 原版 pst-labo: \psclip{ \rput{\angle}{<轮廓>} \psframe<液面不旋转> }
    % pst-labo 的 \psclip 产生两次裁剪 (PostScript clip 是累积的=交集):
    %   1. 裁剪到玻璃形状 (来自 \pscustom)
    %   2. 裁剪到液体矩形 (来自 \psframe(-3,-1)(5, liquid_level))
    % → 物质被裁剪到 玻璃形状 ∩ 液体矩形 = 液面以下
    % TikZ: 先旋转 clip 玻璃路径, 再反向旋转 clip 液体矩形, 然后画液体和物质
    \begin{scope}[rotate=\tkzTiltAngle]
      \begin{scope}
        \tkzGlassClipOptionsSet
        \clip[rounded corners=\tkzGlassClipOptions] \tkzGlassClipPath;
        % 液面: 在旋转的 clip 内, 但反向旋转使液面水平
        % 液面高度需乘 cos(tiltAngle) (原版 pst-labo 的行为)
        \begin{scope}[rotate=-\tkzTiltAngle]
          % 第二次裁剪: 液体矩形 (匹配 pst-labo \psframe 裁剪)
          % 只在有液体时才裁剪; 无液体时不裁剪 (物质可见到玻璃顶部)
          \ifnum\tkzLiquidLevelI>0
            \pgfmathsetmacro{\chemlab@liquidClipY}{3 * \tkzLiquidLevelI / 100 * cos(\tkzTiltAngle)}%
            \clip (-4,-1.5) rectangle (5,\chemlab@liquidClipY);
          \fi
          \ifnum\tkzLiquidLevelI>0
            \tkzTestTubeLiquid{\tkzLiquidLevelI}{\tkzLiquidStyleI}{\tkzTiltAngle}
          \fi
          \ifnum\tkzLiquidLevelII>0
            \tkzTestTubeLiquid{\tkzLiquidLevelII}{\tkzLiquidStyleII}{\tkzTiltAngle}
          \fi
          \ifnum\tkzLiquidLevelIII>0
            \tkzTestTubeLiquid{\tkzLiquidLevelIII}{\tkzLiquidStyleIII}{\tkzTiltAngle}
          \fi
          \tkzSubstance
          \tkzSolid
        \end{scope}
      \end{scope}
      % 标签
      \ifTkzLabel
        \node[fill=VeryLightGray,draw=none,inner sep=2pt] at (0,2) {\tkzNumber};
      \fi
      % 玻璃器皿轮廓
      \tkzGlassDraw
      % 顶部附件 (straightTube / elbowTube / uTube / stopper)
      \ifTkzTestTubeStraightTube
        \begin{scope}[shift={(0,3.5)}]\tkzTestTubeStraightTube\end{scope}
      \else
        \ifTkzTestTubeElbowTube
          \begin{scope}[shift={(0,3.5)}]\tkzTestTubeElbowTube\end{scope}
        \else
          \ifTkzTestTubeUTube
            \begin{scope}[shift={(0,3.5)}]\tkzTestTubeUTube\end{scope}
          \else
            \ifTkzTestTubeStopper
              \begin{scope}[shift={(0,3.5)}]\tkzTestTubeStopper\end{scope}
            \fi
          \fi
        \fi
      \fi
      % 试管夹
      \ifTkzTestTubeClamp
        \begin{scope}[shift={(0,3.2)}]\tkzTestTubeClamp\end{scope}
      \fi
    \end{scope}
  \end{scope}
  \endgroup
}

% ============================================================
% 8. \tkzHeatTube  加热试管
% ============================================================
% pstChauffageTube: 保存用户参数, 传递给 \tkzTestTube
% pst-labo 用 \psset 全局继承; tikz-chemlab 用 \tkzTestTube[noDefaults] 跳过重置
\newif\ifchemlab@noDefaults
\newcommand{\tkzHeatTube}[1][]{%
  \begingroup
  \tikzchemlabdefaults
  \pgfkeys{/tikz-chemlab/.cd, #1}%
  \chemlab@noDefaultstrue
  \begin{scope}
    % pst-labo bbox: tubeOnly=(-2,-5)(5,4), 其他=(-3,-5)(10,4)
    % 不使用 \useasboundingbox — pst-labo 的 pspicture 不裁剪内容,
    % 内容可以超出 bbox (如 elbowTube 弯管延伸到 x>5)
    % 让 TikZ 自动检测完整内容 bbox
    \ifTkzBubbling
      % 上方试管/烧瓶: 继承用户设置的 glassType, 倾斜-60°, 有弯管
      \edef\chemlab@tmpGlassType{\tkzGlassTypeName}%
      % pst-labo: \rput(2.5,2.5){\tkzTestTube[elbowTube=true,tiltAngle=-60]}
      % elbowTube=true → bbox (-1,0)(6,4) center (2.5,2)
      % → 局部原点 (瓶底) 世界 (2.5-2.5, 2.5-2) = (0, 0.5)
      \begin{scope}[shift={(0,0.5)}]
        \tkzTestTube[elbowTube=true,tiltAngle=-60,noDefaults=true]
      \end{scope}
      % 下方试管: pst-labo \rput(8.2,-3){\tkzTestTube[tiltAngle=30,...]}
      % tiltAngle>0 → bbox (-3,0)(1,4) center (-1,2)
      % → 局部原点 (试管底) 世界 (8.2-(-1), -3-2) = (9.2, -5)
      \begin{scope}[shift={(9.2,-5)}]
        \tkzTestTube[liquidLevel1=70,glassType=tube,elbowTube=false,
                       tiltAngle=30,
                       liquidStyle1=Champagne,
                       substance={\tkzChampagneBubbles[10]},clamp=false]
      \end{scope}
    \else
      % 非barbotage: 上方器皿继承 glassType
      \edef\chemlab@tmpGlassType{\tkzGlassTypeName}%
      \ifcase\tkzGlassType
        % glassType=tube
        \ifTkzTubeOnly
          % tubeOnly + tube: pst-labo \rput(1.5,2.5)
          % tiltAngle<0 → bbox (-1,0)(3,4) center (1,2) → 世界 (0.5, 0.5)
          \begin{scope}[shift={(0.5,0.5)}]\tkzTestTube[tiltAngle=-60,noDefaults=true]\end{scope}
        \else
          \ifTkzTestTubeElbowTube
            % elbowTube: pst-labo \rput(2,2.5)
            % elbowTube=true → bbox (-1,0)(6,4) center (2.5,2) → 世界 (-0.5, 0.5)
            \begin{scope}[shift={(-0.5,0.5)}]\tkzTestTube[tiltAngle=-60,noDefaults=true]\end{scope}
          \else
            % 无tubeCoude: pst-labo \rput(1,2.5)
            % tiltAngle<0 → bbox (-1,0)(3,4) center (1,2) → 世界 (0, 0.5)
            \begin{scope}[shift={(0,0.5)}]\tkzTestTube[tiltAngle=-60,noDefaults=true]\end{scope}
          \fi
        \fi
      \or
        % glassType=roundFlask (无论 tubeOnly 与否, pst-labo 都是 \rput(2,2.5))
        % tiltAngle<0 → bbox (-1,0)(3,4) center (1,2) → 局部原点世界 (2-1, 2.5-2) = (1, 0.5)
        \begin{scope}[shift={(1,0.5)}]\tkzTestTube[tiltAngle=-60,noDefaults=true]\end{scope}
      \else
        % 其他 glassType: pst-labo \rput(1,2.5) → 世界 (0, 0.5)
        \begin{scope}[shift={(0,0.5)}]\tkzTestTube[tiltAngle=-60,noDefaults=true]\end{scope}
      \fi
    \fi
    \ifTkzBunsenBurner
      % pst-labo: \rput(1,-5){\pst@BecBunsen} (无 pspicture, 实测原点对齐)
      % 灯底位于世界 (1, -5)
      \begin{scope}[shift={(1,-5)}]\tkzBunsenBurner\end{scope}
    \fi
  \end{scope}
  \chemlab@noDefaultsfalse
  \endgroup
}

% ============================================================
% 9. \tkzHeatFlask  加热烧瓶
% ============================================================
\newcommand{\tkzHeatFlask}[1][]{%
  \tikzchemlabdefaults
  \pgfkeys{/tikz-chemlab/.cd, #1}%
  % pst-labo: 每个 \tkzHeatFlask 有独立 pspicture, 不跨实例裁剪
  % tikz-chemlab: 不设全局 clip, 让所有内容自由绘制
  % pst-labo: \tkzRoundFlask 中 glassType=0 (tube) 也画 \pst@roundFlask
  % tikz-chemlab: 强制 glassType=0 → 1 (roundFlask)
  \begin{scope}[xshift=\chemlab@xshift]
    % 不设 \useasboundingbox, 让 TikZ 自动检测完整内容 bbox
    % 物体间相对位置由 shift 控制, 与 pst-labo 一致
    \ifTkzBunsenBurner\def\cnth{-5}\else\def\cnth{1}\fi
    % 气体回收装置
    \ifTkzGasCollection
      \ifTkzTestTubeCurvedTube
        \begin{scope}[shift={(8,-5)}]
          \tkzCrystallizingDish\relax\tkzTestTubeRack
          \begin{scope}[shift={(0,0.7)}]\tkzInvertedTube\relax\tkzBubbles\end{scope}
        \end{scope}
      \else\ifTkzTestTubeCurvedTubeShort
        \TkzBunsenBurnerfalse
        \begin{scope}[shift={(6,1)}]
          \tkzCrystallizingDish\relax\tkzTestTubeRack
          \begin{scope}[shift={(0,0.7)}]\tkzInvertedTube\relax\tkzBubbles\end{scope}
        \end{scope}
      \else\ifTkzTestTubeDoubleTube
        \TkzBunsenBurnerfalse
        \begin{scope}[shift={(6.2,1)}]
          \tkzCrystallizingDish\relax\tkzTestTubeRack
          \begin{scope}[shift={(0,0.7)}]\tkzInvertedTube\relax\tkzBubbles\end{scope}
        \end{scope}
      \fi\fi\fi
    \fi
    % 烧瓶 + 液体 + 导管
    % pst-labo: \rput(1,3){\tkzRoundFlask} 默认使用 bbox 中心对齐 (ref=c)
    % \tkzRoundFlask 的 pspicture bbox = (-2,0)(2,4), 中心 (0,2)
    % 所以局部原点 (瓶底) 实际位于世界 (1, 3-2) = (1, 1)
    \begin{scope}[shift={(1,1)}]
      % 液体 (clip 到烧瓶内部 ∩ 液体矩形, 匹配 pst-labo \psclip 双重裁剪)
      \begin{scope}
        \tkzGlassClipOptionsSet
        \clip[rounded corners=\tkzGlassClipOptions] \tkzGlassClipPath;
        % 第二次裁剪: 液体矩形 (匹配 pst-labo \psframe 裁剪)
        % 只在有液体时才裁剪; 无液体时不裁剪 (物质可见到玻璃顶部)
        \ifnum\tkzLiquidLevelI>0
          \pgfmathsetmacro{\chemlab@liquidClipY}{3 * \tkzLiquidLevelI / 100 * cos(\tkzTiltAngle)}%
          \clip (-4,-1.5) rectangle (5,\chemlab@liquidClipY);
        \fi
        \ifnum\tkzLiquidLevelI>0
          \tkzTestTubeLiquid{\tkzLiquidLevelI}{\tkzLiquidStyleI}{\tkzTiltAngle}
        \fi
        \tkzSubstance
        \tkzSolid
      \end{scope}
      % 玻璃器皿轮廓 (根据 glassType 选择: roundFlask/erlenmeyer/beaker 等)
      \tkzGlassDraw
      % 导管 (在瓶口 y=3.5 处) — 匹配 pst-labo \tkzRoundFlask 的全部导管类型
      \ifTkzTestTubeStraightTube
        \begin{scope}[shift={(0,3.5)}]\tkzTestTubeStraightTube\end{scope}
      \else\ifTkzTestTubeElbowTube
        \begin{scope}[shift={(0,3.5)}]\tkzTestTubeElbowTube\end{scope}
      \else\ifTkzTestTubeUTube
        \begin{scope}[shift={(0,3.5)}]\tkzTestTubeUTube\end{scope}
      \else\ifTkzTestTubeUTubeLong
        \begin{scope}[shift={(0,3.5)}]\tkzTestTubeUTubeLong\end{scope}
      \else\ifTkzTestTubeCurvedTube
        \begin{scope}[shift={(0,3.5)}]\tkzCurvedTube\end{scope}
      \else\ifTkzTestTubeCurvedTubeShort
        \begin{scope}[shift={(0,3.5)}]\tkzCurvedTubeShort\end{scope}
      \else\ifTkzTestTubeStopper
        \begin{scope}[shift={(0,3.5)}]\tkzTestTubeStopper\end{scope}
      \else\ifTkzTestTubeDoubleTube
        \begin{scope}[shift={(0,3.5)}]\tkzTestTubeDoubleTube\end{scope}
      \fi\fi\fi\fi\fi\fi\fi\fi
    \end{scope}
    \ifTkzBubbling
      % pst-labo: \rput(5.3,-3){\tkzTestTube[...]}, tiltAngle=0 默认, glassType=tube
      % bbox (-1.5,0)(1.5,4) center (0,2) → 局部原点 (试管底) 世界 (5.3, -5)
      \begin{scope}[shift={(5.3,-5)}]
        \tkzTestTube[liquidLevel1=70,glassType=tube]
      \end{scope}
    \fi
    \ifTkzBunsenBurner
      % pst-labo: \rput(1,-5){\pst@BecBunsenGrille} (无 pspicture, 实测原点对齐)
      \begin{scope}[shift={(1,-5)}]\tkzBunsenBurnerGrill\end{scope}
    \fi
  \end{scope}
  % 自动推进原点 (匹配 pst-labo \pspicture 的布局宽度)
  % 使用全局 dimen 累加, scope 的 xshift 会正确应用
  % pst-labo bbox 宽度: curvedTubeShort/DoubleTube = 11, 其他 = 6
  \ifTkzTestTubeUTubeLong
    % uTubeLong: 不推进
  \else\ifTkzTestTubeUTube
    % uTube: 不推进
  \else\ifTkzTestTubeCurvedTubeShort
    \advance\chemlab@xshift by 11cm %  bbox (-2,0)(9,8), width = 11
  \else\ifTkzTestTubeDoubleTube
    \advance\chemlab@xshift by 11cm %  bbox (-2,0)(9,8), width = 11
  \else
    \advance\chemlab@xshift by 6cm %  bbox (-2,\cnth)(4,6), width = 6
  \fi\fi\fi\fi
}

% ============================================================
% 10. \tkzRoundFlask  烧瓶 (高级命令)
% ============================================================
\newcommand{\tkzRoundFlask}[1][]{%
  \tikzchemlabdefaults
  \pgfkeys{/tikz-chemlab/.cd, #1}%
  % pst-labo: \tkzRoundFlask 中 glassType=0 (tube) 也画 \pst@roundFlask (case 0 = roundFlask)
  % tikz-chemlab: 强制 glassType=0 → 1 (roundFlask), 避免画试管轮廓
  \begin{scope}[xshift=\chemlab@xshift]
    % 不设 \useasboundingbox, 让 TikZ 自动检测完整内容 bbox
    % bbox 坐标使用局部坐标系, 会随 scale 自动缩放
    % 物体间相对位置由 shift 控制, 与 pst-labo 一致
    % 液体
    \begin{scope}[rotate=\tkzTiltAngle]
      \begin{scope}
        \tkzGlassClipOptionsSet
        \clip[rounded corners=\tkzGlassClipOptions] \tkzGlassClipPath;
        % 第二次裁剪: 液体矩形 (匹配 pst-labo \psframe 裁剪)
        % 只在有液体时才裁剪; 无液体时不裁剪 (物质可见到玻璃顶部)
        \ifnum\tkzLiquidLevelI>0
          \pgfmathsetmacro{\chemlab@liquidClipY}{3 * \tkzLiquidLevelI / 100 * cos(\tkzTiltAngle)}%
          \clip (-4,-1.5) rectangle (5,\chemlab@liquidClipY);
        \fi
        \ifnum\tkzLiquidLevelI>0
          \tkzTestTubeLiquid{\tkzLiquidLevelI}{\tkzLiquidStyleI}{\tkzTiltAngle}
        \fi
        \ifnum\tkzLiquidLevelII>0
          \tkzTestTubeLiquid{\tkzLiquidLevelII}{\tkzLiquidStyleII}{\tkzTiltAngle}
        \fi
        \ifnum\tkzLiquidLevelIII>0
          \tkzTestTubeLiquid{\tkzLiquidLevelIII}{\tkzLiquidStyleIII}{\tkzTiltAngle}
        \fi
        \tkzSubstance
        \tkzSolid
      \end{scope}
      % 标签
      \ifTkzLabel
        \node[fill=VeryLightGray,draw=none,inner sep=2pt] at (0,1.5) {\tkzNumber};
      \fi
      % 玻璃器皿轮廓
      \tkzGlassDraw
      % 顶部附件 (按优先级)
      \ifTkzTestTubeStraightTube
        \begin{scope}[shift={(0,3.5)}]\tkzTestTubeStraightTube\end{scope}
      \else\ifTkzTestTubeElbowTube
        \begin{scope}[shift={(0,3.5)}]\tkzTestTubeElbowTube\end{scope}
      \else\ifTkzTestTubeUTube
        \begin{scope}[shift={(0,3.5)}]\tkzTestTubeUTube\end{scope}
      \else\ifTkzTestTubeUTubeLong
        \begin{scope}[shift={(0,3.5)}]\tkzTestTubeUTubeLong\end{scope}
      \else\ifTkzTestTubeCurvedTube
        \begin{scope}[shift={(0,3.5)}]\tkzCurvedTube\end{scope}
      \else\ifTkzTestTubeCurvedTubeShort
        \begin{scope}[shift={(0,3.5)}]\tkzCurvedTubeShort\end{scope}
      \else\ifTkzTestTubeStopper
        \begin{scope}[shift={(0,3.5)}]\tkzTestTubeStopper\end{scope}
      \else\ifTkzTestTubeDoubleTube
        \begin{scope}[shift={(0,3.5)}]\tkzTestTubeDoubleTube\end{scope}
      \else\ifTkzTestTubeCondenserBulbs
        \tkzRoundFlaskReflux
        \begin{scope}[shift={(0,3.5)}]\tkzCondenserBulbs\end{scope}
        \tkzHeaterRoundFlask
      \fi\fi\fi\fi\fi\fi\fi\fi\fi
    \end{scope}
  \end{scope}
  % 自动推进原点 (匹配 pst-labo \pspicture 的布局宽度)
  % pst-labo bbox: condenserBulbs = (-2,0)(2,12), 其他 = (-2,0)(2,4), width = 4
  \advance\chemlab@xshift by 4cm %  bbox width = 4
}

% ============================================================
% 11. \tkzChampagneBubbles  香槟气泡
% ============================================================
% 原版 pst-labo 源码:
%   \setrandim{\BubbleX}{-1pt}{1pt}    -- BulleX ∈ [-1, 1] (纯数字, cm)
%   \setrandim{\BubbleY}{0pt}{2pt}     -- BulleY ∈ [0, 2] (纯数字)
%   Y 坐标 = \pointless\BubbleY * liquidLevel1/100  (乘以液位百分比)
%   半径 = {\BubbleY} (带 pt 单位, 1pt = 0.0353 cm)
%   \pst@dimA 分级放大: BulleY<1pt → ×3, 1pt≤BulleY<2pt → ×2
%
% 实测原版渲染图:
%   气泡为黑色边空心圆 (\pscircle, fillstyle=none)
%   半径 0.05~0.175 单位 (比源码计算 0~0.071 大, 可能是渲染缩放)
%   气泡团中心 (世界): (7.89, -0.975)
\newcommand{\tkzChampagneBubbles}[1][25]{%
  \foreach \iBubble in {1,...,#1}{%
    % BulleX ∈ [-1, 1] (纯数字, cm)
    \pgfmathsetmacro{\BubbleX}{-1 + 2*rnd}%
    % BulleY ∈ [0, 2] (纯数字)
    \pgfmathsetmacro{\BubbleY}{2*rnd}%
    % Y 坐标 = BulleY * liquidLevel1/100 (匹配源码逻辑)
    \pgfmathsetmacro{\BubbleYpos}{\BubbleY * \tkzLiquidLevelI / 100}%
    % dimA 分级放大 (匹配源码逻辑)
    \pgfmathsetmacro{\BubbleRand}{rnd}%
    \ifdim\BubbleRand pt<0.5pt
      % BulleY < 1pt (50%): dimA = BulleY * 3
      \pgfmathsetmacro{\BubbleR}{\BubbleY * 0.0875}%
    \else
      % 1pt <= BulleY < 2pt (50%): dimA = BulleY * 2
      \pgfmathsetmacro{\BubbleR}{\BubbleY * 0.0875}%
    \fi
    \begin{scope}[rotate=\tkzTiltAngle]
      \draw[draw=black, line width=0.6pt] (\BubbleX,\BubbleYpos) circle (\BubbleR);
    \end{scope}
  }%
}

% ============================================================
% 12. \tkzBubbles  倒置管中气泡
% ============================================================
\newcommand{\tkzBubbles}[1][25]{%
  \pgfgettransformentries{\clab@aa}{\clab@bb}{\clab@cc}{\clab@dd}{\clab@xx}{\clab@yy}%
  \pgfmathsetmacro{\clab@scale}{(\clab@aa + \clab@dd) / 2}%
  \foreach \iBubble in {1,...,#1}{%
    % pst-labo: \setrandim{BulleX}{-0.1pt}{0.1pt} → \pointless 去掉pt, 数值直接用作局部坐标
    \pgfmathsetmacro{\BubbleX}{-0.1 + 0.2*rnd}%
    \pgfmathsetmacro{\BubbleY}{0.15 + 2.85*rnd}%
    % pst-labo: dimh = BulleY * (3 if <1pt, 2 if <2pt, 1 if >=2pt)
    % dotscale = dimh, dot radius = dotscale * dotsize/2
    % dotsize default = 2pt + 2*linewidth = 2pt + 2*0.8pt = 3.6pt → radius = dimh * 1.8pt
    % 转换为局部坐标: local = pt / (72 * scale)
    \pgfmathparse{\BubbleY<1?3:(\BubbleY<2?2:1)}%
    \let\BubbleScale\pgfmathresult
    \pgfmathsetmacro{\BubbleDimh}{\BubbleY * \BubbleScale}%
    \pgfmathsetmacro{\BubbleR}{1.8 * \BubbleDimh / (72 * \clab@scale)}%
    % Y position = BulleY * 1.5 (局部坐标, 匹配 pst-labo \pointless\BubbleY 1.5 mul)
    \pgfmathsetmacro{\BubbleYpos}{\BubbleY * 1.5}%
    \begin{scope}[rotate=\tkzTiltAngle]
      % pst-labo: dotstyle=o (open circle = hollow, black outline)
      \draw[draw=black, line width=0.4pt, fill=white] (\BubbleX,\BubbleYpos) circle (\BubbleR);
    \end{scope}
  }%
}

% ============================================================
% 13. \tkzFilaments  絮状物
% ============================================================
\newcommand{\tkzFilaments}[2][5]{%
  \pgfgettransformentries{\fil@aa}{\fil@bb}{\fil@cc}{\fil@dd}{\fil@xx}{\fil@yy}%
  \pgfmathsetmacro{\fil@scale}{(\fil@aa + \fil@dd) / 2}%
  \foreach \iFil in {1,...,#1}{%
    % pst-labo: \setrandim{\pst@FilamentWidth}{0pt}{0.05pt} → \pointless = [0, 0.05)
    % linewidth=\pointless\pst@FilamentWidth 解释为 N × unit (PSTricks 规则)
    % 实际线宽 = [0, 0.05) × scale cm (scale 对应 unit)
    \pgfmathsetmacro{\fw}{0.05 * \fil@scale * rnd}%
    \pgfmathsetmacro{\myRand}{rnd}%  ONE random value per filament (匹配 pst-labo MyRand)
    \begin{scope}[rotate=\tkzTiltAngle]
      % pst-labo \pscurve 通过4个点的平滑曲线 (不是贝塞尔):
      %   P1: (MyRand-1, 2*MyRand-1)     = ([-1,0), [-1,1))
      %   P2: (MyRand+0.2, MyRand+2)     = ([0.2,1.2), [2,3))
      %   P3: (MyRand, MyRand-1)         = ([0,1), [-1,0))
      %   P4: (MyRand+0.5, MyRand+4)     = ([0.5,1.5), [4,5))
      \draw[line width=\fw cm,#2] plot[smooth,tension=0.5] coordinates {
        ({\myRand-1},{2*\myRand-1})
        ({\myRand+0.2},{\myRand+2})
        ({\myRand},{\myRand-1})
        ({\myRand+0.5},{\myRand+4})
      };
    \end{scope}
  }%
}

% ============================================================
% 14. \tkzIronNail  铁钉
% ============================================================
% 单个铁钉形状 (单位 0.5)
\newcommand{\tkzNails}{%
  % pst-labo 实际 SLW=0.71135pt=0.025cm (从 xdv 提取)
  % \psarc[unit=0.5,linewidth=0.1] 的 \psline 继承同一 linewidth
  \draw[scale=0.5,line width=0.025cm,blue]
    ([shift={(0,0.5)}]190:0.4) arc[start angle=190,end angle=260,radius=0.4];
  \draw[scale=0.5,line width=0.025cm,blue] (-0.28,0.2) -- (0.5,1);
}

\newcommand{\tkzIronNail}[1][60]{%
  \pgfgettransformentries{\clab@aa}{\clab@bb}{\clab@cc}{\clab@dd}{\clab@xx}{\clab@yy}%
  \pgfmathsetmacro{\clab@scale}{(\clab@aa + \clab@dd) / 2}%
  \foreach \iNail in {1,...,#1}{%
    % pst-labo: \setrandim{ClouX}{-1cm}{1cm} → 物理坐标, 需除以 scale 转换为局部坐标
    \pgfmathsetmacro{\NailX}{(-1 + 2*rnd) / \clab@scale}%
    \pgfmathsetmacro{\NailY}{(-0.2 + 0.6*rnd) / \clab@scale}%
    \pgfmathsetmacro{\NailAngle}{-90 + 180*rnd}%
    \begin{scope}[rotate=\tkzTiltAngle]
      \begin{scope}[shift={(\NailX,\NailY)},rotate=\NailAngle]
        \tkzNails
      \end{scope}
    \end{scope}
  }%
}

% ============================================================
% 15. \tkzCopperTurnings  铜屑
% ============================================================
% 单个铜屑形状 (单位 0.5)
\newcommand{\tkzCopper}{%
  % pst-labo 实际 SLW=0.28458pt=0.01cm (从 xdv 提取)
  \draw[scale=0.5,line width=0.01cm,Copper]
    (0.25,0.25) .. controls (0,0) and (-0.25,0.25) .. (0.25,0.75)
    .. controls (0,0.5) and (-0.25,0.75) .. (0.25,1.25)
    .. controls (0,1) and (-0.25,1.25) .. (0.25,1.75);
}

\newcommand{\tkzCopperTurnings}[1][30]{%
  \pgfgettransformentries{\clab@aa}{\clab@bb}{\clab@cc}{\clab@dd}{\clab@xx}{\clab@yy}%
  \pgfmathsetmacro{\clab@scale}{(\clab@aa + \clab@dd) / 2}%
  \foreach \iT in {1,...,#1}{%
    \pgfmathsetmacro{\TX}{(-1 + 2*rnd) / \clab@scale}%
    \pgfmathsetmacro{\TY}{(-0.2 + 0.5*rnd) / \clab@scale}%
    \pgfmathsetmacro{\TAngle}{-180 + 360*rnd}%
    \begin{scope}[rotate=\tkzTiltAngle]
      \begin{scope}[shift={(\TX,\TY)},rotate=\TAngle]
        \tkzCopper
      \end{scope}
    \end{scope}
  }%
}

% ============================================================
% 16. \tkzZincGranules  锌粒
% ============================================================
% 单个锌粒形状 (单位 0.5)
\newcommand{\tkzZinc}{%
  \filldraw[scale=0.5,fill=LightGray,draw=black]
    plot[smooth,tension=0.5] coordinates {(0,0)(-0.25,0.25)(0,0.5)
      (0.25,0.75)(0.5,0.5)(0.15,0.4)(0,0)} -- cycle;
}

\newcommand{\tkzZincGranules}[1][25]{%
  \pgfgettransformentries{\clab@aa}{\clab@bb}{\clab@cc}{\clab@dd}{\clab@xx}{\clab@yy}%
  \pgfmathsetmacro{\clab@scale}{(\clab@aa + \clab@dd) / 2}%
  \foreach \iG in {1,...,#1}{%
    \pgfmathsetmacro{\GX}{(-1 + 2*rnd) / \clab@scale}%
    \pgfmathsetmacro{\GY}{(-0.2 + 0.5*rnd) / \clab@scale}%
    \pgfmathsetmacro{\GAngle}{-180 + 360*rnd}%
    \begin{scope}[rotate=\tkzTiltAngle]
      \begin{scope}[shift={(\GX,\GY)},rotate=\GAngle]
        \tkzZinc
      \end{scope}
    \end{scope}
  }%
}

% ============================================================
% 17. \tkzBalls  3D 球状物
% ============================================================
\newcommand{\tkzBalls}[1][50]{%
  \foreach \iB in {1,...,#1}{%
    \pgfmathsetmacro{\BX}{-1.5 + 3*rnd}%
    \pgfmathsetmacro{\BY}{3*rnd}%
    \pgfmathsetmacro{\BR}{0.01 + 0.1*rnd}%
    % 按 Y 位置调整半径
    \pgfmathparse{\BY>2?3:(\BY>1?2:1)}%
    \let\Bscale\pgfmathresult
    \pgfmathsetmacro{\BRscaled}{\BR * \Bscale}%
    \begin{scope}[rotate=\tkzTiltAngle]
      % pst-labo Ball3D: GradientCircle=true
      % gradbegin=CMYK(0,0.1,0.95,0.1) → 实际渲染为 RGB(0,0.1,0.95)=蓝色 (固定颜色)
      % gradend=white
      % 径向渐变: 中心=white, 边缘=蓝色
      \definecolor{billeBlue}{rgb}{0,0.1,0.95}%
      \shade[inner color=white, outer color=billeBlue] (\BX,\BY) circle (\BRscaled);
    \end{scope}
  }%
}

% ============================================================
% 18. \tkzBubblesThreeD  2D 气泡
% ============================================================
\newcommand{\tkzBubblesThreeD}[2][50]{%
  \foreach \iB in {1,...,#1}{%
    \pgfmathsetmacro{\BX}{-1.5 + 3*rnd}%
    \pgfmathsetmacro{\BY}{3*rnd}%
    \pgfmathsetmacro{\BR}{0.01 + 0.1*rnd}%
    \pgfmathparse{\BY>2?3:(\BY>1?2:1)}%
    \let\Bscale\pgfmathresult
    \pgfmathsetmacro{\BRscaled}{\BR * \Bscale}%
    \begin{scope}[rotate=\tkzTiltAngle]
      \fill[#2,line width=0.2pt] (\BX,\BY) circle (\BRscaled);
    \end{scope}
  }%
}

\endinput
